\documentclass[12pt]{article}
\usepackage[leqno]{amsmath}
\usepackage{amsfonts}
\usepackage{float}
\usepackage{bm}
\usepackage{cite}
\usepackage{enumerate}
\usepackage{titlesec}
\usepackage{float,subfigure}
\usepackage{graphicx}
\usepackage{geometry}
\usepackage{color}
\usepackage{yhmath}
\usepackage{listings}
\usepackage{xcolor}
\usepackage[thmmarks,amsmath]{ntheorem}
\usepackage{mathrsfs,amssymb}
\usepackage{multirow}
\usepackage{multicol}
\usepackage{bbm}
%\usepackage[toc]{multitoc}
\usepackage{fancyhdr}
\usepackage[colorlinks,
            linkcolor=blue,
            anchorcolor=blue,
            citecolor=red
            ]{hyperref}
\usepackage{indentfirst}
%\usepackage{newtxtext}
\usepackage{extarrows}
\usepackage[nottoc]{tocbibind}
\setlength{\topskip}{0.5cm}
\setlength{\parindent}{2em}
%\usepackage{mathtime}
%\usepackage{palatino}
%\usepackage[framed,numbered,autolinebreaks,useliterate]{mcode}
\geometry{left=1in,right=1in,top=0.7in,bottom=1in}

\titleformat{\section}{\centering\large\bfseries}{\S\,\thesection\,}{1em}{}

\usepackage{ebgaramond} % Use ebgaramond Font
\pdfmapfile{+mtpro2.map} 
\usepackage[lite]{mtpro2} % Load New times Roman Pro Math font

%\usepackage{xeCJK}
%\newcommand{\fontsong}{SimSun}
%\newcommand{\fonthei}{Adobe Heiti Std}
%\newcommand{\fontkai}{FangSong}
%\setCJKmainfont[BoldFont={\fonthei}, ItalicFont={\fontkai}]{\fontsong}
\fancyhead[R]{\rightmark}
\pagestyle{fancy}
\lhead{\vspace{0pt}\includegraphics[height=3.45ex]{tamu_H.pdf}}
%\headsep=5pt

\DeclareMathAlphabet{\mathcmr}{OT1}{cmr}{m}{n}
\SetMathAlphabet{\mathcmr}{bold}{OT1}{cmr}{bx}{n}

% 你要的那个粗体 P
\renewcommand{\P}{\bm{\mathcmr{P}}}
\newcommand{\EE}{\bm{\mathcmr{E}}}

%\def\CJK@null{\kern\CJKnullspace\hskip \f@size \p@\kern\CJKnullspace}
%\defaultfontfeatures{Mapping=tex-text} % use TeX --
%\newcommand{\fontsong}{SimSong}
%\newcommand{\fonthei}{Adobe Heiti Std}
%\newcommand{\fontkai}{Adobe Kaiti Std}
%\newcommand{\fontfs}{Adobe Fangsong Std}
%\newcommand{\fontli}{Adobe Kaiti Std}
%\newcommand{\fontyou}{Adobe Kaiti Std}
%\newcommand{\tabincell}[2]{\begin{tabular}{@{}#1@{}}#2\end{tabular}}
%
%\setCJKmainfont[BoldFont={\fonthei}, ItalicFont={\fontkai}]{\fontsong}
%\setCJKsansfont{\fonthei}
%\setCJKmonofont{\fontkai}
%\setCJKfamilyfont{song}{\fontsong}
%\setCJKfamilyfont{hei}{\fonthei}
%\setCJKfamilyfont{fs}{\fontfs}
%\setCJKfamilyfont{kai}{\fontkai}
%\setCJKfamilyfont{li}{\fontli}
%\setCJKfamilyfont{you}{\fontyou}

\renewcommand{\b}{\mathscr{B}}
\renewcommand{\i}{\bm{1}}
\renewcommand{\sb}{\subset}
\renewcommand{\(}{\left(}%定义随公式大小变动的括号
\renewcommand{\)}{\right)}%定义随公式大小变动的括号
\renewcommand{\[}{\left[}%定义随公式大小变动的括号
\renewcommand{\]}{\right]}%定义随公式大小变动的括号
\renewcommand{\r}{\mathbb{R}}
\renewcommand{\ae}{\mathrm{a.s.}}
\renewcommand{\d}{\mathrm{d}}
\renewcommand{\c}{\mathbb{C}}
\newcommand{\cc}{\mathfrak{C}}
\renewcommand{\a}{\mathscr{A}}
\renewcommand{\l}{\mathscr{L}}
\newcommand{\ct}{\mathcal{T}}
\newcommand{\cp}{\mathcal{P}}
\newcommand{\z}{\mathbb{Z}}
\newcommand{\p}{\mathbb{P}}
\newcommand{\q}{\mathbb{Q}}
\newcommand{\n}{\mathbb{N}}
\newcommand{\x}{\mathscr{X}}
\newcommand{\y}{\mathscr{Y}}
\newcommand{\f}{\mathscr{F}}
\newcommand{\st}{\mathrm{\ s.t.\ }}
\newcommand{\e}{\mathscr{E}}
\newcommand{\ri}{\rightarrow}
\newcommand{\li}{\lim\limits_}
\newcommand{\var}{{\rm Var}}
\newcommand{\pa}{\partial}
\newcommand{\g}{\mathscr{G}}
\newcommand{\pe}{\mathscr{P}}

\def\Xint#1{\mathchoice	{\XXint\displaystyle\textstyle{#1}}	{\XXint\textstyle\scriptstyle{#1}}{\XXint\scriptstyle\scriptscriptstyle{#1}}{\XXint\scriptscriptstyle\scriptscriptstyle{#1}}	\!\int}\def\XXint#1#2#3{{\setbox0=\hbox{$#1{#2#3}{\int}$}	\vcenter{\hbox{$#2#3$}}\kern-.5\wd0}}\def\ddashint{\Xint=}\def\dashint{\Xint-}%\dashint

\newcommand\theref[1]{Theorem~\ref{#1}}
\newcommand\corref[1]{Proposition~\ref{#1}}
\newcommand\lemref[1]{Lemma~\ref{#1}}

\setcounter{MaxMatrixCols}{14}
\allowdisplaybreaks[4]

{\theoremstyle{plain}
\theoremheaderfont{\upshape\bfseries}
\theorembodyfont{\itshape}
\newtheorem{lemma}{Lemma}[section]}

{\theoremstyle{plain}
\theoremheaderfont{\normalfont \bfseries}
\theorembodyfont{\itshape}
\newtheorem{question}{QUESTION}[section]}

{\theoremstyle{plain}
\theoremheaderfont{\upshape\bfseries}
\theorembodyfont{\normalfont}
\newtheorem{definition}{Definition}[section]}

{\theoremstyle{plain}
\theoremheaderfont{\bfseries}
\theorembodyfont{\normalfont}
\newtheorem{example}{Example}[section]}

{\theoremstyle{plain}
\theoremheaderfont{\bfseries}
\theorembodyfont{\normalfont}
\newtheorem{axiom}{Axiom}}

{\theoremstyle{plain}
\theoremheaderfont{\upshape\bfseries}
\theorembodyfont{\itshape}
\newtheorem{theorem}{Theorem}[section]}

{\theoremstyle{plain}
\theoremheaderfont{\upshape\bfseries}
\theorembodyfont{\itshape}
\newtheorem{proposition}{Proposition}[section]}

{\theoremstyle{plain}
\theoremheaderfont{\bfseries}
\theorembodyfont{\normalfont}
\newtheorem{remark}{Remark}[section]}

{\theoremstyle{plain}
\theoremheaderfont{\upshape\bfseries}
\theorembodyfont{\itshape}
\newtheorem{corollary}{Corollary}[section]}

{\theoremstyle{plain}
\theoremheaderfont{\bfseries}
\theorembodyfont{\normalfont}
\newtheorem*{ana}{Analysis}}

{\theoremstyle{nonumberplain}
\theoremheaderfont{\itshape}
\theorembodyfont{\normalfont}
\theoremsymbol{$\square$}
\newtheorem{proof}{Proof.}}

\title{
\includegraphics[scale=0.35]{tamu.pdf}\protect\\
{\itshape Selections} of {\bfseries Inspirations}
}

\author{Han-bai Wang\quad Typeset by \LaTeX}
\date{\sffamily 2022 Spring}

\begin{document}
  \maketitle
  \pdfbookmark[1]{Con}{toc}
  \begin{abstract}
  	The significance of probability theory is to describe the regularity shown by a large number of random factors.
  	
  	\hfill --- Kolmogorov
  \end{abstract}
  \tableofcontents
  \newpage
  \section{Set Theory}
  \begin{question}
   Suppose $\{A_i\}_{i\in I}$ and $\{B_j\}_{j\in I}$ are two set sequences. Then
   \begin{equation*}
    \begin{split}
     \(\bigcap_{i\in I} A_i\) \cap \(\bigcap_{j\in I} B_j\)\subset \bigcap_{i\in J}(A_i\cap B_i),\\
     \(\bigcup_{i\in I} A_i\) \cup \(\bigcup_{j\in I} B_j\)\subset \bigcup_{i\in J}(A_i\cup B_i).
    \end{split}
   \end{equation*}
  \end{question}
  \begin{proof}
   We start from the first inequality.
   \begin{equation*}
    \(\bigcap_{i\in I} A_i\) \cap \(\bigcap_{j\in I} B_j\)=\bigcap_{i\in I} \(A_i \cap \bigcap_{j\in I} B_j\)\subset \bigcap_{i\in J}(A_i\cap B_i).
   \end{equation*}
   The second inequality follows from
   \begin{equation*}
    \(\bigcup_{i\in I} A_i\) \cup \(\bigcup_{j\in I} B_j\)=\bigcup_{i\in I} \(A_i \cap \bigcap_{j\in I} B_j\)
   \end{equation*}
  \end{proof}

  \begin{question}
   Suppose $\{A_n\}$ is a set sequence. Then
   \begin{equation*}
    \(\bigcup_{n=1}^m A_n\)'=\bigcup_{n=1}^m A_n',\quad \overline{\bigcup_{n=1}^m A_n}=\bigcup_{n=1}^m \overline{ A_n}.
   \end{equation*}
  \end{question}

  \begin{question}
   Suppose $A$ is a countable set, then $\a:=\{F\subset A: F\text{ is finite}  \}$ is countable.
  \end{question}
  \begin{proof}
   Wirte $A=\{a_m\}$. We use $|F|$ to denote the cardinal number of $F$. Then we have the following decomposition:
   \begin{equation*}
    \a=\sum_{n=0}^\infty \{F\in\a:|F|=n\}=:\sum_{n=0}^\infty \a_n.
   \end{equation*}
   So we only need to verify $\a_n$ is at most countable for each $n\geq 0$. However, we know for each $n\geq 0$,
   \begin{equation*}
    \varphi_n: \a_n\to \n^n:\{a_{m_k}\}_{k=1}^n\to (m_1,\cdots,m_k ).
   \end{equation*}
   It is easy to verify $\varphi_n$ is a one-to-one mapping. So $\a_n$ is countable and thus $\a$ is countable.
  \end{proof}

  \section{Analysis}
  \begin{question}
   Prove the following inequality:
   \begin{equation*}
    |x\wedge C-y\wedge C|\leq |x-y|,\quad |x\vee C-y\vee C|\leq |x-y|,
   \end{equation*}
   where $x,y,C\in\r$.
  \end{question}
   \begin{proof}
    \begin{enumerate}
     \item \begin{enumerate}
     \item When $x,y\leq  C$, then $|x\wedge C-y\wedge C|= |x-y|$.
     \item When $x,y\geq C$, then $|x\wedge C-y\wedge C|=|C-C|=0\leq |x-y|$.
     \item Without loss of generality, we assume $x\leq C<y$, then $|x\wedge C-y\wedge C|=|x-C|\leq |x-y|$.
    \end{enumerate}
    In conclusion, the first inequality holds.
     \item \begin{enumerate}
     \item When $x,y\leq  C$, then $|x\vee C-y\vee C|= |C-C|=0\leq |x-y|$.
     \item When $x,y\geq C$, then $|x\vee C-y\vee C|= |x-y|$.
     \item Without loss of generality, we assume $x\leq C<y$, then $|x\vee C-y\vee C|=|C-y|\leq |x-y|$.
    \end{enumerate}
    In conclusion, the second inequality holds.
    \end{enumerate}
   \end{proof}
   
   
   \begin{question}
	For any nonnegative (nonpositive) measurable function $f(x)\in\b[a,\infty)$ and $t \geq a$, we have
	\begin{equation*}
		\int_a^{b\wedge t}f(x)dx\leq (\geq)\int_a^b f(x\wedge t )dt. 
	\end{equation*}
  \end{question}
  \begin{proof}
  	If $a\leq b\leq t$, then the above two sides are actually equal. If $b>t$, then 
  	\begin{equation*}
  		\begin{split}
  			\int_a^b f(x\wedge t )dt&=\int_a^t f(x)dt+f(t)(t-b)\\
  			&=\int_a^{b\wedge t} f(x)dt+f(t)(t-b)\\
  			&\geq (\leq )\int_a^{b\wedge t}f(x)dx.
  		\end{split}
  	\end{equation*}
  	So the desired results follows.
  \end{proof}
  
  \begin{question}
   Suppose $f,g:U\to\r$ and $U\subset\r^n$ is an open set. Assume exists $C>0$ such that $f(x)+g(x)\leq C$ for each $x\in U$, then
   \begin{equation*}
    \sup_{x\in U}f(x)+\sup_{x\in U}g(x)\leq 2C.
   \end{equation*}
   \begin{proof}
    The inequality $f(x)+g(x)\leq C,\forall x\in U$ yields
    \begin{equation*}
     f(x)\leq C,\quad g(x)\leq C,\quad\forall x\in U.
    \end{equation*}
    So $\sup_{x\in U}f(x)\leq C$ and $\sup_{x\in U}g(x)\leq C$. This is what we desire.
   \end{proof}
  \end{question}
  \begin{question}
   Suppose $f:\r\to\r$ is generalized intergral, i.e.,
   \begin{equation*}
    \int_1^\infty f(x)dx=\li{\ell\to\infty}\int_1^\ell f(x)dx\in\r.
   \end{equation*}
    Use an example to show $f(x)\not\to0,x\to\infty$ might happen.
  \end{question}
  \begin{proof}
   Consider
   \begin{equation*}
    f(x)=\frac{\sin x^2}{x},\quad f'(x)=2\cos x^2- \frac{\sin x^2}{x^2}.
   \end{equation*}
   Notice $f'\in C^\infty[1,\infty)$, so $\forall\ell> 1, f'\in R[1,\ell]$. And
   \begin{equation*}
    \begin{split}
     \int_1^\infty f'(x)dx=&\li{\ell\to\infty}\int_1^\ell f'(x)dx=\li{\ell\to\infty}\[f(\ell)-f(1) \]\\
     =&\li{\ell\to\infty}\(\frac{\sin\ell^2}{\ell}-\sin 1 \)=-\sin 1\in\r.
    \end{split}
   \end{equation*}
   However
   \begin{equation*}
    \li{x\to\infty}f'(x)=\li{x\to\infty}\(2\cos x^2- \frac{\sin x^2}{x^2}\)
   \end{equation*}
   does not exist.
  \end{proof}

  \begin{question}
   Suppose $f$ is differentiable on $[a,b]$ ($f$ has the left and right derivatives on endpoints). If $\forall\varepsilon>0,\exists\delta>0,$ when $0<|h|<\delta,$ for all $x\in[a,b]$, we have
   \begin{equation*}
    \bigg|\frac{f(x+h)-f(x)}{h}-f'(x)\bigg|<\varepsilon,
   \end{equation*}
   then we say $f$ is uniformly differentiable on $[a,b]$. Prove $f'\in C[a,b]$ iff $f$ is uniformly differentiable.
  \end{question}
  \begin{proof}
   For all $x\in[a,b]$, write
   \begin{equation*}
    f_h(x):=\frac{f(x+h)-f(x)}{h},
   \end{equation*}
   then $f_h\in C[a,b]$.
   \begin{enumerate}
    \item If $f'\in C[a,b]$, then $f'$ is uniformly continuous, then $\forall\varepsilon>0,\exists\delta>0,\forall \xi,x\in[a,b],|\xi-x |<\delta,$ then $|f'(\xi)-f'(x)|<\varepsilon$. For all $x\in[a,b]$, when $0<|h|<\delta$, by Lagrange theorem of mean, exist $\xi$ between $x$ and $x+h$,
    \begin{equation*}
     \bigg|\frac{f(x+h)-f(x)}{h}-f'(x)\bigg|=|f'(\xi )-f'(x)|<\varepsilon,
    \end{equation*}
    which means $f$ is uniformly differentiable.
    \item If $f$ is uniformly differentiable, then $\forall \varepsilon>0,\exists\delta_1>0$, when $|h|<\delta_1, \forall x\in[a,b]$, we have
    \begin{equation*}
     |f_h(x)-f'(x)|<\varepsilon/3.
    \end{equation*}
    Since $f_h$ is continuous, $\forall\varepsilon>0,\exists \delta>0$, when $|\Delta|<\delta$, we have
    \begin{equation*}
     |f_h(x+\Delta )-f_h(x )|<\varepsilon/3.
    \end{equation*}
    Finally, $\forall\varepsilon>0,\exists\delta>0,$ when $|\Delta |<\delta$, we obtain
    \begin{equation*}
     \begin{split}
      &|f'(x+\Delta )-f'(x)|\\
      \leq &|f'(x+\Delta )-f_h(x+\Delta )|+|f_h(x+\Delta )-f_h(x )|+|f_h(x)-f'(x)|\\
      <&\varepsilon/3+\varepsilon/3+\varepsilon/3=\varepsilon,
     \end{split}
    \end{equation*}
    which implies $f'\in C[a,b]$.
   \end{enumerate}
  \end{proof}

  \begin{question}
   Suppose $\sum_i a_i$ is a real series and $\sum_j A_j$ is a real series obtained by adding some brackets in some items of $\sum_i a_i$. If for every $A_j$, the addends of $A_j$ are all positive or all negative, then prove $\sum_i a_i$ converges iff $\sum_j A_j$ converges, and when they converge, we have $\sum_i a_i=\sum_j A_j$.
  \end{question}
  \begin{proof}
   \begin{enumerate}
    \item Notice $\sum_{j=1}^n A_j,n\geq 1$ is just a sub-sequence of $\sum_{i=1}^n a_i,n\geq 1$, if $\sum_i a_i$ converges, then $\sum_j A_j$ must converges and $\sum_j A_j=\sum_i a_i$.
    \item On the contrary, if $\sum_j A_j$ converges, then $\forall \varepsilon>0,\exists N_1\in\n,\forall n>N_1,p\in\n$, we have
    \begin{equation*}
     \bigg|\sum_{j=n}^{n+p}A_j\bigg|<\varepsilon.
    \end{equation*}
    We write $m_j$ to denote the number of items bracketed in $A_j,j\geq 1$ and write $a_i\in A_j$ to denote $a_j$ is an addend of $A_j$. Consider $N_2:=\sum_{j=1}^{N_1}m_j$,  then $\forall n>N_2,p\in\n,\exists \alpha,\beta \in\n,\st$ $a_n\in A_\alpha$ and $a_{n+p}\in A_\beta $. Write $n_\alpha :=\min\{n:a_n\in A_\alpha \}$ and $n_\beta:=\max\{n:a_n\in A_\beta\}$. Observe $\alpha>N_1$, then $\min\{|A_\alpha |,|A_\beta |,|\sum_{j=\alpha}^\beta A_j| \} <\varepsilon$. Whence
    \begin{equation*}
     \begin{split}
      \bigg|\sum_{i=n}^{n+p}a_i\bigg|&=\bigg|\sum_{j=\alpha}^\beta A_j-\sum_{i=n_\alpha}^{n-1}a_i-\sum_{i=n+p+1}^{n_\beta}a_i \bigg|\\
      &\leq \bigg|\sum_{j=\alpha}^\beta A_j\bigg|+\bigg|\sum_{i=n_\alpha}^{n-1}a_i\bigg|+\bigg|\sum_{i=n+p+1}^{n_\beta}a_i \bigg|\\
      &\leq \bigg|\sum_{j=\alpha}^\beta A_j\bigg|+ |A_\alpha|+|A_\beta|< 3\varepsilon,
     \end{split}
    \end{equation*}
    where the last $\leq $ comes from all the addends of $A_j$ are all positive or all negative. So this formula means $\sum_ia_i$ converges. Notice $\sum_{j=1}^n A_j,n\geq 1$ is just a sub-sequence of $\sum_{i=1}^n a_i,n\geq 1$, if $\sum_i a_i$ converges, then $\sum_j A_j$ must converges and $\sum_j A_j=\sum_i a_i$.
   \end{enumerate}
  \end{proof}

  \begin{question}
   Suppose $f_0:D\to \r$ is an increasing function, where $D$ is a dense subset on $\r$. We define
   \begin{equation*}
    f(x):=\inf_{x<t\in D}f_0(t),\quad \forall x\in \r.
   \end{equation*}
   Prove the uniformly continuity of $f_0$ implies the uniformly continuity of $f$.
  \end{question}
  \begin{proof}
   Since $f_0$ is uniformly continuous, then $\forall\varepsilon>0,\exists\delta>0$, for all $t_1,t_2\in D$ and $|t_1-t_2|<3\delta$, we have
   \begin{equation*}
    |f_0(t_1)-f_0(t_2)|<\varepsilon/2.
   \end{equation*}
    For all $x_1,x_2\in \r$ and $|x_1-x_2|<\delta$, we can find a sequence $\{t_1^{(n)} \}$ and $\{t_2^{(n)}\}$ in $D$, such that $t_j^{(n)}\searrow x_j,j=1,2$. Exist $N\in\n_+$, for all $n>N$, we have $|t_j^{(n)}-x_j|<\delta,j=1,2$. This implies $\forall n>N$,
    \begin{equation*}
     |t_1^{(n)}-t_2^{(n)}|\leq |t_1^{(n)}-x|+|x-y|+|y-t_2^{(n)}|<3\delta.
    \end{equation*}
    Hence for all $n>N$,
    \begin{equation*}
     |f_0(t_1^{(n)})-f_0(t_2^{(n)})|<\varepsilon/2.
    \end{equation*}
    Let $n\to\infty$ in the above formula and observe the fact that
    \begin{equation*}
     f(x)=\inf_{x<t\in D}f_0(t)=\lim_{n\to\infty}f_0(t_1^{(n)}),
    \end{equation*}
    then we will obtain
    \begin{equation*}
     |f(x_1)-f(x_2)|\leq  \varepsilon/2<\varepsilon .
    \end{equation*}
    This provides the desired result.
  \end{proof}

  \begin{question}
   Suppose $\text{\rmfamily \bfseries x}=\text{\rmfamily \bfseries f}(\theta_1,\cdots,\theta_{n-1} )$ is a regular surface on $\r^n$. Then  exist a coordinate transformation $\text{\rmfamily \bfseries y}=\varphi(\text{\rmfamily \bfseries x})$, such that the surface is $y_n=0$ under the new coordinate.
  \end{question}
  \begin{proof}
   Since $\text{\rmfamily \bfseries x}=\text{\rmfamily \bfseries f}(\theta_1,\cdots,\theta_{n-1} )$ is a regular surface, then
   \begin{equation*}
    \begin{pmatrix}
     \partial_1 f_1 & \partial_2 f_1 & \cdots & \partial_{n-1} f_1\\
     \partial_1 f_2 & \partial_2 f_2 & \cdots & \partial_{n-1} f_2\\
     \partial_1 f_3 & \partial_2 f_3 & \cdots & \partial_{n-1} f_3\\
     \vdots & \vdots & \vdots & \vdots\\
     \partial_1 f_{n-1} & \partial_2 f_{n-1} & \cdots & \partial_{n-1} f_{n-1}\\
     \partial_1 f_n & \partial_2 f_n & \cdots & \partial_{n-1} f_n\\
    \end{pmatrix}_{n\times (n-1)}
   \end{equation*}
   has the rank of $n-1$. Write $\text{\rmfamily \bfseries x}'=(x_1,\cdots,x_{n-1} )$ and $\text{\rmfamily \bfseries f}'=(f_1,\cdots,f_{n-1} )$, then $\text{\rmfamily \bfseries x}=\text{\rmfamily \bfseries f}'(\theta_1,\cdots,\theta_{n-1})$.
   Suppose
   \begin{equation*}
    \frac{\pa(f_1,\cdots,f_{n-1} )}{\partial( \theta_1,\cdots,\theta_{n-1}) } =\begin{pmatrix}
     \partial_1 f_1 & \partial_2 f_1 & \cdots & \partial_{n-1} f_1\\
     \partial_1 f_2 & \partial_2 f_2 & \cdots & \partial_{n-1} f_2\\
     \partial_1 f_3 & \partial_2 f_3 & \cdots & \partial_{n-1} f_3\\
     \vdots & \vdots & \vdots & \vdots\\
     \partial_1 f_{n-1} & \partial_2 f_{n-1} & \cdots & \partial_{n-1} f_{n-1}\\
    \end{pmatrix}_{(n-1)\times (n-1)}
   \end{equation*}
   has the rank $n-1$.  Then $\text{\rmfamily \bfseries f}^{-1}$ exists and $(\theta_1,\cdots,\theta_{n-1} )=\text{\rmfamily \bfseries f}^{-1}(\text{\rmfamily \bfseries x}' )$. We can plug this into $x_n=f_n(\theta_1,\cdots,\theta_{n-1} )$ to obtain
   \begin{equation*}
    x_n=f_n(\text{\rmfamily \bfseries f}^{-1}(\text{\rmfamily \bfseries x}' ) ):=h(x_1,\cdots,x_{n-1} ).
   \end{equation*}
   So the surface has another expression: $x_n=h(x_1,\cdots,x_{n-1})$ with parameters $(x_1,\cdots,x_{n-1} )$.

   Consider the following mapping:
   \begin{equation*}
    \varphi: \left\{
    \begin{split}
     y_1&=x_1\\
     y_2&=x_2\\
     &\vdots\\
     y_n&=x_n-h(x_1,\cdots,x_{n-1} )
    \end{split}\right.
   \end{equation*}
   and
   \begin{equation*}
     \det D\varphi(\text{\rmfamily \bfseries x} )=\det
    \begin{pmatrix}
     1 & 0 & \cdots & 0 & 0\\
     0 & 1 & \cdots & 0 & 0\\
     \vdots & \vdots & & \vdots & \vdots\\
     0 & 0 & \cdots & 1 &0\\
     -\partial_1 h & -\partial_2 h & \cdots & -\partial_{n-1} h & 1
    \end{pmatrix}
    =1\neq 0.
   \end{equation*}
   Thus this is indeed a coordinate transformation and the surface in this new coordinate is
   \begin{equation*}
    y_n=0,
   \end{equation*}
   which is what we desire.
  \end{proof}

  \begin{question}
   Suppose $\{f_n\}$ is a sequence of measurable functions on $(\Omega,\f )$.  Show
   \begin{equation*}
    \li{n\to\infty}\(\sup_{k\leq n}f_k\)=\sup_{n\in\n}f_n.
   \end{equation*}
  \end{question}
  \begin{proof}
   Since $\forall n\in\n,\sup_{k\leq n}f_k\leq \sup_nf_n$, let $n\to\infty$, we will obtain
   \begin{equation*}
    \li{n\to\infty}\(\sup_{k\leq n}f_k\)\leq \sup_{n\in\n}f_n.
   \end{equation*}
   On the contrary, for all $x\in\Omega$, exists a sequence $\{n_\ell\},\st f_{n_k}(x)\uparrow \sup_{n\in\n}f(x)$. Moreover
   \begin{equation*}
    \sup_{k\leq n_k}f_k(x)\geq f_{n_k}(x).
   \end{equation*}
   Let $\ell\to\infty$ in the above inequality, we will obtain
   \begin{equation*}
    \sup_{n\in\n}f_n(x)\geq \li{\ell\to\infty}\(\sup_{k\leq n_\ell}f_k(x)\)\geq\li{\ell\to\infty} f_{n_\ell}(x)=\sup_{n\in\n}f_n(x).
   \end{equation*}
   Because of the monotonicity of $\sup_{k\leq n}f_k$ with respect to $n\in\n$, then $\sup_{k\leq n}f_k$ and $\sup_{k\leq n_\ell}f_k$ share the same limit. Thus
   \begin{equation*}
    \li{n\to\infty}\(\sup_{k\leq n}f_k\)=\sup_{n\in\n}f_n.
   \end{equation*}
  \end{proof}

  \begin{question}
   Given function $f(x,y):X\times Y\to \r$, where $X,Y$ are two spaces. Then we have
   \begin{equation*}
    \begin{split}
     \sup_{y\in Y} \sup_{x\in X}f(x,y)=\sup_{x\in X} \sup_{y\in Y}f(x,y),\\
     \inf_{y\in Y} \inf_{x\in X}f(x,y)=\inf_{x\in X} \inf_{y\in Y}f(x,y).
    \end{split}
   \end{equation*}
  \end{question}
  \begin{proof}
   We just prove the first equality, since we can apply the first equality to $g(x,y):=-f(x,y)$. At first, for each $y\in Y$, we have
   \begin{equation*}
    \sup_{x\in X}f(x,y)\leq \sup_{x\in X} \[\sup_{y\in Y}f(x,y)\].
   \end{equation*}
   Let $\sup_{y\in Y}\cdot$ at both sides to obtain
   \begin{equation*}
    \sup_{y\in Y} \sup_{x\in X}f(x,y)\leq \sup_{x\in X} \sup_{y\in Y}f(x,y)
   \end{equation*}
   Then the symmetry between $x$ and $y$ guarantees the inverse inequality also holds. So we have completed the proof.
  \end{proof}

  \begin{question}
   Suppose $(E,d)$ is a metric space and $A\subset E$. Then the function $x\mapsto d(x,A)$ is continuous.
  \end{question}
  \begin{proof}
   For all $x,y,z\in E$, we have
   \begin{equation*}
    d(x,z)\leq d(x,y)+d(y,z).
   \end{equation*}
   Take $\inf_{z\in A}\cdot $ in both side of the above inequality:
   \begin{equation*}
    \begin{split}
     d(x,A)&=\inf_{z\in A}d(x,z)\leq \inf_{z\in A}[d(x,y)+d(y,z)]\\
     &=d(x,y)+\inf_{z\in A}d(y,z)=d(x,y)+d(y,A),
   \end{split}
   \end{equation*}
   which implies
   \begin{equation*}
    d(x,A)-d(y,A)\leq d(x,y).
   \end{equation*}
   Interchange the positions of $x$ and $y$, we will obtain
   \begin{equation*}
    |d(x,A)-d(y,A)|\leq d(x,y).
   \end{equation*}
   This is what we desire.
  \end{proof}

  \begin{question}
   Suppose $(E,d)$ is a metric space and $A,B,C\subset E$. Then we have the following triangle inequalities:
   \begin{equation*}
    \begin{split}
     d(x,A)&\leq d(x,y)+d(y,A),\\
     d(A,B)&\leq d(A,y)+d(y,B).
    \end{split}
   \end{equation*}
   Moreover,
   \begin{equation*}
    \begin{split}
     d(x,y)\leq d(x,A)+d(A,y),\\
     d(x,B)\leq d(x,A)+d(A,B),\\
     d(A,C)\leq d(A,B)+d(B,C).
    \end{split}
   \end{equation*}
    will not hold.
  \end{question}
  \begin{proof}
   The previous {\bfseries QUESTION} has shown $d(x,A)\leq d(x,y)+d(y,A)$ holds. We let $\inf_{x\in B}\cdot $ in this inequality to obtain
   \begin{equation*}
    d(A,B)=\inf_{x\in B}d(A,x)\leq\inf_{x\in B}\[d(A,y)+d(y,x)\]=d(A,y)+d(y,B).
   \end{equation*}
   Next we use examples to show the last three inequalities will not hold. Consider $E=\r^2,A=B(0,1),B=B(2/3,1),C=B(7/3,1)$ and $x=(-1,0),y=(2/3,0)$. Then
   \begin{equation*}
    \begin{split}
     \frac{5}{3} =d(x,y)>d(x,A)+d(A,y)=0+0=0,\\
     \frac{1}{3} =d(x,B)> d(x,A)+d(A,B)=0+0=0,\\
     \frac{1}{3} =d(A,C)> d(A,B)+d(B,C)=0+0=0.
    \end{split}
   \end{equation*}
  \end{proof}

  \begin{question}
   For any right (left) limiting function $f:\r\to \r$, the right (left) limit function $g:\r\to \r,x\mapsto f(x+)$ ($h:\r\to \r,x\mapsto f(x-)$) is right (left) continuous. Moreover, if $f$ has left (right) limit pointwise, then $g(x-)$ ($h(x+)$) also exist. Furthermore, we have $g(x-)=f(x-)$ ($h(x+)=f(x+)$)
  \end{question}
  \begin{proof}
   \begin{enumerate}
    \item We just prove the assertion for the right limit function. In fact, by the definition of $g$, we have: $\forall\varepsilon>0,\exists\delta>0$, when $0<y-x<2\delta$, then
   \begin{equation*}
    |f(y)-f(x+)|<\varepsilon/2.
   \end{equation*}
   For any $0<z-x<\delta$, we can find a sequence $\{y_n\}$ satisfying $0<y_n-x<2\delta$ for all $n\in\n_+$ and $y_n\downarrow z$. Replace $y$ with $y_n$ and let $n\to\infty$ in the above inequality to obtain
   \begin{equation*}
    |g(z)-g(x)|=|f(z+)-f(x+)|\leq\varepsilon/2<\varepsilon.
   \end{equation*}
   This means $g$ is right continuous.
    \item For each $x\in\r$, by the definition of right limit, exists $\Delta_s>0$, such that
    \begin{equation*}
     |f(x-s+\Delta_s )-f((x-s)+)|<\varepsilon/3.
    \end{equation*}
    Similarly, exists $\Delta_h>0$, such that
    \begin{equation*}
     |f(x-h+\Delta_h )-f((x-h)+)|<\varepsilon/3.
    \end{equation*}
    Since $f$ has left limit pointwise, consider the given $x\in\r$, then exists $\delta>0$, such that for each $s',h'>0,|s'-h'|<\delta$, we have
    \begin{equation*}
     |f(x-s')-f(x-h')|<\varepsilon/3.
    \end{equation*}
    Thus we can let $s,h\in\r_+$ and $|s-h|<\delta$. Moreover, we let $\Delta_s$ and $\Delta_h$ small enough to satisfy
    \begin{equation*}
     0<\Delta_s=\Delta_h< \min\{s,h\}.
    \end{equation*}
    Replace $s'$ with $s-\Delta_s$ and replace $h'$ with $h-\Delta_h$, we will obtain $|s'-h'|=|s-h|<\delta$ and thus
    \begin{equation*}
     \begin{split}
      &|g(x-s)-g(x-h)|=|f((x-s)+)-f((x-h)+)|\\
      \leq& |f((x-s)+)-f(x-s+\Delta_s ) |+|f(x-s+\Delta_s )-f(x-h+\Delta_h ) |\\
      &+|f(x-h+\Delta_h )- f((x-h)+)|\\
      <&\varepsilon/3+\varepsilon/3+\varepsilon/3=\varepsilon.
     \end{split}
    \end{equation*}
    By {\sffamily Cauchy's criterion}, we know $\lim_{z\to x-}g(z)$ exists.
    \item By the definition of $f(x-)$, for each $\varepsilon>0$, exists $\delta>0$, for all $0<\Delta<\delta$, we have
    \begin{equation*}
     |f(x-\Delta )-f(x-)|<\varepsilon/2.
    \end{equation*}
    By the definition of $g$ and the definition of $f(x+)$ at $x-\Delta $, for each $\varepsilon>0,\exists \delta_\Delta >0$, for all $0<\Delta' <\delta_\Delta$, we have
    \begin{equation*}
     |g(x-\Delta)-f(x-\Delta+\Delta' )|=|f((x-\Delta)+)-f(x-\Delta+\Delta' )|<\varepsilon/2.
    \end{equation*}
    For any $0<\Delta<\delta$, we pick $0<\Delta'<\delta_\Delta\wedge \Delta $. Thus $0<\Delta-\Delta'<\delta$. Combine the above two formulas to obtain
    \begin{equation*}
     \begin{split}
      &|g(x-\Delta )-f(x-)|= |f((x-\Delta)+ )-f(x-)|\\
      \leq & |f((x-\Delta)+ )-f(x-\Delta+\Delta' ) |+|f(x-\Delta+\Delta' )-f(x-)|\\
      < &\varepsilon/2+\varepsilon/2=\varepsilon.
     \end{split}
    \end{equation*}
    This shows $g(x-)$ exists and $g(x-)=f(x-)$.
   \end{enumerate}
  \end{proof}

  \begin{question}
   Suppose $u:U\to \r$  and $U\subset\r^n$ is an open set. Prove $u$ can be continuously extended to $\partial U$, namely, $u\in C(\bar{U} )$ iff $u$ is uniformly continuous on every bounded subset of $U$.
  \end{question}
  \begin{proof}
   \begin{enumerate}
    \item If $u\in C(\bar{U})$, then for all bounded subset $V\subset U$, we have $u\in C(\bar{V})$. Notice $\bar{V}$ is a compact subset, so $u$ is uniformly continuous on $\bar{V}$, as well as $V$.
    \item For all $x\in\partial U$, consider the $\gamma$-neighborhood of $x$, denoted by $B(x,\gamma )\cap U$. Obviously, $B(x,\gamma )\cap U$ is a bounded subset of $U$. Hence $u$ is uniformly continuous on $B(x,\gamma )\cap U$. This means: $\forall\varepsilon>0,\exists \delta>0$, for each $y,z\in B(x,\gamma )\cap U$ and $|y-z|<\delta$, then $|u(y)-u(z)|<\varepsilon$.
    So $\forall\varepsilon>0,\exists\delta>0$, for each $y,z\in B(x,\delta/2 )\cap U$, we have $|y-z|<\delta$, which leads to $|u(y)-u(z)|<\varepsilon$. By Cauchy's criterion, we know $\lim_{U\ni y\to x}u(y)$ exists. Hence $u$ can be continuously extended to $\partial U$.
   \end{enumerate}
  \end{proof}

  \begin{question}
   Given function $f:\r\to \r,x\mapsto |x|^\alpha$ for each $\alpha\in\r$. Then calculate its derivative.
  \end{question}
  \begin{proof}
   Notice for each $\alpha\in\r$, we have the following equality:
   \begin{equation*}
    \lim_{x\to0}\frac{(1+x)^\alpha-1 }{x}=\alpha.
   \end{equation*}
   Then
   \begin{enumerate}
    \item If $x\neq 0$, then for each $\alpha\in\r$, $f(x)$ has definition. Observe when $h>0$ small enough, we can guarantee $x+h\neq 0$. So $f(x+h)$ has definition as well. Then
    \begin{equation*}
     \begin{split}
      f'(x)&=\lim_{h\to0}\frac{f(x+h)-f(x)}{h}=\lim_{h\to0}\frac{|x+h|^\alpha -|x|^\alpha}{h}\\
      &=\lim_{h\to0}|x|^\alpha \frac{|1+h/x|^\alpha -1}{h}\\
      &=\lim_{h\to0}|x|^{\alpha-1}\mathrm{sgn}(x) \frac{|1+h/x|^\alpha -1}{h/x}.
     \end{split}
    \end{equation*}
    Observe $h\to0$ implies $h/x\to 0$ and when $h$ small enough, we have $1+h/x>0$. Thus we can remove the absolute value sign and utilize the above equiality to obtain
    \begin{equation*}
     \begin{split}
      f'(x)&=\lim_{h\to0}|x|^{\alpha-1}\mathrm{sgn}(x) \frac{(1+h/x)^\alpha -1}{h/x}\\
      &=\alpha |x|^{\alpha-1}\mathrm{sgn}(x).
     \end{split}
    \end{equation*}
    \item If $x=0$, then $f(x)$ has definition only for $\alpha>0$. So
    \begin{equation*}
     \begin{split}
      f'(0)&=\lim_{h\to0}\frac{f(h)-f(0)}{h}=\lim_{h\to0}\frac{|h|^\alpha-0}{h}\\
      &= \lim_{h\to0}\(\mathrm{sgn}(h)\frac{|h|^\alpha-0}{|h|}\)\\
      &=\lim_{h\to0}\(\mathrm{sgn}(h)|h|^{\alpha-1}\).
     \end{split}
    \end{equation*}
    \begin{enumerate}
     \item If $\alpha\in(0,1)$, then
     \begin{equation*}
      \mathrm{sgn}(h)=\pm1,\quad \lim_{h\to0}|h|^{\alpha-1}=\infty.
     \end{equation*}
     So the above limit does not exist.
     \item If $\alpha=1$, then
     \begin{equation*}
      \lim_{h\to0}\(\mathrm{sgn}(h)|h|^{\alpha-1}\)=\lim_{h\to0}\mathrm{sgn}(h).
     \end{equation*}
     obviously, this limit does not exist as well.
     \item If $\alpha\in(1,\infty)$, then
     \begin{equation*}
      f'(0)=\lim_{h\to0}\(\mathrm{sgn}(h)|h|^{\alpha-1}\)=0.
     \end{equation*}
    \end{enumerate}
   \end{enumerate}
   In conclusion, we the following results:
   \begin{table}[h]
    \centering
    \begin{tabular}{|c|c|c|c|}
     \hline
     $\bm{f(x)=|x|^\alpha}$ & $\alpha\leq 0$ & $0<\alpha\leq 1$ & $\alpha>1$ \\
     \hline
     $x\neq 0$ &\multicolumn{3}{|c|}{ $f'(x)=\alpha |x|^{\alpha-1}\mathrm{sgn}(x)$}  \\
     \hline
     $x=0$ & $f(0)$ is not defined  & $f'(0)$ does not exist & $f'(0)=0$\\
     \hline
    \end{tabular}
   \end{table}
  \end{proof}

  \begin{question}
   Suppose $D$ is a countable dense subset of $\r$. Then for each $x\in\r$, exists $\{y_n\}$ and $\{z_n\}$ such that $y_n\uparrow x$ and $z_n\downarrow x$.
  \end{question}
  \begin{proof}
   Suppose we do not have such $\{y_n\}$, such that $y_n\uparrow x$. Then exists $\delta>0$, such that $ (x-\delta,x]\cap D =\emptyset$. Select and fix $z\in (x-\delta,x)$. We can find a neighborhood of $z$, denoted as $U(z)$, such that $U(z)\subset (x-\delta,x]$. So $U(z)\cap D=\emptyset$. This contradicts the density of $D$ in $\r$.
  \end{proof}

  \begin{question}
   Suppose $D=\{r_n\}$ is a dense subset of $[a,b]$ and $f:[a,b]\to\r$ is right continuous. Write $D_n=[a,b]\cap \{r_1,\cdots,r_n,b\}$. Then
   \begin{equation*}
    \sup_{x\in U} f(x)=\lim_{n\to\infty}\sup_{s\in D_n}f(s).
   \end{equation*}
  \end{question}
  \begin{proof}
   Similar to {\bfseries QUESTION 2.10}, we can easily verify
   \begin{equation*}
    \lim_{n\to\infty}\sup_{s\in D_n}f(s)=\sup_{s\in D\cup \{b\}}f(s).
   \end{equation*}
   So we only need to prove $\sup_{x\in U} f(x)=\sup_{x\in D\cup \{b\}}f(x)$.
    By the definition of $\sup_{x\in U} f(x)$, we know exists $\{x_n\}\subset U$ such that $f(x_n)\uparrow \sup_{x\in U} f(x)$. By the right continuity of $f$, for each $n\geq 1$, we can find $s_n\in D\cup\{b\}$, such that  $x_n\leq s_n$ and $|f(x_n)-f(s_n)|<1/n$. Thus we have
    \begin{equation*}
     \begin{split}
      |f(s_n)-\sup_{x\in U} f(x)|&\leq |f(s_n)-f(x_n) |+|f(x_n)-\sup_{x\in U} f(x)|\\
      &<1/n+|f(x_n)-\sup_{x\in U} f(x)|.
     \end{split}
    \end{equation*}
    Let $n\to\infty$, we will obtain $f(s_n)\to \sup_{x\in U} f(x)$. Therefore
    \begin{equation*}
     \sup_{x\in D\cup \{b\}}f(s)\geq \lim_{n\to\infty}f(s_n)=\sup_{x\in U} f(x).
    \end{equation*}
    It is obvious $\sup_{x\in D\cup \{b\}}f(s)\leq \sup_{x\in U} f(x)$.
  \end{proof}

  \begin{question}
   Consider the Lebesgue measure space $(\r,\b(\r),dx)$. The Dominance convergence theorem implies: if $f_n\to f$ uniformly, then for each compact set $A\in\b(\r)$, we have
   \begin{equation*}
    \int_Af_n(x)dx\to \int_A f(x)dx.
   \end{equation*}
  \end{question}
  \begin{proof}
   Since $f_n\to f$ uniformly, then $\forall\varepsilon>0$, exists $N\geq 1$, when $n>N$, for each $x\in\r$, we have $|f_n(x)-f(x)|<\varepsilon$. Because the limit has nothing to do with the first finite items, we can omit the first $N$ items and $|g_n(x)|:=|f_n(x)-f(x)|$ is dominated by the integrable function $\varepsilon$. Thus the {\sffamily Dominance convergence theorem} tells us
   \begin{equation*}
    \left|\int_A f_n(x)dx-\int_Af(x)dx\right|\leq \int_A|f_n(x)-f(x)|dx\to0.
   \end{equation*}
  \end{proof}

  \begin{question}
   Suppose $x\in\r$ and $|x|\leq 1$. Let $P_0(x)=0$ and
   \begin{equation*}
    P_{n+1}(x)=P_n(x)+\frac{1}{2}\[x^2-P_n(x)^2 \],\quad n\geq 0.
   \end{equation*}
   Then $P_n(x)\uparrow |x|$.
  \end{question}
  \begin{proof}
   First we verify $P_n(x)\leq |x|$ for each $n\geq 0$. We utilize induction on $n.$ When $n=0$, the assertion naturally holds. Assume $P_n(x)\leq |x|$, then
   \begin{equation*}
    P_{n+1}(x)=P_n(x)+\frac{1}{2}\[x^2-P_n(x)^2 \]=-\frac{1}{2}\[P_n(x)\]^2+P_n(x)+\frac{x^2}{2}
   \end{equation*}
   is a quadratic function with respect to the variable $P_n(x)$ and interval $P_n(x)\in[-|x|,|x|]$. It is easy to know the quadratic function is increasing in $[-|x|,|x|]$. So
   \begin{equation*}
    P_{n+1}(x)=P_n(x)+\frac{1}{2}\[x^2-P_n(x)^2 \]\leq P_{n+1}(x)\big|_{P_n(x)=|x|}= |x|.
   \end{equation*}
   Till now, we have proved the assertion. Next, since $P_n(x)\leq |x|$ for each $n\geq 0$, then
   \begin{equation*}
    P_{n+1}(x)-P_n(x)=\frac{1}{2}\[x^2-P_n(x)^2 \]\geq 0,
   \end{equation*}
   which concludes that $\{P_n(x)\}$ is an increasing sequence with upper bound $|x|$. So the limit $z(x)=\lim_{n\to\infty}P_n(x)$ exists. Thus we can let $n\to\infty$ to obtain
   \begin{equation*}
    z(x)=z(x)+\frac{1}{2}\left\{x^2-\[z(x)\]^2\right\}\Longrightarrow z(x)=\pm|x|.
   \end{equation*}
   Notice $\{P_n(x)\}$ is an increasing sequence, thus
   \begin{equation*}
    z(x)=\lim_{n\to\infty}P_n(x)\geq P_1(x)=\frac{x^2}{2}>-|x|.
   \end{equation*}
   So $z(x)=|x|$.
  \end{proof}

  \begin{question}
   Suppose $f(x):\r\to\r$ is a nondecreasing and continuous function along a dense subset $\Lambda $ of $\r$. Namely, for each $x,y\in\Lambda$ and $x<y$, we have $f(x)\leq f(y)$. Then $f(x)$ is nondecreasing in $\r$.
  \end{question}
  \begin{proof}
   For each $x<y\in\r$, exists $\Lambda\ni x_n\uparrow x$ and $\Lambda\ni y_n\downarrow y$. Therefore, the continuity of $f(x)$ yields
   \begin{equation*}
    f(y)-f(x)=\lim_{n\to\infty}[f(y_n)-f(x_n)]\geq 0.
   \end{equation*}
   So $f(x)$ is nondecreasing in $\r$.
  \end{proof}

  \begin{question}
   Suppose for each $n\geq 1,f_n(x):[0,T]\to\r$ is a nondecreasing function along the finite sequence $(x_i^n:0\leq i\leq p_n)$, which satisfies $0=x_0^n<x_1^n<\cdots< x_{p_n}^{n}=T$ and $\max_{0\leq i\leq p_n}|x_{i+1}^{n}-x_{i}^{n}|\to 0,n\to\infty$. \ Besides, we have $(x_i^n:0\leq i\leq p_n)\subset (x_i^{n+1}:0\leq i\leq p_{n+1})$ for each $n\geq 1$.
   Assume $f(x):=\lim_{n\to\infty}f_n(x)$ exists pointwise and is continuous, then $f(x):[0,T]\to\r$ is a nondecreasing function.
  \end{question}
  \begin{proof}
   First we show $f(x)$ is a nondecreasing function along the countable set $\Lambda:= \cup_n(x_i^n:0\leq i\leq p_n)$. In fact, for each $x_i^n,x_j^m\in\Lambda$ and $x_i^n<x_j^m$, without loss of generality, assume $n\geq m$, then exists $0\leq j'\leq p_n$, such that $x_j^m=x_{j'}^n$. Thus we have $x_i^n\leq x_{j'}^n=x_j^m$. Utilize the monotonicity of $f_n(x)$ on $(x_i^n:0\leq i\leq p_n)$, we have
   \begin{equation*}
    f(x_i^n)=\lim_{k\to\infty}f_k(x_i^n)\leq \lim_{k\to\infty}f_k(x_{j'}^n)=f(x_j^m).
   \end{equation*}
   This shows $f(x)$ is a nondecreasing function along the countable set $\Lambda$.

   Then we show $\Lambda$ is a dense subset of $[0,T]$. In fact, for any $x\in [0,T]$, notice $\max_{0\leq i\leq p_n}|x_{i+1}^{n}-x_{i}^{n}|\to 0,n\to\infty$, then for any $\varepsilon>0$, exists $n\geq 1$ such that $\max_{0\leq i\leq p_n}|x_{i+1}^{n}-x_{i}^{n}|<\varepsilon$. This implies exists $x_i^n\in (x_i^n:0\leq i\leq p_n)$ satisfying $|x-x_i^n|<\varepsilon$. So $\Lambda$ is a dense subset of $[0,T]$.

   Finally, by the previous {\bfseries QUESTION}, we know $f(x):[0,T]\to\r$ is a nondecreasing function.
  \end{proof}

  \begin{question}
   Given two continuous functions $f,g:\r\to\r$. Suppose when $x\to x_0$, we have $f(x)>g(x)+o(1)$. Then $f(x_0)>g(x_0)$.
  \end{question}
  \begin{proof}
   If $f(x_0)=g(x_0)$, consider $f(x)-g(x)$. We have
   \begin{equation*}
    \lim_{x\to x_0}[f(x)-g(x)]=f(x_0)-g(x_0)=0,
   \end{equation*}
   which means $f(x)-g(x)=o(1)\Longrightarrow f(x)=g(x)+o(1)$ when $x\to x_0$. This contradicts  $f(x)>g(x)+o(1)$.
  \end{proof}
  
  \begin{question}
    Given $f(x,y):\r^2\to\r$ and define $U=(x'-\alpha ,x'+\alpha )\times(y'-\beta,y'+\beta)=:U_x\times U_y$. Assume  $\lim_{x\to x'}f(x,y)=g(y)$ and $\lim_{y\to y'}f(x,y)=h(x)$ exist in the region $U$. And the repeated limits 
    \begin{equation*}
    	\lim_{x\to x'}\lim_{y\to y'}f(x,y)\quad\mathrm{and} \quad\lim_{y\to y'}\lim_{x\to x'}f(x,y)
    \end{equation*}
     also exists in the region $U$. If one of the above limit is uniform convergence, then we have
    \begin{equation*}
      \lim_{x\to x'}h(x)=\lim_{y\to y'}g(y).
    \end{equation*}
  \end{question}
  \begin{proof}
    Because of the similarity, we just postulate  $\lim_{x\to x'}f(x,y)=g(y)$ is a uniformly convergence. Then $\forall\varepsilon>0$, exists $\delta\in(0,\alpha)$. Such that for any $y\in U_y$ and any $|x-x'|<\delta$,
    \begin{equation*}
      |f(x,y)- g(y)|<\varepsilon/3.
    \end{equation*}
    Moreover, by the definition of limit, exists $\delta'\in(0,\delta)$ such that for any $|y-y'|<\delta'$,
    \begin{equation*}
    	\Big|\lim_{y\to y'}g(y)- g(y)\Big|<\varepsilon/3\quad\mathrm{and} \quad |f(x,y)- h(x)|<\varepsilon/3.
    \end{equation*}
    Therefore, for any $|x-x'|<\delta'$, we have
    \begin{equation*}
      \begin{split}
         \Big|\lim_{y\to y'}g(y)-h(x)\Big|&\leq \Big|\lim_{y\to y'}g(y)- g(y)\Big|+|g(y)- f(x,y)|+|f(x,y)- h(x)| \\
           &<\varepsilon/3+\varepsilon/3+\varepsilon/3=\varepsilon.
      \end{split}
    \end{equation*}
    This means $\lim_{x\to x'}h(x)=\lim_{y\to y'}g(y)$.
  \end{proof}
  
  \begin{question}
  	Given a complex sequence $\{z_n\}\subset\c$. Suppose $\{z_n\}$ is not convergent (in this question, $z_n\to\infty$ is also a convergence sequence), then exists two convergent sub-sequences $\{z_{n'}\}$ and $\{z_{n''}\}$ with disparate limit.
  \end{question}
  \begin{proof}
  	For the given sequence $\{z_n\}$, the compactness of $\c\cup\{\infty\}$ implies there is a convergent sub-sequences $\{z_{n'}\}$, say $z_{n'}\to z_1\in \c\cup\{\infty\}$. We claim exists $\varepsilon>0$, such that there are infinite items in $B(z_1,\varepsilon)^c$. If not, for any $\varepsilon>0$, there are finite items in $B(z_1,\varepsilon)^c$, then $z_n\to z_1$. This contradict $\{z_n\}$ is not convergent. So select such $\varepsilon>0$, consider $\{z_n\}\cap B(z_1,\varepsilon)^c$. The compactness of $\c\cup\{\infty\}$ shows again there is a convergent sub-sequences $\{z_{n''}\}$. obviously, the limits of $\{z_{n'}\}$ and $\{z_{n''}\}$ are different.
  \end{proof}
  
  \begin{question}
  	Given $n=3m,m\in\n_+$ and $i,j,k\in\n,i+j+k=n$. Then 
  	\begin{equation*}
  		\binom{n}{i\ j\ k}\leq \binom{n}{m\ m\ m}.
  	\end{equation*}
  \end{question}
  \begin{proof}
  	Notice
  	\begin{equation*}
  		\binom{n}{i\ j\ k}=\frac{n!}{i!j!k!}.
  	\end{equation*}
  	In order to verify the assertion, we just need to verify $(m!)^3\leq i!j!k!$. In fact, consider the restraint $i,j,k\in\n$ and $i+j+k=n$, without loss of generality, we have three situations:
  	\begin{enumerate}
  		\item $i=m+n_1,j=m+n_2,k=m-n_1-n_2$, where $n_1,n_2\in\n_+$. In this situation, we have
  			\begin{equation*}
  				\begin{split}
  					&\frac{(m!)^3}{i!j!k!} \\
  					=&\frac{ m(m-1)\cdots(m-n_1-n_2+1)}{(m+n_1)(m+n_1-1)\cdots(m+1)(m+n_2)(m+n_2-1)\cdots(m+1)}.
  				\end{split}
  	    	\end{equation*}
  	    	Observe the numerator and the denominator both has $n_1+n_2$ items and each item of numerator is smaller than each item of denominator. So $(m!)^3\leq i!j!k!$ holds.
  		\item $i=m+n_1+n_2,j=m-n_1,k=m-n_2$, where $n_1,n_2\in\n_+$. In this situation, we have

  		\begin{equation*}
  			\frac{(m!)^3}{i!j!k!}=\frac{ m(m-1)\cdots(m-n_1+1)m(m-1)\cdots(m-n_2+1)}{(m+n_1+n_2)(m+n_1+n_2-1)\cdots(m+1)}.
  		\end{equation*}
  		Observe the numerator and the denominator both has $n_1+n_2$ items and each item of numerator is smaller than each item of denominator. So $(m!)^3\leq i!j!k!$ holds.
  		\item $i=m,j=m+n_1,k=m-n_1$, where $n_1\in\n_+$. The same argument will also gives us $(m!)^3\leq i!j!k!$.
  	\end{enumerate}
  \end{proof}
  
  \begin{question}
  	Given $n\in\n_+$ and $i,j,k\in\n,i+j+k=n$. Then 
  	\begin{equation*}
  		\binom{n}{i\ j\ k}\leq \frac{n!}{\Gamma^3(n/3+1)} .
  	\end{equation*}
  \end{question}
  \begin{proof}
  	Observe $\lg\Gamma(x)$ is a convex function, then
  	\begin{equation*}
  		\begin{split}
  			\lg \binom{n}{i\ j\ k}&=\lg \frac{n!}{i!j!k!}=\lg n!-(\lg i!+\lg j!+\lg k!)\\
  			&=\lg n!-3\times\frac{1}{3} (\lg \Gamma(i+1) +\lg \Gamma(j+1)+\lg \Gamma(k+1))\\
  			&\leq \lg n!-3\lg\Gamma(\frac{i+j+k+3}{3} )\\
  			&=\lg n!-3\lg\Gamma(\frac{n}{3}+1 )\\
  			&=\lg \frac{n!}{\Gamma^3(n/3+1)},
  		\end{split}
  	\end{equation*}
  	which is equivalent to the assertion we want to prove. When $n=3m,m\in\n_+$, we will obtain the conclusion of the previous {\bfseries QUESTION}.
  \end{proof}
  
  \begin{question}
  	Prove the following equality:
  	\begin{equation*}
  		x_0+\sum_{n=2}^\infty x_{n-1}\prod_{i=0}^{n-2}(1-x_i)=1-\prod_{n=0}^\infty (1-x_n),
  	\end{equation*}
  	where the left side is convergent.
  \end{question}
  \begin{proof}
  	Note
  	\begin{equation*}
  		\begin{split}
  			&1-\(x_0+\sum_{n=2}^\infty x_{n-1}\prod_{i=0}^{n-2}(1-x_i)\)\\
  			=&1-x_0-x_1(1-x_0)-\sum_{n=3}^\infty x_{n-1}\prod_{i=0}^{n-2}(1-x_i)\\
  			=&(1-x_0)\(1-x_1-\sum_{n=3}^\infty x_{n-1}\prod_{i=1}^{n-2}(1-x_i)\)\\
  			=&(1-x_0)(1-x_1)\(1-x_2-\sum_{n=4}^\infty x_{n-1}\prod_{i=2}^{n-2}(1-x_i)\)\\
  			=&\cdots =\prod_{n=0}^m (1-x_n)\(1-x_{m+	1}-\sum_{n=m+3}^\infty x_{n-1}\prod_{i=m+1}^{n-2}(1-x_i)\) \\
  			=&\prod_{n=0}^\infty (1-x_n).
  		\end{split}
  	\end{equation*}
  	The last equality holds since the convergence of 
  	\begin{equation*}
  		\sum_{n=m+3}^\infty x_{n-1}\prod_{i=m+1}^{n-2}(1-x_i) \quad \text{and}\quad \sum_{n=2}^\infty x_{n-1}\prod_{i=0}^{n-2}(1-x_i)
  	\end{equation*}
  	are the same. Thus the the desired result follows.
  \end{proof}
  
  \begin{question}
  	For any real sequence $\{a_n\}_{n\geq 1}$, we have the following ine	qualities:
  	\begin{equation*}
  		\varliminf_{n\to\infty} \frac{a_{n+1}}{a_n}\leq \varliminf_{n\to\infty}\sqrt[n]{a_n}\leq \varlimsup_{n\to\infty}\sqrt[n]{a_n}\leq \varlimsup_{n\to\infty} \frac{a_{n+1}}{a_n}.
  	\end{equation*}
  \end{question}
  \begin{proof}
  	The inequality in the middle is trivial, we just prove the first and the final ones.
  	\begin{enumerate}
  		\item Write $K:=\varliminf_{n\to\infty} a_{n+1}/a_n$, then for any $\varepsilon>0$, exists $N\in\n$ such that when $n>N$,
  		\begin{equation*}
  			\frac{a_{n}}{a_{n-1}}>K-\varepsilon\Longrightarrow a_{n}>(K-\varepsilon)a_{n-1}>\cdots>(K-\varepsilon)^{n-1}a_1.
  		\end{equation*}
		Based on this, we have
		\begin{equation*}
			\begin{split}
				\sqrt[n]{a_n}&>(K-\varepsilon)^{1-1/n}\sqrt[n]{a_1}\\
				&\Longrightarrow \varliminf_{n\to\infty}\sqrt[n]{a_n}\geq\lim_{n\to\infty}(K-\varepsilon)^{1-1/n}\sqrt[n]{a_1}=K-\varepsilon.
			\end{split}
		\end{equation*} 
		Let $\varepsilon\to0$ will complete the proof.
  		\item Wirte $L:=\varlimsup_{n\to\infty}a_{n+1}/a_n$, then for any $\varepsilon>0$, exists $N\in\n$ such that when $n>N$,
  		\begin{equation*}
  			\frac{a_{n}}{a_{n-1}}<L+\varepsilon\Longrightarrow a_{n}<(L+\varepsilon)a_{n-1}<\cdots<(L+\varepsilon)^{n-1}a_1.
  		\end{equation*}
		Based on this, we have
		\begin{equation*}
			\begin{split}
				\sqrt[n]{a_n}&<(L+\varepsilon)^{1-1/n}\sqrt[n]{a_1}\\
				&\Longrightarrow \varlimsup_{n\to\infty}\sqrt[n]{a_n}\leq\lim_{n\to\infty}(L+\varepsilon)^{1-1/n}\sqrt[n]{a_1}=L+\varepsilon.
			\end{split}
		\end{equation*} 
		Let $\varepsilon\to0$ will complete the proof.
  	\end{enumerate}
  \end{proof}
  
  \begin{question}
  	Suppose $f:\r\to\r$ is a uniformly continuous function and the generalized integral 
  	\begin{equation*}
  		\int_0^\infty f(x)dx
  	\end{equation*}
  	exists, then $\lim_{x\to\infty}f(x)=0$.
  \end{question}
  \begin{proof}
  	Since the generalized integral exists, then for any $\varepsilon>0$, exists $M>0$, when $x>M$ and for any $\Delta>0$,
  	\begin{equation*}
  		\left|\int_x^{x+\Delta} f(y)dy\right|<\varepsilon/2.
  	\end{equation*}
  	Besides, the uniform continuity of ensure for the $\varepsilon>0$ mentioned above, exists $\delta>0$. When $|x-y|<\delta$, 
  	\begin{equation*}
  		|f(x)-f(y)|<\varepsilon/2.
  	\end{equation*}
  	Let $\Delta<\delta$, then for any $x>M$, we have
  	\begin{equation*}
  		 \begin{split}
  		 	|f(x)|&\leq \frac{1}{\Delta} \int_x^{x+\Delta}|f(x)-f(y)|dy+ \left|\int_x^{x+\Delta} f(y)dy\right| \\
  		 	&<\frac{1}{\Delta}\Delta \varepsilon/2+\varepsilon/2=\varepsilon.
  		 \end{split}
  	\end{equation*}
  \end{proof}
  
  \begin{question}
  	Suppose $f:\r\to\r$ is an increasing function and the generalized integral 
  	\begin{equation*}
  		\int_0^\infty f(x)dx
  	\end{equation*}
  	exists, then $\lim_{x\to\infty}xf(x)=0$.
  \end{question}
  \begin{proof}
  	If $\lim_{x\to\infty}xf(x)\neq 0$, then exists a sequence $\{x_n\}_{n\geq 1}$ such that $\lim_{n\to\infty}x_nf(x_n)=c\neq 0$. Note $f(x)$ is increasing, then 
  	\begin{equation*}
  		\begin{split}
  			\exists\{x_n\}_{n\geq 1},&\text{ such that } \lim_{n\to\infty}x_nf(x_n)=c\neq 0\\
  			&\Longleftrightarrow \lim_{x\to\infty}xf(x)=c\neq 0\\
  			&\Longrightarrow \frac{f(x)}{1/x}=c\neq 0.
  		\end{split}
  	\end{equation*}
  	Therefore, we know the convergence property of
  	\begin{equation*}
  		\int_0^\infty f(x)dx\in\r \text{ and }\int_0^\infty \frac{1}{x} dx=\infty
  	\end{equation*}
  	are identical. This is a contradiction.
  \end{proof}
  
  \begin{question}
  	Suppose $\{A_t\}_{t\geq 0}\subset2^\r$ is a decreasing set sequence. Define 
  	\begin{equation*}
  		\alpha_t:=\inf\{x:x\in A_t\}\text{ and }A_{t-}=\bigcap_{0\leq s<t}A_s.
  	\end{equation*}
  	Then for any $0\leq s<t$, we have
  	\begin{equation*}
  		\lim_{s\uparrow t}\alpha_s \leq \alpha_{t-}:=\inf\{x:x\in A_{t-}\}.
  	\end{equation*}
  \end{question}
  \begin{proof}
  	Since $\{A_t\}_{t\geq 0}$ is decreasing, then it is trivial to find $t\mapsto \alpha_t$ is increasing (not necessarily strict). First, since $A_{t-}\subset A_s$ for any $0\leq s<t$, then we can readily obtain
  	\begin{equation*}
  		\alpha_s\leq \alpha_{t-},\quad \forall s<t\Longrightarrow \lim_{s\uparrow t}\alpha_s \leq \alpha_{t-}.
  	\end{equation*} 
  	Now we give an example to show the strict inequality might happen. Consider 
  	\begin{equation*}
  		A_t=
  		\begin{cases}
  			\(1,6-t\)\cup\(6,\infty\)&,\mbox{if }0\leq t<5, \\
  			\(6,\infty\)&,\mbox{if }t\geq 5.
  		\end{cases}
  	\end{equation*}
  	So $\{A_t\}_{t\geq 0}$ is decreasing and 
  	\begin{equation*}
  		A_{5-}=\bigcap_{0\leq t<5}\(1,6-t\)\cup\(6,\infty\)=  (6,\infty).
  	\end{equation*}
  	Thus $\alpha_t=1$ for any $0\leq t<5$ and $\alpha_{5-}=\alpha_5=6$ implies
  	\begin{equation*}
  		\lim_{t\uparrow 5}\alpha_t=1<6=\alpha_{5-}.
  	\end{equation*}
  \end{proof}
  
  \begin{question}
  	If $f(x)\in C_0(\r^n)$, then $f(x)$ is uniformly continuous.
  \end{question}
  \begin{proof}
  	For any $\varepsilon>0$, since $f(x)\to0,x\to\infty$, then exists $R>0$ satisfying $|f(x)|<\varepsilon/3,\forall x\in B(0,R)^c$. Since $\overline{ B(0,R)}$ is compact, then $f(x)$ is uniformly continuous on $\overline{ B(0,R)}$. Thus exists $\delta>0$, for any $|x-y|<\delta$ and $x,y\in \overline{ B(0,R)}$, we have $|f(x)-f(y)|<\varepsilon/3$. Then we claim for any $|x-y|<\delta$ and $x,y\in\r^n$, we have $|f(x)-f(y)|<\varepsilon$. In fact,
  	\begin{enumerate}
  		\item If $x,y\in \overline{ B(0,R)}$, then the above argument shows $|f(x)-f(y)|<\varepsilon/3<\varepsilon$.
  		\item If $x\in \overline{ B(0,R)}$ and $y\in \overline{ B(0,R)}^c\subset B(0,R)^c$. Then we can find $z\in \partial B(0,R)=\overline{ B(0,R)}\cap B(0,R)^c$ and $|x-z|+|z-y|=|x-y|$. Thus $|x-z|<\delta$ and
  		\begin{equation*}
  			\begin{split}
  				|f(x)-f(y)|&\leq |f(x)-f(z)|+|f(z)-f(y)|\\
  				&< \varepsilon/3 +|f(z)|+|f(y)|<\varepsilon.
  			\end{split}
  		\end{equation*} 
  		\item If $x,y \in B(0,R)^c$, then $|f(x)-f(y)|\leq |f(x)|+|f(y)|<2\varepsilon/3<\varepsilon$.
  	\end{enumerate}
  \end{proof}
  
  \begin{question}
  	Given $f(x,y):\r^m\times\r^n\to\r$. If $f(\cdot,y)$ is uniformly continuous with respect to $y$ and $f(x,\cdot)$ is continuous, then $f(x,y)$ is continuous with respect to $(x,y)$.
  \end{question}
  \begin{proof}
  	For any $\varepsilon>0$, since $f(\cdot,y)$ is uniformly continuous with respect to $y$, then for the given $x\in\r^n$, exists $\delta_x>0$, such that for any $y\in\r^n$, we have 
  	\begin{equation*}
  		|f(x',y)-f(x,y)|<\varepsilon/2,\quad \forall |x'-x|<\delta_x.
  	\end{equation*}
  	And note $f(x,\cdot)$ is continuous, then for the given $x\in\r^n$, exists $\delta_{x,y}>0$ such that 
  	\begin{equation*}
  		|f(x,y')-f(x,y)|<\varepsilon/2,\quad \forall |y'-y|<\delta_{x,y}.
  	\end{equation*}
  	So when $|x-x'|<\delta_x$ and $|y'-y|<\delta_{x,y}$ we have
  	\begin{equation*}
  		|f(x',y')-f(x,y)|\leq |f(x',y')-f(x,y')|+|f(x,y')-f(x,y)|<\varepsilon.
  	\end{equation*}
  	So the claim is complete.
  \end{proof}
  
  \begin{question}
  	Given a complex measurable function $h=f+ig$. Then exists a real measurable function $H$ such that $|h|\leq H$ iff exists two measurable function $F$ and $G$ such that $|f|\leq F$ and $|g|\leq G$. 
  \end{question}
  \begin{proof}
  	Observe 
  	\begin{equation*}
  		|f|\vee |g|\leq  |h|=\sqrt{f^2+g^2}\leq |f|+|g|.
  	\end{equation*}
  	If $|h|\leq H$, then $|f|\vee |g|\leq  |h|\leq H$. Let $F=G=H$ will be fine. On the contrary, if $|f|\leq F$ and $|g|\leq G$, then $|h|\leq |f|+|g|\leq F+G$. Let $H=F+G$ will be fine.
  \end{proof}
  
  \begin{question}
  	Define function $f:\r^d\to \r$,
  	\begin{equation*}
  		f(x)=\prod_{k=1}^d \frac{\sin x_k}{x_k}.
  	\end{equation*}
  	Then we have $f(x)=1-|x|^2/6+O(|x|^4)$ as $|x|\to0$.
  \end{question}
  \begin{proof}
  	As $|x_k|\to 0$, we have
  	\begin{equation*}
  		\frac{\sin x_k}{x_k}=\frac{1}{x_k}\(x_k-\frac{1}{6}x_k^3+O(|x_k|^5) \)=1-\frac{1}{6}x_k^2+O(|x_k|^4).
  	\end{equation*}
  	Then as $|x|\to 0$, which implies $|x_k|\to 0$ for any $1\leq k\leq d$, we have
  	\begin{equation*}
  		\begin{split}
  			\prod_{k=1}^d \frac{\sin x_k}{x_k}&=\prod_{k=1}^d\(1-\frac{1}{6}x_k^2+O(|x_k|^4)\)\\
  			&=1-\frac{1}{6} \sum_{k=1}^d x_k^2+\frac{1}{36}\sum_{i,j=1}^d(x_ix_j)^2-\frac{1}{6}\sum_{i,j=1}^dx_i^2O(|x_j|^4) \cdots
  		\end{split}
  	\end{equation*}
  	By the mean value inequality, note
  	\begin{equation*}
  		\frac{(x_ix_j)^2}{|x|^4}\leq \frac{1}{|x|^4}\(\sqrt{\frac{x_i^2+x_j^2}{2} }\)^4\leq \frac{1}{|x|^4}\(\sqrt{\frac{1}{2}\sum_{k=1}^dx_k^2 }\)^4=\frac{1}{4}
  	\end{equation*}
  	and
  	\begin{equation*}
  		\frac{x_i^2O(|x_j|^4)}{|x|^4}\leq \frac{|x|^2O(|x_j|^4)}{|x_j|^4}\leq M|x|^2.
  	\end{equation*}
  	So $(x_ix_j)^2=O(|x|^4)$ and $x_i^2O(|x_j|^4)=O(|x|^4)$. And the result follows.
  \end{proof}
  
  \begin{question}
  	Given two non negative series satisfying $\sum_n a_n=\infty$ and $\sum_nb_n=\infty$. If $\lim_n a_n=\ell> 0$ (or equivalently $\lim_n b_n=\ell> 0$), then $\sum_na_n\wedge b_n=\infty$. Exemplify $\sum_na_n\wedge b_n=\infty$ might not happen when $\lim_n a_n=\ell>0$ is not correct. 
  \end{question}
  \begin{proof}
  	Write $c_n:=a_n\wedge b_n$. If there are infinite items of $\{c_{n}\}_{n\geq 1}$ satisfying $c_{n_k}=a_{n_k},k\geq 1$, note $a_{n_k}\to\ell>0,k\to\infty$, then 
  	\begin{equation*}
  		\sum_n a_n\wedge b_n =\sum_nc_n\geq \sum_kc_{n_k}=\sum_k a_{n_k}=\infty.
  	\end{equation*} 
  	If there are just finite items of $\{c_{n}\}_{n\geq 1}$ satisfying $c_n=a_n$, then there is a $N\geq 1$ such that for any $n\geq N$, we have $c_n=b_n$. Therefore
  	\begin{equation*}
  		\sum_n a_n\wedge b_n =\sum_nc_n\geq \sum_{n\geq N}c_{n}=\sum_{n\geq N}b_{n}=\infty.
  	\end{equation*}
  	\begin{enumerate}
  		\item When $\lim_na_n$ does not exists. Consider 
  	\begin{equation*}
  		a_n=
  		\begin{cases}
  			1/2^n,&\mbox{if $n$ is even},\\
  			1/n,&\mbox{if $n$ is odd}
  		\end{cases}
  		\text{ and }
  		b_n=
  		\begin{cases}
  			1/n,&\mbox{if $n$ is even},\\
  			1/2^n,&\mbox{if $n$ is odd}.
  		\end{cases}
  	\end{equation*}
  	Thus $\sum_na_n=\infty$ and $\sum_n b_n=\infty$. However $a_n\wedge b_n=1/2^n$ yields $\sum_na_n\wedge b_n=\sum_n1/2^n<\infty$.
  		\item When $\lim_na_n=0$. Consider
  		\begin{equation*}
  		a_n=1/n
  		\text{ and }
  		b_n=
  		\begin{cases}
  			n,&\mbox{if $n=2^k,k\geq 1$},\\
  			0,&\mbox{otherwise}.
  		\end{cases}
  	\end{equation*}
  	Thus $\sum_na_n=\infty$ and $\sum_n b_n=\infty$. However
  	\begin{equation*}
  		a_n\wedge b_n=
  		\begin{cases}
  			1/n,&\mbox{if $n=2^k,k\geq 1$},\\
  			0,&\mbox{otherwise}.
  		\end{cases}
  	\end{equation*}
  	 Therefore $\sum_na_n\wedge b_n=\sum_n1/2^n<\infty$.
  	\end{enumerate}
  \end{proof}
  
  \begin{question}
  	Given an increasing function $g:\r_+\to\r$ satisfying $g(x)\to\infty,x\to\infty$ and define 
  	\begin{equation*}
  		f(x):=\int_0^x g(u)du,\quad x\in\r_+.
  	\end{equation*}
  	Then $f$ is increasing, convex and satisfies $f(x)/x\to\infty,x\to\infty$.
  \end{question}
  \begin{proof}
  	The increasing property of $f$ is obvious. For any two nonnegative real numbers $x<y$, we have
  	\begin{equation*}
  		g(x)(y-x)\leq f(y)-f(x)=\int_x^yg(u)du\leq g(y)(y-x).
  	\end{equation*}
  	Thus for any three nonnegative real numbers $x<y<z$, we have
  	\begin{equation*}
  		\frac{f(y)-f(x)}{y-x}\leq g(y)\leq \frac{f(z)-f(y)}{z-y},
  	\end{equation*}
  	which implies $f$ is convex. Note that for any $0\leq y<x$,
  	\begin{equation*}
  		\frac{f(x)}{x}=\frac{f(x)-f(y)}{x}+\frac{f(y)}{x}\geq g(y)\frac{x-y}{x}+\frac{f(y)}{x}.
  	\end{equation*}
  	Let $x\to\infty$, we will obtain 
  	\begin{equation*}
  		\lim_{x\to\infty}\frac{f(x)}{x}\geq \lim_{x\to\infty}\[g(y)\frac{x-y}{x}+\frac{f(y)}{x}\]=g(y).
  	\end{equation*}
  	Let $y\to\infty$ and apply $g(y)\to\infty,y\to\infty$, we will obtain the desired results.
  \end{proof}
  
  \begin{question}
  	Suppose $f:\r\to\r$ is an increasing function. Define
  	\begin{equation*}
  		g(y):=\inf\{x\in\r:f(x)\geq y\},\quad h(y):=\sup\{x\in\r:f(x)<y\}  \quad y\in\r
  	\end{equation*}
  	and
  	\begin{equation*}
  		\eta(y):=\inf\{x\in\r:f(x)> y\},\quad \theta(y):=\sup\{x\in\r:f(x)\leq y\}  \quad y\in\r.
  	\end{equation*}
  	Then we have 
  	\begin{enumerate}
  		\item $g=h$ and $\eta=\theta$.
  		\item $g\leq \eta$ (or $h\leq \theta$).
  		\item $g(y)>x$ implies $y>f(x)$ (or $y\leq f(x)$ implies $g(y)\leq x$). If furthermore $f$ is right-continuous, then $g(y)>x$ iff $y>f(x)$ (or $g(y)\leq x$ iff $y\leq f(x)$).
  		\item $\eta(y)< x$ implies $y< f(x)$ (or $y\geq f(x)$ implies $\eta(y)\geq x$). If furthermore $f$ is left-continuous, then $\eta(y)<x$ iff $y<f(x)$ (or $\eta(y)\geq x$ iff $y\geq f(x)$).
  	\end{enumerate}
  \end{question}
  \begin{proof}
  	\begin{enumerate}
  		\item Note for any $z_1\in \{x\in\r:f(x)\geq y\}$ and $z_2\in \sup\{x\in\r:f(x)<y\}$, we have $z_1\geq z_2$ by the increasing of $f$. Thus $g(y)\geq h(y)$. If $g(y)>h(y)$, then there exists $h(y)<z<g(y)$. Using the definition of $g$ and $h$, we know $h(y)<z\Longrightarrow f(z)\geq y$ and $z<g(y)\Longrightarrow f(x)<y$. And $y\leq f(z)<y$ is a contradiction. In conclusion, $g=h$. A similar argument yields $\eta=\theta$.
  		\item The proof of $g\leq \eta$ (or $h\leq \theta$) is trivial.
  		\item Since $g(y)>x$, by the definition of infimum, for every $z$ such that $f(z)\geq y$, we have $f(x)<y$. On the contrary, if $y>f(x)$, the increasing property of $f$ shows $g(y)\geq x$. However, we claim $g(y)=x$ is impossible. In fact, if $g(y)=x$, then there exists $(z_n^+,n\geq 1)$ such that $f(z_n^+)\geq y,z_n^+>x$ for any $n\geq 1$ and $z_n^+\downarrow x,n\to\infty$. The right-continuity of $f$ yields $f(x)\geq y$, which contradicts $y>f(x)$.
  		\item The definition of $\eta$ and the increasing property of $f$ shows $f(x)>y$. On the contrary, if $y<f(x)$, then $\eta(y)\leq x$. We claim $\eta(y)=x$ is impossible. In fact, if $\eta(y)=x$, then there exists $(z_n^-,n\geq 1)$ such that $f(z_n^-)\leq y,z_n^-<x$ for any $n\geq 1$ and $z_n^-\uparrow x,n\to\infty$. The left-continuity of $f$ yields $f(x)\leq y$, which contradicts $y<f(x)$.
  	\end{enumerate}
  \end{proof}
  
  \begin{question}
  	Suppose $f:\r\to\r$ is a decreasing function. Define
  	\begin{equation*}
  		g(y):=\inf\{x\in\r:f(x)\leq y\},\quad h(y):=\sup\{x\in\r:f(x)>y\}  \quad y\in\r
  	\end{equation*}
  	and
  	\begin{equation*}
  		\eta(y):=\inf\{x\in\r:f(x)< y\},\quad \theta(y):=\sup\{x\in\r:f(x)\geq y\}  \quad y\in\r.
  	\end{equation*}
  	Then we have 
  	\begin{enumerate}
  		\item $g=h$ and $\eta=\theta$.
  		\item $g\leq \eta$ (or $h\leq \theta$).
  		\item $g(y)>x$ implies $y<f(x)$ (or $y\geq f(x)$ implies $g(y)\leq x$). If furthermore $f$ is right-continuous, then $g(y)>x$ iff $y<f(x)$ (or $g(y)\leq x$ iff $y\geq f(x)$).
  		\item $\eta(y)< x$ implies $y> f(x)$ (or $y\leq f(x)$ implies $\eta(y)\geq x$). If furthermore $f$ is left-continuous, then $\eta(y)<x$ iff $y>f(x)$ (or $\eta(y)\geq x$ iff $y\leq f(x)$).
  	\end{enumerate}
  \end{question}
  \begin{proof}
  	Apply a similar argument of the previous QUESTION.
  \end{proof}
  
  \begin{question}
  	Let $c$ be an increasing right-continuous function from $\r_+$ into $\bar{\r}_+$. Define
  	\begin{equation*}
  		a(u):=\inf\{t\in\r_+:c(t)>u\},\quad u\in\r_+,
  	\end{equation*}
  	with the usual convention that $\inf\emptyset=\infty$. Then
  	\begin{enumerate}
  		\item the function $a:\r_+\to \bar{\r}_+$ is also increasing and right-continuous, and that
  		\begin{equation*}
  			c(t):=\inf\{u\in\r_+:a(u)>t\},\quad t\in\r_+.
  		\end{equation*}
  		Thus $a$ and $c$ are right-continuous 'functional inverse' of each other.
  		\item d
  	\end{enumerate}
  \end{question}
  \begin{proof}
\begin{enumerate}
\item The increasing property of $a$ and $c$ are both obvious.
\item Given $u\in\r_+$. By the definition of $a(u)$, there exist $(t_n)_{n\geq 1}$ such that $c(t_n) >u$ and $t_n\downarrow a(u)$ as $n\to\infty$. The right-continuity of $c$ yields $c (a(u)) =\lim_{n\to\infty}c(t_n) \geq u$.
\begin{enumerate}
\item If $c (a(u)) >u$, then we can find $v\in\r_+$ such that $c(a(u)) >v>u$. On one side, the increasing property of $a$ shows $a(v) \geq a(u)$. On another side, $c(a(u)) >v$ shows $a(v) \leq a(u) $. Thus $a (u)=a(v)$. In short, we have shown that, for any $\varepsilon>0$, there exists $\delta=c(a(u))-u>0$. When $u$ satisfies $0<v-u<\delta$, then $a(v)-a (u) =0<\varepsilon$.
\item If $c(a(u)) =u$, there exist $(t_n)_{n\geq 1}$ such that $c(t_n)>u$ and $t_n\downarrow a(u)$ as $n\to\infty$. Thus for any $\varepsilon>0$, there exists $N\in\z_+$ such that $0<t_n-a (u) <\varepsilon$ for any $n\geq N$. For any $u<v<c(t_n)$, the definition of $a$ yields $a(v)\leq t_n$. Thus $0<a (v)-a(u) \leq t_n-a (u) <\varepsilon$. In short, we have shown that, for any $\varepsilon>0$, there exists $\delta=c(t_n)-u>0$. When $v$ satisfies $0<v-u<\delta$, then $0<a(v)-a(u)<\varepsilon$.
\end {enumerate}
The above arguments together show $a$ is right-continuous.
\item
Write $c' (t) :=\inf\{u\in\r_+:a(u)>t\}$. Then what we need to show is $c=c'$. For any $u\in\r_+$ such that $a(u) >t$, the definition of $a$ and the increasing property of $c$ yield $c(t) \leq u$. Thus $c(t) \leq c'(t)$. On the contrary, applying the definition of $c'$, we have $a(c'(t) -\varepsilon)\leq t$ for any $\varepsilon>0$. The definition of $a$ again shows $c(t) \geq c'(t)-\varepsilon$. Let $\varepsilon\to 0$ to obtain $c(t) \geq c'(t)$.
Therefore $c=c'$.
\end {enumerate}
\end{proof}

\begin{question}
	We say a function $f:\r\to\r$ is convex if for any $x,y\in\r$ and $\theta\in[0,1]$, we have
	\begin{equation*}
		f(\theta x+(1-\theta)y)\leq \theta f(x)+(1-\theta)f(y).\tag{$*$}
	\end{equation*}
	Prove the following statements are all equivalent to the convexity:
	\begin{enumerate}
		\item For any $x,y\in\r$, we have
		\begin{equation*}
			f\(\frac{x+y}{2} \)\leq \frac{f(x)+f(y)}{2}.
		\end{equation*}
		\item For any $n\in\z_+$ and $x_1,\cdots,x_n\in\r$, we have
		\begin{equation*}
			f\(\frac{1}{n}\sum_{k=1}^nx_k \)\leq \frac{1}{n}\sum_{k=1}^nf(x_k).\tag{$*$}
		\end{equation*}
		\item For any $x_1<x_2<x_3$, we have
		\begin{equation*}
			\frac{f(x_2)-f(x_1)}{x_2-x_1}\leq \frac{f(x_3)-f(x_2)}{x_3-x_2}.
		\end{equation*}
		
	\end{enumerate}
	Furthermore, if $f\in C^1(\r)$, the convexity is also equivalent to
	\begin{enumerate}
		\item[4.]For any $x,y\in\r$, we have
		\begin{equation*}
			f(y)\geq f(x)+f'(x)(y-x).
		\end{equation*}
		\item[5.]$f'(x)$ is increasing on $\r$.
	\end{enumerate}
	If $f\in C^2(\r)$, the convexity is also equivalent to
	\begin{enumerate}
		\item[6.] $f''(x)>0$ for any $x\in\r$.
	\end{enumerate}
\end{question}
\begin{proof}
	\begin{enumerate}
		\item $1\Leftrightarrow 2$. We just need to show $1\Rightarrow2$. We use induction on $n$ to show 
		\begin{equation*}
			f\(\frac{1}{2^n}\sum_{k=1}^{2^n}x_k \)\leq \frac{1}{2^n}\sum_{k=1}^{2^n}f(x_k),\quad \forall n\in\z_+.
		\end{equation*}
		When $n=1$, it is obviously true. Suppose the inequality holds for $n=\ell$. When $n=\ell+1$, using the hypothesis assumption, we have
		\begin{equation*}
			\begin{split}
				f\(\frac{1}{2^{\ell+1}}\sum_{k=1}^{2^{\ell+1}}x_k \)&=f\( \frac{\frac{1}{2^\ell} \sum_{k=1}^{2^\ell}x_k+\frac{1}{2^\ell} \sum_{k=2^\ell+1}^{2^{\ell+1}}x_k}{2}\)\\
				&\leq \frac{1}{2}\[f\(\frac{1}{2^\ell} \sum_{k=1}^{2^\ell}x_k \)+f\(\frac{1}{2^\ell} \sum_{k=2^\ell+1}^{2^{\ell+1}}x_k\)\]\\
				&\leq \frac{1}{2^{\ell+1}}\sum_{k=1}^{2^{\ell+1}}f(x_k).
			\end{split}
		\end{equation*}
		Then we show that, if $(*)$ holds for $n+1$, then $(*)$ also holds for $n$. Let $t=\frac{1}{n}\sum_{k=1}^{n}x_k$, then $nt=\sum_{k=1}^{n}x_k\Longrightarrow t=\frac{1}{n+1}(\sum_{k=1}^{n}x_k+t)$. Thus
		\begin{equation*}
			f(t)=f\(\frac{1}{n+1}\(\sum_{k=1}^{n}x_k+t\)\)\leq \frac{1}{n+1}\[\sum_{k=1}^nf(x_k)+f(t)\].
		\end{equation*}
		This is equivalent to
		\begin{equation*}
			f\(\frac{1}{n}\sum_{k=1}^{n}x_k\)=f(t)\leq \frac{1}{n}\sum_{k=1}^nf(x_k).
		\end{equation*}
		Finally, the above arguments show $(*)$ holds for any $n\in\z_+$.
		\item $(*)\Rightarrow1$. We just need to let $\theta=1/2$. 
		\item $2\Rightarrow(*)$. First we assume $\lambda=n/m\in\q$, where $n,m\in\z_+,n\leq m$ and $(n,m)=1$. Then 
		\begin{equation*}
			\begin{split}
				f\(\theta x+(1-\theta)y\)&=f\(\frac{n}{m} x+\frac{m-n}{m}y \)\\
				&=f\(\frac{nx+(m-n)y}{m}\)\\
				&\overset{2}{\leq} \frac{1}{m}\[nf(x)+(m-n)y\]\\
				&=\theta f(x)+(1-\theta)f(y).
			\end{split}
		\end{equation*}
		Since the convex has a left derivative and a right derivative, it must be left continuous and right continuous, i.e., continuous. Next, for any $\theta\in[0,1]$, we can find a rational sequence $(\theta_n)_{n\geq 1}\subset[0,1]\cap\q$ satisfying $\theta_n\to\theta$. So
		\begin{equation*}
			\begin{split}
				f\(\theta x+(1-\theta)y\)&=\lim_{n\to\infty}f\(\theta_n x+(1-\theta_n)y\)\\
				&\leq\lim_{n\to\infty}\theta_n f(x)+(1-\theta_n)f(y) \\
				&=\theta f(x)+(1-\theta)f(y).
			\end{split}
		\end{equation*}
		\item $(*)\Leftrightarrow 3$. For $\Rightarrow$, let $\theta=\frac{x_3-x_2}{x_3-x_1}$, then $x_2=\theta x_1+(1-\theta)x_3$, then 
		\begin{equation*}
			\begin{split}
				f(x_2)&=f(\theta x_1+(1-\theta)x_3)\leq \theta f(x_1)+(1-\theta)f(x_3)\\
				&=\frac{x_3-x_2}{x_3-x_1}f(x_1)+\frac{x_2-x_1}{x_3-x_1}f(x_3).
			\end{split}
		\end{equation*}
		This is equivalent to either of the following inequality:
		\begin{equation*}
			\begin{split}
				[(x_3-x_2)+(x_2-x_1)]f(x_2)\leq (x_3-x_2)f(x_1)+(x_2-x_1)f(x_3)\\
					\Longleftrightarrow \frac{f(x_2)-f(x_1)}{x_2-x_1}\leq \frac{f(x_3)-f(x_2)}{x_3-x_2}
			\end{split}
		\end{equation*}
		For $\Leftarrow$, define $x_1:=x,x_2:=\theta x+(1-\theta)y, x_3:=y$. Then $\theta=\frac{x_3-x_2}{x_3-x_1}$. We can reverse the whole process above to obtain
		\begin{equation*}
			f(x_2)\leq \frac{x_3-x_2}{x_3-x_1}f(x_1)+\frac{x_2-x_1}{x_3-x_1}f(x_3).
		\end{equation*}
		However, this is equivalent to
		\begin{equation*}
			f(\theta x+(1-\theta)y)=\theta f(x_1)+(1-\theta)f(x_3)=\theta f(x)+(1-\theta)f(y).
		\end{equation*}
		\item $3\Rightarrow 4$. Without loss of generality, we assume $y<x$. Arbitrarily pick $y<z<x$, then we have
		\begin{equation*}
			\frac{f(x)-f(y)}{x-y}\leq \frac{f(x)-f(z)}{x-z}.
		\end{equation*}
		Let $z\to x-$, we have
		\begin{equation*}
			\begin{split}
				\frac{f(x)-f(y)}{x-y}\leq f'(x)&\Longleftrightarrow f(x)\leq f(y)+f'(x)(x-y)\\
				&\Longleftrightarrow f(y)\geq f(x)+f'(x)(y-x).
			\end{split}
		\end{equation*}
		\item $4\Rightarrow (*)$. Let $z=\theta x+(1-\theta)y$, then we have
		\begin{equation*}
			\left\{
			\begin{split}
				f(x)\geq f(z)+f'(z)(x-z)\\
				f(y)\geq f(z)+f'(z)(y-z).
			\end{split}\right.
		\end{equation*}
		Plug $z=\theta x+(1-\theta)y$ to the above inequalities, then this is equivalent to 
		\begin{equation*}
			\left\{
			\begin{split}
				f(x)&\geq f(z)+(1-\theta)f'(z)(x-y)\\
				f(y)&\geq f(z)+\theta f'(z)(y-x).
			\end{split}\right.
		\end{equation*}
		This is again equivalent to
		\begin{equation*}
			\left\{
			\begin{split}
				\theta f(x)&\geq \theta f(z)+\theta(1-\theta)f'(z)(x-y)\\
				(1-\theta)f(y)&\geq (1-\theta)f(z)+\theta(1-\theta) f'(z)(y-x).
			\end{split}\right.
		\end{equation*}
		Add these two formulas together to obtain
		\begin{equation*}
			\theta f(x)+(1-\theta)f(y)\geq f(z)=f(\theta x+(1-\theta)y).
		\end{equation*}
		\item $3\Leftrightarrow 5$. For $\Rightarrow$, for any $x,y\in\r$ and $x<y$, we have
		\begin{equation*}
			\frac{f(x)-f(x-h)}{h}\leq \frac{f(x)-f(y)}{x-y}\leq \frac{f(y+h)-f(y)}{h}.
		\end{equation*}
		Let $h\to0$, we have 
		\begin{equation*}
			f'(x)\leq \frac{f(x)-f(y)}{x-y}\leq f'(y).
		\end{equation*}
		For $\Leftarrow$, apply the mean value theorem of differential, we have
		\begin{equation*}
			\frac{f(x_2)-f(x_1)}{x_2-x_1}=f'(\xi),\quad \frac{f(x_3)-f(x_2)}{x_3-x_2}=f'(\eta),
		\end{equation*}
		where $x_1\leq \xi\leq x_2\leq \eta\leq x_3$. The increasing property of $f'$ yields the desired result.
		\item $5\Leftrightarrow6$. We just need to note that, if $f\in C^2(\r)$, then $f'$ is increasing iff $(f')'=f''\geq 0$.
	\end{enumerate}
\end{proof}

\begin{question}
	A function $f:\r^n\to\r$ is called semi-upper continuous at $x=x_0$ if $\varlimsup_{x\to x_0}f(x)\leq f(x_0)$. Prove the following three definitions are equivalent: 
	\begin{enumerate}
		\item $f$ is semi-upper continuous at $x=x_0$.
		\item For any $\varepsilon>0$, there exists $\delta>0$, such that $f(x)<f(x_0)+\varepsilon$ for any $x\in B(x_0,\delta)$.
		\item The upper contour set $U_w(x_0):=\{x\in\r^n:f(x)\geq f(x_0)\}$ is closed.
		\item The strict lower contour set $L_s(x_0):=\{x\in\r^n:f(x)< f(x_0)\}$ is open.
	\end{enumerate}
\end{question}
\begin{proof}
	\begin{enumerate}
		\item $1\Leftrightarrow 2$. First, we show $1\Rightarrow 2$. Recall the definition of the upper limit:
		\begin{equation*}
			\varlimsup_{x\to x_0}f(x)=\lim_{r\to0}\sup_{x\in B(x_0,r)}f(x)=\inf _{r>0}\sup_{x\in B(x_0,r)}f(x).
		\end{equation*}
		Thus for any $\varepsilon>0$, there exists $\delta>0$, for any $0<r<\delta$, we have
		\begin{equation*}
			\sup_{x\in B(x_0,r)}f(x)<\varlimsup_{x\to x_0}f(x)+\varepsilon\leq f(x_0)+\varepsilon.
		\end{equation*}
		Thus
		\begin{equation*}
			f(x)\leq f(x_0)+\varepsilon,\quad \forall x\in B(x_0,r)
		\end{equation*}
		As for $2\Rightarrow 1$, since for any $\varepsilon>0$, there exists $\delta>0$, such that $f(x)<f(x_0)+\varepsilon,\forall x\in B(x_0,\delta)$, we have
		\begin{equation*}
			\sup_{x\in B(x_0,r)}f(x)\leq f(x_0)+\varepsilon.
		\end{equation*}
		Furthermore, 
		\begin{equation*}
			\varlimsup_{x\to x_0}f(x)=\inf _{r>0}\sup_{x\in B(x_0,r)}f(x)\leq f(x_0)+\varepsilon.
		\end{equation*}
		Letting $\varepsilon\to0+$ will yield the result.
		\item $3\Leftrightarrow4$. This is obvious.
		\item $2\Leftrightarrow 4$. First, we show $2\Rightarrow 4$. For any $x\in L_s(x_0)$, we have $f(x)<f(x_0)$. Let $\varepsilon=[f(x_0)-f(x)]/2$. There exists $\delta>0$ such that 
		\begin{equation*}
			f(y)<f(x)+\varepsilon<f(x_0),\quad \forall x\in B(x,\delta).
		\end{equation*}
		This means $B(x,\delta)\subset L_s(x_0)$ and thus $L_s(x_0)$ is open. As for $4\Rightarrow 2$, consider 
	\end{enumerate}
\end{proof}

\begin{question}
	Given a random variable $X$ on the probability space $(\Omega,\g,\p)$ and valued on the measurable space $(E,\e)$. Suppose $(F,\f)$ is another measurable space and $f:E\times F\to \r$ is measurable, then
	\begin{equation*}
		\sup_{y\in F}\p [f(X,y)]\leq\p \[\sup_{y\in F}f(X,y)\].
	\end{equation*}
\end{question}
\begin{proof}
	On the one hand, for any $y\in F$,
	\begin{equation*}
		\p [f(X,y)]\leq \p \[\sup_{y\in F}f(X,y)\].
	\end{equation*}
	Taking $\sup_{y\in F}\cdot$ on two sides of the previous formula yields the inequality. On the other hand, consider $\Phi(x)=$. We have 
	\begin{equation*}
		\p[f(X,y)]=1/y-1/y=0,\quad \forall y\neq 0.
	\end{equation*}
	Therefore $\sup_{y\in F}\p [f(X,y)]=0$. However, 
	\begin{equation*}
		\sup_{y\in F}f(X,y)=\sup_{y\in F}X/y=+\infty\Longrightarrow \p \[\sup_{y\in F}f(X,y)\]=+\infty.
	\end{equation*}
	The proof is completed.
\end{proof}

\begin{question}
	Given an increasing function $\Phi:\r\to \r$ and $A\subset \r$. Then
	\begin{equation*}
		\sup_{x\in A}\Phi(x) \leq\Phi\(\sup_{x\in A}x\) .
	\end{equation*}
\end{question}
\begin{proof}
	On the one hand, for any $x\in A$,
	\begin{equation*}
		\Phi (x)\leq \Phi\(\sup_{x\in A}x\).
	\end{equation*}
	Taking $\sup_{x\in A}\cdot$ on two sides of the previous formula yields the inequality. On the other hand, consider an increasing function $\Phi(x)=\mathbbm{1}_{[1,\infty)}(x)$ and $A=(-\infty,1)$. Then 
	\begin{equation*}
		\sup_{x<1}\mathbbm{1}_{[1,\infty)}(x)=0<1=\mathbbm{1}_{[1,\infty)}(1)=\mathbbm{1}_{[1,\infty)}\(\sup_{x<1}x \).
	\end{equation*}
	The proof is completed.  
\end{proof}

\begin{question}
	Given a differentiable function $f:\r^n\to\r$ satisfying $|f(x)-f(y)|\leq L||x-y||$, namely, $L$-Lipschitz with respect to Euclidian norm. Then $||\nabla f||\leq L$.
\end{question}
\begin{proof}
	Since $f$ is differentiable, we have for any $y\in \r^n$,
	\begin{equation*}
		|f(x)-f(y)|=|\nabla f(y)\cdot(x-y)+o(||x-y||)|\leq L||x-y||.
	\end{equation*}
	Divide both sides by $||x-y||$ to obtain
	\begin{equation*}
		\frac{|\nabla f(y) \cdot(x-y)+o(||x-y||)|}{||x-y||}\leq L.
	\end{equation*}
	Replacing $x$ with $x_n=\nabla f(y) /n+y$ and letting $n\to\infty$, we will have $x_n\to y$ and
	\begin{equation*}
		\left|\nabla f(y) \cdot \frac{\nabla f(y) }{||\nabla f(y) ||}+o(1) \right|\to ||\nabla f(y) || \leq L.
	\end{equation*}
\end{proof}

\begin{question}
	For any c\'adl\'ag function $f\in D[0,1]$ and any $\varepsilon>0$, there must exist finite points $0=x_0<x_1<\cdots<x_{n-1}<x_n=1$ such that $\sup_{x,y\in[x_i,x_{i+1})}|f(x)-f(y)|<\varepsilon$ for any $ 0\leq i\leq n-1$ (we say $f$ is finitely continuous in $[0,1]$).
\end{question}
\begin{proof}
	If not, we divide $[0,1]$ into $[0,1/2)\cup [1/2,1]$. Then $f$ is not finitely continuous in either $[0,1/2)$ or $[1/2,1]$, say $A_1:=[0,1/2)$. We can divide $[0,1/2)$ into $[0,1/4)\cup(1/4,1/2)$ again and $f$ is not finitely continuous in either $[0,1/4)$ or $[1/4,1/2)$, the corresponding interval is denoted by $A_2$. Applying induction, we can get a decreasing sequence $\{A_n\}_{n\in\z}$ and $\mathrm{diam}(A_n)\to 0$. The {\itshape Nested Interval Theorem} shows there exists a unique $x^*\in \cap_n A_n$. Without loss of generality, we assume $x^*\in(0,1)$. Since $f$ is c\'adl\'ag at $x^*$, for the given $\varepsilon>0$, there exists $\delta>0$, such that $|f(x)-f(x^*)|<\varepsilon/2,\forall x-x^*<2\delta$ and $|f(x)-f(x^*-)|<\varepsilon/2,\forall x^*-x<2\delta$. Therefore, letting $x^*-\delta =x_0<x_1=x^*<x+\delta$, we have
	\begin{equation*}
		\sup_{x,y\in[x_i,x_{i+1})}|f(x)-f(y)|<\varepsilon,\quad i=0,1.
	\end{equation*}
	So $f$ is finitely continuous in $[x^*-\delta,x^*+\delta]$. However, when $n$ big enough, we can select $A_n$ such that $A_n\subset [x^*-\delta,x^*+\delta]$. The definition of $A_n$ shows $f$ is not finitely continuous in $A_n$, which is a contradiction.
\end{proof}

\begin{question}
  	Assume $f\in C[0,1]$ and $f\geq 0$. What are the limits of 
  	\begin{equation}\label{1}
  		\[\int_0^1 f(x)^\varepsilon dx\]^{1/\varepsilon}
  	\end{equation}
  	when $\varepsilon\to 0,\varepsilon\to 1$ and $\varepsilon\to \infty$?
  \end{question}
  \begin{proof}
  	\begin{enumerate}[(1)]
  		\item Assume $\varepsilon\to 0$. Now 
  		\begin{equation*}
  			\int_0^1 f(x)^\varepsilon dx\to 1,\quad 1/\varepsilon\to \infty.
  		\end{equation*}
  		This is a $1^\infty$ type indeterminate form. Take $\log(\cdot)$ and apply the L'H\^ospital rule to obtain
  		\begin{equation*}
  			\lim_{\varepsilon\to 0}\frac{1}{\varepsilon}\log\(\int_0^1 f(x)^\varepsilon dx\)=\lim_{\varepsilon\to 0}\frac{\int_0^1f(x)^{\varepsilon}\log f(x)dx}{\int_0^1 f(x)^\varepsilon dx}=\int_0^1\log f(x)dx.
  		\end{equation*}
  		So,
  		\begin{equation*}
  			\[\int_0^1 f(x)^\varepsilon dx\]^{1/\varepsilon}\to \exp\left\{\int_0^1\log f(x)dx\right\},\quad \varepsilon\to0.
  		\end{equation*}
  		This is the Cobb-Douglas case.
  		\item Assume $\varepsilon\to1$. Obviously,
  		\begin{equation*}
  			\[\int_0^1 f(x)^\varepsilon dx\]^{1/\varepsilon}\to \int_0^1 f(x)dx,\quad \varepsilon\to1.
  		\end{equation*}
  		This is a perfectly substitutable case.
  		\item Assume $\varepsilon\to\infty$. Since $f\in C[0,1]$, we can define
  		\begin{equation*}
  			\sup_{y\in[0,1]}f(y)=\max_{y\in[0,1]}f(y)=:\bar{f} 
  		\end{equation*}
  		Moreover, write $z=\arg\max_{y\in[0,1]}f(y)$. For any small enough $\gamma>0$, we can find $\delta(\gamma)>0$ such that 
  		\begin{equation}\label{2}
  			f(x)>\bar{f}-\gamma>0,\quad \forall x\in[z-\delta(\gamma),z+\delta(\gamma)],
  		\end{equation} 
  		On one side, obviously, 
  		\begin{equation*}
  			\[\int_0^1 f(x)^\varepsilon dx\]^{1/\varepsilon}\leq \[\int_0^1 \bar{f}^\varepsilon dx\]^{1/\varepsilon}=\bar{f}.
  		\end{equation*}
  		Letting $\varepsilon\to \infty$ yields
  		\begin{equation}\label{3}
  			\varlimsup_{\varepsilon\to\infty}\[\int_0^1 f(x)^\varepsilon dx\]^{1/\varepsilon}\leq \bar{f}.
  		\end{equation}
  		On the other side, using \eqref{2},
  		\begin{equation*}
  			\[\int_0^1 f(x)^\varepsilon dx\]^{1/\varepsilon}\geq \[\int_{z-\delta(\gamma)}^{z+\delta(\gamma)} (\bar{f}-\gamma)^\varepsilon dx\]^{1/\varepsilon} =(\bar{f}-\gamma) [2\delta(\gamma)]^{1/\varepsilon}.
  		\end{equation*}
  		Let $\varepsilon\to\infty$ on both sides, we will get
  		\begin{equation*}
  			\varliminf_{\varepsilon\to\infty}\[\int_0^1 f(x)^\varepsilon dx\]^{1/\varepsilon}\geq \bar{f}-\gamma.
  		\end{equation*}
  		Let $\gamma\to 0$, we can get 
  		\begin{equation}\label{4}
  			\varliminf_{\varepsilon\to\infty}\[\int_0^1 f(x)^\varepsilon dx\]^{1/\varepsilon}\geq \bar{f}.
  		\end{equation}
		Combine \eqref{3} and \eqref{4} to obtain
		\begin{equation*}
			\bar{f}\leq \varliminf_{\varepsilon\to\infty}\[\int_0^1 f(x)^\varepsilon dx\]^{1/\varepsilon}\leq \varlimsup_{\varepsilon\to\infty}\[\int_0^1 f(x)^\varepsilon dx\]^{1/\varepsilon}\leq \bar{f}.
		\end{equation*}
		This implies 
		\begin{equation*}
			\[\int_0^1 f(x)^\varepsilon dx\]^{1/\varepsilon}\to \bar{f},\quad \varepsilon\to\infty.
		\end{equation*}
		This is the Leontief case.
  	\end{enumerate}
  \end{proof}
  
  \begin{question}
  	If $f:E\to \r^d$ is a concave function, where $E$ is a convex linear space, then the local maximum is the global maximum.
  \end{question}
  \begin{proof}
  	If $x\in E$ is the local maximum, then there exists a neighborhood of $x$, say $U_x$, such that $f(x)\geq f(y),\forall y\in U_x$. If $x$ is not a global maximum, there exists $x'\in E$ such that $f(x')>f(x)$. For any $\theta\in (0,1)$ small enough, we have $\theta x'+(1-\theta)x\in U_x$. And, the definition of concavity yields
  	\begin{equation*}
  		f(\theta x'+(1-\theta)x)\geq \theta f(x')+(1-\theta)f(x)>f(x).
  	\end{equation*}
  	This contradicts the fact that $x$ is a local maximum.
  \end{proof}
  
  \begin{question}
Let $f:\mathbb{R}^n\to(-\infty,+\infty]$ be an extended real-valued function. Recall the \emph{epigraph} and \emph{sublevel sets}:
\begin{equation*}
\mathrm{epi}\,f := \{(x,t)\in\mathbb{R}^{n+1}: f(x)\le t\}, 
\qquad
L_\alpha := \{x\in\mathbb{R}^n : f(x)\le \alpha\}\ (\alpha\in\mathbb{R}).
\end{equation*}
Prove the following equivalences.
\begin{enumerate}
\item[(A)] $f$ is lower semicontinuous (l.s.c.), i.e., $f(x)\le \liminf_{k\to\infty} f(x_k)$ for every sequence $x_k\to x$, if and only if $\mathrm{epi}\,f$ is closed in $\mathbb{R}^{n+1}$.
\item[(B)] $f$ is lower semicontinuous if and only if every sublevel set $L_\alpha$ is closed in $\mathbb{R}^n$.
\end{enumerate}
\end{question}

\begin{proof}
We work in the Euclidean topology, so sequential arguments suffice.

\medskip
\noindent\textbf{(A) ``$\Rightarrow$'':} Assume $f$ is l.s.c. Let $(x_k,t_k)\in \mathrm{epi}\,f$ with $(x_k,t_k)\to (x,t)$. Since $f(x_k)\le t_k$ for all $k$,
\begin{equation*}
f(x)\ \le\ \liminf_{k\to\infty} f(x_k)\ \le\ \liminf_{k\to\infty} t_k\ =\ t,
\end{equation*}
hence $(x,t)\in \mathrm{epi}\,f$. Therefore $\mathrm{epi}\,f$ is closed.

\medskip
\noindent\textbf{(A) ``$\Leftarrow$'':} Assume $\mathrm{epi}\,f$ is closed. Let $x_k\to x$ and set $s:=\liminf_{k\to\infty} f(x_k)$. Choose a subsequence $k_j$ with $f(x_{k_j})\to s'$ where $s'=s$ (or $s'=+\infty$). Then
\begin{equation*}
(x_{k_j},\,f(x_{k_j}))\in \mathrm{epi}\,f,\qquad (x_{k_j},\,f(x_{k_j}))\to (x,\,s').
\end{equation*}
By closedness, $(x,s')\in\mathrm{epi}\,f$, hence $f(x)\le s'=\liminf_{k\to\infty} f(x_k)$. Thus $f$ is l.s.c.

\medskip
\noindent\textbf{(B) ``$\Rightarrow$'':} Assume $f$ is l.s.c. Fix $\alpha\in\mathbb{R}$ and take $x_k\in L_\alpha$ with $x_k\to x$. Then
\begin{equation*}
f(x)\ \le\ \liminf_{k\to\infty} f(x_k)\ \le\ \alpha,
\end{equation*}
so $x\in L_\alpha$. Hence each $L_\alpha$ is closed.

\medskip
\noindent\textbf{(B) ``$\Leftarrow$'':} Assume every $L_\alpha$ is closed. Let $x_k\to x$ and set $s:=\liminf_{k\to\infty} f(x_k)$. If $f(x)>s$, choose $\alpha$ with $s<\alpha<f(x)$. Then $f(x_k)\le\alpha$ holds for infinitely many $k$, so a subsequence $x_{k_j}\in L_\alpha$ converges to $x$. Closedness of $L_\alpha$ implies $x\in L_\alpha$, i.e., $f(x)\le\alpha<f(x)$, a contradiction. Thus $f(x)\le \liminf_{k\to\infty} f(x_k)$, and $f$ is l.s.c.
\end{proof}

\begin{question}[Limsup representation]
For any function $f : X \to \mathbb{R} \cup \{-\infty\}$ and any point
$x_0 \in X$, define 
\begin{equation*}
    \varlimsup_{x\to x_0}f(x):=\sup\left\{\varlimsup_{n\to\infty} f(x_n):x_n\to x_0,x_n\neq x_0 \right\}
\end{equation*}
and
\begin{equation*}
M(\delta)
:=
\sup\{ f(x) : 0 < d(x,x_0) < \delta\},
\qquad \delta > 0.
\end{equation*}
Then
\begin{equation*}
\varlimsup_{x \to x_0} f(x)
=
\inf_{\delta > 0} M(\delta).
\end{equation*}
\end{question}

\begin{proof}
The family $\{M(\delta)\}_{\delta>0}$ is nonincreasing in $\delta$ by
construction, hence the infimum $\inf_{\delta>0} M(\delta)$ coincides
with the limit $\lim_{\delta \downarrow 0} M(\delta)$.

For the inequality
\begin{equation*}
\varlimsup_{x \to x_0} f(x)
\leq
\inf_{\delta > 0} M(\delta),
\end{equation*}
let $(x_n)$ be any sequence with $x_n \to x_0$.
Given $\delta > 0$, we have $d(x_n,x_0) < \delta$ for all sufficiently
large $n$, hence
\begin{equation*}
f(x_n) \leq M(\delta)
\quad\text{for all large } n.
\end{equation*}
It follows that
\begin{equation*}
\varlimsup_{n\to\infty} f(x_n) \leq M(\delta)
\quad\text{for every } \delta>0,
\end{equation*}
and taking the infimum over $\delta>0$ yields
\begin{equation*}
\varlimsup_{n\to\infty} f(x_n)
\leq
\inf_{\delta>0} M(\delta).
\end{equation*}
Since the left-hand side is an arbitrary subsequential limsup along
$x_n \to x_0$, we obtain
\begin{equation*}
\varlimsup_{x \to x_0} f(x)
\leq
\inf_{\delta>0} M(\delta).
\end{equation*}

For the reverse inequality, fix $n \in \mathbb{N}$ and put
$\delta_n := 1/n$.
If the set $\{x : 0<d(x,x_0)<\delta_n\}$ is nonempty, choose
$x_n$ such that
\begin{equation*}
0 < d(x_n,x_0) < \delta_n
\quad\text{and}\quad
f(x_n) \geq M(\delta_n) - \frac{1}{n}.
\end{equation*}
If the set is empty, one can ignore this $n$ or choose any point
$x_n \neq x_0$; this does not affect the argument.
By construction, $x_n \to x_0$ as $n\to\infty$ and
\begin{equation*}
\varlimsup_{n\to\infty} f(x_n)
\geq
\varlimsup_{n\to\infty} \bigl(M(\delta_n) - \tfrac{1}{n}\bigr)
=
\lim_{n\to\infty} M(\delta_n)
=
\inf_{\delta>0} M(\delta).
\end{equation*}
Thus
\begin{equation*}
\varlimsup_{x \to x_0} f(x)
\geq
\inf_{\delta>0} M(\delta).
\end{equation*}
Combining the two inequalities proves the claim.
\end{proof}




\begin{theorem}[Equivalent characterizations of upper semicontinuity]
\label{thm:usc-equivalences}
Let $(X,d)$ be a metric space and let $f : X \to \mathbb{R} \cup \{-\infty\}$.
Let $x_0 \in X$. We say that $f$ is \emph{upper semicontinuous at $x_0$} if
for every $\varepsilon > 0$ there exists $\delta > 0$ such that
\begin{equation*}
d(x,x_0) < \delta
\quad\Longrightarrow\quad
f(x) \leq f(x_0) + \varepsilon.
\end{equation*}
We say that $f$ is \emph{upper semicontinuous} on $X$ if it is upper
semicontinuous at every point $x_0 \in X$.
The following assertions are equivalent:
\begin{itemize}
  \item[(1)] $f$ is upper semicontinuous on $X$;
  \item[(2)] for every $x_0 \in X$,
  \begin{equation*}
  \varlimsup_{x\to x_0} f(x) \leq f(x_0);
  \end{equation*}
  \item[(3)] for every $a \in \mathbb{R}$, the set
  \begin{equation*}
  U_a := \{x \in X : f(x) < a\}
  \end{equation*}
  is open in $X$;
  \item[(4)] the hypograph
  \begin{equation*}
  \operatorname{hypo}(f)
  :=
  \{(x,t) \in X \times \mathbb{R} : t \leq f(x)\}
  \end{equation*}
  is a closed subset of the product metric space $X \times \mathbb{R}$.
\end{itemize}
\end{theorem}

\begin{proof}
\emph{(1) $\Rightarrow$ (2).}
Let $(x_n)$ be any sequence with $x_n \to x_0$.
By upper semicontinuity at $x_0$, for every $\varepsilon>0$ there
exists $\delta>0$ such that
\begin{equation*}
d(x,x_0) < \delta
\quad\Longrightarrow\quad
f(x) \leq f(x_0) + \varepsilon.
\end{equation*}
For all sufficiently large $n$ we have $d(x_n,x_0)<\delta$, hence
\begin{equation*}
f(x_n) \leq f(x_0)+\varepsilon
\quad\text{for all large } n.
\end{equation*}
Therefore
\begin{equation*}
\varlimsup_{n\to\infty} f(x_n)
\leq
f(x_0)+\varepsilon.
\end{equation*}
Since $\varepsilon>0$ is arbitrary, we obtain
\begin{equation*}
\varlimsup_{n\to\infty} f(x_n)
\leq f(x_0)
\end{equation*}
for every sequence $x_n \to x_0$.
By the definition of $\varlimsup_{x\to x_0} f(x)$ as the supremum of
subsequential limsups along convergent sequences, this implies
\begin{equation*}
\varlimsup_{x\to x_0} f(x) \leq f(x_0).
\end{equation*}

\medskip

\emph{(2) $\Rightarrow$ (1).}
Suppose (ii) holds and assume by contradiction that $f$ is not upper
semicontinuous at $x_0$.
Then there exists $\varepsilon_0>0$ such that for every $\delta>0$
we can find some $x$ with
\begin{equation*}
d(x,x_0) < \delta
\quad\text{and}\quad
f(x) > f(x_0) + \varepsilon_0.
\end{equation*}
In particular, we can construct a sequence $(x_n)$ with
\begin{equation*}
d(x_n,x_0) < \frac{1}{n}
\quad\text{and}\quad
f(x_n) > f(x_0)+\varepsilon_0
\quad\text{for all } n.
\end{equation*}
Then $x_n\to x_0$ and
\begin{equation*}
\varlimsup_{n\to\infty} f(x_n)
\geq
f(x_0)+\varepsilon_0,
\end{equation*}
which contradicts (ii).
Hence $f$ must be upper semicontinuous at $x_0$. Then we can apply this argument to every point $x_0\in X$.

\medskip

\emph{(1) $\Rightarrow$ (3).}
Fix $a\in\mathbb{R}$ and let
\begin{equation*}
U_a := \{x\in X : f(x) < a\}.
\end{equation*}
We show that $U_a$ is open.
Let $x_0 \in U_a$.
Then $f(x_0) < a$. Put
\begin{equation*}
\varepsilon := \frac{a - f(x_0)}{2} > 0.
\end{equation*}
By upper semicontinuity at $x_0$, there exists $\delta>0$ such that
\begin{equation*}
d(x,x_0) < \delta
\quad\Longrightarrow\quad
f(x) \leq f(x_0) + \varepsilon < a.
\end{equation*}
Hence
\begin{equation*}
B(x_0,\delta) \subset U_a.
\end{equation*}
Thus $x_0$ is an interior point of $U_a$ and $U_a$ is open.

\medskip

\emph{(3) $\Rightarrow$ (1).}
Assume that for every $a\in\mathbb{R}$, $U_a := \{f<a\}$ is open.
We prove upper semicontinuity at an arbitrary point $x_0\in X$.

Let $x_0 \in X$ and let $(x_n)$ be any sequence with $x_n \to x_0$.
We show that
\begin{equation*}
\varlimsup_{n\to\infty} f(x_n) \leq f(x_0).
\end{equation*}
Assume by contradiction that
\begin{equation*}
\varlimsup_{n\to\infty} f(x_n) > f(x_0).
\end{equation*}
Then there exists a subsequence $(x_{n_k})$ and a number
$y > f(x_0)$ such that
\begin{equation*}
f(x_{n_k}) \to y \quad\text{as } k\to\infty.
\end{equation*}
Set
\begin{equation*}
a := \frac{f(x_0)+y}{2},
\end{equation*}
so that $f(x_0) < a < y$.
Since $f(x_0) < a$, we have $x_0 \in U_a$.
By assumption $U_a$ is open, so there exists $\delta>0$ such that
\begin{equation*}
B(x_0,\delta) \subset U_a.
\end{equation*}
Since $x_{n_k} \to x_0$, we have $x_{n_k} \in B(x_0,\delta)$ for all
sufficiently large $k$, hence
\begin{equation*}
f(x_{n_k}) < a
\quad\text{for all large } k.
\end{equation*}
This contradicts $f(x_{n_k}) \to y > a$.
Therefore
\begin{equation*}
\varlimsup_{n\to\infty} f(x_n) \leq f(x_0)
\end{equation*}
for every sequence $x_n \to x_0$. Take $\sup$ with respect to all the sequences $(x_n)_{n\geq 1}$ satisfying $x_n\to x_0,x_n\neq x_0$.
By the equivalence definition of $\varlimsup$ and what we have proved in (2), $f$ is upper semicontinuous at
$x_0$.
Since $x_0$ was arbitrary, $f$ is upper semicontinuous on $X$.

\medskip

\emph{(1) $\Rightarrow$ (4).}
Assume that $f$ is upper semicontinuous on $X$.
We show that $\operatorname{hypo}(f)$ is closed in $X\times\mathbb{R}$.

Let $(x_n,t_n)$ be a sequence in $\operatorname{hypo}(f)$ such that
$(x_n,t_n) \to (x,t)$ in $X\times\mathbb{R}$.
Then $x_n \to x$ in $X$ and $t_n \to t$ in $\mathbb{R}$, and since
$(x_n,t_n)\in \operatorname{hypo}(f)$ we have
\begin{equation*}
t_n \leq f(x_n)
\quad\text{for all } n.
\end{equation*}
By upper semicontinuity at $x$,
\begin{equation*}
\varlimsup_{n\to\infty} f(x_n) \leq f(x).
\end{equation*}
On the other hand,
\begin{equation*}
\lim_{n\to\infty}t_n=\varlimsup_{n\to\infty} t_n \leq \varlimsup_{n\to\infty} f(x_n).
\end{equation*}
Combining these inequalities yields
\begin{equation*}
t = \lim_{n\to\infty} t_n
\leq
\varlimsup_{n\to\infty} f(x_n)
\leq f(x).
\end{equation*}
Therefore $(x,t)\in\operatorname{hypo}(f)$.
Since every convergent sequence in $\operatorname{hypo}(f)$ converges
to a point in $\operatorname{hypo}(f)$, the set
$\operatorname{hypo}(f)$ is closed.

\medskip

\emph{(4) $\Rightarrow$ (3).}
Assume that $\operatorname{hypo}(f)$ is closed.
Fix $a\in\mathbb{R}$ and consider the set
\begin{equation*}
U_a := \{x\in X : f(x) < a\}.
\end{equation*}
We prove that $U_a$ is open.

Let $x_0\in U_a$, so that $f(x_0) < a$.
Suppose, by contradiction, that $x_0$ is not an interior point of
$U_a$.
Then for every $m\in\mathbb{N}$ there exists $x_m\in X$ such that
\begin{equation*}
d(x_m,x_0) < \frac{1}{m}
\quad\text{and}\quad
x_m \notin U_a,
\end{equation*}
that is,
\begin{equation*}
f(x_m) \geq a
\quad\text{for all } m.
\end{equation*}
For each $m$, define $t_m := a$.
Then
\begin{equation*}
t_m = a \leq f(x_m),
\end{equation*}
so $(x_m,t_m)\in\operatorname{hypo}(f)$ for all $m$.
Moreover,
\begin{equation*}
(x_m,t_m) \to (x_0,a) \quad\text{as } m\to\infty.
\end{equation*}
Since $\operatorname{hypo}(f)$ is closed, we must have
$(x_0,a)\in\operatorname{hypo}(f)$, which means
\begin{equation*}
a \leq f(x_0).
\end{equation*}
This contradicts $f(x_0) < a$.
Hence our assumption was false, and $x_0$ is an interior point of
$U_a$.
Since $x_0\in U_a$ was arbitrary, the set $U_a$ is open.
\end{proof}

\begin{question}
	Prove the following inequality:
	\begin{equation*}
		\sin (\sin x)<\cos (\cos x),\quad \forall x\in\r.
	\end{equation*}
\end{question}
\begin{proof}
	Since $\sin x+\cos x=\sqrt{2}\(\frac{\sqrt{2}}{2}\sin x+\frac{\sqrt{2}}{2}\cos x\)=\sqrt{2}\sin \(x+\frac{\pi}{4}\)\leq \sqrt{2}<\frac{\pi}{2}$, and $x\mapsto \sin x$ is strictly increasing in $[-1,1]$, we have
	\begin{equation*}
		\sin (\sin x)<\sin \(\frac{\pi}{2}-\cos x\)=\cos (\cos x).
	\end{equation*}
\end{proof}

\begin{question}
	Prove the limit $\lim_{n\to \infty} \sin n $ does not exist.
\end{question}
\begin{proof}
	We prove this by contradiction. If the limit exists, denoted by $\lim_{n\to \infty} \sin n=a$, we can let $n\to\infty$ in the following equality 
	\begin{equation*}
		\sin (n+1)+\sin (n-1)=2\sin n\cos 1
	\end{equation*}
	to obtain $2a=2a\cos 1$. If $a\neq 0$, we have $\cos 1=1$, yielding a contradiciton. If $a=0$, we let $n\to \infty$ in the following equality
	\begin{equation*}
		\sin (n+1)-\sin (n-1)=2\cos n\sin 1
	\end{equation*}
	to obtain $0=2\lim_{n\to\infty}\cos n \sin 1\Longrightarrow \cos n\to 0,n\to\infty$. Letting $n\to\infty$ in the following equality 
	\begin{equation*}
		\sin^2 n+\cos^2 n =1,
	\end{equation*}
	we will get $0=1$, which is also a contradiction.
\end{proof}


\begin{question}[Deterministic $o(n^{-1/2})$ equivalence]
Let $(E,\|\cdot\|)$ be a normed linear space and let $(b_n)_{n\ge1}$ be a
\emph{deterministic} sequence in $E$. The following are equivalent:
\begin{enumerate}
\item $\|b_n\|=o(n^{-1/2})$ as $n\to\infty$;
\item There exists a nonnegative real sequence $(a_n)_{n\ge1}$ with $a_n=o(1)$
such that for all $n$,
\begin{equation*}
	\|b_n\|\ \le\ a_n\,n^{-1/2}.
\end{equation*}
\end{enumerate}
\end{question}

\begin{proof}
\emph{(1 $\Rightarrow$ 2)} Assume $\|b_n\|=o(n^{-1/2})$. Define
\begin{equation*}
	a_n:=n^{1/2}\,\|b_n\|\ \ge\ 0.
\end{equation*}
Then $a_n\to0$ and $\|b_n\|=a_n\,n^{-1/2}$, hence in particular
$\|b_n\|\le a_n\,n^{-1/2}$ for all $n$.

\smallskip
\noindent
\emph{(2 $\Rightarrow$ 1)} If there exists $a_n\ge0$ with $a_n=o(1)$ and
$\|b_n\|\le a_n\,n^{-1/2}$ for all $n$, then multiplying both sides by $n^{1/2}$
gives
\begin{equation*}
	n^{1/2}\,\|b_n\|\ \le\ a_n\ \longrightarrow\ 0,
\end{equation*}
so $\|b_n\|=o(n^{-1/2})$.
\end{proof}

\begin{remark}
The statement is specific to \emph{deterministic} sequences. For random
sequences $(X_n)$, the statement $\|X_n\|=o_P(n^{-1/2})$ is \emph{not}
equivalent to the existence of a deterministic $(a_n)$ with
$\|X_n\|\le a_n n^{-1/2}$ almost surely for all $n$; a high-probability
formulation is the correct analogue:
$\exists\,a_n\downarrow0$ such that
$\mathbb P\!\big(\|X_n\|\le a_n n^{-1/2}\big)\to1$.
\end{remark}

\begin{question}[(log-concavity of $F$ and $1-F$)]
	Let $f:\mathbb{R}\to[0,\infty)$ be a one-dimensional probability density
function. Assume that $f>0$ on the interior of its support, that $f$ is
continuously differentiable there, and that $f$ is log-concave, i.e.
$\log f$ is a concave function on the interior of the support.  Define
\begin{equation*}
F(x) := \int_{(-\infty,x]} f(t)\,dt,
\qquad
1-F(x) := \int_{[x,\infty)} f(t)\,dt .
\end{equation*}
Then $F$ and $1-F$ are log-concave on the interior of the support.
\end{question}


\begin{proof}

Because $f$ is log-concave, there exists a concave function
$\varphi:\mathbb{R}\to[-\infty,\infty)$ such that $f(x)=e^{\varphi(x)}.$
For any $c>0$ the super-level set
\begin{equation*}
A_c := \{x\in\mathbb{R} : f(x)\ge c\} = \{x : \varphi(x)\ge \log c\}
\end{equation*}
is convex, since $\varphi$ is concave and hence its super-level sets are
convex. In one dimension this means that $A_c$ is an interval (possibly
empty, a bounded interval, a half-line or all of $\mathbb{R}$).

Suppose that $\varlimsup_{x\to+\infty} f(x) > 0$. Then there exist $c>0$
and a sequence $x_n\to+\infty$ such that $f(x_n)\ge c$ for all $n$.
Hence $A_c$ is unbounded to the right. Because $A_c$ is a convex subset
of $\mathbb{R}$ that is unbounded to the right, there exists $R\in
\mathbb{R}$ such that
\begin{equation*}
[R,\infty) \subset A_c .
\end{equation*}
Consequently
\begin{equation*}
\int_{\mathbb{R}} f(x)\,dx
\;\ge\;
\int_R^\infty f(x)\,dx
\;\ge\;
\int_R^\infty c\,dx
= \infty ,
\end{equation*}
which contradicts that $f$ is a probability density. Therefore $\varlimsup_{x\to+\infty} f(x)\le 0$.
Since $f(x)\ge 0$, we obtain $\lim_{x\to+\infty} f(x) = 0 .$
The same argument applied to $x\to-\infty$ yields $\lim_{x\to-\infty} f(x) = 0$.

On the interior of the support we have $F'(x)=f(x)$ and $F''(x)=f'(x)$.
Moreover $\log f$ is concave and twice differentiable there, so
\begin{equation*}
(\log f)'(x) = \frac{f'(x)}{f(x)}
\end{equation*}
is nonincreasing as a function of $x$.

For any $t\le x$ the monotonicity of $f'/f$ gives
\begin{equation*}
\frac{f'(t)}{f(t)} \ge \frac{f'(x)}{f(x)}.
\end{equation*}
Multiplying by $f(t)f(x)>0$ yields
\begin{equation*}
f'(t) f(x) \ge f'(x) f(t).
\end{equation*}
Integrating both sides from $t=-\infty$ to $t=x$ and using the tail
limit from Step~1, we obtain
\begin{equation*}
\int_{-\infty}^x f'(t) f(x)\,dt
\;\ge\;
\int_{-\infty}^x f'(x) f(t)\,dt
\quad\Longrightarrow\quad
f(x)^2 \ge f'(x) F(x).
\end{equation*}
Now compute the second derivative of $\log F$:
\begin{equation*}
(\log F)'(x) = \frac{F'(x)}{F(x)} = \frac{f(x)}{F(x)},
\end{equation*}
\begin{equation*}
(\log F)''(x)
= \frac{F''(x)F(x)-F'(x)^2}{F(x)^2}
= \frac{f'(x)F(x)-f(x)^2}{F(x)^2}
\le 0,
\end{equation*}
where the inequality follows from $f'(x)F(x)-f(x)^2\le 0$ proved above.
Thus $\log F$ is concave on the interior of the support, i.e.\ $F$ is
log-concave there.

Define the survival function
\begin{equation*}
G(x) := 1-F(x) = \int_x^\infty f(t)\,dt.
\end{equation*}
Then $G'(x)=-f(x)$ and $G''(x)=-f'(x)$, so
\begin{equation*}
(\log G)'(x) = \frac{G'(x)}{G(x)} = -\frac{f(x)}{G(x)} ,
\end{equation*}
\begin{equation*}
(\log G)''(x)
= \frac{G''(x)G(x)-G'(x)^2}{G(x)^2}
= \frac{-f'(x)G(x)-f(x)^2}{G(x)^2}.
\end{equation*}
We again use the monotonicity of $f'/f$, now for $t\ge x$:
\begin{equation*}
\frac{f'(t)}{f(t)} \le \frac{f'(x)}{f(x)} \quad\Rightarrow\quad
f'(t) f(x) \le f'(x) f(t).
\end{equation*}
Integrating from $t=x$ to $t=\infty$ and using the tail limit, we obtain
\begin{equation*}
\int_x^\infty f'(t) f(x)\,dt
\;\le\;
\int_x^\infty f'(x) f(t)\,dt
\quad\Longrightarrow\quad
- f(x)^2 \le f'(x) G(x).
\end{equation*}
Rewriting this inequality as
\begin{equation*}
f'(x)G(x)+f(x)^2 \ge 0
\end{equation*}
and substituting into the formula for $(\log G)''(x)$ gives
\begin{equation*}
(\log G)''(x)=\frac{-f'(x)G(x)-f(x)^2}{G(x)^2}\le 0.
\end{equation*}
Hence $\log G$ is concave on the interior of the support, so $G(x)=1-F(x)$
is log-concave there.
\end{proof}

\begin{question}
Let $\psi:[0,\infty)\to[0,\infty)$ be a nondecreasing, convex function with
$\psi(0)=0$ and $\lim_{x\to\infty}\psi(x)=\infty$.
Define the generalized inverse by
\begin{equation*}
\psi^{-1}(t)
:=\inf\{x\ge 0:\ \psi(x)\ge t\},\qquad t\ge 0.
\end{equation*}
Show that for all $t\ge 0$,
\begin{equation*}
\psi^{-1}(2t)\le 2\,\psi^{-1}(t).
\end{equation*}
\end{question}

\begin{proof}
Since $\psi$ is convex and satisfies $\psi(0)=0$, for any $\lambda\in[0,1]$
and $x\ge 0$ we have
\begin{equation*}
\psi(\lambda x)
\le \lambda \psi(x).
\end{equation*}
Applying this inequality with $\lambda=\tfrac12$ yields
\begin{equation*}
\psi(x/2)\le \tfrac12 \psi(x),
\end{equation*}
which is equivalent to
\begin{equation*}
\psi(2u)\ge 2\psi(u)
\qquad\text{for all }u\ge 0.
\end{equation*}

Fix $t\ge 0$ and set
\begin{equation*}
r:=\psi^{-1}(t)
=\inf\{x\ge 0:\ \psi(x)\ge t\}.
\end{equation*}
By the definition of the infimum and the monotonicity of $\psi$,  for every $\varepsilon>0$ there exists
$u=r+\varepsilon$ such that
\begin{equation*}
\psi(u)\ge t.
\end{equation*}
Applying the inequality above gives
\begin{equation*}
\psi(2u)\ge 2\psi(u)\ge 2t.
\end{equation*}
Hence $2u$ belongs to the set $\{x\ge 0:\psi(x)\ge 2t\}$, and therefore
\begin{equation*}
\psi^{-1}(2t)
\le 2u
=2(r+\varepsilon).
\end{equation*}
Since $\varepsilon>0$ is arbitrary, letting $\varepsilon\downarrow 0$ yields
\begin{equation*}
\psi^{-1}(2t)\le 2r=2\,\psi^{-1}(t),
\end{equation*}
as claimed.
\end{proof}

\begin{question}
Let $\psi:[0,\infty)\to[0,\infty]$ be a nondecreasing function. For $t\ge 0$,
define the generalized inverses
\begin{equation*}
a(t):=\inf\{x\ge 0:\psi(x)\ge t\},\qquad
b(t):=\inf\{x\ge 0:\psi(x)> t\},
\end{equation*}
and also
\begin{equation*}
c(t):=\sup\{x\ge 0:\psi(x)\le t\},\qquad
d(t):=\sup\{x\ge 0:\psi(x)< t\}.
\end{equation*}
Prove the following statements.
\begin{enumerate}[(1)]
\item The complementary-pair identities hold:
\begin{equation*}
a(t)=d(t),\qquad b(t)=c(t),\qquad t\ge 0.
\end{equation*}
\item The one-sided limit relations hold:
\begin{equation*}
b(t)=a(t+):=\lim_{s\downarrow t}a(s),\qquad
a(t)=b(t-):=\lim_{s\uparrow t}b(s),\qquad t\ge 0.
\end{equation*}
\item The threshold (separation) properties hold: for every $t\ge 0$ and $x\ge 0$,
\begin{equation*}
x<a(t)\ \Rightarrow\ \psi(x)<t,\qquad
x>a(t)\ \Rightarrow\ \psi(x)\ge t,
\end{equation*}
and
\begin{equation*}
x<b(t)\ \Rightarrow\ \psi(x)\le t,\qquad
x>b(t)\ \Rightarrow\ \psi(x)>t.
\end{equation*}
\item When $\psi$ is right-continuous or left-continuous, then for every $t\ge 0$ and $x\ge 0$,
\begin{equation*}
x\ge a(t)\ \Longleftrightarrow\ \psi(x)\ge t\text{ or }x\le b(t)\ \Longleftrightarrow\ \psi(x)\le t
\end{equation*}
respectively.
\item $a$ is left--continuous,
and that $b$ is right--continuous, i.e.
\begin{equation*}
\lim_{s\uparrow t} a(s)=a(t),
\qquad
\lim_{s\downarrow t} b(s)=b(t),
\qquad t\ge 0.
\end{equation*}
\end{enumerate}
\end{question}

\begin{proof}
\textbf{(1)}
Fix $t\ge 0$ and define
\begin{equation*}
A_t:=\{x\ge 0:\psi(x)\ge t\},\qquad B_t:=\{x\ge 0:\psi(x)<t\}.
\end{equation*}
Since $\psi$ is nondecreasing, $A_t$ is an upper set: if $x\in A_t$ and $y\ge x$,
then $\psi(y)\ge \psi(x)\ge t$, hence $y\in A_t$.
Therefore, there exists $\alpha\in[0,\infty]$ such that $A_t$ is either
$[\alpha,\infty)$ or $(\alpha,\infty)$, and consequently
\begin{equation*}
B_t=[0,\infty)\setminus A_t
\end{equation*}
is either $[0,\alpha)$ or $[0,\alpha]$. In either case,
\begin{equation*}
a(t)=\inf A_t=\alpha=\sup B_t=d(t).
\end{equation*}
Similarly, letting
\begin{equation*}
C_t:=\{x\ge 0:\psi(x)>t\},\qquad D_t:=\{x\ge 0:\psi(x)\le t\},
\end{equation*}
the set $C_t$ is an upper set and hence is either $[\beta,\infty)$ or
$(\beta,\infty)$ for some $\beta\in[0,\infty]$, while $D_t$ is either
$[0,\beta)$ or $[0,\beta]$. Thus
\begin{equation*}
b(t)=\inf C_t=\beta=\sup D_t=c(t).
\end{equation*}

\medskip
\noindent \textbf{(2)}
First we prove $b(t)=a(t+)$. Let $s>t$. Then
\begin{equation*}
\{x:\psi(x)\ge s\}\subseteq \{x:\psi(x)>t\},
\end{equation*}
so taking infima yields $a(s)\ge b(t)$. Hence
\begin{equation*}
a(t+)=\inf_{s>t}a(s)\ge b(t).
\end{equation*}
For the reverse inequality, fix $\varepsilon>0$ and set $x_\varepsilon:=b(t)+\varepsilon$.
By definition of $b(t)$, we have $\psi(x_\varepsilon)>t$. Define
\begin{equation*}
s_\varepsilon:=\frac{\psi(x_\varepsilon)+t}{2},
\end{equation*}
so that $t<s_\varepsilon<\psi(x_\varepsilon)$, hence $\psi(x_\varepsilon)\ge s_\varepsilon$ and
\begin{equation*}
a(s_\varepsilon)=\inf\{x:\psi(x)\ge s_\varepsilon\}\le x_\varepsilon=b(t)+\varepsilon.
\end{equation*}
Since $s_\varepsilon>t$, it follows that
\begin{equation*}
a(t+)=\inf_{s>t}a(s)\le a(s_\varepsilon)\le b(t)+\varepsilon.
\end{equation*}
Letting $\varepsilon\downarrow 0$ gives $a(t+)\le b(t)$, hence $a(t+)=b(t)$.

Next we prove $a(t)=b(t-)$. Let $s<t$. Then
\begin{equation*}
\{x:\psi(x)>s\}\supseteq \{x:\psi(x)\ge t\},
\end{equation*}
so taking infima yields $b(s)\le a(t)$. Hence
\begin{equation*}
b(t-)=\sup_{s<t}b(s)\le a(t).
\end{equation*}
For the reverse inequality, fix $\varepsilon>0$ and set $y_\varepsilon:=a(t)-\varepsilon$.
By definition of $a(t)$, we have $\psi(y_\varepsilon)<t$. Define
\begin{equation*}
s_\varepsilon:=\frac{\psi(y_\varepsilon)+t}{2},
\end{equation*}
so that $\psi(y_\varepsilon)<s_\varepsilon<t$, hence $\psi(y_\varepsilon)\le s_\varepsilon$ and
\begin{equation*}
b(s_\varepsilon)=\sup\{x:\psi(x)\le s_\varepsilon\}\ge y_\varepsilon=a(t)-\varepsilon,
\end{equation*}
where we used the identity $b(s)=c(s)=\sup\{x:\psi(x)\le s\}$ from Step 1.
Since $s_\varepsilon<t$, it follows that
\begin{equation*}
b(t-)=\sup_{s<t}b(s)\ge b(s_\varepsilon)\ge a(t)-\varepsilon.
\end{equation*}
Letting $\varepsilon\downarrow 0$ gives $b(t-)\ge a(t)$, hence $b(t-)=a(t)$.

\medskip \noindent \textbf{(3)}
Fix $t\ge 0$. If $x<a(t)$ and $\psi(x)\ge t$, then $x\in\{y:\psi(y)\ge t\}$, implying
\begin{equation*}
a(t)=\inf\{y:\psi(y)\ge t\}\le x,
\end{equation*}
a contradiction. Hence $x<a(t)$ implies $\psi(x)<t$.

If $x>a(t)$ and $\psi(x)<t$, then by monotonicity, for every $y\le x$ we have
$\psi(y)\le\psi(x)<t$, so no $y\le x$ can satisfy $\psi(y)\ge t$. Therefore,
\begin{equation*}
a(t)=\inf\{y:\psi(y)\ge t\}\ge x,
\end{equation*}
a contradiction. Hence $x>a(t)$ implies $\psi(x)\ge t$.

If $x<b(t)$ and $\psi(x)>t$, then $x\in\{y:\psi(y)>t\}$, implying
\begin{equation*}
b(t)=\inf\{y:\psi(y)>t\}\le x,
\end{equation*}
a contradiction. Hence $x<b(t)$ implies $\psi(x)\le t$.

If $x>b(t)$ and $\psi(x)\le t$, then by monotonicity, for every $y\le x$ we have
$\psi(y)\le\psi(x)\le t$, so no $y\le x$ can satisfy $\psi(y)>t$. Therefore,
\begin{equation*}
b(t)=\inf\{y:\psi(y)>t\}\ge x,
\end{equation*}
a contradiction. Hence $x>b(t)$ implies $\psi(x)>t$.

\medskip
\noindent \textbf{(4)}
Assume $\psi$ is right-continuous. First we show $\psi(a(t))\geq t$. Indeed, by the definition of $a(t)$, there exists a sequence $(x_n)$ satisfying $x_n\downarrow a(t)$ and $\psi(x_n)\geq t$. The right-continuity of $\psi$ shows $\psi(a(t))=\lim_{n\to\infty}\psi(x_n)\geq t$. Besides, it is easy to see $a(t)$ is non-decreasing. From $x\ge a(t)$ we have $\psi(x)\ge \psi(a(t))\geq t$. On the contrary, if $\psi(x)\geq t$, it is easy to get $x\geq a(t)$ by the definition of $a(t)$.

Now, assume $\psi$ is left-continuous. First we show $\psi(b(t))\leq t$. We have shown $b(t)=c(t)$ in (1), so we just need to show $\psi(c(t))\leq t$. By the definition of $c(t)$, there exists $(x_n)$ satisfying $x_n\uparrow c(t)$ and $\psi(x_n)\leq t$. The left-continuity of $\psi$ shows $\psi(c(t))=\lim_{n\to\infty}\psi(x_n)\leq t$. Besides, it is easy to see $b(t)$ is non-decreasing. From $x\le b(t)$ we have $\psi(x)\le \psi(b(t))\leq t$. On the contrary, if $\psi(x)\geq t$, it is easy to get $x\leq c(t)=b(t)$ by the definition of $c(t)$.

\medskip
\noindent \textbf{(5)}
Fix $t\ge 0$ and let $(s_n)_{n\ge1}$ be any sequence such that
\begin{equation*}
s_n\uparrow t \quad\text{and}\quad s_n<t.
\end{equation*}
Since $a$ is nondecreasing, the sequence $a(s_n)$ is nondecreasing and therefore
admits a limit
\begin{equation*}
L:=\lim_{n\to\infty} a(s_n)\le a(t).
\end{equation*}
Assume, for contradiction, that $L<a(t)$.
Choose $x$ such that
\begin{equation*}
L<x<a(t).
\end{equation*}
By the definition of $a(t)$, the inequality $x<a(t)$ implies
\begin{equation*}
\psi(x)<t.
\end{equation*}
Since $s_n\uparrow t$, there exists $N$ such that $s_n>\psi(x)$ for all $n\ge N$.
Thus $\psi(x)<s_n$, and using the representation
\begin{equation*}
a(s_n)=\sup\{y\ge 0:\psi(y)<s_n\},
\end{equation*}
we obtain
\begin{equation*}
a(s_n)\ge x \quad\text{for all } n\ge N.
\end{equation*}
Taking limits yields $L\ge x$, contradicting $x>L$.
Therefore $L=a(t)$, and $a$ is left--continuous at $t$.

Finally, we prove the right continuity of $b$.
Fix $t\ge 0$ and let $(s_n)_{n\ge1}$ be any sequence such that
\begin{equation*}
s_n\downarrow t \quad\text{and}\quad s_n>t.
\end{equation*}
Since $b$ is nondecreasing, the sequence $b(s_n)$ is nonincreasing and hence
converges to some limit
\begin{equation*}
L:=\lim_{n\to\infty} b(s_n)\ge b(t).
\end{equation*}
Suppose, for contradiction, that $L>b(t)$.
Choose $x$ such that
\begin{equation*}
b(t)<x<L.
\end{equation*}
By the definition of $b(t)$, the inequality $x>b(t)$ implies
\begin{equation*}
\psi(x)>t.
\end{equation*}
Since $s_n\downarrow t$, there exists $N$ such that $s_n<\psi(x)$ for all $n\ge N$.
Thus $\psi(x)>s_n$, and from
\begin{equation*}
b(s_n)=\inf\{y\ge 0:\psi(y)>s_n\}
\end{equation*}
we obtain
\begin{equation*}
b(s_n)\le x \quad\text{for all } n\ge N.
\end{equation*}
Taking limits gives $L\le x$, contradicting $x<L$.
Hence $L=b(t)$, and $b$ is right--continuous at $t$.
\end{proof}

\begin{question}
Let $A \subset \mathbb{R}^n$ and let $f: A \to \mathbb{R}$.
\begin{enumerate}
    \item Show that if $f$ is uniformly continuous on $A$, then $f$ admits a unique continuous extension to $\overline{A}$.
    \item Give an example to show that a continuous function on $A$ need not admit a continuous extension to $\overline{A}$.
\end{enumerate}
\end{question}

\begin{proof}
\textbf{1.}
Assume that $f$ is uniformly continuous on $A$.
Fix $x_0 \in \overline{A}$. Choose a sequence $(x_k)\subset A$ with $x_k\to x_0$.
Since $(x_k)$ converges in $\mathbb{R}^n$, it is Cauchy. By uniform continuity, $(f(x_k))$ is Cauchy in $\mathbb{R}$:
for every $\varepsilon>0$ there exists $\delta>0$ such that for all $x,y\in A$,
\begin{equation*}
|x-y|<\delta \quad \Rightarrow \quad |f(x)-f(y)|<\varepsilon,
\end{equation*}
and since $(x_k)$ is Cauchy, there exists $N$ such that $|x_k-x_m|<\delta$ for all $k,m\ge N$, hence
\begin{equation*}
|f(x_k)-f(x_m)|<\varepsilon \quad \text{for all } k,m\ge N.
\end{equation*}
Completeness of $\mathbb{R}$ yields $\lim_{k\to\infty} f(x_k)$ exists. Define
\begin{equation*}
\tilde f(x_0):=\lim_{k\to\infty} f(x_k).
\end{equation*}


Let $(y_k)\subset A$ be another sequence with $y_k\to x_0$.
We claim
\begin{equation*}
\lim_{k\to\infty} f(x_k)=\lim_{k\to\infty} f(y_k).
\end{equation*}
Assume for contradiction that these limits differ. Then there exists $\varepsilon_0>0$ and infinitely many indices $k$ such that
\begin{equation*}
|f(x_k)-f(y_k)|\ge \varepsilon_0.
\end{equation*}
Since $x_k\to x_0$ and $y_k\to x_0$, we have $|x_k-y_k|\to 0$. In particular, for the $\delta>0$ corresponding to
$\varepsilon_0$ in uniform continuity, there exists $K$ such that for all $k\ge K$,
\begin{equation*}
|x_k-y_k|<\delta.
\end{equation*}
Uniform continuity then implies for all $k\ge K$,
\begin{equation*}
|f(x_k)-f(y_k)|<\varepsilon_0,
\end{equation*}
contradiction. Hence $\lim f(x_k)=\lim f(y_k)$ and $\tilde f$ is well-defined.
If $x\in A$, take the constant sequence $x_k\equiv x$ to get $\tilde f(x)=f(x)$.


Fix $x_0\in \overline{A}$ and let $x_j\to x_0$ with $x_j\in \overline{A}$.
We show $\tilde f(x_j)\to \tilde f(x_0)$.
Let $\varepsilon>0$ and choose $\delta>0$ from uniform continuity so that for all $u,v\in A$,
\begin{equation*}
|u-v|<\delta \quad \Rightarrow \quad |f(u)-f(v)|<\varepsilon/3.
\end{equation*}
Since $x_j\to x_0$, there exists $J$ such that for all $j\ge J$,
\begin{equation*}
|x_j-x_0|<\delta/3.
\end{equation*}
For each fixed $j\ge J$, since $x_j\in\overline{A}$, pick $a_j\in A$ with
\begin{equation*}
|a_j-x_j|<\delta/3.
\end{equation*}
Also pick $a_0\in A$ with
\begin{equation*}
|a_0-x_0|<\delta/3.
\end{equation*}
Then for $j\ge J$,
\begin{equation*}
|a_j-a_0|
\le |a_j-x_j|+|x_j-x_0|+|x_0-a_0|
< \delta/3+\delta/3+\delta/3=\delta,
\end{equation*}
so by uniform continuity,
\begin{equation*}
|f(a_j)-f(a_0)|<\varepsilon/3.
\end{equation*}

Next, by the definition of $\tilde f(x_j)$ as a limit along sequences in $A$ converging to $x_j$, we can choose $a_j$ above
so that additionally
\begin{equation*}
|\tilde f(x_j)-f(a_j)|<\varepsilon/3,
\end{equation*}
and similarly choose $a_0$ so that
\begin{equation*}
|\tilde f(x_0)-f(a_0)|<\varepsilon/3.
\end{equation*}
(Indeed, if $(x_{j,k})\subset A$ with $x_{j,k}\to x_j$, then $f(x_{j,k})\to \tilde f(x_j)$, so take $a_j=x_{j,k}$ for $k$ large.)

Therefore for all $j\ge J$,
\begin{equation*}
|\tilde f(x_j)-\tilde f(x_0)|
\le |\tilde f(x_j)-f(a_j)|+|f(a_j)-f(a_0)|+|f(a_0)-\tilde f(x_0)|
< \varepsilon.
\end{equation*}
Hence $\tilde f(x_j)\to \tilde f(x_0)$, so $\tilde f$ is continuous at $x_0$. Since $x_0$ was arbitrary, $\tilde f$ is continuous
on $\overline{A}$.

If $g:\overline{A}\to\mathbb{R}$ is another continuous extension of $f$, then for any $x_0\in\overline{A}$ and any sequence
$(x_k)\subset A$ with $x_k\to x_0$,
\begin{equation*}
g(x_0)=\lim_{k\to\infty} g(x_k)=\lim_{k\to\infty} f(x_k)=\tilde f(x_0),
\end{equation*}
so $g=\tilde f$.

\bigskip
\noindent
\textbf{2. Example.}
Let $A=(0,1)$ and define $f(x)=1/x$. Then $f$ is continuous on $(0,1)$, but
\begin{equation*}
\lim_{x\to 0^+,\, x\in A} f(x)=+\infty,
\end{equation*}
so there is no real number $L$ that could serve as the value of a continuous extension at $0$. Hence, $f$ admits no continuous
extension to $\overline{A}=[0,1]$.

\end{proof}

\section{Algebra}
\begin{question}
	Given two positive definite matrices $A$ and $B$. If $A-B$ is nonnegative definite, we write $A\geq B$. Prove:
	\begin{enumerate}[(a)]
		\item The order $\geq$ is steady under congruence transformation, i.e., for any invertible matrix $U$,
		\begin{equation*}
			A\geq B\Longleftrightarrow U'AU\geq U'BU.
		\end{equation*}
		\item If $B=I$, then the order $\geq$ is steady under inverse operation, i.e., 
		\begin{equation*}
			A\geq I\Longleftrightarrow I\geq A^{-1}.
		\end{equation*}
		\item The order $\geq$ is steady under inverse operation, i.e., 
		\begin{equation*}
			A\geq B\Longleftrightarrow B^{-1}\geq A^{-1}.
		\end{equation*}
	\end{enumerate} 
\end{question}
\begin{proof}
	\begin{enumerate}[(a)]
		\item Since $A\geq B$ iff $v'(A-B)v\geq 0$ for any $v\in\r^n$. Since each $v\in\r^n$, there exists $u\in\r^n$ such that $v=Uu$. The latter one is equivalent to $u'U'(A-B)Uu=u'(U'AU-U'BU)u\geq 0$ for any $u\in\r^n$. So, $U'AU\geq U'BU$.
		\item Noting $A$ is positive definite, there exists positive definite matrix $A^{1/2}$ such that $A=A^{1/2}A^{1/2}\Longleftrightarrow A^{-1}=A^{-1/2}A^{-1/2}$. Moreover, we can find
		\begin{equation*}
			I-A^{-1}=A^{-1/2}(A-I)A^{-1/2}=A^{-1/2}AA^{-1/2}-A^{-1/2}IA^{-1/2}.
		\end{equation*}
		Letting $U=A^{-1/2}$ and applying (a), we know $I\geq A^{-1}$.
		\item Applying (a) and (b), we have
		\begin{equation*}
			\begin{split}
				A\geq B&\Longleftrightarrow B^{-1/2}AB^{-1/2}\geq I\\
				&\Longleftrightarrow I\geq B^{1/2}A^{-1}B^{1/2}\Longleftrightarrow B^{-1}=B^{-1/2}IB^{-1/2}\geq A^{-1}.
			\end{split}
		\end{equation*}
	\end{enumerate}
\end{proof}

\begin{question}
	Suppose $A\in \r^{m\times m}$ is positive definite and $B\in \r^{m\times n}$ has rank $n$, then
	\begin{equation*}
		(B'B)^{-1}B'AB(B'B)^{-1}-(B'A^{-1}B)^{-1}
	\end{equation*}
	is a positive definite matrix.
\end{question}
\begin{proof}
	Based on part (c) of previous QUESTION, we know what we need to do is equivalent to prove
	\begin{equation*}
		B'A^{-1}B - B'B(B'AB)^{-1}B'B
	\end{equation*}
	is positive semi-definite. Then, we can find
	\begin{equation*}
		\begin{split}
			&B'A^{-1}B - B'B(B'AB)^{-1}B'B\\
			&\qquad =B'A^{-1/2}[I-A^{1/2}B(B'AB)^{-1}B'A^{1/2}]A^{-1/2}B\\
			&\qquad =B'A^{-1/2}\Omega A^{-1/2}B.
		\end{split}
	\end{equation*}
	We can easily verify that $\Omega=\Omega'$ and $\Omega=\Omega^2$. So,
	\begin{equation*}
		\begin{split}
			&B'A^{-1}B - B'B(B'AB)^{-1}B'B\\
			&\qquad =B'A^{-1/2}\Omega'\Omega A^{-1/2}B=(\Omega A^{-1/2}B)'(\Omega A^{-1/2}B),
		\end{split}
	\end{equation*}
	which shows the above matrix is positive semi-definite.
\end{proof}

\begin{question}
	For any $x\in\mathbb{R}^n$, the rank-one matrix $A:=xx^\top$ satisfies
\begin{equation*}
	\|A\|_{F}=\|x\|_2^{\,2}.
\end{equation*}
\end{question}

\begin{proof}
By definition of the Frobenius norm,
\begin{equation*}
\|xx^\top\|_{F}^{2}
=\operatorname{tr}\!\big((xx^\top)^\top(xx^\top)\big)
=\operatorname{tr}(xx^\top xx^\top).
\end{equation*}
Since $x^\top x$ is a scalar, we may regroup and commute it with matrices:
\begin{equation*}
\operatorname{tr}(xx^\top xx^\top)
=\operatorname{tr}\!\big(x(x^\top x)x^\top\big)
=\operatorname{tr}\!\big((x^\top x)\,xx^\top\big)
=(x^\top x)\,\operatorname{tr}(xx^\top).
\end{equation*}
Using the cyclic property of the trace, $\operatorname{tr}(xx^\top)=\operatorname{tr}(x^\top x)$, and
$\operatorname{tr}(x^\top x)=x^\top x$ because $x^\top x$ is a $1\times 1$ matrix. Hence
\begin{equation*}
\|xx^\top\|_{F}^{2}
=(x^\top x)\,\operatorname{tr}(x^\top x)
=(x^\top x)\,(x^\top x)
=(x^\top x)^2
=\|x\|_2^{\,4}.
\end{equation*}
Taking square roots on both sides gives $\|xx^\top\|_{F}=\|x\|_2^{\,2}$.
\end{proof}

\begin{question}
	For any normed linear space $(V,||\cdot||)$, the norm function $f\mapsto ||f||$ is convex. 
\end{question}
\begin{proof}
	For any $t\in [0,1]$ and $f,g\in V$, we have
	\begin{equation*}
		||tf+(1-t)g||\leq ||tf||+||(1-t)g||=t||f||+(1-t)||g||.
	\end{equation*}
\end{proof}

\begin{question}
	For any $A\in M_n(\r)$, define
	\begin{equation*}
		||A||_{2\to 2}:=\sup_{||x||_2\neq 0}\frac{||Ax||_2}{||x||_2}.
	\end{equation*}
	Then 
	\begin{itemize}
		\item $||A||_{2\to 2}=\sqrt{\sigma_{\max}(A)}$, where $\sigma_{\max}(A)$ is the biggest singular value of $A$. 
		\item Moreover, when $A$ is symmetric, $||A||_{2\to 2}=|\lambda_{ \max}(A)|$, where $\lambda_{ \max}(A)$ is the biggest eigenvalue of $A$. 
		\item When $A=BB'$, where $B\in M_{n\times k}(\r)$, we have $||A||_{2\to 2}=||B||_{2\to 2}^2$. 
		\item When $A=xx'$, where $x\in \r^n$, we have $||A||_{2\to 2}=||x||_2^2$.
	\end{itemize}
\end{question}
\begin{proof}
Throughout the proof, we adopt the convention
\begin{equation*}
\sigma_{\max}(A):=\lambda_{\max}(AA^\top).
\end{equation*}
i.e.\ the (squared) largest singular value is the largest eigenvalue of $AA^\top$.
\begin{itemize}
	\item By the definition of the operator norm and the Euclidean inner product,
\begin{equation*}
\|A\|_{2\to 2}^2=\sup_{\|x\|_2=1}\|Ax\|_2^2=\sup_{\|x\|_2=1} x^\top A^\top A x.
\end{equation*}
The variational characterization of the largest eigenvalue yields
\begin{equation*}
\sup_{\|x\|_2=1} x^\top A^\top A x=\lambda_{\max}(A^\top A).
\end{equation*}
It is a standard fact that $A^\top A$ and $AA^\top$ have the same nonzero eigenvalues (with the same multiplicities). Hence
\begin{equation*}
\|A\|_{2\to 2}^2=\lambda_{\max}(A^\top A)=\lambda_{\max}(AA^\top)=\sigma_{\max}(A),
\end{equation*}
and therefore $\|A\|_{2\to 2}=\sqrt{\sigma_{\max}(A)}$.
	\item If $A$ is symmetric, then $A^\top A=A^2$, so
\begin{equation*}
\|A\|_{2\to 2}^2=\lambda_{\max}(A^2)=\max_{1\le i\le n}\lambda_i(A)^2.
\end{equation*}
Taking square roots gives
\begin{equation*}
\|A\|_{2\to 2}=\max_{1\le i\le n}|\lambda_i(A)|.
\end{equation*}
In particular, if $A\succeq 0$, then $\lambda_i(A)\ge 0$ and $\|A\|_{2\to 2}=\lambda_{\max}(A)$.
	\item If $A=BB^\top$ with $B\in\mathbb{R}^{n\times k}$, then $A\succeq 0$, so by Item 2,
\begin{equation*}
\|A\|_{2\to 2}=\lambda_{\max}(A)=\lambda_{\max}(BB^\top).
\end{equation*}
On the other hand,
\begin{equation*}
\|B\|_{2\to 2}^2=\lambda_{\max}(B^\top B)=\lambda_{\max}(BB^\top),
\end{equation*}
since $B^\top B$ and $BB^\top$ have the same nonzero eigenvalues. Hence $\|A\|_{2\to 2}=\|B\|_{2\to 2}^2$.
	\item If $A=xx^\top$ with $x\in\mathbb{R}^n$, then $A\succeq 0$ and $\mathrm{rank}(A)=1$. We have
\begin{equation*}
AA^\top=(xx^\top)(xx^\top)=x(x^\top x)x^\top,
\end{equation*}
so the only nonzero eigenvalue of $AA^\top$ (hence of $A$) is $x^\top x=\|x\|_2^2$. Therefore,
\begin{equation*}
\|xx^\top\|_{2\to 2}=\lambda_{\max}(A)=\|x\|_2^2.
\end{equation*}
This also follows by applying Item~3 with $B=x$ viewed as an $n\times 1$ matrix. Or, since $xx^\top$ is symmetric,
\begin{equation*}
	\begin{split}
		||A||_{2\to 2}&=\lambda_{\max}(A)= \sup_{||y||_2=1}y^\top(xx^\top) y\\
		&=\sup_{||y||_2=1}(y^\top x)(x^\top y)=\sup_{||y||_2=1}(x^\top y)^2\\
		&=\sup_{||y||_2=1}||x||_2^2\, ||y||_2^2=||x||_2^2,
	\end{split}
\end{equation*}
where $y^\top x=x^\top y\in\r$ and $x^\top y\leq ||x||_2\,||y||_2$ by Cauchy-Schewarz inequality.
\end{itemize} 
\end{proof}





\begin{question}
	Let $A\in\mathbb{R}^{m\times n}$. Then their nonzero eigenvalues (counting geometric multiplicities) coincide. Namely, let $\lambda>0$. Denote the $\lambda$-eigenspace of a square matrix $M$ by
\begin{equation*}
	\mathcal{E}_\lambda(M):=\{v\neq 0: Mv=\lambda v\}\cup\{0\}.
\end{equation*}
Then
\begin{equation*}
	\dim \mathcal{E}_\lambda(A^\top A)=\dim \mathcal{E}_\lambda(AA^\top).
\end{equation*}
\end{question}
\begin{proof}
Let $x\in\mathbb{R}^n$ satisfy $A^\top A x=\lambda x$ with $\lambda>0$. Then
\begin{equation*}
\|Ax\|_2^2=x^\top A^\top A x=\lambda\,\|x\|_2^2>0,
\end{equation*}
so $y:=Ax\neq 0$. Moreover,
\begin{equation*}
AA^\top y=AA^\top Ax=A(A^\top A x)=A(\lambda x)=\lambda (Ax)=\lambda y,
\end{equation*}
showing $y$ is an eigenvector of $AA^\top$ with eigenvalue $\lambda$.
Conversely, if $AA^\top y=\lambda y$ with $\lambda>0$, then
\begin{equation*}
\|A^\top y\|_2^2=y^\top AA^\top y=\lambda\,\|y\|_2^2>0,
\end{equation*}
so $x:=A^\top y\neq 0$ and
\begin{equation*}
A^\top A x=A^\top A A^\top y=A^\top(AA^\top y)=A^\top(\lambda y)=\lambda x.
\end{equation*}
Thus positive eigenvalues and their algebraic multiplicities match between $A^\top A$ and $AA^\top$. 

This defines two linear maps
\begin{equation*}
T:\ \mathcal{E}_\lambda(A^\top A)\to \mathcal{E}_\lambda(AA^\top),\qquad T(x)=Ax
\end{equation*}
and
\begin{equation*}
S:\ \mathcal{E}_\lambda(AA^\top)\to \mathcal{E}_\lambda(A^\top A),\qquad S(y)=\tfrac{1}{\lambda}A^\top y.
\end{equation*}
Finally,
\begin{equation*}
\begin{split}
	S(T(x))=\tfrac{1}{\lambda}A^\top(Ax)=\tfrac{1}{\lambda}(A^\top A x)=x,\\
T(S(y))=A\big(\tfrac{1}{\lambda}A^\top y\big)=\tfrac{1}{\lambda}(AA^\top y)=y,
\end{split}
\end{equation*}
so $T$ and $S$ are mutual inverses. Hence $T$ is a linear isomorphism between the two eigenspaces, and therefore
\begin{equation*}
\dim \mathcal{E}_\lambda(A^\top A)=\dim \mathcal{E}_\lambda(AA^\top).
\end{equation*}
\end{proof}

 \begin{question}[Linear independence of functions and the second moment matrix]
Let $r(X)=(r_1(X),\dots,r_k(X))' \in \mathbb{R}^k$
be a vector of real-valued functions of the random variable $X$. Assume that each
$r_j(X)$ is square-integrable so that the matrix $\e[r(X)r(X)']$
exists.
Show that the following two statements are equivalent:

\begin{enumerate}
\item The functions $r_1,\dots,r_k$ are linearly independent almost surely, i.e.
\begin{equation*}
\sum_{j=1}^k a_j r_j(X)=0 \ \text{a.s.} \quad \Rightarrow \quad a_1=\cdots=a_k=0.
\end{equation*}

\item The matrix
\begin{equation*}
\e[r(X)r(X)']
\end{equation*}
has full rank.
\end{enumerate}
\end{question}

\begin{proof}
Let $M:=\e[r(X)r(X)']$. First note that for any vector $a\in\mathbb{R}^k$,
\begin{equation*}
a'Ma
= a'\e[r(X)r(X)']a
= \e[a'r(X)r(X)'a]
= \e[(a'r(X))^2].
\end{equation*}

\medskip
\noindent
\textbf{(1) $\Rightarrow$ (2).}
Suppose the functions $r_1,\dots,r_k$ are linearly independent almost surely.
Assume by contradiction that $M$ is not full rank. Then there exists a vector
$a\neq 0$ such that
\begin{equation*}
Ma=0.
\end{equation*}
Multiplying by $a'$ gives
\begin{equation*}
a'Ma=0.
\end{equation*}
Using the identity above, we obtain
\begin{equation*}
E[(a'r(X))^2]=0.
\end{equation*}
Since $(a'r(X))^2\ge 0$ almost surely, this implies
\begin{equation*}
a'r(X)=0 \quad \text{a.s.}
\end{equation*}
By linear independence of the functions $r_1,\dots,r_k$, this implies $a=0$,
which contradicts $a\neq 0$. Hence $M$ must have full rank.

\medskip
\noindent
\textbf{(2) $\Rightarrow$ (1).}
Suppose that $M$ has full rank. Let $a\in\mathbb{R}^k$ satisfy
\begin{equation*}
a'r(X)=0 \quad \text{a.s.}
\end{equation*}
Squaring and taking expectations gives
\begin{equation*}
E[(a'r(X))^2]=0.
\end{equation*}
Using the identity above, we obtain
\begin{equation*}
a'Ma=0.
\end{equation*}
Since $M$ has full rank and is symmetric positive semidefinite, then $M$ is positive definite. Thus
\begin{equation*}
a'Ma=0 \quad \Rightarrow \quad a=0.
\end{equation*}
Thus
\begin{equation*}
\sum_{j=1}^k a_j r_j(X)=0 \ \text{a.s.} \quad \Rightarrow \quad a=0,
\end{equation*}
which proves that the functions $r_1,\dots,r_k$ are linearly independent almost surely.
\end{proof}

  \section{Topology}
  \begin{question}
   Suppose $X$ and $Y$ are two metric spaces. Find a linear operator $A:X\to Y$ which is a close mapping but not an open mapping.
  \end{question}
  \begin{proof}
   Consider $X=\r$ and $Y=\r^2$ with Euclidean topology, then we claim the following linear operator $A$ is a close mapping but not an open mapping:
   \begin{equation*}
    A:\r\to\r^2,~x\mapsto\begin{pmatrix}
     \hspace{0.15cm} x \hspace{0.15cm}\\
     \hspace{0.15cm} 0 \hspace{0.15cm}
    \end{pmatrix}.
   \end{equation*}
   The linearity of $A$ is obvious. For all set $T\subset \r$, we have
   \begin{equation*}
    A(T)=T\times\{0\}.
   \end{equation*}
   When $T$ is a close set, then $T\times\{0\}$ is a closed set in $\r^2$. When $T$ is an open set, then $T\times\{0\}$ is a non-open and non-closed set in $\r^2$. These facts indicate that $A$  is a close mapping but not an open mapping.
  \end{proof}

  \begin{question}
   Suppose $(E,d)$ and $(F,\rho )$ are two metric spaces and $f:E\to F$ is continuous. Given $A\subset F$, then $\partial f^{-1}(A)\subset f^{-1}(\partial A)$. Moreover, show $\partial f^{-1}(A)\supset f^{-1}(\partial A)$ is not correct.
  \end{question}
  \begin{proof}
   \begin{enumerate}
    \item Since $x\in \partial f^{-1}(A)$ iff exist $\{x_n\}\subset f^{-1}(A)$ and $\{x_n'\}\subset[f^{-1}(A)]^c=f^{-1}(A^c)$, such that $x_n\to x$ and $x_n'\to x$. Note the continuity of $f$, so  exist $\{y_n\}\subset A$ and $\{y_n'\}\subset A^c$ such that $y_n=f(x_n)\to f(x)$ and $y_n'=f(x_n')\to f(x)$. Therefore $f(x)\in \partial A\Longleftrightarrow x\in f^{-1}(\partial A )$.
    \item Consider $E=F=\r,f(x)=x^2$ and $A=[0,1]$. Then $\partial f^{-1}[0,1]=\{-1,1\}\subset \{-1,0,1\}=f^{-1}(\partial[0,1] )$. So  $\partial f^{-1}(A)\supset f^{-1}(\partial A)$ is not correct.
   \end{enumerate}
  \end{proof}
  
  \begin{question}
  	Given two metric spaces $(E,d)$. Then for any $A,B\subset E$, we have $\partial(A\cap B)\cup \partial(A\cup B)\subset \partial A\cup\partial B$. Moreover, generally $\partial(A\cap B)\cap \partial(A\cup B)\neq \emptyset$.
  \end{question}
  \begin{proof}
  	For any $x\in \partial(A\cap B)$, exist $\{y_n\}_{n\geq 1}\subset A\cap B$ and $\{z_n\}_{n\geq 1}\subset (A\cap B)^c=A^c\cup B^c$ such that $y_n\to x$ and $z_n\to x$. There are infinite items of $\{z_n\}_{n\geq 1}$ in $A^c$ or $B^c$. Without loss of generality, we can pick a subsequence $\{z_{n_k}\}_{k\geq 1}\subset A^c$. Thus 
  	\begin{equation*}
  		A\ni y_{n_k}\to x,\quad A^c\ni z_{n_k}\to x,\quad k\to\infty.
  	\end{equation*} 
  	So $x\in \partial A\subset \partial A\cup\partial B$ and therefore $\partial(A\cap B)\subset \partial A\cup\partial B$. The similar argument also shows $\partial(A\cup B)\subset \partial A\cup\partial B$. So $\partial(A\cap B)\cup \partial(A\cup B)\subset \partial A\cup\partial B$.
  	
  	Consider $A=B(-2,3)\subset\r^2$ and $B=B(2,3)\subset\r^2$. Then the intersection point $(x,y)$ are determined by
  	\begin{equation*}
  		\left\{
  		\begin{split}
  			(x+2)^2+y^2=9,\\
  			(x-2)^2+y^2=9.
  		\end{split}\right.
  		\Longrightarrow (x,y)=(0,\pm\sqrt{5}).
  	\end{equation*}
  	So $\partial(A\cap B)\cap \partial(A\cup B)=\{(0,\sqrt{5}),(0,-\sqrt{5})\}\neq \emptyset$.
  \end{proof}
  
  \begin{question}
  	Suppose $(E,d)$ is a metric space and $\{K_n\}_{n\geq 1}\subset E$ is a sequence of compact sets satisfying $K_n\uparrow E$. If $K\subset E$ is a compact, use example to show the fact: there exists $m\geq 1$, such that $K\subset K_m$ might not hold.
  \end{question}
  \begin{proof}
  	\textcolor{red}{\sffamily Way 1.} Consider the subspace $E:=[0,1]$ of $\r$ with subspace topology. Then $E$ itself is a compact space. Consider $K_n=[0,1-1/n]\cup\{1\}$ for each $n\geq 1$. Then $\{K_n\}_{n\geq 1}$ is an increasing compact set sequence satisfying
  	\begin{equation*}
  		\bigcup_{n\geq 1}K_n=[0,1)\cup\{1\}=E.
  	\end{equation*} 
  	Consider the compact set $K:=E$. Obviously there are no $E_n$ satisfying $K\subset E_n$. 
  	
  	\textcolor{red}{\sffamily Way 2.} Consider $E:=\mathbb{S}^2$. Denote $\r^2\cup\{\infty\}$ the one-point compactification of $\r^2$. So $E=\mathbb{S}^2\cong \r^2\cup\{\infty\}$. Write the homeomorphic mapping as $\varphi:\mathbb{S}^2\to \r^2\cup\{\infty\}$. By the definition of one-point compactification, we know $K_n:=\varphi(\{x\in\r^2:|x|\leq  n\})\cup\{N\}$ is an increasing compact set sequence of $\mathbb{S}^2$, where $N$ is the north pole satisfying 
  	\begin{equation*}
  		\bigcup_{n\geq 1}K_n=\varphi\(\bigcup_{n\geq 1} \{x\in\r^2:|x|\leq  n\}\)\cup\{N\}=\mathbb{S}^2.
  	\end{equation*}
  	Apply the definition of one-point compactification again, consider the compact set $K:=\varphi(\{x\in\r^2:|x|\geq 10 \})\cup\{N\}$. Obviously there are no $E_n$ satisfying $K\subset E_n$.
  \end{proof}
  
  \begin{question}
  	Let $x$ be a point in the normed linear space $(E,||\cdot||)$. For any $\delta>0$, we have $B(x+y,\delta)=B(x,\delta)+y$. 
  \end{question}
  \begin{proof}
  	For any $z\in B(x+y,\delta)$, write $z=(z-y)+y=:v+y$. Note 
  	\begin{equation*}
  		||v-x||=||(v+y)-(x+y)||=||z-(x+y)||<\delta.
  	\end{equation*}
  	So $v\in B(x,\delta)$. The decomposition $z=v+y$ yields that $z\in B(x,\delta)+y$.
  	 On the contrary, for any $z\in B(x,\delta)+y$, there exists $u\in B(x,\delta)$ such that $z=u+y$. Note
  	\begin{equation*}
  		||z-(x+y)||=||(u+y)-(x+y)||=||u-x||<\delta.
  	\end{equation*}
  	So $z\in B(x+y,\delta)$. 
  \end{proof}
  
  \begin{question}
  	Let $A,B$ be two subsets of the normed linear space $(E,||\cdot||)$. Let $C:=A+B:=\{x+y:x\in A,y\in B\}$. Then $(\mathring{A}+B)\cup(A+\mathring{B})\subset \mathring{C}$.
  \end{question}
  \begin{proof}
  	Let $x\in \mathring{A}$, there exists $\delta>0$ such that $B(x,\delta)\subset A$. Fix $y\in B$. We claim $x+y\in \mathring{C}$. In fact, the previous {\bfseries QUESTION} shows $B(x+y,\delta)=B(x,\delta)+y\subset A+B=C$. Similarly, we can show $A+\mathring{B}\subset \mathring{C}$. Hence $(\mathring{A}+B)\cup(A+\mathring{B})\subset \mathring{C}$.
  \end{proof}
  
  \begin{question}
  	Assume $A\subset\r^n$ is convex. If $x\in \mathrm{Int}(A)$ and $y\in A$, then $\theta x+(1-\theta)y\in \mathrm{Int}(A)$ for any $\theta\in(0,1)$.
  \end{question}
  \begin{proof}
  	Since $x\in \mathrm{Int}(A)$, there exists $r>0$ such that $B(x,r)\subset A$. For any $x'\in B(x,r)$, there exists $\varepsilon\in[0,r]$ and a direction vector $\beta$ satisfying $||\beta||_{2,n}=1$, such that $x'=x+\varepsilon\beta$. The convexity of $A$ yields
  	\begin{equation*}
  		\begin{split}
  			\gamma x'+(1-\gamma)y&=\gamma(x+\varepsilon\beta)+(1-\gamma)y\\
  			&=\gamma x+(1-\gamma)y+\gamma\varepsilon\beta\in A,\quad \forall \gamma\in[0,1],\forall\varepsilon\in[0,r].
  		\end{split}\tag{$*$}
  	\end{equation*}
  	For any $\theta\in (0,1)$, consider $z=\theta x+(1-\theta)y$. We claim $B(z,\theta r)\subset A$. In fact, for any $z'\in B(z,\theta r)$, there exists $\delta\in[0,\theta r]$ and a direction vector $\beta$ satisfying $||\beta||_{2,n}=1$, such that 
  	\begin{equation*}
  		\begin{split}
  			z'&=z+\delta\beta=\theta x+(1-\theta)y+\delta\beta\\
  			&=\theta x+(1-\theta)y+\theta\(\frac{\delta}{\theta}\) \beta,\quad \frac{\delta}{\theta}\in[0,r].
  		\end{split}
  	\end{equation*}
  	Noticing $z'$ has the form just as ($*$), thus $z'\in A$. Namely, $B(z,\theta r)\subset A$. Finally, we have shown $\theta x+(1-\theta)y\in \mathrm{Int}(A)$ for any $\theta\in(0,1)$.
  \end{proof}

  \begin{question}[Continuity $\Leftrightarrow$ preservation of net limits]
	A \emph{directed set} is a nonempty set $A$ equipped with a reflexive and transitive relation $\le$ such that for every $\alpha,\beta\in A$ there exists $\gamma\in A$ with $\gamma \ge \alpha$ and $\gamma \ge \beta$. 
	
	Let $f:X\to Y$ be a map between topological spaces. The map $f$ is continuous if and only if for every net $(x_\alpha)$ in $X$, i.e., a function $x_\alpha : A \to X$, where $(A,\le)$ is a directed set,  and every $x\in X$,
\begin{equation*}
x_\alpha \to x \quad\Rightarrow\quad f(x_\alpha)\to f(x),
\end{equation*}
where the convergence of net in the topology space is defined as follows: A net $(x_\alpha)_{\alpha\in A}$ converges to $x\in X$, written as $x_\alpha \to x$,
if for every open neighborhood $U$ of $x$ (i.e. every $U\in\tau$ with $x\in U$), there exists $\alpha_0\in A$ such that 
\begin{equation*}
\alpha \ge \alpha_0 \;\Rightarrow\; x_\alpha \in U.
\end{equation*}
\end{question}


\begin{proof}
	\emph{($\Rightarrow$).}  
Assume $f$ is continuous, and let $x_\alpha\to x$ in $X$.  
Take any open neighborhood $W$ of $f(x)$ in $Y$. Since $f$ is continuous, $f^{-1}(W)$ is an open neighborhood of $x$ in $X$.  
By convergence of the net, there exists $\alpha_0$ such that
\begin{equation*}
\alpha \ge \alpha_0 \;\Rightarrow\; x_\alpha \in f^{-1}(W),
\end{equation*}
hence $f(x_\alpha)\in W$ for all $\alpha\ge\alpha_0$.  
This is exactly the statement that $f(x_\alpha)\to f(x)$.

\medskip

\emph{($\Leftarrow$).}  
Assume that $f$ preserves limits of all nets. Suppose, for contradiction, that $f$ is not continuous at some point $x\in X$. Then there exists an open set $V\subset Y$ with $f(x)\in V$ such that $f^{-1}(V)$ is not a neighborhood of $x$, i.e.\ every neighborhood $U$ of $x$ satisfies
\begin{equation*}
U\not\subset f^{-1}(V).
\end{equation*}
Equivalently, for every open $U\ni x$ there exists $x_U\in U$ with $f(x_U)\notin V$. 

Let $\mathcal{N}_x$ denote the set of all open neighborhoods of $x$. Make $\mathcal{N}_x$ a directed set by the order
\begin{equation*}
U\le W \iff U\supset W.
\end{equation*}
(Intersections of neighborhoods are again neighborhoods, so this is indeed directed.)

Define a net $(x_U)_{U\in\mathcal{N}_x}$ by choosing $x_U\in U$ with $f(x_U)\notin V$ for each $U$. Then $x_U\in U$ for each $U$ by construction.

We claim that $x_U\to x$ in $X$.  
Let $W$ be any neighborhood of $x$. Take $\alpha_0:=W$. If $U\ge\alpha_0$ in the above order, then $U\subset W$, hence $x_U\in U\subset W$.  
So $x_U\to x$.

By the assumption that $f$ preserves net limits, we must have $f(x_U)\to f(x)$ in $Y$. Since $V$ is a neighborhood of $f(x)$, there is $U_0$ such that
\begin{equation*}
U\ge U_0 \;\Rightarrow\; f(x_U)\in V.
\end{equation*}
But by construction, $f(x_U)\notin V$ for every $U$. This is a contradiction.  
Hence $f$ must be continuous.
\end{proof}




\begin{question}[Continuity $\Leftrightarrow$ preservation of sequential limits]
	Let $X$ be first-countable, i.e., for every $x\in X$ there exists a countable family of open sets $\{U_n(x)\}_{n\in\mathbb{N}}$, called a \emph{countable local base} (or neighborhood base) at $x$, such that for every open $U\ni x,\ \exists n\in\mathbb{N}$ with $U_n(x)\subset U$. And let $Y$ be any topological space. A map $f:X\to Y$ is continuous if and only if for every sequence $(x_n)$ in $X$ and $x\in X$,
\begin{equation*}
x_n \to x \quad\Rightarrow\quad f(x_n)\to f(x).
\end{equation*}
\end{question}

\begin{proof}
	\emph{($\Rightarrow$).} Same argument as for nets, with $A=\mathbb{N}$.

\medskip

\emph{($\Leftarrow$).} Assume $f$ preserves limits of all sequences. Suppose $f$ is not continuous at $x\in X$. Then there exists an open set $V\subset Y$ with $f(x)\in V$ such that $f^{-1}(V)$ is not a neighborhood of $x$. Thus, for every neighborhood $U\ni x$ we can choose $y_U\in U$ with $f(y_U)\notin V$.

Since $X$ is first-countable, there exists a countable local base $\{B_n\}_{n\in\mathbb{N}}$ at $x$. We can assume (by replacing $B_n$ with $\bigcap_{k=1}^n B_k$) that $B_1\supset B_2\supset\cdots$ and that $\{B_n\}$ is still a local base. For each $n$ choose $x_n\in B_n$ with $f(x_n)\notin V$.

Then one checks from the local base property that $x_n\to x$ in $X$, but $f(x_n)$ never enters $V$, so $f(x_n)\not\to f(x)$. This contradicts the preservation of sequential limits. Hence $f$ must be continuous.
\end{proof}


\begin{question}[Closure characterized by nets]
	Let $(X,\tau)$ be a topological space and $A\subset X$. Then for $x\in X$,
\begin{equation*}
x\in \overline{A}
\quad\Longleftrightarrow\quad
\exists\ \text{a net }(x_\alpha)_{\alpha\in A_0}\subset A\ \text{such that }x_\alpha\to x.
\end{equation*}
If $X$ is first-countable, then 
\begin{equation*}
x\in \overline{A}
\quad\Longleftrightarrow\quad
\exists\ \text{a sequence }(x_n)_{n\geq 1}\subset A\ \text{such that }x_n\to x.
\end{equation*}
Also, give an example to show that closure may not be characterized by sequence if $X$ is not first-countable.
\end{question}

\begin{proof}
	\begin{itemize}
		\item \emph{($\Leftarrow$).}  
If $x_\alpha\to x$ and $x_\alpha\in A$ for all $\alpha$, then for every neighborhood $U$ of $x$ there exists $\alpha_0$ such that $\alpha\ge\alpha_0$ implies $x_\alpha\in U$. Thus $U\cap A\neq\varnothing$, so $x\in\overline{A}$.

\medskip

\emph{($\Rightarrow$).}  
Assume $x\in\overline{A}$. Let $\mathcal{N}_x$ be the set of all open neighborhoods of $x$, ordered by reverse inclusion:
\begin{equation*}
U\le V \iff U\supset V.
\end{equation*}
Then $(\mathcal{N}_x,\le)$ is a directed set. For each $U\in\mathcal{N}_x$, the assumption $x\in\overline{A}$ implies $U\cap A\neq\varnothing$, so choose $x_U\in U\cap A$. The net $(x_U)_{U\in\mathcal{N}_x}$ lies in $A$.

We claim $x_U\to x$. Let $W$ be any neighborhood of $x$. Take $U_0:=W$. If $U\ge U_0$ (i.e.\ $U\subset W$), then $x_U\in U\subset W$. Hence $x_U\to x$.
		\item Based on the definition of first-countability, there exists a countable local base $\{U_n(x)\}_{n\geq 1}$. Following the same argument, we can find a sequence $(x_n)_{n\geq 1}$ satisfying $x_n\in U_n(x)\cap A$ for any $n\geq 1$. Then $x_n\to x$.
		\item Here is the example:
		Let $\{(X_i,\tau_i)\}_{i\in I}$ be a family of topological spaces, where $I$ is an arbitrary (possibly uncountable) index set.  
Define the Cartesian product
\begin{equation*}
X := \prod_{i\in I} X_i.
\end{equation*}
The \emph{product topology} on $X$ is the coarsest (smallest) topology $\tau_{\mathrm{prod}}$ such that every projection $\pi_i$ is continuous.  
Equivalently, $\tau_{\mathrm{prod}}$ is generated by the following basis $\mathcal{B}$:
\begin{equation*}
\mathcal{B} := \left\{
\prod_{i\in I} U_i \;:\;
U_i\in\tau_i\ \text{for all }i,\ \text{and }U_i=X_i\text{ for all but finitely many }i
\right\}.
\end{equation*}
Let $I$ be an uncountable set and consider $X := \{0,1\}^I$ with the product topology, where each factor $\{0,1\}$ is given the discrete topology. Define
\begin{equation*}
A := \left\{
x\in X : \{i\in I : x_i = 1\}\ \text{is finite}
\right\}.
\end{equation*}

\begin{itemize}
	\item First we show $\overline{A} = X$.

A basic open set in the product topology has the form
\begin{equation*}
U = \prod_{i\in I} U_i,
\end{equation*}
where $U_i\subset\{0,1\}$ is open for each $i$, and $U_i=\{0,1\}$ for all but finitely many $i$. Let $F\subset I$ be the finite set of coordinates where $U_i\neq\{0,1\}$.

Given any $x=(x_i)_{i\in I}\in X$ and any basic neighborhood $U$ of $x$, we have $x_i\in U_i$ for all $i\in F$. Define $y=(y_i)$ by
\begin{equation*}
y_i :=
\begin{cases}
x_i, & i\in F, \\
0,   & i\notin F.
\end{cases}
\end{equation*}
Then $y$ has at most $|F|<\infty$ many coordinates equal to $1$, so $y\in A$, and clearly $y\in U$. Therefore every basic neighborhood $U$ intersects $A$, so $\overline{A} = X$.

	\item Then we show that any sequential limit can only admit at most countable non-zero elements. Let $(x^{(n)})_{n\in\mathbb{N}}\subset A$ be a sequence converging to some $x\in X$. For each $n$, define the finite set
\begin{equation*}
S_n := \{ i\in I : x^{(n)}_i = 1 \}.
\end{equation*}
Let
\begin{equation*}
S := \bigcup_{n=1}^\infty S_n.
\end{equation*}
Then $S$ is a countable union of finite sets, hence $S$ is at most countable.

On the other hand, convergence in the product topology is coordinatewise: for each $i\in I$, the sequence $(x^{(n)}_i)$ converges to $x_i$ in the discrete space $\{0,1\}$. Thus
\begin{equation*}
x_i = 1 \;\Rightarrow\; \exists N\ \text{such that }n\ge N \Rightarrow x^{(n)}_i = 1.
\end{equation*}
In particular, if $x_i=1$ then $i\in S_n$ for all $n\ge N$, hence $i\in S$.  

Let $T := \{ i\in I : x_i = 1 \}$ be the support of $x$. The above shows $T \subset S$, so $T$ is at most countable. 
	\item Let $x^\ast\in X$ be the point with all coordinates equal to $1$. Then $\{i : x^\ast_i = 1\} = I$, which is uncountable by assumption. Therefore $x^\ast\in \overline{A}=X$ and there is no sequence $(x^{(n)})\subset A$ such that $x^{(n)}\to x^\ast$ based on what we have just proved.
	\item Finally, we show $X$ is not first-countable at $x^*$. Suppose, for contradiction, that there exists a countable local base $\{U_n\}_{n\in\mathbb{N}}$ at $x^\ast$. Since $x^*\in X=\overline{A}$, the same argument as the 2nd part of the proof shows that there exists a sequence $(x_n)_{n\geq 1}$ such that $x_n\to x^*$. However, as we have just shown, the support of $\lim x_n=x^*$ is at most countable, which is a contradiction.
\end{itemize}
	\end{itemize}
\end{proof}

\begin{question}
	Let $X$ be a topological space and let $K\subset X$. If $K$ is compact, then every net $(x_\alpha)_{\alpha\in A}$ with $x_\alpha\in K$ for all $\alpha$ admits a subnet that converges to a point of $K$.
\end{question}


\begin{proof}
	A point $x\in X$ is a \emph{cluster point} of the net $(x_\alpha)$ if for every neighborhood $U$ of $x$ and every index $\alpha_0\in A$, there exists $\alpha\ge \alpha_0$ such that $x_\alpha\in U$. 
	
	Let $(x_\alpha)_{\alpha\in A}$ be a net in $K$. We claim that $(x_\alpha)$ has at least one cluster point in $K$.
Suppose, for contradiction, that $(x_\alpha)$ has no cluster point in $K$.
Then for each $x\in K$ there exists a neighborhood $U_x$ of $x$ and an index $\alpha_x\in A$ such that for all $\alpha\ge \alpha_x$ we have $x_\alpha\notin U_x$.
The family $\{U_x:x\in K\}$ is an open cover of $K$, hence by compactness there exist $x_1,\dots,x_m\in K$ with
\begin{equation}
K\subset \bigcup_{i=1}^m U_{x_i}.
\end{equation}
Since $A$ is directed, there exists $\alpha^\ast\in A$ such that $\alpha^\ast\ge \alpha_{x_i}$ for every $i=1,\dots,m$.
For any $\alpha\ge \alpha^\ast$, we then have $x_\alpha\notin U_{x_i}$ for all $i$, hence $x_\alpha\notin \bigcup_{i=1}^m U_{x_i}$, contradicting $x_\alpha\in K$ and the inclusion above.
Therefore, $(x_\alpha)$ has a cluster point $x\in K$.

Let $x\in K$ be a cluster point of $(x_\alpha)$.
Let $\mathcal N(x)$ denote the set of neighborhoods of $x$, and define
\begin{equation}
B:=\{(U,\alpha_0):U\in\mathcal N(x),\ \alpha_0\in A\}.
\end{equation}
Equip $B$ with the directed order given by $(U_1,\alpha_1)\le (U_2,\alpha_2)$ if $U_2\subset U_1$ and $\alpha_2\ge \alpha_1$.
By the cluster point property, for each $(U,\alpha_0)\in B$ there exists $\alpha(U,\alpha_0)\ge \alpha_0$ such that $x_{\alpha(U,\alpha_0)}\in U$.
Define $\phi:B\to A$ by $\phi(U,\alpha_0)=\alpha(U,\alpha_0)$ and define a net $(y_\beta)_{\beta\in B}$ by $y_{(U,\alpha_0)}:=x_{\phi(U,\alpha_0)}$.
Then $(y_\beta)$ is a subnet of $(x_\alpha)$ because $\phi$ is cofinal: for any $\alpha_0\in A$, take $\beta_0=(X,\alpha_0)$, and for any $\beta=(U,\alpha)\ge \beta_0$ we have $\alpha\ge \alpha_0$ and hence $\phi(\beta)\ge \alpha\ge \alpha_0$. So, $(y_\beta)_{\beta\in B}$ is a valid subnet of $(x_\alpha)_{\alpha\in A}$.

Finally, we show $y_\beta\to x$.
Let $U\in\mathcal N(x)$ be any neighborhood. Take $\beta_0=(U,\alpha_0)$ for an arbitrary $\alpha_0\in A$.
If $\beta=(V,\alpha)\ge \beta_0$, then $V\subset U$, so
\begin{equation}
y_\beta=x_{\phi(V,\alpha)}\in V\subset U.
\end{equation}
Thus for every neighborhood $U$ of $x$, there exists $\beta_0$ such that $y_\beta\in U$ for all $\beta\ge \beta_0$, proving $y_\beta\to x$.
\end{proof}

\begin{question}
  	Let $X$ and $Y$ be topological spaces and let $f:X\to Y$ be a continuous mapping.
Assume that $A\subset X$ is dense in $X$, that is, $\overline{A}=X$.
Prove that $f(A)$ is dense in $f(X)$.
Equivalently, show that for every nonempty open set $U\subset Y$ satisfying $U\cap f(X)\neq\varnothing$, one has
\begin{equation*}
U\cap f(A)\neq\varnothing.
\end{equation*}
  \end{question}
  

\begin{proof}
Let $U\subset Y$ be a nonempty open set such that $U\cap f(X)\neq\varnothing$.
Choose $y\in U\cap f(X)$.
Then there exists $x\in X$ with $f(x)=y$.

Since $f$ is continuous and $U$ is open, the preimage $f^{-1}(U)$ is an open
subset of $X$.
Moreover, $x\in f^{-1}(U)$, hence $f^{-1}(U)\neq\varnothing$.

Because $A$ is dense in $X$, every nonempty open subset of $X$ intersects $A$.
Therefore,
\begin{equation*}
A\cap f^{-1}(U)\neq\varnothing.
\end{equation*}
Choose $a\in A\cap f^{-1}(U)$.
Then $f(a)\in U$ and $f(a)\in f(A)$, which implies
\begin{equation*}
U\cap f(A)\neq\varnothing.
\end{equation*}
Since $U$ was arbitrary, $f(A)$ is dense in $f(X)$.
\end{proof}

\begin{question}[Metrics generating the same topology via sequences]
Let $d_1$ and $d_2$ be two metrics on the same set $X$. Assume that for every sequence $(x_n)$ in $X$ and every $x\in X$,
\begin{equation*}
x_n \to x \text{ in } (X,d_1)\quad \Longleftrightarrow\quad x_n \to x \text{ in } (X,d_2).
\end{equation*}
Show that $d_1$ and $d_2$ generate the same topology on $X$, i.e., $\tau(d_1)=\tau(d_2)$.
\end{question}

\begin{proof}
We first record a sequential characterization of openness in metric spaces.

\noindent\textbf{Claim.} Let $(X,d)$ be a metric space and $U\subset X$. Then $U$ is open (with respect to $d$) if and only if for every $x\in U$ and every sequence $x_n\to x$ in $(X,d)$, there exists $N$ such that $x_n\in U$ for all $n\ge N$.

\noindent\emph{Proof of the claim.}
If $U$ is open and $x\in U$, then there exists $r>0$ such that the open ball $B_d(x,r)\subset U$. If $x_n\to x$, then for some $N$ we have $d(x_n,x)<r$ for all $n\ge N$, hence $x_n\in B_d(x,r)\subset U$ for all $n\ge N$.

Conversely, assume the sequential property holds but $U$ is not open. Then there exists $x\in U$ such that for every $m\in\mathbb{N}$, $B_d(x,1/m)\not\subset U$. Hence we can choose $x_m\in B_d(x,1/m)\setminus U$. Then $d(x_m,x)<1/m$, so $x_m\to x$, but $x_m\notin U$ for all $m$, contradicting the sequential property. Therefore $U$ is open. This proves the claim.

Now let $U\in\tau(d_1)$ be arbitrary. To show $U\in\tau(d_2)$, fix any $x\in U$ and any sequence $x_n\to x$ in $(X,d_2)$. By the hypothesis on sequential convergence, we also have $x_n\to x$ in $(X,d_1)$. Since $U$ is $d_1$-open, the claim implies that $x_n\in U$ eventually. Applying the claim again (now for $(X,d_2)$), we conclude that $U$ is $d_2$-open. Hence $\tau(d_1)\subset \tau(d_2)$. By symmetry, $\tau(d_2)\subset \tau(d_1)$, so $\tau(d_1)=\tau(d_2)$.
\end{proof}


\begin{question}[A Borel probability measure need not be tight on a non-Polish space]
Give an explicit topological space $X$ and a Borel probability measure $\mu$ on $X$ such that $\mu$ is not tight, i.e., there exists $\varepsilon>0$ for which no compact set $K\subset X$ satisfies $\mu(K)\ge 1-\varepsilon$.
\end{question}

\begin{proof}
Let $X$ be an uncountable set equipped with the \emph{co-countable topology}:
a set $U\subset X$ is open if and only if either $U=\varnothing$ or $X\setminus U$ is countable.
Let $\mathcal{B}(X)$ be the corresponding Borel $\sigma$-algebra.

Define $\mu:\mathcal{B}(X)\to[0,1]$ by
\begin{equation*}
\mu(A)=
\begin{cases}
0, & \text{if $A$ is countable},\\
1, & \text{if $X\setminus A$ is countable}.
\end{cases}
\end{equation*}
This is well-defined on $\mathcal{B}(X)$ (in particular, all countable and co-countable sets are Borel), and it is a probability measure: $\mu(\varnothing)=0$, $\mu(X)=1$, and countable additivity holds because a countable union of countable sets is countable, so one cannot obtain a co-countable set as a disjoint union of countably many countable sets.

Next we show that every compact subset of $X$ is finite. Suppose $K\subset X$ is infinite. For each $x\in K$, the set $U_x=X\setminus\{x\}$ is open since $\{x\}$ is countable. Then $\{U_x:x\in K\}$ is an open cover of $K$. Any finite subcollection $\{U_{x_1},\dots,U_{x_n}\}$ satisfies
\begin{equation*}
U_{x_1}\cup\cdots\cup U_{x_n}=X\setminus\{x_1,\dots,x_n\},
\end{equation*}
which misses the finite set $\{x_1,\dots,x_n\}\subset K$, hence it cannot cover $K$. Therefore $K$ is not compact. This proves that every compact set in $X$ must be finite.

Finally, any finite set is countable, so for every compact $K\subset X$ we have $\mu(K)=0$. Taking $\varepsilon=1/2$, there is no compact $K$ with $\mu(K)\ge 1-\varepsilon=1/2$. Hence $\mu$ is not tight.
\end{proof}

\begin{question}
Let $(X,\tau)$ be a topological space and let $Y\subset X$ be equipped with the subspace topology $\tau_Y$. 
Prove for $K\subset Y$, we have $K$ is compact in $(Y,\tau_Y)$ iff $\exists\, C\subset X$ compact in $(X,\tau)$ such that $ K=C\cap Y$
\end{question}

\begin{proof}
We prove the two implications.

\noindent\emph{($\Rightarrow$)} Assume that $K\subset Y$ is compact in $(Y,\tau_Y)$. 
We claim that $K$ is compact in $(X,\tau)$. 
Let $\{U_i\}_{i\in I}$ be an open cover of $K$ in $X$, i.e., each $U_i\in\tau$ and $K\subset\bigcup_{i\in I}U_i$. 
Then $\{U_i\cap Y\}_{i\in I}$ is a family of open sets in $Y$ (by definition of the subspace topology), and it still covers $K$:
\begin{equation*}
K\subset \bigcup_{i\in I}(U_i\cap Y).
\end{equation*}
Since $K$ is compact in $Y$, there exist indices $i_1,\dots,i_n$ such that
\begin{equation*}
K\subset \bigcup_{k=1}^n (U_{i_k}\cap Y)\subset \bigcup_{k=1}^n U_{i_k}.
\end{equation*}
Hence $\{U_{i_k}\}_{k=1}^n$ is a finite subcover of $K$ in $X$, proving that $K$ is compact in $X$.
Now take $C:=K$. Then $C$ is compact in $X$ and $K=C\cap Y$ (since $K\subset Y$), as desired.

\noindent\emph{($\Leftarrow$)} Conversely, assume that $C\subset X$ is compact in $(X,\tau)$ and set $K:=C\cap Y$. 
Let $\{V_j\}_{j\in J}$ be an open cover of $K$ in $Y$. 
By the definition of the subspace topology, for each $j$ there exists an open set $U_j\in\tau$ such that
\begin{equation*}
V_j = U_j \cap Y.
\end{equation*}
Since $\{V_j\}_{j\in J}$ covers $K$, we have $K\subset \bigcup_{j\in J}(U_j\cap Y)\subset \bigcup_{j\in J}U_j$. 
Define $U_0:=X\setminus K$, which is open in $X$. Then $\{U_0\}\cup\{U_j\}_{j\in J}$ is an open cover of $C$ in $X$, because $C\subset K\cup (X\setminus K)$. 
By compactness of $C$, there exist $j_1,\dots,j_m$ such that
\begin{equation*}
C \subset U_0 \cup \bigcup_{k=1}^m U_{j_k}.
\end{equation*}
Intersecting both sides with $Y$ and using $K=C\cap Y$ and $U_0\cap Y = (X\setminus K)\cap Y = Y\setminus K$, we obtain
\begin{equation*}
K \subset \bigcup_{k=1}^m (U_{j_k}\cap Y)=\bigcup_{k=1}^m V_{j_k}.
\end{equation*}
Thus a finite subcollection of $\{V_j\}_{j\in J}$ covers $K$, proving that $K$ is compact in $Y$.

Combining the two implications, $K\subset Y$ is compact in the subspace topology if and only if $K=C\cap Y$ for some compact $C\subset X$.
\end{proof}

\begin{question}
Let $\phi : X \to Y$ be a \emph{perfect} map between topological spaces, i.e.,
\begin{itemize}
  \item $\phi$ is continuous;
  \item $\phi$ is a closed map;
  \item for every $y \in Y$, the fiber $\phi^{-1}(\{y\})$ is compact.
\end{itemize}
Show that for every compact set $K \subset Y$, the inverse image $\phi^{-1}(K) \subset X$ is compact.
\end{question}

\begin{proof}
Let $K \subset Y$ be compact. Let $\mathcal U$ be an arbitrary open cover of $\phi^{-1}(K)$ in $X$.
We show that $\mathcal U$ admits a finite subcover.

Fix $y \in K$. Since $\phi^{-1}(\{y\})$ is compact and
\begin{equation*}
	\phi^{-1}(\{y\}) \subset \phi^{-1}(K) \subset \bigcup_{U \in \mathcal U} U,
\end{equation*}
there exist finitely many sets $U_{y,1}, \dots, U_{y,n(y)} \in \mathcal U$ such that
\begin{equation*}
\phi^{-1}(\{y\}) \subset \bigcup_{i=1}^{n(y)} U_{y,i}.
\end{equation*}
Define
\begin{equation*}
W_y := \bigcup_{i=1}^{n(y)} U_{y,i}.
\end{equation*}
Then $W_y$ is an open neighborhood of the fiber $\phi^{-1}(\{y\})$. Indeed, consider the closed set $X \setminus W_y$. Since $\phi$ is a closed map, $\phi(X \setminus W_y)$ is closed in $Y$. Define $V_y := Y \setminus \phi(X \setminus W_y)$.
Then $V_y$ is an open neighborhood of $y$ in $Y$. Indeed, if $y \notin V_y$, then
$y \in \phi(X \setminus W_y)$, so there exists $x \in X \setminus W_y$ with $\phi(x) = y$,
which contradicts $\phi^{-1}(\{y\}) \subset W_y$.

Moreover, by construction, $\phi^{-1}(V_y) \subset W_y$. The collection $\{V_y : y \in K\}$ is an open cover of $K$. Since $K$ is compact,
there exist points $y_1, \dots, y_m \in K$ such that
\begin{equation*}
K \subset \bigcup_{j=1}^m V_{y_j}.
\end{equation*}
Taking inverse images, we obtain
\begin{equation*}
\phi^{-1}(K)
\subset \bigcup_{j=1}^m \phi^{-1}(V_{y_j})
\subset \bigcup_{j=1}^m W_{y_j}.
\end{equation*}
Each $W_{y_j}$ is a finite union of elements of $\mathcal U$, hence
$\phi^{-1}(K)$ is covered by finitely many sets from $\mathcal U$.

Therefore $\phi^{-1}(K)$ is compact.
\end{proof}

 \section{Complex Analysis}
 \begin{question}
 	For a complex function $f(z)=u(x,y)+iv(x,y)$, we have
 	\begin{equation*}
 		\overline{\partial_{\overline{z}}f}=\partial_{z}\overline{f}.
 	\end{equation*}
 	Thus we have
 	\begin{equation*}
 		\partial^2_{z\overline{z}}f=\overline{\partial^2_{z\overline{z}}\overline{f}}.
 	\end{equation*}
 \end{question}
 \begin{proof}
 	Note 
 	\begin{equation*}
 		\partial_z=(\partial_x-i\partial_y)/2,\quad \partial_{\overline{z}}=(\partial_x+i\partial_y)/2.
 	\end{equation*}
 	Thus we have 
 	\begin{equation*}
 		\overline{\partial_{\overline{z}}f}=\overline{ (\partial_x+i\partial_y)(u+iv)/2}=(\partial_x-i\partial_y)(u-iv)/2=\partial_{z}\overline{f}.
 	\end{equation*}
 	Therefore $\overline{\partial_{\overline{z}}f}=\partial_{z}\overline{f}$. Finally,
 	\begin{equation*}
 		\partial^2_{z\overline{z}}f=\partial_z(\partial_{\overline{z}}f)=\partial_z\(\overline{ \partial_{z}\overline{f}}\)=\overline{ \partial_{\overline{z}}(\partial_z\overline{f})}=\overline{\partial^2_{z\overline{z}}\overline{f}}.
 	\end{equation*}
 \end{proof}
  
  \begin{question}
  	Let $z=re^{i\theta}$ and $f(z)=u(r,\theta)+iv(r,\theta)$. Prove the Cauchy-Riemann equation is equivalent to
  	\begin{equation*}
  		\partial_ru=\partial_\theta v/r,\quad \partial_rv=-\partial_\theta u/r.\tag{$*$}
  	\end{equation*}
  \end{question}
  \begin{proof}
 	Note $x=r\cos \theta$ and $y=r\sin \theta$. Then
 	\begin{equation*}
 		\left\{
 		\begin{split}
 			\partial_ru&=\cos\theta\partial_xu+\sin\theta\partial_yu,\\
 			\partial_\theta u&=-r\sin\theta\partial_xu+r\cos\theta\partial_yu.
 		\end{split}\right.
 	\end{equation*}
 	Similarly, we can replace all the $u$ in above formula into $v$ to obtain the formula for $v$. Apply the Cauchy-Riemann equation 
 	\begin{equation*}
 		\partial_xu=\partial_yv,\quad \partial_yu=-\partial_xv,
 	\end{equation*}
 	we will get
 	\begin{equation*}
 		\left\{
 		\begin{split}
 			\partial_ru&=\cos\theta\partial_yv-\sin\theta\partial_xv=\partial_\theta v/r,\\
 			\partial_rv&=-\cos\theta\partial_yu+\sin\theta\partial_xu=-\partial_\theta u/r.
 		\end{split}\right.
 	\end{equation*}
 	On the contrary, if $(*)$ holds, we will obtain
 	\begin{equation*}
 		\left\{
 		\begin{split}
 			\partial_ru&=\cos\theta\partial_xu+\sin\theta\partial_yu=-\sin\theta\partial_xv+\cos\theta\partial_yv=\partial_\theta v/r ,\\
 			\partial_rv&=\cos\theta\partial_xv+\sin\theta\partial_yv=\sin\theta\partial_xu-\cos\theta\partial_yu=-\partial_\theta u/r.
 		\end{split}\right.
 	\end{equation*}
 	It is equivalent to
 	\begin{equation*}
 		\left\{\begin{split}
 			(\partial_xu-\partial_yv)\cos\theta+(\partial_yu+\partial_xv)\sin\theta&=0,\\
 			(\partial_xu-\partial_yv)(-\sin\theta)+(\partial_yu+\partial_xv)\cos\theta&=0.
 		\end{split}\right.
 	\end{equation*}
 	Observe the homogeneous equations
 	\begin{equation*}
 		\begin{pmatrix}
 			\cos\theta & \sin\theta\\
 			-\sin\theta & \cos\theta
 		\end{pmatrix}
 		\begin{pmatrix}
 			x\\
 			y
 		\end{pmatrix}
 		=
 		\begin{pmatrix}
 			0\\
 			0
 		\end{pmatrix}
 	\end{equation*}
 	has only zero solution. So we must have 
 	\begin{equation*}
 		\partial_xu-\partial_yv=0,\quad \partial_yu+\partial_xv=0.
 	\end{equation*}
 	Namely, the Cauchy-Riemann equation holds.
 \end{proof}
 
 \begin{question}
 	Let $r=re^{i\theta}$. Prove
 	\begin{equation*}
 		\partial_{\overline{z}}=e^{i\theta}(\partial_r+i/r\partial_\theta)/2,\quad \partial_z=e^{-i\theta}(\partial_r-i/r\partial_\theta)/2.
 	\end{equation*}
 	Hence for any $f\in H(D)$, we have 
 	\begin{equation*}
 		\begin{split}
 			f'=\partial_z f&= r/z(\partial_ru+i\partial_ru)=r/z(\partial_\theta v/r+i\partial_ru)\\
 			&=r/z(\partial_ru-i\partial_\theta u/r)=r/z(\partial_\theta v/r-i\partial_\theta u/r).
 		\end{split}
 	\end{equation*}
 \end{question}
 \begin{proof}
 	In the previous question we have obtained 
 	\begin{equation*}
 		\partial_r=\cos\theta\partial_x+\sin\theta\partial_y,\quad \partial_\theta=-r\sin\theta\partial_x+r\cos\theta\partial_y.
 	\end{equation*}
 	Thus
 	\begin{equation*}
 		\begin{split}
 			e^{i\theta}(\partial_r+i/r\partial_\theta)/2&=(\cos\theta+i\sin\theta)(\partial_r+i/r\partial_\theta)/2\\
 			&=[(\cos\theta\partial_r-\sin\theta/r\partial_\theta)+i(\cos\theta/r\partial_\theta+\sin\theta\partial_r)]/2\\
 			&=(\partial_x+i\partial_y)/2=\partial_{\overline{z}}
 		\end{split}
 	\end{equation*}
 	and
 	\begin{equation*}
 		\begin{split}
 			e^{-i\theta}(\partial_r-i/r\partial_\theta)/2&=(\cos\theta-i\sin\theta)(\partial_r-i/r\partial_\theta)/2\\
 			&=[(\cos\theta\partial_r-\sin\theta/r\partial_\theta)-i(\cos\theta/r\partial_\theta+\sin\theta\partial_r)]/2\\
 			&=(\partial_x-i\partial_y)/2=\partial_{z}.
 		\end{split}
 	\end{equation*}
 	Therefore, we have
 	\begin{equation*}
 		\begin{split}
 			f'&=\partial_zf=e^{-i\theta}(\partial_r-i/r\partial_\theta)(u+iv)/2\\
 			&=e^{-i\theta}[(\partial_ru+\partial_\theta v/r)+i(\partial_rv-\partial_\theta u/r)]/2.
 		\end{split}
 	\end{equation*}
 	Apply the polar coordinate form Cauchy-Riemann equation 
 	\begin{equation*}
 		\partial_ru=\partial_\theta v/r,\quad \partial_rv=-\partial_\theta u/r,
 	\end{equation*}
 	we will get
 	\begin{equation*}
 		\begin{split}
 			f'=\partial_z f&= r/z(\partial_ru+i\partial_ru)=r/z(\partial_\theta v/r+i\partial_ru)\\
 			&=r/z(\partial_ru-i\partial_\theta u/r)=r/z(\partial_\theta v/r-i\partial_\theta u/r).
 		\end{split}
 	\end{equation*}
 \end{proof}
 
 \begin{question}
 	Suppose $\gamma: I\to\c$ is a rectifiable curve and $\varphi$ is holomorphic on $\gamma$. Write $\Gamma=\varphi(\gamma)$. Prove 
 	\begin{enumerate}
 		\item $\Gamma$ is also a rectifiable curve. 
 		\item If $f$ is continuous on $\Gamma$, then
 		\begin{equation*}
 			\int_\Gamma f(w)dw=\int_\gamma f(\varphi(z))\varphi'(z)dz.
 		\end{equation*}
 		\item If $\gamma$ is a smooth curve (i.e. $\gamma\in C^1(I)$ and $\gamma'\not\equiv 0$), then $\Gamma$ is also a smooth curve and the above formula also holds.
 	\end{enumerate}
 \end{question}
 \begin{proof}
 	\begin{enumerate}
 		\item Since $\varphi$ is holomorphic on $\gamma$ and $\gamma$ is compact, then exists $M>0$ such that $|\varphi'(w)|\leq M,\forall w\in\gamma(I)$. So for any division $0=t_0<t_1<\cdots<t_{n}<t_{n+1}=1$, write $\gamma_k^{k+1}$ for the sub-curve $\gamma:[t_k,t_{k+1}]\to\c$ and write $|\gamma|$ for the length of curve $\gamma$, then we have
 		\begin{equation*}
 			\begin{split}
 				\sum_{k=0}^n |\Gamma(t_{k+1})&-\Gamma(t_k)|=\sum_{k=0}^n |\varphi(\gamma (t_{k+1}))-\varphi(\gamma (t_{k}))|\\
 				&=\sum_{k=0}^n \left|\int_{\gamma(t_{k})}^{\gamma(t_{k+1})}\varphi'(z)dz\right|\leq \sum_{k=0}^n \int_{\gamma_{k}^{k+1}}|\varphi'(z)|ds\\
 				&\leq \sum_{k=0}^nM|\gamma_{k}^{k+1}|=M|\gamma|.
 			\end{split}
 		\end{equation*}
 		Therefore $\Gamma$ is also rectifiable. 
 		\item For any $\varepsilon>0$, there exists $\delta>0$ and a division of $I$ as $0=t_0<t_1<\cdots<t_{n-1}<t_n=1$ satisfying $\lambda:=\max_{1\leq k\leq n}|t_k-t_{k-1}|<\delta$ such that
 		\begin{equation*}
 			\Bigg|\int_\Gamma f(w)dw-\sum_{k=1}^nf(\varphi(\gamma(t_{k-1})))[\varphi(\gamma(t_k) )-\varphi(\gamma(t_{k-1}) )]\Bigg|<\varepsilon
 		\end{equation*}
 		and
 		\begin{equation*}
 			\begin{split}
 				\Bigg|\int_\gamma f(\varphi(z))&\varphi'(z)dz\\
 				&-\sum_{k=1}^nf(\varphi(\gamma(t_{k-1})))\varphi'(\gamma(t_{k-1}))[\gamma(t_k)-\gamma(t_{k-1})]\Bigg|<\varepsilon.
 			\end{split}
 		\end{equation*}
 		Moreover, we can let $\delta>0$ small enough, apply the uniform continuity of $\varphi'$ on $\gamma(I)$, such that for any $1\leq k\leq n$,
 		\begin{equation*}
 			|\varphi'(w_1)-\varphi'(w_2) |<\varepsilon,\quad \forall  w_1,w_2\in \gamma_{k-1}^k.
 		\end{equation*}
 		Therefore
 		\begin{equation*}
 			\begin{split}
 				&\Bigg|\sum_{k=1}^nf(\varphi(\gamma(t_{k-1})))[\varphi(\gamma(t_k) )-\varphi(\gamma(t_{k-1}) )]\\
 				&\qquad\qquad -\sum_{k=1}^nf(\varphi(\gamma(t_{k-1})))\varphi'(\gamma(t_{k-1}))[\gamma(t_k)-\gamma(t_{k-1})]\Bigg|\\
 				= &\Bigg|\sum_{k=1}^nf(\varphi(\gamma(t_{k-1})))\int_{\gamma_{k-1}^k}\varphi'(z)dz\\
 				&\qquad\qquad -\sum_{k=1}^nf(\varphi(\gamma(t_{k-1})))\int_{\gamma_{k-1}^k}\varphi'(\gamma(t_{k-1}) )dz\Bigg|\\
 				\leq &\sum_{k=1}^nf(\varphi(\gamma(t_{k-1})))\int_{\gamma_{k-1}^k}|\varphi'(z)-\varphi'(\gamma(t_{k-1}))| ds\\
 				<&\varepsilon\sum_{k=1}^n|f(\varphi(\gamma(t_{k-1})))||\gamma_{k-1}^k|\\
 				\leq & \(\max_{z\in \Gamma}|f(z)||\gamma |\)\varepsilon. 
 			\end{split}
 		\end{equation*}
 		So we have obtained the desired result.
 		\item $\Gamma\in C^1(I)$ is Obvious. Moreover, since $\varphi$ is holomorphic on $\gamma$, by the uniqueness of holomorphic function, $\varphi'(w)\not\equiv 0,\forall w\in\gamma(I)$. Therefore $\Gamma'(t)=\varphi'(\gamma(t))\gamma'(t)\not\equiv 0,\forall t\in I$. So $\Gamma$ is also a smooth curve. Then
 	\begin{equation*}
 		\begin{split}
 			\int_\Gamma f(w)dw&=\int_I f(\Gamma(t))\Gamma'(t)dt=\int_If(\varphi(\gamma(t)))\varphi'(\gamma(t))\gamma'(t)dt\\
 			&=\int_\gamma f(\varphi(z))\varphi'(z)dz.
 		\end{split}
 	\end{equation*}
 	\end{enumerate}
 \end{proof}

  \section{Probability}
  \begin{question}
   Suppose $\xi,\eta$ are two identically distributed random variables on $(\Omega ,\f,\p )$ and sub $\sigma$-algebra $\a\subset\f$. Use an example to show
   \begin{equation*}
    \p(\xi|\a )=\p(\eta|\a )
   \end{equation*}
   is not true. However, if $\xi=\eta,\p$-a.s., then $\p(\xi|\a )=\p(\eta|\a )$.
  \end{question}
  \begin{proof}
   Consider $\xi\sim N(0,1)$ and let $\eta=-\xi$, then $\eta\sim N(0,1)$. Moreover, let $\a=\sigma(\xi)$, we will have
   \begin{equation*}
    \p(\xi|\xi )=\xi\neq\eta=-\xi=-\p(\xi|\xi)=\p(-\xi|\xi)=\p(\eta|\xi ).
   \end{equation*}
   This is a smart example. If $\xi=\eta,\p$-a.s., then for any $A\in\a$, we have $\p(\xi1_A)=\p(\eta1_A)$. This shows $\p(\xi|\a )=\p(\eta|\a )$.
  \end{proof}

  \begin{question}
   Consider the probability space $(\Omega,\f,\p)$ and the sub $\sigma$-algebra $\a_1,\a_2$. Suppose $\xi$ is a random variable. Use an example to show
   \begin{equation*}
    \p\[\p(\xi|\a_1 )|\a_2\]=\p(\xi|\a_1\cap\a_2 )
   \end{equation*}
   is not true.
  \end{question}
  \begin{proof}
   Consider $\Omega=\{0,1,2\},\f=2^\Omega,\xi=\i_{\{0\}}$,
   \begin{equation*}
    \a_1=\{\{1\},\{0,2\},\Omega ,\emptyset \},\quad \a_2=\{\{2\},\{0,1\},\Omega ,\emptyset \},
   \end{equation*}
   and $\p(\{0\})=1/2,\p(\{1\})=\p(\{2\})=1/4$. Observe $\a_1\cap\a_2=\{\Omega,\emptyset \}$, then $\xi$ is independent of $\a_1\cap \a_2$ and hence
   \begin{equation*}
    \p(\xi|\a_1\cap\a_2)=\p(\xi)=1/2.
   \end{equation*}
   Moreover, notice $\{1\}$ and $\{0,2\}$ are two atoms in $\a_1$, so
   \begin{equation*}
    \begin{split}
     \p(\xi|\a_1)&=\p(\xi|\{1\} )\mathbbm{1}_{\{1\}}+\p(\xi|\{0,2\})\mathbbm{1}_{\{0,2\}}\\
     &=\p(\{0\}|\{1\})\mathbbm{1}_{\{1\}}+\p(\{0\}|\{0,2\} )\mathbbm{1}_{\{0,2\}}\\
     &=0\mathbbm{1}_{\{1\}}+\frac{\p(\{0\})}{\p(\{0,2\})}\mathbbm{1}_{\{0,2\}}\\
     &=\frac{1/2}{1/2+1/4}\mathbbm{1}_{\{0,2\}}=2/3\mathbbm{1}_{\{0,2\}}.
    \end{split}
   \end{equation*}
   Note $\{2 \}$ and $\{0,1\}$ are two atoms in $\a_2$, so
   \begin{equation*}
    \begin{split}
     \p\[\p(\xi|\a_1 )|\a_2\]&=\p\(2/3\mathbbm{1}_{\{0,2\}}\big|\a_2 \)\\
     &=2/3\[\p\(\{0,2\}|\{2\} \)\mathbbm{1}_{\{2\}}+\p\(\{0,2\}|\{0,1\} \)\mathbbm{1}_{\{0,1\}}\]\\
     &=\frac{2}{3} \(\mathbbm{1}_{\{2\}}+\frac{\p(\{0\})}{\p(\{0,1\})}\mathbbm{1}_{\{0,1\}} \)\\
     &=\frac{2}{3} \(\mathbbm{1}_{\{2\}}+\frac{1/2}{1/2+1/4}\mathbbm{1}_{\{0,1\}} \)\\
     &=\frac{2}{3} \(\mathbbm{1}_{\{2\}}+\frac{2}{3}\mathbbm{1}_{\{0,1\}} \)\\
     &=\frac{2}{3} \mathbbm{1}_{\{2\}}+\frac{4}{9}\mathbbm{1}_{\{0,1\}}.
    \end{split}
   \end{equation*}
   Hence
   \begin{equation*}
    \p(\xi|\a_1\cap\a_2)=\frac{1}{2}\neq \frac{2}{3} \mathbbm{1}_{\{2\}}+\frac{4}{9}\mathbbm{1}_{\{0,1\}}=\p\[\p(\xi|\a_1 )|\a_2\].
   \end{equation*}
  \end{proof}

  \begin{question}
   \begin{enumerate}
    \item Consider the probability space $(\Omega,\f,\p)$. Given $A,B$ and $C\in\f$, if $A$ is independent to $B,C$ or $C$ is independent to $A,B$, then $A$ and $C$ are independent conditional on $B$.
    \item Suppose conditional on $B,A$ and $C$ are conditionally independent, namely, $\p(A\cap C|B)=\p(A|B)\p(C|B)$, then $\p(A|B\cap C)=\p(A|B)$.
   \end{enumerate}
  \end{question}
  \begin{proof}
   \begin{enumerate}
    \item Without loss of generality, we assume $A$ is independent to $B,C$. Then we have
    \begin{equation*}
     \p(A\cap B|C )=\frac{\p(A\cap B\cap C)}{\p(C)}=\frac{\p(A)\p(B\cap C)}{\p(C)}
    \end{equation*}
    and
    \begin{equation*}
     \begin{split}
      \p(A|C)\p(B|C)&=\frac{\p(A\cap C)\p(B\cap C)}{[\p(C)]^2}=\frac{\p(A)\p(C)\p(B\cap C)}{[\p(C)]^2}\\
      &=\frac{\p(A)\p(B\cap C)}{\p(C)}.
     \end{split}
    \end{equation*}
    \item We have
   \begin{equation*}
    \begin{split}
     \p(A|B\cap C)&=\frac{\p(A\cap B\cap C)}{\p(B\cap C)}=\frac{\p(A\cap C|B)\p(B)}{\p(B\cap C)}\\
     &=\frac{\p(A|B)\p(C|B)\p(B)}{\p(B\cap C)}=\p(A|B).
    \end{split}
   \end{equation*}
   \end{enumerate}
  \end{proof}

  \begin{question}
   Suppose $A$ is a subset on $\r^n$. If $D\subset A$ and $dx(D)=dx(A)$, where $dx$ is the Lebesgue measure, then $D$ is dense in $A$.
  \end{question}
  \begin{proof}
   For all $x\in A$ and for all $\varepsilon>0$, consider the neighborhood $B(x,\varepsilon )\subset A$. We claim there must exist $a\in A$, such that $a\in B(x,\varepsilon )$. If not, then $D\subset B(x,\varepsilon )^c$. Hence $dx(D)\leq dx(A)-n\alpha(n)\varepsilon^n<dx(A)$, which is a contradiction.
  \end{proof}

  \begin{question}
   Let
   \begin{equation*}
    \begin{split}
     G_1&=\(\frac{1}{3},\frac{2}{3} \),\\
     G_2&=\(\frac{1}{3^2},\frac{2}{3^2} \)\bigcup\(\frac{2}{3} +\frac{1}{3^2},\frac{2}{3}+ \frac{2}{3^2} \) ,\\
     &\vdots\\
     G_n&=\bigcup\left\{ \(\sum_{k=1}^{n-1} \frac{a_k}{3^k}+\frac{1}{3^n} ,\sum_{k=1}^{n-1} \frac{a_k}{3^k}+\frac{2}{3^n} \): a_k=0 \text{ or }2 \right\} .\\
    \end{split}
   \end{equation*}
   Then $G:=\cup_n G_n$ is an open subset of $[0,1]$. The closed set $C=[0,1]-G$ is called the {\normalfont \itshape Cantor set}. Let
   \begin{equation*}
    F_0(u)=\sum_{k=1}^{n-1}\frac{a_k}{2^k}+\frac{1}{2^n},\quad \sum_{k=1}^{n-1}\frac{a_k}{3^k}+\frac{1}{3^n}<u<\sum_{k=1}^{n-1}\frac{a_k}{3^k}+\frac{2}{3^n}.
   \end{equation*}
   We can extend $F$ to a probability distribution function on $\r$ by setting
   \begin{equation*}
    F(x):=
    \begin{cases}
     1&,\mbox{if }x\geq 1\\
     \inf\{F_0(y):y\in G \text{ and }y>x\} &,\mbox{if }x<1.\\
    \end{cases}
   \end{equation*}
   Show $F$ is uniformly continuous on $\r$.
  \end{question}
  \begin{proof}
   We first prove $F_0$ is uniformly continuous on $G$. Observe $F_0$ is increasing on $G$ and for all $n\in\n$, the distances between each pair of adjacent open intervals in $\cup_{k=1}^n G_k$ equal to $1/3^{n}$. By the definition of $F_0$, the variation in each closed set on $[0,1]-\cup_{k=1}^nG_k$ is $1/2^n$. Thus for every $x,y\in [0,1]$ and $|x-y|<1/3^n$, there are three situations might happens:
   \begin{enumerate}
    \item $x$ and $y$ all lie in one of open intervals in $\cup_{k=1}^n G_k$. Since $F_0$ stays constant in the open interval, we have $|F_0(x)-F_0(y)|=0$.
    \item $x$ lies in one of open intervals in $\cup_{k=1}^n G_k$ and $y$ lies in one of closed intervals in $[0,1]-\cup_{k=1}^n G_k$, we have $|F_0(x)-F_0(y)|<1/2^n$.
    \item $x$ and $y$ are all lie in one of closed intervals in $[0,1]-\cup_{k=1}^n G_k$, we also have $|F_0(x)-F_0(y)|<1/2^n$.
   \end{enumerate}
   In short, we know if $x,y\in G$ and $|x-y|<1/3^n$, then $|F_0(x)-F_0(y)|<1/2^n$, which means $F$ is uniformly continuous on $G$. Since $dx(G)=1$, by {\bfseries QUESTION 1.4} and {\bfseries QUESTION 3.3}, we know $G$ is dense in $[0,1]$ and $F$ is uniformly continuous on $\r$.
  \end{proof}

  \begin{question}
   Suppose $\b$ is the Borel $\sigma$-algebra on $\r$. Then $\b[0,1]=\b\cap [0,1]$.
  \end{question}
  \begin{proof}
   Write $\mathscr{O}$ for all the open set on $\r$, then $\b=\sigma(\mathscr{O})$. We can rewrite $\b[0,1]=\b\cap [0,1]$ as $\sigma(\mathscr{O}\cap [0,1] )=\sigma(\mathscr{O})\cap [0,1].$
   Observe $\mathscr{O}\cap [0,1]\subset \sigma(\mathscr{O})\cap [0,1]$ and $\sigma(\mathscr{O})\cap [0,1]$ is a $\sigma$-algebra on $[0,1]$, therefore $\sigma(\mathscr{O}\cap [0,1] )\subset\sigma(\mathscr{O})\cap [0,1]$. On the contrary, consider
   \begin{equation*}
    \a=\{A\in\sigma(\mathscr{O}):A\cap[0,1]\in  \sigma(\mathscr{O}\cap [0,1] ) \}.
   \end{equation*}
   Obviously $\mathscr{O}\subset\a$ and it is easy to verify $\a$ is a $\sigma$-algebra. Thus $\forall A\in \sigma(\mathscr{O}),A\cap[0,1]\in \sigma(\mathscr{O}\cap [0,1] )$, i.e., $\sigma(\mathscr{O}\cap [0,1] )\supset \sigma(\mathscr{O})\cap [0,1]$.
  \end{proof}

  \begin{question}
   Suppose $(\Omega,\f )$ and $(E,\mathscr{E})$ are two measurable spaces and $\{X_i\}_{i=1}^n$ are measurable mappings with respect to $\f$ and $\mathscr{E}$. Consider the measurable vector $(X_1,\cdots,X_n ):(\Omega,\f )\to(E^n,\e^n)$, then we have
   \begin{equation*}
    \sigma(X_1,\cdots,X_n)=\sigma\(\bigcup_{i=1}^n\sigma(X_i) \).
   \end{equation*}
  \end{question}
  \begin{proof}
   For each $1\leq i\leq n$ and $B_i\in\e$, we have
   \begin{equation*}
    \{X_i\in B_i\}=\{(X_1,\cdots,X_n)\in E\times\cdots\times B_i\times\cdots\times E \}.
   \end{equation*}
   Then $\sigma(X_i)\subset\sigma(X_1,\cdots,X_n)\Longrightarrow \cup_{i=1}^n \sigma(X_i)\subset\sigma(X_1,\cdots,X_n)$. Thus we have $\sigma(X_i,1\leq i\leq n)\subset \sigma(X_1,\cdots,X_n)$. On the contrary, the above equality means
   \begin{equation*}
    \{(X_1,\cdots,X_n)\in B_1\times\cdots\times B_n \}=\bigcap_{i=1}^n\{X_i\in B_i\}\in \sigma\(\bigcup_{i=1}^n\sigma(X_i) \).
   \end{equation*}
   Consider the following set:
   \begin{equation*}
    \a:=\{B\in\e^n:\{(X_1,\cdots,X_n)\in B\}\in \sigma(X_i,1\leq i\leq n)\}.
   \end{equation*}
   It is easy to prove $\a$ is a $\sigma$-algebra and the previous formula shows all the measurable rectangles are contained in $\a$. Then $\a=\e^n$, namely, $\sigma(X_1,\cdots,X_n)\subset \sigma(X_i,1\leq i\leq n)$.
  \end{proof}

  \begin{question}
   Suppose $(X_t)_{t\in I}$ is a sequence of random variables from the measurable space $(\Omega,\f )$ to $(E,\e)$, where $I$ is a index set. Then
   \begin{equation*}
    \begin{split}
     &\sigma(X_t:t\in I)\\
     =&\sigma\(\{\{X_{t_1}\in A_1,\cdots,X_{t_n}\in A_n\}:A_1,\cdots,A_n \in \e,(t_1,\cdots,t_n )\in S(I) \} \)\\
     =&:\a.
    \end{split}
   \end{equation*}
   where $S(I)$ denotes all the ordered finite subset of $I$.
  \end{question}
  \begin{proof}
   Define $\mathscr{C}=\{A_1\times\cdots\times A_n:A_1,\cdots,A_n \in \e \}$, then $\sigma(\mathscr{C})=\e^n$. Obviously, for each $A_1,\cdots,A_n \in \e,(t_1,\cdots,t_n )\in S(I)$, we have
   \begin{equation*}
    \{X_{t_1}\in A_1,\cdots,X_{t_n}\in A_n\}=\bigcap_{k=1}^n\{X_{t_k}\in A_k\}\in \sigma(X_t:t\in I).
   \end{equation*}
   Thus $\a\subset \sigma(X_t:t\in I)$. On the contrary, for each $t\in I$, obviously $\sigma(X_t)\subset \a\Longrightarrow \cup_{t\in I}\sigma(X_t)\subset\a$ and therefore $\sigma(X_t:t\in I)\subset\a$.
  \end{proof}
  
  \begin{question}
  	Suppose $(X_t^i)_{t\in I},i=1,\cdots,d$ are $d$ sequences of random variables from the measurable space $(\Omega,\f )$ to $(E,\e)$. Then 
  	\begin{equation*}
    \begin{split}
     &\sigma(X_t^i:t\in I,i=1,\cdots,d)\\
     =&\sigma\(\left\{\bigcup_{i=1}^d \{X_{t_1}^i\in A_1^i,\cdots,X_{t_n}^i\in A_n^i\}:A_1^i,\cdots,A_n^i \in \e,(t_1,\cdots,t_n )\in S(I) \right\} \).
    \end{split}
   \end{equation*}
  \end{question}
  \begin{proof}
  	The proof is similar to the previous question.
  \end{proof}

  \begin{question}
   Suppose $(\Omega,\f)$ and $(E,\mathscr{E})$ are two measurable spaces. Then $X:(\Omega,\f )\to (E,\mathscr{E})$ is measurable iff $f(X):(\Omega,\f)\to (\r,\b(\r))$ is measurable for each $f\in\mathrm{b}\mathscr{E}$.
  \end{question}
  \begin{proof}
   If $f(X)\in \f$ for each $f\in\mathrm{b}\mathscr{E}$, then for all $A\in \mathscr{E}$, choose $f=\i_{A}$, we have $f(X)=\i_{\{X\in A\}}\in\f$. Since  $\{1\}\in\b(\r)$, then
   \begin{equation*}
    \{X\in A\}=\{\i_{\{X\in A\}}=1\}\in \f,
   \end{equation*}
   which means $X\in\f$. On the contrary, if $X\in\f$, obviously $f(X)\in\f$ for each $f\in\mathrm{b}\mathscr{E}$.
  \end{proof}

  \begin{question}
   Suppose $\f$ and $\g$ are two $\sigma$ algebras on $\Omega$. Then $\f\subset\g$ iff for each $F\in\mathrm{b} \f$, we have $F\in\mathrm{b}\g$.
  \end{question}
  \begin{proof}
   If $\f\subset\g$, then obviously  $F\in\mathrm{b} \f$ implies $F\in\mathrm{b}\g$. On the contrary, if for each $F\in\mathrm{b} \f$, we have $F\in\mathrm{b}\g$, we can choose $F=1_A$ for each $A\in\f$. Then $1_A\in \mathrm{b}\g$ and thus $A\in\g$. So $\f\subset\g$.
  \end{proof}

  \begin{question}
   Suppose $(\Omega,\f,\p )$ is a probability space and $X$ is an integral random variable on $(\Omega,\f,\p )$. Assume $\g\subset\f$ is also a $\sigma$-algebra, then $Z=\p(X|\g)$ iff $Z\in\g$ and for each measurable function $F\in\mathrm{b} \g$,
   \begin{equation*}
    \p(XF)=\p(ZF).
   \end{equation*}
  \end{question}
  \begin{proof}
   If the above equality holds, we can choose $F=1_A$ for each $A\in\g$, the $Z=\p(X|\g)$ follows. On the contrary, consider the following set:
   \begin{equation*}
    \l :=\{F\in\mathrm{b} \g:\p(XF)=\p(ZF) \}.
   \end{equation*}
   By the definition of $\p(X|\g)$, we know $\{1_A:A\in\g\}\subset\l$. If $F,G\in\l$ and $\alpha,\beta\in\r$, we have
   \begin{equation*}
    \begin{split}
     \p[\p(X|\g)(\alpha F&+\beta G) ]=\alpha \p[\p(X|\g)F]+\beta\p[\p(X|\g)G]\\
     &=\alpha \p(XF)+\beta\p(XG)=\p[X(\alpha F+\beta G)].
    \end{split}
   \end{equation*}
   Moreover, if $\{F_n\}\subset\l, F_n\geq 0,F\in\mathrm{b}\g$ and $F_n\uparrow F$, then exists $M>0$ such that $|F|\leq M$. Observe
   $|\p(X|\g)F|\leq M\p(|X||\g)$ and
   \begin{equation*}
    \p[M\p(|X||\g)]=M\p(|X|)<\infty.
   \end{equation*}
   Therefore by {\sffamily Dominance convergence theorem}, we have
   \begin{equation*}
    \p\[\p(X|\g)F\]=\lim_{n\to\infty} \p[\p(X|\g)F_n]=\lim_{n\to\infty}\p(XF_n)=\p(XF).
   \end{equation*}
   In conclusion, the {\sffamily Monotone class theorem of the function form} means $\l$ contains all the $\mathrm{b}\g$ measurable functions.
  \end{proof}

  \begin{question}
   Suppose $(\Omega,\f,\p )$ is a probability space and $X,Y$ are two  random variables on $(E,\e)$ and $(G,\g)$ respectively. Then $X$ and $Y$ are independent iff for each $f\in\mathrm{b} \e,g\in\mathrm{b} \g$, we have
   \begin{equation*}
    \p[f(X)g(Y)]=\p[f(X)]\p[g(Y)].
   \end{equation*}
  \end{question}
  \begin{proof}
   We can let $f=1_A$ and $g=1_B$, where $A\in\e,B\in\g$. So the definition of independence between $X$ and $Y$ guarantees the above equality holds. On the contrary, if $X$ and $Y$ are independent, then $f(X)$ and $g(Y)$ are independent, the above equality naturally holds. Thus the theorem is completed.
  \end{proof}

  \begin{question}
   Suppose $(\Omega,\f,\p )$ is a complete measure space. If $A,B\in\f$ and $\p(A\Delta B)=0$, then $\p(X;A)=\p(X;B)$.
  \end{question}
  \begin{proof}
   We first prove: if $\p(N)=0$, then for each $U\in\f,\p(X;U\cup N)=\p(X;U)$. In fact,
   \begin{equation*}
    \int_UX^\pm d\p \leq \int_{U\cup N}X^\pm d\p\leq \int_UX^\pm d\p+\int_NX^\pm d\p=\int_UX^\pm d\p.
   \end{equation*}
   Thus $\p(X^\pm;U)=\p(X^\pm,U\cup N)$ and furthermore we have proved the assertion. Finally, observe $\p(A\Delta B)=0$, then we have
   \begin{equation*}
    \begin{split}
     \int_{A}Xd\p&=\int_{(A\Delta B)\Delta B}Xd\p=\int_{(A\Delta B)\cup B}Xd\p-\int_{(A\Delta B)\cap B}Xd\p=\int_BX d\p.
    \end{split}
   \end{equation*}
   The last equality holds because $\p\{(A\Delta B)\cap B \}=0$.
  \end{proof}

  \begin{question}
   Suppose $\{\mu_n\}$ is a measure sequence on the measure space $(\Omega ,\f )$. Define $\mu:=\sum_n\mu_n$ as follows:
   \begin{equation*}
    \mu(A):=\sum_{n=1}^\infty\mu_n(A),\quad \forall A\in\f.
   \end{equation*}
   Show $\mu$ is a well-defined measure. Besides, for each $f\in\f$, if $\mu(f)$ makes sense, then for each $n\geq 1,\mu_n(f)$ makes sense and we have
   \begin{equation*}
    \int fd\mu=\sum_{n=1}^\infty\int fd\mu_n.
   \end{equation*}
  \end{question}
  \begin{proof}
   Given $A\in\f$, for each $n\geq 1$, we know $\mu_n(A)\geq 0$. So $\mu(A)$ either converge to a finite positive number or equals to $\infty$. Thus $\mu$ is well-defined. Moreover, for each disjoint sets $\{A_i\}\subset\f$, by the definition of $\mu$, notice $\mu_n$ is a measure and the {\sffamily Fubini's theorem}, we have
   \begin{equation*}
    \mu\(\sum_{i=1}^\infty A_i\)=\sum_{n=1}^\infty\mu_n\(\sum_{i=1}^\infty A_i\)=\sum_{n=1}^\infty\sum_{i=1}^\infty\mu_n(A_i)=\sum_{i=1}^\infty\sum_{n=1}^\infty\mu_n(A_i)=\sum_{i=1}^\infty\mu(A_i).
   \end{equation*}
   Thus $\mu$ is a measure. Next we prove the second assertion. When $f=\i_A$, where $A\in\f$, we have
   \begin{equation*}
    \int\i_Ad\mu=\mu(A)=\sum_{n=1}^\infty\mu_n(A)=\sum_{n=1}^\infty\int\i_Ad\mu_n.
   \end{equation*}
   Actually, this is the definition of $\mu$. When $f=\sum_{i=1}^j\alpha_i\i_{A_i}$, where $\alpha_i\geq 0$ and $A_i\in\f$ for each $i=1,\cdots,j$, we can use the previous equality to obtain
   \begin{equation*}
    \begin{split}
     \int\sum_{i=1}^j\alpha_i\i_{A_i}d\mu&=\sum_{i=1}^j\alpha_j\int \i_{A_i}d\mu=\sum_{i=1}^j\alpha_j\sum_{n=1}^\infty\int\i_{A_i}d\mu_n\\
     &=\sum_{n=1}^\infty\sum_{i=1}^j\alpha_j\int\i_{A_i}d\mu_n=\sum_{n=1}^\infty\int\sum_{i=1}^j\alpha_j\i_{A_i}d\mu_n.
    \end{split}
   \end{equation*}
   When $f\in\mathrm{p}\f$, we can find positive simple function $\{f_i\}$ satisfying $f_i\uparrow f$. Hence use {\sffamily Monotonicity convergence theorem} twice, we will acquire
   \begin{equation*}
    \int fd\mu=\lim_{i\to\infty}\int f_id\mu=\lim_{i\to\infty}\sum_{n=1}^\infty\int f_id\mu_n=\sum_{n=1}^\infty\lim_{i\to\infty}\int f_id\mu_n=\sum_{n=1}^\infty\int fd\mu_n.
   \end{equation*}
   When $f\in\f$, since $\mu(f)$ makes sense, then either $\mu(f^+)<\infty$ or $\mu(f^-)<\infty$. Without loss of generality, we assume $\mu(f^-)<\infty$, then
   \begin{equation*}
    \mu(f^-)=\int f^-d\mu=\sum_{n=1}^\infty \int f^-d\mu_n<\infty\Longrightarrow \int f^-d\mu_n<\infty,\quad\!\! \forall n\geq 1.
   \end{equation*}
   Thus
   \begin{equation*}
    \begin{split}
     \int fd\mu&=\int f^+d\mu-\int f^-d\mu=\sum_{n=1}^\infty\int f^+d\mu_n-\sum_{n=1}^\infty\int f^-d\mu_n\\
     &=\sum_{n=1}^\infty\(\int f^+d\mu_n-\int f^-d\mu_n \)=\sum_{n=1}^\infty\int fd\mu_n.
    \end{split}
   \end{equation*}
   Till now, we have completed the proof.
  \end{proof}

  \begin{question}
   Suppose $\f_1=\sigma(\g_1)$ and $\f_2=\sigma(\g_2)$. Then $\f_1\times\f_2=\sigma(\{G_1\times G_2:G_1\in\g_1,G_2\in\g_2 \})$.
  \end{question}
  \begin{proof}
    Notice $\f_1\times\f_2=\sigma(\{F_1\times F_2:F_1\in\f_1,F_2\in\f_2 \})$ and $\g_1\subset\f_1,\g_2\subset\f_2$, then $\f_1\times\f_2\supset \sigma(\{G_1\times G_2:G_1\in\g_1,G_2\in\g_2 \})$. On the contrary, for all $F_1\in\f_1$, we have
    \begin{equation*}
     F_1\times\Omega_2\subset\sigma(\g_1)\times\Omega_2\subset \sigma(\{G_1\times G_2:G_1\in\g_1,G_2\in\g_2 \}),
    \end{equation*}
    where $\sigma(\g_1)\times\Omega_2=\sigma(\{G_1\times \Omega _2:G_1\in\g_1 \})$. Similarly, for each $F_2\in\f_2$,
    \begin{equation*}
     \Omega_1\times F_2\subset\Omega_1\times \sigma(\g_2)\subset \sigma(\{G_1\times G_2:G_1\in\g_1,G_2\in\g_2 \}).
    \end{equation*}
    Thus $F_1\times F_2=(F_1\times\Omega_2)\cap (\Omega_1\times F_2)\subset \sigma(\{G_1\times G_2:G_1\in\g_1,G_2\in\g_2 \})$. This means
    \begin{equation*}
     \{F_1\times F_2:F_1\in\f_1,F_2\in\f_2 \}\subset \sigma(\{G_1\times G_2:G_1\in\g_1,G_2\in\g_2 \})
    \end{equation*}
    and furthermore
    \begin{equation*}
     \begin{split}
      \f_1\times\f_2&=\sigma(\{F_1\times F_2:F_1\in\f_1,F_2\in\f_2 \})\\
      &\subset \sigma(\{G_1\times G_2:G_1\in\g_1,G_2\in\g_2 \}).
     \end{split}
    \end{equation*}
    We finish the proof.
  \end{proof}

  \begin{question}
   Suppose $\xi:(\Omega,\f )\to(E,\mathscr{E} )$ is a random variable. Then for each $U\in\mathscr{E}$, the restriction mapping $\xi:(\xi^{-1}(U),\f\cap\xi^{-1}(U)  )\to(U ,\mathscr{E}\cap U )$ is also a random variable.
  \end{question}
  \begin{proof}
   For each $A\in\mathscr{E}$, we have
   \begin{equation*}
    \xi^{-1}(A\cap U)=\xi^{-1}(A)\cap \xi^{-1}(U)\in\f\cap\xi^{-1}(U).
   \end{equation*}
   So the restriction mapping $\xi:(\xi^{-1}(U),\f\cap\xi^{-1}(U)  )\to(U ,\mathscr{E}\cap U )$ is also a random variable.
  \end{proof}

  \begin{question}
   Suppose $(\Omega,\f,\p )$ is a probability space and $\{X_n\}$ is a random variable sequence. Then $X_n\Rightarrow X$ does not implies $\p(X_n)\to \p(X)$.
  \end{question}
  \begin{proof}
   Consider the following random variable sequence:
   \begin{equation*}
    X_n=
    \begin{cases}
     n^2&,\mbox{with probability }1/n\\
     0&,\mbox{with probability }1-1/n.\\
    \end{cases}
   \end{equation*}
   Since for each $\varepsilon\in(0,1)$, we have
   \begin{equation*}
    \p(|X_n|>\varepsilon )=1/n\to 0,\quad n\to\infty.
   \end{equation*}
   Thus $X_n\overset{\p}{\to }0$. Thus we can find a subsequence $\{X_{n_k}\}$ such that $X_{n_k}\overset{\mathrm{a.s.}}{\to}0$. However, we know for each $k\geq 1$,
   \begin{equation*}
    \p(X_{n_k})=\int X_{n_k}d\p=\frac{n^2_k}{n_k}=n_k.
   \end{equation*}
   Thus $\p(X_{n_k})\to\infty \neq \p(0)=0$.
  \end{proof}

  \begin{question}
   Suppose $\{X_n\}$ and $\{Y_n\}$ are two sequences of random variable. If $X_n$ and $Y_n$ are independent for each $n\geq 1$ and $X_n\overset{\mathrm{a.s.}}{ \to} X,Y_n\overset{\mathrm{a.s.}}{ \to} Y$. Then $X$ and $Y$ are independent.
  \end{question}
  \begin{proof}
   We can utilize the characteristic function and {\sffamily Dominance convergence theorem} to obtain
   \begin{equation*}
    \begin{split}
     \p\[e^{i(t_1X+t_2Y)}\]&=\lim_{n\to\infty}\p\[e^{i(t_1X_n+t_2Y_n)}\]=\lim_{n\to\infty}\p\(e^{it_1X_n}\)\p\(e^{it_2Y_n}\)\\
     &=\p\(e^{it_1X}\)\p\(e^{it_2Y}\).
    \end{split}
   \end{equation*}
   Thus $X$ and $Y$ are independent.
  \end{proof}

  \begin{question}
   Suppose $\mu$ is a measure on $(\r,\b(\r))$ and $dx$ is the Lebesgue measure. Define $I_n:=n\mathbbm{1}_{[0,1/n]}$ and
   \begin{equation*}
    I_n*\mu(x)=\int_\r I_n(t)\mu(x-dt).
   \end{equation*}
   Then discuss: under what conditions the weak convergence $(I*\mu)dx\Rightarrow \mu(dx)$ holds?
  \end{question}
  \begin{proof}
   Since
   \begin{equation*}
    \begin{split}
     I_n*\mu(x)&=\int_\r I_n(t)\mu(x-dt)=\int_\r n\mathbbm{1}_{[0,1/n ]}(t)\mu(x-dt)\\
     &=n\mu(x-[0,1/n])=n\mu[x-1/n,x ].
    \end{split}
   \end{equation*}
   First we claim $I_n*\mu\in\b(\r)$. In fact,
   \begin{equation*}
    I_n*\mu(x)=n\mu[x-1/n,x ]=n\[F_\mu(x)-F_\mu((x-1/n)-) \],
   \end{equation*}
   where $F_\mu$ is the distribution function of $\mu$, namely $F_\mu(x)=\mu(x_0,x]$ for a $x_0\in\r$ such that $F_\mu(x)<\infty $ for each $x\in\r$. Since $F_\mu$ is a non-decreasing function, so $F_\mu\in \b(\r)$. Moreover, for each $x\in\r$,
   \begin{equation*}
    F_\mu(x-)=\lim_{y\uparrow x}F_\mu(y)=\lim_{\q\ni q\uparrow x}F_\mu(q) \in\b(\r).
   \end{equation*}
   Thus $F_\mu((x-1/n)-)\in\b(\r)$ and therefore $I_n*\mu\in\b(\r)$.
   \sffamily Suppose $\mu$ satisfying for each $f\in\mathrm{b}C(\r)$,
   \begin{equation*}
     \int_{\r}\sup_{x\in[t,t+1]}|f(x)|\mu(dt)<\infty.
   \end{equation*}
   Then
   we have
   \begin{equation*}
    \begin{split}
     \int_\r f(x)I_n&*\mu(x)dx=\int_\r nf(x)\int_{[x-1/n,x ]}\mu(dt)dx\\
     &=\int_\r\(\int_{[t,t+1/n]}nf(x)dx\)\mu(dt)\\
     &=\int_\r\(\dashint_{[t,t+1/n]}f(x)dx\)\mu(dt)\\
     &\to \int_\r f(t)\mu(dt),\quad n\to\infty.
    \end{split}
   \end{equation*}
   \begin{enumerate}
    \item The second equality holds since
   \begin{equation*}
    \begin{split}
     \int_\r n|f(x)|&\int_{[x-1/n,x ]}\mu(dt)dx=\int_\r\(\int_{[t,t+1/n]}n|f(x)|dx\)\mu(dt)\\
     &\leq \int_\r \sup_{x\in[t,t+1]}|f(x)|\mu(dt)<\infty.
    \end{split}
   \end{equation*}
   Then we can use {\sffamily Fubini's theorem}.
    \item The last limit holds since
    \begin{equation*}
     \dashint_{[t,t+1/n]}f(x)dx\leq \sup_{x\in[t,t+1]}|f(x)|
    \end{equation*}
    and $\sup_{x\in[t,t+1]}|f(x)|$ is integrable with respect to $\mu$. So we can utilize the {\sffamily Dominance convergence theorem} and the {\sffamily integral mean value theorem}.
   \end{enumerate}
   In conclusion, if we suppose $\mu$ satisfying for each $f\in \mathrm{b}C(\r)$,
   \begin{equation*}
    \int_{\r}\sup_{x\in[t,t+1]}|f(x)|\mu(dt)<\infty,
   \end{equation*}
   then the weak convergence $(I*\mu)dx\Rightarrow \mu(dx)$ holds. Actually, we claim the above requirement is equivalent to $\mu(\r)<\infty$. In fact, we can choose $1\in \text{b}C(\r)$, then
   \begin{equation*}
    \int_{\r}\sup_{x\in[t,t+1]}1\mu(dt)=\mu(\r)<\infty.
   \end{equation*}
   The inverse direction is obvious.
  \end{proof}

  \begin{question}
   Suppose $(\Omega,\f,\p )$ and $(E,\e)$ are two probability spaces and $X,Y$ are two random variables with respect to $\f$ and $\e$. Assume $\g\subset\f$ is a sub $\sigma$-algebra and $X\in\g$, then for each $f\in\mathrm{b}\e^2$, we have
   \begin{equation*}
    \p[f(X,Y)|\g]=\p[f(x,Y)|\g]|_{x=X}.
   \end{equation*}
   Moreover, if $Y$ is independent to $\g$, then we have
   \begin{equation*}
    \p[f(X,Y)|\g]=\p[f(x,Y)]|_{x=X}.
   \end{equation*}
   Especially, if we let $F(x)=\p(x\leq Y)$, then we have
   \begin{equation*}
    \p[F(X)]=\p(X\leq Y).
   \end{equation*}
  \end{question}
  \begin{proof}
   When $f(x,y)=\mathbbm{1}_{A\times B}(x,y)$, where $A,B\in\e$, we have
   \begin{equation*}
    \begin{split}
     \p[f(X,Y)|\g]&=\p[\mathbbm{1}_{\{X\in A\}}\mathbbm{1}_{\{Y\in B\}}|\g]=\mathbbm{1}_{\{X\in A\}}\p[\mathbbm{1}_{\{Y\in B\}}|\g]\\
     &=\mathbbm{1}_{A}(x)\p[\mathbbm{1}_{\{Y\in B\}}|\g]|_{x=X}=\p[\mathbbm{1}_{A}(x)\mathbbm{1}_{\{Y\in B\}}|\g]|_{x=X}\\
     &=\p[f(x,Y)|\g]|_{x=X}.
    \end{split}
   \end{equation*}
   Consider the following set
   \begin{equation*}
    \l:=\{f(x,y)\in\mathrm{b}\e^2:\p[f(X,Y)|\g]=\p[f(x,Y)|\g]|_{x=X} \}
   \end{equation*}
   and write $\a:=\{\mathbbm{1}_{A\times B}(x,y):A,B\in\e\}$. Then we know $\a\subset\l$. Obviously, $1\in \l$, for each $f,g\in\l$ and $\alpha,\beta\in\r $, we have
   \begin{equation*}
    \begin{split}
     \p[\alpha f(X,Y)+\beta &g(X,Y)|\g]=\alpha \p[ f(X,Y)|\g]+\beta \p[g(X,Y)|\g]\\
     &=\alpha \p[f(x,Y)|\g]|_{x=X}+\beta \p[g(x,Y)|\g]|_{x=X}\\
     &=\p[\alpha f(x,Y)+\beta g(x,Y)|\g]|_{x=X}.
    \end{split}
   \end{equation*}
   So $\l$ is a linear space. Moreover, for $\{f_n\}\subset\l,f\in\mathrm{b}\e$ and $0\leq f_n\uparrow f$, by {\sffamily Monotone convergence theorem}, we have
   \begin{equation*}
    \begin{split}
     \p[f(X,Y)|\g]&=\lim_{n\to\infty}\p[f_n(X,Y)|\g]=\lim_{n\to\infty}\p[f_n(x,Y)|\g]|_{x=X}\\
     &=\p[f(x,Y)|\g]|_{x=X}.
    \end{split}
   \end{equation*}
   Define $\mathscr{C}:=\{A\times B:A,B\in\e\}$.  The {\sffamily Monotone class theorem of function form} gives $\l$ contains all the $\sigma(\mathscr{C})=\e^2$ measurable mappings. The second assertion is trivial.

   Observe $X$ and $Y$ are independent, then
   \begin{equation*}
    F(X)=\p(x\leq Y)|_{x=X}=\p(\mathbbm{1}_{\{X\leq Y\}}|X).
   \end{equation*}
   Therefore,
   \begin{equation*}
    \p[F(X)]=\p[\p(\mathbbm{1}_{\{X\leq Y\}}|X)]=\p(X\leq Y).
   \end{equation*}
  \end{proof}


  \begin{question}
   Suppose $X$ is a random variable on the probability space $(\Omega,\f,\p )$  and $\a,\g$ are two  sub $\sigma$-algebras of $\f$. Assume $\g$ and $\sigma(\a\cup\sigma(X) )$ are independent, then we have
   \begin{equation*}
    \p[X|\sigma(\a\cup\g ) ]=\p(X|\a ).
   \end{equation*}
  \end{question}
  \begin{proof}
   For each $A\in\a,G\in\g$, we have
   \begin{equation*}
    \p(X\mathbbm{1}_{A\cap G})=\p(X1_A)\p(G)=\p[\p(X|\a)1_A]\p(G)=\p[\p(X|\a)\mathbbm{1}_{A\cap G}].
   \end{equation*}
   Define the $\pi$-system $\mathscr{C}:=\{A\cap G:A\in\a,G\in\g\}$, then we have $\sigma(\mathscr{C})=\sigma( \a\cup \g)$. So the {\sffamily Monotone class theorem} yields that the above formula holds for each $B\in \sigma( \a\cup \g)$.
   Obviously $\p(X|\a )\in \sigma( \a\cup \g)$, so $\p[X|\sigma(\a\cup\g ) ]=\p(X|\a )$.
  \end{proof}

  \begin{question}
   Suppose $(\Omega,\f,\p )$ is a probability space and $\a,\g$ are two independent sub $\sigma$-algebras of $\f$. Given two random variables $X\in\a,Y\in\g$ and two $\sigma$-algebra $\a'\subset\a,\g'\subset\g$, then we have
   \begin{equation*}
    \p[XY|\sigma( \a'\cup \g')]=\p(X|\a')\p(Y|\g').
   \end{equation*}
  \end{question}
  \begin{proof}
   For each $A\in\a',G\in\g'$, observe $X1_A\in \a,Y1_G\in\g$ and the independence between $\a,\g$, we have
   \begin{equation*}
    \begin{split}
     \p(XY1&_{A \cap G})=\p(X1_A Y1_G)=\p(X1_A)\p(Y1_G)\\
     =&\p[\p(X|\a')1_A]\p[\p(Y|\g')1_G]=\p[\p(X|\a')\p(Y|\g')\mathbbm{1}_{A\cap G}].
    \end{split}
   \end{equation*}
   Define the $\pi$-system $\mathscr{C}:=\{A\cap G:A\in\a',G\in\g'\}$, then we have $\sigma(\mathscr{C})=\sigma( \a'\cup \g')$. So the {\sffamily Monotone class theorem} yields that the above formula holds for each $B\in \sigma( \a'\cup \g')$.
   Obviously $\p(X|\a')\p(Y|\g')\in \sigma( \a'\cup \g')$, so $\p[XY|\sigma( \a'\cup \g')]=\p(X|\a')\p(Y|\g')$.
  \end{proof}

  \begin{question}
   Suppose $X,Y$ are two random variables on the probability space $(\Omega ,\f,\p )$ and $\g$ is a sub $\sigma$-sigma algebra of $\f$. Write $\mathscr{N}$ for all the $\p$-zero set and define the enlargement $\sigma $-algebra $\g^*=\sigma(\g\cup\mathscr{N} )$. Assume $X=Y,\p$-a.s. and $\mathbb{E}X$ exists, then we have
   \begin{equation*}
    \p(X|\g)=\p(Y|\g^*),\quad \p\text{\rmfamily -a.s.}
   \end{equation*}
  \end{question}
  \begin{proof}
   First we have $\p(X|\g)\in\g\subset\g^*$ and $\p(Y|\g^*)\in\g^*$. Moreover, for each $G^*\in\g^*$, exist $G\in\g$ and $N\in\mathscr{N}$ such that $G^*=G\Delta N$. So
   \begin{equation*}
    \p[\p(X|\g)\mathbbm{1}_{G^*}]=\p[\p(X|\g)\mathbbm{1}_{G}]=\p(X1_G)=\p(Y1_G)=\p(Y\mathbbm{1}_{G^*}).
   \end{equation*}
   This yields $\p(X|\g)=\p(Y|\g^*),\p$-a.s..
  \end{proof}

  \begin{question}
   Prove $\b(\r_+)=\sigma( \cup_{t> 0}\b[0,t))$.
  \end{question}
  \begin{proof}
   Obviously $\sigma( \cup_{t> 0}\b[0,t))\subset\b(\r_+)$. On the contrary, for each $0\leq a<b<\infty$, then exists $t\geq 0$ such that $(a,b]\subset [0,t)$. So $(a,b]\in \b[0,t)\subset \sigma( \cup_{t> 0}\b[0,t))$. Therefore
     \begin{equation*}
      \b(\r_+)=\sigma(\{(a,b]:0\leq a<b<\infty\})\subset \sigma( \cup_{t> 0}\b[0,t)).
     \end{equation*}
     Thus $\b(\r_+)=\sigma( \cup_{t> 0}\b[0,t))$ holds.
  \end{proof}

  \begin{question}
   Give an example to show $X$ and $Y$ are uncorrelated but not independent.
  \end{question}
  \begin{proof}
   Consider $\xi$ is a random variable satisfying $\p(\xi=1)=1/2$ and $\p(\xi=-1)=1/2$. Moreover, let $X$ is a random variable independent to $\xi$. Define $Y=\xi X$, the we show $X$ and $Y$ are uncorrelated but not independent. In fact,
   \begin{equation*}
    \p(XY)=\p(\xi X^2)=\p(\xi)\p(X^2)=0\Longleftrightarrow \mathrm{Cov}(XY)=0.
   \end{equation*}
   However, obviously $X$ and $Y=\xi X$ are not independent in general.
  \end{proof}

  \begin{question}
   Consider the measurable space $(E,\b(E) )$, where $E$ is a metric space and $\b(E)$ is the Borel $\sigma$ algebra of $E$. When $A\in\b(E)$ is open or closed, give some examples to show $1_A$ can be a pointwise limit of Lipschitz continuous function sequence.
  \end{question}
  \begin{proof}
   \begin{enumerate}
    \item When $A$ is an open set.
    \begin{enumerate}
     \item Consider the following continuous function sequence:
    \begin{equation*}
     F_n(x)=[nd(x,A^c)]\wedge 1,\quad x\in E,n\geq 1.
    \end{equation*}
    Since $A^c$ is a closed subset, then $d(x,A^c)>0$ iff $x\in A$. Then it is easy to find $F_n\uparrow 1_A$. Besides, $F_n$ is Lipschitz continuous. We just need to notice {\bfseries QUESTION 2.1}, we have
    \begin{equation*}
     \begin{split}
      |F_n(x)-F_n(y)|&=|[nd(x,A^c)]\wedge 1-[nd(y,A^c)]\wedge 1|\\
      &\leq |nd(x,A^c)-nd(y,A^c) |=nd(x,y).
     \end{split}
    \end{equation*}
     \item Consider the following continuous function sequence:
     \begin{equation*}
      F_n(x)=[nd(x,A^c)-1]^+\wedge 1,\quad x\in E,n\geq 1.
     \end{equation*}
     Since $A^c$ is a closed subset, then $d(x,A^c)>0$ iff $x\in A$. Then it is easy to find $F_n\uparrow 1_A$. Besides, $F_n$ is Lipschitz continuous. We just need to notice {\bfseries QUESTION 2.1}, we have
    \begin{equation*}
     \begin{split}
      |F_n(x)-F_n(y)|&=|[nd(x,A^c)-1]^+\wedge 1-[nd(x,A^c)-1]^+\wedge 1|\\
      &\leq |[nd(x,A^c)-1]^+-[nd(t,A^c)-1]^+ |\\
      &\leq |nd(x,A^c)-nd(y,A^c)| =nd(x,y).
     \end{split}
    \end{equation*}
    \end{enumerate}
    \item When $A$ is a closed set.
    \begin{enumerate}
     \item Consider the following continuous function sequence:
    \begin{equation*}
     F_n(x)=[1-nd(x,A)]^+,\quad x\in E,n\geq 1.
    \end{equation*}
    Since $A$ is a closed subset, then $d(x,A)=0$ iff $x\in A$. Then it is easy to find $F_n\uparrow 1_A$. Besides, $F_n$ is Lipschitz continuous. We just need to notice {\bfseries QUESTION 2.1}, we have
    \begin{equation*}
     \begin{split}
      |F_n(x)-F_n(y)|&=|[1-nd(x,A)]^+-[1-nd(y,A)]^+|\\
      &\leq |nd(x,A)-nd(y,A) |=nd(x,y).
     \end{split}
    \end{equation*}
     \item Consider the following continuous function sequence:
    \begin{equation*}
     F_n(x)=\frac{1}{1+nd(x,A)} ,\quad x\in E,n\geq 1.
    \end{equation*}
    Since $A$ is a closed subset, then $d(x,A)=0$ iff $x\in A$. Then it is easy to find $F_n\uparrow 1_A$. Besides, $F_n$ is Lipschitz continuous. In fact, we have
    \begin{equation*}
     \begin{split}
      |F_n(x)-F_n(y)|&=\left|\frac{1}{1+nd(x,A)}-\frac{1}{1+nd(y,A)}\right|\\
      &=\left|\frac{n[d(y,A)-d(x,A)]}{[1+nd(x,A)][1+nd(y,A)]} \right|\\
      & \leq |nd(x,A)-nd(y,A) |=nd(x,y).
     \end{split}
    \end{equation*}
    \end{enumerate}
   \end{enumerate}
  \end{proof}

  \begin{question}
   Prove a $\sigma$-algebra $\f$ is either finite or uncountable.
  \end{question}
  \begin{proof}
   We prove this by contradiction. Assume $|\f|=\aleph_0$. Define
  	\begin{equation*}
  		\g:=\{A\in\f:\text{there is no }B_1, B_2\subsetneq A\text{ such that }A=B_1\cup B_2\backslash\{\emptyset\} .
  	\end{equation*}
  	Obviously $|\g|\leq \aleph_0$, say $\g=\{A_k\}_{k\geq 1}$. We claim $A_i\cap A_j=\emptyset$ for any $i\neq j$. We prove this by contradiction. If this is not the case, we have $A_i\cap A_j\neq \emptyset$.
  	\begin{enumerate}
  		\item Without loss of generality, if $A_i\cap A_j=A_i$, then 
  		\begin{equation*}
  			A_j=A_i\cup(A_j\backslash A_i).
  		\end{equation*}
  		Noting $A_i\neq A_j$ and $A_i,A_j\backslash A_j\in\f$, this contradict the fact $A_j\in\g$.
  		\item If $A_i\cap A_j\neq A_i$ and $A_i\cap A_j\neq A_j$, we have
  		\begin{equation*}
  			A_i=(A_i\cap A_j)\cup (A_i\backslash A_j).
  		\end{equation*}
  		Noting $A_i\cap A_j\in\f$ and $A_i\backslash A_j\in\f$, this contradict the fact $A_i\in\g$.
  	\end{enumerate}
  	In conclusion, we have shown the elements in $\g$ are mutually disjoint. 
  	
  	Now, we consider several situations:
  	\begin{enumerate}
  		\item If $\sigma(\g)=\f$, noting $|\f|=\aleph_0$, we must have $|\g|=\aleph_0$. since the elements in $\g$ are mutually disjoint, the well-known result shows 
  		\begin{equation*}
  			\f=\sigma(\g)=\left\{\sum_{k\in I}A_k:I\subset \z_+ \right\}.
  		\end{equation*}
  		So we can construct a bijection:
  		\begin{equation*}
  			\varphi : 2^{\z_+}\to \f,I\mapsto \sum_{k\in I}A_k,
  		\end{equation*}
  		where $2^{\z_+}$ is the set of all the subset of $\z_+$. This implies $|\f|=|\sigma(\g)|=\aleph$, which contradicts $|\f|=\aleph_0$.
  		\item If $\sigma(\g)\subsetneq \f$, there must exists $A\in\f$ such that $A\notin \sigma(\g)$. Since $A\notin\g$, there exist $A^{(1)}_i\in \f$ and $A^{(1)}_i\subsetneq A$ $(i=1,2)$ such that $A=A_1^{(1)}\cup A_2^{(1)}$. Noting $A\notin \sigma(\g)$, so there exists $n_1=1$ or $2$ such that $A^{(1)}_{n_1}\notin \sigma(\g)$ and thus $A^{(1)}_{n_1}\notin \g$. So there exist $A^{(2)}_i\in \f$ and $A^{(2)}_i\subsetneq A^{(1)}_{n_1}$ $(i=1,2)$ such that $A^{(1)}_{n_1}=A_1^{(2)}\cup A_2^{(2)}$. Again, $A^{(1)}_{n_1}\notin \sigma(\g)$ yields that there exists $n_2=1$ or $2$ such that $A^{(2)}_{n_2}\notin\sigma(\g)$. Using induction, we can obtain a sequence $\{A^{(k)}_{n_k}\}_{k\geq 1}\subset\f$ satisfying
  		\begin{equation*}
  			A^{(0)}_{n_0}:=A\supsetneq A^{(1)}_{n_1}\supsetneq A^{(2)}_{n_2}\supsetneq \cdots \supsetneq A^{(k)}_{n_k}\cdots.
  		\end{equation*}
  		Define for any $k\geq 1$, 
   \begin{equation*}
    C_k:=A_{n_{k-1}}^{(k-1)}-A_{n_{k}}^{(k)}.
   \end{equation*}
   Then $\{C_k\}_{k\geq 1}$ is a disjoint set sequence. Since $\{C_k\}_{k\geq 1}\subset \f$, then 
   \begin{equation*}
   	\sigma\(\{C_k\}_{k\geq 1}\)=\left\{\sum_{k\in I}C_k:I\subset \z_+ \right\}\subset \f
   \end{equation*}
   A similar argument will yield contradiction again. 
  	\end{enumerate}
  	In conclusion, $|\f|\neq \aleph_0$ for any $\sigma$-algebra $\f$.

   
   
   \textcolor{red}{A fake proof}. If $\g$ is countable, denoted as $\g=\{A_n\}_{n\geq 1}$. Since $\cup_{n}A_n\in\g$, then exists $n_0\geq 1$ such that $\cup_{n}A_n=A_{n_0}$. Note $A_{n_0}^c\in\g$, then exists $n_1\geq 1$ such that $A_{n_0}^c=A_{n_1}$. So
   	\begin{equation*}
   		A_{n_0}^c=A_{n_1}\subset \bigcup_{n\geq 1}A_n=A_{n_0}.
   	\end{equation*}
   	This is a contradiction. \textcolor{red}{ This proof is not grounded because $\cup_n A_n=A_{n_0}=\Omega$ and $A_{n_1}=A_{n_0}^c=\emptyset$. So $A_{n_0}^c=\emptyset\subset A_{n_0}=\Omega$ naturally true and there is no contradiction.}
  \end{proof}

  \begin{question}
   Suppose $D$ is a dense subset of $\r$. If the distribution function $F_n(x)\to F(x)$ for each $x\in D$ and $F$ is continuous, then
   \begin{equation*}
    \sup_{x\in\r} |F_n(x)-F(x)|\to 0,\qquad n\to\infty.
   \end{equation*}
   Thus $F_n(x)\overset{\mathrm{w}}{\to}F(x)$ naturally.
  \end{question}
  \begin{proof}
   At first, we show if $\exists M\in D\cap[0,\infty),\st F_n(x)\to F(x),\forall x\in D\cap[-M,M]$, then $F_n$ converges to $F$ uniformly on $D\cap[-M,M]$. In fact, since $F$ is uniformly continuous on $[-M,M]$, then $\forall\varepsilon>0,\exists\delta>0,\forall x,y\in[-M,M]$ and $|x-y|<\delta$, then $|F(x)-F(y)|<\varepsilon$. Hence we can separate $[-M,M]$ into $n$ subintervals: $-M=x_0<x_1<\cdots<x_n=M$ satisfying $\{x_j\}_{j=0}^n\subset D\cap[-M,M]$ and $\max_{0\leq j\leq n-1}|x_{j+1}-x_j|<\delta$. Moreover, since $F_n(x_j)\to F(x_j),\forall j\in\{0,\cdots,n\}$, then $\exists N_1\in\n,\forall n\geq N_1$, we have $|F_n(x_j)-F(x_j)|<\varepsilon,\forall j\in\{0,\cdots,n \}$. When $n\geq N_1$, for each $x\in [-M,M]$, observe $\exists j\in\{0,\cdots,n \},\st x\in[x_n,x_{n+1}]$, then
   \begin{equation*}
    \begin{split}
     |F_n(x&)-F(x)|\\
     \leq& |F_n(x)-F_n(x_n)|+|F_n(x_n)-F(x_n)|+|F(x_n)-F(x)|\\
     < &|F_n(x_{n+1})-F_n(x_n)|+\varepsilon+\varepsilon\\
     \leq &|F_n(x_{n+1})-F(x_{n+1})|+|F(x_{n+1})-F(x_n) |+|F(x_n)-F_n(x_n)|+2\varepsilon\\
     < &\varepsilon+\varepsilon+\varepsilon+2\varepsilon=5\varepsilon.
    \end{split}
   \end{equation*}
   This implies $F_n$ converges to $F$ uniformly on $[-M,M]$.

   Next, notice $\lim_{x\to-\infty}F(x)=0,\lim_{x\to\infty}F(x)=1$. Then $\exists M\in D\cap [0,\infty) ,\st F(-M)<\varepsilon$ and $1-F(M)<\varepsilon$. Note $F_n(\pm M)\to F(\pm M)$, then $\exists N_2\in\n,\forall n\geq N_2$, then $|F_n(\pm M)-F(\pm M)|<\varepsilon$. So when $n\geq N_2$, we have
   \begin{equation*}
    F_n(-M)\leq |F_n(-M)-F(-M)|+F(-M)<\varepsilon+\varepsilon=2\varepsilon
   \end{equation*}
   and
   \begin{equation*}
    1-F_n(M)\leq 1-F(M)+|F(M)-F_n(M)|<\varepsilon+\varepsilon=2\varepsilon.
   \end{equation*}
   Whence $\forall x\leq -M$, we have
   \begin{equation*}
    |F_n(x)-F(x)|\leq F_n(x)+F(x)\leq F_n(-M)+F(-M)<2\varepsilon+\varepsilon=3\varepsilon.
   \end{equation*}
   And $\forall x\geq M$, we have
   \begin{equation*}
    \begin{split}
     |F_n(x)&-F(x)|\leq |F_n(x)-1|+|1-F(x)|\\
     \leq& \[1-F_n(M)\]+\[1-F(M)\]<2\varepsilon+\varepsilon=3\varepsilon.
    \end{split}
   \end{equation*}
   This implies $F_n$ converges to $F$ uniformly on $[-M,M]^c$.

   In conclusion $F_n$ converges to $F$ uniformly on $\r$, namely, $\sup_{x\in\r}|F_n(x)-F(x)|\to 0$. As for the last assertion, since $F(x)$ is continuous on $\r$, then for each $x\in\r$, we have
   \begin{equation*}
    |F_n(x)-F(x)|\leq \sup_{x\in\r}|F_n(x)-F(x)|\to 0,
   \end{equation*}
   which implies $F_n(x)\overset{\mathrm{w}}{\to}F(x)$.
  \end{proof}

  \begin{question}
   Suppose $X,Y$ are two $(E,\e)$ random variables on the probability space $(\Omega,\f,\p )$. Assume $\g\subset \f$ is a sub $\sigma$ algebra and for each $f,g\in\mathrm{b}\e$,
   \begin{equation*}
    \p[f(X)g(Y)|\g]=\p[f(X)]\p[g(Y)].
   \end{equation*}
   Then $X$ and $Y$ are independent. Especially,
   assume $\{C_k\}_{k\geq 1}$ is a partition of $\Omega$ and $\p(C_k)>0,\forall k\geq 1$. If for each $k\geq 1$ and $A,B\in\e$,
   \begin{equation*}
    \p(X\in A,Y\in B|C_k)=\p(X\in A)\p(Y\in B),
   \end{equation*}
   then $X$ and $Y$ are independent.
  \end{question}
  \begin{proof}
   Take expectation on both sides of $\p[f(X)g(Y)|\g]=\p[f(X)]\p[g(Y)]$ to obtain
   \begin{equation*}
    \p[f(X)g(Y)]=\p\{\p[f(X)g(Y)|\g]\}=\p[f(X)]\p[g(Y)].
   \end{equation*}
   This implies $X$ and $Y$ are independent. Especially, we let $f=1_A,g=1_B$ and $\g=\sigma(C_k:k\geq 1)$. Then $\p(X\in A,Y\in B|C_k)=\p(X\in A)\p(Y\in B)$ is equivalent to $\p[f(X)g(Y)|\g]=\p[f(X)]\p[g(Y)]$. So the result follows.
  \end{proof}

  \begin{question}
   Suppose $X$ is an exponential distribution random variable with parameter $\mu$ and $Y$ is an independent variable on the probability space $(\Omega,\f,\p )$. Then for each $t>0$, we have
   \begin{equation*}
    \p(X>t+Y|X>Y,Y\geq 0)=\p(X>t).
   \end{equation*}
  \end{question}
  \begin{proof}
   When $Y$ is a nonnegative constant, this is the well-known memoryless property. When $Y$ is a random variable, write its distribution as $\nu(dx)$. Observe for each $t\geq 0$, the independence between $X$ and $Y$ impies
   \begin{equation*}
    \begin{split}
     \p(X>t+Y,Y\geq 0)&=\p(\mathbbm{1}_{\{X>t+Y\}}\mathbbm{1}_{\{Y\geq 0\}})\\
     &=\p[\p(\mathbbm{1}_{\{X>t+Y\}}\mathbbm{1}_{\{Y\geq 0\}}|Y)]\\
     &=\p[\mathbbm{1}_{\{Y\geq 0\}}\p(\mathbbm{1}_{\{X>t+Y\}}|Y)]\\
     &=\p[\mathbbm{1}_{\{Y\geq 0\}}\p(\mathbbm{1}_{\{X>t+y\}}|Y)|_{y=Y}]\\
     &=\p[\mathbbm{1}_{\{Y\geq 0\}}\p(\mathbbm{1}_{\{X>t+y\}})|_{y=Y}]\\
     &=\p[\mathbbm{1}_{\{Y\geq 0\}}\p(X>t+y)|_{y=Y}]\\
     &=\int_{-\infty}^\infty \mathbbm{1}_{\{y\geq 0\}}\p(X>t+y)\mu(dy)\\
     &=\int_0^\infty\p(X>t+y)\mu(dy).
    \end{split}
   \end{equation*}
   Therefore, we have
   \begin{equation*}
    \begin{split}
     \p(X>&t+Y|X>Y,Y\geq 0)=\frac{\p(X>t+Y,Y\geq 0)}{\p(X>Y,Y\geq 0)}\\
     &=\int_0^\infty  \p(X>t+y)\nu( dy)\bigg/\int_0^\infty \p(X>y)\nu( dy) \\
     &=\int_0^\infty   e^{-\mu(t+y)} \nu( dy)\bigg/\int_0^\infty  e^{-\mu y} \nu( dy)\\
     &=e^{-\mu t}\int_0^\infty  e^{-\mu y} \nu( dy)\bigg/\int_0^\infty  e^{-\mu y} \nu( dy) \\
     &=e^{-\mu t}=\p(X>t).
    \end{split}
   \end{equation*}
   This shows a {\sffamily strong} version of {\sffamily memoryless property}.
  \end{proof}

  \begin{question}
   Suppose $X$ and $Y$ are two $(E,\e)$ random variables on the probability space $(\Omega,\f,\p )$. If exists a distribution $\nu(dx)$ on $(E,\e)$ such that for each $A\in\e$, we have $\p(X\in A|Y)=\nu(A)$, then $X$ and $Y$ are independent. Moreover, the distribution of $X$ is $\nu(dx)$.
  \end{question}
  \begin{proof}
   Take the expectation on both sides of $\p(X\in A|Y)=\nu(A)$, we have $\p(X\in A)=\nu(A)$. Therefore, the distribution of $X$ is $\nu(dx)$. Since for each $B\in\e$, we have
   \begin{equation*}
    \begin{split}
     \p(X\in A, Y\in B)&=\p[\p(X\in A, Y\in B|Y)]\\
     &=\p[\mathbbm{1}_{\{Y\in B\}}\p(X\in A|Y)]\\
     &=\p[\mathbbm{1}_{\{Y\in B\}}\nu(A)]\\
     &=\nu(A)\p(Y\in B).
    \end{split}
   \end{equation*}
   The desired result follows.
  \end{proof}

  \begin{question}
   Given a sequence of integrable i.i.d. random variables $\{X_n\}_{n\geq 1}:(\Omega,\f,\p )\to (\r,\b(\r))$. Then for each sequence of random times $\{T_m\}_{m\geq 1}:\Omega \to \z_+$, where $T_1<T_2<\cdots <T_m<\cdots $, we have
   \begin{equation*}
    \p\(\sum_{m=1}^\infty X_{T_m}=\infty\)=1.
   \end{equation*}
  \end{question}
  \begin{proof}
   For any subsequence $\{n_m\}_{m\geq 1}$ of $\z_+$, we know $\{X_{n_m}\}_{m\geq 1}$ is also an integrable i.i.d. random variable sequence. Thus, by the {\sffamily strong law of large numbers}, we have
   \begin{equation*}
    \p\(\sum_{m=1}^\infty X_{n_m}=\infty\)=1.
   \end{equation*}
   Write $N$ for the exceptional set with $\p(N)=0$. Then
   \begin{equation*}
    \sum_{m=1}^\infty X_{n_m}(\omega )=\infty,\quad\forall \omega\in N^c.
   \end{equation*}
    Given any $\omega\in N^c$, notice $\{T_m(\omega )\}_{m\geq 1}$ is a subsequence of $\z_+$, then
   \begin{equation*}
    \sum_{m=1}^\infty X_{T_m(\omega )}(\omega )=\infty,
   \end{equation*}
   which implies
   \begin{equation*}
    \p\(\sum_{m=1}^\infty X_{T_m}=\infty\)=1.
   \end{equation*}
  \end{proof}

  \begin{question}
   Suppose $F$ is a distribution function of the random variable $X$, i.e., $F(x)=\p(-\infty,x] $ and define:
   \begin{equation*}
    F^{-1}(u)=\inf\{x\in\r:F(x)\geq u \}.
   \end{equation*}
   Then $F(x)\geq u$ iff $x\geq F^{-1}(u)$.
  \end{question}
  \begin{proof}
   By the definition of infimum, we know $F(x)\geq u$ implies $x\geq F^{-1}(u)$. On the contrary, by the right continuity of $F$ and the definition of infimum, exists $\{x_n\}\subset\r$ such that $F(x_n)\geq u$ and $x_n\downarrow F^{-1}(u)$. Moreover, observe $F$ is increasing, then $F(x)\geq F(F^{-1}(u))=\lim_{n\to\infty}F(x_n)\geq u.$.
  \end{proof}

  \begin{question}
   Suppose $T>0$. Assume $F:[0,T]\to \r$ is left continuous and $F(0)=0$. Then the signed measure $\mu$ induced by $F$ has no atoms, where $\mu[0,t]=F(t),\forall t\in[0,T]$.
  \end{question}
  \begin{proof}
   {\itshape Actually, $F:[0,T]\to \r$ must be continuous. In fact, the left continuity is given. The right continuity is obtained by noticing $F(t)=\mu[0,t]$.}

   For each $t\in(0,T]$, consider $[0,t-1/n]$ for $n>1/t$. Then we have
   \begin{equation*}
    \begin{split}
     \mu(\{t\})&=\lim_{n\to\infty}\mu(t-1/n,t]=\mu[0,t]-\lim_{n\to\infty}\mu[0,t-1/n] \\
     &=F(t)-\lim_{n\to\infty}F(t-1/n)=F(t)-F(t-)=0.
    \end{split}
   \end{equation*}
   For $t=0$, we have $\mu(\{0\})=\mu[0,0]=F(0)=0$. Therefore $\mu(\{t\})=0$ for each $t\in[0,T]$. Namely, the signed measure $\mu$ induced by $F$ has no atoms.
  \end{proof}

  \begin{question}
   Suppose $F:\r_+\to \r$ is a right continuous non-decreasing function. Then exists a unique $\sigma$-finite measure $\mu$ on $(\r_+,\b(\r_+))$ such that for each $T>0,\mu|_{[0,T]}$ is the measure induced by $F|_{[0,T]}$, denoted as $\mu_F|_{[0,T]}$ (since all the definitions on finite interval $[0,T]$ is already done).
  \end{question}
  \begin{proof}Define
   \begin{equation*}
    \mathscr{C}=\{(a,b]:0<a\leq b\leq \infty \}\cup\{[0,c]:c\geq 0\}.
   \end{equation*}
   Then $\mathscr{C}$ is a semi-set algebra on $\r_+$ and $\sigma(\mathscr{C})=\b(\r_+)$. Define
   \begin{equation*}
    \begin{split}
     \mu(\r_+)&:=\lim_{t\to\infty}F(t)-F(0),\quad \text{the limit exists since $F$ is non-decreasing,}\\
     \mu(a,b]&:=F(b)-F(a),\quad \mu[0,c]:=F(c)-F(0).
    \end{split}
   \end{equation*}
   Then $\mu$ is indeed a finite measure on $\mathscr{C}$. So the {\sffamily Measure extension theorem} shows exists a unique measure $\mu$ on $(\r_+,\b(\r_+))$. Moreover $\mu$ is $\sigma$-finite obviously. Since for each $T>0$, consider
   \begin{equation*}
    \mathscr{C}_T=\{(a,b]:0<a\leq b\leq T_1\}\cup\{[0,c]:0\leq c\leq T_1\}.
   \end{equation*}
    Then $\forall A\in\mathscr{C}_T$, we have $\mu(A)=\mu_{F}|_{[0,T]}(A)$. The uniqueness of {\sffamily Measure extension theorem} still yields $\mu|_{[0,T]}=\mu_{F}|_{[0,T]}$.
  \end{proof}

  \begin{question}
   Suppose $\Omega$ is the sample space and $\mu$ is a measure on semi-set algebra $\a$. Denote also by $\mu$ the extension of $\mu$ to $\sigma(\a)$. Define the inner measure as follows:
   \begin{equation*}
    \mu_*(A)=\sup \left\{\sum_n \mu(A_n):\{A_n\}\subset\a\text{ and }\bigcup_n A_n\subset A \right\}.
   \end{equation*}
   Then
   \begin{equation*}
    \mu^*(A)=\inf\{\mu(B):A\subset B,B\in\sigma(\a) \}=:\lambda ^*(A)
   \end{equation*}
   and
   \begin{equation*}
    \mu_*(A)=\sup\{\mu(C):C\subset A,C\in\sigma(\a) \}=:\lambda_*(A).
   \end{equation*}
   Moreover, the infimum and supremum are both achieved. Finally, we have $A$ is $\mu^*$-measurable iff $\mu^*(A)=\mu_*(A)$.
  \end{question}
  \begin{proof}
   For each $A\subset \Omega$, by the definition of the outer measure, for each $\varepsilon>0$, exists $\{A_n\}\subset\a$ and $A\subset\cup_n A_n$ such that
   \begin{equation*}
    \sum_{n}\mu(A)<  \mu^*(A)+\varepsilon.
   \end{equation*}
   Note $\cup_n A_n\in\sigma(\a)$, then
   \begin{equation*}
    \lambda^*(A)\leq \mu\(\bigcup_n A_n \)\leq \sum_{n}\mu(A)<  \mu^*(A)+\varepsilon.
   \end{equation*}
   The arbitrariness of $\varepsilon$ shows $\lambda^*(A)\leq \mu^*(A)$. On the contrary, for each $B\in\sigma(\a),A\subset B$, we have
   \begin{equation*}
    \mu^*(A)\leq \mu^*(B)=\mu(B).
   \end{equation*}
   Take $\inf_{A\subset B,B\in\sigma(\a)}$ on both sides of the above inequality yields $\mu^*(A)\leq \lambda^*(A)$. In conclusion, we have $\mu^*(A)=\lambda^*(A)$. Let $\varepsilon=1/k$, then exists $\{A_n^k\}\subset \a$ cover $A$ and
   \begin{equation*}
    \mu^*(A)\leq \mu\(\bigcup_n A_n^k \)<\mu^*(A)+1/k.
   \end{equation*}
   Then consider $B=\cap_k\cup_n A_n^k$, then obviously
   \begin{equation*}
    \mu^*(A)=\mu(B).
   \end{equation*}
   The results for the inner measure follows similarly.

   If $A$ is $\mu^*$-measurable, then obviously $\mu^*(A)=\mu_*(A)$. On the contrary, if $\mu^*(A)=\mu_*(A)$, since the above infimum and supremum are both achieved, then we can choose $A_1,A_2\in\sigma(\a),A_1\subset A\subset A_2$ and $\mu(A_1)=\mu(A_2)$. For each $E\subset\Omega$, choose $B\in\sigma(\a)$ such that $\mu^*(E)=\mu(B)$. Then
   \begin{equation*}
    \begin{split}
     \mu^*(A\cap E)&+\mu^*(A^c\cap E)\leq \mu(A_2\cap B)+\mu(A_1^c\cap B)\\
     &\leq \mu(B)+\mu(A_2\backslash A_1 )=\mu(B)=\mu^*(E).
    \end{split}
   \end{equation*}
   So $A$  is $\mu^*$-measurable.
  \end{proof}
  
  \begin{question}
    If $(X_n)_{n\geq 1}$ is a random variable sequence on the probability space $(\Omega,\f,\p )$. Assume $\g\subset\f$ and $X_n\overset{L^p}{\to} X$, where $p\geq 1$. Then
    \begin{equation*}
      \p(X_n|\g)\overset{L^p}{\to} \p(X|\g).
    \end{equation*}
  \end{question}
  \begin{proof}
    Since we have 
    \begin{equation*}
      \begin{split}
         \p[|\p(X_n|\g)- \p(X|\g)|^p] & \leq \p[\p(|X_n-X|^p|\g)] \\
           & =\p(|X_n-X|^p)\to 0,\quad n\to\infty.
      \end{split}
    \end{equation*}
    So $\p(X_n|\g)\overset{L^p}{\to} \p(X|\g).$
  \end{proof}
  
  \begin{question}
  	Presume $X\geq 0$ is a random variable on the probability space $(\Omega,\f,\p)$, then we have 
  	\begin{equation*}
  		\p(X)=\int_0^\infty \p(X> x)dx=\int_0^\infty [1-F(x)]dx.
  	\end{equation*}
  	Prove when $X$ is not nonnegative, we have the following expression:
  	\begin{equation*}
  		\p(X)=\int_0^\infty [1-F(x)]dx+\int_{-\infty}^0 F(x)dx.
  	\end{equation*}
  \end{question}
  \begin{proof}
  	We have
  	\begin{equation*}
  		\begin{split}
  			\p(X)&=\int Xd\p=\int \mathbbm{1}_{\{X\geq 0\}}Xd\p+\int \mathbbm{1}_{\{X< 0\}}Xd\p\\
  			&=\int \mathbbm{1}_{\{X\geq 0\}}\int_0^\infty \mathbbm{1}_{\{x<X\}}dx d\p+\int \mathbbm{1}_{\{X< 0\}}\int_{-\infty}^0\mathbbm{1}_{\{x\geq X\}} d\p\\
  			&=\int \int_0^\infty \mathbbm{1}_{\{x<X\}}dx d\p+\int \int_{-\infty}^0\mathbbm{1}_{\{x\geq X\}} dxd\p\\
  			&=\int_0^\infty \p(X>x)dx+\int_{-\infty}^0 \p(X\leq x)dx\\
  			&=\int_0^\infty [1-F(x)]dx+\int_{-\infty}^0 F(x)dx.
  		\end{split}
  	\end{equation*}
  \end{proof}
  
  \begin{question}
   Consider measurable space $(\r^n,\b^n)$. Then $\mu_k\overset{v}{\to}\mu$ iff for all the continuous intervals $I$ of $\mu$, $\mu_k(I)\to \mu(I)$.
  \end{question}
  \begin{proof}
   \begin{enumerate}[\itshape (1)]
    \item Sufficiency. Suppose for all the continuous intervals $I$ of $\mu$, $\mu_k(I)\to \mu(I)$, then for all $f\in C_0(\r^n)$, exist $M>0,\st |f(x)|\leq M$. Note $C(\mu)$ is dense in $\r^n$, we can find a continuous interval $I$ of $\mu$, such that $\mathrm{supp}(f)\subset I$. Moreover, notice $f$ is uniformly continuous on $I$, then $\forall\varepsilon>0,\exists \delta>0,\forall |x-y|<\delta$, then $|f(x)-f(y)|<\varepsilon$. Thus we can find a finite partition of $I$, say $\{I_\ell\}_{\ell=1}^n$, such that every $I_\ell$ is the continuous interval of $\mu$, the radius of $I_k$ is lease than $\delta$ so that $\forall x,y\in I_\ell,|f(x)-f(y)|<\varepsilon$. Define the following simple function:
    \begin{equation*}
     g=\sum_{\ell=1}^nf(\alpha_\ell)\i_{I_\ell},
    \end{equation*}
    where $\alpha_\ell\in I_\ell$. Then $|f(x)-g(x)|<\varepsilon,\forall x\in I$ and
    \begin{equation*}
     \begin{split}
      &\bigg|\int fd\mu_k-\int fd\mu\bigg|=\bigg|\int_I fd\mu_k-\int_I fd\mu\bigg|\\
      \leq &\bigg|\int_I fd\mu_k-\int_I gd\mu_k\bigg|+\bigg|\int_I gd\mu_k-\int_I gd\mu\bigg|+\bigg|\int_I gd\mu-\int_I fd\mu\bigg|\\
      \leq &\int_I|f-g|d\mu_k+\sum_{\ell=1}^nf(\alpha_\ell)|\mu_k(I_\ell)-\mu(I_\ell)|+\int_I|f-g|d\mu\\
      <&\varepsilon\[\mu_k(I)+\mu(I)\]+M\sum_{\ell=1}^n|\mu_k(I_\ell)-\mu(I_\ell)|.
     \end{split}
    \end{equation*}
    Let $k\to\infty$ and $\varepsilon\to0$ to obtain
    \begin{equation*}
     \varlimsup_{k\to\infty}\bigg|\int fd\mu_k-\int fd\mu\bigg|=0.
    \end{equation*}
    \item Necessity. For all continuous interval $I$ of $\mu$, write $I=[a,b]$. Then $\forall\varepsilon>0,\exists\delta>0$ small enough, such that $a+\delta<b-\delta$ and $\mu(U)<\varepsilon$, where
    \begin{equation*}
     U=(a-\delta,a+\delta )\cup (b-\delta,b+\delta ).
    \end{equation*} 
     We can construct two functions $f,g\in C_0(\r^n),\st f\leq \i_I\leq g$ and $|f-g|\leq I_U$. Thus
     \begin{equation*}
      \int fd\mu_k \leq \mu_k(I)\leq \int gd\mu_k
     \end{equation*}
     and
     \begin{equation*}
      \int fd\mu \leq \mu(I)\leq \int gd\mu.
     \end{equation*}
     Combine the above two formula to obtain
     \begin{equation*}
      |\mu_k(I)-\mu(I)|\leq \min\bigg\{\Big| \int gd\mu-\int fd\mu_k\Big|,\Big|\int gd\mu_k-\int fd\mu\Big| \bigg\}.
     \end{equation*} 
     Let $k\to\infty$ to obtain
     \begin{equation*}
      \begin{split}
       \varlimsup_{k\to\infty}|\mu_k(I)-\mu(I)|\leq&\min\bigg\{\Big|\int gd\mu-\int fd\mu\Big|,\Big|\int gd\mu-\int fd\mu\Big| \bigg\}\\
       =&\bigg|\int gd\mu-\int fd\mu\bigg|\leq \int |f-g|d\mu\\
       \leq &\int \i_Ud\mu=\mu(U)<\varepsilon.
      \end{split}
     \end{equation*}
     Let $\varepsilon\to0$, we will obtain 
     \begin{equation*}
      \varlimsup_{k\to\infty}|\mu_k(I)-\mu(I)|=0.
     \end{equation*}
   \end{enumerate}
  \end{proof}
  
  \begin{question}
  	Given finite measure sequence $\{\mu_n\}_{n\in\n}$ and $\mu$ on measurable space $(\Omega ,\f)$. Prove the following assertions.
  	\begin{enumerate}
  		\item The uniformly convergence: 
  		\begin{equation*}
  			\sup_{A\in\f }|\mu_n(A)-\mu(A)|\to0,\quad n\to\infty,
  		\end{equation*} 
  		is equivalent to 
  		\begin{equation*}
  			\sup_{f\in\f,|f|\leq 1}|\mu_n(f)-\mu(f)|\to 0,\quad n\to\infty.
  		\end{equation*}
  		\item The strong convergence: for any $A\in\f$,
  		\begin{equation*}
  			|\mu_n(A)-\mu(A)|\to0,\quad n\to\infty,
  		\end{equation*} 
  		is equivalent to for any $f\in \mathrm{b}\f$
  		\begin{equation*}
  			|\mu_n(f)-\mu(f)|\to 0,\quad n\to\infty.
  		\end{equation*}
  	\end{enumerate}
  \end{question}
  \begin{proof}
  	\begin{enumerate}
  		\item For the sufficiency, since $1_A\in\f,|1_A|\leq 1$, then
  		\begin{equation*}
  			\begin{split}
  				\sup_{A\in\f }|\mu_n(A)-\mu(A)|&=\sup_{A\in\f }|\mu_n(1_A)-\mu(1_A)|\\
  				&\leq \sup_{f\in\f,|f|\leq 1}|\mu_n(f)-\mu(f)|\to 0.
  			\end{split}
  		\end{equation*}
  		For the necessity, firstly let $A=\Omega$. Then $\mu_n(\Omega)\to\mu(\Omega)$ implies exists $M>0$ such that 
  		\begin{equation*}
  			\sup_{n\in\n}\mu_n(\Omega)\vee \mu(\Omega)\leq M.
  		\end{equation*}
  		For any $f\in \f,|f|\leq 1$, we can find increasing simple functions:
  		\begin{equation*}
  			f_m=\sum_{k=0}^{2^m}k/2^m\mathbbm{1}_{\{k/2^m<f\leq (k+1)/2^m  \}}=:\sum_{k=0}^{2^m}k/2^m\mathbbm{1}_{A_k}\uparrow f,\quad m\to\infty
  		\end{equation*}
  		and $||f_m-f||_\infty\leq 1/2^m$. Therefore 
  		\begin{equation*}
  			\begin{split}
  				|\mu_n(f)&-\mu(f)|\\
  				&\leq |\mu_n(f)-\mu_n(f_m)|+|\mu_n(f_m)-\mu(f_m)|+|\mu(f_m)-\mu(f)|\\
  				&\leq \mu_n(|f-f_m|)+\sum_{k=0}^{2^m}k/2^m|\mu_n(A_k)-\mu(A_k)|+\mu(|f-f_m|)\\
  				&\leq 1/2^m M+\sup_{A\in\f} |\mu_n(A)-\mu(A)|\sum_{k=0}^{2^m}k/2^m+1/2^m M.
  			\end{split}
  		\end{equation*}
  		It follows from the above formula that
  		\begin{equation*}
  			\sup_{f\in\f,|f|\leq 1}|\mu_n(f)-\mu(f)|\leq M/2^{m-1}+\sup_{A\in\f}|\mu_n(A)-\mu(A)|\sum_{k=0}^{2^m}k/2^m.
  		\end{equation*}
  		Let $n\to\infty$ first, and then let $m\to\infty$, we have 
  		\begin{equation*}
  			\sup_{f\in\f,|f|\leq 1}|\mu_n(f)-\mu(f)|\to 0,\quad n\to\infty.
  		\end{equation*}
  		\item The proof of the second assertion is nearly a copy of the first one.
  	\end{enumerate}
  \end{proof}
  
  \begin{question}
  	Given a measurable space $(E,\b(E))$, where $E$ is a metric space with metric $d(\cdot,\cdot)$, and $\{\mu_n\}_{n\geq 1},\mu$ are uniformly bounded measure on $(E,\b(E))$. Then the following statements are all equivalent.
  	\begin{enumerate}
  		\item $\mu_n\overset{\mathrm{w}}{\to}\mu,n\to\infty$.
  		\item $\mu_n(B)\to\mu(B)$ for all $\mu$-continuous closed set $B\in\b(E)$.
  		\item $\mu_n(C)\to\mu(C)$ for all $\mu$-continuous open set $C\in\b(E)$.
  	\end{enumerate}
  \end{question}
  \begin{proof}
  	
  	\textcolor{red}{\sffamily Way 1.} 
  	First we show $2\Rightarrow 3$. Since $E$ is naturally the $\mu$-continuous closed/open set. Then $\mu_n(E)\to\mu(E)$. For any $\mu$-continuous open set $C\in\b(E)$, we have $\mu(\partial C)=0$. Note $B:=C^c$ is closed and $\partial B=\partial C$. So $\mu(\partial B)=\mu(\partial C)=0$. Therefore
  	\begin{equation*}
  		\begin{split}
  			\mu_n(C)&=\mu_n(E)-\mu_n(C^c)=\mu_n(E)-\mu_n(B)\\
  			&\to \mu(E)-\mu(B)=\mu(E)-\mu(C^c)=\mu(C).
  		\end{split}
  	\end{equation*}
  	Similarly we can prove $3\Rightarrow 2$. Then we show $1\Leftrightarrow 2(3)$. If $\mu_n\overset{\mathrm{w}}{\to}\mu$, then $\mu_n(A)\to\mu(A)$ for all $\mu$-continuous open/closed set $A\in\b(E)$ naturally holds. We just prove the opposite. If $\mu_n(B)\to\mu(B)$ for all $\mu$-continuous open/closed set $B\in\b(E)$, then for any $\mu$-continuous set $A$, we have
  	\begin{equation*}
  		\varlimsup_{n\to\infty} \mu_n(A)\leq \varlimsup_{n\to\infty}\mu_n(\overline{A})=\lim_{n\to\infty}\mu_n(\overline{A})= \mu(\overline{A})=\mu(A)
  	\end{equation*}
  	and
  	\begin{equation*}
  		\varliminf_{n\to\infty} \mu_n(A)\geq \varliminf_{n\to\infty}\mu_n(\mathring{A})=\lim_{n\to\infty}\mu_n(\mathring{A})= \mu(\mathring{A})=\mu(A).
  	\end{equation*}
  	Thus $\lim_{n\to\infty}\mu_n(A)$ exists and equals to $\mu(A)$. So $\mu_n\overset{w}{\to}\mu,n\to\infty$.
  	
    \textcolor{red}{\sffamily Way 2.} First we show $1\Leftrightarrow 2$. If $\mu_n\overset{\mathrm{w}}{\to}\mu$, then $\mu_n(A)\to\mu(A)$ for all $\mu$-continuous closed set $A\in\b(E)$ naturally holds. We just prove the opposite. For any open set $A\in\b(E)$, define
  	\begin{equation*}
  		K_r:=\{x\in E:d(x,A^c)\geq r\},\quad r>0.
  	\end{equation*}
  	Then $K_r$ is closed, $K_r\subset A$ for any $r>0$ and $K_r\uparrow A,r\to 0$. Now we claim there exists $\{r_n\}\subset (0,\infty)$ such that $\{K_{r_n}\}_{n\geq 1}$ are all $\mu$-continuous sets. In fact, since we have proved $\partial f^{-1}(A)\subset f^{-1}(\partial A)$ for any $A\in\b(E)$ in the topology part, then 
  	\begin{equation*}
  		\partial K_r\subset \{x\in E:d(x,A^c)=r\}.\tag{$*$}
  	\end{equation*}
  	We say 
  	\begin{equation*}
  		\{r>0:\mu(\partial K_r)>0\}=\bigcup_{m\geq 1}\{r>0:\mu(\partial K_r)>1/m\}
  	\end{equation*}
  	is at most countable. If not, there exists $m\geq 1$, such that $\{r>0:\mu(\partial K_r)>1/m\}$ is uncountable. So we can arbitrarily choose a sequence $\{s_n\}_{n\geq 1}\subset \{r>0:\mu(\partial K_r)>1/m\}$. Moreover, formula $(*)$ shows for different $r$, the boundary $\partial K_r$ is disjoint. So
  	\begin{equation*}
  		\mu\(\sum_{n\geq 1}\partial K_{s_n} \)=\sum_{n\geq 1}\mu(\partial K_{s_n})>\sum_{n\geq 1}1/m=\infty.
  	\end{equation*}
  	This contradicts the fact that  $\mu$ is bounded. So $\{r>0:\mu(\partial K_r)>0\}$ is at most countable. Therefore we can choose $\{r_k\}_{k\geq 1}\in \{r>0:\mu(\partial K_r)=0\}$ and $r_k\downarrow 0,k\to\infty$. So $\{K_{r_k}\}_{k\geq 1}$ are naturally all $\mu$-continuous sets. Finally,
  	\begin{equation*}
  		\varliminf_{n\to\infty}\mu_n(A)\geq \varliminf_{n\to\infty}\mu_n(K_{r_k})=\lim_{n\to\infty}\mu_n(K_{r_k})=\mu(K_{r_k}).
  	\end{equation*}
  	Let $k\to\infty$ in above formula to obtain
  	\begin{equation*}
  		\varliminf_{n\to\infty}\mu_n(A)\geq \mu\(\bigcup_{k\geq 1}K_{r_k} \)=\mu(A).
  	\end{equation*}
  	This means $\mu_n\overset{\mathrm{w}}{\to}\mu,n\to\infty$.
  	
  	 So the proof is complete.
  \end{proof}
  
  \begin{question}
  	Suppose $\{a_i\}_{i\in\n}$ and $\{a_i^{(n)}\}_{i\in\n},n\in\n$ are probability distributions. Then 
  	\begin{equation*}
  		\lim_{n\to\infty}a_i^{(n)}= a_i,\forall i\in\n\Longleftrightarrow \lim_{n\to\infty}\sum_{i=0}^\infty s^ia_i^{(n)}=\sum_{i=0}^\infty s^i a_i,\quad \forall s\in(0,1).
  	\end{equation*}
  \end{question}
  \begin{proof}
  	The necessity part can be simply obtained by using dominance convergence theorem. We just prove the sufficiency part. We use the induction on $i$ to verify this part.
  	\begin{enumerate}
  		\item When $i=0$, since
  	\begin{equation*}
  		\begin{split}
  			\sum_{i=0}^\infty s^ia_i^{(n)}-\sum_{i=0}^\infty s^i a_i&=\sum_{i=1}^\infty s^i(a_i^{(n)}-a_i)+(a_0^{(n)}-a_0)\\
  			&\leq \sum_{i=1}^\infty s^i+(a_0^{(n)}-a_0)\\
  			&=\frac{s}{1-s}+(a_0^{(n)}-a_0).
  		\end{split}
  	\end{equation*}
  	Then
  	\begin{equation*}
  		0=\varliminf_{n\to\infty}\(\sum_{i=0}^\infty s^ia_i^{(n)}-\sum_{i=0}^\infty s^i\)\leq \frac{s}{1-s}+\varliminf_{n\to\infty} (a_0^{(n)}-a_0).
  	\end{equation*}
  	Let $s\to0+$ to obtain
  	\begin{equation*}
  		0\leq \varliminf_{n\to\infty} (a_0^{(n)}-a_0).
  	\end{equation*}
  	Interchange the position of $a_0^{(n)}$ and $a_0$, we will get 
  	\begin{equation*}
  		\varliminf_{n\to\infty} (a_0^{(n)}-a_0)= \varlimsup_{n\to\infty} (a_0^{(n)}-a_0)=0,
  	\end{equation*}
  	which is equivalent to $\lim_{n\to\infty}a_0^{(n)}=a_0$.
  		\item Suppose when $i\leq k$, we have $\lim_{n\to\infty}a_i^{(n)}=a_i$. When $i=k+1$, the induction hypothesis ensure 
  		\begin{equation*}
  			\lim_{n\to\infty}\sum_{i=k+1}^\infty s^ia_i^{(n)}=\sum_{i=k+1}^\infty s^i a_i
  		\end{equation*}
  		and furthermore, 
  		\begin{equation*}
  			\lim_{n\to\infty}\sum_{i=k+1}^\infty s^{i-(k+1)}a_i^{(n)}=\sum_{i=k+1}^\infty s^{i-(k+1)} a_i.
  		\end{equation*}
  		Then the similar argument shows $\lim_{n\to\infty}a_{k+1}^{(n)}=a_{k+1}$.
  		\item Hence the argument is correct for all $i\in\n$.
  	\end{enumerate}
  \end{proof}
  
  \begin{question}
  	A random variable $X$ on probability space $(\Omega,\f,\p)$ is distributed as a constant $c$ iff $X\overset{\ae}{=} c$. 
  \end{question}
  \begin{proof}
  	Obviously, if $X\overset{\ae}{=} c$, then $X\sim c$. On the contrary, if $X\sim c$, then we have for any $n\geq 1$,
  	\begin{equation*}
  		\p(c-1/n\leq X\leq c+1/n)=\p(c-1/n\leq c\leq c+1/n)=1.
  	\end{equation*}
  	Then let $\varepsilon\to 0+$, we will obtain
  	\begin{equation*}
  		\p(X=c)=1,
  	\end{equation*}
  	which is equivalent to $X\overset{\ae}{=} c$.
  \end{proof}
  
  \begin{question}[The generalization of MCT]
  	Given measurable functions $\{f_t\}_{t\in I}$ on measure space $(E,\e,\mu)$, where $I=\n,I=\r_+$ (or $I=[a,b]$). If $\{f_t\}_{t\in I}$ is monotone and exists $s\in I$ (or $s\in [a,b)$) such that $f_s\in L^1(\mu)$, then 
  	\begin{equation*}
  		\int_E f_td\mu\to \int_E fd\mu,\quad t\to\infty \text{ (or }t\to b).
  	\end{equation*}
  \end{question}
  \begin{proof}
  	If $\{f_t\}_{t\in I}$ is increasing, then $f_s\leq f_t$ for any $I\ni t\geq s$. Then the normal MCT gives us the desired results. If $\{f_t\}_{t\in I}$ is decreasing, then $f_s\geq f_t\Longrightarrow -f_s\leq -f_t$ for any $I\ni t\geq s$. Then the normal MCT also gives us the desired results.
  \end{proof}
  
  \begin{question}
   	Suppose $(\Omega,\f,\p)$ is a probability space and $A\notin\f$. Let $\f_1:=\sigma(\f\cup\{A\})$. Prove $\p$ can be extended to $(\Omega,\f_1)$.
   \end{question}
   \begin{proof}
   	\begin{enumerate}
   		\item For any $D\in\f$, we define
   	\begin{equation*}
   		\widetilde{\p}(A\cap D):=\inf\{\p(E):A\cap D\subset E,E\in\f\}.
   	\end{equation*}
   	 Next we show exists $E_D\in\f,A\cap D\subset E_D$ such that $\p(E_D)=\widetilde{\p}(A\cap D)$. In fact, by the definition of infimum, exists $\{E_n\}_{n\geq 1}\subset\f$ satisfying $A\cap D\subset E_n$ for any $n\geq 1$ and $\p(E_n)\downarrow \widetilde{\p}(A\cap D)$. Define $E_D:=\varliminf_{n\to\infty} E_n\in\f$, then 
   	 \begin{equation*}
   	 	E_D=\varliminf_{n\to\infty} E_n=\bigcup_{n\geq 1}\bigcap_{k\geq n}E_k\supset \bigcup_{n\geq 1}\bigcap_{k\geq n} (A\cap D)=A\cap D
   	 \end{equation*}
   	 and
   	 \begin{equation*}
   	 	\widetilde{\p}(A\cap D)\leq \p(E_D)=\lim_{n\to\infty}\p\bigg(\bigcap_{k\geq n}E_k \bigg)\leq \lim_{n\to\infty}\p(E_n)=\widetilde{\p}(A\cap D).
   	 \end{equation*}
   	 This shows $\p(E_D)=\widetilde{\p}(A\cap D)$. We can additionally presume $E_D\subset D$. In fact, consider $E_D\cap D$ is enough.
   		\item Moreover, we show for any $(A\cap D_1)\cap (A\cap D_2)=\emptyset$, we have 
   	 \begin{equation*}
   	 	\widetilde{\p}[(A\cap D_1)+(A\cap D_2)]=\widetilde{\p}[A\cap (D_1\cup D_2)]=\widetilde{\p}(A\cap D_1)+\widetilde{\p}(A\cap D_2).
   	 \end{equation*}
   	 In fact, exist $E_1,E_2\in\f$ such that $\p(E_1)=\widetilde{\p}(A\cap D_1)$ and $\p(E_2)=\widetilde{\p}(A\cap D_2)$. 
   	 We claim 
   	 \begin{equation*}
   	 	\p(E_1\cap E_2)=0.
   	 \end{equation*}
   	  Otherwise, since $(A\cap D_1)\cap(A\cap D_2)=\emptyset$ and $A\cap D_i\subset E_i,i=1,2$, then 
   	  \begin{equation*}
   	  	\f\ni E_1\cap E_2\subset [A\cap (D_1\cup D_2)]^c
   	  \end{equation*}
   	  and $\p(E_1\cap E_2)>0$. Therefore $\hat{ E}_i:=E_i-E_1\cap E_2$ satisfying 
   	  \begin{equation*}
   	  	A\cap D_i\subset \hat{ E}_i,\qquad\p(\hat{ E}_i)=\p (E_i)-\p(E_1\cap E_2)<\p(E_i),\quad i=1,2,
   	  \end{equation*}
   	  which contradicts the definition of $E_i,i=1,2$. Since we have proved $\p(E_1\cap E_2)=0$, then
   	  \begin{equation*}
   	  	\p(E_1\cup E_2)=\p(E_1)+\p(E_2)-\p(E_1\cap E_2)=\p(E_1)+\p(E_2).
   	  \end{equation*}
   	  
   	  
   	 Since $A\cap D_1\subset E_1$ and $A\cap D_2\subset E_2$, then
   	 \begin{equation*}
   	 	A\cap (D_1\cup D_2)\subset E_1\cup E_2.
   	 \end{equation*}
   	 Therefore
   	 \begin{equation*}
   	 	\begin{split}
   	 		\widetilde{\p}[A\cap (D_1\cup D_2)]&\leq \widetilde{\p}(E_1\cup E_2)=\p(E_1)+\p(E_2)\\
   	 		&=\widetilde{\p}(A\cap D_1)+\widetilde{\p}(A\cap D_2).
   	 	\end{split} 
   	 \end{equation*}
   	 On the contrary, exists $E_{12}\in\f$ such that $A\cap (D_1\cup D_2)\subset E_{12}\subset D_1\cup D_2$ and $\p(E_{12})=\widetilde{\p}[A\cap (D_1\cup D_2)]$. Consider $E_i':=E_{12}\cap D_i,i=1,2$. Since $A\cap (D_1\cup D_2)\subset E_{12}$, then $A\cap D_i\subset E'_{i},i=1,2$ and $E_{12}=E'_{1}\cup E'_{2}$. We claim $\p(E'_{1}\cap E'_{2})=0$. In fact, if $\p(E'_{1}\cap E'_{2})>0$, since $(A\cap D_1)\cap(A\cap D_2)=\emptyset$, then we must have
   	 \begin{equation*}
   	 	E'_{1}\cap E'_{2}\subset [A\cap (D_1\cup D_2)]^c.
   	 \end{equation*}
   	 Consider $E_{12}'=E_{12}-E_1'\cap E_2'$, then $A\cap (D_1\cup D_2)\subset E_{12}'$ and
   	 \begin{equation*}
   	 	\p(E_{12}')=\p(E_{12})-\p(E_1'\cap E_2')<\p(E_{12}),
   	 \end{equation*}
   	 which contradicts the definition of $E_{12}$. So $\p(E_{12})=\p(E'_{{1}})+\p(E'_{{2}})$ and therefore
   	 \begin{equation*}
   	 	\begin{split}
   	 		\widetilde{\p}[A\cap (D_1\cup D_2)]&=\p(E_{12})=\p(E'_{{1}})+\p(E'_{{2}})\\
   	 		&\geq \p(E_{{1}})+\p(E_{{2}})=\widetilde{\p}(A\cap D_1)+\widetilde{\p}(A\cap D_2).
   	 	\end{split}
   	 \end{equation*}
   	 The above two formulas show for any $D_1,D_2\in\f,(A\cap D_1)\cap (A\cap D_2)=\emptyset$,
   	 \begin{equation*}
   	 	\widetilde{\p}[A\cap (D_1\cup D_2)]=\widetilde{\p}(A\cap D_1)+\widetilde{\p}(A\cap D_2).
   	 \end{equation*}
   		\item Next we show for any disjoint sequence $\{A\cap D_n\}_{n\geq 1}$, we have
   		\begin{equation*}
   			\widetilde{\p}\(\sum_{n\geq 1}A\cap D_n\)=\sum_{n\geq 1}\widetilde{\p}(A\cap D_n).
   		\end{equation*}
   		We have shown exists $E_{\geq 1}\in\f$ such that 
   		\begin{equation*}
   			\sum_{n\geq 1}A\cap D_n\subset E_{\geq 1}\subset \bigcup_{n\geq 1}D_n\text{ and }\p(E_{\geq 1})=\widetilde{\p}\(\sum_{n\geq 1}A\cap D_n\)
   		\end{equation*}
   		 and exist $\{E_n\}_{n\geq 1}\subset\f$ such that $A\cap D_n\subset E_n\subset D_n$,
   		\begin{equation*}
   			\widetilde{\p}(A\cap D_n)=\p(E_n)\text{ and }\p(E_n\cap E_m)=0,\quad \forall n\neq m\geq 1.
   		\end{equation*}
   		Observe $\sum_{n\geq 1} A\cap D_n \subset \cup_{n\geq 1}E_n$, then
   		\begin{equation*}
   			\widetilde{\p}\(\sum_{n\geq 1}A\cap D_n\)=\p(E_{\geq 1})\leq \p\(\bigcup_{n\geq 1}E_n\)\leq \sum_{n\geq 1}\p(E_n)= \sum_{n\geq 1}\widetilde{\p}(A\cap D_n)
   		\end{equation*}
   		Define $E_n':=E_{\geq 1}\cap D_n$, then $A\cap D_n\subset E_n'\subset D_n$ for any $n\geq 1$. We claim $\p(E'_{i}\cap E'_{j})=0$ for any $i\neq j\geq 1$. In fact, if $\p(E'_{i}\cap E'_{j})>0$, since $\cap_{n\geq 1}A\cap D_n=\emptyset$, then we must have
   	 \begin{equation*}
   	 	E'_{i}\cap E'_{j}\subset [(A\cap D_i)+(A\cap D_j)]^c\text{ and } E'_{i}\cap E'_{j}\subset D_i\cap D_j.
   	 \end{equation*}
   	 This must yields
   	 \begin{equation*}
   	 	E'_{i}\cap E'_{j}\subset\[\sum_{n\geq 1}A\cap D_n\]^c.
   	 \end{equation*}
   	 Consider $\hat{E}_{\geq 1}=E_{\geq 1}-E_i'\cap E_j'$, then $\sum_{n\geq 1}A\cap D_n\subset \hat{E}_{\geq 1}$ and
   	 \begin{equation*}
   	 	\p(\hat{E}_{\geq 1})=\p(E_{\geq 1})-\p(E_i'\cap E_j')<\p(E_{\geq 1}),
   	 \end{equation*}
   	 which contradicts the definition of $E_{\geq 1}$. So we have proved $\p(E'_{i}\cap E'_{j})=0$ for any $i\neq j\geq 1$.
   	 
   	 Define $H_1=E_1'$ and $H_n=E_n'/ \cup_{k=1}^{n-1} E_k'$ for $n\geq 2$. Then $\{H_n\}_{n\geq 1}\subset\f$ is disjoint and $\p(H_n)=\p(E_n')$. In fact, we have for any $n\geq 2$,
   		\begin{equation*}
   			\p\(E_n'\cap \bigcup_{k=1}^{n-1} E_k'\)=\p\[\bigcup_{k=1}^{n-1} (E_n'\cap E_k')\]\leq \sum_{k=1}^{n-1}\p(E_n'\cap E_k')=0.
   		\end{equation*}
   		Then for any $n\geq 2$,
   		\begin{equation*}
   			\p(H_n)=\p\(E_n'\bigg/ \bigcup_{k=1}^{n-1} E_k'\)=\p(E_n')-\p\(E_n'\cap \bigcup_{k=1}^{n-1} E_k'\)=\p(E_n').
   		\end{equation*}
   		Therefore
   		\begin{equation*}
   			\begin{split}
   				\sum_{n\geq 1}\p(E_n')&=\sum_{n\geq 1}\p(H_n)=\p\(\sum_{n\geq 1}H_n\)=\p\(\bigcup_{n\geq 1}E_n'\)\\
   				&=\p(E_{\geq 1})=\widetilde{\p}\(\sum_{n\geq 1}A\cap D_n\).
   			\end{split}
   		\end{equation*}
   		Finally, we have
   		\begin{equation*}
   			\sum_{n\geq 1}\widetilde{\p}(A\cap D_n)=\sum_{n\geq 1}\p(E_n)\leq \sum_{n\geq 1}\p(E_n')= \widetilde{\p}\(\sum_{n\geq 1}A\cap D_n\).
   		\end{equation*}
   		Combine with the former formula, we have proved 
   		\begin{equation*}
   			\widetilde{\p}\(\sum_{n\geq 1}A\cap D_n\)=\sum_{n\geq 1}\widetilde{\p}(A\cap D_n).
   		\end{equation*}
   		\item Define
   	 \begin{equation*}
   	 	\widetilde{\p}(A^c\cap D):=\sup\left \{\p(F):A^c\cap D\supset F,F\in\f\right\}.
   	 \end{equation*}
   	 We want to show exists $F_D\in\f$ such that $F_D\subset A^c\cap D$ and $F_D=\widetilde{\p}(A^c\cap D)$.
   	 Observe 
   	 \begin{equation*}
   	 	A^c\cap D\supset F \Longleftrightarrow D-(A^c\cap D)\subset D-F 
   	 	\Longleftrightarrow A\cap D\subset D-F
   	 \end{equation*}
   	 and $F=D-(D-F)\ (F\subset D)$. Then 
   	 \begin{equation*}
   	 	\begin{split}
   	 		\widetilde{\p}(A^c\cap D)&=\sup\left \{\p(D)-\p(D-F):A\cap D\subset D-F,F\in\f\right\}\\
   	 		&=\sup\big(\p(D)- \{\p(D-F):A\cap D\subset D-F,F\in\f\}\big)\\
   	 		&=\p(D)-\inf \{\p(D-F):A\cap D\subset D-F,F\in\f\}\\
   	 		&=\p(D)-\inf \{\p(E):A\cap D\subset E,E\in\f\}\\
   	 		&=\p(D)-\widetilde{\p}(A\cap D).
   	 	\end{split}
   	 \end{equation*}
   	 Select $E_D\in\f$ such that $A\cap D\subset E_D\subset D$ and $\p(E_D)=\widetilde{\p}(A\cap D)$. Then 
   	 \begin{equation*}
   	 	\widetilde{\p}(A^c\cap D)=\p(D)-\widetilde{\p}(A\cap D)=\p(D)-\p(E_D)=\p(D-E_D).
   	 \end{equation*}
   	 Then $F_D:=D-E_D\subset A^c\cap D$ is what we want. Besides, we still want to show for any $D_1,D_2\in\f,(A^c\cap D_1)\cap (A^c\cap D_2)=\emptyset$,
   	 \begin{equation*}
   	 	\widetilde{\p}[A^c\cap (D_1\cup D_2)]=\widetilde{\p}(A^c\cap D_1)+\widetilde{\p}(A^c\cap D_2).
   	 \end{equation*} 
   	 In fact, exists $F_1,F_1\in\f$ satisfying $F_i\subset A^c\cap D_i$ and $\p(F_i)=\widetilde{\p}(A^c\cap D_i)$ for $i=1,2$. Observe $F_1\cup F_2\subset A^c\cap (D_1\cup D_2)$, then $(A^c\cap D_1)\cap (A^c\cap D_2)=\emptyset$ yields 
   	 \begin{equation*}
   	 	\begin{split}
   	 		\widetilde{\p}[A^c\cap (D_1\cup D_2)]&\geq \p(F_1\cup F_2)=\p(F_1)+\p(F_2)\\
   	 		&=\widetilde{\p}(A^c\cap D_1)+\widetilde{\p}(A^c\cap D_2).
   	 	\end{split}
   	 \end{equation*}
   	 On the contrary, exists $F_{12}\in\f$ satisfying $F_{12}\subset A^c\cap (D_1\cup D_2)$ and $\p(F_{12})=\widetilde{\p}[A^c\cap (D_1\cup D_2)]$. Define $F_{{i}}':=F_{12}\cap D_i\subset A^c\cap D_i$ for $i=1,2$. Then $(A^c\cap D_1)\cap (A^c\cap D_2)=\emptyset$ yields
   	 \begin{equation*}
   	 	\begin{split}
   	 		\widetilde{\p}[A^c\cap (D_1+D_2)]&=\p(F_{12})=\p(F'_{1})+\p(F'_{2})\\
   	 		&\leq \p(F_1)+\p(F_2)=\widetilde{\p}(A^c\cap D_1)+\widetilde{\p}(A^c\cap D_2).
   	 	\end{split}
   	 \end{equation*}
   	 So we have proved for any $D_1,D_2\in\f,(A^c\cap D_1)\cap (A^c\cap D_2)=\emptyset$,
   	 \begin{equation*}
   	 	\widetilde{\p}[A^c\cap (D_1\cup D_2)]=\widetilde{\p}(A^c\cap D_1)+\widetilde{\p}(A^c\cap D_2).
   	 \end{equation*} 
   	 And similar argument shows for any disjoint $\{A^c\cap D_n\}_{n\geq 1}$, we have
   	 \begin{equation*}
   			\widetilde{\p}\(\sum_{n\geq 1}A^c\cap D_n\)=\sum_{n\geq 1}\widetilde{\p}(A^c\cap D _n).
   		\end{equation*}
   	\item Note $\f_1=\sigma(\f\cup\{A\})=\{(A\cap D_1)+ (A^c\cap D_2):D_1,D_2\in\f\}$, then define
   	 \begin{equation*}
   	 	\widetilde{\p}[(A\cap D_1)+ (A^c\cap D_2)]:=\widetilde{\p}(A\cap D_1)+\widetilde{\p}(A^c\cap D_2).
   	 \end{equation*}
   	 We show $\widetilde{\p}(\cdot)$ is $\sigma$-additive. For any disjoint sets $\{(A\cap D_1^n)+ (A^c\cap D_2^n)\}_{n\geq 1}$, it is easy to find $\{A\cap D_1^n\}_{n\geq 1}$ and $\{A\cap D_2^n\}_{n\geq 1}$ are two sequences of disjoint sets. The former discussion yields
   	 \begin{equation*}
   	 	\begin{split}
   	 		&\widetilde{\p}\left\{\sum_{n\geq 1}\[(A\cap D_1^n)+ (A^c\cap D_2^n)\]\right\}\\
   	 		&\qquad\quad =\widetilde{\p}\[A\cap\(\bigcup_{n\geq 1} D_1^n\)+ A^c\cap\(\bigcup_{n\geq 1} D_2^n\) \]\\\\
   	 		&\qquad\quad =\widetilde{\p}\[A\cap\(\bigcup_{n\geq 1} D_1^n\)\]+\widetilde{\p}\[ A^c\cap\(\bigcup_{n\geq 1} D_2^n\) \]\\
   	 		&\qquad\quad = \widetilde{\p}\[\sum_{n\geq 1}(A\cap D_1^n)\]+\widetilde{\p}\[\sum_{n\geq 1}(A^c\cap D_2^n) \]\\
   	 		&\qquad\quad =\sum_{n\geq 1}[\widetilde{\p} (A\cap D_1^n)+\widetilde{\p}(A^c\cap D_2^n)]\\
   	 		&\qquad\quad =\sum_{n\geq 1}\widetilde{\p}\[(A\cap D_1^n)+ (A^c\cap D_2^n)\].
   	 	\end{split}
   	 \end{equation*}
   	\end{enumerate}
   	Till now, $(\Omega,\f,\p)$ has been extended to $(\Omega,\f_1,\widetilde{\p})$.
   \end{proof}
   
   \begin{question}
   	Given a probability space $(\Omega,\f,\p)$. Suppose $\pi$-system $\a$ satisfies $\Omega\in\a$ and $\f=\sigma(\a)$. Define $\mathcal{D}$ the linear span of $\{1_A:A\in\a\}$. Then $\mathcal{D}$ is dense in $L^p(\Omega,\f,\p)$ for $0<p<\infty$.
   \end{question}
   \begin{proof}
   	Define
   	\begin{equation*}
   		\mathscr{L}=\{f\in L^p(\Omega,\f,\p): \text{exists $\{g_n\}_{n\geq 1}\subset \mathcal{D}$ such that $||g_n-f||_{L^p}\to0$}\}.
   	\end{equation*}
   	First it is easy to know $\{1_A:A\in\a\}\subset \mathscr{L}$. We next shows $\mathscr{L}$ is a $\mathscr{L}$-system. 
   	\begin{enumerate}
   		\item Since $\Omega\in\a$, then $1=1_\Omega\in\mathscr{L}$.
   		\item For any $f_1,f_2\in \mathscr{L}$, then there exists $\{g_{i,n}\}_{n\geq 1}\subset \mathcal{D}$ such that $||g_{i,n}-f_i||_{L^p}\to 0,i=1,2$. Therefore for any $\alpha,\beta\in \r$, we have $\{\alpha g_{1,n}+\beta g_{2,n}\}_{n\geq 1}\subset\mathcal{D}$ and 
   		\begin{equation*}
   			\begin{split}
   				||\alpha g_{1,n}&+\beta g_{2,n}-(\alpha f_{1}+\beta f_{2})||_{L^p}\\
   				&\leq \alpha ||g_{1,n}-f_1||_{L^p}+ \beta ||g_{2,n}-f_2||_{L^p}\to0.
   			\end{split}
   		\end{equation*}
   		So $\alpha f_{1}+\beta f_{2}\in \mathscr{L}$.
   		\item For any $\{f_i\}_{i\geq 1}\subset\mathscr{L},f\in \mathscr{L}$ and $0\leq f_i\uparrow f$, there exists $\{g_{i,n}\}_{n\geq 1}\subset \mathcal{D}$ such that $||g_{i,n}-f_i||_{L^p}\to 0,n\to\infty,i\geq 1$. So for any $i\geq 1$, exists $n_i\geq 1$ such that $||g_{i,n_i}-f_i||_{L^p}<1/i$. Moreover, $0\leq f_n\uparrow f$ means $0\leq f_n^p\uparrow f^p$. Monotoncity convergence theorem yields $||f_n||_{L^p}\to ||f||_{L^p}$. So $f_n^p\uparrow f^p$ and $||f_n||_{L^p}\to ||f||_{L^p}$ yield $||f_n-f||_{L^p}\to0$. Finally, consider $\{g_{i,n_i}\}_{i\geq 1}\subset \mathcal{D}$, we have 
   		\begin{equation*}
   			\begin{split}
   				||g_{i,n_i}-f||&\leq ||g_{i,n_i}-f_i||_{L^p}+||f_i-f||_{L^p}\\
   				&<1/i+||f_i-f||_{L^p}\to 0,\quad i\to\infty.
   			\end{split}
   		\end{equation*}
   		So $f\in \mathscr{L}$.
   	\end{enumerate}
   	Therefore $\mathscr{L}$ is a $\mathscr{L}$-system and $\{1_A:A\in\a\}\subset \mathscr{L}$. The function form monotonicity class theorem yields $\mathscr{L}=L^p(\Omega,\f,\p)$.
   \end{proof}
   
   
   \begin{question}
   	Consider the measure space $(\r^d,\b(\r^d),\mu)$. Define the dual of $\mu$ as $\widetilde{\mu}(A):=\mu(-A)$ for $A\in \b(\r^d)$ and the symmetrization of $\mu$ as $\mu^\sharp:=\mu*\widetilde{\mu}$. Then
   	\begin{enumerate}
   		\item For any $f\in\b(\r^d)$ and $\mu(f)$ exists, we have
   		\begin{equation*}
   			\int f(x)\widetilde{\mu}(dx)=\int f(-x)\mu(dx).
   		\end{equation*}
   		\item Consider product measure $\mu\times\mu$. For any $g\in \b(\r^{2d})$ and $\int gd(\mu\times\mu)$ exists, we have
   		\begin{equation*}
   			\iint g(x,y)\mu(dx)\mu(dy)=\iint g(y,x)\mu(dx)\mu(dy).\tag{$*$}
   		\end{equation*}
   		\item $\mu^\sharp$ is symmetric. Namely, $\mu^\sharp(A)=\mu^\sharp(-A)$ for $A\in \b(\r^d)$.
   	\end{enumerate}
   \end{question}
   \begin{proof}
   	\begin{enumerate}
   		\item For $f=1_A$, where $A\in \b(\r^d)$, we can easily find the assertion is correct. For the general case, we can just apply the standard 'simple function', 'nonnegative measurable function' and 'general measurable function' procedure to obtain the desired result.
   		\item For any $A,B\in\b(\r^d)$, we have
   		\begin{equation*}
   			\iint \mathbbm{1}_{A\times B}(x,y)\mu(dx)\mu(dy)=\mu(A)\mu(B)=\iint \mathbbm{1}_{A\times B}(y,x)\mu(dx)\mu(dy).
   		\end{equation*}
   		Write $\g:=\{C\in\b(\r^{2d}):,g=1_C \text{ satisfies} (*)\}$ and $\a:=\{A\times B:A,B\in\b(\r^d)\}$. What we have proved shows $\pi$ system $\a\subset\g$. Apply monotone convergence theorem, we will easily find $\g$ is a $\lambda$ system. So $\g=\sigma(\a)=\b(\r^{2d})$. Use the standard procedure, we will find the assertion is correct for all $g\in\b(\r^{2d})$.
   		\item For any $A\in\b(\r^d)$, apply Fubini's theorem and the previous two small questions to obtain
   		\begin{equation*}
   			\begin{split}
   				\mu^\sharp(A)&=\iint 1_A(x+y)\mu(dx)\widetilde{\mu}(dy)\\
   				&=\int \mu(A-y)\widetilde{\mu}(dy)=\int \mu(A+y)\mu(dy)\\
   				&=\iint 1_A(x-y)\mu(dx)\mu(dy)\\
   				&=\iint 1_A(y-x)\mu(dx)\mu(dy)\\
   				&=\iint \mathbbm{1}_{-A}(x-y)\mu(dx)\mu(dy)\\
   				&=\mu^\sharp(-A).
   			\end{split}
   		\end{equation*}
   	\end{enumerate}
   \end{proof}
   
   \begin{question}
   	Given a distribution $\mu$ on measurable space $([0,\infty),\b[0,\infty))$. Write its Laplace transformation as 
   	\begin{equation*}
   		L_\mu(u)=\int_{[0,\infty)}e^{-ux}\mu(dx),\quad u\geq 0.
   	\end{equation*}
   	For two distributions $\mu_1$ and $\mu_2$ on $([0,\infty),\b[0,\infty))$, if $L_{\mu_1}(u)=L_{\mu_2}(u)$, then $\mu_1=\mu_2$. Namely, distributions and its Laplace transformations are mutually uniquely determined.
   \end{question}
   \begin{proof}
   	Since $\mu_1,\mu_2$ are only determined on $\b[0,\infty)$, then $|e^{-ux}|\leq 1,u\geq 0$ and thus $L_\mu(u)$ is well-defined. Moreover, write the region $D:=\{w\in\c:\text{Re}\, w< 0\}$, then $\overline{D}=\{w\in\c:\text{Re}\, w\leq 0\}$. We find 
   	\begin{equation*}
   		|e^{zx}|=e^{(\mathrm{Re}\,z)x}\leq 1,\quad \forall z\in D.
   	\end{equation*}
   	 So
   	\begin{equation*}
   		f_j(z):=\int_{[0,\infty)}e^{zx}\mu_j(dx),\quad j=1,2
   	\end{equation*} 
   	is still well-defined. Besides, for any $n\geq 1$, there exists $M_n>0$ satisfies
   	\begin{equation*}
   		|x^ne^{zx}|=x^ne^{(\mathrm{Re}\,z)x}\leq M_n.
   	\end{equation*}
   	By the dominated convergence theorem, we know $f_j(z)\in H(D)\cap C(\overline{D})$. Let $z=-u<0$, we have
   	\begin{equation*}
   		f_1(-u)=L_{\mu_1}(u)=L_{\mu_2}(u)=f_2(-u),\quad \forall u\geq 0.
   	\end{equation*}
   	Hence $f_1(z)$ and $f_2(z)$ coincide on $(-\infty,0]$. Apply the unique determination theorem for analytic functions, we have $f_1(z)=f_2(z)$ on $D$. Note $f_j(z)\in C(\overline{D})$, then $f_1(z)=f_2(z)$ on $\overline{D}$. Let $z=it,t\in\r$, we have 
   	\begin{equation*}
   		\varphi_{\mu_1}(t)=f_1(it)=f_2(it)=\varphi_{\mu_2}(t),\quad t\in\r,
   	\end{equation*}
   	where $\varphi_{\mu_j}(t)$ is the characteristic function of $\mu_j$. So $\mu_1=\mu_2$.
   \end{proof}
   
   \begin{question}
   	Given two measures $\mu,\nu$ on probability space $(\Omega,\f,\p)$. Assume $\mu\ll\nu$. Then $\nu\ll\mu$ iff $d\mu/d\nu\neq 0,\nu$-$\mathrm{a.e.}$.
   \end{question}
   \begin{proof}
   	$\Rightarrow$. If $\nu\ll\mu$, then for any $A\in\f$, we have
   	\begin{equation*}
   		\mu(A)=\int_A\frac{d\mu}{d\nu}d\nu=\int_A\frac{d\mu}{d\nu}\frac{d\nu}{d\mu}d\mu.
   	\end{equation*}
   	This shows 
   	\begin{equation*}
   		\int_A\(1-\frac{d\mu}{d\nu}\frac{d\nu}{d\mu}\)d\mu\equiv 0,\quad \forall A\in\f\Longrightarrow \frac{d\mu}{d\nu}\frac{d\nu}{d\mu}\equiv 1,\quad \mu\mathrm{-a.e.}.
   	\end{equation*}
   	Thus $d\mu/d\nu\neq 0,\mu\text{-a.e.}$ and by the mutual absolute continuity between $\mu$ and $\nu$, we have $d\mu/d\nu\neq 0,\nu\text{-a.e.}$.
   	
   	$\Leftarrow$. If $\nu\not\ll\mu$, then exists $B\in\f$ such that $\mu(B)=0$ but $\nu(B)>0$. Moreover, it is easy to find $d\mu/d\nu\geq 0,\nu$-$\mathrm{a.e.}$. So $d\mu/d\nu\neq 0,\nu$-$\mathrm{a.e.}$ means $d\mu/d\nu> 0,\nu$-$\mathrm{a.e.}$. Thus
   	\begin{equation*}
   		0=\mu(B)=\int_B\frac{d\mu}{d\nu}d\nu>0.
   	\end{equation*}
   	This contradiction yields the desired result.
   \end{proof}
   
   \begin{question}
   	For finite measures $\{\mu_n\}_{n\geq 1},\mu$ on metric space $(E,\b(E))$. Assume $\mu_n\overset{\mathrm{w}}{\to}\mu$, then for any $\mu$-continuous set $B$, we have $\mu_n(B\cap \cdot)\overset{\mathrm{w}}{\to}\mu(B\cap \cdot)$. It is also equivalent to
   	\begin{equation*}
   		\int_B f(x)\mu_n(dx)\to \int_Bf(x)\mu(dx),\quad n\to\infty.
   	\end{equation*}
   \end{question}
   \begin{proof}
   	Since we have proved $\partial(A\cup B)\subset\partial A\cup\partial B$. Then for any $\mu$-continuous set $A$ and the given $\mu$-continuous set $B$, we have
   	\begin{equation*}
   		\mu[\partial (A\cup B)]\leq \mu(\partial A\cup\partial B)\leq \mu(\partial A)+\mu(\partial B)=0.
   	\end{equation*}
   	Thus $A\cup B$ is also a $\mu$-continuous set for any $\mu$-continuous set $A$. Therefore 
   	\begin{equation*}
   		\mu_n(B\cap A)\to \mu(B\cap A),\quad n\to\infty.
   	\end{equation*}
   	And the last assertion is obvious.
   \end{proof}
   
   \begin{question}
   	Given a metric measurable space $(E,\b(E))$ with metric $d$. If $E$ is Polish, then the diagonal $D=\{(x,x):x\in E\}\in \b^2(E)=\b(E^2)$.
   \end{question}
   \begin{proof}
   	Since $E$ is separable, then there exists a countable dense set $\{q_n\}_{n\geq 1}\subset E$. We claim $\{(q_n,q_n)\}_{n\geq 1}\subset D$ is also a countable dense set of $D$. In fact, by the definition of the metric on the product metric space, for any $x\in E$ and $\varepsilon>0$, there exists $q_x\in \{q_n\}_{n\geq 1}$ satisfying $d(x,q_x)<\varepsilon/2$. Then
   	\begin{equation*}
   		d((x,x),(q_x,q_x))=2d(x,q_x)<\varepsilon.
   	\end{equation*}
   	And $\{(q_n,q_n)\}_{n\geq 1}$ is obviously countable. So $\{(q_n,q_n)\}_{n\geq 1}\subset D$ is a countable dense set of $D$. Define
   	\begin{equation*}
   		A_m:=\bigcup_{n\geq 1}B((q_n,q_n),1/m)\in \b^2(E)=\b(E^2),\quad m\geq 1.
   	\end{equation*}
   	We claim $D\subset A_m$ for any $m\geq 1$. In fact, for any $(x,x)\in D$, since $\{(q_n,q_n)\}_{n\geq 1}\subset D$ is a countable dense set of $D$, then exists $(q_x,q_x)\in \{(q_n,q_n)\}_{n\geq 1}$ such that $d((x,x),(q_x,q_x))<1/m$. Namely, 
   	\begin{equation*}
   		(x,x)\in B((q_x,q_x),1/m)\subset \bigcup_{n\geq 1}B((q_n,q_n),1/m)=A_m.
   	\end{equation*}
   	Thus $D\subset A_m$ for any $m\geq 1$ and therefore $D\subset \cap_{m\geq 1}A_m$. Next we show $\cap_{m\geq 1}A_m\subset D$. In fact, Given any 
   	\begin{equation*}
   		(x,y)\in \bigcap_{m\geq 1}A_m=\bigcap_{m\geq 1}\bigcup_{n\geq 1}B((q_n,q_n),1/m),
   	\end{equation*}
   	then for any $m\geq 1$, there exists $n_m\geq 1$, such that $(x,y)\in B((q_{n_m},q_{n_m}),1/m)$. So $d(x,y)\leq 2/m$ for any $m\geq 1$. Letting $m\to\infty$ yields $d(x,y)=0\Leftrightarrow x=y$. That is to say, $(x,y)\in D$. In conclusion, $D=\cap_{m\geq 1}A_m\in \b^2(E)=\b(E^2)$.
   \end{proof}
  
  \begin{question}
  	Given a real valued random variable $X$ on probability space $(\Omega,\f,\p)$. Suppose the $\sigma$-algebra $\g\subset \f$. Then the following claims are all equivalent (all the equalities below are in the sense of $\p$-a.s).
  	\begin{enumerate}
  		\item $X$ is independent to $\g$.
  		\item $\p[f(X)|\g]=\p[f(X)]$ for any $f\in\mathrm{b}\b(\r)$.
  		\item $\p(e^{i\theta X}|\g)=\p(e^{i\theta X})$.
  		\item $\p(e^{i\theta X}1_A)=\p(e^{i\theta X})\p(A)$ for any $A\in\g$.
  	\end{enumerate}
  \end{question}
  \begin{proof}
  	\begin{enumerate}
  		\item $1\Leftrightarrow 2$. For any $f\in\mathrm{b}\b(\r)$ and $A\in\g$, note $X$ is independent to $\g$, we have 
  		\begin{equation*}
  			\p[f(X)1_A]=\p[f(X)]\p(A)=\p\{\p[f(X)]1_A\}.
  		\end{equation*}
  		This shows $\p[f(X)|\g]=\p[f(X)]$. On the contrary, since $\p[f(X)|\g]=\p[f(X)]$, then 
  		\begin{equation*}
  			\begin{split}
  				\p[f(X)1_A]&=\p\{\p[f(X)1_A|\g]\}=\p\{1_A\p[f(X)|\g]\}\\
  				&=\p\{1_A\p[f(X)]\}=\p[f(X)]\p(A).
  			\end{split}
  		\end{equation*}
  		So $X$ is independent to $\g$.
  		\item $3\Leftrightarrow 4$. The proof is essentially the same as $1\Leftrightarrow 2$.
  		\item $1(2)\Leftrightarrow 3(4)$. First $2\Longrightarrow 3$ is obvious. On the contrary, we show $3\Rightarrow 1$. Note $\r$ is a polish space, then the regular conditional distribution of $X$ exists. Namely, 
  		\begin{equation*}
  			\mu_X^\g(dx):=\p(X\in dx|\g)
  		\end{equation*}
  		is a measure. And we have
  		\begin{equation*}
  			\p(e^{i\theta X}|\g)=\int e^{i\theta x} \mu_X^\g(dx),.
  		\end{equation*}
  		Since we have $\p(e^{i\theta X}|\g)=\p(e^{i\theta X})$, then
  		\begin{equation*}
  			\int e^{i\theta x} \mu_X^\g(dx)=\p(e^{i\theta X}|\g)=\p(e^{i\theta X})=\int e^{i\theta x} \mu_X(dx).
  		\end{equation*}
  		Since the characteristic function and the distribution are mutually uniquely determined, then we have $\mu_X^\g(dx)=\mu_X(dx)$. Namely 
  		\begin{equation*}
  			\p(X\in A|\g)=\p(X\in A),\quad \forall A\in\b(\r),
  		\end{equation*}
  		which implies $X$ is independent to $\g$.
  	\end{enumerate}
  \end{proof}
  
  \begin{question}
  	Given a random variable $X$ on probability space $(\Omega ,\f,\p)$ and valued on $(\bar\r_+,\b(\bar \r_+))$, where $\bar\r_+=\r_+\cup\{+\infty\}$ Consider its Laplace trasformation $\hat f(r)=\p (e^{-rX}),r>0$. Then we have 
  	\begin{equation*}
  		\lim_{r\downarrow 0}\hat f(r)=\p(X<+\infty).
  	\end{equation*}
  \end{question}
  \begin{proof}
  	Note 
  	\begin{equation*}
  		\begin{split}
  			e^{-rX}&=e^{-rX}\mathbbm{1}_{[X<+\infty]}+e^{-r\cdot \infty}\mathbbm{1}_{[X=+\infty]}\\
  			&=e^{-rX}\mathbbm{1}_{[X<+\infty]}\uparrow \mathbbm{1}_{[X<+\infty]},\quad r\downarrow 0.
  		\end{split}
  	\end{equation*}
  	Using the dominant convergence theorem to obtain
  	\begin{equation*}
  		\lim_{r\downarrow 0}\hat f(r)=\lim_{r\downarrow 0}\p (e^{-rX})=\p(X<+\infty).
  	\end{equation*}
  \end{proof}
  
  \begin{question}
  	Given random variables $(\xi_n)_{n\geq 1},(\eta_n)_{n\geq 1}$ and $\zeta$ on probability space $(\Omega,\f,\p)$. Assume $\xi_n-\eta_n\overset{\p}{\to} 0$ and $\xi_n\overset{\mathrm{w}}{\to}\zeta$, then $\eta_n\overset{\mathrm{w}}{\to}\zeta$.
  \end{question}
  \begin{proof}
  	For any $x\in C(F_\zeta)$, note $C^c(F_\zeta)$ is at most countable, then $C(F_\zeta)$ is dense in $\r$. Thus we can choose $\varepsilon>0$ such that $x+\varepsilon\in C(F_\zeta)$. Therefore for ant $k\geq 1$,
  	\begin{equation*}
  		\begin{split}
  			\p(\eta_n\leq x)&=\p(\eta_n\leq x,|\xi_n-\eta_n|<\varepsilon)+\p(\eta_n\leq x,|\xi_n-\eta_n|\geq \varepsilon)\\
  			&\leq \p(\xi_n-\varepsilon\leq x,|\xi_n-\eta_n|<\varepsilon)+\p(\eta_n\leq x,|\xi_n-\eta_n|\geq \varepsilon)\\
  			&=\p(\xi_n\leq x+\varepsilon)-\p(\xi_n-\varepsilon\leq x,|\xi_n-\eta_n|\geq \varepsilon)\\
  			&\quad+\p(\eta_n\leq x,|\xi_n-\eta_n|\geq \varepsilon).
  		\end{split}
  	\end{equation*}
  	Let $n\to\infty$ and note $\p(|\xi_n-\eta_n|\geq \varepsilon)\to0,n\to\infty$ for each $k\geq 1$, then we have
  	\begin{equation*}
  		\varlimsup_{n\to\infty}\p(\eta_n\leq x)\leq \p(\zeta\leq x+\varepsilon).
  	\end{equation*}
  	Let $\varepsilon\to0$, we have 
  	\begin{equation*}
  		\varlimsup_{n\to\infty}\p(\eta_n\leq x)\leq \p(\zeta\leq x).
  	\end{equation*}
  	On the other side, similarly, we can choose $\delta>0$ such that $x-\delta\in C(F_\zeta)$, so
  	\begin{equation*}
  		\begin{split}
  			\p(\eta_n\leq x)&=\p(\eta_n\leq x,|\xi_n-\eta_n|<\delta)+\p(\eta_n\leq x,|\xi_n-\eta_n|\geq \delta)\\
  			&\geq \p(\xi_n+\delta\leq x,|\xi_n-\eta_n|<\delta)+\p(\eta_n\leq x,|\xi_n-\eta_n|\geq \delta)\\
  			&=\p(\xi_n\leq x-\delta)-\p(\xi_n+\delta\leq x,|\xi_n-\eta_n|\geq \delta)\\
  			&\quad+\p(\eta_n\leq x,|\xi_n-\eta_n|\geq \delta).
  		\end{split}
  	\end{equation*}
  	Let $n\to\infty$ to obtain
  	\begin{equation*}
  		\varliminf_{n\to\infty}\p(\eta_n\leq x)\geq \p(\zeta\leq x-\delta).
  	\end{equation*}
  	Let $\delta\to0$ and note $x\in C(F_\zeta)$, we have
  	\begin{equation*}
  		\varliminf_{n\to\infty}\p(\eta_n\leq x)\geq \p(\zeta<x)=\p(\zeta\leq x).
  	\end{equation*}
  	The above argument yields
  	\begin{equation*}
  		F_{\eta_n}(x)=\p(\eta_n\leq x)\to \p(\zeta\leq x)=F_\zeta(x),\quad \forall x\in C(F_\zeta).
  	\end{equation*}
  \end{proof}
  
  \begin{question}
  	Given random variables $(X_n)_{n\geq 1},(Y_n)_{n\geq 1},X,Y$ on probability space $(\Omega,\f,\p)$. Assume $X_n\overset{\p}{\to} X$ and $Y_n\overset{\mathrm{w}}{\to}Y$, then the following four claims hold:
  	\begin{enumerate}
  		\item $X_n\overset{\p}{\to}X$ implies $X_n\overset{\mathrm{w}}{\to}X$.
  		\item If $X=C\in\r$ is a constant, then $X_n\overset{\p}{\to}C$ is equivalent to $X_n\overset{\mathrm{w}}{\to}C$.
  		\item If $X_n\overset{\p}{\to}C$ is a constant, then $X_n+Y_n\overset{\mathrm{w}}{\to}C+Y$.
  		\item If $X_n\overset{\p}{\to}C$ is a constant, then $X_nY_n\overset{\mathrm{w}}{\to}CY$.
  	\end{enumerate} 
  \end{question}
  \begin{proof}
  	\begin{enumerate}
  		\item Apply the above {\bfseries QUESTION} with $\xi_n=X,\eta_n=X_n$ and $\zeta=X$.
  		\item The first sub-question shows that $X_n\overset{\p}{\to}C$ implies $X_n\overset{\mathrm{w}}{\to}C$, we only need to show the converse implication. For any $\varepsilon>0$, we have
  		\begin{equation*}
  			\begin{split}
  				\p(|X_n-C|\geq \varepsilon)&=\p(X_n\geq C+\varepsilon)+\p(X_n\leq C-\varepsilon)\\
  				&=1-\p(X_n< C+\varepsilon)+\p(X_n\leq C-\varepsilon).
  			\end{split}
  		\end{equation*}
  		Note $F_C(x)$ has only one discontinuous point $x=C$, thus $C+\varepsilon,C-\varepsilon\in C(F_C)$. Let $n\to\infty$, we will get
  		\begin{equation*}
  			\begin{split}
  				\varlimsup_{n\to\infty}\p(|X_n-C|\geq \varepsilon)&=1-\lim_{n\to\infty}\p(X_n< C+\varepsilon)+\lim_{n\to\infty}\p(X_n\leq C-\varepsilon)\\
  				&=1-F_C(C+\varepsilon)+F_C(C-\varepsilon)=1-1+0=0.
  			\end{split}
  		\end{equation*}
  		\item Since we have $(X_n+Y_n)-(C+Y_n)=X_n-C\overset{\p}{\to}0$. Moreover we can get $C+Y_n\overset{\mathrm{w}}{\to}C+Y$ from $Y_n\overset{\mathrm{w}}{\to}Y$. In fact, it can easily be found that $C(F_{C+Y})=C(F_Y)+C:=\{y+C:y\in C(F_Y)\}$. Therefore, for any $x\in C(F_{C+Y})$, we have $x-C\in C(F_Y)$. And
  		\begin{equation*}
  			\begin{split}
  				\p(C+Y_n\leq x)&=\p(Y_n\leq x-C)\\
  				&\to \p(Y\leq X-C)=\p(C+Y\leq x),\quad\forall x\in C(F_{C+Y})
  			\end{split}
  		\end{equation*}
  		So we can apply the previous {\bfseries QUESTION} with $\xi_n=C+Y_n,\eta_n=X_n+Y_n$ and $\zeta=C+Y$.
  		\item Since we have proved in the 2nd sub-question that, convergence in probability to a constant is equivalent to convergence in distribution, we only need to use $\xi_n\overset{\p}{\to} \zeta$ and $\xi_n\overset{\p}{\to}0$ to show $\xi_n\eta_n\overset{\p}{\to}0$. However, this is obviously true, since the function $f(x,y)=xy$ is continuous in $\r^2$ and $\p$ is a finite measure. We only need to prove this assertion for $C\neq 0$.
  		
  		Suppose $C>0$. For any $x\in C(F_{CY})$ and $x\geq 0$, note $C(F_Y)$ is dense in $\r$, we can choose $\varepsilon>0$ such that $C-\varepsilon>0$ and $\frac{x}{C-\varepsilon}\in C(F_Y)$. Then for any $k\geq 1$, 
  		\begin{equation*}
  			\begin{split}
  				&\p(X_nY_n\leq x)=\p(X_nY_n\leq x,|X_n-C|< \varepsilon)+\p(X_nY_n\leq x,|X_n-C|\geq  \varepsilon)\\
  				=& \p(X_nY_n\leq x,Y_n\geq 0,|X_n-C|< \varepsilon)+\p(X_nY_n\leq x,Y_n< 0,|X_n-C|< \varepsilon)\\
  				&+\p(X_nY_n\leq x,|X_n-C|\geq  \varepsilon)\\
  				\leq &\p[(C-\varepsilon)Y_n\leq x,Y_n\geq 0,|X_n-C|< \varepsilon]+\p(Y_n<0,|X_n-C|< \varepsilon)+o(1)\\
  				=&\p\(0\leq Y_n\leq \frac{x}{C-\varepsilon}\)+\p(Y_n<0)+o(1)\\
  				=&\p\(Y_n\leq \frac{x}{C-\varepsilon} \)+o(1).
  			\end{split}
  		\end{equation*}
  		Let $n\to\infty$, we have
  		\begin{equation*}
  			\varlimsup_{n\to\infty}\p(X_ny_n\leq x)\leq \p\(Y\leq \frac{x}{C-\varepsilon}\).
  		\end{equation*}
  		Let $\varepsilon\to0$ and we will obtain
  		\begin{equation*}
  			\varlimsup_{n\to\infty}\p(X_ny_n\leq x)\leq \p(Y\leq x/C)=\p(CY\leq x).
  		\end{equation*}
  		If $x\in C(F_{CY})$ and $x< 0$, we can choose $\varepsilon>0$ such that $\frac{x}{C+\varepsilon}\in C(F_Y)$. Then 
  		\begin{equation*}
  			\begin{split}
  				\p(X_nY_n\leq x)&=\p(X_nY_n\leq x,Y_n< 0,|X_n-C|< \varepsilon)+o(1)\\
  				&\leq \p[(C+\varepsilon)Y_n\leq x,Y_n<0,|X_n-C|< \varepsilon]+o(1)\\
  				&=\p\(Y_n\leq \frac{x}{C+\varepsilon} \)+o(1).
  			\end{split}
  		\end{equation*}
  		Let $n\to\infty$, we have
  		\begin{equation*}
  			\varlimsup_{n\to\infty}\p(X_ny_n\leq x)\leq \p\(Y\leq \frac{x}{C+\varepsilon}\).
  		\end{equation*}
  		Let $\varepsilon\to 0$ and note $C(F_{CY})=C\cdot C(F_Y)=\{Cy:y\in C(F_Y)\}$. Therefore for all $x\in C(F_{CY})$, we have $x/C\in C(F_Y)$. Then we will again obtain
  		\begin{equation*}
  			\varlimsup_{n\to\infty}\p(X_nY_n\leq x)\leq \p(Y< x/C)=\p(Y\leq x/C)=\p(CY\leq x).
  		\end{equation*}
  		In short, we have $\varlimsup_{n\to\infty} F_{X_nY_n}(x)\leq F_{CY}(x)$ for any $x\in C(F_{CY})$.
  		
  		Similarly, if $C>0$ and $x\geq 0$, we can select $\delta>0$ such that $\frac{x}{C+\delta}\in C(F_{Y})$. So
  		\begin{equation*}
  			\begin{split}
  				\p(X_nY_n\leq x)&\geq \p[(C+\delta)Y_n\leq x,Y_n\geq 0]+\p(Y_n<0)+o(1)\\
  				&=\p\(Y_n\leq \frac{x}{C+\delta} \)+o(1).
  			\end{split}
  		\end{equation*}
  		Let $n\to\infty$ to find
  		\begin{equation*}
  			\varliminf_{n\to\infty}\p(X_nY_n\leq x)\geq \p\(Y\leq \frac{x}{C+\delta }\).
  		\end{equation*}
  		Let $\delta\to 0$ to obtain
  		\begin{equation*}
  			\varliminf_{n\to\infty}\p(X_nY_n\leq x)\geq \p(Y<x/C)=\p(Y\leq  x/C)=\p(CY\leq x).
  		\end{equation*}
  		If $x\in C(F_{CY})$ and $x<0$, we can select $\delta>0$ such that $\frac{x}{C-\delta}\in C(F_{Y})$, therefore
  		\begin{equation*}
  			\p(X_nY_n\leq x)\geq \p\(Y_n\leq \frac{x}{C-\delta} \)+o(1).
  		\end{equation*}
  		Let $n\to\infty$ and then let $\delta\to0$, we will get
  		\begin{equation*}
  			\varliminf_{n\to\infty}\p(X_nY_n\leq x)\geq \p(CY\leq x).
  		\end{equation*}
  		
  		In short, we have gotten $\varliminf_{n\to\infty} F_{X_nY_n}(x)\geq F_{CY}(x)$ for any $x\in C(F_{CY})$.  
  		
  		Hence when $C>0$, we have shown $F_{X_nY_n}(x)\to F_{CY}(x)$ for any $x\in C(F_{CY})$.
  		
  		When $C<0$, we have $-X_n\overset{\p}{\to}-C>0$. So the above discuss yields $-X_nY_n\overset{\mathrm{w}}{\to}-CY$. However, using the characteristic function will give us
  		 \begin{equation*}
  		 	\mathbb{E}e^{i(-t)X_nY_n}=\mathbb{E}e^{it(-X_nY_n)}\to \mathbb{E}e^{it(-CY)}=\mathbb{E}e^{i(-t)CY},\quad t\in\r,
  		 \end{equation*}
  		 which shows $X_nY_n\overset{\mathrm{w}}{\to}CY$.
  	\end{enumerate}
  \end{proof}
  
  \begin{question}
  	Given a Poisson process $N=(N_t)_{t\geq 0}$ and i.i.d. random variables $(Y_i)_{i\geq 1}$ on the probability space $(\Omega,\f,\p)$. Then the compound Poisson process
  	\begin{equation*}
  		X_t:=\sum_{i=1}^{N_t}Y_i,\quad t\geq 0
  	\end{equation*}
  	has independent increment.
  \end{question}
  \begin{proof}
  	For any $0<s<t$, we have to show the following random variables are independent, 
  	\begin{equation*}
  		X_s:=\sum_{i=1}^{N_s}Y_i,\quad X_t-X_s=\sum_{i=N_s+1}^{N_t}Y_i,
  	\end{equation*}
  	where $\sum_{i=m}^n\cdot=0$ for any $m>n$. If $N_t=N_s$, we know $X_t-X_s=0$, and therefore the theorem is trivially true. If $N_t>N_s$, for any $f,g\in\b(\r)$, we have
  	\begin{equation*}
  		\begin{split}
  			\p[f(X_s)&g(X_t-X_s)]=\p\{\p[f(X_s)g(X_t-X_s)|N]\}\\
  			&=\p\left\{\p\[\left.f\(\sum_{i=1}^{k_s}Y_i\)g\(\sum_{i=k_s+1}^{k_t}Y_i\)\]\right|_{k_s=N_s,k_t=N_t}\right\}\\
  			&=\p\left\{\p\[\left.f\(\sum_{i=1}^{k_s}Y_i\)\]\p\[g\(\sum_{i=k_s+1}^{k_t}Y_i\)\]\right|_{k_s=N_s,k_t=N_t}\right\}\\
  			&=:\p[F(N_s)G(N_s,N_t)].
  		\end{split}
  	\end{equation*}
  	Observe $(Y_i)_{i\geq 1}$ is an i.i.d. random variable sequence, then 
  	\begin{equation*}
  		G(k_s,k_t)=\p\[g\(\sum_{i=k_s+1}^{k_t}Y_i\)\],\quad k_s\leq k_t
  	\end{equation*}
  	only depends on the deference $k_t-k_s\in\z_+$. Thus we can write $\widetilde{G}(k_t-k_s):=G(k_s,k_t)$. Finally, we can apply the independent increment property of $(N_t)_{t\geq 0}$ to obtain
  	\begin{equation*}
  		\begin{split}
  			\p[f(X_s)&g(X_t-X_s)]=\p[F(N_s)\widetilde{G}(N_t-N_s)]\\
  			&=\p[F(N_s)]\p[\widetilde{G}(N_t-N_s)]\\
  			&=\p\left\{\p\[\left.f\(\sum_{i=1}^{k_s}Y_i\)\]\right|_{k_s=N_s}\right\}\p\left\{\left.\p\[g\(\sum_{i=k_s+1}^{k_t}Y_i\)\]\right|_{k_s=N_s,k_t=N_t}\right\}\\
  			&=\p\{\p[f(X_s)|N]\}\p\{\p[g(X_t-X_s)]|N\}\\
  			&=\p[f(X_s)]\p[g(X_t-X_s)].
  		\end{split}
  	\end{equation*}
  \end{proof}
  
  \begin{question}
  	Given a probability space $(\Omega,\f,\p)$ and $A,B\in\f$. Then we have $|\p(AB)-\p(A)\p(B)|\leq 1/4$.
  \end{question}
  \begin{proof}
  	Note $\p(AB)-\p(A)\p(B)=\mathrm{Cov}(1_A,1_B)$. So
  	\begin{equation*}
  		\begin{split}
  			|\mathrm{Cov}(1_A,1_B)|&\leq \sqrt{\mathrm{Var}(1_A)\mathrm{Var}(1_B)}\\
  			&\leq \sqrt{\[\p(A)-\p^2(A)\]\[\p(B)-\p^2(B)\]}\\
  			&\leq \sqrt{(1/4)^2}=1/4,
  		\end{split}
  	\end{equation*}
  	where the last inequality comes from the quadratic function $f(x)=x-x^2\leq 1/4$ for any $x\in\r$.
  \end{proof}
  
  \begin{question}
  	Given a probability space $(\Omega,\f,\p)$. Suppose $A\in\f$ satisfies $\p(A)>0$. Define $\q(B):=\p(B|A)$. Prove for any $F\in\mathrm{b}\f$ and sub-$\sigma$ algebra $\mathscr{C}$,
  	\begin{equation*}
  		\q(F|\mathscr{C})=\frac{\p(F1_A|\mathscr{C})}{\p(A|\mathscr{C})},\quad \q\text{-}\ae.
  	\end{equation*}
  	\begin{proof}
  		First we show $\p(A|\mathscr{C})\neq 0,\q$-$\ae$. Define $\Lambda:=\{\p(A|\mathscr{C})=0\}\in \mathscr{C}$. We have
  		\begin{equation*}
  			0=\p[1_\Lambda \p(A|\mathscr{C})]=\p(1_\Lambda 1_A )=\p(\Lambda\cap A).
  		\end{equation*}
  		Thus $\q(\Lambda)=\p(\Lambda|A)=\p(\Lambda\cap A)/\p(A)=0$. This implies $\p(A|\mathscr{C})\neq 0,\q$-$\ae$.
  		
  		Applying the monotonic class theorem, we can show 
  		\begin{equation*}
  			\int Fd\q=\frac{1}{\p(A)}\int F1_Ad\p,\quad F\in \mathrm{b}\f.
  		\end{equation*}
  		For any $G\in\mathrm{b}\mathscr{C}$, we have
  		\begin{equation*}
  			\begin{split}
  				\int \frac{\p(F1_A|\mathscr{C})}{\p(A|\mathscr{C})} Gd\q&=\frac{1}{\p(A)}\int \frac{\p(F1_A|\mathscr{C})}{\p(A|\mathscr{C})} G1_Ad\p\\
  				&=\frac{1}{\p(A)}\int \frac{\p(F1_A|\mathscr{C})}{\p(A|\mathscr{C})} G\p(A|\mathscr{C})d\p\\
  				&=\frac{1}{\p(A)}\int G\p(F1_A|\mathscr{C})d\p\\
  				&=\frac{1}{\p(A)}\int GF1_Ad\p\\
  				&=\int GFd\q.
  			\end{split}
  		\end{equation*}
  		This implies the statement holds.
  	\end{proof}
  \end{question}
  
  \begin{question}
  	Given a probability space $(\Omega,\f,\p)$. Suppose $\mathscr{D}=\sigma(A_i,i\geq 1)$, where $A_i$ is a partition of $\Omega$ and $\mathscr{C}\vee \mathscr{D}:=\sigma(\mathscr{C}\cup \mathscr{D})$. For any $B\in\f$,
  	\begin{equation}
  		\p(B|\mathscr{C}\vee \mathscr{D})=\sum_{n\geq 1}\frac{\p(A_n\cap B|\mathscr{C})}{\p(A_n|\mathscr{C})}\mathbbm{1}_{A_n}.\tag{$*$}
  	\end{equation}
  \end{question}
  \begin{proof}
  	If there exists $m\geq 1$ such that $\p(A_m)=0$, then $\p[\p(A_m|\mathscr{C})]=\p(A_m)=0$ and $\p(A_m|\mathscr{C})\geq 0$ yield $\p(A_m|\mathscr{C})=0,\p$-$\ae$. So 
  	\begin{equation*}
  		\frac{\p(A_m\cap B|\mathscr{C})}{\p(A_m|\mathscr{C})}=\infty.
  	\end{equation*}
  	Noticing $\p(A_m)=0$, we can define
  	\begin{equation*}
  		\frac{\p(A_m\cap B|\mathscr{C})}{\p(A_m|\mathscr{C})}\mathbbm{1}_{A_m}=c\mathbbm{1}_{A_m},\quad \forall c\in \r\cup\{\infty\}.
  	\end{equation*}
  	In short, the right side of $(*)$ is well-defined.
  	
  	 For any $C\in \mathscr{C}$ and $D\in \mathscr{D}$, there exist $\{A_{n_k}\}_{k=1}^m\subset \{A_i\}_{i\geq 1}$ such that $D=\sum_{k=1}^m A_{n_k}$. Suppose without loss of generality $\{A_{n_k}\}_{k=1}^{\ell}$ are $\p$-null sets, where $1\leq\ell\leq m$. Applying the result to the previous question, we have 
  	\begin{equation*}
  		\begin{split}
  			\int \sum_{n\geq 1}\frac{\p(A_n\cap B|\mathscr{C})}{\p(A_n|\mathscr{C})}&\mathbbm{1}_{A_n} \mathbbm{1}_{C\cap D}d\p = \int \sum_{k=1}^m\frac{\p(A_{n_k}\cap B|\mathscr{C})}{\p(A_{n_k}|\mathscr{C})}\mathbbm{1}_{A_{n_k}} \mathbbm{1}_{C}d\p\\
  			&=\sum_{k=\ell+1}^m\int \frac{\p(A_{n_k}\cap B|\mathscr{C})}{\p(A_{n_k}|\mathscr{C})}\mathbbm{1}_{A_{n_k}} \mathbbm{1}_{C}d\p\\
  			&=\sum_{k=\ell+1}^m\int \frac{\p(A_{n_k}\cap B|\mathscr{C})}{\p(A_{n_k}|\mathscr{C})}\p(A_{n_k}|\mathscr{C}) \mathbbm{1}_{C}d\p\\
  			&=\sum_{k=\ell+1}^m\int \p(A_{n_k}\cap B|\mathscr{C}) \mathbbm{1}_{C}d\p\\
  			&=\sum_{k=\ell+1}^m\int \mathbbm{1}_{A_{n_k}\cap B} \mathbbm{1}_{C}d\p\\
  			&=\int \sum_{k= 1}^m\mathbbm{1}_{A_{n_k} }1_B \mathbbm{1}_{C}d\p\\
  			&=\int\sum_{n\geq 1} \mathbbm{1}_{A_n}1_B \mathbbm{1}_{C}1_Dd\p\\
  			&=\int1_B \mathbbm{1}_{C\cap D}d\p.
  		\end{split}
  	\end{equation*}
  	This implies the statement holds.
  \end{proof}
  
  \begin{question}
  	Suppose $X$ is a random variable on the probability space $(\Omega,\f,\p)$ and $\g,\a\subset\f$ are two sub $\sigma$-algebras. Then $\p(X|\g\vee \a)=\p(X|\g)$ iff $X$ and $\a$ are conditional independent with respect to $\g$.
  \end{question}
  \begin{proof}
  	For any $G\in\g,A\in\a$, we have
  	\begin{equation*}
  		\begin{split}
  			\p(1_G1_AX)&=\p[1_G\p(1_AX|\g)]=\p[1_G\p(A|\g)\p(X|\g)]\\
  			&=\p[1_G1_A\p(X|\g)].
  		\end{split}
  	\end{equation*}
  	where the second equality holds because $X$ and $\a$ are conditional independent with respect to $\g$. The monotone class theorem yields $\p(1_BX)=\p[1_B\p(X|\g)]$, where $B\in \g\vee\a$. So $\p(X|\g)=\p(X|\g\vee\a)$. On the contrary, if $\p(X|\g\vee \a)=\p(X|\g)$, then for any $G\in\g,A\in\a$,
  	\begin{equation}
  		\p(1_G1_AX)=\p[1_G1_A\p(X|\g)]=\p[1_G\p(A|\g)\p(X|\g)]. \tag{$*$}
  	\end{equation}
  	So $(*)$ shows $\p(A|\g)\p(X|\g)=\p(1_AX|\g)$, which means $X$ and $\a$ are conditional independent with respect to $\g$.
  \end{proof}
  
  \begin{question}
  	There is a random variable $X$ on $(\Omega,\f,\p)$ such that $\p( |X|^\alpha)=\infty$ for any $\alpha>0$. 
  \end{question}
  \begin{proof}
  	Consider a continuous random variable $X$ with density 
  	\begin{equation*}
  		f(x)=
  		\begin{cases}
  			\frac{1}{2|x|\log^2 |x|},&\mbox{if }|x|>e,\\
  			0,&\mbox{otherewise}.
  		\end{cases}
  	\end{equation*}
  	Because
  	\begin{equation*}
  		\begin{split}
  			\int_{-\infty}^\infty f(x)dx&=\int_{[|x|>e]}\frac{dx}{2|x|\log^2 |x|}=\int_e^\infty \frac{dx}{x\log^2 x}\\
  			&=\int_e^\infty \frac{d\log x}{\log^2 x}=\left.\frac{1}{\log x}\right|_{\infty}^e=1,
  		\end{split}
  	\end{equation*}
  	$f(x)$ is indeed a density. And for any $\alpha>0$, we have $\frac{x^\alpha}{\log^2 x}\to \infty,x\to\infty$. So there exists $C>0$ such that for any $\frac{x^\alpha}{\log^2 x}\geq 1$ for any $x>C$. Whence
  	\begin{equation*}
  		\begin{split}
  			\p (|X|^\alpha)&= \int_{[|x|>e]}\frac{|x|^\alpha}{2|x|\log^2 |x|}dx= \int_e^\infty \frac{x^{\alpha}}{x\log^2 x}dx\\
  			&\geq \int_e^\infty x^{-1}dx=\infty. 
  		\end{split}
  	\end{equation*}
  \end{proof}
  
  
  \begin{question}
  	There is two random variables $X,Y$ on $(\Omega,\f,\p)$ such that $X,Y$ are related but there exists some $f\in\mathrm{b}\b$ such that $\p[f(X)|Y]=\p[f(X)]$.
  \end{question}
  \begin{proof}
  	Consider $\p(Y=\pm 1)=1/2$ and $X=-Y$. Obviously, $X,Y$ are related. Let $f(x)=x^2$. We have $\p[f(X)|Y]=\p(X^2|Y)=1=\p(X^2)$.
  \end{proof}
  
  \begin{question}
  	Assume $X, Y, D$ are three random variables on the probability space $(\Omega,\f,\p)$ and $\p|Y|<\infty$. If $\p[f(X)|D]\neq \p[f(X)]$ for any non-constant $f\in\b$ satisfying $\p[|f(X)|]<\infty$, and $Y,D$ are independent, then $\p(Y|X,D)=\p(Y|X)$ implies $\p(Y|X)=\p Y$.
  \end{question}
  \begin{proof}
  	Noting $\p[\p(Y|X,D)|D]=\p(Y|D)$, combining the fact $\p(Y|X,D)=\p(Y|X)$ and $Y,D$ are independent, we have
  	\begin{equation*}
  		\p[\p(Y|X)|D]=\p[\p(Y|X,D)|D]=\p(Y|D)=\p Y.
  	\end{equation*}
  	Since $\p[f(X)|D]\neq \p[f(X)]$ for any non-constant $f\in\b$ satisfying $\p[|f(X)|]<\infty$, and $\p Y=\p[\p(Y|X)]$, we have $\p(Y|X)=\p Y$ is a constant. 
  \end{proof}
  
  \begin{question}
  	Assume $X, Y$ are two random variables on the probability space $(\Omega,\f,\p)$ and $A\in\b$. Prove
  	\begin{equation*}
  	\p[\p(Y|X)|X\in A]=\p(Y|X\in A).
  \end{equation*}
  \end{question}
  \begin{proof}
  	We just need to note
  	\begin{equation*}
  		\begin{split}
  			\p[\p(Y|X)|X\in A]&=\frac{1}{\p(X\in A)}\int \p(Y|X)\mathbbm{1}_{[X\in A]}d\p\\
  			&=\frac{1}{\p(X\in A)}\int Y\mathbbm{1}_{[X\in A]}d\p\\
  			&=\p(Y|X\in A).
  		\end{split}
  	\end{equation*}
  \end{proof}
  
  \begin{question}
  	Assume $X, Y, Z$ are three random variables on the probability space $(\Omega,\f,\p)$ and $A\in\b$. Prove
  	\begin{equation*}
  		\p[\p(X|Y,Z\in A)|Z\in A]=\p(X|Z\in A).
  	\end{equation*}
  \end{question}
  \begin{proof}
  	Since
  	\begin{equation*}
  		\p(X|Y,Z\in A)=\frac{\p(X\mathbbm{1}_{[Z\in A]}|Y)}{\p(Z\in A|Y)}.
  	\end{equation*}
  	The definition of conditional expectation yields
  	\begin{align*}
  		\p[\p(X|Y,Z\in A)|Z\in A]&=\frac{1}{\p(Z\in A)} \int  \p(X|Y,Z\in A)\mathbbm{1}_{[Z\in A]}d\p\\
  			&=\frac{1}{\p(Z\in A)} \int  \frac{\p(X\mathbbm{1}_{[Z\in A]}|Y)}{\p(Z\in A|Y)}\mathbbm{1}_{[Z\in A]}d\p\\
  			&=\frac{1}{\p(Z\in A)} \int  \frac{\p(X\mathbbm{1}_{[Z\in A]}|Y)}{\p(Z\in A|Y)}\p(Z\in A|Y)d\p\\
  			&=\frac{1}{\p(Z\in A)} \int  \p(X\mathbbm{1}_{[Z\in A]}|Y)d\p\\
  			&=\frac{1}{\p(Z\in A)} \p[  \p(X\mathbbm{1}_{[Z\in A]}|Y)]\\
  			&=\frac{1}{\p(Z\in A)} \p(X\mathbbm{1}_{[Z\in A]})\\
  			&=\p(X|Z\in A).
  	\end{align*}
  \end{proof}
  
  \begin{question}
  	Suppose $(\Omega,\f)$ is a measurable space. Assume $\g\subset \f$ is a sub $\sigma$ algerba and $A\in\f$. Prove
  	\begin{equation*}
  		\g\vee\{A\}:=\sigma(\g\cup\{A\})=\{(B_1\cap A)\cup (B_2\cap A^c):B_1,B_2\in\g\}.
  	\end{equation*}
  \end{question}
  \begin{proof}
  	Write $\a:=\{(B_1\cap A)\cup (B_2\cap A^c):B_1,B_2\in\g\}$. Obviously $\a\subset \g\vee\{A\}$. We only need to prove $\a$ is a $\sigma$ algebra. First, $\Omega\in\a$ is obvious. Next, for any $(B_1\cap A)\cup (B_2\cap A^c)\in \a$, we have
  	\begin{equation*}
  		\begin{split}
  			[(B_1\cap A)\cup (B_2\cap A^c)]^c&=(B_1^c\cup A^c)\cap (B_2^c\cup A)\\
  			&=[(B_1^c\cup A^c)\cap B_2^c]\cup [(B_1^c\cup A^c)\cap A]\\
  			&=[(B_1^c\cap B_2^c)\cup (A^c\cap B_2^c)]\cup (B_1^c\cap A).
  		\end{split}
  	\end{equation*} 
  	We claim $B_1^c\cap B_2^c\subset (A^c\cap B_2^c)\cup (B_1^c\cap A)$. In fact, since $A$ and $A^c$ is a partition of $\Omega$, for any $\omega\in B_1^c\cap B_2^c$, we have either $w\in A$ or $w\in A^c$. Without loss of generality, we assume $\omega\in A^c$. Because $\omega\in B_1^c\cap B_2^c$, we have $\omega \in A^c\cap B_1^c\cap B_2^c\subset A^c\cap B_2^c\subset (A^c\cap B_2^c)\cup (B_1^c\cap A)$. So, the claim is true. The formula above turns out to be
  	\begin{equation*}
  		[(B_1\cap A)\cup (B_2\cap A^c)]^c=(A^c\cap B_2^c)\cup (B_1^c\cap A)\in \a.
  	\end{equation*}
  	Finally, for any sequence $\{(B_{1,k}\cap A)\cup (B_{2,k}\cap A^c)\}_{k\geq 1}\subset\a$, we have
  	\begin{equation*}
  		\begin{split}
  			\bigcup_{k\geq 1}[(B_{1,k}\cap A)&\cup (B_{2,k}\cap A^c)]=\[\bigcup_{k\geq 1}(B_{1,k}\cap A)\]\cup \[\bigcup_{k\geq 1}(B_{2,k}\cap A^c)\]\\
  			&=\[\(\bigcup_{k\geq 1}B_{1,k}\)\cap A\]\cup \[\(\bigcup_{k\geq 1}B_{2,k}\)\cap A^c\]\in\a.
  		\end{split}
  	\end{equation*} 
  	In conclusion, $\a$ is a $\sigma$ algebra and therefore $\g\vee\{A\}=\a$.
  \end{proof}
  
  \begin{question}
  	Suppose $(\Omega,\f)$ is a measurable space. Assume $\g\subset \f$ is a sub $\sigma$ algerba and $A\in\f$. For any measurable function $H\in \g\vee\{A\}$, there exist $F_1,F_2\in \g$ such that $H=F_1\mathbbm{1}_A+F_2\mathbbm{1}_{A^c}$.
  \end{question}
  \begin{proof}
  	Define 
  	\begin{equation*}
  		\Lambda:=\{H\in \g\vee\{A\}:\text{there exist $F_1,F_2\in \g$ such that $H=F_1\mathbbm{1}_A+F_2\mathbbm{1}_{A^c}$}\}.
  	\end{equation*}
  	Based on the previous QUESTION, we know all the indicator functions of $\g\vee\{A\}$ are in $\Lambda$. For any $H,H'\in\Lambda$, there exist $F_1,F_1',F_2,F_2'\in\g$ such that
  	\begin{equation*}
  		H=F_1\mathbbm{1}_A+F_2\mathbbm{1}_{A^c},\quad H'=F_1'\mathbbm{1}_A+F_2'\mathbbm{1}_{A^c}.
  	\end{equation*}
  	Therefore, for any $\alpha,\beta\in\r$,
  	\begin{equation*}
  		\alpha H+\beta H'=(\alpha F_1+\beta F_1')\mathbbm{1}_A+(\alpha F_2+\beta F_2')\mathbbm{1}_{A^c}\in\Lambda.
  	\end{equation*}
  	If $H\in\g\vee\{A\}$ and there exists a sequence $\{H_n\}_{n\geq 1}\subset\Lambda$ satisfying $0\leq H_n=F_{1,n}\mathbbm{1}_A+F_{2,n}\mathbbm{1}_{A^c}\uparrow H$, defining $F_1:=\varlimsup F_{1,n},F_2:=\varlimsup F_{2,n}\in \g$, we have $H=F_{1}\mathbbm{1}_A+F_{2}\mathbbm{1}_{A^c}$. Actually, noticing $A$ and $A^c$ are a partition of $\Omega$, then the pointwise convergence implies $F_{1,n}\mathbbm{1}_A\uparrow F_1\mathbbm{1}_A$ and $F_{2,n}\mathbbm{1}_A\uparrow F_2\mathbbm{1}_A$. 
  	
  	Based on the {\itshape monotonic class theorem of functions}, we know $\Lambda=\g\vee\{A\}$.
  \end{proof}
  
  \begin{question}
  	Suppose $X$ and $Y$ are two random variables on the probability space $(\Omega,\f,\p)$, and $A\in\f$. We have
  	\begin{equation*}
  		\mathbb{E}[X|\sigma(Y) \vee \{A\}]=\mathbb{E}(X|Y,A)\mathbbm{1}_A+\mathbb{E}(X|Y,A^c)\mathbbm{1}_{A^c},
  	\end{equation*} 
  	where $\mathbb{E}(X|Y,A):=\frac{\mathbb{E}(X\mathbbm{1}_A|Y)}{\p(A|Y)}\mathbbm{1}_{[\p(A|Y)>0]}$.
  \end{question}
  \begin{proof}
  	We only need to prove $\mathbb{E}[X|\sigma(Y) \vee \{A\}]\mathbbm{1}_A=\mathbb{E}(X|Y,A)\mathbbm{1}_A$. Because $\mathbb{E}[X|\sigma(Y) \vee \{A\}]\in \sigma(Y) \vee \{A\}$ and $A\in \sigma(Y) \vee \{A\}$, we have $\mathbb{E}[X|\sigma(Y) \vee \{A\}]\mathbbm{1}_A\in \sigma(Y) \vee \{A\}$. Similarly, since $\mathbb{E}(X|Y,A)\in \sigma(Y)\subset \sigma(Y) \vee \{A\}$, we also have $\mathbb{E}(X|Y,A)\mathbbm{1}_A\in \sigma(Y) \vee \{A\}$. Next, we only need to show 
  	\begin{equation*}
  		\int \mathbb{E}[X|\sigma(Y) \vee \{A\}]\mathbbm{1}_AHd\p=\int \frac{\mathbb{E}(X\mathbbm{1}_A|Y)}{\p(A|Y)}\mathbbm{1}_{[\p(A|Y)>0]}\mathbbm{1}_AHd\p,
  	\end{equation*}
  	for any $H\in \sigma(Y) \vee \{A\}$. Based on the previous QUESTION, we can find $H_1,H_2\in \sigma(Y)$ such that $H=H_1\mathbbm{1}_A+H_2\mathbbm{1}_{A^c}$. Therefore
  	\begin{align*}
  		&\int \frac{\mathbb{E}(X\mathbbm{1}_A|Y)}{\p(A|Y)}\mathbbm{1}_{[\p(A|Y)>0]}\mathbbm{1}_AHd\p\\
  			=&\int \frac{\mathbb{E}(X\mathbbm{1}_A|Y)}{\p(A|Y)}\mathbbm{1}_{[\p(A|Y)>0]}\mathbbm{1}_A(H_1\mathbbm{1}_A+H_2\mathbbm{1}_{A^c})d\p\\
  			=&\int \frac{\mathbb{E}(X\mathbbm{1}_A|Y)}{\p(A|Y)}\mathbbm{1}_{[\p(A|Y)>0]}H_1\mathbbm{1}_Ad\p\\
  			=&\int \frac{\mathbb{E}(X\mathbbm{1}_A|Y)}{\p(A|Y)}\mathbbm{1}_{[\p(A|Y)>0]}H_1\p(A|Y)d\p\\
  			=&\int \mathbb{E}(X\mathbbm{1}_A|Y)\mathbbm{1}_{[\p(A|Y)>0]}H_1d\p\\
  			=&\int X\mathbbm{1}_A\mathbbm{1}_{[\p(A|Y)>0]}H_1d\p\\
  			=&\int X\mathbbm{1}_A\mathbbm{1}_{[\p(A|Y)>0]}Hd\p\\
  			=&\int \mathbb{E}[X|\sigma(Y) \vee \{A\}]\mathbbm{1}_A\mathbbm{1}_{[\p(A|Y)>0]}Hd\p.
  	\end{align*}
  	Next, we show $\{\p(A|Y)=0\}\subset A^c,\p$-a.s.. In fact,
  	\begin{equation*}
  		\begin{split}
  			\p(A\cap \{\p(A|Y)=0\})&=\mathbb{E}[\p(A\cap \{\p(A|Y)=0\}|Y)]\\
  			&=\mathbb{E}[\mathbbm{1}_{[\p(A|Y)=0]}\p(A|Y)]=0,
  		\end{split}
  	\end{equation*}
  	which means $\{\p(A|Y)=0\}\subset A^c,\p$-a.s.. So, $A\subset \{\p(A|Y)>0\},\p$-a.s. and therefore
  	\begin{equation*}
  		\begin{split}
  			&\int \frac{\mathbb{E}(X\mathbbm{1}_A|Y)}{\p(A|Y)}\mathbbm{1}_{[\p(A|Y)>0]}\mathbbm{1}_AHd\p\\
  			=&\int \mathbb{E}[X|\sigma(Y) \vee \{A\}]\mathbbm{1}_A\mathbbm{1}_{[\p(A|Y)>0]}Hd\p\\
  			=&\int \mathbb{E}[X|\sigma(Y) \vee \{A\}]\mathbbm{1}_AHd\p.
  		\end{split}
  	\end{equation*}
  \end{proof}
  
  \begin{question}
  	Suppose $X$ is a random variable on $(\Omega,\f,\p)$. Given sub $\sigma$-field $\g\subset\f$, and measurable set $A\in\f$. The following conditional law of iterated expectation 
  	\begin{equation*}
  		\p[\p(X|\g)|A]=\p(X|A)\tag{$*$}
  	\end{equation*}
  	holds if at least one of the following conditions holds:
  	\begin{enumerate}[(1)]
  		\item $X\in\g$.
  		\item $A\in\g$.
  		\item $X$ and $A$ are independent conditional on $\g$.
  	\end{enumerate}
  \end{question}
  \begin{proof}
  	\begin{enumerate}
  		\item If (1) holds, then $\p(X|\g)=X$. So, $(*)$ naturally holds.
  		\item If (2) holds, we have
  		\begin{equation*}
  			\begin{split}
  				\p[\p(X|\g)|A]&=\frac{\p[\p(X|\g)\mathbbm{1}_A]}{\p(A)}=\frac{\p[\p(X\mathbbm{1}_A|\g)]}{\p(A)}\\
  				&=\frac{\p(X\mathbbm{1}_A)}{\p(A)}=\p(X|A).
  			\end{split}
  		\end{equation*}
  		\item If (3) holds, we have
  		\begin{equation*}
  			\begin{split}
  				\p[\p(X|\g)|A]&=\frac{\p[\p(X|\g)\mathbbm{1}_A]}{\p(A)}=\frac{\p\{\p[\p(X|\g)\mathbbm{1}_A|\g]\}}{\p(A)}\\
  				&=\frac{\p[\p(X|\g)\p(A|\g)]}{\p(A)}=\frac{\p[\p(X\mathbbm{1}_A |\g)]}{\p(A)}\\
  				&=\frac{\p(X\mathbbm{1}_A)}{\p(A)}=\p(X|A).
  			\end{split}
  		\end{equation*} 
  	\end{enumerate}
  \end{proof}
  
  \begin{question}
  	Given random variables $X,Y$ on the probability space $(\Omega,\f,\p)$, a measurable function $f\in\f$ and set $A\in\f$. Suppose $A$ and $X$ are independent conditional on $Y$. Then, $A$ and $X$ are independent conditional on $f(Y)$ if $\p(A|Y)\in \sigma(f(Y))$.
  \end{question}
  \begin{proof}
  	$A$ and $X$ are independent conditional on $Y$ is equivalent to
  	\begin{equation*}
  		\p(A|X,Y)=\p(A|Y).
  	\end{equation*}
  	Since $\p(A|Y)\in \sigma(f(Y))$ and $\sigma(f(Y))\subset \sigma(Y)$, we have
  	\begin{equation*}
  		\begin{split}
  			\p[A|X,f(Y)]&=\p[\p(A|X,Y)|X,f(Y)]\\
  			&=\p[\p(A|Y)|X,f(Y)]=\p(A|Y)\\
  			&=\p[\p(A|Y)|f(Y)]=\p[A|f(Y)].
  		\end{split}
  	\end{equation*}
  	Thus, $A$ and $X$ are independent conditional on $f(Y)$.
  \end{proof}
  
  \begin{question}
  	Suppose $(X_n)_{n\geq 1}$ is a r.v. sequence on the probability space $(\Omega,\f,\p)$ and valued in $(E,\e)$. Prove one part of the Pormanteau theorem: $\p f(X_n)\to \p f(X)$ for any $f\in C_b(E)$ iff $\varliminf_{n\to\infty}\p f(X_n)\geq \p f(X)$ for any $f\in C_+(E)$, where $C_+(E)$ is the set of all the non-negative and continuous functions.
  \end{question}
  \begin{proof}
  	\begin{enumerate}[(1)]
  		\item If $\varliminf_{n\to\infty}\p f(X_n)\geq \p f(X)$ for any $f\in C_+(E)$, for any $f\in C_b(E)$, there exists $m,M\in\r$ such that $m\leq f\leq M$. So, $M-f\geq 0,f-m\geq 0$ and therefore
  		\begin{equation*}
  			\begin{split}
  				&\varliminf_{n\to\infty}\p [M-f(X_n)]\geq \p [M-f(X)]\\
  				\Longleftrightarrow\quad & M-\varlimsup_{n\to\infty} \p f(X_n)\geq M-\p f(X)\\
  				\Longleftrightarrow\quad & \varlimsup_{n\to\infty} \p f(X_n)\leq \p f(X).
  			\end{split} 
  		\end{equation*}
  		Similarly, 
  		\begin{equation*}
  			\begin{split}
  				&\varliminf_{n\to\infty}\p [f(X_n)-m]\geq \p [f(X)-m]\\
  				\Longleftrightarrow\quad & \varliminf_{n\to\infty} \p f(X_n)\geq \p f(X).
  			\end{split}
  		\end{equation*}
  		Combine the above two equalities to obtain
  		\begin{equation*}
  			Pf(X)\leq \varliminf_{n\to\infty} \p f(X_n)\leq \varlimsup_{n\to\infty} \p f(X_n)\leq \p f(X),
  		\end{equation*}
  		which means $\p f(X_n)\to \p f(X)$.
  		\item If $\p f(X_n)\to \p f(X)$ for any $f\in C_b(E)$, for any  $f\in C_+(E)$, we know $f\wedge M\in C_b(E)$, where $M>0$. So,
  		\begin{equation*}
  			\lim_{n\to\infty}\p [f(X_n)\wedge M]=\p[f(X)\wedge M].
  		\end{equation*}
  		Observe $f\geq f\wedge M$, therefore
  		\begin{equation*}
  			\varliminf_{n\to\infty}\p f(X_n)\geq \lim_{n\to\infty}\p [f(X_n)\wedge M]=\p[f(X)\wedge M].
  		\end{equation*}
  		Let $M\to\infty$ and apply the monotone convergence theorem to obtain
  		\begin{equation*}
  			\varliminf_{n\to\infty}\p f(X_n)\geq \p f(X).
  		\end{equation*}
  	\end{enumerate}
  \end{proof}
  
  \begin{question}
  	Suppose $(X_n)_{n\geq 1}$ is a non-negative r.v. sequence on the probability space $(\Omega,\f,\p)$. If $X_n\overset{\mathrm{d}}{\to}X$, then $X$ is also non-negative.
  \end{question}
  \begin{proof}
  	Since $X_n$ is non-negative, we have $\p(X_n<0)=0$. By the Portmanteau theorem, we know 
  	\begin{equation*}
  		0=\varliminf_{n\to\infty}\p(X_n<0)\geq \p(X<0).
  	\end{equation*}
  	This implies $\p(X<0)=0$.
  \end{proof}
  
  \begin{question}
  	Suppose $(X_n)_{n\geq 1}$ and $X$ are r.v.s in the probability space $(\Omega,\f,\p)$. If $X_n\overset{L^r}{\to}X,n\to\infty$, then $X^r_n\overset{L^1}{\to}X^r,n\to\infty$.
  \end{question}
	\begin{proof}
		Since $X_n\overset{L^r}{\to}X,n\to\infty$, we know that $X_n\overset{\p}{\to}X,n\to\infty$ and
		\begin{equation*}
			\int_A |X_n|^rd\p\text{ is integral absolutely continuous.}
		\end{equation*}
		Since $x\mapsto |x|^r$ is a continuous function, $|X_n|^r\overset{\p}{\to}|X|^r,n\to\infty$. Moreover,
		\begin{equation*}
			\int_A |X_n|^rd\p=\int_A |X_n|^rd\p\text{ is integral absolutely continuous.}
		\end{equation*}
		So, $X^r_n\overset{L^1}{\to}X^r,n\to\infty$.
	\end{proof}
	
	\begin{question}
		If $(X_n)_{n\geq 1}$ and $(Y_n)_{n\geq 1}$ are two tight r.v. sequences in the probability space $(\Omega,\f,\p)$ and valued in $(E,\b(E))$, then $\{(X_n,Y_n)\}_{n\geq 1}$ is also a tight r.v. sequences in the product probability space $(\Omega^2,\f^2,\p^2)$.
	\end{question}
	\begin{proof}
		For any $\varepsilon>0$, there exist compact sets $M_\varepsilon,V_\varepsilon\subset E$ such that
		\begin{equation*}
			\sup_{n\geq 1}\p(X_n\in M_\varepsilon^c)<\varepsilon,\quad \sup_{n\geq 1}\p(X_n\in V_\varepsilon^c)<\varepsilon.
		\end{equation*}
		Therefore
		\begin{equation*}
			\begin{split}
				\sup_{n\geq 1}\p[(X_n,Y_n)&\in (M_\varepsilon\times V_\varepsilon) ^c]=\sup_{n\geq 1}\p(X_n\in M_\varepsilon^c \text{ or }Y_n\in V_\varepsilon^c)\\
				&\leq \sup_{n\geq 1}\p(X_n\in M_\varepsilon^c)+\sup_{n\geq 1}\p(X_n\in V_\varepsilon^c)<2\varepsilon,
			\end{split}
		\end{equation*}
		where $M_\varepsilon\times V_\varepsilon\subset E^2$ is also compact in the topology of product topology. So, $\{(X_n,Y_n)\}_{n\geq 1}$ is also a compact r.v. sequences in the product probability space $(\Omega^2,\f^2,\p^2)$.
	\end{proof}
	
	\begin{question}
		Suppose $\mu,\nu$ are two measures on the measurable space $(\Omega,\f)$ satisfying $\mu\ll \nu$, and $\lambda,\#$ are two measures on the measurable space $(E,\e)$ satisfying $\lambda\ll \#$. Then $\mu\times\lambda\ll \nu\times\#$ on the measurable space $(\Omega\times E,\f\times \e)$. 
	\end{question}
	\begin{proof}
		For any $A\in \f\times\e$ such that $(\nu\times\#)(A)=0$, we have
		\begin{equation*}
			\int_E\nu(A_x)\#(dx)=0\Longrightarrow \nu(A_x)=0,\quad \#\text{ - a.s.,}
		\end{equation*}
		where $A_x$ is the section of $A$ at $x$. Since $\mu\ll \nu$ and $\lambda\ll \#$, it follows from the above expression that 
		\begin{equation*}
			\mu(A_x)=0,\quad \lambda\text{ - a.s..}
		\end{equation*}
		Therefore
		\begin{equation*}
			(\mu\times\lambda)(A)=\int_E\mu(A_x)\lambda(dx)=0,
		\end{equation*}
		which implies $\mu\times\lambda\ll \nu\times\#$ on the measurable space $(\Omega\times E,\f\times \e)$. 
	\end{proof}
	
	\begin{question}
  	Suppose $X_n\overset{\p}{\to} X$. There exists a sequence $\{\varepsilon_n\}_{n\geq 1}$ satisfying $\varepsilon_n\downarrow 0,n\to\infty$, such that $\p(|X_n-X|>\varepsilon_n)\to0,n\to\infty$.
  \end{question}
  \begin{proof}
  	Since $X_n\overset{\p}{\to} X$, we have
  	\begin{equation*}
  		\p(|X_n-X|>\varepsilon)\to0,n\to\infty,\quad \text{for any }\varepsilon>0.
  	\end{equation*}
  	Given any decreasing sequence $\gamma_n\downarrow 0,n\to\infty$, we can construct another decreasing sequence $\{\varepsilon_n\}_{n\geq 1}$ as follows:
  	\begin{itemize}
  		\item Let $\varepsilon_1:=\gamma_1$ and $N_1:=1$.
  		\item For $\gamma_k$, where $k\geq 2$, since $\p(|X_n-X|>\gamma_k)\to0,n\to\infty$, there exists $\ell_k\in\z_+$ such that 
  		\begin{equation*}
  			\p(|X_n-X|>\gamma_k)\leq 1/k,\quad \forall n\geq \ell_k.
  		\end{equation*}
  		Define $N_k:=\max\{N_{k-1},\ell_k\}+1$ and $\varepsilon_n:=\gamma_{k-1},\forall N_{k-1}\leq n<N_k-1$.
  	\end{itemize}
  	
  	Based on this construction, the sequence $\{\varepsilon_n\}_{n\geq 1}$ satisfies $\varepsilon_n\to0,n\to\infty$. Moreover, for any $n\in\z_+$, there exists $k(n)\in\z_+$ such that $N_{k(n)}\leq n<N_{k(n)}-1$ and 
  	\begin{equation}\label{1}
  		\p(|X_n-X|>\varepsilon_n)\leq 1/k(n).
  	\end{equation}
  	Noting the integer sequence $\{N_k\}$ is strictly increasing to $\infty$, we know $k(n)\uparrow \infty,n\to\infty$. Therefore, letting $n\to\infty$ in \eqref{1}, we will get the desired result.
  	\end{proof}
  	
  	\begin{question}
  		Suppose $X_n$ and $Y_n$ are two r.v.s on the probability space $(\Omega,\f,\p)$. Then $X_n\leq Y_n+o_p(1)$ is equivalent to 
  		\begin{equation*}
  			\p(X_n>Y_n+\varepsilon)\to 0,\quad n\to\infty,\quad \forall \varepsilon>0. \tag{$*$}
  		\end{equation*}
  	\end{question}
  	\begin{proof}
  		Suppose $X_n\leq Y_n+o_p(1)$, there exists $Z_n=o_p(1)$, such that $X_n\leq Y_n+Z_n$. Then, for any $\varepsilon>0$,
  		\begin{equation*}
  			\p(X_n>Y_n+\varepsilon)\leq \p(Z_n>\varepsilon)\leq \p(|Z_n|>\varepsilon)\to0,\quad n\to\infty.
  		\end{equation*}
  		On the contrary, if ($*$) holds, then define $Z_n:=(X_n-Y_n)^+$, we have 
  		\begin{equation*}
  			X_n=Y_n+(X_n-Y_n)\leq Y_n+(X_n-Y_n)^+=Y_n+Z_n
  		\end{equation*}
  		and
  		\begin{equation*}
  			\begin{split}
  				\p(|Z_n|>\varepsilon)&=\p[(X_n-Y_n)^+>\varepsilon]=\p(X_n-Y_n>\varepsilon)\\
  				&=\p(X_n>Y_n+\varepsilon)\to0,\quad n\to\infty,\quad \forall \varepsilon>0.
  			\end{split}
  		\end{equation*}
  		So, $Z_n=o_p(1)$ and $X_n\leq Y_n+Z_n$.
  	\end{proof}
  	
  	\begin{question}
  		Suppose $X_n$ and $Y_n$ are two r.v.s on the probability space $(\Omega,\f,\p)$. If $\p(X_n\leq Y_n)\to 1,n\to\infty$, then $X_n\leq Y_n+o_p(1)$.
  	\end{question}
	\begin{proof}
		Define $Z_n=(X_n-Y_n)^+$, then
		\begin{equation*}
			\begin{split}
				\p(|Z_n|>\varepsilon)&=\p[(X_n-Y_n)^+>\varepsilon]=\p(X_n-Y_n>\varepsilon)\\
				&=\p(X_n>Y_n+\varepsilon)\leq \p(X_n>Y_n)\\
				&=1-\p(X_n\leq Y_n)\to0,\quad n\to\infty.
			\end{split}
		\end{equation*}
		Thus $Z_n=o_p(1)$ and obviously $X_n\leq Y_n+Z_n$.
	\end{proof}
	
	\begin{question}
		Suppose $X_n,Y_n$ and $Z_n$ are three positive r.v.s on the probability space $(\Omega,\f,\p)$. Suppose $X_n=O_p(Y_n)+O_p(Z_n)$, then $X_n=O_p(Y_n+Z_n)$.
	\end{question}
	\begin{proof}
		For any $\forall \varepsilon>0$, there exists $M_1^\varepsilon>0$ and $M_2^\varepsilon>0$ such that
		\begin{equation*}
			\sup_{n\geq 1}\p\(\frac{O_p(Y_n)}{Y_n}> M_1^\varepsilon \)<\varepsilon/2\quad \text{ and }\quad \sup_{n\geq 1}\p\(\frac{O_p(Z_n)}{Z_n}> M_2^\varepsilon \)<\varepsilon/2.
		\end{equation*}
		Define $M:=2(M_1^\varepsilon\vee M_2^\varepsilon)>0$, then
		\begin{equation*}
			\begin{split}
				\p\(\frac{X_n}{Y_n+Z_n} >M\)&=\p\(\frac{O_p(Y_n)+O_p(Z_n)}{Y_n+Z_n} >M\)\\
				&=\p\(\frac{O_p(Y_n)}{Y_n+Z_n}+\frac{O_p(Z_n)}{Y_n+Z_n} >M\)\\
				&\leq \p\(\frac{O_p(Y_n)}{Y_n}+\frac{O_p(Z_n)}{Z_n} >M\)\\
				&\leq \p\(\frac{O_p(Y_n)}{Y_n}>\frac{M}{2}\)+\p\(\frac{O_p(Z_n)}{Z_n}>\frac{M}{2}\)<\varepsilon.
			\end{split}
		\end{equation*}
		So, $X_n=O_p(Y_n+Z_n)$
	\end{proof}
	
	\begin{question}
		Suppose $X,Y$ and $T$ are three r.v.s on the probability space $(\Omega,\f,\p)$ and valued on the measurable space $(E,\e)$. If $Y\perp T|X$ and $T\perp X$, then $Y\perp X$.
	\end{question}
	\begin{proof}
		Since $Y\perp T|X$, for any $f,g\in \mathrm{b}\e$,
		\begin{equation*}
			\begin{split}
				\mathbb{E}[f(Y)g(T)|X]=\mathbb{E}[f(Y)|X]\mathbb{E}[g(T)|X]=\mathbb{E}[f(Y)|X]\mathbb{E}[g(T)].
			\end{split}
		\end{equation*}
		Taking the expectation on both sides, we will get
		\begin{equation*}
			\mathbb{E}[f(Y)g(T)]=\mathbb{E}[f(Y)]\mathbb{E}[g(T)],
		\end{equation*}
		which implies $Y\perp X$.
	\end{proof}
	
	% LaTeX snippet for Exercise 4.1.5. Inline math: $...$; display math: equation* only

\begin{question}
	Let $X$ and $Y$ be independent and identically distributed random variables with $\mathbb E|X|<\infty$. Show that $\mathbb E[|X+Y|]\ge \mathbb E[|X-Y|]$, where equality holds if and only if the distribution of $X$ is symmetric around $0$.
\end{question}
\begin{proof}
	The proof uses the Dirichlet integral identity
\begin{equation*}
\int_{-\infty}^{+\infty} \frac{1-\cos(a x)}{x^2}\,dx = \pi |a|.
\end{equation*}
Therefore, by Tonelli's theorem,
\begin{equation*}
\begin{split}
	&\mathbb E\big[|X+Y|-|X-Y|\big]\\
=\quad &\frac{1}{\pi}\,\mathbb E\[\int_{-\infty}^{+\infty} \frac{1-\cos\big((X+Y)t\big)}{t^2}\,dt\; -\; \int_{-\infty}^{+\infty} \frac{1-\cos\big((X-Y)t\big)}{t^2}\,dt\]\\
=\quad & \frac{1}{\pi}\int_{-\infty}^{+\infty} \mathbb E\bigg[\frac{\cos\big((X-Y)t\big)-\cos\big((X+Y)t\big)}{t^2}\bigg] dt.
\end{split}
\end{equation*}
Using the trigonometric identity $\cos u-\cos v=-2\sin\frac{u+v}{2}\,\sin\frac{u-v}{2}$ with $u=(X-Y)t$ and $v=(X+Y)t$, we obtain
\begin{equation*}
\cos\big((X-Y)t\big)-\cos\big((X+Y)t\big) = 2\sin(Xt)\sin(Yt).
\end{equation*}
Hence, by independence and identical distribution of $X$ and $Y$,
\begin{equation*}
\begin{split}
	\mathbb E\big[|X+Y|-|X-Y|\big]
&= \frac{2}{\pi}\int_{-\infty}^{+\infty} \mathbb E\left[\frac{\sin(Xt)\sin(Yt)}{t^2}\right] dt\\
&= \frac{2}{\pi}\int_{-\infty}^{+\infty} \frac{1}{t^2}\,\mathbb E[\sin(Xt)]\,\mathbb E[\sin(Yt)]\,dt\\
&= \frac{2}{\pi}\int_{-\infty}^{+\infty} \left(\frac{\mathbb E[\sin(Xt)]}{t}\right)^{\!2} dt\ \ge\ 0.
\end{split}
\end{equation*}
This proves $\mathbb E[|X+Y|] \ge \mathbb E[|X-Y|]$. If equality holds, then necessarily $\mathbb E[\sin(Xt)]=0$ for all $t\in\mathbb R$. Since the characteristic function is $\varphi_X(t)=\mathbb E[e^{itX}]=\mathbb E[\cos(tX)]+i\,\mathbb E[\sin(tX)]$, this means $\varphi_X$ is real for all $t$, hence the distribution of $X$ is symmetric about $0$. The converse is immediate.
\end{proof}

\begin{question}[Equivalences for stochastic boundedness]
Let $\{X_n\}_{n\ge1}$ be random variables taking values in $\mathbb{R}^d$ (for any fixed norm $\|\cdot\|$ on $\mathbb{R}^d$). The following statements are equivalent:

\begin{enumerate}
\item[(1)] For every $\varepsilon>0$ there exists $M<\infty$ such that
\begin{equation*}
\sup_{n\ge1}\,\mathbb{P}\big(\|X_n\|>M\big)\le \varepsilon.
\end{equation*}
Equivalently, $\lim_{M\to\infty}\sup_{n}\mathbb{P}(\|X_n\|>M)=0$.

\item[(2)] For every deterministic sequence $\{\ell_n\}_{n\ge1}$ with $\ell_n\to\infty$,
\begin{equation*}
\mathbb{P}\big(\|X_n\|>\ell_n\big)\ \longrightarrow\ 0.
\end{equation*}

\item[(3)] For every $\varepsilon>0$ there exists $M<\infty$ such that
\begin{equation*}
\limsup_{n\to\infty}\mathbb{P}\big(\|X_n\|>M\big)\le \varepsilon.
\end{equation*}
\end{enumerate}
In particular, these are all equivalent to saying that $\{X_n\}$ is bounded in probability (stochastically bounded, i.e.\ $X_n=O_{\mathbb{P}}(1)$).
\end{question}

\begin{proof}
\textbf{(1) $\Rightarrow$ (2).}
Fix $\varepsilon>0$ and a sequence $\ell_n\to\infty$. By (1), choose $M<\infty$ with
\begin{equation*}
\sup_{n}\mathbb{P}(\|X_n\|>M)\le\varepsilon.
\end{equation*}
For all large $n$, $\ell_n\ge M$, hence
\begin{equation*}
\mathbb{P}(\|X_n\|>\ell_n)\le \mathbb{P}(\|X_n\|>M)\le \varepsilon,
\end{equation*}
so $\mathbb{P}(\|X_n\|>\ell_n)\to0$.

\smallskip
\textbf{(2) $\Rightarrow$ (1).}
Argue by contradiction. Suppose (1) fails. Then there exists $\varepsilon_0>0$ such that
\begin{equation*}
\lim_{M\to\infty}\sup_{n}\mathbb{P}(\|X_n\|>M)\ \ge\ \varepsilon_0.
\end{equation*}
Define $g_N(M):=\sup_{n\ge N}\mathbb{P}(\|X_n\|>M)$ for $N\in\mathbb{N}$. If for some $N$ we had $\lim_{M\to\infty}g_N(M)=0$, then the tail family $\{X_n:n\ge N\}$ would be tight; combined with the finiteness of $\{X_1,\dots,X_{N-1}\}$ this would yield (1), a contradiction. Hence for every $N$,
\begin{equation*}
\limsup_{M\to\infty} g_N(M)\ \ge\ \varepsilon_0.
\end{equation*}
Let $\eta:=\varepsilon_0/2$. We now construct strictly increasing sequences $M_k\uparrow\infty$ and $n_k\uparrow\infty$ recursively. Set $N_1:=1$. Choose $M_1$ large with $g_{N_1}(M_1)\ge \eta$, and then pick $n_1\ge N_1$ such that $\mathbb{P}(\|X_{n_1}\|>M_1)\ge\eta$.
Given $(n_{k-1},M_{k-1})$, set $N_k:=n_{k-1}+1$. Choose $M_k>M_{k-1}$ with $g_{N_k}(M_k)\ge\eta$, and then pick $n_k\ge N_k$ so that $\mathbb{P}(\|X_{n_k}\|>M_k)\ge\eta$.

Define a deterministic, piecewise-constant sequence
\begin{equation*}
\ell_n:=\begin{cases}
M_1, & n<n_1,\\
M_k, & n_k\le n<n_{k+1}\quad (k\ge1).
\end{cases}
\end{equation*}
Then $\ell_n\to\infty$ and, for all $k$,
\begin{equation*}
\mathbb{P}(\|X_{n_k}\|>\ell_{n_k})
=\mathbb{P}(\|X_{n_k}\|>M_k)\ \ge\ \eta,
\end{equation*}
which contradicts (2). Thus (1) must hold.

\smallskip
\textbf{(1) $\Rightarrow$ (3).}
Immediate, since for any $M$,
\begin{equation*}
\limsup_{n\to\infty}\mathbb{P}(\|X_n\|>M)\ \le\ \sup_{n}\mathbb{P}(\|X_n\|>M).
\end{equation*}

\smallskip
\textbf{(3) $\Rightarrow$ (1).}
Fix $\varepsilon>0$. By (3), choose $M_\infty<\infty$ with
\begin{equation*}
\limsup_{n\to\infty}\mathbb{P}(\|X_n\|>M_\infty)\le \varepsilon/2,
\end{equation*}
so there exists $N$ such that
\begin{equation*}
\sup_{n\ge N}\mathbb{P}(\|X_n\|>M_\infty)\le \varepsilon/2.
\end{equation*}
For each fixed $n<N$, the sets $\{\|X_n\|>t\}$ decrease to the empty set as $t\to\infty$, hence $\mathbb{P}(\|X_n\|>t)\downarrow 0$. Thus we can choose $M_0<\infty$ with
\begin{equation*}
\max_{1\le n<N}\mathbb{P}(\|X_n\|>M_0)\le \varepsilon/2.
\end{equation*}
Let $M:=\max\{M_\infty,M_0\}$. Then
\begin{equation*}
\begin{split}
	\sup_{n}\mathbb{P}(\|X_n\|>M)
&=\max\Big\{\max_{n<N}\mathbb{P}(\|X_n\|>M),\ \sup_{n\ge N}\mathbb{P}(\|X_n\|>M)\Big\}\\
&\le \varepsilon/2+\varepsilon/2=\varepsilon,
\end{split}
\end{equation*}
which proves (1).

\smallskip
Combining the implications, we obtain the equivalence of (1), (2), and (3).
\end{proof}

\begin{question}[Equivalence of stochastic orders under comparable normalizations]
Let $\{X_n\}_{n\ge1}$ be random variables, and let $\{a_n\},\{b_n\},\{c_n\}$ be deterministic sequences of positive real numbers. 

\begin{enumerate}
\item[(i)] Suppose there exist constants $0<m\le M<\infty$ and $N\in\mathbb{N}$ such that for all $n\ge N$,
\begin{equation*}
m a_n \le c_n \le M a_n.
\end{equation*}
Show that
\begin{equation*}
X_n = O_{\mathbb{P}}(c_n b_n)
\quad \Longleftrightarrow \quad
X_n = O_{\mathbb{P}}(a_n b_n).
\end{equation*}

\item[(ii)] Suppose instead that
\begin{equation*}
a_n \le c_n \le a_n + 1 \quad \text{for all } n.
\end{equation*}
Show that the condition in (i) holds (and hence the equivalence above is valid) under either of the following additional assumptions:
\begin{enumerate}
\item[(a)] $a_n \to \infty$;
\item[(b)] $\inf_n a_n \ge \delta > 0$.
\end{enumerate}
\end{enumerate}
\end{question}

\begin{proof}
\textbf{(i)}
Recall that $X_n = O_{\mathbb{P}}(r_n)$ means: for every $\varepsilon>0$, there exist $K_\varepsilon<\infty$ and $N_\varepsilon$ such that for all $n\ge N_\varepsilon$,
\begin{equation*}
\mathbb{P}(|X_n| > K_\varepsilon r_n) < \varepsilon.
\end{equation*}

\smallskip
\emph{($\Rightarrow$)} Suppose $X_n = O_{\mathbb{P}}(c_n b_n)$. Then for every $\varepsilon>0$, there exist $K_\varepsilon$ and $N_1$ such that for all $n\ge N_1$,
\begin{equation*}
\mathbb{P}(|X_n| > K_\varepsilon c_n b_n) < \varepsilon.
\end{equation*}
Using $c_n \le M a_n$ for $n\ge N$, we obtain
\begin{equation*}
K_\varepsilon c_n b_n \le K_\varepsilon M a_n b_n.
\end{equation*}
Hence, for all $n\ge \max\{N,N_1\}$,
\begin{equation*}
\mathbb{P}(|X_n| > K_\varepsilon M a_n b_n)
\le 
\mathbb{P}(|X_n| > K_\varepsilon c_n b_n)
< \varepsilon,
\end{equation*}
which implies $X_n = O_{\mathbb{P}}(a_n b_n)$.

\smallskip
\emph{($\Leftarrow$)} Suppose $X_n = O_{\mathbb{P}}(a_n b_n)$. Then for every $\varepsilon>0$, there exist $L_\varepsilon$ and $N_2$ such that for all $n\ge N_2$,
\begin{equation*}
\mathbb{P}(|X_n| > L_\varepsilon a_n b_n) < \varepsilon.
\end{equation*}
Using $c_n \ge m a_n$ for $n\ge N$, we have
\begin{equation*}
a_n \le \frac{1}{m} c_n.
\end{equation*}
Thus,
\begin{equation*}
L_\varepsilon a_n b_n \le \frac{L_\varepsilon}{m} c_n b_n.
\end{equation*}
Therefore, for all $n\ge \max\{N,N_2\}$,
\begin{equation*}
\mathbb{P}\Big(|X_n| > \tfrac{L_\varepsilon}{m} c_n b_n\Big)
\le 
\mathbb{P}(|X_n| > L_\varepsilon a_n b_n)
< \varepsilon,
\end{equation*}
which yields $X_n = O_{\mathbb{P}}(c_n b_n)$.

\smallskip
This proves the equivalence.

\bigskip
\noindent
\textbf{(ii)}
\smallskip
\emph{(a) If $a_n \to \infty$.}

Then for all $n$,
\begin{equation*}
1 \le \frac{c_n}{a_n} \le 1 + \frac{1}{a_n}.
\end{equation*}
Since $a_n \to \infty$, we have $\frac{1}{a_n} \to 0$, hence for all sufficiently large $n$,
\begin{equation*}
1 \le \frac{c_n}{a_n} \le 2.
\end{equation*}
Thus $c_n \asymp a_n$, i.e.,
\begin{equation*}
a_n \le c_n \le 2 a_n,
\end{equation*}
so the condition in (i) holds.

\smallskip
\noindent
\emph{(b) If $\inf_n a_n \ge \delta > 0$.}

Then
\begin{equation*}
c_n \le a_n + 1 \le \Big(1 + \frac{1}{\delta}\Big)a_n.
\end{equation*}
Combining with $a_n \le c_n$, we obtain
\begin{equation*}
a_n \le c_n \le C a_n,
\end{equation*}
where $C = 1 + \frac{1}{\delta} < \infty$. Hence again $c_n \asymp a_n$, and the condition in (i) holds.

\smallskip
Therefore, under either (a) or (b), we may apply part (i) to conclude
\begin{equation*}
X_n = O_{\mathbb{P}}(c_n b_n)
\quad \Longleftrightarrow \quad
X_n = O_{\mathbb{P}}(a_n b_n).
\end{equation*}
\end{proof}

\begin{question}[Uniform integrability vs.\ tightness under moment bounds]
Let $(X_n)_{n\ge1}$ be real-valued random variables on a common probability space.
Assume that for some $r>0$,
\begin{equation*}
  M_r:=\sup_{n\ge1}\ \mathbb E|X_n|^r<\infty.
\end{equation*}
Prove that:
\begin{enumerate}
  \item[(a)] If $r>0$, then $(X_n)$ is tight.
  \item[(b)] If $r>1$, then $(X_n)$ is uniformly integrable.
  \item[(c)] Uniform integrability implies tightness.
  \item[(d)] The condition $r=1$ alone does \emph{not} imply uniform integrability: construct a counterexample.
\end{enumerate}
\end{question}

\begin{proof}
\textbf{(a) $r>0\Rightarrow$ tight.}
Applying Markov inquality with $r>0$ and taking the supremum over $n$ gives
\begin{equation*}
  \sup_n \mathbb P(|X_n|>K)
  \le \frac{\sup_n \mathbb E|X_n|^r}{K^r}
  \le \frac{M_r}{K^r}\xrightarrow[K\to\infty]{}0,
\end{equation*}
which is the definition of tightness.

\medskip
\textbf{(b) $r>1\Rightarrow$ UI.}
Let $r>1$ and $r^\ast:=\frac{r}{r-1}$ be the Hölder conjugate. For each $n$ and $K>0$,
\begin{equation*}
  \mathbb E\big[|X_n|\mathbbm{1}_{\{|X_n|>K\}}\big]
  \le \big(\mathbb E|X_n|^r\big)^{1/r}\,
      \big(\mathbb P(|X_n|>K)\big)^{1/r^\ast}
  \le M_r^{1/r}\,\Big(\frac{M_r}{K^r}\Big)^{1/r^\ast}
  = \frac{M_r}{K^{\,r-1}}.
\end{equation*}
Taking $\sup_n$ and then $K\to\infty$ yields
\begin{equation*}
  \lim_{K\to\infty}\ \sup_n \mathbb E\big[|X_n|\mathbbm{1}_{\{|X_n|>K\}}\big]=0,
\end{equation*}
so $(X_n)$ is uniformly integrable.

\emph{Alternative (de la Vallée Poussin).} For $r>1$, the convex function $\Phi(t)=t^r$ satisfies $\Phi(t)/t=t^{r-1}\to\infty$ as $t\to\infty$, and
\begin{equation*}
  \sup_n \mathbb E\Phi(|X_n|)=\sup_n \mathbb E|X_n|^r\le M_r<\infty,
\end{equation*}
hence $(X_n)$ is UI.

\medskip
\textbf{(c) UI $\Rightarrow$ tight.}
For any $K>0$ and any $n$,
\begin{equation*}
  K\,\mathbbm{1}_{\{|X_n|>K\}}\le |X_n|\,\mathbbm{1}_{\{|X_n|>K\}},
\end{equation*}
so taking expectations and then $\sup_n$,
\begin{equation*}
  K\,\sup_n \mathbb P(|X_n|>K)\ \le\ \sup_n \mathbb E\big[|X_n|\mathbbm{1}_{\{|X_n|>K\}}\big].
\end{equation*}
Letting $K\to\infty$ and using UI, the right-hand side tends to $0$, hence $\sup_n \mathbb P(|X_n|>K)\to 0$, i.e.\ tightness.

\medskip
\textbf{(d) Counterexample for $r=1$.}
Let $X_n:=n\,\mathbbm{1}_{A_n}$ where $\mathbb P(A_n)=1/n$ (the events $A_n$ can be arbitrary). Then
\begin{equation*}
  \mathbb E|X_n|=n\cdot \mathbb P(A_n)=1\quad\text{for all }n,
\end{equation*}
but for any $K\ge0$,
\begin{equation*}
  \sup_n \mathbb P(|X_n|>K)=\sup_{n>K}\frac{1}{n}=\frac{1}{\lfloor K\rfloor+1}\xrightarrow[K\to\infty]{}0,
\end{equation*}
so $(X_n)$ is tight. However,
\begin{equation*}
  \sup_n \mathbb E\big[|X_n|\mathbbm{1}_{\{|X_n|>K\}}\big]
  =\sup_{n>K} n\cdot \mathbb P(A_n)=\sup_{n>K} 1=1,
\end{equation*}
which does not converge to $0$ as $K\to\infty$. Hence $(X_n)$ is \emph{not} uniformly integrable. This shows that $r=1$ alone is insufficient for UI.
\end{proof}

\begin{question}[Weak convergence implies $O_P(1)$ (tightness)]
Let $X_n \Rightarrow X$ in $\mathbb{R}^d$. Prove that $(X_n)$ is tight, i.e., for every $\varepsilon>0$ there exists $M<\infty$ such that
\begin{equation*}
\sup_{n\ge 1}\,\mathbb{P}(\|X_n\|>M) \le \varepsilon .
\end{equation*}
Equivalently, $X_n = O_P(1)$.
\end{question}

\begin{proof}
For $M>0$, define the closed set
\begin{equation*}
F_M := \{x\in\mathbb{R}^d : \|x\|\ge M\}.
\end{equation*}
By the Portmanteau theorem,
\begin{equation*}
\limsup_{n\to\infty}\mathbb{P}(X_n\in F_M) \le \mathbb{P}(X\in F_M)
= \mathbb{P}(\|X\|\ge M).
\end{equation*}
Since $\|X\|<\infty$ almost surely, we have $\mathbb{P}(\|X\|\ge M)\to 0$ as $M\to\infty$. Hence, given any $\varepsilon>0$, choose $M$ large enough so that $\mathbb{P}(\|X\|\ge M)\le \varepsilon/2$. Then there exists $N$ such that for all $n\ge N$,
\begin{equation*}
\mathbb{P}(\|X_n\|\ge M)=\mathbb{P}(X_n\in F_M)
\le \mathbb{P}(X\in F_M)+\varepsilon/2 \le \varepsilon .
\end{equation*}
Therefore $\sup_{n\ge N}\mathbb{P}(\|X_n\|>M)\le \varepsilon$, which is tightness. In particular, $X_n=O_P(1)$.
\end{proof}

\begin{question}
Let $X$ and $Y$ be (marginally) normal random variables.
\begin{enumerate}
\item[(a)] Must the sum $X+Y$ be normal?
\item[(b)] Can $\var(X+Y)$ be smaller than both $\var(X)$ and $\var(Y)$?
\end{enumerate}
\end{question}

\begin{proof}
\emph{(a) Not necessarily.} 
It is not enough that $X$ and $Y$ are each normal; dependence can make $X+Y$ non-normal.
Construct the following counterexample: let $X\sim N(0,1)$ and let $Z\in\{-1,+1\}$ be independent of $X$ with $\p(Z=1)=\p(Z=-1)=\tfrac12$. Define $Y:=ZX$.
Then $Y\mid\{Z=1\}\sim N(0,1)$ and $Y\mid\{Z=-1\}\sim N(0,1)$, hence $Y\sim N(0,1)$ marginally. However,
\begin{equation*}
X+Y=
\begin{cases}
2X,& Z=1 \ \ (\text{prob. }1/2),\\
0,& Z=-1 \ \ (\text{prob. }1/2),
\end{cases}
\end{equation*}
so $X+Y$ has an atom at $0$ of mass $1/2$ and is therefore not normal. 

On the other hand, a sufficient (indeed, necessary and sufficient) condition for any linear combination $aX+bY$ to be normal is that $(X,Y)$ be jointly normal. In particular, if $X$ and $Y$ are independent normals, then $X+Y$ is normal with
\begin{equation*}
\mathbb E(X+Y)=\mathbb E(X)+\mathbb E(Y), 
\qquad 
\var(X+Y)=\var(X)+\var(Y),
\end{equation*}
which can be verified, for instance, from the product of moment generating functions.

\medskip
\emph{(b) Yes, it can be smaller under negative correlation (but not if $X$ and $Y$ are independent).}
In general,
\begin{equation*}
\var(X+Y)=\var(X)+\var(Y)+2\,\mathrm{Cov}(X,Y).
\end{equation*}
If $X$ and $Y$ are independent, then $\mathrm{Cov}(X,Y)=0$ and hence
\begin{equation*}
\var(X+Y)=\var(X)+\var(Y)\ge \max\{\var(X),\var(Y)\},
\end{equation*}
with equality only if one variance is zero. 

If $X$ and $Y$ are (say, jointly normal and) negatively correlated with correlation $\rho\in[-1,0)$, writing $\var(X)=\sigma_X^2$, $\var(Y)=\sigma_Y^2$, we have
\begin{equation*}
\var(X+Y)=\sigma_X^2+\sigma_Y^2+2\rho\,\sigma_X\sigma_Y.
\end{equation*}
Choosing $\rho$ sufficiently negative can make $\var(X+Y)$ smaller than both $\sigma_X^2$ and $\sigma_Y^2$. For example, with $\sigma_X^2=\sigma_Y^2=1$ and $\rho=-0.9$,
\begin{equation*}
\var(X+Y)=1+1+2(-0.9)=0.2<1.
\end{equation*}
In the extreme case $\rho=-1$ and $\sigma_X=\sigma_Y$, we have $Y=-X$ almost surely and hence $\var(X+Y)=0$.
\end{proof}




\begin{proposition}[Conditional Strong Law of Large Numbers]
Let $(\Omega,\mathcal F,\mathbb P)$ be a probability space and let $\mathcal G\subset\mathcal F$ be a sub-$\sigma$-algebra.
Suppose $(Y_i)_{i\ge1}$ are integrable random variables such that \emph{conditionally on $\mathcal G$} they are i.i.d., and
\begin{equation*}
\mathbb E\left[\,|Y_1|\,\middle|\,\mathcal G\right]<\infty\quad\mathrm{a.s.}
\end{equation*}
Then
\begin{equation*}
\frac{1}{n}\sum_{i=1}^n Y_i \ \xrightarrow{\mathrm{a.s.}}\ \mathbb E\left[Y_1\,\middle|\,\mathcal G\right].
\end{equation*}
\end{proposition}

\begin{proof}
(Disintegration + Kolmogorov SLLN.) Assume a regular conditional probability $\mathbb P(\cdot\,|\,\mathcal G)(\omega)$ exists (e.g., on standard Borel spaces).
For $\omega$ in a full-measure set, define $\mathbb P_\omega(\cdot):=\mathbb P(\cdot\,|\,\mathcal G)(\omega)$ and
\begin{equation*}
m(\omega):=\mathbb E\left[Y_1\,\middle|\,\mathcal G\right](\omega).
\end{equation*}
By the hypothesis ``i.i.d. given $\mathcal G$'', under $\mathbb P_\omega$ the sequence $(Y_i)$ is i.i.d. with mean $m(\omega)$ and is integrable.
Hence, by Kolmogorov's SLLN,
\begin{equation*}
\frac{1}{n}\sum_{i=1}^n Y_i \ \xrightarrow{\text{a.s. under }\mathbb P_\omega}\ m(\omega).
\end{equation*}
Therefore $\mathbb P_\omega\!\left(\lim_{n\to\infty} n^{-1}\sum_{i=1}^n Y_i=m(\omega)\right)=1$ for a.e. $\omega$.
Integrating both sides (Fubini/Tonelli) yields
\begin{equation*}
\mathbb P\!\left(\lim_{n\to\infty}\frac{1}{n}\sum_{i=1}^n Y_i=\mathbb E\left[Y_1\,\middle|\,\mathcal G\right]\right)=1,
\end{equation*}
which is the claim.
\end{proof}


\begin{example}[Random environment: $Y_i\mid \Lambda \stackrel{iid}{\sim}N(\Lambda,1)$]
Let $\Lambda$ be a real random variable with $\mathbb E|\Lambda|<\infty$.
Conditional on $\Lambda$, let $(Y_i)_{i\ge1}$ be i.i.d.\ with distribution $N(\Lambda,1)$.
Equivalently, write
\begin{equation*}
Y_i=\Lambda+\varepsilon_i,\qquad \varepsilon_i\stackrel{iid}{\sim}N(0,1),\ \ \varepsilon_i\ \text{independent of }\Lambda.
\end{equation*}
So, $Y_i$ is not independent if $\Lambda\neq \mathrm{const.}\ \p$-a.s.. And, 
\begin{equation*}
\frac{1}{n}\sum_{i=1}^n Y_i \ \xrightarrow{\mathrm{a.s.}}\ \Lambda .
\end{equation*}
\end{example}

\begin{proof}
By construction, given $\Lambda$ the $Y_i$ are independent and $Y_i\mid\Lambda\sim N(\Lambda,1)$, so they are i.i.d.\ given $\mathcal G=\sigma(\Lambda)$.
Unconditionally, for $i\neq j$,
\begin{equation*}
\operatorname{Cov}(Y_i,Y_j)
=\mathbb E\left[\operatorname{Cov}(Y_i,Y_j\,\middle|\,\Lambda)\right]
+\operatorname{Cov}\!\left(\mathbb E[Y_i\mid\Lambda],\mathbb E[Y_j\mid\Lambda]\right)
=0+\operatorname{Var}(\Lambda).
\end{equation*}
If $\operatorname{Var}(\Lambda)>0$, then $\operatorname{Cov}(Y_i,Y_j)>0$ and hence $Y_i,Y_j$ cannot be independent.

For the last statement, 
\begin{itemize}
	\item Proof 1 (decomposition). Write $\bar Y_n=\frac{1}{n}\sum_{i=1}^n Y_i=\Lambda+\bar\varepsilon_n$ with $\bar\varepsilon_n=\frac{1}{n}\sum_{i=1}^n \varepsilon_i$.
By Kolmogorov's SLLN, $\bar\varepsilon_n\to 0$ almost surely (indeed, both conditionally on $\Lambda$ and unconditionally).
Hence $\bar Y_n\to \Lambda$ almost surely.
	\item Proof 2 (conditional SLLN). Let $\mathcal G=\sigma(\Lambda)$. Then $Y_i$ are i.i.d.\ given $\mathcal G$ with
\begin{equation*}
\mathbb E[Y_1\mid\mathcal G]=\Lambda,\qquad \mathbb E\left[\,|Y_1|\,\middle|\,\mathcal G\right]<\infty.
\end{equation*}
Apply the Conditional SLLN (the proposition) to deduce $\bar Y_n\to \mathbb E[Y_1\mid\mathcal G]=\Lambda$ almost surely.
\end{itemize}
\end{proof}

\begin{question}
	Suppose $X$ is a random variable in the probability space $(\Omega,\f,\p)$, $\a\subset \f$ is a sub $\sigma$-algebra, and $X\in L^r$ for any $r\in (0,\infty)$.  Then $||\mathbb{E} (X|\a)||_r\leq ||X||_r$. When $r=2$, the inequality is the well-known \emph{triangle inequality}.
\end{question}
\begin{proof}
	The H\"older's inequality yields
	\begin{equation*}
		\begin{split}
			||\mathbb{E} (X|\a)||_r&=\mathbb{E}[|\mathbb{E}(X|\a)|^r]^{1/r}\leq \mathbb{E}[\mathbb{E}(|X||\a)^r]^{1/r}\\
			&\leq \mathbb{E}[\mathbb{E}(|X|^r|\a)]^{1/r}=\mathbb E(|X|^r)^{1/r}=||X||_r.
		\end{split}
	\end{equation*}
\end{proof}

\begin{question}
Let $(X_n)$ and $(Y_n)$ be $\r^d$-valued random variables. If
\begin{equation*}
Y_n - X_n \xrightarrow{\p} 0
\quad\text{and}\quad
X_n \xrightarrow{\mathrm{d}} X,
\end{equation*}
then
\begin{equation*}
Y_n \xrightarrow{\mathrm{d}} X.
\end{equation*}
\end{question}

\begin{proof}
By the Portmanteau lemma, it is enough to show that for every bounded Lipschitz function $f:\mathbb{R}\to\mathbb{R}$,
\begin{equation*}
\mathbb{E}[f(Y_n)] \to \mathbb{E}[f(X)] .
\end{equation*}
Fix such an $f$. There exist $M<\infty$ and $K<\infty$ such that $|f(u)|\le M$ for all $u$ and
\begin{equation*}
|f(u)-f(v)| \le K\,|u-v| \quad\text{for all }u,v\in\mathbb{R}.
\end{equation*}
For any $\varepsilon>0$,
\begin{equation*}
\begin{aligned}
&|\mathbb{E}[f(Y_n)]-\mathbb{E}[f(X_n)]|
\le \mathbb{E}[\,|f(Y_n)-f(X_n)|\,] \\
=& \mathbb{E}[\,|f(Y_n)-f(X_n)|\,\mathbbm{1}_{[|Y_n-X_n|\le\varepsilon]}]
   + \mathbb{E}[\,|f(Y_n)-f(X_n)|\,\mathbbm{1}\{|Y_n-X_n|>\varepsilon\}] \\
\le& K\,\varepsilon + 2M\,\mathbb{P}(|Y_n-X_n|>\varepsilon).
\end{aligned}
\end{equation*}
Hence, by the triangle inequality,
\begin{equation*}
\begin{split}
	|\mathbb{E}[f(Y_n)]-\mathbb{E}[f(X)]|
&\le |\mathbb{E}[f(Y_n)]-\mathbb{E}[f(X_n)]|
 + |\mathbb{E}[f(X_n)]-\mathbb{E}[f(X)]|\\
&\le K\varepsilon + 2M\,\mathbb{P}(|Y_n-X_n|>\varepsilon) + |\mathbb{E}[f(X_n)]-\mathbb{E}[f(X)]|.
\end{split}
\end{equation*}
Taking $n\to\infty$, the second term vanishes because $Y_n-X_n\xrightarrow{\p}0$, and the third term vanishes by the Portmanteau lemma since $X_n\xrightarrow{\mathrm{d}}X$. Therefore,
\begin{equation*}
\limsup_{n\to\infty}\big|\mathbb{E}[f(Y_n)]-\mathbb{E}[f(X)]\big| \le K\varepsilon .
\end{equation*}
Because $\varepsilon>0$ was arbitrary, the limit is $0$, so $\mathbb{E}[f(Y_n)]\to\mathbb{E}[f(X)]$ for all bounded Lipschitz $f$. Another application of the Portmanteau lemma yields $Y_n\xrightarrow{\mathrm{d}}X$, as claimed.
\end{proof}


\begin{theorem}
Let $(X_n)$ and $(Y_n)$ be $\r^d$-valued random variables. Suppose $X_n \xrightarrow{\mathrm{d}} X$ and $Y_n \xrightarrow{\p} c$, where $c$ is a constant. Then
\begin{equation*}
 (X_n,\,Y_n) \xrightarrow{\mathrm{d}} (X,\,c).
\end{equation*}
\end{theorem}

\begin{proof}
 Show that $(X_n,\,c)\xrightarrow{\mathrm{d}}(X,\,c)$.
By the portmanteau lemma, it suffices to verify that for every bounded continuous $f:\mathbb{R}^2\to\mathbb{R}$,
\begin{equation*}
 \mathbb{E}\!\left[f(X_n,\,c)\right] \longrightarrow \mathbb{E}\!\left[f(X,\,c)\right].
\end{equation*}
Fix such $f$ and define the bounded continuous function $g:\mathbb{R}\to\mathbb{R}$ by $g(x):=f(x,\,c)$. Since $X_n\xrightarrow{\mathrm{d}}X$, another application of the portmanteau lemma yields
\begin{equation*}
 \mathbb{E}\!\left[f(X_n,\,c)\right]=\mathbb{E}[g(X_n)] \longrightarrow \mathbb{E}[g(X)]=\mathbb{E}\!\left[f(X,\,c)\right],
\end{equation*}
Therefore $(X_n,\,c)\xrightarrow{\mathrm{d}}(X,\,c)$.
Since
\begin{equation*}
 \|(X_n,\,Y_n)-(X_n,\,c)\| \asymp \|Y_n-c\| \xrightarrow{\p} 0,
\end{equation*}
applying the previous QUESTION, we have
\begin{equation*}
 (X_n,\,Y_n) \xrightarrow{\mathrm{d}} (X,\,c).
\end{equation*}
\end{proof}

\begin{question}
Let $(D,d)$ be a metric space and let $\mathcal B_{\mathrm{ball}}(D)$ denote the ball $\sigma$-field on $D$, i.e. the smallest $\sigma$-field containing all open balls $B(s,r)=\{x\in D:d(x,s)<r\}$. Fix a sequence $(s_i)_{i\ge1}\subset D$ and define
\begin{equation*}
\Phi:D\to \mathbb R^\infty,\qquad 
\Phi(x)=(d(x,s_1),d(x,s_2),\ldots).
\end{equation*}
If $g:\mathbb R^\infty\to\mathbb R$ is Borel measurable (with respect to the product topology on $\mathbb R^\infty$), show that
\begin{equation*}
f(x):=g(d(x,s_1),d(x,s_2),\ldots)=g(\Phi(x))
\end{equation*}
is ball measurable on $D$, i.e. $f$ is $\mathcal B_{\mathrm{ball}}(D)/\mathcal B(\mathbb R)$-measurable.
\end{question}

\begin{proof}
Let $\pi_i:\mathbb R^\infty\to\mathbb R$ be the $i$th coordinate projection, $\pi_i((x_j)_{j\ge1})=x_i$. Equip $\mathbb R^\infty=\prod_{i\ge1}\mathbb R$ with the product topology, and let $\mathcal B(\mathbb R^\infty)$ be its Borel $\sigma$-field. In the countable product case,
\begin{equation*}
\mathcal B(\mathbb R^\infty)
=\bigotimes_{i\ge1}\mathcal B(\mathbb R)
=\sigma\big(\pi_i^{-1}(B): i\ge1,\ B\in\mathcal B(\mathbb R)\big).
\end{equation*}
Therefore, to prove that $\Phi$ is $\mathcal B_{\mathrm{ball}}(D)/\mathcal B(\mathbb R^\infty)$-measurable, it suffices to prove that each coordinate map $\pi_i\circ\Phi$ is $\mathcal B_{\mathrm{ball}}(D)/\mathcal B(\mathbb R)$-measurable.

Fix $i\ge1$. For all $x\in D$ we have $(\pi_i\circ\Phi)(x)=d(x,s_i)$. For any $r\in\mathbb R$,
\begin{equation*}
\{x\in D: d(x,s_i)<r\}
=
\begin{cases}
\emptyset, & r\le 0,\\
B(s_i,r), & r>0.
\end{cases}
\end{equation*}
Hence $\{x:d(x,s_i)<r\}\in\mathcal B_{\mathrm{ball}}(D)$ for every $r$, so $d(\cdot,s_i)$ is $\mathcal B_{\mathrm{ball}}(D)/\mathcal B(\mathbb R)$-measurable. Since this holds for every $i$, it follows that $\Phi$ is $\mathcal B_{\mathrm{ball}}(D)/\mathcal B(\mathbb R^\infty)$-measurable.

Finally, since $g$ is $\mathcal B(\mathbb R^\infty)/\mathcal B(\mathbb R)$-measurable and $f=g\circ\Phi$, the composition rule yields that $f$ is $\mathcal B_{\mathrm{ball}}(D)/\mathcal B(\mathbb R)$-measurable. Thus $f$ is ball measurable.
\end{proof}



\begin{question}
	Suppose $(A_n)_{n\geq 1}$ and $(B_n)_{n\geq 1}$ are two sequences of arbitrary sets (probably not measurable) satisfying $\p^*(A_n)=0,\p_*(B_n)=1$ for any $n\geq 1$. Then 
\begin{equation*}
	\p^*\Big(\bigcup_{n\geq1} A_n\Big)=0,\qquad \p_*\Big(\bigcap_{n\geq 1}B_n\Big)=1.
\end{equation*}
\end{question}

\begin{proof}
	For the first equality, we just need to note that by the $\sigma$-subadditivity,
\begin{equation*}
	\p^*\Big(\bigcup_{n\geq 1}A_n\Big)\leq \sum_{n\geq 1}\p^*(A_n)=0.
\end{equation*}
For the second, since by definition, $\p_*(B)=1-\p^*(B^c)$ for any set $B\subset \Omega$, we have $\p^*(B_n^c)=0$ for any $n\geq 1$.
\begin{equation*}
	\p_*\Big(\bigcap_{n\geq1} B_n\Big)=1-\p^*\Big(\bigcup_{n\geq1} B^c_n\Big)=1,
\end{equation*}
where the last equality follows from the first equality.
\end{proof}

\begin{question}[Separability of a random element in a countable product space]
Let $(D_i,d_i)$ be metric spaces for $i\in\mathbb N$, and let
\begin{equation*}
D := \prod_{i=1}^\infty D_i
\end{equation*}
be endowed with the product topology. Let $X:\Omega\to D$ be a Borel measurable
map and write $X_i := \pi_i\circ X$ for its $i$th coordinate, where
$\pi_i:D\to D_i$ is the coordinate projection. Show that $X$ is separable iff $X_i$ is separable $\forall i\in\mathbb N$.
\end{question}

\begin{proof}
We prove both implications.

\medskip
\noindent
\textbf{($\Rightarrow$)} Suppose $X$ is separable. Then there exists a separable set
$S\subset D$ such that $\mathbb P(X\in S)=1$. Fix $i\in\mathbb N$ and set $S_i := \pi_i(S)\subset D_i.$
\paragraph{Claim:} The continuous image of a separable set is separable. 

{\itshape Proof:} If topology space $T$ is separable, there exist a dense subset $A\subset T$. Next, suppose $f$ is the continuous mapping. We show $f(A)$ is a dense subset of $f(T)$. For any $y\in f(T)$ and a corresponding open neighborhood $y\in U$, there exists $x\in T$ such that $f(x)=y$ and $x\in f^{-1}(U)$ is open. Because $A$ is dense, there exists $x'\in A$ satisfying $x'\in f^{-1}(U)$, which is equivalent to $f(x')\in U$. Noting $f(x')\in f(A)$, we can conclude that $f(A)$ is dense in $f(T)$.  

Since $\pi_i$ is continuous, the image of a separable set is separable, hence $S_i$
is separable in $D_i$. Moreover, because $X(\omega)\in S$ implies $X_i(\omega)=\pi_i(X(\omega))\in \pi_i(S)=S_i$, we have $\{X\in S\}\subset \{X_i\in S_i\}$. So,
\begin{equation*}
1=\mathbb P(X\in S)\le \mathbb P(X_i\in S_i)\le 1,
\end{equation*}
so $\mathbb P(X_i\in S_i)=1$. Hence $X_i$ is separable.

\medskip
\noindent
\textbf{($\Leftarrow$)} Suppose that for every $i\in\mathbb N$ there exists a separable
set $S_i\subset D_i$ such that $\mathbb P(X_i\in S_i)=1$. Define
\begin{equation*}
S := \prod_{i=1}^\infty S_i \subset D.
\end{equation*}
We claim that $S$ is separable in the product topology. Indeed, for each $i$ choose
a countable dense subset $A_i\subset S_i$, and fix points $a_i^0\in A_i$. Let $A$ be
the set of sequences $(a_1,a_2,\dots)$ such that $a_i\in A_i$ for $1\le i\le N$ and
$a_i=a_i^0$ for $i>N$ for some $N$. Then $A$ is countable (a countable union of finite
products of countable sets). To see that $A$ is dense in $S$, let
$U=\prod_{i=1}^\infty U_i$ be a basic open neighborhood in $S$; then $U_i=S_i$ for all
but finitely many $i$, and for those finitely many indices we can pick
$a_i\in A_i\cap U_i$ since $A_i$ is dense in $S_i$. Setting $a_i=a_i^0$ for the remaining
coordinates yields a point in $A\cap U$. Hence $S$ is separable.

Finally, since the intersection of countably many outer-probability-one events has
outer probability one, we have
\begin{equation*}
\mathbb P(X\in S)
=
\mathbb P\Big(\bigcap_{i=1}^\infty \{X_i\in S_i\}\Big)
=1.
\end{equation*}
Therefore, $X$ is separable.

\medskip
\noindent
The proof is complete.
\end{proof}

\begin{question}
Let $Y_n\sim \mathcal N(\mu_n,\sigma_n^2)$ be a sequence of (possibly non-centered) normal random variables.

\smallskip
\noindent (a) Assume $Y_n \Rightarrow Y$ for some random variable $Y$. Prove that there exist $\mu\in\mathbb R$ and $\sigma^2\ge 0$ such that $\mu_n\to \mu$ and $\sigma_n^2\to \sigma^2$, and moreover $Y\sim \mathcal N(\mu,\sigma^2)$ (with the convention that $\sigma^2=0$ means $Y\equiv \mu$ a.s.).

\smallskip
\noindent (b) Specialize to the case $Y_n \Rightarrow 0$ (i.e., the limit is the constant $0$). Prove that $\mu_n\to 0$ and $\sigma_n^2\to 0$.

\smallskip
\noindent (c) In the setting of (b), fix any $p>0$ and prove that $\mathbb E|Y_n|^p\to 0$.
\end{question}

\begin{proof}
Write the characteristic function of $Y_n$ as
\begin{equation*}
\varphi_n(u):=\mathbb E e^{iuY_n}
=\exp\!\Big(iu\mu_n-\tfrac12 u^2\sigma_n^2\Big),\qquad u\in\mathbb R.
\end{equation*}
If $Y_n\Rightarrow Y$, then by Lévy's continuity theorem $\varphi_n(u)\to \varphi(u):=\mathbb E e^{iuY}$ for every $u\in\mathbb R$.

\smallskip
\noindent\textbf{(a) Convergence of $\sigma_n^2$ and identification of $\sigma^2$.}
Fix any $u\neq 0$. Taking moduli and using continuity of $z\mapsto |z|$,
\begin{equation*}
|\varphi_n(u)|
=\exp\!\Big(-\tfrac12 u^2\sigma_n^2\Big)
\longrightarrow
|\varphi(u)|.
\end{equation*}
Since the map $x\mapsto \exp(-\tfrac12 u^2 x)$ is continuous and strictly decreasing on $[0,\infty)$, the existence of the limit forces $\sigma_n^2$ to converge. Indeed, if $\sigma_n^2$ did not converge, there would exist two subsequences $\sigma_{n_k}^2\to a$ and $\sigma_{m_k}^2\to b$ with $a\neq b$, which would imply
\begin{equation*}
|\varphi_{n_k}(u)|\to e^{-\frac12u^2 a}
\quad\text{and}\quad
|\varphi_{m_k}(u)|\to e^{-\frac12u^2 b},
\end{equation*}
contradicting the fact that $|\varphi_n(u)|$ has a unique limit $|\varphi(u)|$.
Hence there exists $\sigma^2\ge 0$ such that $\sigma_n^2\to \sigma^2$, and necessarily
\begin{equation*}
|\varphi(u)|=\exp\!\Big(-\tfrac12 u^2\sigma^2\Big)\qquad\text{for all }u\in\mathbb R.
\end{equation*}

Or we can simply choose $u_0\in\r$. Then the continuity of $x\mapsto -\frac{2}{u^2}\log x$ yields
\begin{equation*}
	-\frac{2}{u_0^2}\log |\varphi_n(u_0)|=\sigma_n^2\to -\frac{2}{u_0^2}\log |\varphi(u_0)|=:\sigma^2.
\end{equation*}
So, $\sigma_n^2$ converges to $\sigma^2$.

\smallskip
\noindent\textbf{(a) Convergence of $\mu_n$ and identification of $\mu$.}
For each $u\in\mathbb R$,
\begin{equation*}
e^{iu\mu_n}
=\varphi_n(u)\exp\!\Big(\tfrac12 u^2\sigma_n^2\Big)
\longrightarrow
g(u):=\varphi(u)\exp\!\Big(\tfrac12 u^2\sigma^2\Big).
\end{equation*}
In particular, $e^{iu\mu_n}$ converges for each $u$, so it has a unique limit $g(u)$.
If $\mu_n$ did not converge, there would exist two subsequences $\mu_{n_k}\to \mu'$ and $\mu_{m_k}\to \mu''$ with $\mu'\neq \mu''$. Then for each fixed $u$,
\begin{equation*}
e^{iu\mu_{n_k}}\to e^{iu\mu'}
\quad\text{and}\quad
e^{iu\mu_{m_k}}\to e^{iu\mu''},
\end{equation*}
but both limits must equal the unique limit $g(u)$; hence $e^{iu\mu'}=e^{iu\mu''}$ for all $u$.
Indeed, let $d:=\mu'-\mu''\neq 0$. Choose $u:=\pi/|d|$. Then $ud=\pm\pi$ and thus
\begin{equation*}
e^{iu\mu'}=e^{iu\mu''}e^{iud}=-e^{iu\mu''}\neq e^{iu\mu''},
\end{equation*}
a contradiction. Therefore $\mu_n$ converges to some $\mu\in\mathbb R$.

\smallskip
\noindent\textbf{(a) Identification of the limit distribution.}
Passing to the limit in the explicit form of $\varphi_n(u)$ yields, for every $u$,
\begin{equation*}
\varphi(u)=\lim_{n\to\infty}\varphi_n(u)
=\exp\!\Big(iu\mu-\tfrac12 u^2\sigma^2\Big).
\end{equation*}
This is exactly the characteristic function of $\mathcal N(\mu,\sigma^2)$ (with $\sigma^2=0$ interpreted as the degenerate law at $\mu$). Hence $Y\sim \mathcal N(\mu,\sigma^2)$ and $\mu_n\to\mu$, $\sigma_n^2\to\sigma^2$.

\smallskip
\noindent\textbf{(b) The special case $Y_n\Rightarrow 0$.}
Applying (a) with $Y\equiv 0$ gives $\mu=0$ and $\sigma^2=0$, hence $\mu_n\to 0$ and $\sigma_n^2\to 0$.

\smallskip
\noindent\textbf{(c) Convergence of $p$th absolute moments in the case $Y_n\Rightarrow 0$.}
Let $Z\sim \mathcal N(0,1)$. Then $Y_n\stackrel{d}{=}\mu_n+\sigma_n Z$.
Fix any $p>0$. Using the inequality $|a+b|^p\le 2^{(p-1)^+}(|a|^p+|b|^p)$ (valid for all $p>0$), we obtain in all cases $p>0$ that
\begin{equation*}
|\,\mu_n+\sigma_n Z\,|^p
\le 2^{(p-1)^+}\big(|\mu_n|^p+\sigma_n^p|Z|^p\big).
\end{equation*}
Taking expectations and using $\mathbb E|Z|^p<\infty$ for every $p>0$,
\begin{equation*}
\mathbb E|Y_n|^p
=\mathbb E|\mu_n+\sigma_n Z|^p
\le 2^{(p-1)^+}\Big(|\mu_n|^p+\sigma_n^p\,\mathbb E|Z|^p\Big).
\end{equation*}
By (b), $|\mu_n|^p\to 0$ and $\sigma_n^p\to 0$, hence the right-hand side converges to $0$, proving $\mathbb E|Y_n|^p\to 0$.
\end{proof}


\begin{question}[Egorov's theorem: an ``if and only if'' on finite measure spaces]
Let $(X,\mathcal A,\mu)$ be a \emph{finite} measure space, i.e.\ $\mu(X)<\infty$, and let
$(f_n)_{n\ge1}$ and $f$ be measurable real-valued functions on $X$. Prove that if $f_n\to f$ almost everywhere on $X$ iff $f_n\to f$ almost uniformly on $X$.
\end{question}

\begin{proof}
\textbf{(a) Egorov.}
Assume $f_n\to f$ almost everywhere. For $k,n\in\mathbb N$, define
\begin{equation*}
E_{k,n}:=\bigcup_{m\ge n}\Big\{x\in X:\ |f_m(x)-f(x)|>\frac{1}{k}\Big\}.
\end{equation*}
Then $(E_{k,n})_{n\ge1}$ is decreasing in $n$ and
\begin{equation*}
\bigcap_{n=1}^\infty E_{k,n}
=\varlimsup_{n\to\infty}E_{k,n}=
\Big\{x\in X:\ |f_m(x)-f(x)|>\frac{1}{k}\ \text{i.o.}\Big\}.
\end{equation*}
Since $f_n(x)\to f(x)$ for almost all $x$, we have
\begin{equation*}
\mu\!\left(\bigcap_{n=1}^\infty E_{k,n}\right)=0
\qquad\text{for each fixed } k.
\end{equation*}
By continuity from above of $\mu$ (using $\mu(X)<\infty$ and $E_{k,n}\downarrow \cap_n E_{k,n}$), we have $\mu(E_{k,n})\downarrow 0,n\to\infty$.
Fix $\varepsilon>0$. For each $k$, choose $n_k$ such that $\mu(E_{k,n_k})<\frac{\varepsilon}{2^k}$.
Let $A:=\bigcup_{k=1}^\infty E_{k,n_k}.$
Then
\begin{equation*}
\mu(A)\le \sum_{k=1}^\infty \mu(E_{k,n_k})<\varepsilon.
\end{equation*}
We claim that $f_n\to f$ uniformly on $X\setminus A$. Indeed, fix $\eta>0$ and choose $k$ with
$1/k<\eta$. If $x\in X\setminus A$, then $x\notin E_{k,n_k}$, hence for all $m\ge n_k$, $|f_m(x)-f(x)|\le \frac{1}{k}<\eta.$
Therefore,
\begin{equation*}
\sup_{x\in X\setminus A}|f_m(x)-f(x)|\le \frac{1}{k}<\eta
\qquad \text{for all } m\ge n_k,
\end{equation*}
which proves uniform convergence on $X\setminus A$. Since $\mu(A)<\varepsilon$ and $\varepsilon$ was
arbitrary, this is almost uniform convergence.

\medskip\noindent
\textbf{(b) Almost uniform $\Rightarrow$ almost everywhere.}
Assume $f_n\to f$ almost uniformly. For each $m\in\mathbb N$, apply the definition with
$\varepsilon=2^{-m}$ to obtain $A_m\in\mathcal A$ such that $\mu(A_m)<2^{-m}$
and $\sup_{x\in X\setminus A_m}|f_n(x)-f(x)|\to 0,n\to\infty$.
Define $B:=\bigcap_{m=1}^\infty A_m.$
Then $B\subseteq A_m$ for every $m$, hence
\begin{equation*}
\mu(B)\le \mu(A_m)<2^{-m}\qquad \forall m,
\end{equation*}
which implies $\mu(B)=0$.
Fix $x\in X\setminus B$. Then there exists $m_0$ such that $x\notin A_{m_0}$, i.e.\ $x\in X\setminus A_{m_0}$.
By uniform convergence on $X\setminus A_{m_0}$,
\begin{equation*}
|f_n(x)-f(x)|
\le
\sup_{y\in X\setminus A_{m_0}}|f_n(y)-f(y)|
\to 0,
\end{equation*}
so $f_n(x)\to f(x)$. Since this holds for all $x\in X\setminus B$ and $\mu(B)=0$, we conclude that
$f_n\to f$ almost everywhere.
\end{proof}

\begin{question}
Let $(X_n)_{n\ge1}$ be a sequence of random variables and let $X$ be a random variable,
all defined on the same probability space.
Show that the following statements are equivalent:
\begin{itemize}
	\item $X_n \xrightarrow{\mathbb P} X$.
	\item for every subsequence, there exists a further subsequence $(X_{n_{k_\ell}})$ such that $X_{n_{k_\ell}} \xrightarrow{\mathrm{a.s.}} X$.
\end{itemize}
\end{question}

\begin{proof}
We prove the two implications separately.

\medskip
\noindent\textbf{($\Rightarrow$)}
Assume that $X_n \xrightarrow{\mathbb P} X$.
Let $(X_{n_k})$ be an arbitrary subsequence.
By convergence in probability, for each $\ell\ge1$ we can choose $k_\ell$ large enough such that
\begin{equation*}
\mathbb P(|X_{n_{k_\ell}}-X|>\tfrac{1}{\ell})
<2^{-\ell}.
\end{equation*}
This defines a further subsequence $(X_{n_{k_\ell}})$ of $(X_{n_k})$.
Since
\begin{equation*}
\sum_{\ell=1}^\infty 
\mathbb \mathbb P(|X_{n_{k_\ell}}-X|>\tfrac{1}{\ell})
<\infty,
\end{equation*}
the Borel--Cantelli lemma implies that
\begin{equation*}
|X_{n_{k_\ell}}-X|>\tfrac{1}{\ell}
\end{equation*}
occurs only finitely many times almost surely.
Hence $X_{n_{k_\ell}}\xrightarrow{\mathrm{a.s.}}X$.

\medskip
\noindent\textbf{($\Leftarrow$)}
Assume that for every subsequence $(X_{n_k})$ there exists a further subsequence
$(X_{n_{k_\ell}})$ such that $X_{n_{k_\ell}}\xrightarrow{\mathrm{a.s.}}X$.
Suppose, by contradiction, that $X_n\not\xrightarrow{\mathbb P}X$.
Then there exist $\varepsilon>0$, $\delta>0$, and a subsequence $(X_{n_k})$ such that
\begin{equation*}
\mathbb P(|X_{n_k}-X|>\varepsilon)\ge\delta
\quad\text{for all }k.
\end{equation*}
Define the events
\begin{equation*}
A_\ell:=\{|X_{n_{k_\ell}}-X|>\varepsilon\},
\qquad \ell\ge1,
\end{equation*}
for any further subsequence $(X_{n_{k_\ell}})$ of $(X_{n_k})$.
Then $\mathbb P(A_\ell)\ge\delta$ for all $\ell$.
For each $m\ge1$,
\begin{equation*}
\mathbb P\Big(\bigcup_{\ell\ge m} A_\ell\Big)
\ge \mathbb P(A_m)\ge\delta.
\end{equation*}
By continuity of probability from above,
\begin{equation*}
\mathbb P\Big(\varlimsup_{\ell\to\infty} A_\ell\Big)
=
\lim_{m\to\infty}
P\Big(\bigcup_{\ell\ge m} A_\ell\Big)
\ge\delta>0.
\end{equation*}
Hence the event $\varlimsup_{\ell\to\infty}A_\ell$ has positive probability, which means that
\begin{equation*}
|X_{n_{k_\ell}}-X|>\varepsilon
\end{equation*}
occurs infinitely often with positive probability.
Therefore $(X_{n_{k_\ell}})$ cannot converge to $X$ almost surely.
This contradicts the assumption that every subsequence admits an almost surely convergent
subsequence.
Consequently,
\begin{equation*}
X_n \xrightarrow{\mathbb P} X.
\end{equation*}
\end{proof}

\begin{question}
Does convergence in probability of a sequence $(X_n)_{n\ge1}$ to a random variable $X$
imply convergence in probability of the tail supremum $\sup_{k\ge n} X_k$
to $X$?
\end{question}

\begin{proof}
The answer is negative in general. We construct a counterexample.

Let $(X_n)_{n\ge1}$ be a sequence of independent random variables defined by
\begin{equation*}
X_n =
\begin{cases}
1, & \text{with probability } \frac{1}{n},\\[4pt]
0, & \text{with probability } 1-\frac{1}{n}.
\end{cases}
\end{equation*}

First we show $X_n \xrightarrow{\mathbb P} 0$. Fix $\varepsilon\in(0,1)$. Then
\begin{equation*}
\mathbb P(|X_n-0|>\varepsilon)
=
\mathbb P(X_n=1)
=
\frac{1}{n}
\;\xrightarrow[]{n\to\infty}\;0.
\end{equation*}
Hence $X_n \xrightarrow{\mathbb P} 0$.

Next, we show $\sup_{k\ge n} X_k$ does not converge in probability to $0$. Observe that
\begin{equation*}
\sup_{k\ge n} X_k = 0
\quad\Longleftrightarrow\quad
X_k=0 \text{ for all } k\ge n.
\end{equation*}
By independence,
\begin{equation*}
\mathbb P\Big(\sup_{k\ge n} X_k = 0\Big)
=
\prod_{k=n}^{\infty}\left(1-\frac{1}{k}\right).
\end{equation*}

Using the inequality $\log(1-x)\le -x$ for $x\in(0,1)$, we obtain
\begin{equation*}
\sum_{k=n}^{\infty}\log\!\left(1-\frac{1}{k}\right)
\le
-\sum_{k=n}^{\infty}\frac{1}{k}
=
-\infty,
\end{equation*}
since the harmonic series diverges. Therefore,
\begin{equation*}
\prod_{k=n}^{\infty}\left(1-\frac{1}{k}\right)=0.
\end{equation*}

Consequently, for any $\varepsilon>0$
\begin{equation*}
\p\Big(\big|\sup_{k\ge n} X_k-0\big|>\varepsilon\Big)= \mathbb P\Big(\sup_{k\ge n} X_k = 1\Big)=1
\quad \text{for all } n.
\end{equation*}
Thus 
\begin{equation*}
\sup_{k\ge n} X_k \xrightarrow{\mathbb P} 1 \neq 0.
\end{equation*}
\end{proof}

\begin{question}
  	Let $r>0$ and let $X\in L^r(\p)$. Then
\begin{equation*}
x^r \mathbb{P}(|X|>x) \longrightarrow 0
\quad \text{as } x \to \infty.
\end{equation*}
Namely, $\p(|X|>x)=o(x^{-r})$ as $x\to\infty$.
  \end{question}

\begin{proof}
By Fubini's theorem,
\begin{equation*}
\mathbb{E}X^r
=\int_\Omega\int_0^X r t^{r-1} \, dtd\p=
\int_0^\infty r t^{r-1} \mathbb{P}(X>t)\, dt.
\end{equation*}
Because $\mathbb{E}X^r < \infty$, the tail integral
\begin{equation*}
F(x)
:=
\int_x^\infty r t^{r-1} \mathbb{P}(X>t)\, dt
\end{equation*}
satisfies $F(x) \longrightarrow 0$ as $x \to \infty.$

Since $t \mapsto \mathbb{P}(X>t)$ is nonincreasing, for $t \in [x,2x]$ we have
\begin{equation*}
\mathbb{P}(X>t)
\ge
\mathbb{P}(X>2x).
\end{equation*}
Therefore,
\begin{equation*}
\begin{split}
	F(x)
&\ge
\int_x^{2x}
r t^{r-1}
\mathbb{P}(X>t)
\, dt
\ge
\mathbb{P}(X>2x)
\int_x^{2x}
r t^{r-1} \, dt\\
&=\mathbb{P}(X>2x)[(2x)^r - x^r]
=(2^r - 1)\mathbb{P}(X>2x)x^r.
\end{split}
\end{equation*}
Since $F(x) \to 0$, it follows that
\begin{equation*}
x^r \mathbb{P}(X>2x) \longrightarrow 0.
\end{equation*}
Replacing $2x$ by $x$ yields
\begin{equation*}
x^r \mathbb{P}(|X|>x) \longrightarrow 0,
\end{equation*}
which completes the proof.
\end{proof}

\begin{question}
Let $X$ be a real-valued random variable such that
\begin{equation*}
\p(|X|>t)=o(t^{-2}) \quad \text{as } t\to\infty.
\end{equation*}
Show that
\begin{equation*}
\mathbb E [ |X|  \mathbbm 1\{|X|>t\} ]
=
o(t^{-1})
\quad \text{as } t\to\infty.
\end{equation*}
\end{question}

\begin{proof}
Let $Y:=|X|\ge 0$. We first establish the identity
\begin{equation*}
\mathbb E [Y  \mathbbm 1\{Y>t\} ]
=
t \p(Y>t)
+
\int_t^\infty \p(Y>u) du.
\end{equation*}
Indeed, observe that
\begin{equation*}
Y  \mathbbm 1\{Y>t\}
=
t  \mathbbm 1\{Y>t\}
+
(Y-t)  \mathbbm 1\{Y>t\}.
\end{equation*}
For the second term, using the layer-cake representation for nonnegative random variables,
\begin{equation*}
\mathbb E[(Y-t)  \mathbbm 1\{Y>t\}]
=
\int_0^\infty
 \mathbbm 1\{(Y-t) \mathbbm 1\{Y>t\}>u\} du.
\end{equation*}
Since $(Y-t) \mathbbm 1\{Y>t\}>u$ if and only if $Y>t+u$, we obtain
\begin{equation*}
\mathbb E[(Y-t)  \mathbbm 1\{Y>t\}]
=
\int_0^\infty  \mathbbm 1\{Y>t+u\} du=\int_t^\infty \p(Y>u) du.
\end{equation*}

Since $\p(Y>t)=o(t^{-2})$, as $t\to\infty$,
then for every $\varepsilon>0$, there exists $T$ such that for all $u\ge T$,
\begin{equation*}
\p(Y>u)\le \varepsilon u^{-2}.
\end{equation*}
For $t\ge T$, we bound the two terms separately:
\begin{equation*}
t\p(Y>t)\le \varepsilon t^{-1},
\end{equation*}
and
\begin{equation*}
\int_t^\infty \p(Y>u) du
\le
\int_t^\infty \varepsilon u^{-2} du
=
\varepsilon t^{-1}.
\end{equation*}
Hence
\begin{equation*}
\mathbb E[Y \mathbbm 1\{Y>t\}]
\le
2\varepsilon t^{-1}
\quad \forall t\geq T.
\end{equation*}
Since $\varepsilon>0$ is arbitrary, it follows that
\begin{equation*}
\mathbb E[Y \mathbbm 1\{Y>t\}]
=
o(t^{-1}),
\end{equation*}
which completes the proof.
\end{proof}

\begin{question}
Let $X$ be a nonnegative random variable. Define
\begin{equation*}
\|X\|_{2,\infty}^2
:=
\sup_{t>0} t^2 \p(X>t).
\end{equation*}
Show that
\begin{equation*}
\|X\|_{2,\infty}^2
\;\le\;
\sup_{t>0} t\,\mathbb E [X \mathbbm 1\{X>t\} ]
\;\le\;
2\,\|X\|_{2,\infty}^2.
\end{equation*}
\end{question}

\begin{proof}

\textbf{First inequality.}
It is easy to find
\begin{equation*}
\mathbb E [X \mathbbm 1\{X>t\} ]
\;\ge\; \mathbb E [t \mathbbm 1\{X>t\}]=
t\,\p(X>t).
\end{equation*}
Multiplying both sides by $t$ yields
\begin{equation*}
t\,\mathbb E [X \mathbbm 1\{X>t\} ]
\;\ge\;
t^2 \p(X>t).
\end{equation*}
Taking the supremum over $t>0$ gives
\begin{equation*}
\|X\|_{2,\infty}^2
=
\sup_{t>0} t^2 \p(X>t)
\;\le\;
\sup_{t>0} t\,\mathbb E [X \mathbbm 1\{X>t\} ].
\end{equation*}

\medskip
\noindent
\textbf{Second inequality.}
For any $t>0$, use the identity
\begin{equation*}
\mathbb E [X \mathbbm 1\{X>t\} ]
=
\int_0^\infty \p (X \mathbbm 1\{X>t\}>x )\,dx.
\end{equation*}
When $x\leq t$, we can easily check $\{X \mathbbm 1\{X>t\}>x\}=\{X>t\}$. When $x>t$, we have $\{X \mathbbm 1\{X>t\}>x\}=\{X>x\}$. So
split the integral at $t$:
\begin{equation*}
\mathbb E [X \mathbbm 1\{X>t\} ]
=
\int_0^t \p(X>t)\,dx
+
\int_t^\infty \p(X>x)\,dx.
\end{equation*}
The first term equals
\begin{equation*}
\int_0^t \p(X>t)\,dx
=
t \p(X>t).
\end{equation*}
For the second term, use the definition of $\|X\|_{2,\infty}$:
\begin{equation*}
\p(X>x)
\le
\frac{\|X\|_{2,\infty}^2}{x^2},
\qquad x>0.
\end{equation*}
Hence
\begin{equation*}
\int_t^\infty \p(X>x)\,dx
\le
\|X\|_{2,\infty}^2
\int_t^\infty x^{-2}\,dx
=
\frac{\|X\|_{2,\infty}^2}{t}.
\end{equation*}
Combining the two bounds gives
\begin{equation*}
\mathbb E [X \mathbbm 1\{X>t\} ]
\le
t \p(X>t)
+
\frac{\|X\|_{2,\infty}^2}{t}.
\end{equation*}
Multiplying by $t$ yields
\begin{equation*}
t\,\mathbb E [X \mathbbm 1\{X>t\} ]
\le
t^2 \p(X>t)
+
\|X\|_{2,\infty}^2
\le
2\,\|X\|_{2,\infty}^2.
\end{equation*}
Taking the supremum over $t>0$ completes the proof.
\end{proof}

\begin{question}[Gaussian conditioning, Brownian bridges, and $P$-Brownian bridges]
Let $(X,Y)$ be a jointly Gaussian random vector with $\mathbb E[X]=0$, $\mathbb E[Y]=0$, and $\mathrm{Var}(Y)>0$.  

\begin{enumerate}
\item Show that
\begin{equation*}
(X \mid Y=0)\ \stackrel{d}{=}\ X-\frac{\mathrm{Cov}(X,Y)}{\mathrm{Var}(Y)}\,Y.
\end{equation*}

\item Let $W=\{W(t):t\in[0,1]\}$ be standard Brownian motion. Prove that the following three definitions of the standard Brownian bridge on $[0,1]$ are equivalent (in the sense of equality of all finite-dimensional distributions, and hence equality in law as processes):
\begin{enumerate}
\item[(i)] $B(t)\ \stackrel{d}{=}\ (W(t)\mid W(1)=0)$ for each $t\in[0,1]$;
\item[(ii)] $B(t):=W(t)-tW(1)$ for $t\in[0,1]$;
\item[(iii)] $B$ is a centered Gaussian process with covariance
\begin{equation*}
\mathrm{Cov}\big(B(s),B(t)\big)=s\wedge t-st,\qquad s,t\in[0,1].
\end{equation*}
\end{enumerate}

\item Let $P$ be a probability measure on $(\mathcal X,\mathcal A)$ and let $\mathcal F\subset L^2(P)$ be a class of measurable functions. Define the $P$-Brownian bridge indexed by $\mathcal F$, denoted $\mathbb G_P=\{\mathbb G_P(f):f\in\mathcal F\}$, and prove that the following three definitions are equivalent:
\begin{enumerate}
\item[(i)] (Conditional definition) $\mathbb G_P(f)\ \stackrel{d}{=}\ (\mathbb W_P(f)\mid \mathbb W_P(1)=0)$ for all $f\in\mathcal F$, where $\mathbb W_P$ is the $P$-Brownian motion (isonormal Gaussian process) on $L^2(P)$;
\item[(ii)] (Projection definition) $\mathbb G_P(f):=\mathbb W_P(f)-Pf\cdot \mathbb W_P(1)$ for all $f\in\mathcal F$;
\item[(iii)] (Covariance characterization) $\mathbb G_P$ is a centered Gaussian process with covariance
\begin{equation*}
\mathrm{Cov}\big(\mathbb G_P(f),\mathbb G_P(g)\big)=P(fg)-Pf\,Pg,\qquad f,g\in\mathcal F.
\end{equation*}
\end{enumerate}
Finally, show that if $1\in\mathcal F$ and $P$ is the uniform distribution on $[0,1]$ with $\mathcal F=\{1\{x\le t\}:t\in[0,1]\}$, then $\{\mathbb G_P(1\{x\le t\}):t\in[0,1]\}$ is the standard Brownian bridge.
\end{enumerate}
\end{question}

\begin{proof}
\textbf{1.}
Let $X,Y$ be jointly Gaussian with mean zero and $\mathrm{Var}(Y)>0$. Define
\begin{equation*}
a:=\frac{\mathrm{Cov}(X,Y)}{\mathrm{Var}(Y)},\qquad Z:=X-aY.
\end{equation*}
Since $(X,Y)$ is jointly Gaussian and $(Z,Y)$ is obtained from $(X,Y)$ by a linear map, $(Z,Y)$ is also jointly Gaussian. Compute
\begin{equation*}
\mathrm{Cov}(Z,Y)=\mathrm{Cov}(X-aY,Y)=\mathrm{Cov}(X,Y)-a\,\mathrm{Var}(Y)=0.
\end{equation*}
For jointly Gaussian random variables, zero covariance implies independence; hence $Z$ is independent of $Y$. Therefore, for any Borel set $A\subset\mathbb R$,
\begin{equation*}
\mathbb P(Z\in A\mid Y=0)=\mathbb P(Z\in A).
\end{equation*}
On the event $\{Y=0\}$, we have $X=Z+aY=Z$. Thus
\begin{equation*}
\mathbb P(X\in A\mid Y=0)=\mathbb P(Z\in A),
\end{equation*}
which is exactly
\begin{equation*}
(X\mid Y=0)\ \stackrel{d}{=}\ Z=X-\frac{\mathrm{Cov}(X,Y)}{\mathrm{Var}(Y)}\,Y.
\end{equation*}

\medskip
\noindent
\textbf{2.}
\emph{(i) $\Leftrightarrow$ (ii).}
Fix $k\in\mathbb N$ and $t_1,\dots,t_k\in[0,1]$. Let
\begin{equation*}
X:=\big(W(t_1),\dots,W(t_k)\big)^\top,\qquad Y:=W(1).
\end{equation*}
The vector $(X,Y)$ is jointly Gaussian with mean zero. Moreover,
\begin{equation*}
\mathrm{Var}(Y)=\mathrm{Var}(W(1))=1,
\end{equation*}
and for each $j$,
\begin{equation*}
\mathrm{Cov}(W(t_j),W(1))=t_j\wedge 1=t_j.
\end{equation*}
Hence
\begin{equation*}
\mathrm{Cov}(X,Y)=\big(t_1,\dots,t_k\big)^\top.
\end{equation*}
Applying Part 1 componentwise in vector form (with $\mathrm{Var}(Y)=1$) yields
\begin{equation*}
(X\mid Y=0)\ \stackrel{d}{=}\ X-\mathrm{Cov}(X,Y)\,Y
=\big(W(t_1)-t_1W(1),\dots,W(t_k)-t_kW(1)\big)^\top.
\end{equation*}
Therefore, the finite-dimensional distributions of $\{W(t)\mid W(1)=0\}$ coincide with those of $\{W(t)-tW(1)\}$, proving (i) $\Leftrightarrow$ (ii) in the sense of equality of all finite-dimensional distributions.

\emph{(ii) $\Rightarrow$ (iii).}
Define $B(t):=W(t)-tW(1)$. Since $B$ is a linear transform of the Gaussian process $W$, the process $B$ is centered Gaussian. For $s,t\in[0,1]$,
\begin{equation*}
\mathrm{Cov}\big(B(s),B(t)\big)
=\mathrm{Cov}\big(W(s)-sW(1),\,W(t)-tW(1)\big).
\end{equation*}
Expand bilinearly:
\begin{equation*}
\begin{split}
	\mathrm{Cov}\big(B(s),B(t)\big)
=&\mathrm{Cov}(W(s),W(t))
-t\,\mathrm{Cov}(W(s),W(1))\\
&-s\,\mathrm{Cov}(W(1),W(t))
+st\,\mathrm{Var}(W(1)).
\end{split}
\end{equation*}
Using $\mathrm{Cov}(W(s),W(t))=s\wedge t$, $\mathrm{Cov}(W(s),W(1))=s$, $\mathrm{Cov}(W(1),W(t))=t$, and $\mathrm{Var}(W(1))=1$, we obtain
\begin{equation*}
\mathrm{Cov}\big(B(s),B(t)\big)=s\wedge t-ts-st+st=s\wedge t-st,
\end{equation*}
which is exactly (iii).

\emph{(iii) $\Rightarrow$ (ii).}
Let $B$ be a centered Gaussian process with covariance $s\wedge t-st$. Consider the process $\widetilde B(t):=W(t)-tW(1)$ constructed from Brownian motion $W$. By the previous step, $\widetilde B$ is a centered Gaussian with the same covariance kernel. For any $k$ and $t_1,\dots,t_k$, the vectors
\begin{equation*}
(B(t_1),\dots,B(t_k))\quad\text{and}\quad(\widetilde B(t_1),\dots,\widetilde B(t_k))
\end{equation*}
are both mean-zero multivariate normal with the same covariance matrix $(t_i\wedge t_j-t_it_j)_{i,j}$. Hence they have the same distribution. Therefore, $B$ and $\widetilde B$ have identical finite-dimensional distributions, establishing (iii) $\Rightarrow$ (ii).

\medskip
\noindent
\textbf{3.}
Let $\mathbb W_P=\{\mathbb W_P(h):h\in L^2(P)\}$ be a centered Gaussian process such that for all $h,g\in L^2(P)$,
\begin{equation*}
\mathrm{Cov}\big(\mathbb W_P(h),\mathbb W_P(g)\big)=P(hg).
\end{equation*}
In particular, since $1\in L^2(P)$, we have $\mathrm{Var}(\mathbb W_P(1))=P(1^2)=1$.

\emph{(i) $\Leftrightarrow$ (ii).}
Fix $f_1,\dots,f_k\in\mathcal F$ and set
\begin{equation*}
X:=\big(\mathbb W_P(f_1),\dots,\mathbb W_P(f_k)\big)^\top,\qquad Y:=\mathbb W_P(1).
\end{equation*}
Then $(X,Y)$ is jointly Gaussian with mean zero. Moreover, for each $j$,
\begin{equation*}
\mathrm{Cov}( \mathbb W_P(f_j),\,\mathbb W_P(1)) = P(f_j\cdot 1)=Pf_j,
\end{equation*}
so
\begin{equation*}
\mathrm{Cov}(X,Y)=\big(Pf_1,\dots,Pf_k\big)^\top,
\qquad
\mathrm{Var}(Y)=\mathrm{Var}(\mathbb W_P(1))=1.
\end{equation*}
Applying Part 1 in vector form gives
\begin{equation*}
\begin{split}
	(X\mid Y=0)\ &\stackrel{d}{=}\ X-\mathrm{Cov}(X,Y)\,Y\\
&=\big(\mathbb W_P(f_1)-Pf_1\,\mathbb W_P(1),\dots,\mathbb W_P(f_k)-Pf_k\,\mathbb W_P(1)\big)^\top.
\end{split}
\end{equation*}
Hence the conditional process $\{\mathbb W_P(f)\mid \mathbb W_P(1)=0:f\in\mathcal F\}$ has the same finite-dimensional distributions as the projected process $\{\mathbb W_P(f)-Pf\,\mathbb W_P(1):f\in\mathcal F\}$, proving (i) $\Leftrightarrow$ (ii).

\emph{(ii) $\Rightarrow$ (iii).}
Define $\mathbb G_P(f):=\mathbb W_P(f)-Pf\cdot\mathbb W_P(1)$. This is a linear transform of a centered Gaussian process, hence centered Gaussian. For $f,g\in\mathcal F$,
\begin{equation*}
\mathrm{Cov}\big(\mathbb G_P(f),\mathbb G_P(g)\big)
=\mathrm{Cov}\big(\mathbb W_P(f)-Pf\,\mathbb W_P(1),\,\mathbb W_P(g)-Pg\,\mathbb W_P(1)\big).
\end{equation*}
Expand bilinearly:
\begin{equation*}
\begin{split}
	\mathrm{Cov}\big(\mathbb G_P(f),\mathbb G_P(g)\big)
=&\mathrm{Cov}(\mathbb W_P(f),\mathbb W_P(g))
-Pg\,\mathrm{Cov}(\mathbb W_P(f),\mathbb W_P(1))\\
&-Pf\,\mathrm{Cov}(\mathbb W_P(1),\mathbb W_P(g))
+Pf\,Pg\,\mathrm{Var}(\mathbb W_P(1)).
\end{split}
\end{equation*}
Using $\mathrm{Cov}(\mathbb W_P(f),\mathbb W_P(g))=P(fg)$, $\mathrm{Cov}(\mathbb W_P(f),\mathbb W_P(1))=Pf$, $\mathrm{Cov}(\mathbb W_P(1),\mathbb W_P(g))=Pg$, and $\mathrm{Var}(\mathbb W_P(1))=1$, we obtain
\begin{equation*}
\mathrm{Cov}\big(\mathbb G_P(f),\mathbb G_P(g)\big)=P(fg)-Pg\,Pf-Pf\,Pg+Pf\,Pg=P(fg)-Pf\,Pg,
\end{equation*}
which is (iii).

\emph{Step 3.4: (iii) $\Rightarrow$ (ii).}
Let $\mathbb G_P$ be a centered Gaussian process on $\mathcal F$ with covariance $P(fg)-Pf\,Pg$. Let $\widetilde{\mathbb G}_P(f):=\mathbb W_P(f)-Pf\cdot\mathbb W_P(1)$ constructed from some $P$-Brownian motion $\mathbb W_P$ as in Step 3.1. By Step 3.3, $\widetilde{\mathbb G}_P$ is centered Gaussian with the same covariance kernel. Therefore, for any finite $f_1,\dots,f_k\in\mathcal F$, the vectors
\begin{equation*}
\big(\mathbb G_P(f_1),\dots,\mathbb G_P(f_k)\big)
\quad\text{and}\quad
\big(\widetilde{\mathbb G}_P(f_1),\dots,\widetilde{\mathbb G}_P(f_k)\big)
\end{equation*}
are both mean-zero multivariate normal with the same covariance matrix $(P(f_if_j)-Pf_i\,Pf_j)_{i,j}$. Hence, they have the same distribution, implying equality of all finite-dimensional distributions and establishing (iii) $\Rightarrow$ (ii). 

\medskip
Assume $P$ is the uniform distribution on $[0,1]$ and $\mathcal F=\{f_t:t\in[0,1]\}$ with $f_t(x)=1\{x\le t\}$. Then for $s,t\in[0,1]$,
\begin{equation*}
Pf_t=\int_0^1 1\{x\le t\}\,dx=t,
\qquad
P(f_sf_t)=\int_0^1 1\{x\le s\wedge t\}\,dx=s\wedge t.
\end{equation*}
Hence the covariance kernel in Part 3(iii) becomes
\begin{equation*}
\mathrm{Cov}\big(\mathbb G_P(f_s),\mathbb G_P(f_t)\big)=s\wedge t-st,
\end{equation*}
which matches the covariance characterization of the standard Brownian bridge in Part 2(iii). Therefore $\{\mathbb G_P(1\{x\le t\}):t\in[0,1]\}$ is (in law) the standard Brownian bridge.
\end{proof}

\begin{question}[$P$-Brownian bridge as an orthogonal projection in an isonormal Gaussian space]
Let $(\mathcal X,\mathcal A,P)$ be a probability space and let $H=L^2(P)$ with inner product $\langle f,g\rangle:=P(fg)$. Let $\mathbb W_P=\{\mathbb W_P(h):h\in H\}$ be an isonormal Gaussian process on $H$, i.e.
\begin{equation*}
\mathbb E[\mathbb W_P(h)]=0,\qquad 
\mathbb E[\mathbb W_P(h)\mathbb W_P(g)]=\langle h,g\rangle,\qquad h,g\in H.
\end{equation*}
Fix a class $\mathcal F\subset H$ and define the $P$-Brownian bridge indexed by $\mathcal F$ by
\begin{equation*}
\mathbb G_P(f):=\mathbb W_P(f)-Pf\cdot \mathbb W_P(1),\qquad f\in\mathcal F,
\end{equation*}
where $1$ denotes the constant function on $\mathcal X$.
Let $H_0:=\{h\in H:Ph=0\}$ and let $\Pi_0:H\to H_0$ be the orthogonal projection onto $H_0$.
Prove that:
\begin{enumerate}
\item (Projection structure in $H$) For every $h\in H$,
\begin{equation*}
\Pi_0 h = h-Ph\cdot 1,
\end{equation*}
and $H= \mathrm{span}\{1\}\oplus H_0$ with $H_0=\mathrm{span}\{1\}^\perp$.

\item (Bridge as projection of the isonormal process) For every $f\in\mathcal F$,
\begin{equation*}
\mathbb G_P(f)=\mathbb W_P(\Pi_0 f)=\mathbb W_P(f-Pf).
\end{equation*}
Moreover, $\mathbb G_P$ is a centered Gaussian process indexed by $\mathcal F$ with covariance
\begin{equation*}
\mathrm{Cov}\big(\mathbb G_P(f),\mathbb G_P(g)\big)=P(fg)-Pf\,Pg,\qquad f,g\in\mathcal F.
\end{equation*}

\item (Isometric embedding on the orthogonal complement) The restriction $\mathbb W_P|_{H_0}$ is isonormal on $H_0$ and hence an isometric embedding of $H_0$ into $L^2(\Omega)$:
\begin{equation*}
\|\mathbb W_P(h)\|_{L^2(\Omega)}=\|h\|_{L^2(P)},\qquad h\in H_0
\end{equation*}
and $\mathbb W_P|_{H_0} : H_0 \hookrightarrow L^2(\Omega)$.
Consequently, $\mathbb G_P$ is the image under $\mathbb W_P|_{H_0}$ of the centered class $\{f-Pf:f\in\mathcal F\}\subset H_0$.
\end{enumerate}
\end{question}

\begin{proof}
\textbf{(1)}
First note that $1\in H$ and $\|1\|_H^2=\langle 1,1\rangle=P(1)=1$. Define
\begin{equation*}
H_0:=\{h\in H:\langle h,1\rangle=Ph=0\}.
\end{equation*}
Since the map $h\mapsto \langle h,1\rangle$ is continuous linear on $H$, the set $H_0=\ker(h\mapsto \langle h,1\rangle)$ is a closed subspace of $H$.

For any $h\in H$, define
\begin{equation*}
h_0:=h-Ph\cdot 1.
\end{equation*}
Then $Ph_0=Ph-Ph\cdot P1=Ph-Ph=0$, so $h_0\in H_0$. Also, the remaining component is $Ph\cdot 1\in \mathrm{span}\{1\}$. Hence every $h\in H$ admits the decomposition
\begin{equation*}
h = Ph\cdot 1 + (h-Ph\cdot 1)
\in \mathrm{span}\{1\}+H_0.
\end{equation*}
To see the sum is orthogonal, let $u=\alpha\cdot 1\in \mathrm{span}\{1\}$ and $v\in H_0$. Then
\begin{equation*}
\langle u,v\rangle = \alpha \langle 1,v\rangle = \alpha Pv = 0.
\end{equation*}
Thus $\mathrm{span}\{1\}\perp H_0$, so the above decomposition is orthogonal and
\begin{equation*}
H=\mathrm{span}\{1\}\oplus H_0,
\qquad
H_0=\mathrm{span}\{1\}^\perp.
\end{equation*}
In a Hilbert space, the orthogonal projection onto a closed subspace is unique. Since $h_0\in H_0$ and $h-h_0=Ph\cdot 1\in \mathrm{span}\{1\}=H_0^\perp$, it follows that $h_0$ is exactly the orthogonal projection of $h$ onto $H_0$, i.e.
\begin{equation*}
\Pi_0 h = h-Ph\cdot 1.
\end{equation*}
This proves part 1.

\medskip
\noindent
\textbf{2.}
Fix $f\in\mathcal F$. By definition of $\mathbb G_P$ and linearity of $\mathbb W_P$ on $H$,
\begin{equation*}
\mathbb G_P(f)=\mathbb W_P(f)-Pf\cdot \mathbb W_P(1)=\mathbb W_P(f-Pf\cdot 1)=\mathbb W_P(\Pi_0 f).
\end{equation*}
Hence $\mathbb G_P$ is a linear image of the centered Gaussian family $\mathbb W_P$, so $\mathbb G_P$ is centered Gaussian on $\mathcal F$.

Now take $f,g\in\mathcal F$. Using isonormality and Step 1,
\begin{equation*}
\mathrm{Cov}\big(\mathbb G_P(f),\mathbb G_P(g)\big)
=\mathbb E\big[\mathbb W_P(f-Pf)\,\mathbb W_P(g-Pg)\big]
=\langle f-Pf,\ g-Pg\rangle.
\end{equation*}
Expand the inner product:
\begin{equation*}
\langle f-Pf,\ g-Pg\rangle
=P\big((f-Pf)(g-Pg)\big)
=P(fg)-Pf\,Pg.
\end{equation*}
This proves the covariance characterization in part 2, as well as the identity $\mathbb G_P(f)=\mathbb W_P(\Pi_0 f)=\mathbb W_P(f-Pf)$.

\medskip
\noindent
\textbf{3.}
Let $h,g\in H_0$. Since $H_0\subset H$ and $\mathbb W_P$ is isonormal on $H$,
\begin{equation*}
\mathbb E[\mathbb W_P(h)\mathbb W_P(g)]=\langle h,g\rangle.
\end{equation*}
Therefore $\mathbb W_P|_{H_0}$ is isonormal on $H_0$. In particular, for any $h\in H_0$,
\begin{equation*}
\|\mathbb W_P(h)\|_{L^2(\Omega)}^2=\mathbb E[\mathbb W_P(h)^2]=\langle h,h\rangle=\|h\|_{L^2(P)}^2,
\end{equation*}
so $\mathbb W_P|_{H_0} : H_0 \hookrightarrow L^2(\Omega)$.

Finally, for each $f\in\mathcal F$, the centered function $f-Pf$ belongs to $H_0$ by definition. Step 2 shows
\begin{equation*}
\mathbb G_P(f)=\mathbb W_P(f-Pf),
\end{equation*}
so $\mathbb G_P$ is exactly the image under the isonormal embedding $\mathbb W_P|_{H_0}$ of the centered class $\{f-Pf:f\in\mathcal F\}\subset H_0$. This proves 3.
\end{proof}  

\begin{question}
Let $(D,d)$ be a metric space. Define
\begin{equation*}
BL_1 := \{ f : D \to \mathbb{R} \; : \; \|f\|_\infty \le 1, \ \mathrm{Lip}(f) \le 1 \}.
\end{equation*}
The bounded Lipschitz distance between probability measures $P,Q$ on $D$ is
\begin{equation*}
d_{BL}(P,Q) := \sup_{f \in BL_1} \left| \int f \, dP - \int f \, dQ \right|.
\end{equation*}
Show that for all $x,y \in D$,
\begin{equation*}
d_{BL}(\delta_x,\delta_y) = d(x,y) \wedge 2.
\end{equation*}
\end{question}

\begin{proof}
For Dirac measures we have
\begin{equation*}
d_{BL}(\delta_x,\delta_y)
=
\sup_{f \in BL_1} |f(x) - f(y)|.
\end{equation*}

\medskip
\noindent
\textbf{Upper bound.}
Let $f \in BL_1$. Since $\mathrm{Lip}(f) \le 1$,
\begin{equation*}
|f(x)-f(y)| \le d(x,y).
\end{equation*}
Since $\|f\|_\infty \le 1$,
\begin{equation*}
|f(x)-f(y)| \le |f(x)| + |f(y)| \le 2.
\end{equation*}
Therefore,
\begin{equation*}
|f(x)-f(y)| \le d(x,y) \wedge 2.
\end{equation*}
Taking the supremum over $BL_1$ gives
\begin{equation*}
d_{BL}(\delta_x,\delta_y)
\le
d(x,y) \wedge 2.
\end{equation*}

\medskip
\noindent
\textbf{Lower bound.}
For any fixed $x\in D$, define
\begin{equation*}
f(z) := (d(x,z) - 1) \wedge 1.
\end{equation*}
First, we verify $f \in BL_1$.
Since $d(x,z) \ge 0$, we have $d(x,z)-1 \ge -1$, and truncation at $1$ yields
\begin{equation*}
-1 \le f(z) \le 1,
\end{equation*}
so $\|f\|_\infty \le 1$.

The map $z \mapsto d(x,z)$ is $1$-Lipschitz, and the function
$t \mapsto t \wedge 1$ is also $1$-Lipschitz on $\mathbb{R}$.
Hence $f$ is $1$-Lipschitz:
\begin{equation*}
|f(z) - f(w)| \le d(z,w).
\end{equation*}
Thus $f \in BL_1$.

Now compute:
\begin{equation*}
f(x) = (d(x,x) - 1) \wedge 1 = (-1) \wedge 1 = -1.
\end{equation*}
If $d(x,y) \le 2$, then $d(x,y)-1 \le 1$, so $f(y) = d(x,y) - 1,$
and therefore
\begin{equation*}
|f(x)-f(y)|
=
|-1 - (d(x,y)-1)|
=
d(x,y).
\end{equation*}
If $d(x,y) \ge 2$, then $d(x,y)-1 \ge 1$, so $f(y) = 1,$
and hence
\begin{equation*}
|f(x)-f(y)|
=
|-1 - 1|
=
2.
\end{equation*}
In conclusion,
\begin{equation*}
|f(x)-f(y)| = d(x,y) \wedge 2.
\end{equation*}
Thus
\begin{equation*}
d_{BL}(\delta_x,\delta_y)
\ge
d(x,y) \wedge 2.
\end{equation*}
\end{proof}

\begin{question}[Pinsker inequality]
Let $(\Omega,\mathcal F)$ be a measurable space and let $P,Q$ be
probability measures on $(\Omega,\mathcal F)$. Define the total
variation distance
\begin{equation*}
\|P-Q\|_{\mathrm{TV}}
=
\sup_{A\in\mathcal F}|P(A)-Q(A)|.
\end{equation*}
Then the Kullback--Leibler divergence (relative entropy) satisfies
\begin{equation*}
D(P\|Q)\ge 2\,\|P-Q\|_{\mathrm{TV}}^{\,2}.
\end{equation*}
\end{question}

\begin{proof}
If $P\not\ll Q$, then $D(P\|Q)=+\infty$ and the inequality holds
trivially. Hence assume $P\ll Q$ and let
\begin{equation*}
f=\frac{dP}{dQ}.
\end{equation*}
Then
\begin{equation*}
D(P\|Q)=\int f\log f\,dQ.
\end{equation*}

\medskip
\noindent
\textbf{Step 1: Monotonicity of relative entropy.}

Let $\mathcal G\subset\mathcal F$ be a sub-$\sigma$-algebra.
Denote by $P_{\mathcal G}$ and $Q_{\mathcal G}$ the restrictions of
$P$ and $Q$ to $\mathcal G$. Then
\begin{equation*}
\frac{dP_{\mathcal G}}{dQ_{\mathcal G}}
=
\e_Q[f\mid\mathcal G]
\qquad Q\text{-a.s.}
\end{equation*}
Hence
\begin{equation*}
D(P_{\mathcal G}\|Q_{\mathcal G})
=
\int \phi(\e_Q[f\mid\mathcal G])\,dQ,
\qquad
\phi(x)=x\log x.
\end{equation*}
Since $\phi$ is convex, Jensen's inequality gives
\begin{equation*}
\phi(\e_Q[f\mid\mathcal G])
\le
\e_Q[\phi(f)\mid\mathcal G].
\end{equation*}
Integrating both sides yields
\begin{equation*}
D(P_{\mathcal G}\|Q_{\mathcal G})
\le
D(P\|Q).
\end{equation*}

\medskip
\noindent
\textbf{Step 2: Hahn decomposition and reduction to a binary
$\sigma$-algebra.}

Let $\nu=P-Q$. By the Hahn decomposition theorem there exists a set
$A\in\mathcal F$ such that
\begin{equation*}
\|P-Q\|_{\mathrm{TV}}
=
P(A)-Q(A).
\end{equation*}
Define
\begin{equation*}
p=P(A), \qquad q=Q(A).
\end{equation*}
Let $\mathcal G=\sigma(A)=\{\varnothing,A,A^c,\Omega\}$. Then
$P_{\mathcal G}$ and $Q_{\mathcal G}$ correspond to Bernoulli
measures with parameters $p$ and $q$, and therefore
\begin{equation*}
D(P_{\mathcal G}\|Q_{\mathcal G})
=
p\log\frac{p}{q}
+
(1-p)\log\frac{1-p}{1-q}.
\end{equation*}
By Step 1,
\begin{equation*}
D(P\|Q)
\ge
D(P_{\mathcal G}\|Q_{\mathcal G}).
\end{equation*}

\medskip
\noindent
\textbf{Step 3: Binary inequality.}

Define
\begin{equation*}
F(p)
=
p\log\frac{p}{q}
+
(1-p)\log\frac{1-p}{1-q}.
\end{equation*}
A direct computation shows
\begin{equation*}
F'(p)
=
\log\frac{p}{q}
-
\log\frac{1-p}{1-q},
\qquad
F''(p)=\frac1p+\frac1{1-p}.
\end{equation*}
Since $p(1-p)\le 1/4$, we have
\begin{equation*}
F''(p)\ge 4.
\end{equation*}
Define
\begin{equation*}
G(p)=F(p)-2(p-q)^2.
\end{equation*}
Then
\begin{equation*}
G''(p)=F''(p)-4\ge 0,
\end{equation*}
so $G$ is convex. Moreover
\begin{equation*}
G(q)=0,
\qquad
G'(q)=0.
\end{equation*}
Hence $G(p)\ge 0$ for all $p$, implying
\begin{equation*}
F(p)\ge 2(p-q)^2.
\end{equation*}
\end{proof}


  \section{Stochastic Processes}

  \begin{question}
   Suppose $(X_n)_{n\geq 0}$ is a Markov chain. If the state space $E$ has at least two states, then the following statement is not true:
   \begin{equation*}
    \p(X_{n+1}=j|X_n\in A,X_{n-1}=i_{n-1},\cdots,X_0=i_0 )=\p(X_{n+1}=j|X_n\in A),
   \end{equation*}
   where $A\subset E$. Use an example to show this.
  \end{question}
  \begin{proof}
   Consider the $(p,q)$-biased random walk $\{X_n\}$, where $0<q<p<1$ and the initial distribution is $\p(X_0=0)=\p(X_0=2)=1/2$. Let $A=\{1,3\}$. Then we have
   \begin{equation*}
    \p(X_2=2|X_1\in A,X_0=0)=\p(X_2=2|X_1=1)=p,
   \end{equation*}
   and
   \begin{equation*}
    \begin{split}
     \p(X_1=1)&=\p_0(X_1=1)\p(X_0=0)+\p_2(X_1=1)\p(X_0=2)\\
     &= p/2+q/2=1/2,\\
     \p(X_1=3)&=\p_0(X_1=3)\p(X_0=0)+\p_2(X_1=3)\p(X_0=2)= p/2.\\
    \end{split}
   \end{equation*}
   So
   \begin{equation*}
    \begin{split}
     &\p(X_2=2|X_1\in A)=\frac{\p\[\{X_2=2\}\cap (\{X_1=1\}+\{X_1=3\} )\]}{\p(\{X_1=1\}+\{X_1=3\} )}\\
     =&\frac{\p(X_2=2,X_1=1)+\p(X_2=2,X_1=3)}{\p(X_1=1)+\p(X_1=3)}\\
     =&\frac{\p(X_2=2|X_1=1)\p(X_1=1)+\p(X_2=2|X_1=3)\p(X_1=3)}{\p(X_1=1)+\p(X_1=3)}\\
     =&\frac{p/2+qp/2}{1/2+p/2}=\frac{p(1+q)}{1+p}.
    \end{split}
   \end{equation*}
   Hence
   \begin{equation*}
    \p(X_2=2|X_1\in A,X_0=0)\neq \p(X_2=2|X_1\in A).
   \end{equation*}
  \end{proof}

  \begin{question}
   Use an example to show exists an adapted stochastic process $(X_t:t\in I)$ with respect to $(\f_t)$ which is not measurable. Namely, $I\times\Omega\to E,(t,\omega)\mapsto X_t(\omega )$ is not measurable.
  \end{question}
  \begin{proof}
   Let $I=\r_+$. Choose an unmeasurable subset $A\subset [0,1]$ and a measurable set $F\in\f$. Consider the following stochastic process
   \begin{equation*}
    (\r_+\times\Omega,\b(\r_+)\times\f)\to (\r,\b(\r)),\quad(t,\omega )\mapsto \i_{A\times F}(t,\omega)=:X_t(\omega )
   \end{equation*}
    and the filtration $(\f_t)=(\f)$. For any $t\geq 0$, if $t\in A$, then $X_t(\omega)=\i_F(\omega)\in\f$. If $t\notin A$, then $X_t(\omega )=0\in\f$. So $X$ is an adapted stochastic process. However, $\{(t,\omega ):X_t(\omega )=1\}=A\times F\notin \b(\r_+)\times\f$. So $X$ is not measurable.
  \end{proof}

  \begin{question}
   Use an example to show exists a measurable stochastic process $(X_t:t\in I)$ which is not adapted to $(\f_t)$.
  \end{question}
  \begin{proof}
   Let $I=\z_+$. Choose a measurable set $F\in\f$ ($F\neq\Omega,\emptyset$). Consider the following stochastic process
   \begin{equation*}
    (\z_+\times\Omega,\b(\z_+)\times\f)\to (\r,\b(\r)),\quad(t,\omega )\mapsto \i_{\{0\}\times F}(t,\omega)=:X_t(\omega )
   \end{equation*}
   and the filtration $(\f_t)=(\{\Omega,\emptyset\} )$. Since
   \begin{equation*}
    \sigma((t,\omega )\mapsto X_t(\omega ))=\{\z_+\times\Omega,\emptyset, \{0\}\times F\}\subset \b(\z_+)\times\f.
   \end{equation*}
   Thus $X$ is measurable. However for $t=0$, we know $X_0(\omega )=\i_F(\omega )$ and
   \begin{equation*}
    \sigma(\omega\mapsto X_0(\omega ))=\{\Omega,\emptyset,F \}\not\subset \f_0=\{\Omega,\emptyset\}.
   \end{equation*}
   So $X$ is not adapted to $(\f_t)$.
  \end{proof}

  \begin{question}
   Use an example to show exists an adapted stochastic process $(X_t:t\in I)$ with respect to $(\f_t)$ which is not progressively measurable.
  \end{question}
  \begin{proof}
   We still use the example in {\bfseries QUESTION 4.2}. We have proved $X$ is adapted to $(\f_t)=(\f)$. However, for any $t>1$, we have
   \begin{equation*}
    \{(t,\omega ):X_t(\omega )=1\}\cap([0,t]\times\Omega)  =A\times F\notin \b[0,t]\times\f.
   \end{equation*}
   So we have finished the proof.
  \end{proof}

  \begin{question}
   Use an example to show exists a measurable stochastic process $(X_t:t\in I)$ which is not progressively measurable.
  \end{question}
  \begin{proof}
   We still use the example in {\bfseries QUESTION 4.3}. We have proved $X$ is measurable. However for $t=0$, we know
   \begin{equation*}
    \begin{split}
     &\sigma((t,\omega )\mapsto X_t(\omega ))\cap (\{0\}\times \Omega)\\
     =&\{\{0\}\times\Omega,\emptyset, \{0\}\times F\}\not\subset \b(\{0\})\times\f_0.
    \end{split}
   \end{equation*}
   So $X$ is not progressively measurable.
  \end{proof}

  \begin{question}
   Use an example to show exists a measurable and adapted stochastic process $(X_t:t\in I)$ with respect to $(\f_t)$ which is not progressively measurable.
  \end{question}
  \begin{proof}
   \textcolor{red}{I still can not find a proper example!!!}
  \end{proof}

  \begin{question}
   Suppose $(\f_t)$ is a filtration. Prove $(\f_{t+})$ is right continuous.
  \end{question}
  \begin{proof}
   Clearly, we have $\f_{t+}\subset \f_{t++}$. On the contrary, for all $A\in\f_{t++}=\cap_{\ell>t}\f_{\ell+}= \cap_{\ell>t}\cap_{s>\ell}\f_s$, then $\forall  s,\ell$ satisfying $s>\ell>t$, we have $A\in \f_s$. Hence for each $\gamma>t$, exists $\ell_\gamma$, such that $\gamma>\ell_\gamma>t$. So $A\in\f_\gamma$. By the arbitrariness of $\gamma$, we obtain $A\in \cap_{\gamma>t}\f_\gamma=\f_{t+}$. In conclusion, $\f_{t++}\subset\f_{t+}$.
  \end{proof}

  \begin{question}
   Suppose $(E,\mathscr{E})$ is measurable space and isomorphic to a complete separable metric space $(F,\b(F))$. Given a Poisson random measure $X:(\Omega,\f,\p )\to (E,\mathscr{E})$ with intensity $\mu$. Then for each measurable functions $\{f_i\}_{i=1}^n$ and disjoint sets $\{A_i\}\subset\mathscr{E}$, if $\{\mu(f_i)\}_{i=1}^n$ all make sense, then $X(f_1\i_{A_1}),\cdots X(f_n\i_{A_n}) $ is independent.
  \end{question}
  \begin{proof}
   When $f_i=\sum_{j=1}^{n_i} \alpha_{ij} \i_{U_{ij}}$, where $\{U_{ij}\}_{j=1}^{n_i}\subset\mathscr{E}$ is disjoint for each $i=1,\cdots,n$, then we have
   \begin{equation*}
    X(f_i;A_i)=\int_{A_i}\sum_{j=1}^{n_i} \alpha_{ij}\i_{U_{ij}}(x)X(dx)=\sum_{j=1}^{n_i}\alpha_{ij}X(A_i\cap U_{ij}).
   \end{equation*}
   By the definition of Poisson random measure, we know
   \begin{equation*}
    \{X(A_i\cap U_{ij} ):\forall i=1,\cdots,n,j=1,\cdots,n_i\}
   \end{equation*}
   are independent. So the random vectors
   \begin{equation*}
    \{(X(A_i\cap U_{i1} ),\cdots,X(A_i\cap U_{in_i} )): \forall i=1,\cdots,n\}
   \end{equation*}
   are independent. Furthermore, we know
   \begin{equation*}
    \left\{X(f_i;A_i)=\sum_{j=1}^{n_i}\alpha_{ij}X(A_i\cap U_{ij}):i=1,\cdots,n \right\}
   \end{equation*}
   are independent. When $\{f_i\}_{i=1}^n\subset\mathrm{p}\mathscr{E}$, we can find simple functions $\{g_{ij}\}$ such that $g_{ij}\uparrow f,j\to\infty$. Then for $\{\theta_i\}_{i=1}^n>0$ and $A\in\mathscr{E}$, by {\sffamily Monotonicity convergence theorem}, we know $X(f_{ij}\i_{A_i})\uparrow X(f_i\i_{A_i})$. Then by {\sffamily Dominance convergence theorem}, we have
   \begin{equation*}
    \begin{split}
     &\p\exp \left\{- \sum_{i=1}^n\theta_iX(f_i\i_{A_i})\right\}=\lim_{j\to\infty}\p\exp \left\{- \sum_{i=1}^n\theta_iX(f_{ij}\i_{A_i})\right\}\\
     &\qquad\qquad\qquad = \lim_{j\to\infty}\p\(\prod_{i=1}^n \exp \left\{- \theta_iX(f_{ij}\i_{A_i})\right\}\)\\
     &\qquad\qquad\qquad =  \lim_{j\to\infty}\prod_{i=1}^n\p\exp \left\{- \theta_iX(f_{ij}\i_{A_i})\right\}\\
     &\qquad\qquad\qquad =\prod_{i=1}^n\p\exp \left\{- \theta_iX(f_{i}\i_{A_i})\right\}
    \end{split}
   \end{equation*}
   By the property of generating function, we know $\{X(f_i\i_{A_i})\}_{i=1}^n$ are independent. When $\{f_i\}_{i=1}^n\in\mathscr{E}$, then
   \begin{equation*}
    \begin{split}
     X(f_i\i_{A_i})&=X(f_i^+\i_{A_i})-X(f_i^-\i_{A_i})\\
     &=X(f\i_{\{f_i>0\}\cap A_i})-X((-f)\i_{\{f_i<0\}\cap A_i}).
    \end{split}
   \end{equation*}
   Utilize what we have proved, we know
   \begin{equation*}
    \{X(f_i^\pm \i_{A_i}):i=1,\cdots,n \}
   \end{equation*}
   are independent. So the random vectors
   \begin{equation*}
    \{(X(f_i^+ \i_{A_i}),X(f_i^- \i_{A_i})):i=1,\cdots,n \}
   \end{equation*}
   are independent. Furthermore, we know
   \begin{equation*}
    \{X(f_i\i_{A_i}):i=1,\cdots,n\}
   \end{equation*}
   are independent.
  \end{proof}

  \begin{question}
   Suppose $T$ is a $(\f_t)_{t\in I}$ stopping time. Then $\{T=\infty\}\in\f_\infty$.
  \end{question}
  \begin{proof}
   Observe for each $n\geq 1,\{T>n\}=\{T\leq n\}^c\in\f_n\subset\f_{\infty}$, then
   \begin{equation*}
    \{T=\infty \}=\bigcap_{n=1}^\infty \{T>n\}\in\f_\infty.
   \end{equation*}
   The proof is completed.
  \end{proof}

  \begin{question}
   Suppose $(\g_t)_{t\in I}$ is the enlargement of the filtration $(\f_t)_{t\in I}$. Then $\g_\infty=\sigma(\f_\infty\cup\mathscr{N} )$.
  \end{question}
  \begin{proof}
   Obviously $\f_\infty\subset  \g_\infty$ and $\mathscr{N}\subset\g_0\subset\g_\infty$, thus $\sigma(\f_\infty\cup\mathscr{N} )\subset\g_\infty$. On the contrary, for each $t\in I$, we have $\g_t=\sigma(\f_t\cup\mathscr{N} )\subset \sigma(\f_\infty\cup\mathscr{N} )$. Thus $\g_\infty=\sigma(\cup_{t\in I}\g_t)\subset \sigma(\f_\infty\cup\mathscr{N} )$.
  \end{proof}

  \begin{question}
   Suppose $(B_t)_{t\geq 0}$ is the Brownian motion on the probability space $(\Omega,\f,\p )$. Then $B_t/t\to 0,\p\mathrm{\text{-}a.s.}$. Therefore
   \begin{equation*}
    X_t=\begin{cases}
     B_t/t&,\mbox{if }t>0,\\
     0&,\mbox{if }t=0.
    \end{cases}
   \end{equation*}
   is a.s. continuous.
  \end{question}
  \begin{proof}
   Since $(B_t)_{t\geq 0}$ is a continuous time martingale and $B_t\sim N(0,t)$, then for each $\varepsilon>0$ and $0<\delta<\gamma $,
   \begin{equation*}
    \begin{split}
     \p\(\sup_{\delta\leq t\leq \gamma }|B_t/t|\geq \varepsilon \)&\leq \p\( \sup_{\delta\leq t\leq \gamma }|B_t|/\delta\geq \varepsilon \)\\
     &\leq \varepsilon^{-2}\p\(B_\gamma^2/\delta^2\)=\frac{\gamma }{\varepsilon^2\delta^2}.
    \end{split}
   \end{equation*}
   Especially, we take $\delta=2^n,\gamma=2^{n+1}$ for each $n\geq 0$. Then we have
   \begin{equation*}
    \sum_n\p\(\sup_{2^n\leq t\leq 2^{n+1} }|B_t/t|\geq \varepsilon \)\leq \sum_n \frac{1}{2^{n-1}\varepsilon^2}<\infty.
   \end{equation*}
   The {\sffamily Borel-Cantelli's lemma} shows
   \begin{equation*}
    \p\(\bigcap_{m=0}^\infty\bigcup_{n=m}^\infty \sup_{2^n\leq t\leq 2^{n+1} }|B_t/t|\geq \varepsilon\)=0,\quad\forall\varepsilon>0.
   \end{equation*}
   This yields $B_t/t\to 0,\p\mathrm{\text{-}a.s.}$.
  \end{proof}

  \begin{question}
   Suppose $(X_t)_{t\in I}$ is a Markov chain on the discrete state space $E$. For each $i\in E$, define $S_i^0=0$ and successively define $S_i^{n}=\inf\{t>S_{i}^{n-1}:X_t=i\}$ for $n\geq 1$. Write $T_i^n=S_i^n-S_i^{n-1}$ for $n\geq 1$. If $\p_i(T_i^1<\infty)=\p_i(S_i^1<\infty)=1$, then $\{T_i^n\}_{n\geq 1}$ is a i.i.d.\!\!\! random variables sequence under the probability measure $\p^i$ and $T_i^k\sim T_i^1$ for each $k\geq 1$.
  \end{question}
  \begin{proof}
   We use the notation of the Markov system. Then $S_i^{n+1}=S_i^n+S_i\circ\theta_{S_i^n}$ and $T_i^{n+1}=S_i\circ\theta_{S_i^n}$ for each $n\geq 0$. For each
  \end{proof}

  \begin{question}
   Suppose $(X_n)_{n\geq 0}$ is a Markov chain on the discrete state space $E$. For each $i\in E$, define
   \begin{equation*}
    \begin{split}
     N_n(i)&=\{\mathrm{the\ number\ of \ visits\ to\  }i \mathrm{\ at\ times}\leq n \},\\
     N(i)&=\{\mathrm{the\ number\ of \ visits\ to\  }i \}.
    \end{split}
   \end{equation*}
    Then $N(i)=\infty, \p_i\text{-}\ae$ iff $i$ is recurrent. Moreover, when $i\in E$ is recurrent, then $N_n(i)\to \infty(n\to\infty),\p_i\text{-}\ae$
  \end{question}
  \begin{proof}
    Define $S_i^0=0$ and successively define $S_i^{n}=\inf\{n>S_{i}^{n-1}:X_n=i\}$ for $n\geq 1$. By the {\sffamily strong Markov property}, we have
    \begin{equation*}
     \p_i(S_i^n<\infty)=[\p_i(S_i^1<\infty)]^n=(f_{ii})^{n},\quad n\geq 1
    \end{equation*}
    and
   \begin{equation*}
    \p_i[N(i)=\infty]=\lim_{n \to\infty}\p_i[N(i)\geq n ]=\lim_{n \to\infty} \p_i(S_i^n <\infty )=\lim_{n \to\infty} (f_{ii})^{n}.
   \end{equation*}
   Therefore $\p_i[N(i)=\infty]=1$ iff $f_{ii}=1$. Namely, $N(i)=\infty, \p_i\text{-}\ae$ iff $i$ is recurrent. Moreover, notice
   \begin{equation*}
    N_n(i)=\sum_{k=1}^n\mathbbm{1}_{\{X_k=i\}}.
   \end{equation*}
   When $i$ is recurrent, we have $N_n(i)\uparrow N(i)=\infty,\p_i\text{-}\ae$.
  \end{proof}

  \begin{question}
   Suppose $(N_t)$ is a Poisson process with parameter $\lambda>0$ on the probability space $(\Omega,\f,\p )$. Define $T_1=\inf\{t\geq 0:N_t=1 \}$, then we have
   \begin{equation*}
    \p(T_1\leq s|N_t=1)=s/t,\quad \forall 0\leq s\leq t.
   \end{equation*}
  \end{question}
  \begin{proof}
   Utilize the independence of $N_s$ and $N_t-N_s$, then we have
   \begin{equation*}
    \begin{split}
     \p(T_1\leq s|N_t=1)&=\p(N_s=1|N_t=1 )=\frac{\p(N_s=1,N_t=1 )}{\p(N_t=1)}\\
     &=\frac{\p(N_s=1,N_t-N_s=0 )}{\p(N_t=1)}\\
     &=\frac{\p(N_s=1)\p(N_t-N_s=0 )}{\p(N_t=1)}\\
     &=\frac{\lambda se^{-\lambda s} e^{-\lambda(t-s)}}{\lambda te^{-\lambda t}}=\frac{s}{t}.
    \end{split}
   \end{equation*}
   This is what we desire.
  \end{proof}

  \begin{question}
   Suppose $T$ is a $(\f_n)_{n\geq 0}$ stopping time on the filtered probability space $(\Omega,\f,(\f_n),\p )$. If the random variable $X\in \f$ is integrable, then
   \begin{equation*}
    \p(X|\f_T)=\sum_{n=0}^\infty\p(X|\f_n)\mathbbm{1}_{\{T=n\}}+\p(X|\f_\infty)\mathbbm{1}_{\{T=\infty\}}.
   \end{equation*}
  \end{question}
  \begin{proof}
   Write $F$ for the right side. For each $k\geq 0$, we have
   \begin{equation*}
    F\mathbbm{1}_{\{T\leq k\}}=\sum_{n=0}^k\p(X|\f_n)\mathbbm{1}_{\{T=n\}}\in\f_k,
   \end{equation*}
   so $F\in\f_T$. Moreover
   \begin{equation*}
    \begin{split}
     &\p\[\bigg|\sum_{n=0}^\infty\p(X|\f_n)\mathbbm{1}_{\{T=n\}}+\p(X|\f_\infty)\mathbbm{1}_{\{T=\infty\}}\bigg|\]\\
     &\qquad \leq\p\bigg[\sum_{n=0}^\infty\p(|X||\f_n)\mathbbm{1}_{\{T=n\}}+\p(|X||\f_\infty)\mathbbm{1}_{\{T=\infty\}}\bigg]\\
     &\qquad=\sum_{n=0}^\infty\p[\p(|X||\f_n)\mathbbm{1}_{\{T=n\}}]+\p[\p(|X||\f_\infty)\mathbbm{1}_{\{T=\infty\}}]\\
     &\qquad=\sum_{n=0}^\infty\p(|X|;T=n)+\p(|X|;T=\infty)=\p(|X|)<\infty.
    \end{split}
   \end{equation*}
   Thus $F$ is also integrable. For each $A\in\f_T$, by the {\sffamily Fubini's theorem},
   \begin{equation*}
    \begin{split}
     \p(F;A)&=\sum_{n=0}^\infty\p[\p(X|\f_n)\mathbbm{1}_{\{T=n\}\cap A}]+\p[\p(X|\f_\infty)\mathbbm{1}_{\{T=\infty\}\cap A}]\\
     &=\sum_{n=0}^\infty\p(X;\{T=n\}\cap A)+\p(X;\{T=\infty\}\cap A]\\
     &=\p(X;A).
    \end{split}
   \end{equation*}
   So $\p(X|\f_T)=F$.
  \end{proof}

  \begin{question}
   Let $M=(M_n)_{n\geq 0}$ be a $(\f_n)$ uniformly integrable martingale and let $M_\infty$ be the a.s. and $L^1$ limit of $M$, where $(\f_n)$ is reinforced. Then for any stopping time $T$, we have
   \begin{equation*}
    \p(M_\infty|\f_T )=M_T,
   \end{equation*}
   where with convention that $M_T=M_\infty$ on the event $\{T=\infty \}$. Therefore,
   \begin{enumerate}
    \item For each stopping times $S,T$ and $S\leq T$, we have
    \begin{equation*}
     \p(M_T|\f_S)=M_S.
    \end{equation*}
    \item For each stopping time $T$, we have $\p(|M_T|)<\infty$ and $\p(M_T)=\p(M_0)$.
   \end{enumerate}
  \end{question}
  \begin{proof}
   \textcolor{blue}{\sffamily Way 1.} Since $M$ is a uniformly integrable martingale, then exists $M_\infty$ such that $M_n\xlongrightarrow[\mathrm{a.s.}]{L^1} M_\infty$ and $\p(|M_\infty|)<\infty$. Moreover, for each $n\geq 0$,
   \begin{equation*}
    \p(M_\infty|\f_n )=M_n.
   \end{equation*}
   Furthermore, for the given $n\geq 0$, by the {\sffamily optional stopping theorem}, we have
   \begin{equation*}
    \p(M_n|\f_{T\wedge n})=M_{T\wedge n}.
   \end{equation*}
   Combine the above two conditions, note $\f_{T\wedge n}\subset\f_n$, we have
   \begin{equation*}
    \p(M_\infty|\f_{T\wedge n} )=\p[\p(M_\infty|\f_n )|\f_{T\wedge n } ]=\p(M_n|\f_{T\wedge n})=M_{T\wedge n}.\tag{$*$}
   \end{equation*}
   It follows from $(*)$ that
   \begin{equation*}
    |M_{T\wedge n}|\leq \p(|M_\infty ||\f_{T\wedge n}).
   \end{equation*}
   Since $\{T\leq n\}\in\f_{T\wedge n}$, then
   \begin{equation*}
    \begin{split}
     \p(|M_{T}|;&\{T\leq n\})=\p(|M_{T\wedge n}|;\{T\leq n\})\\
     &\leq \p[\p(|M_\infty ||\f_{T\wedge n});\{T\leq n\}]=\p(|M_\infty|;\{T\leq n\}).
    \end{split}
   \end{equation*}
   Let $n\to\infty$, the {\sffamily Monotone convergence theorem} yields
   \begin{equation*}
    \p(|M_{T}|;\{T<\infty \})\leq \p(|M_\infty|;\{T<\infty\}).
   \end{equation*}
   Moreover, observe $\p(|M_{T}|;\{T=\infty \})=\p(|M_{\infty}|;\{T=\infty \})$, then
   \begin{equation*}
    \begin{split}
     \p(|M_T|)&=\p(|M_{T}|;\{T<\infty \})+\p(|M_{T}|;\{T=\infty \})\\
     &\leq \p(|M_\infty|;\{T<\infty\})+\p(|M_{\infty}|;\{T=\infty \})\\
     &=\p(|M_\infty|)<\infty.
    \end{split}
   \end{equation*}
   Therefore $\p(|M_T|)<\infty$.

   If $A\in \f_T$, we claim $A\cap \{T\leq n\}\in \f_{T\wedge n}$. In fact, $A\cap \{T\leq n\}\in \f_n$ is obvious. Moreover $A\in\f_T$ and $\{T\leq n\}\in\f_T$ implies $A\cap \{T\leq n\}\in\f_T$. So $A\cap \{T\leq n\}\in \f_T\cap \f_n=\f_{T\wedge n}$. So $(*)$ implies
   \begin{equation*}
    \p(M_\infty;A\cap \{T\leq n\} )=\p(M_{T\wedge n};A\cap \{T\leq n\})=\p(M_T;A\cap \{T\leq n\}).
   \end{equation*}
   Let $n\to\infty$ and utilize the {\sffamily Dominance convergence theorem},
   \begin{equation*}
    \p(M_\infty;A\cap\{T<\infty\} )=\p(M_T;A\cap\{T<\infty\} ).
   \end{equation*}
   Observe $\p(M_\infty;A\cap\{T=\infty\} )=\p(M_T;A\cap\{T=\infty\} )$. then
   \begin{equation*}
    \p(M_\infty;A )=\p(M_T;A ), \quad \forall A\in\f_T.
   \end{equation*}
   Since $M_T\in\f_T$, then $\p(M_\infty|\f_T)=M_T$. Therefore
   \begin{enumerate}
    \item For each stopping times $S,T$ and $S\leq T$, observe $\f_S\subset\f_T$, then
    \begin{equation*}
     \p(M_T|\f_S)=\p[\p(M_\infty|\f_T)|\f_S]=\p(M_\infty|\f_S)=M_S.\tag{$**$}
    \end{equation*}
    \item For each stopping time $T$, we have shown $\p(|M_T|)<\infty$. Use $(**)$, we have
    \begin{equation*}
     \p(M_T)=\p[\p(M_T|\f_0)]=\p(M_0).
    \end{equation*}
   \end{enumerate}
   \textcolor{blue}{\sffamily Way 2.} We just prove $\p(M_\infty|\f_T)=M_T$ since the rest is similar. Observe $(\f_n)$ is reinforced, then $M_\infty\in\f_\infty$. Because $M_\infty$ is integrable, then the previous {\bfseries QUESTION} shows
   \begin{equation*}
    \begin{split}
     \p(M_\infty|\f_T )&=\sum_{n=0}^\infty \p(M_\infty|\f_n)\mathbbm{1}_{\{T=n\}}+\p(M_\infty|\f_\infty)\mathbbm{1}_{\{T=\infty\}}\\
     &=\sum_{n=0}^\infty M_n\mathbbm{1}_{\{T=n\}}+M_\infty\mathbbm{1}_{\{T=\infty\}}=M_T.    \end{split}
   \end{equation*}
   The result is completed.
  \end{proof}

  \begin{question}
   Suppose $X=(X_n)$ and $Y=(Y_n)$ are two irreducible and aperiodic Markov chains on the state space $E$ with common transition probability $P$ and initial distributions $\lambda$ and $\pi$, where $\pi$ is the stationary distribution of $P$. Consider the independent coupling $W_n=(X_n,Y_n)$, which has the transition probability $\tilde{P}$, where $\tilde{ p}_{(i,j)(\gamma ,\ell )}=p_{i\gamma}p_{j,\ell}$, and the initial distribution $\tilde{\mu}_i=\lambda_i \pi_i$. Fix $i\in E$ and define the following stopping time,
   \begin{equation*}
    T:=\inf\{n\geq 1:W_n=(i,i) \}.
   \end{equation*}
   The we define
   \begin{equation*}
    Z_n:=\begin{cases}
     (X_n,Y_n)&\mbox{if }n< T,\\
     (Y_n,X_n)&\mbox{if }n\geq  T.
    \end{cases}
   \end{equation*}
   Prove $Z=(Z_n)$ is a Markov chain with transition probability $\tilde{P}$ and the initial distribution $\tilde{\mu}$.
  \end{question}
  \begin{proof}
   For each $n\geq 1$ and $\{i_k\}_{k=1}^n,\{j_k\}_{k=1}^n\subset E$, without loss of generality, we assume exists $1\leq m\leq n$, such that
   \begin{equation*}
    \{Z_0=(i_0,j_0),\cdots,Z_n=(i_n,j_n)\}\subset\{T=m\}\subset  \{n\geq T\}.
   \end{equation*}
   Therefore, utilize the Markov property of $W_n=(X_n,Y_n)$,
   \begin{equation*}
    \begin{split}
     &\p[Z_0=(i_0,j_0),\cdots,Z_T=(i_T,j_T),\cdots, Z_n=(i_n,j_n) ]\\
     =&\p(X_0=i_0,Y_0=j_0,\cdots, Y_T=i_T,X_T=j_T,\cdots, Y_n=i_n,X_n=j_n )\\
     =&\p(X_0=i_0,Y_0=j_0,\cdots, Y_m=i,X_m=i,\cdots, Y_n=i_n,X_n=j_n )\\
     =&\p[W_0=(i_0,j_0),\cdots,W_m=(i,i),\cdots, W_{n}=(j_{n},i_{n}) ]\\
     =&\tilde{\mu}_{(i_0,j_0)}\tilde{p}_{(i_0,j_0)(i_1,j_1)}\cdots \tilde{p}_{(i_{m-1},j_{m-1})(i,i)}\tilde{p}_{(i,i)(j_{m+1},i_{m+1})}\cdots\tilde{p}_{(j_{n-1},i_{n-1})(j_{n},i_{n})} \\
     =&\lambda_{i_0}\pi_{i_0}p_{i_0i_1}p_{j_0j_1}\cdots p_{i_{m-1}i}p_{j_{m-1}i}p_{ij_{m+1}}p_{ii_{m+1}}\cdots p_{j_{n-1}j_n}p_{i_{n-1}i_n}\\
     =&\lambda_{i_0}\pi_{i_0}p_{i_0i_1}p_{j_0j_1}\cdots p_{i_{m-1}i}p_{j_{m-1}i}p_{ii_{m+1}}p_{ij_{m+1}}\cdots p_{i_{n-1}i_n}p_{j_{n-1}j_n}\\
     =&\tilde{\mu}_{(i_0,j_0)}\tilde{p}_{(i_0,j_0)(i_1,j_1)}\cdots \tilde{p}_{(i_{m-1},j_{m-1})(i,i)}\tilde{p}_{(i,i)(i_{m+1},j_{m+1})}\cdots\tilde{p}_{(i_{n-1},j_{n-1})(i_{n},j_{n})}.
    \end{split}
   \end{equation*}
   This shows $Z=(Z_n)$ is a Markov chain with transition probability $\tilde{P}$ and the initial distribution $\tilde{\mu}$.
  \end{proof}

  \begin{question}
   Suppose $\{S_n\}_{n\geq 1}$ is a sequence of i.i.d. exponential distribution with parameter $\lambda$. Fix $i\geq 1$, define $J_n:=\sum_{k=1}^nS_k$ and
   \begin{equation*}
    \tilde{S}_1=S_{i+1}-(s-J_i),\qquad \tilde{S}_n=S_{i+n}\quad \forall n\geq 2.
   \end{equation*}
   Then conditional on the event $J_i\leq s<J_{i+1}$, we have $\{\tilde{S_n}\}_{n\geq 1}$ are i.i.d. exponential distribution with parameter $\lambda$.
  \end{question}
  \begin{proof}
   Given $k\geq 1$, for each $\{t_\ell\}_{\ell=1}^k\subset\r_+$, utilize the independence between $\{S_n\}_{n\geq 1}$ and {\bfseries QUESTION 4.31}, we have
   \begin{equation*}
    \begin{split}
     \p(\tilde{S}_1> &t_1,\cdots,\tilde{S}_k> t_n|J_i\leq s<J_{i+1})\\
     &=\p[S_{i+1}-(s-J_i)> t_1,\cdots,S_{i+k}> t_k|s-J_i\geq 0,S_{i+1}>s-J_i   ]\\
     &=\frac{\p[S_{i+1}-(s-J_i)> t_1,\cdots,S_{i+k}> t_k,s-J_i\geq 0,S_{i+1}>s-J_i]}{\p(s-J_i\geq 0,S_{i+1}>s-J_i )}\\
     &=\frac{\p[S_{i+1}-(s-J_i)> t_1,s-J_i\geq 0,S_{i+1}>s-J_i]}{\p(s-J_i\geq 0,S_{i+1}>s-J_i )}\\
     &\quad \times \p(S_{i+2}>t_2,\cdots,S_{i+k}> t )\\
     &=\p[S_{i+1}> t_1+(s-J_i)|s-J_i\geq 0,S_{i+1}>s-J_i]\times\prod_{\ell=2}^k\p(S_{i+\ell}>t_\ell)\\
     &=\prod_{\ell=1}^k\p(S_{i+\ell}>t_\ell)=e^{-\lambda\sum_{\ell=1}^k t_\ell}.
    \end{split}
   \end{equation*}
   So the desired result follows.
  \end{proof}

  \begin{question}
   Suppose $(X_n)$ is an irreducible Markov chain with transition matrix $P$ and state space $E$. If $\mu$ is a non-trivial invariant measure of $P$, then $\mu$ is strictly positive.
  \end{question}
  \begin{proof}
   Since $\mu$ is a non-trivial invariant measure of $P$, then exists $i\in E$, such that $\mu_i>0$. Observe $P$ is irreducible, then for any $j\in E$, exists $n\geq 1$ satisfying $p_{ij}^{(n)}>0$. Then utilize $\mu=\mu P^{(n)}$, we have
   \begin{equation*}
    \mu_j=\sum_{k\in E}\mu_k p_{kj}^{(n)}\geq \mu_i p_{ij}^{(n)}>0.
   \end{equation*}
   So $\mu$ is strictly positive.
  \end{proof}

  \begin{question}
    Suppose $(X_t)_{t\geq 0}$ is a stochastic process in the filtered probability space $(\Omega,\f,(\f_t)_{t\geq 0},\p )$ and valued on a metric space $(E,d)$. Define $T=\inf\{t\geq 0:X_t\in A\}$, where $A\in \b(E)$ ensuring $T$ is a $(\f_t)_{t\geq 0}$ optional time. 
    \begin{enumerate}
    	\item If $(X_t)_{t\geq 0}$ has continuous-path, then $X_T\in\partial A$.
    	\item If $(X_t)_{t\geq 0}\in D_E[0,\infty)$, then $X_T\notin \mathrm{Ext}A$. Especially, when $A$ is a skeleton, say $A=\{x\}$, then $X_T=x$.
    	\item Assume $(X_t)_{t\geq 0}\in D_E[0,\infty)$ and $A$ is a open set. If $X_s\in A$, then exists $\delta>0$ satisfying $X_t\in A,\forall t\in[s,s+\delta)$. Especially, if $A=\{x\}$ is an open set, then exists $\delta>0$ satisfying $X_t=x,\forall t\in[s,s+\delta)$
    \end{enumerate}
  \end{question}
  \begin{proof}
    \begin{enumerate}
    	\item Firstly, if $X_T\in \mathring{A}$, notice $\mathring{A}$ is open and $(X_t)_{t\geq 0}$ is left continuous, then exists $0\leq s<T$ such that $X_s\in \mathring{A}$, which contradicts to definition of $T$. Besides, if $X_T\in \text{Ext}A$, notice $\text{Ext}{A}$ is open and $(X_t)_{t\geq 0}$ is right continuous, then exists $\delta>0$ such that $X_s\in \text{Ext}{A}, \forall T\leq s<T+\delta$. Namely, $X_s\notin A$ for any $T\leq s<T+\delta$,  which also contradicts to definition of $T$.
    	\item If $(X_t)_{t\geq 0}\in D_E[0,\infty)$, the similar argument concerning the exterior of $A$ in the first question gives the desired results. Finally, since when $A=\{x\}$, we have $A=\mathring{A}\cup\partial A$. Then $X_T\notin \mathrm{Ext}A\Longrightarrow X_T=x$.
    	\item The last assertion is just the definition.
    \end{enumerate}
  \end{proof}
  
  \begin{question}
  	Discuss the the relationship between the three statements.
  	\begin{enumerate}
  		\item For any $i,j\in E$, exist $n\in\n_+$ and distinctive states $i=i_1,\cdots,i_n=j$ such that $q_{i_1i_2}\cdots q_{i_{n-1}i_n}>0$.
  		\item For any $i,j\in E$, exist $n\in\n_+$ and states $i=i_1,\cdots,i_n=j$ such that $q_{i_1i_2}\cdots q_{i_{n-1}i_n}\neq 0$.
  		\item For any $i,j\in E$, exist $n\in\n_+$ such that $q_{ij}^{(n)}\neq 0$.
  	\end{enumerate}
  \end{question}
  \begin{proof}
  	\begin{enumerate}
  		\item $1\Longleftrightarrow 2$. First, $1\Longrightarrow 2$ is obvious. 
  		
  		$2\Longrightarrow 1$. If $i_1,\cdots, i_{n-1}$ are distinct, then all are done and it is clear that $q_{i_1i_2}\cdots q_{i_{n-1}i_n}>0$. Otherwise, presume $i_{n_1}<\cdots< i_{n_{m_1}}$ are the same, $i_{n_2}<\cdots< i_{n_{m_2}}$ are the same, $\cdots$, and $i_{n_3}<\cdots< i_{n_{m_3}}$ are the same. We can cut all the loops, then the rest part satisfies our requirement.
  		\item $3\Longrightarrow 2$. Observe we have the following expression:
  		\begin{equation*}
  			\sum_{i_1,\cdot,i_{n-1}\in E}q_{ii_1}q_{i_1i_2}\cdots q_{i_{n-1}i_n}\neq 0,
  		\end{equation*}
  		which implies exist states $i=i_1,\cdots,i_n=j$ such that $q_{i_1i_2}\cdots q_{i_{n-1}i_n}\neq 0$.
  		\item However $2\not\Longrightarrow 3$. Consider the following example.
  		\begin{equation*}
  			Q=\begin{pmatrix}
  				-3 & 2 & 1\\
  				2 & -3 & 1\\
  				3 & 12 & -15
  			\end{pmatrix}.
  		\end{equation*}
  		It is trivial to text $Q$ is indeed a $Q$-matrix. We have
  		\begin{equation*}
  			q_{12}^{(2)}=q_{11}q_{12}+q_{12}q_{22}+q_{13}q_{32}=-6-6+12=0.
  		\end{equation*}
  		However, clearly $q_{11}q_{12}\neq 0$, $q_{12}q_{22}\neq 0$ and $q_{13}q_{32}\neq 0$.
  	\end{enumerate}
  \end{proof}
  
  \begin{question}
  	Suppose $P$ satisfyies $\exists i_0\in E$ such that $\sum_{j\in E}p_{i_0j}<1$. Then
  	\begin{enumerate}
  		\item For all $n\in\n$, we have $\sum_{j\in E}p_{i_0j}^{(n)}<1$.
  		\item If exists $m\in\n$ and $i_1\in E$ such that $p_{i_1i_0}^{(m)}>0$, then for any $n>m$, we have $\sum_{j\in E}p_{i_1j}^{(n)}<1$.
  	\end{enumerate}
  \end{question}
  \begin{proof}
  	\begin{enumerate}
  		\item For any  $n\geq 2$, we have
  		\begin{equation*}
  			\begin{split}
  				\sum_{j\in E}p_{i_0j}^{(n)}&=\sum_{j\in E}\sum_{k\in E}p_{i_0k}p_{k j}^{(n-1)}=\sum_{k\in E}p_{i_0k}\sum_{j\in E}p_{k j}^{(n-1)}\\
  				&\leq \sum_{k\in E}p_{i_0k}<1.
  			\end{split}
  		\end{equation*}
  		\item For any $n>m$, by the previous question, we have $\sum_{j\in E}p_{i_0 j}^{(n-m)}<1$. Moreover, notice $p_{i_1i_0}^{(m)}>0$, then 
  		\begin{equation*}
  			\begin{split}
  				\sum_{j\in E}p_{i_1j}^{(n)}&=\sum_{j\in E}\sum_{k\in E}p_{i_1k}^{(m)}p_{k j}^{(n-m)}=\sum_{j\in E}\sum_{k\neq i_0}p_{i_1k}^{(m)}p_{k j}^{(n-m)}+\sum_{j\in E}p_{i_1i_0}^{(m)}p_{i_0 j}^{(n-m)}\\
  				&<\sum_{k\neq i_0}p_{i_0k}^{(m)}+p_{i_0i_0}^{(m)}\leq 1.
  			\end{split}
  		\end{equation*}
  	\end{enumerate}
  \end{proof}
  
  \begin{theorem}[Perron-Frobenius theorem of $Q$]
    Presume $Q$ is irreducible. Then
    \begin{enumerate}
      \item $Q$ has a principal eigenvalue $-\lambda\leq 0$ and two eigenvectors $u$ and $v$ satisfying
          \begin{equation*}
            uQ=-\lambda u,\quad Qv=-\lambda v.
          \end{equation*}
      \item If we restrict $\sum_{i\in E}u_i=1$ and $\sum_{i\in E}u_iv_i=1$, then $u$ and $v$ is unique.
      \item For any other eigenvalue $\lambda'$ of $Q$, we have $|\lambda' |>\lambda$.
      \item Let $P(t)=e^{tQ}$, then
      \begin{equation*}
        p_{ij}(t)\sim e^{-\lambda t}v_iu_j,\quad t\to\infty.
      \end{equation*}
    \end{enumerate}
  \end{theorem}
  \begin{proof}
    Let $c=\max_{i\in E}q_i$ and consider $Q^c=:Q+cI\geq 0$. The P-F theorem of nonnegative irreducible matrix yields
    \begin{enumerate}
     \item $Q^c$ has a principal eigenvalue $\lambda^c> 0$, a positive left eigenvector $u$ and a positive right eigenvector $v$ satisfying
     \begin{equation*}
      uQ^c=\lambda^cu,\quad Q^cv=\lambda^cv.\tag{$*$}
     \end{equation*}
     \item If we restrict $\sum_{i\in E}u_i=1,\sum_{i\in E}u_iv_i=1$, the eigenvectors $u,v$ are unique.
     \item For any other eigenvalues $\lambda^c_0$ of $Q^c$, we have $|\lambda^c_0 |< \lambda^c$.
    \end{enumerate}
    From $(*)$ we have
    \begin{equation*}
     \begin{split}
      u\(Q+cI\)&=\lambda^cu\Longrightarrow uQ=-(c-\lambda^c)u, \\
      \(Q+cI\)v^c&=\lambda^cv\Longrightarrow Qv=-(c-\lambda^c)v. \\
     \end{split}
    \end{equation*}
    Define $\lambda=:c-\lambda^c$. Hence
    \begin{equation*}
     uQ=-\lambda u,\quad Qv=-\lambda v,
    \end{equation*}
    and $u,v$ are unique with the restriction $\sum_{i\in E}u_i=1$ and $\sum_{i\in E}u_iv_i=1$.
    Observe $uQ^c=\lambda^cu$ and $\sum_{j\in E}q^c_{ij}\leq c$, then
    \begin{equation*}
     \sum_{i\in E}u_{i}q^c_{ij}=\lambda^cu_j\Rightarrow\sum_{j\in E}\sum_{i\in E}u_{i}q^c_{ij}=\lambda^c\sum_{j\in E}u_j\Rightarrow c\geq\lambda^c.
    \end{equation*}
    Therefore $\lambda>0$. Note $|\lambda_0^c|<\lambda^c$, then
    \begin{equation*}
     \lambda= c-\lambda^c<c-|\lambda_0^c |\leq |\lambda_0^c-c |.\tag{$**$}
    \end{equation*}
    Since there is a one-to-one correspondence
    \begin{equation*}
     \varphi: \sigma\(Q^c\)\to\sigma\(Q\): \mu\mapsto\mu-c.\tag{$.$}
    \end{equation*}
    Hence $(**)$ and $(.)$ means for any other eigenvalue $\lambda'$ of $Q$, we have $|\lambda' |>\lambda$. Furthermore, for any {\bfseries fixed} $t>0$,
     \begin{equation*}
      \begin{split}
       uP(t)&=u\sum_{n=0}^\infty\frac{Q^{n}t^n}{n!}=\sum_{n=0}^\infty\frac{uQ^{n}t^n}{n!}\\
       &=\sum_{n=0}^\infty\frac{(-\lambda)^n ut^n}{n!}=\(\sum_{n=0}^\infty\frac{(-\lambda)^n  t^n}{n!}\)u=\mathrm{e}^{-\lambda t}u.
      \end{split}
     \end{equation*}
     And symmetrically $P(t)v=\mathrm{e}^{-\lambda t}v$.
     Besides, we can also construct the following one-to-one correspondence
    \begin{equation*}
     \phi: \sigma\(Q\)\to\sigma\(P(t)\): \mu\mapsto\mathrm{e}^{\mu t }.
    \end{equation*}
    Consider $\phi\circ \varphi$, then for any other eigenvalue of $P(t)$, we can write it as  $e^{(\lambda_0^c-c) t}$, where $\lambda_0^c\in\sigma(Q^c)$ and $\lambda_0^c\neq \lambda^c$. Thus
     \begin{equation*}
      |e^{(\lambda_0^c-c) t}|=\mathrm{e}^{(\text{Re}\lambda_0^c-c) t}\leq e^{(|\lambda_0^c|-c) t} <e^{(\lambda^c-c) t}=\mathrm{e}^{-\lambda t}.\tag{..}
     \end{equation*}
     Formula $(..)$ mean $\mathrm{e}^{-\lambda t}$ is the principal eigenvalue of $P(t)$. In conclusion, $P(t)$ has principal eigenpairs $(\mathrm{e}^{-\lambda t},u,v)$.

     For any {\bfseries fixed} $t>0$, the P-F theorem and C-K equation yield
     \begin{equation*}
      p_{ij}(nt)=p_{ij}^{(n)}(t)\sim \mathrm{e}^{-\lambda nt}v_iu_j,\quad n\to\infty.\tag{...}
     \end{equation*}
     Finally, we just need to show $p_{ij}(t)\sim\mathrm{e}^{-\lambda t}v_iu_j(t\to\infty )$. Or equivalently, 
     we just need to prove $\forall i,j\in E$,
     \begin{equation*}
       e^{\lambda t}p_{ij}(t)\to v_iu_j, \quad t\to\infty.
     \end{equation*}
     Define $\varphi_{ij}(t)=e^{\lambda t}p_{ij}(t),\forall t>0$. Then the following properties naturally 
     follow by $(...)$ and $P(t)v=e^{-\lambda t}v$:
     \begin{equation*}
       \forall h>0, \varphi_{ij}(nh)\to v_iu_j(n\to\infty),\qquad \forall t>0, \sum_{j\in E}\varphi_{ij}(t)v_j=v_i.
     \end{equation*}
     Consider the following decomposition:
     \begin{equation*}
       \varphi_{ij}(t)=\sum_{k\in E}\varphi_{ik}(t-nh)\varphi_{kj}(nh),
     \end{equation*}
     where $n=[t/h ]$. Then $t-nh\in[0,h]$ and $t\to\infty$ implies $n\to\infty$. Moreover, observe 
     $\varphi_{ij}(t)$ in bounded in $[0,h]$. Thus for any $i,j\in E$, 
     \begin{equation*}
       \begin{split}
          |\varphi_{ij}(t)-v_iu_j| &=\left|\sum_{k\in E}\varphi_{ik}(t-nh)\[\varphi_{kj}(nh)-v_ku_j\]\right|  \\
            & \leq \sum_{k\in E}\max_{s\in[0,h]}\varphi_{ik}(s)|\varphi_{kj}(nh)-v_ku_j|\to0
            ,\quad t\to\infty.
       \end{split}
     \end{equation*}
     Till now, the proof is completed.
   \end{proof}
   
   \begin{question}
   	Suppose the Markov chain $P$ is irreducible and $i\neq j$.
   	\begin{enumerate}
   		\item If $f_{jj}=1$, then $\sum_{n\in\n}p_{ij}^{(n)}=\infty$ and $f_{ij}=1$.
   		\item If $f_{jj}<1$, then $\sum_{n\in\n}p_{ij}^{(n)}<\infty$. However, we can not determine $f_{ij}$.
   	\end{enumerate} 
   \end{question}
   \begin{proof}
   	\begin{enumerate}
   		\item If $f_{jj}=1$, the results from the following two equalities 
   		\begin{equation*}
   			\sum_{n\in\n}p_{ij}^{(n)}=\frac{f_{ij}}{1-f_{jj}}=\infty
   		\end{equation*}
   		and 
   		\begin{equation*}
   			f_{jj}=p_{jj}+\sum_{k\neq j}p_{jk}f_{kj}=\sum_{k\in E}p_{jk}f_{kj}\Longrightarrow F_j=PF_j,
   		\end{equation*}
   		where $F_j=(f_{kj})_{k\in E}$. For any $i\neq j$, the irreducibility implies exists $m\in \n$ such that $p_{ij}^{(m)}>0$. Therefore, 
   		\begin{equation*}
   			F_j=PF_j\Longrightarrow F_j=P^mF_j\Longrightarrow 1=f_{jj}=\sum_{k\in E}p_{jk}^{(m)}f_{kj}.
   		\end{equation*}
   		However $\sum_{k\in E}p_{jk}^{(m)}=1$, so $f_{ij}=1$.
   		\item If $f_{jj}<1$, The first equality shows $\sum_{n\in\n}p_{ij}^{(n)}<\infty$. We use two example to show $f_{ij}=1$ and $f_{ij}<1$ are both possible. Consider $E=\{1,2,3\}$ and $i=1,j=2$.
   		\begin{enumerate}
   		 	\item  Let 
   		 	\begin{equation*}
   		 		P=
   		 		\begin{pmatrix}
   		 			0 & 1 & 0\\
   		 			1/2 & 0 & 1/2\\
   		 			0 & 0 & 1\\
   		 		\end{pmatrix}
   		 	\end{equation*}
   		 	Observe $3$ is an absorb point, then $f_{12}=p_{12}=1,f_{12}^{(n)}=0(n\geq 2)$. Moreover, $f^{(2)}_{22}=1/2, f_{22}^{(n)}=0(n\neq 2)$. Therefore
   		 	\begin{equation*}
   		 		\begin{split}
   		 			f_{12}&=\sum_{n\in\n}f_{12}^{(n)}=p_{12}+\sum_{n=2}^\infty f_{12}^{(n)}=1+0=1,\\
   		 			f_{22}&=\sum_{n\in\n}f_{12}^{(n)}= f^{(2)}_{22}=1/2\in(0,1).
   		 		\end{split}
   		 	\end{equation*}
   		 	\item Let 
   		 	\begin{equation*}
   		 		P=
   		 		\begin{pmatrix}
   		 			0 & 1/2 & 1/2\\
   		 			1/2 & 0 & 1/2\\
   		 			0 & 0 & 1\\
   		 		\end{pmatrix}
   		 	\end{equation*}
   		 	Observe $3$ is an absorb point, then $f_{12}=p_{12}=1/2,p_{12}^{(n)}=0(n\geq 2)$. Moreover $f_{22}^{(2)}=1/4,f_{22}^{(n)}=0(n\neq 2)$ and $3$ is an absorb point. Therefore
   		 	\begin{equation*}
   		 		\begin{split}
   		 			f_{12}&=\sum_{n\in\n}f_{12}^{(n)}=p_{12}+\sum_{n=2}^\infty f_{12}^{(n)}=1/2+0=1/2,\\
   		 			f_{22}&=\sum_{n\in\n}f_{12}^{(n)}= f^{(2)}_{22}=1/4\in(0,1).
   		 		\end{split}
   		 	\end{equation*}
   		\end{enumerate}
   	\end{enumerate}
   \end{proof}
  
  \begin{question}
  	Suppose $E$ is a countable state space and $P=(p_{ij})_{i,j\in E}$ is an irreducible and aperiodic transition matrix. Suppose $P$ is null-recurrent, then $p_{ij}^{(n)}\to0,n\to\infty$.
  \end{question}
  \begin{proof}
  	The diagonal method provides a sub-sequence $\{n_m\}$ such that 
  	\begin{equation*}
  		v_k:=\lim_{m\to\infty}p_{ik}^{(n_m)},\quad\forall k\in E.
  	\end{equation*}
  	If exists $j\in E$ such that $v_j>0$, then the Fatou's lemma yields:
  	\begin{equation*}
  		\sum_{k}v_k\leq \lim_{m\to\infty}\sum_k p_{ik}^{(n_m)}=1.
  	\end{equation*}
  	Moreover, the dominated convergence theorem yields:
  	\begin{equation*}
  		\lim_{m\to\infty}p_{ik}^{(n_m+1)}=\lim_{m\to\infty}\sum_{\ell}p_{i\ell}p_{\ell k}^{(n_m)}=\sum_{\ell}p_{i\ell}v_k=v_k.
  	\end{equation*}
  	Therefore, use Fatou's lemma again to obtain
  	\begin{equation*}
  		v_k=\lim_{m\to\infty}p_{ik}^{(n_m+1)}=\lim_{m\to\infty}\sum_{\ell}p_{i\ell}^{(n_m)}p_{\ell k}\geq \sum_\ell v_\ell p_{\ell k}.
  	\end{equation*}
  	Add all the $k$ together, we will find the inequality above is actually equality. The irreducibility of $P$ ensures $v_k>0,\forall k\in E$. Observe $\sum_kv_k\leq 1$, then $P$ has a stationary distribution, which is equivalent to the positive recurrence of $P$. The contradiction arise.
  	
  	If $v_k=0,\forall k\in E$, then the dominated convergence theorem yields: $\forall r\geq 1$,
  	\begin{equation*}
  		\lim_{m\to\infty}p_{ik}^{(n_m+r)}=\lim_{m\to\infty}\sum_\ell p_{i\ell}^{(r)}p_{\ell k}^{(n_m)}=\sum_\ell p_{i\ell}^{(r)}\lim_{m\to\infty}p_{\ell k}^{(n_m)}=0.
  	\end{equation*}
  	
  	Next we shows the limit $\lim_{m\to\infty}p_{ik}^{(n_m+r)}$ is uniform with respect to the variable $r$. In fact, the coupling inequality 
  	\begin{equation*}
  		|p_{ik}^{(n)}-p_{jk}^{(n)}|\leq \p(T\geq n),\quad \forall i,j\in E
  	\end{equation*}
  	shows the limit $\lim_{m\to\infty}p_{\ell k}^{(n_m)}$ is uniform with respect to the state $\ell$. Namely, for any $\varepsilon>0$, exists $N\in\n$, for any $m>N$ and $\ell\in E$, we have $p_{\ell k}^{(n_m)}<\varepsilon$. This result together with the C-K equation yield: for any $r\in\n$,
  	\begin{equation*}
  		p_{ik}^{(n_m+r)}=\sum_\ell p_{i\ell}^{(r)}p_{\ell k}^{(n_m)}<\varepsilon\sum_\ell p_{i\ell}^{(r)}=\varepsilon.
  	\end{equation*}
  	Observe the $N$ mentioned above has nothing to do with $\ell\in E$ and $r\in\n$, the uniformity of $r$ in $\lim_{m\to\infty}p_{ik}^{(n_m+r)}$ naturally follows.
  	
  	For any $n\in\n$, exists unique $m\in\n$ such that $n_m\leq n<m_{m+1}$. Therefore we can write $p_{ik}^{(n)}=p_{ik}^{(n_m+r_m)}$, where $r_n=n-n_m$. Observe $n\to\infty$ implies $m\to\infty$, then the uniformity of $r$ in $\lim_{m\to\infty}p_{ik}^{(n_m+r)}$ implies 
  	\begin{equation*}
  		\lim_{n\to\infty}p_{ik}^{(n)}=\lim_{m\to\infty}p_{ik}^{(n_m+r_n)}=0.
  	\end{equation*}
  \end{proof}
  
  \begin{question}
  	If the state space $E$ is finite, and $P$ is random. Then there is at least one recurrent state and there is no null recurrence state.
  \end{question}
  \begin{proof}
  	Since the asymptotic property of non-recurrence has nothing to do with the irreducibility, then if all the states are non-recurrent, then $p_{ij}^{(n)}\to0,i,j\in E$. Therefore,
  	\begin{equation*}
  		1=\lim_{n\to\infty}\sum_jp_{ij}^{(n)}=\sum_j\lim_{n\to\infty}p_{ij}^{(n)}=0,
  	\end{equation*}
  	which is a contradiction. We have two ways to handle the next assertion.
  	\begin{enumerate}
  		\item If $k\in E$ is null-recurrent, consider the irreducible class containing $C_k$ such that $k\in C_k$. So all the states in $C_k$ is null-recurrent and therefore $C_k$ is closed. We have $p_{ij}^{(n)}\to0,i,j\in C_k$. The same argument provides the same contradiction.
  		\item If $k\in E$ is recurrent, we want to shows $k$ is positive recurrent. In fact, in the closed irreducible class $C_k$, exists invariant measure $\mu$ such that $\mu P|_{C_k}=\mu$. Observe $E$ therefore $C_k$ is finite, so the existence of invariant measure is equivalent to the existence of stationary distribution. Therefore $k$ is positive recurrent.
  	\end{enumerate}
  \end{proof}
  
  \begin{question}
  Suppose $P$ is a Markov chain and $\mu$ is its invariant measure. Define $\widetilde{p}_{ij}=\mu_jp_{ji}/\mu_i.$ Prove $\widetilde{P}=(\widetilde{p}_{ij} )$ is a Markov chain. $\widetilde{P}$ is called the dual process of $P$. Please prove:
  \begin{enumerate}
   \item $P$ is irreducible if and only if $\widetilde{P}$ is irreducible.
   \item $P$ is recurrent if and only if $\widetilde{P}$ is recurrent.
   \item $P$ is positive recurrent if and only if $\widetilde{P}$ is positive recurrent.
  \end{enumerate}
 \end{question}
 \begin{proof}
  \begin{enumerate}
   \item 
   First, we use induction on $n$ to verify $\widetilde{p}_{ij}^{(n)}=\mu_jp_{ji}^{(n)}/\mu_i$. When $n=1$, it is exactly the definition.
   Suppose when $n=k-1$, our assertion is true. When $n=k$,
   \begin{equation*}
    \begin{split}
    	\widetilde{p}_{ij}^{(k)}&=\sum_{\tau\in E}\widetilde{p}_{i\tau}^{(k-1)}\widetilde{p}_{\tau j}= \sum_{\tau\in E}\frac{\mu_\tau p_{\tau i}^{(k-1)}}{\mu_i} \frac{\mu_j p_{j\tau}}{\mu_\tau}\\
    	&=\frac{\mu_j}{\mu_i}\sum_{\tau\in E}p_{j\tau}p_{\tau i}^{(k-1)}=\frac{\mu_j}{\mu_i}p_{ji}^{(k)}.
    \end{split}
   \end{equation*} 
   Therefore the assertion is correct. For any $i,j\in E$, exist $n_{ij}\in\n$ such that $p_{ij}^{(n_{ij})}>0$ iff $\widetilde{ p}_{ij}^{(n_{ij})}>0$. So the irreducibility is equivalent.
   \item 
   For any $j\in E$,
   \begin{equation*}
    \sum_{n=0}^\infty \widetilde{p}_{jj}^{(n)}=\sum_{n=0}^\infty\frac{\mu_j}{\mu_j}p_{jj}^{(n)}=\sum_{n=0}^\infty p_{jj}^{(n)}.
   \end{equation*}
   So $P$ is recurrent iff $\widetilde{P}$ is recurrent.
   \item There are three ways to prove this assertion.
   \begin{enumerate}
   \item Notice $\forall j\in E,\widetilde{p}_{jj}^{(n)}=p_{jj}^{(n)}$, so
   \begin{equation*}
    \widetilde{ P}_{jj}(1)=\sum_{n=0}^\infty \widetilde{p}_{jj}^{(n)}=\sum_{n=0}^\infty p_{jj}^{(n)}=P_{jj}(1)
   \end{equation*}
   $\forall \{s_n\}_{n\in\n},$ and $s_n\uparrow 1 $, we have
   \begin{equation*}
    \begin{split}
     \widetilde{ m}_{jj}&=\li{n\to\infty}\(\left.\frac{\d }{\d s} \widetilde{ F}_{jj}(s)\right|_{s=s_n}\)=\li{n\to\infty} \[\frac{\d }{\d s}\left.\(1-\frac{1}{\widetilde{P}_{jj}(s)} \)\right|_{s=s_n}\]\\
     &=\li{n\to\infty} \[\frac{\d }{\d s}\left.\(1-\frac{1}{P_{jj}(s)} \)\right|_{s=s_n}\]=\li{n\to\infty}\(\left.\frac{\d }{\d s}  F_{jj}(s)\right|_{s=s_n}\)=m_{jj}
    \end{split}
   \end{equation*}
   So $P$ is positive recurrent if and only if $\widetilde{P}$ is positive recurrent.
   \item The existence of the dual chain $\widetilde{P}$ implies the existence of invariant measure $\mu$ of $P$. Moreover, 
   \begin{equation*}
   	\sum_{i\in E}\mu_i\widetilde{p}_{ij}=\sum_{i\in E}\mu_i\frac{\mu_jp_{ji}}{\mu_i}=\sum_{i\in E}\mu_jp_{ji}=\mu_i.
   \end{equation*}
   Therefore $\mu$ is also a invariant measure of $\widetilde{P}$. {\itshape By a general conclusion, $P$ is positive recurrent iff $P$ admit a stationary distribution (even $P$ is not irreducible, and $\mu$ is not unique at the same time)}, then 
   \begin{equation*}
   	P \text{ is positive recurrent}\Leftrightarrow \sum_{i\in E}\mu_i<\infty\Leftrightarrow \widetilde{P}  \text{ is positive recurrent}
   \end{equation*}
   \item Observe
   \begin{equation*}
   	\begin{split}
   		\widetilde{ f}_{jj}^{(n)}&=\sum_{j_1,\cdots,j_{n-1}\neq j}\widetilde{p}_{jj_1}\widetilde{p}_{j_1j_2}\cdots \widetilde{p}_{j_{n-1}j}\\
   		&=\sum_{j_1,\cdots,j_{n-1}\neq j}\frac{\mu_{j_1}p_{j_1j}}{\mu_j} \frac{\mu_{j_2}p_{j_2j_1}}{\mu_{j_1}}\cdots \frac{\mu_jp_{jj_{n-1}}}{\mu_{j_{n-1}}}\\
   		&=\sum_{j_1,\cdots,j_{n-1}\neq j}p_{jj_{n-1}}p_{j_{n-1}j_{n-2}}\cdots p_{j_{1}j}^{(n)}\\
   		&=f_{jj}.
   	\end{split}
   \end{equation*}
   Then $\widetilde{m}_{jj}=m_{jj},\forall j\in E$.
   \end{enumerate}
  \end{enumerate}
 \end{proof}
 
 \begin{question}
  Markov chain $(X_n)$ is invertible if and only if any times $ n,m\in\n$ and any states $i,j\in E$, then
  \begin{equation*}
   \p(X_n=i,X_m=j)=\p(X_n=j,X_m=i)
  \end{equation*}
  \begin{enumerate}
   \item The above formula is equivalent to any times $n_1<\cdots<n_m,$ which suffice $n_{j+1}-n_j=n_{m+1-j}-n_{m-j},j=1,\cdots,m-1$ and any state $i_1,\cdots,i_m\in E$, then
   \begin{equation*}
    \p(X_{n_1}=i_1,\cdots,X_{n_m}=i_m )=\p(X_{n_1}=i_m,\cdots,X_{n_m}=i_1 )
   \end{equation*}
   \item $(X_n)$ is invertible if and only if $\exists $  stationary distribution  $(\pi_i),\st \pi_ip_{ij}=\pi_jp_{ji}$ and the initial distribution $\pi_0=\pi$.
  \end{enumerate}
 \end{question}
 \begin{proof}
  \begin{enumerate}
   \item We just need to prove the sufficiency part. Let $n=0,m=1$, then we have
   \begin{equation*}
   	\begin{split}
   		\p(X_0=i)p_{ij}&=\p(X_0=i,X_1=j)\\
   		&=\p(X_0=j,X_1=i)=\p(X_0=j)p_{ji}.
   	\end{split}
   \end{equation*}
   Write $\pi_i:=\p(X_0=i)$ for $i\in E$, then the above formula yields
   \begin{equation*}
   	\pi_ip_{ij}=\pi_jp_{ji}\Longrightarrow \pi P^n=\pi,\quad\forall n\geq 1.\tag{$*$}
   \end{equation*}
   Thus $X_0\sim \pi$ is the reversible (stationary) distribution and thus for any $n\geq 1$,
   \begin{equation*}
   	\p(X_n=j)=\p(X_0=j)=\pi_j,\quad\forall j\in E.
   \end{equation*}
   Therefore
   \begin{equation*}
   	\begin{split}
   		&\p(X_{n_1}=i_1,\cdots,X_{n_m}=i_m )\\
   		=&\p(X_{n_1}=i_1)p_{i_1i_2}^{(n_2-n_1)}p_{i_2i_3}^{(n_3-n_2)}\cdots p_{i_{m-1}i_m}^{(n_m-n_{m-1})}\\
   		=&\pi_{i_1}p_{i_1i_2}^{(n_2-n_1)}p_{i_2i_3}^{(n_3-n_2)}\cdots p_{i_{m-1}i_m}^{(n_m-n_{m-1})}\\
   		=&p_{i_2i_1}^{(n_2-n_1)}\pi_{i_2}p_{i_2i_3}^{(n_3-n_2)}\cdots p_{i_{m-1}i_m}^{(n_m-n_{m-1})}\\
   		=&\cdots= p_{i_2i_1}^{(n_2-n_1)}p_{i_3i_2}^{(n_3-n_2)}\cdots p_{i_{m}i_{m-1}}^{(n_m-n_{m-1})}\pi_{i_m}\\
   		=&\pi_{i_m}p_{i_{m}i_{m-1}}^{(n_m-n_{m-1})}\cdots p_{i_3i_2}^{(n_3-n_2)}p_{i_2i_1}^{(n_2-n_1)} \\
   		=&\p(X_{n_1}=i_m,\cdots,X_{n_m}=i_1).
   	\end{split}
   \end{equation*}
   \item  Formula $(*)$ has proven the necessity part. For then sufficiency, we use induction on $n$ to prove $\forall n\in\n_+,\pi_ip_{ij}^{(n)}=\pi_jp_{ji}^{(n)}$. When $n=1$, this is the requirement of this question. Suppose when $n=k-1$, the assertion is right. When $n=k$,
   \begin{equation*}
    \begin{split}
     \pi_ip_{ij}^{(k)}&=\pi_i\sum_{\tau\in E}p_{i\tau}^{(k-1)}p_{\tau j} =\sum_{\tau\in E}\pi_ip_{i\tau}^{(k-1)}p_{\tau j}=\sum_{\tau\in E}\pi_\tau p_{\tau i}^{(k-1)}p_{\tau j}\\
     &=\sum_{\tau\in E} p_{\tau i}^{(k-1)}\pi_\tau p_{\tau j}=\sum_{\tau\in E} p_{\tau i}^{(k-1)}\pi_j p_{j\tau}=\pi_j \sum_{\tau\in E} p_{j\tau}p_{\tau i}^{(k-1)}\\
     &=\pi_jp_{ji}^{(k)}.
    \end{split}
   \end{equation*}
    In conclusion, our assertion holds for $\forall n\in\n$.
   Since $\pi$ is also a stationary distribution, then $\forall m\in\n, X_m\sim\pi$. So for any $n<m$,
   \begin{equation*}
    \begin{split}
    	\p(X_n=i,X_m=j)&=\p(X_n=i)p_{ij}^{(m-n)}=\pi_i p_{ij}^{(m-n)} \\
    	&=\pi_j p_{ji}^{(m-n)}=\p(X_n=j,X_m=i).
    \end{split}  
   \end{equation*}
  \end{enumerate}
 \end{proof}
  
  \begin{question}
  	Suppose $(X_n)_{n\in\n}$ is an irreducible, aperiodic Markov chain on $\z_+$. Then for any initial distribution $\mu$,
  	\begin{enumerate}
  		\item $(X_n)_{n\in\n}$ is non-recurrent iff $X_n\overset{\ae}{\to}\infty$ 
  		\item $(X_n)_{n\in\n}$ is non-recurrent iff $X_n\overset{\p}{\to}\infty$.
  	\end{enumerate}
  \end{question}
  \begin{proof}
  	\begin{enumerate}
  		\item Since for any $N\in\n$ and any initial distribution $\mu=(\mu_k)_{k\in E}$, we have
  		\begin{equation*}
  			\begin{split}
  				\sum_{n=0}^\infty\p^{\mu}(X_n\leq N)&=\sum_{n=0}^\infty\sum_{j=0}^N \p^\mu(X_n= j)=\sum_{n=0}^\infty\sum_{j=0}^N\sum_{k\in E} \mu_kp_{k j}^{(n)}\\
  				&=\sum_{j=0}^N\sum_{k\in E}\mu_k\sum_{n=0}^\infty\ p_{k j}^{(n)}=\sum_{j=0}^N\sum_{n=0}^\infty p_{\mu j}^{(n)}
  			\end{split}
  		\end{equation*}
  	\end{enumerate}
  \end{proof}
  
  \begin{question}
  	Suppose there are two kinds of genes: $A$ and $B$. And the number of the genes is $2N$. The composition of the next generation's genes is determined by $2N$ times of independent binomial experiments. Namely, if the numbers of $A$ and $B$ are $i$ and $2N-i$, then the probabilities of each experiment with outcome $A$ and $B$ are
  	\begin{equation*}
  		p_i=\frac{i}{2N},\quad q_i=1-p_i.
  	\end{equation*} 
  	Let $X_n$ be the number of $A$ in the $n$th generation. Prove
  	when $0<j<2N$,
  		\begin{equation*}
  			\lim_{n\to\infty}p_{ij}^{(n)}=0,\quad \lim_{n\to\infty}p_{i0}^{(n)}=1-\frac{i}{2N},\quad \lim_{n\to\infty}p_{i,2N}^{(n)}=\frac{i}{2N}.
  		\end{equation*}
  \end{question}
  \begin{proof}
  	Define $T=\inf\{n\in\n_+:X_n=0\text{ or }2N\}=\tau_0^+\wedge \tau_{2N}^+$. Observe $A:=\{1,\cdots,2N-1\}$ is a irreducible class and $P|_A$ is sub-random, then $p_{ij}^{(n)}\to0,\forall i,j\in A$. Moreover, note $0,2N$ are two absorbing points, then 
  	\begin{equation*}
  		p_{i0}^{(n)}=\sum_{m=1}^n f_{i0}^{m}p_{00}^{(n-m)}=\sum_{m=1}^n f_{i0}^{m}.
  	\end{equation*}
  	So $f_{i0}=\lim_{n\to\infty}p_{i0}^{(n)}$ and similarly $f_{i,2N}=\lim_{n\to\infty}p_{i,2N}^{(n)}$.
  	Therefore for any $i\in A$,
  	\begin{equation*}
  		\begin{split}
  			\p^i(T<\infty)&=\p^i(\tau_0^+<\infty)+\p^i(\tau_{2N}^+<\infty)=f_{i0}+f_{i,2N}\\
  			&=\lim_{n\to\infty}\(p_{i0}^{(n)}+p_{i,2N}^{(n)}\)=\lim_{n\to\infty}\(1-\sum_{j=1}^{2N}p_{ij}^{(n)}\)\\
  			&=1-\sum_{j=1}^{2N}\lim_{n\to\infty}p_{ij}^{(n)}=1.
  		\end{split}
  	\end{equation*}
  	We claim $(X_n)_{n\in\n}$ is a martingale. In fact, given $X_n=i$, we know $X_{n+1}\sim B(2N,i/2N)$, then the property of binomial distribution yields
  	\begin{equation*}
  		\p(X_{n+1}|X_{n}=i)=i.
  	\end{equation*}
  	Hence $\p(X_{n+1}|X_n)=X_n$ and $(X_n)_{n\in\n}$ is a martingale. Although $T$ is not bounded, we have $|X_n|\leq 2N$. Therefore the {\itshape Doob's stopping theorem} still yields
  	\begin{equation*}
  		0f_{i0}+2Nf_{i,2N}=\p^i(X_T)=\p^i(X_0)=i.
  	\end{equation*}
  	This formula together with $f_{i0}+f_{i,2N}=1$ show us 
  	\begin{equation*}
  		\lim_{n\to\infty}p_{i0}^{(n)}=f_{i0}=\frac{i}{2N},\quad \lim_{n\to\infty}p_{i,2N}^{(n)}=f_{i,2N}=\frac{i}{2N}.
  	\end{equation*}
  \end{proof}
  
  \begin{question}
  	Suppose $(M_t)_{t\geq 0}$ is a $(\f_t)_{t\geq 0}$ martingale. Given $0\leq r<u$ and $Z\in\f_r$, define
  	\begin{equation*}
  		N_t:=Z(M_{u\wedge t}-M_{r\wedge t}).
  	\end{equation*}
  	Then $(N_t)_{t\geq 0}$ is also a $(\f_t)_{t\geq 0}$ martingale.
  \end{question}
  \begin{proof}
  	Observe for any $s,t>0$, we have $\p(M_t|\f_s)=M_{s\wedge t}$ and $N_t=0,\forall 0\leq t\leq r$. Then for any $0\leq s<t$,
  	\begin{enumerate}
  		\item When $s<r$, we have
  		\begin{equation*}
  			\begin{split}
  				\p(N_t|\f_s)&=\p[Z(M_{u\wedge t}-M_{r\wedge t})|\f_s]\\
  				&=\p\{\p[Z(M_{u\wedge t}-M_{r\wedge t})|\f_r]\f_s\}\\
  				&=\p[Z\p(M_{u\wedge t}-M_{r\wedge t}|\f_r)\f_s]\\
  				&=\p[Z(M_{u\wedge t\wedge r}-M_{r\wedge t\wedge r})|\f_s]\\
  				&=\p[Z(M_{r\wedge t}-M_{r\wedge t})|\f_s]\\
  				&=0=N_s.
  			\end{split}
  		\end{equation*}
  		\item When $s\geq r$, we have
  		\begin{equation*}
  			\begin{split}
  				\p(N_t|\f_s)&=\p[Z(M_{u\wedge t}-M_{r\wedge t})|\f_s]\\
  				&=Z\p(M_{u\wedge t}-M_{r\wedge t}|\f_s)\\
  				&=Z(M_{u\wedge t\wedge s}-M_{r\wedge t\wedge s})\\
  				&=Z(M_{u\wedge t}-M_{r\wedge s})\\
  				&=N_s.
  			\end{split}
  		\end{equation*}
  	\end{enumerate}
  	Thus the proof is completed.
  \end{proof}
  
  \begin{question}
  	Prove the minimum solution of the forward equation and the backward equation are identical.
  \end{question}
  \begin{proof}
  	Consider the following expression of the forward equation and the backward equation:
  	\begin{equation*}
  		\begin{split}
  			x_{ij}&=\sum_{k\neq i}\frac{q_{ik}}{\lambda+q_i}x_{kj}+\frac{\delta_{ij}}{\lambda+q_i},\\
  			y_{ij}&=\sum_{k\neq j}y_{ik}\frac{q_{kj}}{\lambda+q_j}+\frac{\delta_{ij}}{\lambda+q_j}.
  		\end{split}
  	\end{equation*}
  	Consider the second iteration. Write the outcomes of the each step as $\{x_{ij}^{(n)}\}_{i,j\in E}$ and $\{y_{ij}^{(n)}\}_{i,j\in E}$ for the forward equation and the backward equation respectively. We use induction on $n$ to prove this. 
  	\begin{enumerate}
  		\item When $n=1$, then 
  		\begin{equation*}
  			x_{ij}^{(1)}=\frac{\delta_{ij}}{\lambda+q_i}=\frac{\delta_{ij}}{\lambda+q_j}=y_{ij}^{(1)}.
  		\end{equation*}
  		\item When $n\leq \ell$, assume the statement is correct. When $n=\ell+1$, we have
  		\begin{equation*}
  			\begin{split}
  				x_{ij}^{(\ell+1)}&=\sum_{k\neq i}\frac{q_{ik}}{\lambda+q_i}x_{kj}^{(\ell)}=\sum_{k\neq i}\frac{q_{ik}}{\lambda+q_i}y_{kj}^{(\ell)}\\
  				&=\sum_{k\neq i}\frac{q_{ik}}{\lambda+q_i}\sum_{s\neq j}y_{ks}^{(\ell-1)}\frac{q_{s j}}{\lambda+q_j}=\sum_{s\neq j}\sum_{k\neq i}\frac{q_{ik}}{\lambda+q_i}x_{ks}^{(\ell-1)}\frac{q_{s j}}{\lambda+q_j}\\
  				&=\sum_{s\neq j}x_{is}^{(\ell)}\frac{q_{s j}}{\lambda+q_j}=\sum_{s\neq j}y_{is}^{(k)}\frac{q_{s j}}{\lambda+q_j}\\
  				&=y_{ij}^{(\ell+1)}.
  			\end{split}
  		\end{equation*}
  		\item So $x_{ij}^{(n)}=y_{ij}^{(n)}$ for any $n\geq 1$.
  	\end{enumerate}
  	Therefore 
  	\begin{equation*}
  		p^\text{min}_{ij}=x_{ij}=\sum_{n\geq 1}x_{ij}^{(n)}=\sum_{n\geq 1}y_{ij}^{(n)}.
  	\end{equation*}
  \end{proof}
  
  \begin{question}
  	Given a continuous process $(X_t)_{t\geq 0}$ on the probability space $(\Omega,\f,\p)$ and valued on metric space $(E,d)$. Define the stopping time
  	\begin{equation*}
  		T_a:=\inf\{t\geq 0:X_t=a\}.
  	\end{equation*}
  	Then
  	\begin{equation*}
  		\begin{split}
  			\{T_a\leq t\}=&\{\exists 0\leq s\leq t, X_s=a \}\\
  			=&\left\{\inf_{0\leq s\leq t}d(X_s,a)=0\right\}\\
  			=&\left\{\inf_{s\in[0,t]\cap\q}d(X_s,a)=0\right\}
  		\end{split}
  	\end{equation*}
  \end{question}
  \begin{proof}
  	\begin{enumerate}
  		\item We have
  	\begin{equation*}
  		\{T_a\leq t\}=\{T_a<t\}\cup\{T_a=t\}=\{\exists 0\leq s< t, X_s=a \}\cup\{T_a=t\}.
  	\end{equation*}
  	Since $T_a=t$ is equivalent to either of the following statement:
  	\begin{enumerate}
  		\item $X_s\neq 0$ for any $0\leq s<t$ and $X_t=a$.
  		\item there is a sequence $\{t_n\}\subset(t,\infty)$ such that $t_n\downarrow t$ and $X_{t_n}=a$ for any $n\geq 1$. Then the continuity of $(X_t)$ implies 
  	\begin{equation*}
  		X_t=\lim_{n\to\infty}X_{t_n}=a.
  	\end{equation*}
  	\end{enumerate}
  	 Therefore $\{T_a=t\}\subset \{\forall 0\leq s< t, X_s\neq a,X_t=a\}$. Besides, another direction is obvious. So $\{T_a=t\}=\{\forall 0\leq s< t, X_s\neq a,X_t=a\}$ and it is easy to verify
  	 \begin{equation*}
  	 	\begin{split}
  	 		&\{T_a\leq t\}\\
  	 		=&\{\exists 0\leq s< t, X_s=a \}\cup\{\forall 0\leq s< t, X_s\neq a,X_t=a\}\\
  	 		=&\{\exists 0\leq s\leq t, X_s=a \}.
  	 	\end{split}
  	 \end{equation*}
  		\item Now we start to prove the second equality. First, 
  	 \begin{equation*}
  	 	\{\exists 0\leq s\leq t, X_s=a \}\subset \left\{\inf_{0\leq s\leq t}d(X_s,a)=0\right\}
  	 \end{equation*}
  	 is trivial. Ob the contrary, if $\inf_{0\leq s\leq t}d(X_s,a)=0$, then we can find a sequence $\{s_n\}\subset[0,t]$ such that $d(X_{s_n},a)\to 0$. Knowing $[0,t]$ is compact, then we can find a convergent subsequence $\{s_{n_k}\}$, say $s_{n_k}\to s_0\in[0,t]$. Then the continuity of $(X_t)$ and $d(\cdot,\cdot )$ yields
  	 \begin{equation*}
  	 	d(X_{s_0},a)=\lim_{k\to\infty}d(X_{s_{n_k}},a)=0,
  	 \end{equation*} 
  	 which means $X_{s_0}=a$. Therefore 
  	 \begin{equation*}
  	 	\{\exists 0\leq s\leq t, X_s=a \}\supset \left\{\inf_{0\leq s\leq t}d(X_s,a)=0\right\}.
  	 \end{equation*}
  		\item Finally, we started to prove the third equality. We only need to show
  		\begin{equation*}
  			\left\{\inf_{0\leq s\leq t}d(X_s,a)=0\right\}\subset \left\{\inf_{s\in[0,t]\cap\q}d(X_s,a)=0\right\}.
  		\end{equation*}
  		In fact, if $\inf_{0\leq s\leq t}d(X_s,a)=0$, then we can find a sequence $\{s_n\}\subset[0,t]$ such that $d(X_{s_n},a)\to 0$. Utilize the density of $\q $ in $\r$ and the continuity of $(X_t)$, for any $n\geq 1$, we can find $q_n\in\q$ such that $d(X_{q_n},X_{s_n})<1/n$. Thus 
  		\begin{equation*}
  			\begin{split}
  				\inf_{s\in[0,t]\cap\q}d(X_s,a)\leq& \varlimsup_{n\to\infty}d(X_{q_n},a)\\
  				\leq& \varlimsup_{n\to\infty}d(X_{q_n},X_{s_n})+\varlimsup_{n\to\infty}d(X_{s_n},a)=0.
  			\end{split}
  		\end{equation*}
  	\end{enumerate}
  	Thus the whole statement is completed.
  \end{proof}
  
  \begin{question}
  	Consider a Markov chain $(X_n)_{n\geq 0}$ on probability space $(\Omega,\f,\p)$ and valued on countable state space $E$. For each $a\in E$, consider the following stopping times
  	\begin{equation*}
  		T_a=\inf\{n\geq 0:X_n=a\},\quad T_a^+=\inf\{n\geq 1:X_n=a\}.
  	\end{equation*}
  	Then 
  	\begin{equation*}
  		T^+_a(\omega)=
  		\begin{cases}
  			(T_a\circ \theta_1)(\omega)+1=(T_a^+\circ\theta_1)(\omega)+1 &,\mbox{if }\omega\in \{X_1\neq a\},\\
  			(T_a\circ\theta_1)(\omega)+1&,\mbox{if }\omega\in \{X_1= a\}.\\
  		\end{cases}
  	\end{equation*}
  	Moreover, if additionally $\omega\in \{ X_0\neq a\}$, then $T_a^+(\omega)=T_a(\omega)$. Hence $T_a$ also has the above expression.
  \end{question}
  \begin{proof}
  	Since if $\omega\in \{X_1\neq a\}$, then obviously $(T_a\circ \theta_1)(\omega)=(T_a^+\circ\theta_1)(\omega)$. Then we just need to show for any $\forall\omega\in\Omega$,
  	\begin{equation*}
  		T^+_a(\omega)=(T_a\circ \theta_1)(\omega)+1.
  	\end{equation*}
  	If $\omega\in \{X_1= a\}$, then $T^+_a(\omega)=1$ and $(T_a\circ \theta_1)(\omega)=0$. So the above formula holds. If $\omega\in \{X_1\neq a\}$, by the definition of shift operator $\theta_1$, we have the following equality
  	\begin{equation*}
  		T^+_a(\omega)-(T_a^+\circ \theta_1)(\omega)=1.
  	\end{equation*}
  	Moreover, we know $(T_a\circ \theta_1)(\omega)=(T_a^+\circ\theta_1)(\omega),\forall \omega\in \{X_1\neq a\}$. So the proof is completed.
  \end{proof}
  
  \begin{question}
  	Suppose $Q$ is a $Q$ matrix on the state space $E$ and $\sup_{i\in E}|q_{i}|<\infty$. Then we have the following estimation:
  	\begin{equation*}
  		\sup_{i,j\in E}|q_{ij}^{(n)}|\leq  \frac{1}{2} \(2\sup_{i\in E}|q_{i}|\)^n,\quad  \forall n\geq 1.
  	\end{equation*}
  	Especially, when $|E|=m<\infty$, for any $A=(a_{ij})\in M_m(\r)$, define $||A||:=m\sup_{i,j\in E}|a_{ij}|$, then $(M_m(\r),||\cdot||)$ is a Banach algebra.
  \end{question}
  \begin{proof}
  	We use induction on $n$ to prove the first estimation. 
  	\begin{enumerate}
  		\item When $n=1$, obviously $|q_{ij}|\leq \sup_{i\in E}|q_{i}|$.
  		\item Suppose $n=k$, the estimation is correct. When $n=k+1$, we have
  		\begin{equation*}
  			\begin{split}
  				|q_{ij}^{(k+1)}|&=\left|\sum_{k}q_{ik}q_{kj}^{(k)} \right|\leq \sum_{k\neq i}q_{ik}|q_{kj}^{(k)}|+q_i|q_{ij}^{(k)}|\\
  				&\leq \frac{1}{2} \(2\sup_{i\in E}|q_{i}|\)^k\(\sum_{k\neq i}q_{ik}+q_i \)\\
  				&=\frac{1}{2} \(2\sup_{i\in E}|q_{i}|\)^k2q_i\\
  				&\leq \frac{1}{2} \(2\sup_{i\in E}|q_{i}|\)^{k+1}.
  			\end{split}
  		\end{equation*}
  		\item In conclusion, the estimation is correct for all $n\geq 1$.
  	\end{enumerate}
  	Furthermore, we can infer from this conclusion that $e^{tQ}$ is well-defined.
  	
  	Obviously $||\cdot||$ is a norm. For any $A=(a_{ij}),B=(b_{ij})\in M_n(\r)$, we have
  	\begin{equation*}
  		\begin{split}
  			||AB||&=m\sup_{i,j\in E}\left|\sum_{k\in E}a_{ik}b_{kj}\right|\leq m\sum_{k\in E}\sup_{i,j\in E}a_{ij}\sup_{i,j\in E}b_{ij}\\
  			&=m^2\sup_{i,j\in E}a_{ij}\sup_{i,j\in E}b_{ij}=||A||\cdot||B||.
  		\end{split}
  	\end{equation*}
  	Thus $(M_m(\r),||\cdot||)$ is a Banach algebra. Moreover, it provide us a less strict estimation for $|q_{ij}^{(n)}|$ when $|E|=m<\infty$. In fact, 
  	\begin{equation*}
  		m\sup_{i,j\in E}|q_{ij}^{(n)}|=||Q^n||\leq ||Q||^n=\(m\sup_{i,j\in E}q_{ij}\)^n,
  	\end{equation*}
  	which is equivalent to
  	\begin{equation*}
  		\sup_{i,j\in E}|q_{ij}^{(n)}|\leq \frac{1}{m}\(m\sup_{i,j\in E}q_{ij}\)^n.
  	\end{equation*}
  \end{proof}
  
  \begin{question}
  	Consider a symmetry matrix $A\in M_n(\r)$ and write all its eigenpairs as $(\lambda_j,x_j)_{j=1}^n$. Then 
  	\begin{equation*}
  		A=\sum_{j=1}^n \lambda_j x_jx_j^\mathrm{T}=\sum_{j=1}^n \lambda_j \langle \cdot, x_j\rangle x_j.
  	\end{equation*}
  \end{question}
  \begin{proof}
  	Since $A\in M_n(\r)$ is symmetry, then indeed it has $n$ eigenvalues which form a unit standard orthogonal basis, denoted as $\{x_j\}_{j=1}^n$ (column vectors). And we write their corresponding eigenvalues as $\{\lambda_j\}_{j=1}^n$ (note some eigenvalues in $\{\lambda_j\}$ might be repeated). Then exists a orthogonal matrix $P$ such that 
  	\begin{equation*}
  		PAP^\mathrm{T} =
  		\begin{pmatrix}
  			\lambda_1 &  & & \\
  			 & \lambda_2 & & \\
  			 & &\ddots &\\
  			 & & & \lambda_n
  		\end{pmatrix},
  	\end{equation*}
  	where $P=\begin{pmatrix}
  		x_1 & x_2 & \cdots & x_n
  	\end{pmatrix}$. Therefore
  	\begin{equation*}
  		\begin{split}
  			A&=P^\mathrm{T}\begin{pmatrix}
  			\lambda_1 &  & & \\
  			 & \lambda_2 & & \\
  			 & &\ddots &\\
  			 & & & \lambda_n
  		\end{pmatrix}P\\
  		&= \begin{pmatrix}
  			x_1^\mathrm{T}\\
  			x_2^\mathrm{T}\\
  			\vdots\\
  			x_n^\mathrm{T}
  		\end{pmatrix}
  		\begin{pmatrix}
  			\lambda_1 &  & & \\
  			 & \lambda_2 & & \\
  			 & &\ddots &\\
  			 & & & \lambda_n
  		\end{pmatrix}
  		\begin{pmatrix}
  		x_1 & x_2 & \cdots & x_n
  	\end{pmatrix}\\
  	&=\sum_{j=1}^n \lambda_j x_jx_j^\mathrm{T}.
  		\end{split}\tag{$*$}
  	\end{equation*}
  	Moreover, since $\{x_j\}_{j=1}^n$ is a unit standard orthogonal basis on $\r^n$, then for any $y\in\r^n$, exists $\{y_j\}_{j=1}^n$ such that
   \begin{equation*}
   	y=\sum_{j=1}^n y_jx_j.
   \end{equation*}
   Therefore
   \begin{equation*}
   	Ay=\sum_{j=1}^n y_jAx_j=\sum_{j=1}^n \lambda_j y_j x_j=\sum_{j=1}^n \lambda_j \langle y,x_j\rangle x_j,
   \end{equation*}
   which is equivalent to
   \begin{equation*}
   	A=\sum_{j=1}^n \lambda_j \langle \cdot,x_j\rangle x_j.\tag{$**$}
   \end{equation*}
   Compare $(*)$ and $(**)$, we can concludes that
   \begin{equation*}
   	x_jx_j^\mathrm{T}=\langle \cdot,x_j\rangle x_j.
   \end{equation*}
  \end{proof}
  
  \begin{question}
  	Suppose $N=(N_t,t\geq 0)$ is a Poisson process with parameter $\lambda>0$, then the following statement is true:
  	\begin{enumerate}
  		\item $N$ is increasing.
  		\item For any $t\geq 0$, there exists $\delta>0$ such that $N$ is constant in the interval $(t,t+\delta]$.
  		\item Define stopping times recursively by $T_0=0$ and $T_n=\inf\{t>T_{n-1}:N_t-N_{T_{n-1}}\}$ for $n\geq 1$. Then $N_{T_n}=n$ for $n\in\z_+$.
  		\item $(\Delta N(t),t\geq 0)$ takes values in $\{0,1\}$ and for any $t\geq 0$.
  	\end{enumerate}
  \end{question}
  \begin{proof}
  	\begin{enumerate}
  		\item For any $0\leq s<t$, we have $N_t=N_s+(N_t-N_s)$ and $N_t-N_s\sim P(\lambda(t-s))$. So $N_t\geq N_s$ and therefore $N$ is increasing.
  		\item If there exists $t\geq 0$ satisfying for any $\delta>0$, $N_t$ is not constant in the interval $(t,t+\delta]$. Note $N$ is increasing, then there exist $t<t_1< t+\delta$ satisfying $N_{t+\delta}> N_{t_1}$. Since $N_{t+\delta}- N_{t_1}$ is poisson distributed, then $N_{t+\delta}- N_{t_1}\geq 1$. Consider $(t,t+t_1]$, since $N$ is not constant in $(t,t+t_1]$, we can find $t<t_2<t_1$ and $N_{t_1}-N_{t_2}\geq 1$. Do this repeatedly, we will obtain a sequence $\{t_n\}_{n\geq 1}$ satisfying
  		\begin{equation*}
  			t<\cdots<t_n<t_{n-1}<\cdots<t_1<t_0:=t+\delta\text{ and } N_{t_n}-N_{t_{n+1}}\geq 1.
  		\end{equation*}
  		So for any $n\geq 1$, 
  		\begin{equation*}
  			N_{t+\delta}-N_t\geq \sum_{k=0}^n(N_{t_k}-N_{t_{k+1}})\geq n.
  		\end{equation*}
  		Observe $n$ is arbitrary, then $N_{t+\delta}-N_t=\infty$. However we know $N_{t+\delta}-N_t$ is Poisson distributed, so $\p(N_{t+\delta}-N_t=\infty)=0$.
  		\item First we show $\Delta N_t\in\z_+$. Apply the increasing property of $N$, we have
  	\begin{equation*}
  		\Delta N_t=N_t-N_{t-}=\lim_{n\to\infty}(N_t-N_{t-1/n}).
  	\end{equation*}
  	Observe $N_t-N_{t-1/n}\sim P(\lambda/n)$, then $N_t-N_{t-1/n}$ is integer valued and therefore $\Delta N_t$ is integer valued. 
  	\end{enumerate}
  \end{proof}
  
  \begin{question}
  	Suppose $(X_n)_{n\geq 1}$ is an irreducible Markov chain on a countable state space $E$. 
  	\begin{enumerate}
  		\item For $x\neq y\in E$, if there exists $n\geq 1$ satisfies $p_{ij}^{(n)}>0$, then there exists $m\leq n$ such that $\p_i(X_m=j,\tau_i^+>m)>0$. 
  		\item Apply this property to show: if $f_{ii}=1$ for any $i\in E$, then $f_{ij}=1$ for any $i,j\in E$.
  	\end{enumerate}
  \end{question}
  \begin{proof}
  	\begin{enumerate}
  		\item If $\p_i(X_n=j,\tau_i^+>n)>0$, then all is done. If $\p(X_n=j,\tau_i^+>n)=0$, then $\p_i(X_n=j)=\p_i(X_n=j,\tau_i^+\leq n)$ and
  		\begin{equation*}
  			\begin{split}
  				\p_i(X_n&=j,\tau_i^+\leq n)=\p_i(X_n=j,\tau_i^+< n)\\
  				&=\sum_{k=1}^{n-1}\p_i(X_n=j,\tau_i^+= k)=\sum_{k=1}^{n-1}\p_i(X_n=j|\tau_i^+= k)\p_i(\tau_i^+= k)\\
  				&=\sum_{k=1}^{n-1}\p_i(X_{n-k}=j)\p_i(\tau_i^+= k)>0.
  			\end{split}
  		\end{equation*}
  		So there exists $n_1< n$ such that $p_{ij}^{(n_1)}=\p_i(X_{n_1}=j)>0$. Again if $\p_i(X_{n_1}=j,\tau_i^+>n_1)>0$, then all is done. If $\p_i(X_{n_1}=j,\tau_i^+>n_1)=0$, we can repeat the above procedure to find there $n_2< n_1$ such that $p_{ij}^{(n_2)}=\p_i(X_{n_2}=j)>0$. If there exists $1<m<n$ such that $\p_i(X_{m}=j,\tau_i^+>m)=0$, then all is done. If not, we will find 
  		\begin{equation*}
  			\p_i(X_1=j)=\p_i(X_1=j,\tau_i^+>1)>0.
  		\end{equation*}
  		\item  If exists $i\neq j\in E$ such that $f_{ij}<1$, then $\p_i(\tau_j^+=\infty)>0$. Since $(X_n)_{n\geq 1}$ is irreducible, then there exists $n\geq 1$ satisfies $p_{ji}^{(n)}>0$. So what we have proved yields $\p_j(X_m=i,\tau_j^+>m)>0$. Thus
  		\begin{equation*}
  			\begin{split}
  				\p_j(\tau_j^+&=\infty)=\sum_{k\in E}\p_j(X_m=k,\tau_j^+>m,\tau_j^+=\infty)\\
  				&\geq \p_j(X_m=i,\tau_j^+>m,\tau_j^+=\infty)\\
  				&=\p_j(X_m=i,\tau_j^+>m,\tau_j^+\circ\theta_m+m =\infty)\\
  				&=\p_j(\tau_j^+\circ\theta_m+m =\infty|X_m=i,\tau_j^+>m)\p_j(X_m=i,\tau_j^+>m)\\
  				&=\p_j(\tau_j^+\circ\theta_m =\infty|X_m=i,\tau_j^+>m)\p_j(X_m=i,\tau_j^+>m)\\
  				&=\p_i(\tau_j^+=\infty)\p_j(X_m=i,\tau_j^+>m)>0.
  			\end{split}
  		\end{equation*}
  		This contradict the fact $f_{jj}=1$.
  	\end{enumerate}
  \end{proof}
  
  \begin{question}
  	Given a nonnegative supermartingale $(X_n)_{n\geq 0}$ with respect to $(\f_n)_{n\geq 1}$. Then $X_\infty:=\lim_{n\to\infty}X_n$ exists almost surely. Moreover we have $\p(X_\infty|\f_n)\leq X_n$ for any $n\geq 0$ and $\{X_\infty<\infty\}\subset\{X_0<\infty\}$ almost surely.
  \end{question}
  \begin{proof}
  	$X_\infty:=\lim_{n\to\infty}X_n$ exists almost surely is a well-known result. To show $\p(X_\infty|\f_n)\leq X_n$, we use the nonnegativity of $(X_n)_{n\geq 1}$ and the Fatou's lemma. For any $A\in\f_n$ and $m\geq n$,
  	\begin{equation*}
  		\p(X_\infty 1_A)\leq \varliminf_{m\to\infty}\p(X_m1_A)\leq \varliminf_{m\to\infty}\p(X_n1_A)=\p(X_n1_A).
  	\end{equation*}
  	This yields  $\p(X_\infty|\f_n)\leq X_n$. Let $n=0$ and $A=\{X_0\leq M\}$, where $M>0$. Then 
  	\begin{equation*}
  		\p(X_\infty \mathbbm{1}_{\{X_0\leq M\}})\leq \p(X_0 \mathbbm{1}_{\{X_n\leq M\}})\leq M<\infty.
  	\end{equation*}
  	This means $\{X_\infty=\infty\}\cap\{X_0\leq M\}=\emptyset,\p$-a.s. for any $M>0$. Hence
  	\begin{equation*}
  		\{X_\infty=\infty\}\subset \bigcap_{M=1}^\infty\{X_0>M\}=\{X_0=\infty\},\quad \p\text{-a.s.}.
  	\end{equation*} 
  \end{proof}
  
  \begin{question}
  	Suppose $(\xi_n,n\geq 0)$ is a discrete-time Markov process with transition kernel $P(x,dy)$ and initial distribution $\mu(dx)$ on $(E,\e)$. Besides, $(N_t,t\geq 0)$ is an independent Poisson process with parameter $\alpha>0$. Let $X_t:=\xi_{N_t}$. Then the continuous-time process $(X_t,t\geq 0)$ is a Markov process with transition kernel as follows:
  	\begin{equation*}
  		Q_t(x,A):=e^{-\alpha t}\sum_{k\geq 0}\frac{(\alpha t)^k}{k!}P^k(x,A),\quad A\in\e.
  	\end{equation*}
  \end{question}
  \begin{proof}
  	We just need to show that 
  	\begin{equation*}
  		\p[f(X_{s+t})|X_{s_1},\cdots,X_{s_n}]=Q_tf(X_s) 
  	\end{equation*} 
  	for any $n\geq 1,t\geq0,0=s_0\leq s_1\leq \cdots\leq s_n=s$ and $f\in\mathrm{b}\e$. And this is again equivalent to showing
  	\begin{equation*}
  		\p[f(X_{s+t})g(X_{s_1},\cdots,X_{s_n})]=\p[Q_tf(X_s)g(X_{s_1},\cdots,X_{s_n})]
  	\end{equation*}
  	for any $n\geq 1,t\geq0,0=s_0\leq s_1\leq \cdots\leq s_n=s$ and $g\in\mathrm{b}\e^n$. First note by the independency of $(\xi_n,n\geq 0)$ and $(N_t,t\geq 0)$,
  	\begin{align*}
  		&\p[g(X_{s_1},\cdots,X_{s_n})]=\p[g(\xi_{N_{s_1}},\cdots,\xi_{N_{s_n}} )]\\
  		=&\sum_{k_1,\cdots,k_n\geq 0}\p\[g(\xi_{N_{s_1}},\cdots,\xi_{N_{s_n}} );N_{s_r}-N_{s_{r-1}}=k_r,1\leq r\leq n \]\\
  		=&\sum_{k_1,\cdots,k_n\geq 0}\p\[g(\xi_{k_1},\cdots,\xi_{k_1+\cdots+k_n});N_{s_r}-N_{s_{r-1}}=k_r,1\leq r\leq n \] \\
  		=&\sum_{k_1,\cdots,k_n\geq 0}\p[g(\xi_{k_1},\cdots,\xi_{k_1+\cdots+k_n} )]\p(N_{s_r}-N_{s_{r-1}}=k_r,1\leq r\leq n )\\
  		=&\sum_{k_1,\cdots,k_n\geq 0}\int_{E^n}g(x_1,\cdots,x_n)\mu P_{J_n}(dx_1,\cdots,dx_n)\p(N_{s_r}-N_{s_{r-1}}=k_r,1\leq r\leq n ),
  	\end{align*}
  	where $J_n=\{k_1,k_1+k_2,\cdots,k_1+\cdots+ k_n\}$ and
  	\begin{equation*}
  		P_{J_n}(dx_1,\cdots,dx_n)=P_{k_1}(x,dx_1)P_{k_2}(x_1,dx_2)\cdots P_{k_n}(x_{n-1},dx_n).
  	\end{equation*}
  	Thus let $s_{n+1}=s+t$, we have
  	\begin{align*}
  		&\p[f(X_{s+t})g(X_{s_1},\cdots,X_{s_n})]\\
  			=&\sum_{k_1,\cdots,k_n,k_{n+1}\geq 0}\int_{E^{n+1}}f(x_{n+1})g(x_1,\cdots,x_n)\mu P_{J_{n+1}}(dx_1,\cdots,dx_{n+1})\\
  			&\qquad\times \p(N_{s_r}-N_{s_{r-1}}=k_r,1\leq r\leq n+1)\\
  			=&\sum_{k_1,\cdots,k_n,k_{n+1}\geq 0}\int_{E^{n}}P_{k_{n+1}}f(x_{n})g(x_1,\cdots,x_n)\mu P_{J_n}(dx_1,\cdots,dx_n)\\
  			&\qquad\times \p(N_{s_r}-N_{s_{r-1}}=k_r,1\leq r\leq n)\p(N_{s_{n+1}}-N_{s_{n}}=k_{n+1})\\
  			=&\sum_{k_1,\cdots,k_n\geq 0}\int_{E^{n}}\bigg[\sum_{k_{n+1}\geq 0}\p(N_{s_{n+1}}-N_{s_{n}}=k_{n+1})P_{k_{n+1}}f(x_{n})\\
  			&\qquad\times g(x_1,\cdots,x_n)\mu P_{J_n}(dx_1,\cdots,dx_n)\bigg]\p(N_{s_r}-N_{s_{r-1}}=k_r,1\leq r\leq n)\\
  			=&\sum_{k_1,\cdots,k_n\geq 0}\int_{E^{n}}e^{-\alpha t} \sum_{k\geq 0}\frac{(\alpha t)^{k}}{k!} P_{k}f(x_{n})g(x_1,\cdots,x_n)\mu P_{J_n}(dx_1,\cdots,dx_n)\\
  			&\qquad\times \p(N_{s_r}-N_{s_{r-1}}=k_r,1\leq r\leq n)\\
  			=&\p[Q_tf(X_s)g(X_{s_1},\cdots,X_{s_n})].
  	\end{align*}
  	Till now, we have completed the proof.
  \end{proof}
  
  \begin{question}
  	For a $\r^d$-valued L\'evy process $(\Omega,\f,\f_t, X_t,t\geq 0)$, write its transition and convolution semigroups as $(P_t)_{t\geq 0}$ and $(\mu_t)_{t\geq 0}$ respectively. Then  the following statements are both equivalent: for any $(\f_t)_{t\geq 0}$ stopping time $T$ and $f\in\mathrm{b}\b(\r^d)$,
  	\begin{enumerate}
  		\item $\p[f(X_{T+t})\mathbbm{1}_{\{T<\infty\}}|\f_T]=P_tf(X_T)\mathbbm{1}_{\{T<\infty\}}$.
  		\item $\p[f(X_{T+t}-X_T)\mathbbm{1}_{\{T<\infty\}}|\f_T]=P_tf(0)\mathbbm{1}_{\{T<\infty\}}$.
  	\end{enumerate}
  \end{question}
  \begin{proof}
  	If the first formula is satisfied, write $g_y(x):=f(x-y)$, then
  	\begin{equation*}
  		\begin{split}
  			\p[&f(X_{T+t}-X_T)\mathbbm{1}_{\{T<\infty\}}|\f_T]\\
  			&=\p[f(X_{T+t}-y)\mathbbm{1}_{\{T<\infty\}}|\f_T]|_{y=X_T}=\p[g_y(X_{T+t})\mathbbm{1}_{\{T<\infty\}}|\f_T]|_{y=X_T}\\
  			&=P_tg_y(X_T)\mathbbm{1}_{\{T<\infty\}}|_{y=X_T}=\left.\mathbbm{1}_{\{T<\infty\}}\int g_y(x)P_t(X_T,dx)\right|_{y=X_T}\\
  			&=\left.\mathbbm{1}_{\{T<\infty\}}\int g_y(X_T+x)\mu_t(dx)\right|_{y=X_T}\\
  			&=\left.\mathbbm{1}_{\{T<\infty\}}\int f(X_T+x-y)\mu_t(dx)\right|_{y=X_T}=\mathbbm{1}_{\{T<\infty\}}\int f(x)\mu_t(dx)\\
  			&=\mathbbm{1}_{\{T<\infty\}}\int f(x)P_t(0,dx)=P_tf(0)\mathbbm{1}_{\{T<\infty\}}.
  		\end{split}
  	\end{equation*}
  	On the contrary, if the second formula is satisfied, write $h_y(x):=f(x+y)$, then 
  	\begin{equation*}
  		\begin{split}
  			\p[&f(X_{T+t})\mathbbm{1}_{\{T<\infty\}}|\f_T]\\
  			&=\p[f(X_{T+t}-X_T+X_T)\mathbbm{1}_{\{T<\infty\}}|\f_T]\\
  			&=\p[f(X_{T+t}-X_T+y)\mathbbm{1}_{\{T<\infty\}}|\f_T]|_{y=X_T}\\
  			&=\p[h_y(X_{T+t}-X_T)\mathbbm{1}_{\{T<\infty\}}|\f_T]|_{y=X_T}=P_th_y(0)\mathbbm{1}_{\{T<\infty\}}|_{y=X_T}\\
  			&=\left.\mathbbm{1}_{\{T<\infty\}}\int h_y(x)P_t(0,dx)\right|_{y=X_T}=\mathbbm{1}_{\{T<\infty\}}\int f(x+X_T)\mu_t(dx)\\
  			&=\mathbbm{1}_{\{T<\infty\}}\int f(x)P_t(X_T,dx)=P_tf(X_T)\mathbbm{1}_{\{T<\infty\}}.
  		\end{split}
  	\end{equation*}
  \end{proof}

  \section{Functional Analysis}
  \begin{question}
   Suppose $(M,d)$ is a compact metric space, then $C(M)$ with the supremum norm $||\cdot||$ is a Polish space.
  \end{question}
  \begin{proof}
   Obviously, $C(M)$ is a Banach space, so we just need to verify $C(M)$ is separable. In fact, since $M$ is a compact metric space, then $M$ is totally bounded and therefore is separable. Since $M$ is a separable metric space, then $M$ is a $C_2$ topology space, namely, $M$ has a countable topology basis, denoted as $\{U_n\}$. Consider $f_n(x)=d(x,U_n)$, we claim $\{f_n\}\subset C(M)$ can separate points in $M$. In fact, note $M$ is a Hausdorff topology space, then for each $x,y\in M$, exists $U_m,U_k\in\{U_n\}$ such that $x\in U_m,y\in U_k$ and $U_m\cap U_k=\emptyset$. Therefore $f_k(x)=d(x,U_k)=0$ and $f_k(y)=d(y,U_k)>0$. So we have proved $\{f_n\}$ can separate points in $M$. Let 
   \begin{equation*}
   	\g:=\left\{\sum_{n=1}^m q_nf_n:m\in\n,q_n\in\q,f_n\in\{f_n\} \right\}.
   \end{equation*}
   Next, we define
   \begin{equation*}
   	\f:=\left\{\prod_{n=1}^m g_n:m\in\n,g_n\in\g \right\}\cup\{1\},
   \end{equation*}
   where 1 denoted $f(x)\equiv 1$. We can find $\f$ is countable. Then $\bar{\f}$ has the unit element $1$ and can separate points in $E$, where the closure is in the sense of $||\cdot||$. Moreover, it can be found that $\bar{\f}$ is an algebra. The {\sffamily Stone-Weirestrass theorem} yields $\overline{\f}=C(M)$. So the assertion is proved.
  \end{proof}

  \begin{question}
   Suppose $N$ is a normal operator on Hilbert space $H$. For all $\Psi\in B(\sigma(N) )$, define $\Psi(N)$ as follows: for every $x,y\in H$,
   \begin{equation*}
    \langle \Psi(N)x,y\rangle=\int_{\sigma(N)}\Phi(z)m_{x,y}(dz),
   \end{equation*}
   where $m_{x,y}(dz)$ is the spectral measure of $N$ corresponding to $x,y$. Moreover, suppose $(\c,\b(\c),E)$ is the spectral family of $N$, and $\Phi:B_E(\c)\to \mathscr{N}(H)$ is the $*$-isometry corresponding to the spectral family $(\c,\b(\c),E)$. Then
   \begin{equation*}
    \Phi|_{B(\sigma(N) )}:\Psi\mapsto \Psi(N).
   \end{equation*}
  \end{question}
  \begin{proof}
   For all $\Psi\in B(\sigma(N) )$, define
   \begin{equation*}
    \widetilde{\Psi}(z):=
    \begin{cases}
     \Psi(z)&,\mbox{if }z\in\sigma(N)\\
     0&,\mbox{if }z\notin \sigma(N)\\
    \end{cases}.
   \end{equation*}
    Since the support of $E$ is $\sigma(N)$, then $\forall f\in B(\c),f=\widetilde{ f|_{\sigma(N)}},E-a.e.$. As a result, $\forall\Psi\in B(\sigma(N))$, we identify it with $\widetilde{\Psi}\in B(\c )$. When
    \begin{equation*}
     \Psi=\sum_{k=1}^n\alpha_k\i_{A_k}\in\mathscr{S}(\sigma(N) ) ,\quad \alpha_k \subset \c,A_k\in \b(\c)\cap\sigma(N),
    \end{equation*}
     For every $x,y\in H$, observe $(\sigma(N),\b(\c)\cap\sigma(N),m_{x,y} )=(\sigma(N),\b(\c)\cap\sigma(N),E)$, namely, $m_{x,y}=E|_{\b(\c)\cap \sigma(N)}$ and use the definition of $\Phi$ in \cite[Theorem 2.5.29]{ref1}, then
    \begin{equation*}
     \begin{split}
      \langle \Psi(N)x,y \rangle=&\int_{\sigma(N)}\Psi(z)m_{x,y}(dz)=\int_{\sigma(N)}\Psi(z)d\langle E(z)x,y\rangle\\
      =&\int_{\sigma(N)}\sum_{k=1}^n\alpha_k\i_{A_k}d\langle E(z)x,y\rangle=\sum_{k=1}^n\alpha_k\int_{\sigma(N)}\i_{A_k} d\langle E(z)x,y\rangle\\
      =&\sum_{k=1}^n\alpha_k \langle E(A_k)x,y\rangle= \left\langle \sum_{k=1}^n\alpha_k E(A_k)x,y\right \rangle\\
      =&\left\langle \Phi(\Psi ) x,y\right \rangle.
     \end{split}
    \end{equation*}
    Since this formula holds for every $x,y\in H$, then $\Phi(\Psi )=\Psi(N),\forall \Psi\in \mathscr{S}(\sigma(N) ) $. When $\Psi\in B(\sigma(N) )$, exist $\{\Psi_n \}\subset \mathscr{S}(\sigma(N) ),\st \Psi_n\to \Psi$ and $\exists M>0,\st ||\Psi_n||\leq M$. By the definition of $\Phi$ in \cite[Theorem 2.5.29]{ref1}, then $\Phi(\Psi ):=\lim_{n\to\infty}\Phi(\Psi_n )$. By \cite[Theorem 2.5.20 (iii)]{ref1}, we have
    \begin{equation*}
     \langle \Psi(N)x,y \rangle=\li{n\to\infty}\langle \Psi_n(N)x,y \rangle=\li{n\to\infty}\langle \Phi(\Psi_n )x,y \rangle=\langle \Phi(\Psi )x,y \rangle.
    \end{equation*}
    Since the above formula holds for every $x,y\in H$, then $\Phi(\Psi )=\Psi(N),\forall \Psi\in B(\sigma(N) ) $. Namely, we have proved
    \begin{equation*}
    \Phi|_{B(\sigma(N) )}:\Psi\mapsto \Psi(N).
   \end{equation*}
  \end{proof}

  \begin{question}
   Suppose $N$ is a normal operator on Hilbert space $H$. If $\varphi\in C(\sigma(N) ), \Psi\in B(\varphi(\sigma(N) ) )$, then
   \begin{equation*}
    (\Psi\circ\varphi )(N)=\Psi(\varphi(N) ).
   \end{equation*}
  \end{question}
  \begin{proof}
   We have proved: $\forall \varphi\in C(\sigma(N) ), \Psi\in C(\varphi(\sigma(N) ) )$, then $(\Psi\circ\varphi )(N)=\Psi(\varphi(N) )$. When $\Psi\in B(\varphi(\sigma(N) ) )$, we can find a sequence $\{\Psi_n\}\in C(\varphi(\sigma(N) ) )$, $\st \Psi_n\to \Psi,n\to\infty$ and $\exists M>0,\st ||\Psi_n||\leq M$. Hence $\Psi_n\circ\varphi\to \Psi\circ\varphi,||\Psi_n\circ\varphi||\leq M$. Notice $\varphi(N)$ is also a normal operator, by \cite[Theorem 2.5.20 (iii)]{ref1}, $\forall x\in H$,
   \begin{equation*}
    \begin{split}
     (\Psi\circ\varphi )(N)x&=\li{n\to\infty}(\Psi_n\circ\varphi )(N)x=\li{n\to\infty}\Psi_n(\varphi(N) )x\\
     &=\li{n\to\infty}\Psi_n(\varphi(N) )x=\Psi(\varphi(N) )x,
    \end{split}
   \end{equation*}
   which means $(\Psi\circ\varphi )(N)=\Psi(\varphi(N) )$.
  \end{proof}

  \begin{question}
   Suppose $\x$ is a Banach space and $T:\x\supset \mathscr{D}(T)\to \x$ is a closed operator. For every $\lambda\in\sigma_c(T)$, we have $(\lambda I-T )^{-1}$ is not bounded.
  \end{question}
  \begin{proof}
   Suppose $\exists \lambda\in\sigma_c(T),\st (\lambda I-T )^{-1}$ is bounded. By the definition of $\sigma_c(T)$, we have $R(\lambda I-T )\neq\x$ and $\overline{R(\lambda I-T )}=\x$. Hence $\forall y\in \x,\exists \{y_n\}\subset R(\lambda I-T )$, such that $y_n\to y,n\to\infty$. Write $x_n:=(\lambda I-T )^{-1}y_n$. Since $\{y_n\}$ is a Cauchy sequence and  $(\lambda I-T )^{-1}$ is bounded, $\{x_n\}$ is also a Cauchy sequence. Because $\x$ is complete, then $\exists x\in\x,\st x_n\to x,n\to\infty$. Till now, we have
   \begin{equation*}
    x_n\to x,(\lambda I-T )x_n=y_n\to y,n\to\infty
   \end{equation*}
   Moreover, $\lambda I-T$ is also a closed operator, whence $y=(\lambda I-T )x\in R(\lambda I-T )$. This implies $\x=R(\lambda I-T )$, which contradicts to the fact that $R(\lambda I-T )\neq\x$.
  \end{proof}

  \begin{question}
   Suppose $H$ is a Hilbert space and $u_k\rightharpoonup u$. Then $||u||\leq \varliminf_{k\to\infty}||u_k||$.
  \end{question}
  \begin{proof}
   Since for each $f\in H^*$, we have
   \begin{equation*}
    |f(x)|=\lim_{k\to\infty}|f(x_k)|\leq ||f||\varliminf_{k\to\infty}||x_k||.
   \end{equation*}
   Moreover, we have
   \begin{equation*}
    ||x||=\sup_{||f||=1}|f(x)|\leq \sup_{||f||=1}\(||f||\varliminf_{k\to\infty}||x_k||\)=\varliminf_{k\to\infty}||x_k||.
   \end{equation*}
   This is what we want.
  \end{proof}

  \begin{question}
  	Given a Banach space $(B,||\cdot||)$ and a separable subset $E\subset B$. Write $C$ for its countable dense subset of $E$ and $\widetilde{C}$ to be the set of rational multiples of elements of $C$. Then for any $x\in E$ and $\varepsilon>0$ there exists $y\in\widetilde{C}$ satisfying $||y||\leq ||x||$ and $||y-x||<\varepsilon$.
  \end{question}
  \begin{proof}
  	For the given $x\in E$ and $0<\delta<\varepsilon/2$, there exists $y'\in C$ such that $||y'-x||<\delta$. If $||y'||\leq ||x||$, then we let $y:=y'$ and all we need is completed. Otherwise, if $||y'||>||x||$, consider a rational $r\in\q\cap[0,1]$, then we have the following estimation
  	\begin{equation*}
  		\begin{split}
  			||ry'-x||&\leq ||ry'-rx||+||rx-x||\\
  			&=|r|\,||y'-x||+(1-r)||x||\\
  			&\leq \delta +(1-r)||x||.
  		\end{split}
  	\end{equation*}
  	Let $(1-r)||x||<\varepsilon/2$, then we have $r>1-\varepsilon/(2||x||)$. Choose a $r\in(1-\varepsilon/(2||x||),1)\cap\q$, and define $y:=ry'\in \widetilde{C}$. Hence the above estimation yields
  	\begin{equation*}
  		||y-x||\leq \delta +(1-r)||x||<\varepsilon/2+\varepsilon/2=\varepsilon.
  	\end{equation*}
  \end{proof}
  
  \begin{question}
  	Any finite-dimensional metric space $(E,d)$ is compact.
  \end{question}
  \begin{proof}
  	Since every norm on a finite-dimensional space is equivalent, everything boils down to the finite-dimensional Euclidean space. Obviously, it is compact.
  \end{proof}
  
  \begin{question}
  	Any finite-dimensional subspace of a metric space $(E,d)$ is closed.
  \end{question}
  \begin{proof}
  	Let $A\subset E$ be a finite-dimensional subspace with dimension $m<\infty$ and basis $\{x_k\}_{k=1}^ n$. Any element $x\in A$ can be uniquely expressed as $x=\sum_{k=1}^m \alpha_kx_k$. Given a convergent sequence $\{x^n\}_{n\in\z_+}$, say
  	\begin{equation*}
  		x^n=\sum_{k=1}^m \alpha_k^nx_k\to x\in E,\quad n\to\infty,
  	\end{equation*}
  	we want to prove $x\in A$. Firstly, the convergence of $\{x^n\}_{n\in\z_+}$ ensures that $\{x^n\}_{n\in\z_+}$ is bounded. The above QUESTION shows that $A$ is compact. So, there exists a subsequence $\{x^{n_\ell}\}_{\ell\in\z_+}$ such that $x^{n_\ell}\to y\in A,\ell\to\infty$. Note $x^{n_\ell}\to x\in E,\ell\to\infty$, the uniqueness of the limit yields that $x=y\in A$. In conclusion, $A$ is closed.
  \end{proof}

  \begin{question}
Let $H$ be a Hilbert space.

\begin{enumerate}
    \item[(a)] Prove that every finite-dimensional linear subspace $F \subset H$ is complete and hence closed.
    
    \item[(b)] Let $H_1,\dots,H_n \subset H$ be closed subspaces that are pairwise orthogonal. Prove that their direct sum
    $
    H_1 \oplus \cdots \oplus H_n
    $
    is closed.
    
    \item[(c)] Let $(X_i)_{i=1}^n$ be i.i.d.\ random variables and consider
    \begin{equation*}
         \mathcal H
    =
    \left\{
    \sum_{i=1}^n g_i(X_i)
    :
    g_i \in L^2(\mu_X)
    \right\}
    \subset L^2(\Omega,\mathbb P).
    \end{equation*}
    Show that $\mathcal H$ is a closed subspace of $L^2(\Omega,\mathbb P)$.
\end{enumerate}
\end{question}

\begin{proof}

\textbf{(a)} Let $F = \operatorname{span}\{e_1,\dots,e_k\}$ be finite-dimensional. Define $T : \mathbb R^k \to F$ by $ T(a_1,\dots,a_k) = \sum_{j=1}^k a_j e_j.$
Then $T$ is a linear isomorphism. Since all norms on finite-dimensional spaces are equivalent, $F$ is complete because $\mathbb R^k$ is complete. In a Hilbert space, a linear subspace is closed if and only if it is complete. Hence $F$ is closed.

\medskip
\noindent
\textbf{(b)} Let $H_1,\dots,H_n$ be closed subspaces such that $H_i \perp\!\!\!\perp H_j$ for $i\neq j$. Consider a Cauchy sequence
$
x_m = \sum_{i=1}^n h_{i,m},
$
where $h_{i,m} \in H_i$. By orthogonality,
\begin{equation*}
\|x_m - x_\ell\|^2
=
\left\|
\sum_{i=1}^n (h_{i,m}-h_{i,\ell})
\right\|^2
=
\sum_{i=1}^n
\|h_{i,m}-h_{i,\ell}\|^2.
\end{equation*}
Hence $(x_m)$ is Cauchy if and only if each $(h_{i,m})$ is Cauchy in $H_i$. Since each $H_i$ is closed and therefore complete, there exist $h_i \in H_i$ such that $h_{i,m}\to h_i$. Consequently,
$
x_m \to \sum_{i=1}^n h_i
\in
H_1 \oplus \cdots \oplus H_n.
$
Thus, the direct sum is complete and therefore closed.

\medskip
\noindent
\textbf{(c)}
Let $(X_i)_{i=1}^n$ be i.i.d.\ and define
$
\mathcal H_i
=
\{ g(X_i) : g \in L^2(\mu_X) \}.
$
For each $g$, write
\begin{equation*}
g(X_i)
=
\big(g(X_i)-\mathbb E[g(X_i)]\big)
+
\mathbb E[g(X_i)].
\end{equation*}
Define the centered subspace
$
\mathcal H_i^0
=
\{ g(X_i) : \mathbb E[g(X_i)] = 0,g \in L^2(\mu_X) \}.
$
Then
\begin{equation*}
\mathcal H
=
\operatorname{span}\{1\}
\oplus
\left(
\bigoplus_{i=1}^n
\mathcal H_i^0
\right).
\end{equation*}

Each $\mathcal H_i^0$ is closed because the map
$
T_i : L^2(\mu_X) \to L^2(\Omega,\mathbb P),
\;
T_i g = g(X_i),
$
is an isometry, so its image is complete.

Independence implies that for $i\neq j$ and centered elements $U\in\mathcal H_i^0$, $V\in\mathcal H_j^0$,
\begin{equation*}
\langle U,V\rangle
=
\mathbb E[UV]
=
\mathbb E[U]\mathbb E[V]
=
0.
\end{equation*}
Thus the $\mathcal H_i^0$ are pairwise orthogonal. Moreover, $\operatorname{span}\{1\}$ is finite-dimensional and orthogonal to each $\mathcal H_i^0$. By part (a) and part (b), the finite orthogonal direct sum above is closed. Therefore,
$
\mathcal H
$
is a closed subspace of $L^2(\Omega,\mathbb P)$.

\end{proof}
  
  \begin{question}
  	Given two subspaces $V,W$ of the Hilbert space $\x$. Suppose $V\perp W$, then $\overline{V\oplus W}=\overline{V}\oplus \overline{W}$. Besides, give an example to show $\overline{V\oplus W}=\overline{V}\oplus \overline{W}$ is not always true if $V\not\perp W$. 
  \end{question}
  \begin{proof}
  	Since $V\perp W$, the notation $V\oplus W$ is reasonable.
  	
  	Given a convergent sequence $x_n=v_n+w_n,n\geq 1$, where $(v_n)_{n\geq 1}\subset V,(w_n)_{n\geq 1}\subset W$ and $x_n\to x\in \x$, since $\overline{V}$ and $\overline{W}$ are both closed subspace of $\x$, we can get the following decomposition:
  	\begin{equation*}
  		x=P_{\overline{V}}x+P_{\overline{V}^\perp}x,
  	\end{equation*} 
  	where $P_{\overline{V}}$ is the projection operator of $\overline{V}$. Since $V\perp W$, the continuity of the inner product yields $\overline{ V}\perp\overline{ W}$. Therefore $\overline{W}\subset \overline{V}^\perp$. Since $v_n-P_{\overline{V}}x\in \overline{V}$ and $w_n-P_{\overline{V}^\perp}x \in \overline{V}^\perp$, we have
  	\begin{equation*}
  		\begin{split}
  			||x_n-x||&=||(v_n-P_{\overline{V}}x)+(w_n-P_{\overline{V}^\perp}x)||\\
  			&=||v_n-P_{\overline{V}}x||+||w_n-P_{\overline{V}^\perp}x||\to 0,\quad n\to\infty.
  		\end{split}
  	\end{equation*}
  	This is equivalent to
  	\begin{equation*}
  		||v_n-P_{\overline{V}}x||\to0,\quad ||w_n-P_{\overline{V}^\perp}x||\to 0,\quad n\to\infty.
  	\end{equation*}
  	So, $(v_n)_{n\geq 1}$ and $(w_n)_{n\geq 1}$ are both convergent sequences and thus 
  	\begin{equation*}
  		v_n\to P_{\overline{V}}x\in \overline{V},\quad w_n\to P_{\overline{V}^\perp}x\in \overline{W},\Longrightarrow x\in \overline{V}\oplus \overline{W}.
  	\end{equation*}
  	So, $\overline{V\oplus W}\subset\overline{V}\oplus \overline{W}$. 
  	
  	On the contrary, for any $v\in \overline{V}$ and $w\in\overline{W}$, there exist $(v_n)_{n\geq 1}\subset V$ and $(w_n)_{n\geq 1}\subset W$ such that $v_n\to v$ and $w_n\to w$. Therefore, $V\oplus W\ni v_n+w_n\to v+w$. This implies $\overline{V}\oplus \overline{W}\subset \overline{V\oplus W}$.
  	
  	{\bfseries Example:} Consider three subspaces of $\ell^2$: $V=\mathrm{span}\{e_1\},U=\mathrm{span}\{e_1+e_n\}_{n\geq 1}$ and $W=\overline{U}$. It is easy to show
  	\begin{equation*}
  		\begin{split}
  			V&=\{(a_1,0,\cdots):a_1\in\r\},\\
  			W&=\left\{\(\sum_{k\geq 2}a_k,a_2,a_3,\cdots\):\sum_{k\geq 2}a_k<\infty,\sum_{k\geq 2}a_k^2<\infty\right\}.
  		\end{split}
  	\end{equation*}
  	And from the expression above, we can find $V\cap W=\{0\}$. So the notation $V\oplus W$ is reasonable. Obviously, $V$ is closed and $U$ is closed by its definition, so $\overline{V}\oplus \overline{W}=V\oplus W$. However, consider
  	\begin{equation*}
  		w_n=\(\sum_{k= 2}^n \frac{1}{k} ,\frac{1}{2} ,\cdots,\frac{1}{n},0,\cdots \)\in W,\quad v_n=\(-\sum_{k= 2}^n \frac{1}{k},0,\cdots\)\in V.
  	\end{equation*}
  	Noting $\sum_k \frac{1}{k^2}<\infty$, we have 
  	\begin{equation*}
  		\begin{split}
  			&v_n+w_n=\(0,\frac{1}{2} ,\cdots,\frac{1}{n},0,\cdots\)\in V\oplus W\\
  			\text{ and }&v_n+w_n\to x=\(0,\frac{1}{2},\frac{1}{3},\cdots\)\in \ell^2.
  		\end{split}
  	\end{equation*}
  	Since $\sum_k \frac{1}{k}=\infty$, there is no $v\in V$ and $w\in W$ such that $x=v+w$. This means $\overline{V}\oplus \overline{W}=V\oplus W\subsetneq\overline{V\oplus W}$.
  	
  	This partial inclusion is because $V\not\perp W$. 
  \end{proof}
  
  \begin{question}
  	Given a subspace $A$ in the Hilbert space $H$. Suppose $\Lambda\subset A$ is a subspace of $A$. Suppose $\mathscr{A},\b,\lambda$ are three subspaces of the Hilbert space $\mathscr{H}$. Then
  	\begin{enumerate}
  		\item $\mathrm{Proj}_{\Lambda^\perp }(A)\subset A.$
  		\item If $\Lambda$ is a $q$-replicating space $\Lambda=\lambda\times\cdots\times\lambda=\lambda ^q\subset \mathscr{H}^q=H\ (q\geq 2)$, we have 
  	\begin{equation*}
  		\Lambda^\perp=\lambda^\perp \times \cdots\times \lambda^\perp.
  	\end{equation*}
  		\item $\mathrm{Proj}_{\Lambda ^\perp }(A)= A\cap \Lambda^\perp.$ Furthermore, if $A=\mathscr{A}\times \cdots\times \mathscr{A}=\mathscr{A}^q,\Lambda=\lambda\times\cdots\times\lambda=\lambda ^q$, we have
  		\begin{equation*}
  			\mathrm{Proj}_{\Lambda ^\perp }(A)=(\a\cap\lambda)\times\cdots\times(\a\cap\lambda)=(\a\cap\lambda)^q,
  		\end{equation*}
  		which is also a $q$-replicating space.
  	\end{enumerate} 
  	
  \end{question}
  \begin{proof}
  	\begin{enumerate}
  		\item Note $\mathrm{Proj}_{\Lambda^\perp }(A)=\{x-\mathrm{Proj}_{\Lambda}(x):\forall x\in A\}$ and $\mathrm{Proj}_{\Lambda}(x)\in \Lambda \subset A$, the fact that $A$ is closed under addition yields
  	\begin{equation*}
  		x-\mathrm{Proj}_{\Lambda}(x)\in A,\quad \forall x\in A.
  	\end{equation*}
  	Therefore, $\mathrm{Proj}_{\Lambda^\perp }(A)\subset A.$
  		\item We can find that $x=(x_1,\cdots,x_q)\in \Lambda^\perp\subset \mathscr{H}^q=H$ iff 
  	\begin{equation*}
  		\langle x,y\rangle=\sum_{k=1}^q \langle x_k,y_k\rangle=0,\quad \forall y=(y_1,\cdots,y_q)\in \Lambda.\tag{$*$}
  	\end{equation*}
  	Letting $x_{-k}\equiv 0\in \mathscr{H}^{q-1}$ and $x_k$ to vary,  we can find $(*)$ is equivalent to $x_k\perp \lambda$ for any $k=1,\cdots,q$. So, 
  	\begin{equation*}
  		\begin{split}
  			\Lambda^\perp&=\{x=(x_1,\cdots,x_q)\in H:x_k\in \mathscr{A}^\perp,\forall k=1,\cdots,q \}\\
  			&=\lambda^\perp \times \cdots\times \lambda^\perp.
  		\end{split}
  	\end{equation*}
  		\item Since we have proved in 1. that $\mathrm{Proj}_{\Lambda ^\perp }(A)\subset A$ and obviously $\mathrm{Proj}_{\Lambda ^\perp }(A)\subset \Lambda^\perp$, we have $\mathrm{Proj}_{\Lambda ^\perp }(A)\subset A\cap \Lambda^\perp$. On the contrary, since $A\cap \Lambda^\perp\subset \Lambda^\perp$, obviously we have $A\cap \Lambda^\perp=\mathrm{Proj}_{\Lambda ^\perp }(A\cap \Lambda^\perp)\subset \mathrm{Proj}_{\Lambda ^\perp }(A)$. Thus, $\mathrm{Proj}_{\Lambda ^\perp }(A)= A\cap \Lambda^\perp$. The second assertion is obvious.
  	\end{enumerate}
  \end{proof}

\begin{question}[Variance identity iff the norm is Hilbertian]
Let $V$ be a finite-dimensional real or complex vector space (e.g.\ $V=\mathbb{R}^{m\times n}$ with matrices viewed as vectors), equipped with a norm $\|\cdot\|$. The following are equivalent:
\begin{itemize}
\item[(i)] There exists an inner product $\langle\cdot,\cdot\rangle$ on $V$ such that $\|x\|^2=\langle x,x\rangle$ for all $x\in V$.
\item[(ii)] For every $V$-valued random element $X$ with $\mathbb{E}\|X\|^2<\infty$,
\begin{equation*}
\mathbb{E}\|X-\mathbb{E}X\|^2=\mathbb{E}\|X\|^2-\|\mathbb{E}X\|^2.
\end{equation*}
\end{itemize}
\end{question}

\begin{proof}
\textbf{(i) $\Rightarrow$ (ii).}
Write $m:=\mathbb{E}X$. Using the inner product that induces the norm and the bilinearity (sesquilinearity in the complex case) of $\langle\cdot,\cdot\rangle$,
\begin{equation*}
\|X-m\|^2=\langle X-m,X-m\rangle=\langle X,X\rangle-2\operatorname{Re}\langle X,m\rangle+\langle m,m\rangle.
\end{equation*}
Taking expectations,
\begin{equation*}
\mathbb{E}\|X-m\|^2=\mathbb{E}\langle X,X\rangle-2\operatorname{Re}\langle \mathbb{E}X,m\rangle+\langle m,m\rangle
=\mathbb{E}\|X\|^2-\|m\|^2,
\end{equation*}
which is the desired identity.

\medskip
\textbf{(ii) $\Rightarrow$ (i).}
We derive the parallelogram law. Fix arbitrary $u,v\in V$ and let $X$ take the two values $u$ and $v$ with probability $1/2$ each. Then $\mathbb{E}X=(u+v)/2$ and
\begin{equation*}
\mathbb{E}\|X-\mathbb{E}X\|^2=\frac12\Big\|\frac{u-v}{2}\Big\|^2+\frac12\Big\|\frac{v-u}{2}\Big\|^2
=\frac14\|u-v\|^2,
\end{equation*}
while
\begin{equation*}
\mathbb{E}\|X\|^2-\|\mathbb{E}X\|^2=\frac12(\|u\|^2+\|v\|^2)-\Big\|\frac{u+v}{2}\Big\|^2
=\frac12(\|u\|^2+\|v\|^2)-\frac14\|u+v\|^2.
\end{equation*}
By assumption (ii) these two quantities are equal, hence for all $u,v$,
\begin{equation*}
\|u+v\|^2+\|u-v\|^2=2(\|u\|^2+\|v\|^2),
\end{equation*}
the parallelogram identity. By the Jordan--von Neumann theorem, a norm satisfies the parallelogram law if and only if it is induced by an inner product. Concretely, one may define, for real $V$,
\begin{equation*}
\langle u,v\rangle:=\frac14\big(\|u+v\|^2-\|u-v\|^2\big),
\end{equation*}
and for complex $V$,
\begin{equation*}
\langle u,v\rangle:=\frac14\sum_{k=0}^{3} i^k\,\|u+i^k v\|^2,
\end{equation*}
which are inner products satisfying $\|x\|^2=\langle x,x\rangle$ for all $x\in V$. This proves (i).
\end{proof}

\begin{question}
Let $(X,\|\cdot\|)$ be a Banach space and let $K\subset X$.  
Show that the following are equivalent:

\medskip
\noindent (a) $K$ is relatively compact (i.e.\ $\overline K$ is compact in $X$).

 \noindent (b) $K$ is bounded and for every $\delta>0$ there exists a finite-dimensional
  subspace $V\subset X$ such that
  \begin{equation*}
    K \subset V^\delta
    =
    \{x\in X:\operatorname{dist}(x,V)<\delta\}.
  \end{equation*}
  In particular, $K$ is compact if and only if it is closed and satisfies (b).
\end{question}

\begin{proof}
(a) $\Rightarrow$ (b).
If $K$ is relatively compact, then it is totally bounded.
Fix $\delta>0$. By total boundedness, there exist points
$x_1,\dots,x_m\in X$ such that
\begin{equation*}
  K \subset \bigcup_{i=1}^m B(x_i,\delta),
\end{equation*}
where $B(x_i,\delta)=\{x:\|x-x_i\|<\delta\}$.
Let
\begin{equation*}
  V = \operatorname{span}\{x_1,\dots,x_m\},
\end{equation*}
which is finite-dimensional.
For any $x\in K$ there exists $i$ with $\|x-x_i\|<\delta$, and since $x_i\in V$ we have
\begin{equation*}
  \operatorname{dist}(x,V)\le \|x-x_i\|<\delta.
\end{equation*}
Thus $K\subset V^\delta$. Boundedness of $K$ follows from relative compactness.

\medskip
(b) $\Rightarrow$ (a).
Assume that $K$ is bounded and satisfies (b).
Let $R>0$ be such that $K\subset RB_X$, where $B_X=\{x:\|x\|\le1\}$.
Fix $\varepsilon>0$ and choose $\delta=\varepsilon/2$.
By (b) there exists a finite-dimensional subspace $V\subset X$ such that
\begin{equation*}
  K \subset V + \frac{\varepsilon}{2}B_X.
\end{equation*}
For each $x\in K$, choose $v_x\in V$ with $\|x-v_x\|<\varepsilon/2$.
Then
\begin{equation*}
  \|v_x\|\le \|x\|+\|x-v_x\|\le R+\frac{\varepsilon}{2},
\end{equation*}
so all such $v_x$ lie in the bounded set
\begin{equation*}
  A := V\cap \bigl(R+\tfrac{\varepsilon}{2}\bigr)B_X.
\end{equation*}
Since $V$ is finite-dimensional, $A$ is totally bounded.
Hence there exist $a_1,\dots,a_n\in A$ such that
\begin{equation*}
  A \subset \bigcup_{k=1}^n B(a_k,\varepsilon/2).
\end{equation*}
For any $x\in K$, choose $k$ with $\|v_x-a_k\|<\varepsilon/2$.
Then
\begin{equation*}
  \|x-a_k\|\le \|x-v_x\|+\|v_x-a_k\|
  < \frac{\varepsilon}{2}+\frac{\varepsilon}{2}
  = \varepsilon.
\end{equation*}
Thus
\begin{equation*}
  K \subset \bigcup_{k=1}^n B(a_k,\varepsilon),
\end{equation*}
showing that $K$ is totally bounded.
Since $X$ is complete, the totally boundedness is equivalent to relative compactness. Hence, $K$ is relatively compact.

Finally, $K$ is compact if and only if it is relatively compact and closed.
\end{proof}



  \section{Partial Differential Equations}
  \begin{question}
   Suppose $\{u_n\}\subset L^\infty(\r^n)$ is bounded, then exists $u\in L^\infty(\r^n)$ such that $u_n\rightharpoonup v \text{ in }L^2_\mathrm{loc}(\r^n)$.
  \end{question}
  \begin{proof}
   Observe $L^\infty(\r^n)=[L^1(\r^n)]^*$ and $L^1(\r^n)$ is separable. Since $\{u_n\}\subset L^\infty(\r^n)$ is bounded then exists $u\in L^\infty(\r^n)$, such that $w^*$-$\lim_{n\to\infty}u_n=u$. Namely, for each $v\in L^1(\r^n)$, we have
   \begin{equation*}
    \int_{\r^n}u_nvdx\to \int_{\r^n}uvdx.
   \end{equation*}
   For any $V\subset\subset \r^n$ and each $w\in L^2(V)$, by H\"older inequality, we know $w\in L^1(V)$ as well. Moreover, we can extend $w$ to $\r^n$ as follows:
   \begin{equation*}
    v(x)=
    \begin{cases}
     w(x)&,\mbox{if }x\in V\\
     0&,\mbox{if }x\in V^c.
    \end{cases}
   \end{equation*}
   Then $v\in L^1(\r^n)$. Thus, the above limit means
   \begin{equation*}
    \int_{V}u_nwdx =\int_{\r^n}u_nvdx\to \int_{\r^n}uvdx=\int_Vuwdx.
   \end{equation*}
   This shows $u_n\rightharpoonup v \text{ in }L^2_\mathrm{loc}(\r^n)$.
  \end{proof}

  \begin{question}
   Assume $U\subset \r^n$ is a bounded open set. Prove $W^{1,p}(U)\subset\subset L^p(U)$ for each $1\leq p\leq \infty$.
  \end{question}
  \begin{proof}
   \begin{enumerate}
    \item When $1\leq p<n$, this is exactly the conclusion of {\sffamily Rellich-Kondrachov Compactness Theorem}.
    \item When $p=n$. Observe $U$ is bounded, then the H\"older inequality tells us $u\in W^{1,n}(U)$ implies $u\in W^{1,p}(U)$ for each $1\leq p<n$. Observe $p*=\frac{np}{n-p}\uparrow\infty(p\uparrow n)$. So we can find a $p\in[1,n)$ such that $p^*>n$. Utilize {\sffamily Rellich-Kondrachov Compactness Theorem}, we know $W^{1,p}(U)\subset\subset L^p(U)$.
    \item When $n<p\leq\infty$.
    \begin{enumerate}
     \item \textcolor{blue}{\sffamily Way 1.} \textcolor{blue}{\itshape However, this way is suitable only for $n<p<\infty$.} It is similar to the previous situation. For each $u\in W^{1,p}(U)$, we can find $1\leq q<n$ such that $u\in W^{1,q}(U)$. Besides, observe $q*=\frac{nq}{n-q}\uparrow\infty(q\uparrow n)$, so we assume the given $q$ is closed enough to $n$ such that $q^*>p$. Thus the {\sffamily Rellich-Kondrachov Compactness Theorem} gives $W^{1,q}(U)\subset\subset L^p(U)$. However, by {\sffamily H\"older inequality}, we know the embedding mapping $F:W^{1,p}(U)\to W^{1,q}(U),u\mapsto u$ is continuous. So we have $W^{1,p}(U)\subset\subset L^p(U)$.
     \item \textcolor{blue}{\sffamily Way 2.} Given a bounded sequence $\{u_m\}\subset W^{1,p}(U)$. By {\sffamily Morrey's inequality}, we have
     \begin{equation*}
      ||u_m^*||_{C(\bar{U})}+ [u_m^*]_{C^{0,1-n/p}(\bar{U})}\leq C||u_m^*||_{W^{1,p}(U)},
     \end{equation*}
     where $u^*_m$ is the continuous version of $u_m$. Next we verify $u_m^*$ is uniformly and equicontinuous. In fact, the above inequality implies
     \begin{equation*}
      \sup_{n\geq 1}||u_m^*||_{C(\bar{U})}\leq C\sup_{m\geq 1}||u_m^*||_{W^{1,p}(U)}=C'
     \end{equation*}
     and
     \begin{equation*}
      \sup_{m\geq 1}[u_m^*]_{C^{0,1-n/p}(\bar{U})}\leq C\sup_{m\geq 1}||u_m^*||_{W^{1,p}(U)}=C'.
     \end{equation*}
     The first inequality show $\{u_m^*\}$ is uniformly bounded. The second inequality shows
     \begin{equation*}
      \sup_{m\geq 1}|u_m^*(x)-u_m^*(y)|\leq C'||x-y||^{1-n/p},
     \end{equation*}
     which means $\{u^*_m\}$ is equicontinuous. So exists $u\in C(\bar{U} )$ and subsequence $\{m_k\}_{k\geq 1}$ such that $||u_{m_k}^*-u||_{C(\bar{U} )}\to0,k\to\infty$. However, in the bounded open set, $u_{m_k}^*\to u$ in $C(\bar{U} )$ implies $u_{m_k}^*\to u$ in $L^p(U)$. This assertion holds for each $n<p\leq \infty$. Thus $W^{1,p}(U)\subset\subset L^p(U)$.
    \end{enumerate}
   \end{enumerate}
  \end{proof}

  \begin{question}[A corollary of Weak maximum principle for $c\geq 0$] Assume $u\in C^2(U)\cap C(\bar{U})$ and $Lu=0$ in $U$, where $L=-a^{ij}\partial_{ij}+b^i\partial_i+c$ is the elliptic partial differential operators of second order. Assume $c\geq 0$ in $U$, then we have
  \begin{equation*}
   \max_{\bar{U}}|u|=\max_{\partial U}|u|.
  \end{equation*}
  \end{question}
  \begin{proof}
   Since $Lu=0$, then we have
   \begin{equation*}
    \max_{\bar{U}}u\leq \max_{\partial U}u^+\leq \max_{\partial U}|u| ,\quad \max_{\bar{U}}(-u)\leq \max_{\partial U}u^-\leq \max_{\partial U}|u|.
   \end{equation*}
   Hence,
   \begin{equation*}
    \max_{\bar{U}}|u|=\max_{\bar{U}}u\vee \max_{\bar{U}}(-u) \leq \max_{\partial U}|u|
   \end{equation*}
   The converse inequality is obvious.
  \end{proof}

  \begin{question}
   Assume $u\in C^2(U)\cap C(\bar{U})$ and $U$ is connected, open and bounded. Write $M:=\max_{\bar{U}}u$ and assume $C:=\{x\in U|u(x)=M\}\neq\emptyset$. If $u\not\equiv M$ in $U$, so $V:=\{x\in U|u(x)<M\}\neq \emptyset$. Then exists a point $y\in V$ satisfying $d(y,C)<d(y,\partial U )$.
  \end{question}
  \begin{proof}
    Choose and fix $x\in C$ and $z\in V$. Since $U$ is connected and bounded, then exists a path $\gamma:[0,1]\to U$ satisfying $\gamma(0)=x,\gamma(1)=z$ and $\gamma[0,1]$ is a compact subset of $\r^n$. Observe $r:= d(\gamma[0,1],\partial U)>0$, we can find finite open balls $\{B(x_i,r/3 )\}_{i=1}^n$ covering $\gamma[0,1]$, where $x_i\in \gamma[0,1]$ and $B(x_i,r/3 )\cap B(x_{i+1},r/3 )\neq\emptyset$ for each $i$. There must exists $B(x_k,r/3)$ such that $B(x_k,r/3)\cap C\neq\emptyset$ and $B(x_k,r/3)\cap V\neq\emptyset$. We can choose $y\in B(x_k,r/3)\cap V$. So
    \begin{equation*}
     d(y,\partial U)\geq d(\gamma[0,1],\partial U )- d(y,\gamma[0,1] )\geq r-\frac{r}{3}=\frac{2r}{3}>d(y,C).
    \end{equation*}
  \end{proof}

  \begin{question}
   We write $\widetilde{\partial}_i$ for the weak derivative operator with respect to the $i$-th component. Suppose $u\in W^{1,p}(U)$, then $\widetilde{\partial}_{ij}^2u:=\widetilde{\partial}_i(\widetilde{\partial}_j u)$ exists iff $\widetilde{\partial}_{ji}^2u$ exists. In this case, we still have $\widetilde{\partial}_{ij}^2u=\widetilde{\partial}_{ji}^2u$, where $U\subset\r^n$ is an open subset and $1\leq i,j\leq n$.
  \end{question}
  \begin{proof}
   Suppose $\widetilde{\partial}_{ij}^2u$ exists, then for each $\phi\in C^\infty_c(U)$, we have
   \begin{equation*}
    \int_U \widetilde{\partial}_{ij}^2u \phi dx=-\int_U\widetilde{\partial}_{j}u\partial_i\phi dx=\int_Uu \partial_{ji}^2\phi dx=\int_Uu \partial_{ij}^2\phi dx.
   \end{equation*}
   By the definition of the weak derivative, we have $\widetilde{\partial}_{ij}^2u=\widetilde{\partial}_{ji}^2u$.
  \end{proof}

  \begin{question}
   Suppose $u\in W^{k,n}(U)$ and $k> 2$, where $U\subset \r^n$. Then for each $|\alpha|\leq k-2$, the weak derivative $D^\alpha u=\partial^\alpha u$ (the strong derivative). Therefore $u\in C^{k-2,\gamma}(\bar{U})$ for each $0<\gamma<1$.
  \end{question}
  \begin{proof}
   Since $u\in W^{k,n}(U)$, then the {\sffamily Gagliardo-Nirenberg-Sobolev inequality} and the {\sffamily H\"older inequality} shows: for each $n\leq q<\infty$, exists $p\in[1,n)$ such that for each $|\beta|\leq k-1$,
   \begin{equation*}
    ||D^\beta u||_{L^q(U)}\leq C||D^\beta u||_{W^{1,p}(U)}  \leq C||D^\beta u||_{W^{1,n}(U)}\leq C||u||_{W^{k,n}(U)}.
   \end{equation*}
    Thus we have $u\in W^{k-1,q}(U)$ for each $n\leq q<\infty$. As a consequence, the {\sffamily Morrey's inequality} shows  $D^\alpha u\in C^{0,1-n/q}(\bar{U})$ for all $n<q<\infty$ and $|\alpha |\leq k-2$. Next we show $D^\alpha u=\partial^\alpha u$. For each $x\in U$, select a close ball $B(x,\epsilon )\subset U$. Consider the standard modifier $\eta$ and the convolution $D^\alpha u^\varepsilon=\eta_\varepsilon*D^\alpha u$ in $U_\varepsilon$. When $\varepsilon$ small enough, we know $B(x,\epsilon )\subset U_\varepsilon$. Observe $D^\alpha u\in C^{0,1-n/q}(\bar{U})\subset C(\bar{ U})$, so $D^\alpha u^\varepsilon\to D^\alpha u$ uniformly in $B(x,\epsilon )$.  So for any sequence $\{\varepsilon_m \}$ satisfying $\varepsilon_m\to 0$, we have $D^\alpha u^{\varepsilon_{m}}\to D^\alpha u$ uniformly in $B(x,\epsilon )$. Based on the previous discussion, $\{u^{\varepsilon_m}\}$ is a Cauchy sequence in the Banach space $C^{k-2}(B(x,\epsilon ) )$, so we can find $v\in C^{k-2}(B(x,\epsilon ) )$ such that $u^{\varepsilon_m}\to v$ in $C^{k-2}(B(x,\epsilon ) )$. The uniqueness of the limit shows $u=v\in C^{k-2}(B(x,\epsilon ) )$. Thus $D^\alpha u=\partial^\alpha u$ for each $|\alpha |\leq k-2$.
  \end{proof}

  \begin{question}
   Suppose $U\subset\r^n$ and $\partial U\cap \r_n^+\neq \emptyset$. Assume $f\in C^1(\bar{U})$, then for each $x\in\partial U\cap\r_n^+$, we have
   \begin{equation*}
    \partial_{x_n+} f(x):=\lim_{h\to0}\frac{f(x+he_n)-f(x)}{h}=\lim_{h\to 0} \partial_{x_n}f(x+he_n).
   \end{equation*}
  \end{question}
  \begin{proof}
   Write $x=(x',x_n)\in\r^n$ and fix $x'$. Then define $g(x_n):=f(x',x_n)$. We know $g(x_n)$ is continuous on $[0,h]$ and differentiable in $(0,h)$. So the {\sffamily Differential mean value theorem} yields exists $\theta(h)\in(0,1)$ such that
   \begin{equation*}
    \begin{split}
     f(x+he_n)-&f(x)=g(h)-g(0)\\
     =&g'(\theta(h) h)h=\partial_{x_n}f(x+\theta(h) h e_n )h.
    \end{split}
   \end{equation*}
   Observe $f\in C^1(\bar{U})$, then we have
   \begin{equation*}
    \begin{split}
     \partial_{x_n+} f(x)&=\lim_{h\to0}\frac{f(x+he_n)-f(x)}{h}=\lim_{h\to0}\frac{\partial_{x_n}f(x+\theta(h) h e_n )h}{h}\\
     &=\lim_{h\to0}\partial_{x_n}f(x+\theta(h) h e_n )=\lim_{h\to 0} \partial_{x_n}f(x+he_n).
    \end{split}
   \end{equation*}
  \end{proof}

  \begin{question}
   Suppose $U\subset\r^n$ is an open, bounded subset. Consider $u(x)= C,\ \mathrm{a.e.}\  x\in  U$, where $C$ is a constant. Then we have $u\in W^{1,p}(U)$ for each $1\leq p< \infty$ and $Tu=C,\ \mathrm{a.e.}\  x\in \partial U$ provided $1\leq p<\infty$.
  \end{question}
  \begin{proof}
   We can easily verify $u\in W^{1,p}(U)$. Consider another function $v(x)\equiv C,\forall x\in\bar{U}$. Since $v\in W^{1,p}(U)\cap C(\bar{U})$ for each $1\leq p<\infty$, then the {\sffamily Trace theorem} shows $Tv=v|_{\partial U}=C$. Thus
   \begin{equation*}
    ||C-Tu||_{L^p(\partial U )}=||Tv-Tu||_{L^p(\partial U )}\leq C'||v-u||_{W^{1,p}(U)}=0.
   \end{equation*}
   This yields $Tu=C$, a.e. $x\in\partial U$.
  \end{proof}

\section{Econometrics}
\begin{question}
	Consider a univariate regression model
	\begin{equation*}
		y_i=\alpha+\beta x_i+\varepsilon_i.
	\end{equation*}
	If we want to do a weighted OLS with  weight $\sqrt{w_i}$:
	\begin{equation*}
		\sqrt{w_i}y_i=\sqrt{w_i}\alpha+\sqrt{w_i}\beta x_i+\sqrt{w_i}\varepsilon_i,
	\end{equation*}
	we can write it as
	\begin{equation*}
		\begin{split}
			V^{1/2}Y&=\alpha V^{1/2}\bm{1}+ V^{1/2}X\gamma +V^{1/2}\varepsilon\\
			\Longleftrightarrow \quad \widetilde{ Y}&=\alpha\widetilde{ \bm{1}}+\beta\widetilde{ x} +\widetilde{ \varepsilon}.
		\end{split}
	\end{equation*}
	Then, the weighted OLS estimator has the following expression:
	\begin{equation*}
		\widetilde{\gamma}=\frac{\widehat{\mathrm{Cov}}_w(\widetilde{ x},\widetilde{ Y})}{\widehat{\var}_w(\widetilde{ x}) },
	\end{equation*}
	where the subscript $w$ means taking the sample covariance, variance, and mean with respect to the measure $w=(w_i)_{i=1}^n$.
\end{question}
\begin{proof}
	Write $\widetilde{X}:=(\widetilde{ \bm{1}},\widetilde{x})$. Then, the weighted OLS estimate $\widetilde{\gamma}$ is
	\begin{equation*}
		\begin{split}
			\widetilde{\gamma}&=(\widetilde{X}'X)^{-1}\widetilde{X}'\widetilde{Y}\\
			&=\[\begin{pmatrix}
				\widetilde{ \bm{1}}'\\
				\widetilde{x}'
			\end{pmatrix}
			\begin{pmatrix}
				\widetilde{ \bm{1}} & \widetilde{x}
			\end{pmatrix}\]^{-1}\begin{pmatrix}
				\widetilde{ \bm{1}}'\\
				\widetilde{x}'\
			\end{pmatrix}\widetilde{Y}\\\
			&=\[\begin{pmatrix}
				\bm{1}'V^{1/2}\\
				x'V^{1/2}
			\end{pmatrix}
			\begin{pmatrix}
				V^{1/2}\bm{1} & V^{1/2}x
			\end{pmatrix}\]^{-1}\begin{pmatrix}
				\bm{1}'V^{1/2}\\
				x'V^{1/2}
			\end{pmatrix}V^{1/2}Y\\
			&=\begin{pmatrix}
				\bm{1}' V\bm{1} & \bm{1}' Vx\\
				x' V\bm{1} & x'Vx
			\end{pmatrix}^{-1}
			\begin{pmatrix}
				\bm{1}' VY\\
				x' VY
			\end{pmatrix}\\
			&=\begin{pmatrix}
				\sum_i w_i & \sum_iw_ix_i\\
				\sum_iw_ix_i & \sum_iw_ix_i^2
			\end{pmatrix}^{-1}
			\begin{pmatrix}
				\sum_iw_iy_i\\
				\sum_iw_ix_i y_i
			\end{pmatrix}\\
			&=\frac{1}{\sum_i w_i\sum_iw_ix_i^2-\(\sum_iw_ix_i\)^2} \begin{pmatrix}
				\sum_iw_ix_i^2 & -\sum_iw_ix_i\\
				-\sum_iw_ix_i & \sum_i w_i
			\end{pmatrix}
			\begin{pmatrix}
				\sum_iw_iy_i\\
				\sum_iw_ix_i y_i
			\end{pmatrix}.
		\end{split}
	\end{equation*}
	Therefore,
	\begin{equation*}
		\widetilde{\beta}=\frac{\sum_i w_i\sum_iw_ix_i y_i-\sum_iw_ix_i\sum_iw_iy_i}{\sum_i w_i\sum_iw_ix_i^2-\(\sum_iw_ix_i\)^2}.
	\end{equation*}
	Note 
	\begin{equation*}
		\begin{split}
			\widehat{\mathrm{Cov}}_w(\widetilde{ x},\widetilde{ Y})&=\frac{1}{\sum_iw_i}\sum_{i=1}^n w_i\(x_i-\frac{1}{\sum_i w_i} \sum_{j=1}^nw_jx_j\)\(y_i-\frac{1}{\sum_i w_i} \sum_{j=1}^nw_jy_j\)\\
			&=\frac{1}{\sum_iw_i}\[\sum_{i=1}^n w_ix_iy_i-\sum_{i=1}^n w_i x_i\(\frac{1}{\sum_i w_i} \sum_{j=1}^nw_jy_j\)\right.\\
			&\quad \left.-\sum_{i=1}^n w_i y_i\(\frac{1}{\sum_i w_i} \sum_{j=1}^nw_jx_j\)+\frac{1}{\sum_i w_i} \(\sum_{j=1}^nw_jx_j\)\(\sum_{j=1}^nw_jy_j\)\]\\
			&=\frac{1}{\sum_iw_i}\sum_{i=1}^n w_ix_iy_i-\(\frac{1}{\sum_i w_i}\)^2\(\sum_{i=1}^n w_i x_i\)\( \sum_{j=1}^nw_jy_j\).
		\end{split}
	\end{equation*}
	So,
	\begin{equation*}
		\widetilde{\beta}=\frac{\(\sum_iw_i\)^2\widehat{\mathrm{Cov}}_w(\widetilde{ x},\widetilde{ Y})}{\(\sum_iw_i\)^2\widehat{\var}_w(\widetilde{ x})}=\frac{\widehat{\mathrm{Cov}}_w(\widetilde{ x},\widetilde{ Y})}{\widehat{\var}_w(\widetilde{ x}) }.
	\end{equation*}
\end{proof}

\begin{question}
	Given a parameter space $\Theta$ and observable distribution family $\mathcal{P}=\{P_\theta:\theta\in\Theta\}$. Define $G:\Theta\to \mathcal{P},\theta\mapsto P_\theta$. Then
	\begin{itemize}
		\item Model is point identified at true value $\theta_0$ if
		\begin{equation*}
			P_\theta=P_{\theta_0}\Longrightarrow \theta=\theta_0,
		\end{equation*}
		which is equivalent to $G^{-1}(\{P_{\theta_0}\})=\{\theta\}$
		\item Model is globally identified if 
		\begin{equation*}
			P_\theta=P_{\theta'}\Longrightarrow \theta=\theta',\quad \forall \theta,\theta'\in\Theta,
		\end{equation*} 
		which is equivalent to saying $G$ is an injection.
		\item Model is locally identified if $G$ is an injection in some neighborhood of the true value $\theta_0$.
		\item Model is partially identified if $G^{-1}(\{P_{\theta_0}\})\subsetneq \Theta$ and $G^{-1}(\{P_{\theta_0}\})$ is not a singleton.
		\item For a functional parameter $\psi()$
	\end{itemize}
	When there is a functional $F$ such that $\theta=F(P_\theta)$, $\theta$ is globally identified.
\end{question}

\begin{question}[When can regime/path and state representations be interchanged?]
Consider a panel indexed by $i$ and $t\in\{1,\dots,T\}$ with an absorbing treatment.
Let $G_i\in\{1,\dots,T\}\cup\{\infty\}$ denote the first treatment time (with $G_i=\infty$ meaning never treated). 

\medskip
\noindent
\textbf{(A) Fixed-adoption (two-path) case.}
Assume there are only two regimes: $g=\infty$ (never treated) and $g=T_0$ (treatment starts at a common $T_0$).
Show that the following are equivalent:
\begin{itemize}
\item[(i)] (Regime model with no anticipation) There exists $\{Y_{it}(g): g\in\{1,\dots,T\}\cup\{\infty\}\}$ be the regime/path-based potential outcomes, where the \emph{Consistency axiom} holds,
\begin{equation*}
Y_{it}=Y_{it}(G_i)\quad \forall i,t,
\end{equation*}
and \emph{No anticipation assumption} holds,
\begin{equation*}
Y_{it}(T_0)=Y_{it}(\infty)\quad \forall i,\ \forall t<T_0 .
\end{equation*}
\item[(ii)] (State model / switching equation) There exist state-based potential outcomes
$\{Y_{it}(0),Y_{it}(1)\}$ and binary treatment indicator $D_{it}\in\{0,1\}$ such that
\begin{equation*}
Y_{it}=\widetilde Y_{it}(0)(1-D_{it})+\widetilde Y_{it}(1)D_{it}\quad \forall i,t .
\end{equation*}
\end{itemize}

\medskip
\noindent
\textbf{(B) Staggered adoption case.}
Assume there are at least two distinct finite adoption times $g\neq g'$.
Show that, in general, there do \emph{not} exist state-based potential outcomes
$\{Y_{it}(0),Y_{it}(1)\}$ satisfying
\begin{equation*}
Y_{it}=\widetilde Y_{it}(0)(1-D_{it})+\widetilde Y_{it}(1)D_{it}\quad \forall i,t,
\end{equation*}
provided the following additional restriction holds:
\begin{equation*}
\boxed{
Y_{it}(g)=Y_{it}(g') \quad \text{for all } t \text{ and all } g,g'\le t .
}
\end{equation*}
In words, conditional on being treated at time $t$, outcomes must be invariant to the
treatment start date.
\end{question}

\begin{proof}
\textbf{(A)}
The equivalence in the fixed-adoption (two-path) case is established as follows.
Define
\begin{equation*}
\widetilde Y_{it}(0):= Y_{it}(\infty),\quad \forall i,t,
\qquad
\widetilde Y_{it}(1):=
\begin{cases}
Y_{it}(\infty), & t<T_0,\\
Y_{it}(T_0), & t\ge T_0 .
\end{cases}
\end{equation*}
Under no anticipation, this definition is well-posed.
Since $t\geq \infty$ is not true, define the realized treatment status
\begin{equation*}
	D_{it}:=\mathbbm{1}_{[t\ge G_i]}=\mathbbm{1}_{[t\geq T_0,G_i=T_0]}=\mathbbm{1}_{[t\ge T_0]}\mathbbm{1}_{[G_i=T_0]}.
\end{equation*}
So
\begin{equation*}
	\begin{split}
		&\widetilde Y_{it}(0)(1-D_{it})+\widetilde Y_{it}(1)D_{it}\\
		=&Y_{it}(\infty)(1-\mathbbm{1}_{[t\ge T_0]}\mathbbm{1}_{[G_i=T_0]})+[Y_{it}(\infty)\mathbbm{1}_{[t<T_0]}+Y_{it}(T_0)\mathbbm{1}_{[t\geq T_0]}]\mathbbm{1}_{[t\ge T_0]}\mathbbm{1}_{[G_i=T_0]}\\
		=&Y_{it}(\infty)(\mathbbm{1}_{[t< T_0,G_i=T_0]}+\mathbbm{1}_{[G_i=\infty]})+Y_{it}(T_0)\mathbbm{1}_{[t\ge T_0,G_i=T_0]}\\
		=&Y_{it}(G_i)(\mathbbm{1}_{[t< T_0,G_i=T_0]}+\mathbbm{1}_{[G_i=\infty]})+Y_{it}(G_i)\mathbbm{1}_{[t\ge T_0,G_i=T_0]}\\
		=&Y_{it}(G_i),
	\end{split}
\end{equation*}
where the third equality follows from the \emph{Consistency axiom} and \emph{No anticipation assumption}. 

\medskip
Conversely, given a switching equation, i.e., there exists $\{Y_{it}(0),Y_{it}(1)\}$ and binary treatment indicator $D_{it}\in\{0,1\}$ such that
\begin{equation*}
Y_{it}=\widetilde Y_{it}(0)(1-D_{it})+\widetilde Y_{it}(1)D_{it}:=\widetilde Y_{it}(D_{it})\quad \forall i,t .
\end{equation*}
Fix $i$, there are only two kinds of realization of treatment paths $D_i:=(D_{it})_{t=1}^T$: $D_{it}\equiv 0$ and $D_{it}=\mathbbm{1}_{[t\geq T_0]}$. Furthermore, 
\begin{equation*}
	  \{D_{it}\equiv 0\}=\{G_i=\infty\} \text{ and }\{D_{it}=\mathbbm{1}_{[t\geq T_0]}\}=\{G_i=T_0\}. 
\end{equation*} 
So, one can reconstruct two kinds of path-based potential outcomes associated with the above two paths as follows:
\begin{equation*}
Y(G_i)=Y_{it}(\infty):=\widetilde Y_{it}(0)=\widetilde Y_{it}(D_{it})=Y_{it}\text{ for }D_{it}\equiv 0,
\end{equation*}
and
\begin{equation*}
	Y(G_i)=Y_{it}(T_0):=\widetilde Y_{it}(\mathbbm{1}_{[t\geq T_0]})=\widetilde Y_{it}(D_{it})=Y_{it}\text{ for }D_{it}=\mathbbm{1}_{[t\geq T_0]}.
\end{equation*}
The above displays have already shown consistency. When $t<T_0$,
\begin{equation*}
	Y_{it}(T_0)=\widetilde Y_{it}(\mathbbm{1}_{[t\geq T_0]})=\widetilde Y_{it}(0)=Y_{it}(\infty)
\end{equation*}
So, no anticipation assumption holds.

Hence, the regime/path and state representations are equivalent when there are exactly two
treatment paths.

\medskip
\noindent
\textbf{(B)}
Consider now staggered adoption with two distinct finite start dates $g\neq g'$.
Fix a time $t\ge \max\{g,g'\}$.
Then $D_{it}=\mathbbm{1}_{[t\ge G_i]}=1$ for units with $G_i=g$ and for units with $G_i=g'$.
If a state-based representation existed, we would have for all such units
\begin{equation*}
Y_{it}(g)=\widetilde  Y_{it}(1)=Y_{it}(g'),
\end{equation*}
because both share the same current state $D_{it}=1$.
Therefore, a \emph{necessary} condition for the switching equation to hold is
\begin{equation*}
Y_{it}(g)=Y_{it}(g') \quad \text{for all } g,g'\le t .
\end{equation*}



In general staggered DiD settings, treatment effects are allowed to depend on
event time $k=t-g$, so the invariance condition fails.
Consequently, $D_{it}$ is not a sufficient index for the potential outcomes,
and the regime/path representation $\{Y_{it}(g)\}$ cannot be reduced to a
state-based switching equation.
\end{proof}

\begin{question}
Let $(\Omega,\mathcal F,\mathbb P)$ be a probability space. Let $Y(0),Y(1),D,X$ be random variables and define the $\sigma$-algebras
\begin{equation*}
\mathcal Y := \sigma\bigl(Y(0),Y(1)\bigr),\qquad
\mathcal X := \sigma(X),\qquad
\mathcal D := \sigma(D).
\end{equation*}
Assume the \emph{complete randomization} (joint independence) condition
\begin{equation*}
\mathcal D \perp (\mathcal Y \vee \mathcal X),
\qquad
\mathcal Y \vee \mathcal X := \sigma\bigl(Y(0),Y(1),X\bigr).
\end{equation*}
Prove that the \emph{conditional independence} (unconfoundedness) condition holds:
\begin{equation*}
\mathcal Y \perp \mathcal D \mid \mathcal X.
\end{equation*}
\end{question}

\begin{proof}
Assume $\mathcal D \perp (\mathcal Y\vee \mathcal X)$. Let $U$ be bounded and $\mathcal Y$-measurable, and let $V$ be bounded and $\mathcal D$-measurable. Since $V\perp Y\vee \mathcal X$,  by the tower property,
\begin{equation*}
\begin{split}
	\mathbb E[UV\mid \mathcal X]
&=
\mathbb E\!\left[\mathbb E[UV\mid \mathcal Y\vee \mathcal X]\mid \mathcal X\right]=\mathbb E\!\left[U\mathbb E[V\mid \mathcal Y\vee \mathcal X]\mid \mathcal X\right]\\
&=\mathbb E\!\left[U\mathbb E[V]\mid \mathcal X\right]=\mathbb E[V]E[U\mid \mathcal X]=\mathbb E[V]\mathbb E[U]
\end{split}
\end{equation*}
Since $U$ and $V$ were arbitrary bounded $\mathcal Y$- and $\mathcal D$-measurable random variables, this establishes $\mathcal Y \perp \mathcal D \mid \mathcal X$.
\end{proof}

\section{Mathemetical Statistics}

\begin{question}
Let $X_1,\ldots,X_n$ be i.i.d.\ with distribution $U(\theta,\theta+|\theta|)$, where $\theta\neq 0$.
Find the MLE of $\theta$.
\end{question}

\begin{proof}
Write the order statistics $X_{(1)}=\min_i X_i$ and $X_{(n)}=\max_i X_i$.
The density is
\begin{equation*}
f(x;\theta)=\frac{1}{|\theta|}\,\mathbbm{1}_{[\theta\le x\le \theta+|\theta|]},
\end{equation*}
hence the likelihood is
\begin{equation*}
L(\theta)=\frac{1}{|\theta|^n}\,
\mathbbm{1}_{[\theta\le X_{(1)}\ \text{ and }\ X_{(n)}\le \theta+|\theta|]}.
\end{equation*}

\paragraph{Case 1: $\theta>0$.}
Then the support is $[\theta,2\theta]$ and feasibility requires
\begin{equation*}
\theta\le X_{(1)}\quad\text{and}\quad 2\theta\ge X_{(n)}
\ \Longleftrightarrow\ 
\theta\in\big[X_{(n)}/2,\; X_{(1)}\big].
\end{equation*}
On this interval $L(\theta)=\theta^{-n}$ is strictly decreasing in $\theta$, so the maximizer is
\begin{equation*}
\hat\theta_{\,+}= \frac{X_{(n)}}{2},
\qquad
L\!\left(\frac{X_{(n)}}{2}\right)=\left(\frac{2}{X_{(n)}}\right)^{\!n}.
\end{equation*}

\paragraph{Case 2: $\theta<0$.}
Then the support is $[\theta,0]$ and feasibility requires
\begin{equation*}
\theta\le X_{(1)}\quad\text{and}\quad X_{(n)}\le 0
\ \Longleftrightarrow\ 
\theta\in(-\infty,\,X_{(1)}].
\end{equation*}
Here $L(\theta)=(-\theta)^{-n}$ is strictly increasing in $\theta$, so the maximizer is
\begin{equation*}
\hat\theta_{\,-}= X_{(1)},
\qquad
L\!\left(X_{(1)}\right)=\frac{1}{(-X_{(1)})^n}.
\end{equation*}

Since the sample must come from one of the two feasible regimes, the MLE is
\begin{equation*}
\hat\theta_{\mathrm{MLE}}=
\begin{cases}
X_{(1)}, & \text{if } X_{(n)}\le 0 \quad(\text{negative--support case}),\\[4pt]
\dfrac{X_{(n)}}{2}, & \text{if } \dfrac{X_{(n)}}{2}\le X_{(1)} \quad(\text{positive--support case}).
\end{cases}
\end{equation*}
Equivalently, when both candidates are admissible, select the one with the larger likelihood:
\begin{equation*}
\hat\theta_{\mathrm{MLE}}=
\begin{cases}
X_{(1)}, & \text{if } \left(\dfrac{2}{X_{(n)}}\right)^{\!n}\le \dfrac{1}{(-X_{(1)})^n},\\[6pt]
\dfrac{X_{(n)}}{2}, & \text{if } \left(\dfrac{2}{X_{(n)}}\right)^{\!n}> \dfrac{1}{(-X_{(1)})^n}.
\end{cases}
\end{equation*}
\end{proof}


\section{High-Dimensional Statistics}
\begin{question}[Hoeffding's lemma]
Let $Z$ be a random variable such that $Z \in [a,b]$ almost surely. Then for all $\lambda \in \mathbb{R}$,
\begin{equation*}
\mathbb{E}\!\left[\exp\big(\lambda(Z-\mathbb{E}Z)\big)\right]
\le
\exp\!\left(\frac{\lambda^2 (b-a)^2}{8}\right).
\end{equation*}
\end{question}

\begin{proof}
Let $\mu := \mathbb{E}Z$ and $d := b-a$. By convexity of the exponential function, for every $z \in [a,b]$,
\begin{equation*}
e^{\lambda z}
\le
\frac{b-z}{d} e^{\lambda a}
+
\frac{z-a}{d} e^{\lambda b}.
\end{equation*}
Taking expectations and multiplying by $e^{-\lambda\mu}$ yields
\begin{equation*}
\mathbb{E}\!\left[e^{\lambda(Z-\mu)}\right]
\le
(1-\theta)e^{-\lambda \theta d}
+
\theta e^{\lambda (1-\theta)d},
\end{equation*}
where
\begin{equation*}
\theta := \frac{\mu-a}{d} \in [0,1].
\end{equation*}

Define
\begin{equation*}
g(\lambda)
:=
(1-\theta)e^{-\lambda \theta d}
+
\theta e^{\lambda (1-\theta)d},
\qquad
\phi(\lambda)
:=
\log g(\lambda).
\end{equation*}
We compute derivatives explicitly. Set
\begin{equation*}
a_1 := -\theta d,
\qquad
a_2 := (1-\theta)d.
\end{equation*}
Writing $g(\lambda)=u_1(\lambda)+u_2(\lambda)$ with
\begin{equation*}
u_1=(1-\theta)e^{a_1\lambda},
\qquad
u_2=\theta e^{a_2\lambda},
\end{equation*}
we observe
\begin{equation*}
u_i' = a_i u_i,
\qquad
u_i'' = a_i^2 u_i.
\end{equation*}
Hence
\begin{equation*}
g' = a_1 u_1 + a_2 u_2,
\qquad
g'' = a_1^2 u_1 + a_2^2 u_2.
\end{equation*}
Using $\phi'' = \frac{g''}{g} - \left(\frac{g'}{g}\right)^2$, define weights
\begin{equation*}
p_i(\lambda) := \frac{u_i(\lambda)}{g(\lambda)},
\qquad
p_1+p_2=1,
\end{equation*}
and obtain
\begin{equation*}
\phi''(\lambda)
=
\big(p_1 a_1^2 + p_2 a_2^2\big)
-
\big(p_1 a_1 + p_2 a_2\big)^2
=
p_1 p_2 (a_1-a_2)^2.
\end{equation*}
Since $a_1-a_2=-d$, this simplifies to
\begin{equation*}
\phi''(\lambda)
=
d^2 p_1(\lambda)p_2(\lambda).
\end{equation*}
Because $p_1p_2 \le 1/4$ for all $p_1 \in [0,1]$, we conclude
\begin{equation*}
\phi''(\lambda)
\le
\frac{d^2}{4}
\quad \text{for all } \lambda.
\end{equation*}

Finally, note that $\phi(0)=0$ and $\phi'(0)=0$. By Taylor's theorem with remainder, for some $\xi$ between $0$ and $\lambda$,
\begin{equation*}
\phi(\lambda)
=
\frac{1}{2}\phi''(\xi)\lambda^2
\le
\frac{\lambda^2 d^2}{8}.
\end{equation*}
Exponentiating gives
\begin{equation*}
\mathbb{E}\!\left[\exp\big(\lambda(Z-\mu)\big)\right]
\le
\exp\!\left(\frac{\lambda^2 (b-a)^2}{8}\right),
\end{equation*}
which completes the proof.
\end{proof}

\begin{question}[Finite-class Rademacher maximal inequality]
Let $\{\epsilon_i\}_{i=1}^n$ be i.i.d. Rademacher random variables, and let $\{a_{ij}\}_{i\le n, j\le p}$ be fixed real numbers. Define
\begin{equation*}
Z_j := \frac{1}{n}\sum_{i=1}^n \epsilon_i a_{ij}, \quad j=1,\dots,p.
\end{equation*}
Show that
\begin{equation*}
\mathbb{E}_\epsilon \max_{1\le j\le p} |Z_j|
\le
\sqrt{2\log(2p)}\;
\max_{1\le j\le p}
\left(\frac{1}{n^2}\sum_{i=1}^n a_{ij}^2\right)^{1/2}.
\end{equation*}
\end{question}

\begin{proof}
For fixed $j$, write
\begin{equation*}
Z_j = \frac{1}{n}\sum_{i=1}^n \epsilon_i a_{ij}.
\end{equation*}
Using independence and the bound $\mathbb{E}e^{t\epsilon x}=\cosh(tx)\le e^{t^2x^2/2}$, we obtain
\begin{equation*}
\mathbb{E}_\epsilon e^{t Z_j}
=
\prod_{i=1}^n \mathbb{E}_\epsilon e^{t\epsilon_i a_{ij}/n}
\le
\exp\bigg\{\frac{t^2}{2n^2}\sum_{i=1}^n a_{ij}^2\bigg\}.
\end{equation*}
Hence $Z_j$ is sub-Gaussian with variance proxy
\begin{equation*}
\sigma_j^2 := \frac{1}{n^2}\sum_{i=1}^n a_{ij}^2.
\end{equation*}

Define $2p$ variables
\begin{equation*}
Y_1,\dots,Y_{2p} = Z_1,\dots,Z_p,-Z_1,\dots,-Z_p,
\end{equation*}
so that
\begin{equation*}
\max_{1\le j\le p} |Z_j| = \max_{1\le k\le 2p} Y_k.
\end{equation*}
Let $\sigma := \max_{1\le j\le p} \sigma_j.$
Then for all $k$,
\begin{equation*}
\mathbb{E}_\epsilon e^{t Y_k} \le e^{t^2 \sigma^2/2}.
\end{equation*}

For any $t>0$, by monotonicity of the exponential function,
\begin{equation*}
e^{t \max_k Y_k} \le \sum_{k=1}^{2p} e^{t Y_k}.
\end{equation*}
Taking expectations,
\begin{equation*}
\mathbb{E}_\epsilon e^{t \max_k Y_k}
\le
\sum_{k=1}^{2p} \mathbb{E}_\epsilon e^{t Y_k}
\le
2p \, e^{t^2\sigma^2/2}.
\end{equation*}
By Jensen's inequality,
\begin{equation*}
e^{t\, \mathbb{E}_\epsilon \max_k Y_k}
\le
\mathbb{E}_\epsilon e^{t \max_k Y_k}
\le
2p \, e^{t^2\sigma^2/2}.
\end{equation*}
Taking logarithms,
\begin{equation*}
\mathbb{E}_\epsilon \max_k Y_k
\le
\frac{\log(2p)}{t} + \frac{t\sigma^2}{2}.
\end{equation*}
Optimizing over $t>0$ yields
\begin{equation*}
t = \frac{\sqrt{2\log(2p)}}{\sigma},
\end{equation*}
and hence
\begin{equation*}
\mathbb{E}_\epsilon \max_k Y_k
\le
\sigma \sqrt{2\log(2p)}.
\end{equation*}

Substituting $\sigma = \max_j \sigma_j$ gives
\begin{equation*}
\mathbb{E}_\epsilon \max_{1\le j\le p} |Z_j|
\le
\sqrt{2\log(2p)}\;
\max_{1\le j\le p}
\left(\frac{1}{n^2}\sum_{i=1}^n a_{ij}^2\right)^{1/2}.
\end{equation*}
\end{proof}

\section{Lasso}
\begin{question}[Uniqueness of the Lasso fitted values]
Consider the Lasso estimator
\begin{equation*}
\hat\beta \in 
\arg\min_{\beta\in\mathbb{R}^p}
\left\{
\frac{1}{n}\|Y-X\beta\|_2^2+\lambda\|\beta\|_1
\right\},
\end{equation*}
where $X\in\mathbb{R}^{n\times p}$ and $\lambda>0$.

\begin{enumerate}
\item Show that the minimizer $\hat\beta$ need not be unique, in particular when $p>n$. 
\item Show that although $\hat\beta$ may not be unique, the fitted value
\begin{equation*}
\hat u := X\hat\beta
\end{equation*}
is unique.
\item Provide an example to show even $p>n$, it is still possible that $\hat\beta$ is unique.
\end{enumerate}
\end{question}

\begin{proof}
First note that the objective function
\begin{equation*}
F(\beta)
=
\frac{1}{n}\|Y-X\beta\|_2^2+\lambda\|\beta\|_1
\end{equation*}
is convex because both the quadratic loss and the $\ell_1$ norm are convex.

\medskip
\noindent
\textbf{1.}
If $p>n$, the matrix $X^\top X$ is singular. Hence the quadratic loss
\begin{equation*}
\frac{1}{n}\|Y-X\beta\|_2^2
\end{equation*}
is not strictly convex in $\beta$. Consequently, the objective function $F(\beta)$ is not strictly convex. A convex function that is not strictly convex may admit multiple minimizers. Therefore, the Lasso solution $\hat\beta$ need not be unique.

\medskip
\noindent
\textbf{2.}
Suppose $\hat\beta^{(1)}$ and $\hat\beta^{(2)}$ are two minimizers of the Lasso objective. Let
\begin{equation*}
u_1 := X\hat\beta^{(1)}, 
\qquad
u_2 := X\hat\beta^{(2)}.
\end{equation*}

Assume for contradiction that $u_1\neq u_2$. For any $t\in(0,1)$ define
\begin{equation*}
\beta_t
=
t\hat\beta^{(1)}+(1-t)\hat\beta^{(2)},
\qquad
u_t = X\beta_t
=
t u_1+(1-t)u_2.
\end{equation*}

Because the squared loss is strictly convex in its argument, we have
\begin{equation*}
\|Y-u_t\|_2^2
<
t\|Y-u_1\|_2^2+(1-t)\|Y-u_2\|_2^2.
\end{equation*}

Moreover, the $\ell_1$ norm is convex, hence
\begin{equation*}
\|\beta_t\|_1
\le
t\|\hat\beta^{(1)}\|_1+(1-t)\|\hat\beta^{(2)}\|_1.
\end{equation*}

Combining the two inequalities gives
\begin{equation*}
F(\beta_t)
<
tF(\hat\beta^{(1)})+(1-t)F(\hat\beta^{(2)}).
\end{equation*}

Since both $\hat\beta^{(1)}$ and $\hat\beta^{(2)}$ are minimizers, the right-hand side equals the minimal value of $F$. This contradicts optimality. Therefore we must have
\begin{equation*}
X\hat\beta^{(1)} = X\hat\beta^{(2)}.
\end{equation*}

Hence, the fitted value $X\hat\beta$ is unique even though the coefficient vector $\hat\beta$ may not be unique.

\medskip
\noindent
\textbf{3.}
Even when $p>n$, the Lasso minimizer can still be unique. 
We construct an explicit example with $n=2$ and $p=3$.

Let
\begin{equation*}
X=
\begin{pmatrix}
1 & 0 & 1/2\\
0 & 1 & 1/2
\end{pmatrix},
\qquad
Y=
\begin{pmatrix}
2\\
0
\end{pmatrix},
\qquad
\lambda=1 .
\end{equation*}
Consider the candidate solution
\begin{equation*}
\hat\beta=(1,0,0).
\end{equation*}

For the Lasso objective
\begin{equation*}
F(\beta)
=
\frac{1}{n}\|Y-X\beta\|_2^2+\lambda\|\beta\|_1,
\end{equation*}
the Karush--Kuhn--Tucker (KKT) conditions state that $\hat\beta$ is optimal 
if and only if there exists $z\in\mathbb{R}^p$ such that
\begin{equation*}
\frac{2}{n}X^\top(Y-X\hat\beta)=\lambda z,
\end{equation*}
where
\begin{equation*}
z_j\in
\begin{cases}
\{\operatorname{sign}(\hat\beta_j)\}, & \hat\beta_j\neq0,\\
[-1,1], & \hat\beta_j=0 .
\end{cases}
\end{equation*}

Compute the fitted value
\begin{equation*}
X\hat\beta=
\begin{pmatrix}
1\\
0
\end{pmatrix},
\end{equation*}
so the residual is
\begin{equation*}
r:=Y-X\hat\beta
=
\begin{pmatrix}
1\\
0
\end{pmatrix}.
\end{equation*}
Since $n=2$ and $\lambda=1$, the KKT condition becomes
\begin{equation*}
X^\top r=z .
\end{equation*}
Direct computation gives
\begin{equation*}
X^\top r
=
\begin{pmatrix}
1\\
0\\
1/2
\end{pmatrix}.
\end{equation*}
Hence, we may choose
\begin{equation*}
z=X^\top r=
\begin{pmatrix}
1\\
0\\
1/2
\end{pmatrix}.
\end{equation*}
Now check the subgradient conditions:

\begin{itemize}
\item For $j=1$, $\hat\beta_1=1\neq0$, and
\begin{equation*}
z_1=1=\operatorname{sign}(\hat\beta_1).
\end{equation*}

\item For $j=2$, $\hat\beta_2=0$, and
\begin{equation*}
z_2=0\in[-1,1].
\end{equation*}

\item For $j=3$, $\hat\beta_3=0$, and
\begin{equation*}
z_3=\tfrac12\in[-1,1].
\end{equation*}
\end{itemize}
Thus, the KKT conditions are satisfied. Combined with the fact that the objection is convex, $\hat\beta$ is a Lasso minimizer.

Next, we show the Uniqueness. The active set is
\begin{equation*}
A=\{1\}.
\end{equation*}
First note that
\begin{equation*}
\left|\frac{2}{n}X_2^\top r\right|=0<\lambda,
\qquad
\left|\frac{2}{n}X_3^\top r\right|=\tfrac12<\lambda .
\end{equation*}
Hence, for every inactive coordinate $j\notin A$ we have the strict inequality
\begin{equation*}
\left|\frac{2}{n}X_j^\top r\right|<\lambda .
\end{equation*}

Because we have shown that the fitted value $\hat u=X\hat\beta$ is unique, the residual $r$ is also unique. So, $\frac{2}{n}X_j^\top r$ is unique.
Therefore, any minimizer must satisfy
\begin{equation*}
\beta_2=\beta_3=0.
\end{equation*}
Thus, every minimizer must have the form
\begin{equation*}
\beta=(b,0,0).
\end{equation*}
Using the fact that all minimizers have the same fitted value again, we must have
\begin{equation*}
X\beta=X\hat\beta\Longrightarrow bX_1 = X_1.
\end{equation*}
Because $X_1\neq0$, it follows that $b=1$. Hence
\begin{equation*}
\beta=\hat\beta.
\end{equation*}
Therefore, the Lasso minimizer is unique even though $p>n$.
\end{proof}

\section{Semiparametric Model}
\begin{question}
	Given a model $\mathcal{P}$ and its dominating measure $\nu$. Assume $P_0\in \mathcal{P}$ is the true model and observation $Z\sim P_0$. Define its $n$-dimensional regular submodels as follows
	\begin{equation*}
		 \mathcal{P}^{(n)}:=\{P_\gamma=P_{(\gamma_1,\cdots,\gamma_n)} :\gamma_i\in(-\varepsilon_i,\varepsilon_i),1\leq i\leq n\}\subset \mathcal{P}.
	\end{equation*}
	Write $s_{\gamma,\mathcal{P}^{(n)}}$ as the score of submodels $\mathcal{P}^{(n)}$ in the sense of DQM. Then the following definition of tangent space does not depend on $n$:
	\begin{equation*}
			\mathcal{T}_n(\mathcal{P}):=\overline{\left\{h's_{\gamma,\mathcal{P}^{(n)}}(Z) :\mathcal{P}^{(n)}\subset \mathcal{P},h\in\r^n\right\}}^{L^2_0}.
		\end{equation*}
\end{question}
\begin{proof}
	We just need to prove $\ct_1(\cp)=\ct_n(\p),\forall n\geq 2$. 
	
	For any $s_{\ell,\overline{\mathcal{P}}^{(1)}}(Z)$, where 
	\begin{equation*}
		\overline{\mathcal{P}}^{(1)}=\{\overline{ P}_\ell:\ell\in(-\varepsilon_1,\varepsilon_1)\}\subset\cp,
	\end{equation*}
	we can define the $n$-dimensional submodel 
	\begin{equation*}
		\mathcal{P}^{(n)}:=\{P_\gamma=P_{(\ell,\gamma_2,\cdots,\gamma_n)}= \overline{ P}_\ell:\gamma_i\in(-\varepsilon_i,\varepsilon_i),1\leq i\leq n\}\subset \mathcal{P},
	\end{equation*}
	which means that the model $\cp^{(n)}$ will not change with respect to parameters $(\gamma_2,\cdots,\gamma_n)$. We now show $\cp^{(n)}$ is regular, i.e., $\cp^{(n)}$ is DQM at $\gamma=0$ with score $s_{\gamma,\cp^{(n)}}(Z):=(s_{\ell,\overline{\mathcal{P}}^{(1)}}(Z),0,\cdots,0)'$:
	\begin{equation*}
		\begin{split}
			&||\sqrt{p_{\gamma,\cp^{(n)}}}-\sqrt{p_0}-\frac{1}{2}\sqrt{p_0}\gamma' s_{\gamma,\cp^{(n)}}||_{L^2_0(\nu)}\\
			&\qquad = ||\sqrt{p_{\ell,\overline{ \cp}^{(1)}}}-\sqrt{p_0}-\frac{1}{2}\sqrt{p_0}\ell s_{\gamma,\overline{\cp}^{(1)}}||_{L^2_0(\nu)}=o_{L^2_0(\nu)}(|\ell|),
		\end{split}
	\end{equation*}
	where the last equality is because $\overline{\mathcal{P}}^{(1)}$ is DQM at $\ell=0$ with score $s_{\ell,\overline{\mathcal{P}}^{(1)}}(Z)$. Since $|\ell|\leq ||\gamma||$, we have $o_{L^2_0(\nu)}(|\ell|)=o_{L^2_0(\nu)}(||\gamma||)$. Therefore, for any $h=(h_1,\cdots,h_n)\in\r^n$, $h_1s_{\gamma,\overline{\cp}^{(1)}}=h' s_{\gamma,\cp^{(n)}}$, which implies $\cp^{(1)}\subset \cp^{(n)}$.
	
	On the contrary, for any $\widetilde{h}=(\widetilde{h}_1,\cdots,\widetilde{h}_n)$ and $\widetilde{h}'s_{\gamma,\overline{\cp}^{(n)}}(Z)$, where
	\begin{equation*}
		\overline{\cp}^{(n)}=\{\overline{P}_{\gamma}=\overline{P}_{(\gamma_1,\cdots,\gamma_n)}:\gamma_i\in(-\varepsilon_i,\varepsilon_i),1\leq i\leq n\},
	\end{equation*}
	We can let $M>0$ big enough such that $h:=\widetilde{h}/M$ satisfies $h=\widetilde{h}_i/M\in (-\varepsilon_i,\varepsilon_i)$. The model
	\begin{equation*}
		\widetilde{\cp}^{(n)}:=\{\overline{P}_{h}=\overline{P}_{(h_1,\cdots,h_n)}:h_i\in(-\varepsilon_i,\varepsilon_i),1\leq i\leq n\}
	\end{equation*}
	Then
	we can define the one-dimensional submodel
	\begin{equation*}
		\cp^{(1)}:=\{P_\ell=\overline{P}_{\ell h}=\overline{P}_{(\ell h_1,\cdots,\ell h_n)}:\ell\in (-1,1)\}. 
	\end{equation*}
	Write 
	\begin{equation*}
		\widetilde{\cp}^{(n)}:=\{\overline{P}_{\ell h}=\overline{P}_{(\ell h_1,\cdots,\ell h_n)}:\ell h_i\in(-h_i,h_i),1\leq i\leq n\}.
	\end{equation*}
	Obviously, using the DQM property of $\overline{\cp}^{(n)}$ (treat $\ell  h_i$ as a new parameter $\tilde{\gamma}_i$), we know $\widetilde{\cp}^{(n)}$ is also DQM at $\ell h=0$ with score $ s_{\gamma,\overline{\cp}^{(n)}}(Z)$.
	We now show $\cp^{(1)}$ is regular, i.e., $\cp^{(1)}$ is DQM at $\ell=0$ with score $s_{\ell,\cp^{(1)}}(Z)=h's_{\gamma,\overline{\cp}^{(n)}}(Z)$:
	\begin{equation*}
		\begin{split}
			&||\sqrt{p_{\ell,\cp^{(1)}}}-\sqrt{p_0}-\frac{1}{2}\sqrt{p_0}\ell h's_{\gamma,\overline{\cp}^{(n)}}(Z)||_{L^2_0(\nu)}\\
			&\qquad = ||\sqrt{p_{\ell h,\widetilde{\cp}^{(n)}}}-\sqrt{p_0}-\frac{1}{2}\sqrt{p_0}\ell h's_{\gamma,\overline{\cp}^{(n)}}(Z)||_{L^2_0(\nu)}=o_{L^2(\nu)}(||\ell h||).
		\end{split}
	\end{equation*}
	Since $h$ is a constant, $o_{L^2(\nu)}(||\ell h||)=o_{L^2(\nu)}(|\ell|)$. Therefore, $\widetilde{h}'s_{\gamma,\overline{\cp}^{(n)}}(Z)=Mh's_{\gamma,\overline{\cp}^{(n)}}(Z)=Ms_{\ell,\cp^{(1)}}(Z)$, which implies $\cp^{(n)}\subset\cp^{(1)}$.
\end{proof}


\begin{question}[Any IF is score-orthogonal to $\mathcal T_\eta$ at $P_0$]
  Fix the true value $\theta_0=\theta(P_0)$. Let $\mathcal P_\eta$ be the collection of paths that perturb only the nuisance while keeping $\theta$ fixed at $\theta_0$, i.e., $\mathcal P_\eta:=\{P_\varepsilon\in\mathcal P:\ \theta(P_\varepsilon)\equiv \theta_0\}$. 
The \emph{nuisance tangent space} is
\begin{equation*}
\mathcal T_\eta:=\overline{\{s:\ s \text{ is the score of some }P_\varepsilon\in\mathcal P_\eta\}}\ \subset\ \mathcal T.
\end{equation*}
   If $\phi$ is an influence function at $P_0$, then for every $s\in\mathcal T_\eta$,
\begin{equation*}
\mathbb E_{P_0}\big[\phi(W)\,s(W)\big]=\dot\theta_{P_0}(s)=0,
\end{equation*}
\end{question}
\begin{proof}
    This is simply because nuisance-only paths keep $\theta$ fixed. 
\end{proof}

\begin{question}[When both orthogonalities hold]
    Assume the model is DQM at $P_0$.

\noindent \emph{(A) Global EIF selection.}
Let $\psi(W;\theta,\eta)=\phi^\star_{P_{\theta,\eta}}(W)$, i.e.\ the moment function equals the EIF at the law $P_{\theta,\eta}$. Then both
\begin{equation*}
\mathbb E_{P_0}\!\big[\psi(W;\theta_0,\eta_0)\,s(W)\big]=0\quad(\forall s\in\mathcal T_\eta),
\qquad
\partial_\eta\,\mathbb E_{P_0}\psi(W;\theta_0,\eta_0)[h]=0\quad(\forall h)
\end{equation*}
hold simultaneously.

\noindent \emph{(B) Full-rank linear rescaling of EIF.}
More generally, if $\psi(W;\theta,\eta)=M\,\phi^\star_{P_{\theta,\eta}}(W)$ with a fixed full-rank $d\times d$ matrix $M$, then the same two orthogonality properties hold.
\end{question}

\begin{proof}
\emph{(A) Score orthogonality.}
Fix a nuisance-only regular path $\varepsilon\mapsto P_\varepsilon=P_{\theta_0,\eta_\varepsilon}$ with score $s\in\mathcal T_\eta$.
At $P_0$, the EIF identity gives
\begin{equation*}
\mathbb E_{P_0}\!\big[\psi(W;\theta_0,\eta_0)\,s(W)\big]
=\mathbb E_{P_0}\!\big[\phi^\star_{P_0}(W)\,s(W)\big]
=\dot\theta_{P_0}(s)=0.
\end{equation*}
\emph{Functional orthogonality.}
Define
\begin{equation*}
H(\varepsilon):=\mathbb E_{P_\varepsilon}\!\big[\psi(W;\theta_0,\eta_\varepsilon)\big]
=\mathbb E_{P_{\theta_0,\eta_\varepsilon}}\!\big[\phi^\star_{P_{\theta_0,\eta_\varepsilon}}(W)\big]\equiv 0,
\end{equation*}
since for every $P$ the EIF has zero mean under $P$. By DQM, $H$ is differentiable at $0$ and
\begin{equation*}
0=H'(0)=\underbrace{\mathbb E_{P_0}\!\big[\psi(W;\theta_0,\eta_0)\,s(W)\big]}_{=\,0\ \text{by the previous part}}
+\ \partial_\eta\,\mathbb E_{P_0}\psi(W;\theta_0,\eta_0)[\dot\eta].
\end{equation*}
Hence $\partial_\eta\,\mathbb E_{P_0}\psi(W;\theta_0,\eta_0)[\dot\eta]=0$ for every admissible nuisance direction $\dot\eta$, proving functional orthogonality.

\emph{(B)} Replace $\psi$ by $M\phi^\star_{P_{\theta,\eta}}$. The two conditions are preserved because multiplication by a fixed full-rank matrix does not affect zero inner products nor zero derivatives.
\end{proof}

\begin{question}[Calibration matrix $A$ when using (rescaled) EIF]
Let
\begin{equation*}
A:=\partial_\theta\,\mathbb E_{P_0}\psi(W;\theta,\eta_0)\big|_{\theta=\theta_0}.
\end{equation*}
If $\psi(W;\theta,\eta)=\phi^\star_{P_{\theta,\eta}}(W)$, then $A=-I_d$.
If $\psi(W;\theta,\eta)=M\,\phi^\star_{P_{\theta,\eta}}(W)$ with full-rank $M$, then $A=-M$.

Moreover, suppose $(\hat\theta,\hat\eta)$ solves $\mathbb P_n\psi(W;\theta,\eta)=0$ with $\mathbb E_{P_0}\psi(W;\theta_0,\eta_0)=0$, and regularity (e.g.\ cross-fitting, stochastic equicontinuity) ensures the linearization
\begin{equation*}
\sqrt n(\hat\theta-\theta_0)
=-A^{-1}\,\frac{1}{\sqrt n}\sum_{i=1}^n \psi(W_i;\theta_0,\eta_0)+o_P(1).
\end{equation*}
If $\psi=\phi^\star_{P_{\theta,\eta}}$, the estimator $\hat{\theta}$ reach the semiparametric efficiency bound. 
\end{question}

\begin{proof}
\begin{itemize}
    \item Consider a $\theta$-only regular path $\varepsilon\mapsto P_\varepsilon=P_{\theta_\varepsilon,\eta_0}$ with score $s_\theta$ and $\dot\theta:=\left.\tfrac{d}{d\varepsilon}\theta_\varepsilon\right|_{0}\in\mathbb R^d$ arbitrary (by local identification). Define
\begin{equation*}
G(\varepsilon):=\mathbb E_{P_\varepsilon}\big[\phi^\star_{P_{\theta_\varepsilon,\eta_0}}(W)\big]\equiv 0.
\end{equation*}
DQM implies
\begin{equation*}
0=G'(0)=\underbrace{\mathbb E_{P_0}\big[\phi^\star_{P_0}(W)\,s_\theta(W)\big]}_{=\ \dot\theta\ \text{by EIF identity}}
+\ \partial_\theta\,\mathbb E_{P_0}\phi^\star_{P_{\theta,\eta_0}}(W)\big|_{\theta_0}[\dot\theta].
\end{equation*}
The local identification guarantees $\dot\theta\in \r^d$. Therefore $A=\partial_\theta \mathbb E_{P_0}\phi^\star_{P_{\theta,\eta_0}}|_{\theta_0}=-I_d$.
For $\psi=M\phi^\star$, linearity gives $A=M\,(-I_d)=-M$.
    \item By what we just proved, $A=-I_d$ and
\begin{equation*}
\mathrm{IF}_{\hat\theta}(W)=\phi^\star_{P_0}(W),\qquad
\mathrm{Var}\big(\sqrt n(\hat\theta-\theta_0)\big)=\mathrm{Var}_{P_0}\big(\phi^\star_{P_0}\big),
\end{equation*}
which is the semiparametric efficiency bound.
If $\psi=M\phi^\star$, then $A=-M$ and the IF is still $\phi^\star_{P_0}$; efficiency is unchanged.
\end{itemize}
\end{proof}


\begin{remark}[What each orthogonality controls]
\
\begin{itemize}
\item \emph{Score/geometric orthogonality} $\mathbb E_{P_0}[\psi_0 s]=0$ for $s\in\mathcal T_\eta$ kills the ``distribution channel'' part of the pathwise derivative (variations of $P$ along nuisance directions).
\item \emph{Functional (DML) orthogonality} $\partial_\eta \mathbb E_{P_0}\psi=0$ kills the ``function channel'' that arises when plugging $\hat\eta$ for $\eta_0$ under a fixed law $P_0$.
\item In general they are independent; when $\psi$ is the (rescaled) \emph{EIF map}, both hold simultaneously, and moreover $A=-I$ (or $-M$), yielding efficiency.
\item Any IF at $P_0$ (not necessarily efficient) is automatically score-orthogonal to $\mathcal T_\eta$, but may fail functional (DML) orthogonality under a fixed $P_0$ unless it is chosen as a global selection and the above argument applies.
\end{itemize}
\end{remark}

\begin{question}[Uniqueness in $\overline{\mathcal T}$ and efficiency]
    If $\psi\in\overline{\mathcal T}$ also satisfies
\begin{equation*}
\dot\theta_{P_0}(s)=\mathbb E_{P_0}[\psi s]\qquad\forall\, s\in\mathcal T,
\end{equation*}
then $\psi=\phi^\star$ in $L^2(P_0)$, and for any influence function $\tilde\phi$ one has
\begin{equation*}
\mathrm{Var}_{P_0}(\tilde\phi)\ \ge\ \mathrm{Var}_{P_0}(\phi^\star),
\end{equation*}
with equality iff $\tilde\phi=\phi^\star$ a.s.
\end{question}


\begin{proof}
From the two representations, $\mathbb E[(\psi-\phi^\star)s]=0$ for all $s\in\mathcal T$, i.e.\ $\psi-\phi^\star\perp\mathcal T$. Since both lie in $\overline{\mathcal T}$, orthogonality forces $\psi=\phi^\star$. For efficiency, write any IF as $\tilde\phi=\phi^\star+g$ with $g\perp\mathcal T$ (this is the orthogonal decomposition relative to $\overline{\mathcal T}$). Then
\begin{equation*}
\mathrm{Var}(\tilde\phi)=\mathrm{Var}(\phi^\star)+\mathrm{Var}(g)\ \ge\ \mathrm{Var}(\phi^\star),
\end{equation*}
with equality iff $g=0$.
\end{proof}

\begin{proposition}[Near-orthogonal scores yield root-$n$ linearization; EIF gives efficiency]
Suppose $(\hat\theta,\hat\eta)$ solves the empirical moment equation $\mathbb P_n\psi(W;\theta,\eta)=0$ with $\mathbb E\psi(W;\theta_0,\eta_0)=0$. Let
\begin{equation*}
A:=\partial_\theta \mathbb E_{P_0}\psi(W;\theta_0,\eta_0),\qquad
B:=\partial_\eta \mathbb E_{P_0}\psi(W;\theta_0,\eta_0).
\end{equation*}
Assume $A$ is nonsingular, $\|B\|\le \delta_N\to 0$ (near orthogonality), $\|\hat\eta-\eta_0\|=a_N=o(1)$, and $\delta_N a_N=o(n^{-1/2})$. Then
\begin{equation*}
\sqrt n(\hat\theta-\theta_0)
=-A^{-1}\frac{1}{\sqrt n}\sum_{i=1}^n \psi(W_i;\theta_0,\eta_0)+o_P(1).
\end{equation*}
If, furthermore, $\psi(W;\theta_0,\eta_0)=\phi^\star(W)$, then $A=-I_d$ and
\begin{equation*}
\sqrt n(\hat\theta-\theta_0)\ \rightsquigarrow\ N\!\Big(0,\ A^{-1}\,\mathbb E[\phi^\star\phi^{\star\top}]\,(A^{-1})^\top\Big),
\end{equation*}
achieving the semiparametric efficiency bound.
\end{proposition}

\begin{proof}
A first-order Taylor expansion of the sample moment at $(\theta_0,\eta_0)$ yields
\begin{equation*}
0=\mathbb P_n\psi(\hat\theta,\hat\eta)
= \mathbb P_n\psi(\theta_0,\eta_0)
+ A(\hat\theta-\theta_0) + B(\hat\eta-\eta_0) + r_n,
\end{equation*}
with remainder $r_n=o_P(n^{-1/2})$ under standard regularity (e.g., cross-fitting and stochastic equicontinuity). Multiply by $\sqrt n$ and rearrange:
\begin{equation*}
\sqrt n(\hat\theta-\theta_0)
= -A^{-1}\Big\{\sqrt n(\mathbb P_n-\mathbb P)\psi(\theta_0,\eta_0)
+ \sqrt n\,B(\hat\eta-\eta_0)\Big\}+o_P(1).
\end{equation*}
By the near-orthogonality and rate assumptions, $\|\sqrt n\,B(\hat\eta-\eta_0)\|\le \sqrt n\,\delta_N a_N=o(1)$, which gives the asserted linearization. If $\psi=\phi^\star$, the central limit theorem yields the stated Gaussian limit with covariance $A^{-1}\mathbb E(\phi^\star\phi^{\star\top})(A^{-1})^\top$.
\end{proof}

\begin{corollary}[Relationship between IF/EIF and (near) Neyman orthogonality]

\
\begin{itemize}
\item The EIF is orthogonal to the nuisance tangent space: $\phi^\star\perp\mathcal T_\eta$. Thus, using the EIF as a score produces exact Neyman orthogonality in the sense $\partial_\eta \mathbb E\psi=0$ at the truth.
\item Any IF in $\overline{\mathcal T}$ that represents the pathwise derivative coincides with the EIF and attains the minimal asymptotic variance.
\item Near-orthogonal scores behave like the EIF up to $o(1)$ in $L^2(P_0)$; if $\delta_N a_N=o(n^{-1/2})$, they deliver root-$n$ asymptotics, and if the score equals the EIF at the truth, they achieve the semiparametric efficiency bound.
\end{itemize}
\end{corollary}

\begin{question}[Nuisance tangent space, propensity-score projection, and the efficient influence function]
Let $Z=(X,Y)$ denote the full data and suppose the parameter $\gamma_0$ is defined by
\begin{equation*}
\mathbb{E}[\psi(Z,\gamma_0)]=0.
\end{equation*}
The observed data are
\begin{equation*}
W=(X,D,DY),
\end{equation*}
where $D\in\{0,1\}$ indicates whether $Y$ is observed. Assume missing at random:
\begin{equation*}
\mathbb{P}(D=1\mid X,Y)=p_0(X).
\end{equation*}
Define
\begin{equation*}
q(X)=\mathbb{E}[\psi(Z,\gamma_0)\mid X],
\qquad
\Gamma_0=\mathbb{E}\!\left[\partial_\gamma \psi(Z,\gamma_0)\right].
\end{equation*}

Show that:

\begin{enumerate}
\item any full-data score admits the decomposition
\begin{equation*}
s_Z(Z)=u(X)+b(Y,X),
\end{equation*}
with
\begin{equation*}
\mathbb{E}[u(X)]=0,
\qquad
\mathbb{E}[b(Y,X)\mid X]=0;
\end{equation*}

\item under the missing-at-random assumption, any observed-data score admits the decomposition
\begin{equation*}
s_W(W)=u(X)+D\,b(Y,X)+\left(\frac{D}{p_0(X)}-1\right)a(X),
\end{equation*}
for suitable square-integrable functions $u,b,a$ satisfying
\begin{equation*}
\mathbb{E}[u(X)]=0,
\qquad
\mathbb{E}[b(Y,X)\mid X]=0;
\end{equation*}

\item the nuisance tangent space is
\begin{equation*}
\mathcal{T}_{\eta}
=
\overline{
\left\{
u(X)+D\,b(Y,X)+\left(\frac{D}{p_0(X)}-1\right)a(X):
\mathbb{E}\!\left[\psi(Z,\gamma_0)\{u(X)+b(Y,X)\}\right]=0
\right\}
}^{\,L_2(P_0)};
\end{equation*}

\item if 
\begin{equation*}
m_0(W)=\frac{D}{p_0(X)}\psi(Z,\gamma_0),
\end{equation*}
then its projection onto the propensity-score nuisance subspace
\begin{equation*}
\mathcal{T}_{p}
=
\overline{
\left\{
\left(\frac{D}{p_0(X)}-1\right)a(X): a \in L^2(P_X)
\right\}
}^{\,L_2(P_0)}
\end{equation*}
is
\begin{equation*}
\Pi_{\mathcal{T}_{p}} m_0(W)
=
\left(\frac{D}{p_0(X)}-1\right)q(X);
\end{equation*}

\item the efficient influence function is
\begin{equation*}
\phi^*(W)
=
-\Gamma_0^{-1}
\left[
\frac{D}{p_0(X)}\psi(Z,\gamma_0)
-
\left(\frac{D}{p_0(X)}-1\right)q(X)
\right].
\end{equation*}
\item show that the variance of the efficient influence function is
\begin{equation*}
\mathrm{Var}\bigl(\phi^*(W)\bigr)
=
\Gamma_0^{-1}
\mathbb{E}\!\left[
q(X)q(X)'
+
\frac{\mathrm{Var}\bigl(\psi(Z,\gamma_0)\mid X\bigr)}{p_0(X)}
\right]
\Gamma_0^{-1'} .
\end{equation*}
\end{enumerate}
\end{question}

\begin{proof}

\textbf{Part 1: Decomposition of any full-data score.}

Write the full-data density as
\begin{equation*}
f_Z(x,y)=f_X(x)f_{Y\mid X}(y\mid x).
\end{equation*}
We first show, by explicit disturbances, what the full-data score directions look like.

Let $u=u(X)$ satisfy
\begin{equation*}
\mathbb{E}[u(X)]=0,
\qquad
\mathbb{E}[u(X)^2]<\infty.
\end{equation*}
Define the perturbed marginal density
\begin{equation*}
f_{X,t}(x)=f_X(x)\{1+t\,u(x)\},
\end{equation*}
for $|t|$ sufficiently small. Since $\mathbb{E}[u(X)]=0$,
\begin{equation*}
\int f_{X,t}(x)\,dx
=
1+t\int u(x)f_X(x)\,dx
=
1,
\end{equation*}
so $f_{X,t}$ is again a density. Its score contribution is
\begin{equation*}
\left.\frac{\partial}{\partial t}\log f_{X,t}(X)\right|_{t=0}=u(X).
\end{equation*}

Next let $b=b(Y,X)$ satisfy
\begin{equation*}
\mathbb{E}[b(Y,X)\mid X]=0,
\qquad
\mathbb{E}[b(Y,X)^2]<\infty.
\end{equation*}
Define the perturbed conditional density
\begin{equation*}
f_{Y\mid X,t}(y\mid x)=f_{Y\mid X}(y\mid x)\{1+t\,b(y,x)\},
\end{equation*}
for $|t|$ sufficiently small. Since
\begin{equation*}
\int f_{Y\mid X,t}(y\mid x)\,dy
=
1+t\int b(y,x)f_{Y\mid X}(y\mid x)\,dy
=
1+t\,\mathbb{E}[b(Y,X)\mid X=x]
=
1,
\end{equation*}
this is again a conditional density, and its score contribution is
\begin{equation*}
\left.\frac{\partial}{\partial t}\log f_{Y\mid X,t}(Y\mid X)\right|_{t=0}=b(Y,X).
\end{equation*}

Hence the score generated by the disturbed full-data model
\begin{equation*}
f_{Z,t}(x,y)=f_{X,t}(x)f_{Y\mid X,t}(y\mid x)
\end{equation*}
is
\begin{equation*}
s_Z(Z)=u(X)+b(Y,X).
\end{equation*}

Conversely, if $s_Z\in L_2(P_Z)$ is any full-data score, then the conditional-mean decomposition gives
\begin{equation*}
u(X)=\mathbb{E}[s_Z(Z)\mid X],
\qquad
b(Y,X)=s_Z(Z)-\mathbb{E}[s_Z(Z)\mid X].
\end{equation*}
Therefore
\begin{equation*}
s_Z(Z)=u(X)+b(Y,X),
\end{equation*}
with
\begin{equation*}
\mathbb{E}[u(X)]=\mathbb{E}[s_Z(Z)]=0
\end{equation*}
and
\begin{equation*}
\mathbb{E}[b(Y,X)\mid X]
=
\mathbb{E}[s_Z(Z)\mid X]-\mathbb{E}[s_Z(Z)\mid X]
=
0.
\end{equation*}
This proves the stated decomposition.

\medskip
\noindent
\textbf{Part 2: Decomposition of any observed-data score under MAR.}

Under MAR, the observed-data density is
\begin{equation*}
f_W(x,d,dy)
=
f_X(x)\,
p_0(x)^d\{1-p_0(x)\}^{1-d}\,
f_{Y\mid X}(y\mid x)^d .
\end{equation*}
We again identify the score directions by explicit disturbances of each factor.

The disturbances of $f_X$ and $f_{Y\mid X}$ from Part 1 induce the observed-data density
\begin{equation*}
f_{W,t}(x,d,dy)
=
f_{X,t}(x)\,
p_0(x)^d\{1-p_0(x)\}^{1-d}\,
f_{Y\mid X,t}(y\mid x)^d .
\end{equation*}
Taking logarithms and differentiating at $t=0$ yields
\begin{equation*}
\left.\frac{\partial}{\partial t}\log f_{W,t}(W)\right|_{t=0}
=
u(X)+D\,b(Y,X).
\end{equation*}
The factor $D$ appears because the conditional density of $Y$ contributes to the observed-data likelihood only when $D=1$.

Now perturb the propensity score by a testing function $a(X)$ with
\begin{equation*}
\mathbb{E}[a(X)^2]<\infty.
\end{equation*}
Define the submodel
\begin{equation*}
\operatorname{logit} p_t(X)=\operatorname{logit} p_0(X)+t\,a(X),
\end{equation*}
or equivalently
\begin{equation*}
p_t(X)=\Lambda\bigl(\operatorname{logit} p_0(X)+t\,a(X)\bigr),
\qquad
\Lambda(v)=\frac{e^v}{1+e^v}.
\end{equation*}
Then
\begin{equation*}
\dot p(X):=\left.\frac{\partial}{\partial t}p_t(X)\right|_{t=0}
=
p_0(X)\{1-p_0(X)\}a(X).
\end{equation*}
Hence
\begin{equation*}
\left.\frac{\partial}{\partial t}
\Bigl[
D\log p_t(X)+(1-D)\log\{1-p_t(X)\}
\Bigr]\right|_{t=0}
=
D\frac{\dot p(X)}{p_0(X)}-(1-D)\frac{\dot p(X)}{1-p_0(X)}
\end{equation*}
\begin{equation*}
=
D\{1-p_0(X)\}a(X)-(1-D)p_0(X)a(X)
=
\bigl(D-p_0(X)\bigr)a(X).
\end{equation*}
Since
\begin{equation*}
D-p_0(X)=p_0(X)\left(\frac{D}{p_0(X)}-1\right),
\end{equation*}
this score direction can equivalently be written as
\begin{equation*}
\left(\frac{D}{p_0(X)}-1\right)\tilde a(X),
\qquad
\tilde a(X)=p_0(X)a(X).
\end{equation*}
Because $\varepsilon< p_0(X)<1-\varepsilon$ for some $\varepsilon>0$, the arbitrariness of $a$ is equivalent to the arbitrariness of $\tilde a$. We can still rename $\tilde a$ to $a$.

Combining the three disturbances, the observed-data score generated by
\begin{equation*}
f_{W,t}(x,d,dy)
=
f_{X,t}(x)\,
p_t(x)^d\{1-p_t(x)\}^{1-d}\,
f_{Y\mid X,t}(y\mid x)^d
\end{equation*}
is
\begin{equation*}
s_W(W)=u(X)+D\,b(Y,X)+\left(\frac{D}{p_0(X)}-1\right)a(X).
\end{equation*}

Conversely, since the observed likelihood factorizes into an $X$-part, a $D\mid X$-part, and a $Y\mid X,D=1$-part, any regular observed-data score must lie in the $L_2(P_0)$-closure of the span of these three types of score directions. Hence every observed-data tangent direction admits the displayed representation, with
\begin{equation*}
\mathbb{E}[u(X)]=0,
\qquad
\mathbb{E}[b(Y,X)\mid X]=0,
\end{equation*}
and $u,b,a$ square-integrable.

\medskip
\noindent
\textbf{Part 3: Identification of the nuisance tangent space.}

Let $\gamma_t$ denote the parameter value along a regular submodel. Since $\gamma_t$ is defined by
\begin{equation*}
\mathbb{E}_t[\psi(Z,\gamma_t)]=0,
\end{equation*}
differentiate at $t=0$:
\begin{equation*}
0
=
\left.\frac{d}{dt}\mathbb{E}_t[\psi(Z,\gamma_t)]\right|_{t=0}
=
\mathbb{E}\!\left[\psi(Z,\gamma_0)s_Z(Z)\right]
+
\Gamma_0\,\dot\gamma.
\end{equation*}
By Part 1,
\begin{equation*}
s_Z(Z)=u(X)+b(Y,X),
\end{equation*}
hence
\begin{equation*}
0
=
\mathbb{E}\!\left[\psi(Z,\gamma_0)\{u(X)+b(Y,X)\}\right]
+
\Gamma_0\,\dot\gamma.
\end{equation*}
Therefore the parameter does not move along the path, i.e.
\begin{equation*}
\dot\gamma=0,
\end{equation*}
if and only if
\begin{equation*}
\mathbb{E}\!\left[\psi(Z,\gamma_0)\{u(X)+b(Y,X)\}\right]=0.
\end{equation*}
Thus the nuisance tangent space is exactly
\begin{equation*}
\mathcal{T}_{\eta}
=
\overline{
\left\{
u(X)+D\,b(Y,X)+\left(\frac{D}{p_0(X)}-1\right)a(X):
\mathbb{E}\!\left[\psi(Z,\gamma_0)\{u(X)+b(Y,X)\}\right]=0
\right\}
}^{\,L_2(P_0)}.
\end{equation*}

\medskip
\noindent
\textbf{Part 4: Projection onto the propensity-score nuisance subspace.}

Consider
\begin{equation*}
m_0(W)=\frac{D}{p_0(X)}\psi(Z,\gamma_0).
\end{equation*}
We compute its $L^2$ projection onto
\begin{equation*}
\mathcal{T}_{p}
=
\overline{
\left\{
\left(\frac{D}{p_0(X)}-1\right)a(X): a\in L^2(P_X)
\right\}
}^{\,L_2(P_0)}.
\end{equation*}
Let
\begin{equation*}
h_a(W)=\left(\frac{D}{p_0(X)}-1\right)a(X).
\end{equation*}
The projection $h_{a^*}$ is characterized by
\begin{equation*}
\mathbb{E}\!\left[\{m_0(W)-h_{a^*}(W)\}h_c(W)\right]=0
\qquad
\text{for all } c\in L^2(P_X),
\end{equation*}
where
\begin{equation*}
h_c(W)=\left(\frac{D}{p_0(X)}-1\right)c(X).
\end{equation*}
Thus
\begin{equation*}
\mathbb{E}\!\left[
\left\{
\frac{D}{p_0(X)}\psi(Z,\gamma_0)
-
\left(\frac{D}{p_0(X)}-1\right)a^*(X)
\right\}
\left(\frac{D}{p_0(X)}-1\right)c(X)
\right]=0
\end{equation*}
for all $c(X)$.

We now evaluate the two terms. First,
\begin{equation*}
\mathbb{E}\!\left[
\frac{D}{p_0(X)}\psi(Z,\gamma_0)
\left(\frac{D}{p_0(X)}-1\right)
\Bigm|X
\right]
=
\frac{1-p_0(X)}{p_0(X)}\,q(X).
\end{equation*}
Indeed, conditioning on $(X,Y)$,
\begin{equation*}
\mathbb{E}\!\left[
\frac{D}{p_0(X)}\left(\frac{D}{p_0(X)}-1\right)
\Bigm|X,Y
\right]
=\mathbb{E}\!\left[
\frac{D}{p_0(X)}\left(\frac{1}{p_0(X)}-1\right)
\Bigm|X,Y
\right]=
\frac{1-p_0(X)}{p_0(X)},
\end{equation*}
and then taking conditional expectation with respect to $Y$ given $X$, and noting CAR give the claim.

Second,
\begin{equation*}
\mathbb{E}\!\left[
\left(\frac{D}{p_0(X)}-1\right)^2
\Bigm|X
\right]
=
\frac{1-p_0(X)}{p_0(X)}.
\end{equation*}
Therefore the orthogonality condition becomes
\begin{equation*}
\mathbb{E}\!\left[
\frac{1-p_0(X)}{p_0(X)}
\{q(X)-a^*(X)\}
c(X)
\right]=0
\qquad
\text{for all } c(X).
\end{equation*}
Hence
\begin{equation*}
a^*(X)=q(X)
\qquad \text{a.s.}
\end{equation*}
and so
\begin{equation*}
\Pi_{\mathcal{T}_{p}} m_0(W)
=
\left(\frac{D}{p_0(X)}-1\right)q(X).
\end{equation*}

Thus the natural candidate obtained by subtracting the propensity-score nuisance projection from $m_0(W)$ is
\begin{equation*}
m_0(W)-\Pi_{\mathcal{T}_{p}}m_0(W)
=
\frac{D}{p_0(X)}\psi(Z,\gamma_0)
-
\left(\frac{D}{p_0(X)}-1\right)q(X).
\end{equation*}

\medskip
\noindent
\textbf{Part 5: Verification of the efficient influence function.}

Define
\begin{equation*}
\xi(W)
=
\frac{D}{p_0(X)}\psi(Z,\gamma_0)
-
\left(\frac{D}{p_0(X)}-1\right)q(X).
\end{equation*}
We show that $\xi\in \mathcal{T}_{\eta}^{\perp}$.

Take any
\begin{equation*}
s_W(W)=u(X)+D\,b(Y,X)+\left(\frac{D}{p_0(X)}-1\right)a(X)\in\mathcal{T}_{\eta}.
\end{equation*}
Then
\begin{equation*}
\mathbb{E}[\xi(W)s_W(W)]
=
I_1+I_2+I_3,
\end{equation*}
where
\begin{equation*}
I_1
=
\mathbb{E}\!\left[
\left\{
\frac{D}{p_0(X)}\psi(Z,\gamma_0)
-
\left(\frac{D}{p_0(X)}-1\right)q(X)
\right\}
u(X)
\right],
\end{equation*}
\begin{equation*}
I_2
=
\mathbb{E}\!\left[
\left\{
\frac{D}{p_0(X)}\psi(Z,\gamma_0)
-
\left(\frac{D}{p_0(X)}-1\right)q(X)
\right\}
D\,b(Y,X)
\right],
\end{equation*}
and
\begin{equation*}
I_3
=
\mathbb{E}\!\left[
\left\{
\frac{D}{p_0(X)}\psi(Z,\gamma_0)
-
\left(\frac{D}{p_0(X)}-1\right)q(X)
\right\}
\left(\frac{D}{p_0(X)}-1\right)a(X)
\right].
\end{equation*}

For $I_1$, because
\begin{equation*}
\mathbb{E}\!\left[\frac{D}{p_0(X)}\psi(Z,\gamma_0)\mid X\right]=q(X)
\end{equation*}
and
\begin{equation*}
\mathbb{E}\!\left[\frac{D}{p_0(X)}-1\mid X\right]=0,
\end{equation*}
we obtain
\begin{equation*}
I_1
=
\mathbb{E}[q(X)u(X)]
=
\mathbb{E}[\psi(Z,\gamma_0)u(X)].
\end{equation*}

For $I_2$, first note that
\begin{equation*}
\begin{split}
	&\mathbb{E}\!\left[
\left(\frac{D}{p_0(X)}-1\right)q(X)\,D\,b(Y,X)
\right]
	=\mathbb{E}\!\left[
\left(\frac{1}{p_0(X)}-1\right)q(X)\,D\,b(Y,X)
\right]\\
&\qquad=\mathbb{E}\!\left[
\left(\frac{1}{p_0(X)}-1\right)q(X)\,\e(D|X)\e[b(Y,X)|X]
\right] 
=
0.
\end{split}
\end{equation*}
Also,
\begin{equation*}
\mathbb{E}\!\left[\frac{D}{p_0(X)}\psi(Z,\gamma_0)\,D\,b(Y,X)\right]
=
\mathbb{E}[\psi(Z,\gamma_0)b(Y,X)].
\end{equation*}
Hence
\begin{equation*}
I_2=\mathbb{E}[\psi(Z,\gamma_0)b(Y,X)].
\end{equation*}

For $I_3$, based on Part 4,
\begin{equation*}
I_3=0.
\end{equation*}

Combining the three pieces, 
\begin{equation*}
\mathbb{E}[\xi(W)s_W(W)]
=
\mathbb{E}[\psi(Z,\gamma_0)u(X)]
+
\mathbb{E}[\psi(Z,\gamma_0)b(Y,X)]
=
\mathbb{E}\!\left[\psi(Z,\gamma_0)\{u(X)+b(Y,X)\}\right].
\end{equation*}
Because $s_W\in\mathcal{T}_{\eta}$, the defining restriction of $\mathcal{T}_{\eta}$ gives
\begin{equation*}
\mathbb{E}\!\left[\psi(Z,\gamma_0)\{u(X)+b(Y,X)\}\right]=0.
\end{equation*}
Hence
\begin{equation*}
\mathbb{E}[\xi(W)s_W(W)]=0
\qquad
\forall s_W\in\mathcal{T}_{\eta},
\end{equation*}
so
\begin{equation*}
\xi\in \mathcal{T}_{\eta}^{\perp}.
\end{equation*}

Finally, the efficient influence function is obtained by premultiplying by $-\Gamma_0^{-1}$:
\begin{equation*}
\phi^*(W)
=
-\Gamma_0^{-1}\xi(W)
=
-\Gamma_0^{-1}
\left[
\frac{D}{p_0(X)}\psi(Z,\gamma_0)
-
\left(\frac{D}{p_0(X)}-1\right)q(X)
\right].
\end{equation*}
This proves the claim.

\medskip
\noindent
\textbf{Part 6: Variance of the efficient influence function.}

Recall that
\begin{equation*}
\phi^*(W)=-\Gamma_0^{-1}\xi(W),
\end{equation*}
where
\begin{equation*}
\xi(W)
=
\frac{D}{p_0(X)}\psi(Z,\gamma_0)
-
\left(\frac{D}{p_0(X)}-1\right)q(X).
\end{equation*}
Since $\Gamma_0$ is nonrandom,
\begin{equation*}
\mathrm{Var}\bigl(\phi^*(W)\bigr)
=
\Gamma_0^{-1}\,
\mathrm{Var}\bigl(\xi(W)\bigr)\,
\Gamma_0^{-1'}.
\end{equation*}
Thus it suffices to compute $\mathrm{Var}(\xi(W))$.

We first rewrite $\xi(W)$ as
\begin{equation*}
\xi(W)
=
q(X)+\frac{D}{p_0(X)}\bigl(\psi(Z,\gamma_0)-q(X)\bigr).
\end{equation*}
Indeed,
\begin{equation*}
\frac{D}{p_0(X)}\psi(Z,\gamma_0)
-
\left(\frac{D}{p_0(X)}-1\right)q(X)
=
q(X)+\frac{D}{p_0(X)}\bigl(\psi(Z,\gamma_0)-q(X)\bigr).
\end{equation*}

Let
\begin{equation*}
u(Z)=\psi(Z,\gamma_0)-q(X).
\end{equation*}
Then
\begin{equation*}
\mathbb{E}[u(Z)\mid X]=0,
\end{equation*}
and hence
\begin{equation*}
\xi(W)=q(X)+\frac{D}{p_0(X)}u(Z).
\end{equation*}

Now
\begin{equation*}
\mathbb{E}[\xi(W)\mid X]
=
q(X)+\mathbb{E}\!\left[\frac{D}{p_0(X)}u(Z)\Bigm| X\right].
\end{equation*}
Using iterated expectations and the missing-at-random assumption,
\begin{equation*}
\mathbb{E}\!\left[\frac{D}{p_0(X)}u(Z)\Bigm| X\right]
=
\frac{1}{p_0(X)}\mathbb{E}\!\left[\mathbb{E}[D\,u(Z)\mid X,Y]\Bigm|X\right]
=
\frac{1}{p_0(X)}\mathbb{E}[p_0(X)u(Z)\mid X]
=
\mathbb{E}[u(Z)\mid X]
=
0.
\end{equation*}
Therefore
\begin{equation*}
\mathbb{E}[\xi(W)\mid X]=q(X).
\end{equation*}

By the law of total variance,
\begin{equation*}
\mathrm{Var}\bigl(\xi(W)\bigr)
=
\mathbb{E}\!\left[\mathrm{Var}\bigl(\xi(W)\mid X\bigr)\right]
+
\mathrm{Var}\left(\mathbb{E}[\xi(W)\mid X]\right).
\end{equation*}
Since $\mathbb{E}[\xi(W)\mid X]=q(X)$ and $\mathbb{E}[q(X)]=\mathbb{E}[\psi(Z,\gamma_0)]=0$, we have
\begin{equation*}
\mathrm{Var}\left(\mathbb{E}[\xi(W)\mid X]\right)
=
\mathrm{Var}\bigl(q(X)\bigr)
=
\mathbb{E}[q(X)q(X)'].
\end{equation*}

It remains to compute the conditional variance term. Because $q(X)$ is $X$-measurable,
\begin{equation*}
\mathrm{Var}\bigl(\xi(W)\mid X\bigr)
=
\mathrm{Var}\left(\frac{D}{p_0(X)}u(Z)\Bigm|X\right).
\end{equation*}
Also,
\begin{equation*}
\mathbb{E}\!\left[\frac{D}{p_0(X)}u(Z)\Bigm|X\right]=0,
\end{equation*}
so
\begin{equation*}
\mathrm{Var}\left(\frac{D}{p_0(X)}u(Z)\Bigm|X\right)
=
\mathbb{E}\!\left[
\frac{D}{p_0(X)^2}u(Z)u(Z)'
\Bigm|X
\right].
\end{equation*}
Conditioning further on $(X,Y)$,
\begin{equation*}
\mathbb{E}\!\left[
\frac{D}{p_0(X)^2}u(Z)u(Z)'
\Bigm|X,Y
\right]
=
\frac{1}{p_0(X)^2}u(Z)u(Z)'\,\mathbb{E}[D\mid X,Y]
=
\frac{1}{p_0(X)}u(Z)u(Z)'.
\end{equation*}
Taking conditional expectation with respect to $Y$ given $X$,
\begin{equation*}
\mathrm{Var}\bigl(\xi(W)\mid X\bigr)
=
\frac{1}{p_0(X)}\mathbb{E}[u(Z)u(Z)'\mid X]
=
\frac{1}{p_0(X)}\mathrm{Var}\bigl(\psi(Z,\gamma_0)\mid X\bigr).
\end{equation*}
Therefore
\begin{equation*}
\mathbb{E}\!\left[\mathrm{Var}\bigl(\xi(W)\mid X\bigr)\right]
=
\mathbb{E}\!\left[
\frac{\mathrm{Var}\bigl(\psi(Z,\gamma_0)\mid X\bigr)}{p_0(X)}
\right].
\end{equation*}

Combining the two pieces yields
\begin{equation*}
\mathrm{Var}\bigl(\xi(W)\bigr)
=
\mathbb{E}\!\left[
q(X)q(X)'
+
\frac{\mathrm{Var}\bigl(\psi(Z,\gamma_0)\mid X\bigr)}{p_0(X)}
\right].
\end{equation*}
Hence
\begin{equation*}
\mathrm{Var}\bigl(\phi^*(W)\bigr)
=
\Gamma_0^{-1}
\mathbb{E}\!\left[
q(X)q(X)'
+
\frac{\mathrm{Var}\bigl(\psi(Z,\gamma_0)\mid X\bigr)}{p_0(X)}
\right]
\Gamma_0^{-1'}.
\end{equation*}
This proves the claim.
\end{proof}

\begin{question}[Efficient influence function for a linear functional of a regression]
Let $O=(Y,X)$ be a random vector with distribution $P=P_0$, and define
\begin{equation*}
\gamma_P(x) := \mathbb{E}_P[Y \mid X=x].
\end{equation*}
Let the target parameter be
\begin{equation*}
\theta(P) := \mathbb{E}_P\big[m(Z,\gamma_P)\big],
\end{equation*}
where $m(Z,\cdot)$ is linear in its second argument.

Assume that for every regular parametric submodel $\{P_t:|t|<\delta\}$ through $P=P_0$:

\begin{enumerate}
\item[(i)] The joint law of $(Y,X)$ under $P_t$ has density $p_t(y,x)$ with respect to a dominating measure $\nu(dy)\mu(dx)$, and
\begin{equation*}
\dot p_0(y,x):=\left.\frac{d}{dt}p_t(y,x)\right|_{t=0}=p_0(y,x)s(y,x)
\end{equation*}
for some score $s(O)\in L^2_0(P)$.

\item[(ii)] Differentiation under the integral sign is valid for
\begin{equation*}
N_t(x):=\int y\,p_t(y,x)\,\nu(dy),
\qquad
D_t(x):=\int p_t(y,x)\,\nu(dy),
\end{equation*}
and $D_0(x)>0$ for $P_X$-almost every $x$.

\item[(iii)] The regression map is differentiable along the path in $L^2(P_X)$:
\begin{equation*}
\gamma_{P_t}=\gamma_P+t\dot\gamma+r_t,
\qquad
\|r_t\|_{L^2(P_X)}=o(t).
\end{equation*}

\item[(iv)] For every $f\in L^2(P)$,
\begin{equation*}
\left.\frac{d}{dt}\mathbb{E}_{P_t}[f(O)]\right|_{t=0}
=
\mathbb{E}_P[f(O)s(O)].
\end{equation*}

\item[(v)] The map $g\mapsto m(Z,g)$ is a bounded linear operator from $L^2(P_X)$ into $L^2(P)$, and the linear functional
\begin{equation*}
\mathcal L(g):=\mathbb{E}_P[m(Z,g)]
\end{equation*}
is continuous on $L^2(P_X)$. Hence there exists $\alpha_P\in L^2(P_X)$ such that
\begin{equation*}
\mathbb{E}_P[m(Z,g)]
=
\mathbb{E}_P[\alpha_P(X)g(X)]
\qquad
\forall g\in L^2(P_X).
\end{equation*}
\end{enumerate}

Show that
\begin{equation*}
\dot\gamma(x)
=
\mathbb{E}_P\big[(Y-\gamma_P(X))s(O)\mid X=x\big]
\qquad
P_X\text{-a.e.},
\end{equation*}
and that $\theta(P)$ is pathwise differentiable with influence function
\begin{equation*}
\phi(O)
=
m(Z,\gamma_P)-\theta(P)+\alpha_P(X)\{Y-\gamma_P(X)\}.
\end{equation*}
\end{question}

\begin{proof}
Fix a regular parametric submodel $\{P_t\}$ with score $s(O)$, and write
\begin{equation*}
\theta_t:=\theta(P_t)=\mathbb{E}_{P_t}[m(Z,\gamma_{P_t})].
\end{equation*}

\medskip
\textbf{Step 1: Identify the pathwise derivative of $\gamma_{P_t}(x)$.}

By definition,
\begin{equation*}
\gamma_{P_t}(x)
=
\frac{N_t(x)}{D_t(x)},
\qquad
N_t(x):=\int y\,p_t(y,x)\,\nu(dy),
\qquad
D_t(x):=\int p_t(y,x)\,\nu(dy).
\end{equation*}
Since $D_0(x)>0$ for $P_X$-almost every $x$, the quotient rule applies at such $x$:
\begin{equation*}
\left.\frac{d}{dt}\gamma_{P_t}(x)\right|_{t=0}
=
\frac{\dot N_0(x)D_0(x)-N_0(x)\dot D_0(x)}{D_0(x)^2}.
\end{equation*}

By differentiation under the integral sign and the identity $\dot p_0=p_0s$,
\begin{equation*}
\dot N_0(x)
=
\int y\,\dot p_0(y,x)\,\nu(dy)
=
\int y\,p_0(y,x)s(y,x)\,\nu(dy),
\end{equation*}
and
\begin{equation*}
\dot D_0(x)
=
\int \dot p_0(y,x)\,\nu(dy)
=
\int p_0(y,x)s(y,x)\,\nu(dy).
\end{equation*}
Also,
\begin{equation*}
N_0(x)=\int y\,p_0(y,x)\,\nu(dy)=\gamma_P(x)D_0(x).
\end{equation*}
Substituting these identities into the quotient rule gives
\begin{equation*}
\left.\frac{d}{dt}\gamma_{P_t}(x)\right|_{t=0}
=
\frac{\int y\,p_0(y,x)s(y,x)\,\nu(dy)-\gamma_P(x)\int p_0(y,x)s(y,x)\,\nu(dy)}{D_0(x)}.
\end{equation*}
Now use
\begin{equation*}
p_0(y,x)=p_0(y\mid x)D_0(x)
\end{equation*}
to rewrite the numerator:
\begin{equation*}
\left.\frac{d}{dt}\gamma_{P_t}(x)\right|_{t=0}
=
\int (y-\gamma_P(x))\,s(y,x)\,p_0(y\mid x)\,\nu(dy).
\end{equation*}
Therefore,
\begin{equation*}
\left.\frac{d}{dt}\gamma_{P_t}(x)\right|_{t=0}
=
\mathbb{E}_P\big[(Y-\gamma_P(X))s(O)\mid X=x\big]
\qquad
P_X\text{-a.e.}
\end{equation*}

On the other hand, assumption (iii) says that
\begin{equation*}
\gamma_{P_t}=\gamma_P+t\dot\gamma+r_t,
\qquad
\|r_t\|_{L^2(P_X)}=o(t).
\end{equation*}
Hence $\dot\gamma$ is the $L^2(P_X)$ pathwise derivative of $\gamma_{P_t}$, and by uniqueness of the derivative in $L^2(P_X)$ we conclude
\begin{equation*}
\boxed{
\dot\gamma(x)
=
\mathbb{E}_P\big[(Y-\gamma_P(X))s(O)\mid X=x\big]
}
\qquad
P_X\text{-a.e.}
\end{equation*}

\medskip
\textbf{Step 2: Expand $m(Z,\gamma_{P_t})$.}

By assumption (iii),
\begin{equation*}
\gamma_{P_t}-\gamma_P=t\dot\gamma+r_t.
\end{equation*}
Using linearity of $m$ in its second argument,
\begin{equation*}
m(Z,\gamma_{P_t})
=
m(Z,\gamma_P)+t\,m(Z,\dot\gamma)+m(Z,r_t).
\end{equation*}
Because $g\mapsto m(Z,g)$ is bounded from $L^2(P_X)$ to $L^2(P)$,
\begin{equation*}
\|m(Z,r_t)\|_{L^2(P)}
\le C\|r_t\|_{L^2(P_X)}
=o(t).
\end{equation*}
Hence
\begin{equation*}
m(Z,\gamma_{P_t})
=
m(Z,\gamma_P)+t\,m(Z,\dot\gamma)+o(t)
\qquad
\text{in }L^2(P).
\end{equation*}

\medskip
\textbf{Step 3: Expand $\theta_t$.}

Taking expectations under $P_t$,
\begin{equation*}
\theta_t
=
\mathbb{E}_{P_t}[m(Z,\gamma_P)]
+
t\,\mathbb{E}_{P_t}[m(Z,\dot\gamma)]
+
o(t).
\end{equation*}
We treat the two leading terms separately.

First, by assumption (iv),
\begin{equation*}
\left.\frac{d}{dt}\mathbb{E}_{P_t}[m(Z,\gamma_P)]\right|_{t=0}
=
\mathbb{E}_P[m(Z,\gamma_P)s(O)].
\end{equation*}
Second, since $m(Z,\dot\gamma)\in L^2(P)$, we have
\begin{equation*}
\mathbb{E}_{P_t}[m(Z,\dot\gamma)]
=
\mathbb{E}_P[m(Z,\dot\gamma)] + o(1),
\end{equation*}
and therefore
\begin{equation*}
t\,\mathbb{E}_{P_t}[m(Z,\dot\gamma)]
=
t\,\mathbb{E}_P[m(Z,\dot\gamma)] + o(t).
\end{equation*}
Combining these displays,
\begin{equation*}
\theta_t
=
\theta(P)
+
t\Big(
\mathbb{E}_P[m(Z,\gamma_P)s(O)]
+
\mathbb{E}_P[m(Z,\dot\gamma)]
\Big)
+
o(t).
\end{equation*}

Since $\mathbb{E}_P[s(O)]=0$,
\begin{equation*}
\mathbb{E}_P[m(Z,\gamma_P)s(O)]
=
\mathbb{E}_P[(m(Z,\gamma_P)-\theta(P))s(O)].
\end{equation*}

\medskip
\textbf{Step 4: Use the Riesz representation for the second term.}

By the definition of $\alpha_P$,
\begin{equation*}
\mathbb{E}_P[m(Z,\dot\gamma)]
=
\mathbb{E}_P[\alpha_P(X)\dot\gamma(X)].
\end{equation*}
Now substitute the formula for $\dot\gamma$ proved in Step 1:
\begin{equation*}
\mathbb{E}_P[\alpha_P(X)\dot\gamma(X)]
=
\mathbb{E}_P\Big[
\alpha_P(X)\mathbb{E}_P[(Y-\gamma_P(X))s(O)\mid X]
\Big].
\end{equation*}
Applying the tower property,
\begin{equation*}
\mathbb{E}_P[\alpha_P(X)\dot\gamma(X)]
=
\mathbb{E}_P[\alpha_P(X)(Y-\gamma_P(X))s(O)].
\end{equation*}

\medskip
\textbf{Step 5: Identify the influence function.}

Substituting the last two displays into the expansion of $\theta_t$ yields
\begin{equation*}
\left.\frac{d}{dt}\theta(P_t)\right|_{t=0}
=
\mathbb{E}_P\Big[
\big(
m(Z,\gamma_P)-\theta(P)+\alpha_P(X)(Y-\gamma_P(X))
\big)s(O)
\Big].
\end{equation*}
Hence
\begin{equation*}
\phi(O)
=
m(Z,\gamma_P)-\theta(P)+\alpha_P(X)\{Y-\gamma_P(X)\}
\end{equation*}
represents the pathwise derivative.

Finally, $\phi$ is mean zero:
\begin{equation*}
\mathbb{E}_P[\phi(O)]
=
\mathbb{E}_P[m(Z,\gamma_P)-\theta(P)]
+
\mathbb{E}_P[\alpha_P(X)\{Y-\gamma_P(X)\}]
=
0,
\end{equation*}
because
\begin{equation*}
\mathbb{E}_P[Y-\gamma_P(X)\mid X]=0.
\end{equation*}
Therefore $\theta(P)$ is pathwise differentiable with influence function
\begin{equation*}
\boxed{
\phi(O)
=
m(Z,\gamma_P)-\theta(P)+\alpha_P(X)\{Y-\gamma_P(X)\}.
}
\end{equation*}
\end{proof}

\section{Empirical Processes}
\begin{question}[Asymptotic tightness in countable product spaces]
Let $(D_i,d_i)$ be metric spaces for $i\in\mathbb N$, and let
\begin{equation*}
D := \prod_{i=1}^\infty D_i
\end{equation*}
be endowed with the product topology, metrized for instance by
\begin{equation*}
d(x,z)
:= \sum_{i=1}^\infty 2^{-i}(d_i(x_i,z_i)\wedge 1).
\end{equation*}
Let $X_\alpha=(X_{\alpha,1},X_{\alpha,2},\dots)$ be a family of (possibly nonmeasurable)
$D$-valued random elements.

Show that $X_\alpha$ is asymptotically tight in $D$ if and only if, for every
$i\in\mathbb N$, the coordinate process $X_{\alpha,i}$ is asymptotically tight in $D_i$,
where asymptotic tightness is defined via $\delta$-enlargements of compact sets.
\end{question}

\begin{proof}
Recall that for a metric space $(S,d)$ and $A\subset S$, the $\delta$-enlargement is
\begin{equation*}
A^\delta := \{x\in S:\ d(x,A)<\delta\}.
\end{equation*}
We want to show the following two facts:

\smallskip
\noindent
(a) Let $\delta_i\in(0,1]$ satisfy $\sum_i 2^{-i}\delta_i\le\delta$. For arbitrary sets
$K_i\subset D_i$,
\begin{equation*}
\prod_{i=1}^\infty K_i^{\delta_i}
\subset
\(\prod_{i=1}^\infty K_i\)^\delta.
\end{equation*}
Indeed, for $x\in\prod_i K_i^{\delta_i}$ choose $y_i\in K_i$ with
$d_i(x_i,y_i)<\delta_i$. Then $y=(y_i)\in\prod_i K_i$ and
\begin{equation*}
d(x,y)
=\sum_i 2^{-i}(d_i(x_i,y_i)\wedge1)
< \sum_i 2^{-i}\delta_i
\le \delta.
\end{equation*}
Note $y\in \prod_{i=1}^\infty K_i$, we have $d(x,\prod_{i=1}^\infty K_i)\leq d(x,y)<\delta$. Namely, $x\in \(\prod_{i=1}^\infty K_i\)^\delta$.

\smallskip
\noindent
(b) Let $\pi_i:D\to D_i$ denote the coordinate projection. For any $K\subset D$ and
$\delta>0$,
\begin{equation*}
K^\delta \subset \prod_{i=1}^\infty \pi_i(K)^{2^i\delta}.
\end{equation*}
Indeed, if $x\in K^\delta$, there exists $y\in K$ with $d(x,y)<\delta$, which implies
\begin{equation*}
2^{-i}(d_i(x_i,y_i)\wedge1)<\sum_i 2^{-i}(d_i(x_i,y_i)\wedge1)=d(x,y)<\delta.
\end{equation*}
Note $y_i\in \pi_i(K)$, hence $d_i(x_i,\pi_i(K))\leq d_i(x_i,y_i)<2^i\delta$. Namely, $x\in \prod_{i=1}^\infty \pi_i(K)^{2^i\delta}$.

\paragraph{Joint asymptotic tightness $\Longrightarrow$ coordinatewise asymptotic tightness.}

Assume $X_\alpha$ is asymptotically tight in $D$. Then for every $\varepsilon>0$ there
exists a compact $K\subset D$ such that for all $\delta>0$,
\begin{equation*}
\varliminf_\alpha \mathbb P^*(X_\alpha\in K^\delta)\ge 1-\varepsilon.
\end{equation*}
Fix $i\in\mathbb N$. Since $\pi_i$ is continuous, $\pi_i(K)$ is compact in $D_i$.
Applying inclusion (b) with $\delta=\eta/2^i$ yields
\begin{equation*}
\{X_\alpha\in K^{\eta/2^i}\}
\subset
\{X_{\alpha,i}\in \pi_i(K)^\eta\}.
\end{equation*}
Hence
\begin{equation*}
\varliminf_\alpha \mathbb P^*(X_{\alpha,i}\in \pi_i(K)^\eta)
\ge 1-\varepsilon,
\end{equation*}
showing that $X_{\alpha,i}$ is asymptotically tight in $D_i$.

\paragraph{Coordinatewise asymptotic tightness $\Longrightarrow$ joint asymptotic tightness.}

Assume that each $X_{\alpha,i}$ is asymptotically tight in $D_i$. Fix $\varepsilon>0$
and $\delta>0$. For every $i$, choose a compact $K_i\subset D_i$ such that
\begin{equation*}
\varliminf_\alpha \mathbb P^*(X_{\alpha,i}\in K_i^\eta)
\ge 1-\varepsilon 2^{-i}
\quad\text{for all }\eta>0.
\end{equation*}
Set
\begin{equation*}
K:=\prod_{i=1}^\infty K_i,
\end{equation*}
which is compact in $D$.
Choose numbers $\delta_i\in(0,1]$ with $\sum_i 2^{-i}\delta_i\le\delta$ (for example, we can choose $\delta_i\equiv \delta$).
By inclusion (a),
\begin{equation*}
\prod_i K_i^{\delta_i} \subset K^\delta.
\end{equation*}
Therefore,
\begin{equation*}
\mathbb P^*(X_\alpha\in K^\delta)
\ge
\mathbb P^*(\forall i,\ X_{\alpha,i}\in K_i^{\delta_i})
\ge
1-\sum_i \mathbb P^*(X_{\alpha,i}\notin K_i^{\delta_i}).
\end{equation*}
Taking $\varliminf_\alpha$ and using the coordinatewise bounds yields
\begin{equation*}
\varliminf_\alpha \mathbb P^*(X_\alpha\in K^\delta)
\ge 1-\sum_i \varepsilon 2^{-i}
=1-\varepsilon.
\end{equation*}
Since $\varepsilon>0$ is arbitrary, $X_\alpha$ is asymptotically tight in $D$.

\medskip
\noindent
The proof is complete.
\end{proof}

\begin{question}
Let $(T,\rho)$ be a semimetric space and let $\ell^\infty(T)$ be equipped with the supremum norm $\|f\|_\infty=\sup_{t\in T}|f(t)|$.
For a set $K\subset\ell^\infty(T)$ and $\eta>0$, define its $\eta$-enlargement by
\begin{equation*}
K^\eta := \Big\{ f\in \ell^\infty(T): \inf_{g\in K}\|f-g\|_\infty < \eta \Big\}.
\end{equation*}
Assume $K$ is uniformly $\rho$-equicontinuous, i.e., for every $\varepsilon>0$ there exists $r>0$ such that
\begin{equation*}
\sup_{g\in K}\ \sup_{\rho(s,t)<r}|g(s)-g(t)| \le \varepsilon.
\end{equation*}
Show that for every $\varepsilon>0$ and $\eta>0$ there exists $r>0$ such that
\begin{equation*}
\sup_{f\in K^\eta}\ \sup_{\rho(s,t)<r}|f(s)-f(t)| \le \varepsilon + 2\eta.
\end{equation*}
Deduce that if $(X_\alpha)_\alpha$ is a net of $\ell^\infty(T)$-valued random elements that is asymptotically tight with compacts $K_m$ in the sense that for each $m\in\mathbb{N}$,
\begin{equation*}
\varliminf_\alpha \mathbb{P}\big(X_\alpha \in K_m^\eta\big) \ge 1-\frac{1}{m}\quad\text{for all }\eta>0,
\end{equation*}
and if each $K_m$ is uniformly $\rho$-equicontinuous, then $(X_\alpha)_\alpha$ is asymptotically uniformly $\rho$-equicontinuous in probability, i.e., for every $\varepsilon>0$,
\begin{equation*}
\lim_{r\downarrow 0}\ \varliminf_\alpha \mathbb{P}\Big(\sup_{\rho(s,t)<r}|X_\alpha(s)-X_\alpha(t)|\le \varepsilon\Big)=1.
\end{equation*}
\end{question}

\begin{proof}
Fix $\varepsilon>0$ and $\eta>0$. Since $K$ is uniformly $\rho$-equicontinuous, there exists $r>0$ such that
\begin{equation*}
\sup_{g\in K}\ \sup_{\rho(s,t)<r}|g(s)-g(t)| \le \varepsilon.
\end{equation*}
Let $f\in K^\eta$. By definition of $K^\eta$, choose $g\in K$ such that $\|f-g\|_\infty<\eta$. Then for any $s,t\in T$,
\begin{equation*}
\begin{split}
	|f(s)-f(t)|
&\le |f(s)-g(s)| + |g(s)-g(t)| + |g(t)-f(t)|\\
&\le 2\|f-g\|_\infty + |g(s)-g(t)|.
\end{split}
\end{equation*}
Taking the supremum over all $(s,t)$ with $\rho(s,t)<r$ yields
\begin{equation*}
\sup_{\rho(s,t)<r}|f(s)-f(t)|
\le 2\|f-g\|_\infty + \sup_{\rho(s,t)<r}|g(s)-g(t)|
\le 2\eta + \varepsilon.
\end{equation*}
Since $f\in K^\eta$ was arbitrary, we obtain
\begin{equation*}
\sup_{f\in K^\eta}\ \sup_{\rho(s,t)<r}|f(s)-f(t)| \le \varepsilon + 2\eta.
\end{equation*}

For the second claim, fix $\varepsilon>0$ and $m\in\mathbb{N}$. Apply the first part with $K=K_m$ and choose $\eta=\varepsilon/2$. Since $K_m$ is uniformly $\rho$-equicontinuous, there exists $r=r(\varepsilon,m)>0$ such that
\begin{equation*}
\sup_{f\in K_m^{\varepsilon/2}}\ \sup_{\rho(s,t)<r}|f(s)-f(t)| \le \varepsilon + 2\cdot\frac{\varepsilon}{2}=2\varepsilon.
\end{equation*}
Consequently,
\begin{equation*}
\mathbb{P}\Big(\sup_{\rho(s,t)<r}|X_\alpha(s)-X_\alpha(t)|\le 2\varepsilon\Big)
\ge \mathbb{P}\big(X_\alpha \in K_m^{\varepsilon/2}\big).
\end{equation*}
Taking $\varliminf_\alpha$ and using asymptotic tightness with the same compacts $K_m$ gives
\begin{equation*}
\varliminf_\alpha \mathbb{P}\Big(\sup_{\rho(s,t)<r}|X_\alpha(s)-X_\alpha(t)|\le 2\varepsilon\Big)
\ge \varliminf_\alpha \mathbb{P}\big(X_\alpha \in K_m^{\varepsilon/2}\big)
\ge 1-\frac{1}{m}.
\end{equation*}
Since $m$ is arbitrary, this implies that for every $\varepsilon>0$ there exists $r>0$ such that the above probability can be made arbitrarily close to $1$. Equivalently,
\begin{equation*}
\lim_{r\downarrow 0}\ \varliminf_\alpha \mathbb{P}\Big(\sup_{\rho(s,t)<r}|X_\alpha(s)-X_\alpha(t)|\le 2\varepsilon\Big)=1.
\end{equation*}
Renaming $2\varepsilon$ as $\varepsilon$ yields the desired asymptotic uniform $\rho$-equicontinuity in probability.
\end{proof}

\begin{question}
	Given an arbitrary mapping $X$ (perhaps not measurable) and a probability space $(\Omega,\f,\p)$. Then for any $t>0$ and any increasing real-valued function $f\in\f$, we have
	\begin{equation*}
		\{f(X)>t\}^*\subset \{f(X^*)>t\},\quad \p\text{-}\mathrm{a.s.}.
	\end{equation*}
\end{question}
\begin{proof}
	We consider $\mathbbm 1_{\{f(X)>t\}^*}$ and $\mathbbm 1_{\{f(X^*)>t\}}$. Since $X\leq X^*,\p$-a.s. and $f$ is increasing, we have 
	\begin{equation*}
		\{f(X)>t\}\subset \{f(X^*)>t\},\quad \p\text{-}\mathrm{a.s.}.
	\end{equation*}
	Thus $\mathbbm 1_{\{f(X)>t\}}\leq \mathbbm 1_{\{f(X^*)>t\}},\p$-a.s.. Using the property of $A^*$ for any set $A\subset \Omega$, i.e., $(\mathbbm 1_A)^*=\mathbbm 1_{A^*}$, and noticing $\mathbbm 1_{\{f(X^*)>t\}}\in\f$, we have
	\begin{equation*}
		\mathbbm 1_{\{f(X)>t\}^*}=(\mathbbm 1_{\{f(X)>t\}})^*\leq \mathbbm 1_{\{f(X^*)>t\}},\quad \p\text{-}\mathrm{a.s.},
	\end{equation*}
	which yields 
	\begin{equation*}
		\{f(X)>t\}^*\subset \{f(X^*)>t\},\quad \p\text{-}\mathrm{a.s.}.
	\end{equation*}
\end{proof}

\begin{question}
Let $(T,\rho)$ be a compact semimetric space and let $C(T)$ denote the space of all real-valued continuous functions on $T$. Equip $C(T)$ with the uniform norm $\|z\|_T := \sup_{t\in T} |z(t)|$.
For each $t\in T$, define the evaluation (projection) map $\pi_t : C(T) \to \mathbb R, \pi_t(z) = z(t),$
and let
\begin{equation*}
\mathcal P := \sigma(\pi_t : t\in T)
\end{equation*}
be the $\sigma$-field generated by all coordinate projections. Show that
\begin{equation*}
\mathcal B(C(T)) = \mathcal P,
\end{equation*}
where $\mathcal B(C(T))$ denotes the Borel $\sigma$-field on $C(T)$ induced by the uniform norm. Conclude that for a mapping $X:\Omega\to C(T)$,
\begin{equation*}
X \text{ is Borel measurable } \iff \forall t\in T,\ X(t):=\pi_t(X) \text{ is measurable}.
\end{equation*}
\end{question}

\begin{proof}
We first show that $\mathcal P \subset \mathcal B(C(T))$. For each $t\in T$ and any $z_1,z_2\in C(T)$,
\begin{equation*}
|\pi_t(z_1)-\pi_t(z_2)| = |z_1(t)-z_2(t)| \le \|z_1-z_2\|_T.
\end{equation*}
Hence $\pi_t$ is Lipschitz continuous and therefore Borel measurable. It follows that the $\sigma$-field generated by all $\pi_t$ is contained in $\mathcal B(C(T))$, so $\mathcal P \subset \mathcal B(C(T))$.

To prove the reverse inclusion, it suffices to show that every closed ball in $C(T)$ belongs to $\mathcal P$, since closed balls generate the Borel $\sigma$-field of a metric space. Fix $z_0\in C(T)$ and $r>0$. The closed ball centered at $z_0$ with radius $r$ is
\begin{equation*}
\overline B(z_0,r)
= \{z\in C(T): \|z-z_0\|_T \le r\}
= \Big\{z\in C(T): \sup_{t\in T}|z(t)-z_0(t)| \le r\Big\}.
\end{equation*}
Since $T$ is compact and $z-z_0$ is continuous, the function $t\mapsto |z(t)-z_0(t)|$ attains its maximum on $T$. Let $S\subset T$ be a countable dense subset. Then
\begin{equation*}
\sup_{t\in T}|z(t)-z_0(t)| = \sup_{s\in S}|z(s)-z_0(s)|.
\end{equation*}
Consequently,
\begin{equation*}
\overline B(z_0,r)
= \bigcap_{s\in S} \{z\in C(T): |z(s)-z_0(s)| \le r\}.
\end{equation*}
Each set on the right-hand side can be written as
\begin{equation*}
\{z\in C(T): |z(s)-z_0(s)| \le r\}
= \pi_s^{-1}\big([z_0(s)-r,\ z_0(s)+r]\big),
\end{equation*}
which belongs to $\mathcal P$. Since $S$ is countable, $\overline B(z_0,r)$ is a countable intersection of sets in $\mathcal P$, hence $\overline B(z_0,r)\in\mathcal P$. This shows that $\mathcal B(C(T))\subset \mathcal P$.

Combining both inclusions yields $\mathcal B(C(T))=\mathcal P$.

Finally, let $X:\Omega\to C(T)$. If $X$ is Borel measurable, then each composition $\pi_t\circ X = X(t)$ is measurable. Conversely, if $X(t)$ is measurable for every $t\in T$, since $X(t)=\pi_t\circ X$ and $\mathcal B(C(T))=\sigma(\pi_t:t\in T)$ is generate by $(\pi_t)_{t\in T}$, then $X$ is measurable with respect to $\sigma(\pi_t:t\in T)=\mathcal B(C(T))$. This completes the proof.
\end{proof}

\begin{question}
Let $D:=D[a,b]$ denote the space of real-valued c\`adl\`ag functions on the compact interval $[a,b]$.
Equip $D$ with the uniform norm
\begin{equation*}
\|z\|_\infty := \sup_{t\in[a,b]} |z(t)|
\end{equation*}
and the associated uniform metric $d_\infty(z,w):=\|z-w\|_\infty$. Let $\mathcal B_{\mathrm{ball}}$ be the $\sigma$-field generated by all open balls
\begin{equation*}
B(z,r):=\{x\in D:\|x-z\|_\infty<r\}, \qquad z\in D,\ r>0,
\end{equation*}
that is, the Borel $\sigma$-field induced by $d_\infty$.
Let
\begin{equation*}
\mathcal C := \sigma(\pi_t:t\in[a,b]),
\qquad \pi_t(z):=z(t),
\end{equation*}
be the $\sigma$-field generated by the coordinate projections.
Prove that $\mathcal B_{\mathrm{ball}}=\mathcal C$.
\end{question}

\begin{proof}
We prove the two inclusions.

\medskip
\noindent\textbf{Step 1: $\mathcal C\subset\mathcal B_{\mathrm{ball}}$.}
It suffices to show that for every $t\in[a,b]$ and $c\in\mathbb R$, the set
\begin{equation*}
A_{t,c}:=\{z\in D:z(t)>c\}
\end{equation*}
belongs to $\mathcal B_{\mathrm{ball}}$, since such sets generate $\mathcal C$.

Fix $t$ and $c$. For each integer $n\ge1$, define
\begin{equation*}
z_{n,c}(s):=c+(n+n^{-1})\mathbbm{ 1}_{[t,t+n^{-1})}(s),
\qquad s\in[a,b],
\end{equation*}
which is a c\`adl\`ag function, and let
\begin{equation*}
B_n:=\{z\in D:\|z-z_{n,c}\|_\infty<n\}.
\end{equation*}
We claim that
\begin{equation*}
A_{t,c}=\bigcup_{n=1}^\infty B_n.
\end{equation*}

If $z\in B_n$, then
\begin{equation*}
|z(t)-z_{n,c}(t)|<n,
\end{equation*}
and since
\begin{equation*}
z_{n,c}(t)=c+n+n^{-1},
\end{equation*}
it follows that $z(t)>c$. Hence $B_n\subset A_{t,c}$.

Conversely, suppose $z(t)>c$. Set $\Delta:=z(t)-c>0$. By right-continuity of $z$ at $t$, there exists $\delta>0$ such that
\begin{equation*}
z(s)>c+\tfrac{\Delta}{2}, \qquad s\in[t,t+\delta).
\end{equation*}
Choose $n$ so large that
\begin{equation*}
n^{-1}<\min\{\delta,\tfrac{\Delta}{2}\}
\quad\text{and}\quad
n>\|z\|_\infty+|c|+1.
\end{equation*}
A direct check on the intervals $[t,t+n^{-1})$ and its complement shows that
\begin{equation*}
\|z-z_{n,c}\|_\infty<n,
\end{equation*}
hence $z\in B_n$.
This proves the claim and thus $A_{t,c}\in\mathcal B_{\mathrm{ball}}$.
Therefore $\mathcal C\subset\mathcal B_{\mathrm{ball}}$.

\medskip
\noindent\textbf{Step 2: $\mathcal B_{\mathrm{ball}}\subset\mathcal C$.}
It suffices to show that every open ball $B(x,r)$ belongs to $\mathcal C$.

Let
\begin{equation*}
\mathbb D := (\mathbb Q\cap[a,b])\cup\{a,b\},
\end{equation*}
which is a countable dense subset of $[a,b]$.
For $z\in D$, the function $t\mapsto|z(t)-x(t)|$ is c\`adl\`ag, and therefore
\begin{equation*}
\sup_{t\in[a,b]}|z(t)-x(t)|
=
\sup_{t\in\mathbb D}|z(t)-x(t)|.
\end{equation*}
Consequently,
\begin{equation*}
B(x,r)
=
\bigcap_{t\in\mathbb D}
\{z\in D:|z(t)-x(t)|<r\}.
\end{equation*}
The right-hand side is a countable intersection of sets depending on a single
coordinate projection $\pi_t$, hence belongs to $\mathcal C$.
Thus $B(x,r)\in\mathcal C$, and $\mathcal B_{\mathrm{ball}}\subset\mathcal C$.

\medskip
Combining the two inclusions yields
\begin{equation*}
\mathcal B_{\mathrm{ball}}=\mathcal C.
\end{equation*}
\end{proof}

\begin{question}
Equip $D[0,1]$ with the uniform metric
\begin{equation*}
d_\infty(z,w):=\sup_{t\in[0,1]}|z(t)-w(t)|.
\end{equation*}
Define $X:[0,1]\to D[0,1]$ by
\begin{equation*}
X(\omega):=\mathbbm{1}_{[\omega,1]}, \qquad \omega\in[0,1].
\end{equation*}
Let $\mathcal B([0,1])$ be the Borel $\sigma$-field on $[0,1]$.
Let $\mathcal B_{\mathrm{ball}}$ be the $\sigma$-field on $D[0,1]$ generated by all open balls
\begin{equation*}
B(z,r):=\{y\in D[0,1]:\|y-z\|_\infty<r\}.
\end{equation*}
Let $\mathcal B(D[0,1])$ denote the Borel $\sigma$-field on $D[0,1]$ induced by the topology of $d_\infty$.

Prove that $X$ is $\big(\mathcal B([0,1]),\mathcal B_{\mathrm{ball}}\big)$-measurable but not
$\big(\mathcal B([0,1]),\mathcal B(D[0,1])\big)$-measurable.
\end{question}

\begin{proof}
{\bfseries $X$ is ball measurable.} 
By definition, it suffices to show that for every $z\in D[0,1]$ and $r>0$,
the set $X^{-1}\big(B(z,r)\big)$ is a Borel subset of $[0,1]$.

Fix $z\in D[0,1]$ and $r>0$. Consider the function
\begin{equation*}
F(\omega):=\|X(\omega)-z\|_\infty=\sup_{t\in[0,1]}|\mathbbm{1}_{[\omega,1]}(t)-z(t)|.
\end{equation*}
We claim that $F$ is Borel measurable on $[0,1]$.

Let $\mathbb D:=\mathbb Q\cap[0,1]$, which is countable and dense. For each fixed $\omega$,
the map $t\mapsto |\mathbbm{1}_{[\omega,1]}(t)-z(t)|$ is c\`adl\`ag in $t$, hence its supremum
over $[0,1]$ equals its supremum over $\mathbb D$:
\begin{equation*}
F(\omega)=\sup_{t\in[0,1]}|\mathbbm{1}_{[\omega,1]}(t)-z(t)|
=\sup_{t\in\mathbb D}|\mathbbm{1}_{[\omega,1]}(t)-z(t)|.
\end{equation*}
For each fixed $t\in\mathbb D$, the map $\omega\mapsto \mathbbm{1}_{[\omega,1]}(t)$ equals
$\mathbbm{1}_{[0,t]}(\omega)$, hence it is Borel measurable. Therefore
$\omega\mapsto |\mathbbm{1}_{[\omega,1]}(t)-z(t)|$ is Borel measurable for each $t\in\mathbb D$.
Since a supremum of countably many Borel measurable functions is Borel measurable, $F$ is
Borel measurable. Consequently,
\begin{equation*}
X^{-1}\big(B(z,r)\big)=\{\omega\in[0,1]:F(\omega)<r\}\in\mathcal B([0,1]).
\end{equation*}
This proves that $X$ is $\big(\mathcal B([0,1]),\mathcal B_{\mathrm{ball}}\big)$-measurable.

\medskip
\noindent\textbf{$X$ is not Borel measurable.}
For each $s\in[0,1]$ define the open ball
\begin{equation*}
B_s:=B\!\left(\mathbbm{1}_{[s,1]},\tfrac12\right).
\end{equation*}
We first compute its preimage. For $\omega\in[0,1]$,
\begin{equation*}
X(\omega)\in B_s
\Longleftrightarrow
\|\mathbbm{1}_{[\omega,1]}-\mathbbm{1}_{[s,1]}\|_\infty<\tfrac12.
\end{equation*}
The difference $\mathbbm{1}_{[\omega,1]}-\mathbbm{1}_{[s,1]}$ takes only values in $\{-1,0,1\}$,
so the above inequality holds if and only if the difference is identically $0$, i.e.\ $\omega=s$.
Hence
\begin{equation*}
X^{-1}(B_s)=\{s\}.
\end{equation*}

Now let $S\subset[0,1]$ be arbitrary and set
\begin{equation*}
G:=\bigcup_{s\in S} B_s.
\end{equation*}
Since each $B_s$ is open in the uniform metric, $G$ is open, hence $G\in\mathcal B(D[0,1])$.
Moreover,
\begin{equation*}
X^{-1}(G)
=
\bigcup_{s\in S}X^{-1}(B_s)
=
\bigcup_{s\in S}\{s\}
=
S.
\end{equation*}
If $X$ were $\big(\mathcal B([0,1]),\mathcal B(D[0,1])\big)$-measurable, then $X^{-1}(G)$ would
belong to $\mathcal B([0,1])$ for every Borel set $G\subset D[0,1]$. In particular, for the
open set $G$ above we would have $S\in\mathcal B([0,1])$ for every $S\subset[0,1]$, which is
false. Therefore, $X$ is not Borel measurable.
\end{proof}

\begin{question}
Show that:

\begin{enumerate}
\item[(i)] If $I$ is an infinite index set, then $\ell^\infty(I)$ under the supremum norm is not separable.

\item[(ii)] The space $D[0,1]$ under the uniform norm is not separable.

\item[(iii)] (\emph{Borel versus projection $\sigma$-field under separability.})
Let $E$ be a \emph{separable} metric subspace of $\ell^\infty(\mathcal F)$ with metric induced by
\begin{equation*}
\|z\|_{\mathcal F}=\sup_{f\in\mathcal F}|z(f)|.
\end{equation*}
Assume there exists a \emph{countable} subset $\mathcal F_0=\{f_1,f_2,\dots\}\subset\mathcal F$
such that, for every $z,w\in E$,
\begin{equation*}
\|z-w\|_{\mathcal F}=\sup_{m\ge 1}|z(f_m)-w(f_m)|.
\end{equation*}
Let $\pi_m:E\to\mathbb R$ be the coordinate projections $\pi_m(z)=z(f_m)$.
Prove that the Borel $\sigma$-field on $E$ equals the projection (cylindrical) $\sigma$-field:
\begin{equation*}
\mathcal B(E)=\sigma(\pi_1,\pi_2,\dots).
\end{equation*}
\end{enumerate}
\end{question}

\begin{proof}

\textbf{(i) Non-separability of $\ell^\infty(I)$.}
Assume that $I$ is infinite. For each subset $A\subset I$, define
\begin{equation*}
x_A(i)=\mathbbm 1\{i\in A\}, \qquad i\in I.
\end{equation*}
Then $x_A\in \ell^\infty(I)$ and $\|x_A\|_\infty\le 1$. If $A\neq B$, pick $i\in A\triangle B$.
Then $|x_A(i)-x_B(i)|=1$, hence
\begin{equation*}
\|x_A-x_B\|_\infty=1.
\end{equation*}
Thus $\{x_A:A\subset I\}$ is an uncountable $1$-separated subset. A separable metric space
cannot contain an uncountable $\varepsilon$-separated subset for any $\varepsilon>0$.
Therefore,
\begin{equation*}
\ell^\infty(I)\ \text{is not separable.}
\end{equation*}

\medskip
\noindent
\textbf{(ii) Non-separability of $D[0,1]$ under $\|\cdot\|_\infty$.}
For each $t\in(0,1]$, define
\begin{equation*}
f_t(s)=\mathbbm 1\{s\ge t\}, \qquad s\in[0,1].
\end{equation*}
Each $f_t$ is càdlàg, so $f_t\in D[0,1]$. If $t<u$, then for $s\in[t,u)$,
\begin{equation*}
f_t(s)=1,\qquad f_u(s)=0,
\end{equation*}
hence
\begin{equation*}
\|f_t-f_u\|_\infty=1.
\end{equation*}
Therefore $\{f_t:t\in(0,1]\}$ is an uncountable $1$-separated subset of $D[0,1]$, implying
\begin{equation*}
D[0,1]\ \text{is not separable under the uniform norm.}
\end{equation*}

\medskip
\noindent
\textbf{(iii) $\mathcal B(E)$ equals the projection $\sigma$-field under separability and a countable determining set.}
Let
\begin{equation*}
\mathcal G:=\sigma(\pi_1,\pi_2,\dots).
\end{equation*}
Since each $\pi_m$ is continuous (hence Borel measurable), we have
\begin{equation*}
\mathcal G\subset \mathcal B(E).
\end{equation*}
It remains to prove $\mathcal B(E)\subset \mathcal G$. Because $\mathcal B(E)$ is generated by closed balls,
it suffices to show that every closed ball in $E$ belongs to $\mathcal G$.

Fix $z\in E$ and $\varepsilon>0$. By the assumption that $\mathcal F_0$ determines the metric on $E$,
\begin{equation*}
B(z,\varepsilon)
=
\{y\in E:\ \|y-z\|_{\mathcal F}\leq\varepsilon\}
=
\Big\{y\in E:\ \sup_{m\ge 1}|\pi_m(y)-\pi_m(z)|\leq\varepsilon\Big\}.
\end{equation*}
The condition $\sup_{m\ge 1} a_m\leq\varepsilon$ is equivalent to $a_m\leq\varepsilon$ for all $m$, hence
\begin{equation*}
B(z,\varepsilon)
=
\bigcap_{m=1}^\infty
\Big\{y\in E:\ |\pi_m(y)-\pi_m(z)|\leq\varepsilon\Big\}.
\end{equation*}
For each $m$, the set
\begin{equation*}
\Big\{y\in E:\ |\pi_m(y)-\pi_m(z)|<\varepsilon\Big\}
=
\pi_m^{-1}\big((\pi_m(z)-\varepsilon,\ \pi_m(z)+\varepsilon)\big)
\end{equation*}
belongs to $\mathcal G$, and $\mathcal G$ is closed under countable intersections. Therefore
\begin{equation*}
B(z,\varepsilon)\in \mathcal G.
\end{equation*}
Since open balls generate $\mathcal B(E)$, we conclude
\begin{equation*}
\mathcal B(E)\subset \mathcal G.
\end{equation*}
Combining with $\mathcal G\subset\mathcal B(E)$ yields
\begin{equation*}
\mathcal B(E)=\sigma(\pi_1,\pi_2,\dots),
\end{equation*}
as desired.
\end{proof}

\begin{question}{\emph{\bfseries (A.s.\ convergence of a measurable net need not imply convergence in probability)}}
Let $(\Omega,\mathcal F,\p)=([0,1],\mathcal B([0,1]),\lambda)$ be the Lebesgue probability space.
Fix a Borel set $B\subset[0,1]$ with $P(B)=1/2$. Let $D$ be the collection of all finite subsets
of $[0,1]$, directed by inclusion: $F\preceq G$ iff $F\subseteq G$.

For each $F\in D$, define the (measurable) random variable
\begin{equation*}
X_F(\omega):=\mathbf 1_{B\setminus F}(\omega),\qquad \omega\in[0,1],
\end{equation*}
and let $X\equiv 0$. Prove that the net $(X_F)_{F\in D}$ converges to $0$ almost surely (indeed,
pointwise), but does \emph{not} converge to $0$ in probability. Conclude that for general nets,
almost sure convergence (hence outer almost sure convergence) does not imply convergence in
probability, even when all $X_F$ are measurable.
\end{question}

\begin{proof}
Fix $\omega\in[0,1]$.
\begin{itemize}
	\item 
If $\omega\notin B$, then $\omega\notin B\setminus F$ for every $F\in D$, hence
\begin{equation*}
X_F(\omega)=0\qquad\text{for all }F\in D,
\end{equation*}
so $X_F(\omega)\to 0$.
	\item If $\omega\in B$, let $F_0:=\{\omega\}\in D$. For any $F\in D$ with $F\succeq F_0$ (i.e.\ $F\supseteq\{\omega\}$),
we have $\omega\in F$, hence $\omega\notin B\setminus F$, and therefore
\begin{equation*}
X_F(\omega)=\mathbf 1_{B\setminus F}(\omega)=0\qquad\text{for all }F\succeq F_0.
\end{equation*}
By the definition of convergence of a net, this implies $X_F(\omega)\to 0$.
\end{itemize}

Since the above holds for every $\omega\in[0,1]$, we have pointwise convergence
$X_F(\omega)\to 0$. In particular,
\begin{equation*}
X_F \to 0 \qquad \p\text{-}\text{almost surely}.
\end{equation*}

Fix any $F\in D$. Since $F$ is finite, $\p(F)=0$, hence
\begin{equation*}
\p(X_F=1)=\p(\omega\in B\setminus F)=\p(B)-\p(F)=\frac12-0=\frac12.
\end{equation*}
Therefore, for any $\varepsilon\in(0,1)$ (e.g.\ $\varepsilon=1/2$),
\begin{equation*}
\p(|X_F-0|>\varepsilon)=\p(X_F=1)=\frac12
\qquad\text{for all }F\in D.
\end{equation*}
In particular, $\p(|X_F|>\varepsilon)$ does not tend to $0$ along the directed set $D$.
Hence, $X_F$ does not converge to $0$ in probability.
\end{proof}

\begin{question}
	Let $(\Omega,\mathcal{F},P)$ be a probability space and let $Y:\Omega\to\mathbb{R}$ be an arbitrary function (not necessarily measurable).
Denote by $Y^*$ and $Y_*$ the upper and lower measurable envelopes of $Y$, respectively, where $Y_*:=-(-Y)^*$. Prove: $Y^*=-(-Y)_*$.
\end{question}
\begin{proof}
	We just need to use $Y_*:=-(-Y)^*$ and replace $Y$ with $-Y$, then
	\begin{equation*}
		(-Y)_*=-Y^*\Longleftrightarrow Y^*=-(-Y)_*.
	\end{equation*}
\end{proof}

\begin{question}[Control of measurable envelope gaps under Lipschitz transforms]
Let $(\Omega,\mathcal{F},P)$ be a probability space and let $Y:\Omega\to\mathbb{R}$ be an arbitrary function (not necessarily measurable).
Denote by $Y^*$ and $Y_*$ the upper and lower measurable envelopes of $Y$, respectively.

Let $\phi:\mathbb{R}\to\mathbb{R}$ be an $L$--Lipschitz function, i.e.,
\begin{equation*}
|\phi(a)-\phi(b)| \le L |a-b| \quad \text{for all } a,b\in\mathbb{R}.
\end{equation*}

\begin{enumerate}
\item[(i)] Show that the measurable envelope gap of $\phi(Y)$ satisfies
\begin{equation*}
\phi(Y)^* - \phi(Y)_* \le 2L\,(Y^*-Y_*).
\end{equation*}

\item[(ii)] If, in addition, $\phi$ is monotone nondecreasing (or nonincreasing), show that the sharper bound
\begin{equation*}
\phi(Y)^* - \phi(Y)_* \le L\,(Y^*-Y_*)
\end{equation*}
holds.
\end{enumerate}
\end{question}

\begin{proof}
First, for $L$--Lipschitz function $\phi$, we have
\begin{equation*}
	-L|a-b|\leq \phi(a)-\phi(b)\leq L|a-b|\quad \text{for all } a,b\in\mathbb{R}.
\end{equation*}
Then we prove the two parts separately.

\medskip
\noindent
\textbf{Part (i): general Lipschitz case.}

Since $Y_*\le Y \le Y^*$ pointwise, the Lipschitz property implies
\begin{equation*}
\begin{split}
	\phi(Y(\omega))&\leq \phi(Y_*(\omega)) + L\,(Y(\omega)-Y_*(\omega))\\
&\le
\phi(Y_*(\omega)) + L\,(Y^*(\omega)-Y_*(\omega))
\quad \text{for all } \omega\in\Omega.
\end{split}
\end{equation*}
Define
\begin{equation*}
U(\omega)
:=
\phi(Y_*(\omega)) + L\,(Y^*(\omega)-Y_*(\omega)).
\end{equation*}
The function $U$ is measurable because $Y_*$ and $Y^*$ are measurable and $\phi$ is continuous.
Moreover, $U \ge \phi(Y)$ pointwise, hence by the definition of the upper measurable envelope,
\begin{equation*}
\phi(Y)^* \le U.
\end{equation*}

Similarly, by the Lipschitz property,
\begin{equation*}
\begin{split}
	\phi(Y(\omega))&\ge
\phi(Y_*(\omega)) - L\,(Y(\omega)-Y_*(\omega))\\
&\ge
\phi(Y_*(\omega)) - L\,(Y^*(\omega)-Y_*(\omega))\quad \text{for all } \omega\in\Omega,
\end{split}
\end{equation*}
and defining
\begin{equation*}
L(\omega)
:=
\phi(Y_*(\omega)) - L\,(Y^*(\omega)-Y_*(\omega)),
\end{equation*}
we obtain a measurable function $L$ with $L\le \phi(Y)$ pointwise.
By the definition of the lower measurable envelope,
\begin{equation*}
\phi(Y)_* \ge L.
\end{equation*}

Subtracting the two bounds yields
\begin{equation*}
\phi(Y)^* - \phi(Y)_*
\le
\bigl[\phi(Y_*) + L(Y^*-Y_*)\bigr]
-
\bigl[\phi(Y_*) - L(Y^*-Y_*)\bigr]
=
2L\,(Y^*-Y_*),
\end{equation*}
which proves (i).

\medskip
\noindent
\textbf{Part (ii): monotone Lipschitz case (optimal constant).}

Assume now that $\phi$ is monotone nondecreasing.
Since $Y\le Y^*$ and $Y_*\le Y$ pointwise, monotonicity implies
\begin{equation*}
\phi(Y) \le \phi(Y^*),
\qquad
\phi(Y_*) \le \phi(Y).
\end{equation*}
Because $Y^*$ and $Y_*$ are measurable, both $\phi(Y^*)$ and $\phi(Y_*)$ are measurable.
Therefore, by the extremal properties of measurable envelopes,
\begin{equation*}
\phi(Y)^* \le \phi(Y^*),
\qquad
\phi(Y)_* \ge \phi(Y_*).
\end{equation*}
Consequently,
\begin{equation*}
\phi(Y)^* - \phi(Y)_*
\le
\phi(Y^*) - \phi(Y_*).
\end{equation*}

Finally, applying the $L$--Lipschitz property to the pair $(Y^*,Y_*)$ yields
\begin{equation*}
\phi(Y^*) - \phi(Y_*)
\le
L\,(Y^*-Y_*),
\end{equation*}
which proves the sharper bound in (ii) for nondecreasing functions. 

For the nonincreasing function $\phi$, $-\phi$ is nondecreasing. The result for nondecreasing function yields
\begin{equation*}
	(-\phi(Y))^* - (-\phi(Y))_*
\le L\,(Y^*-Y_*). \tag{$*$}
\end{equation*}
By definition, $\phi(Y)_*=-(-\phi(Y))^*$ and $\phi(Y)^*=-(-\phi(Y))_*$. Multiply ($*$) by $-1$ to obtain
\begin{equation*}
	-\phi(Y)_*+\phi(Y)^*\le L(Y_*-Y^*).
\end{equation*}
\end{proof}

\begin{question}
Let $\psi:[0,\infty)\to[0,\infty)$ be a Young function, that is,
$\psi$ is convex, $\psi(0)=0$, and $\psi\not\equiv 0$.
Show that:
\begin{enumerate}
\item $\psi$ is locally Lipschitz at every interior point (so continuous) and $\psi(t)/t$ is nondecreasing (so $\psi$ is nondecreasing) on $[0,\infty)$ and satisfies $\psi(t)\to\infty$ as $t\to\infty$;
\item the Luxemburg functional
\begin{equation*}
	\|X\|_\psi
:=
\inf\Big\{C>0:\ \mathbb E\,\psi(|X|/C)\le 1\Big\}
\end{equation*}
defines a norm on the Orlicz space $L_\psi=\{X:\|X\|_\psi<\infty\}$.
\end{enumerate}
\end{question}

\begin{proof}
\textbf{Step 1: Continuity and Monotonicity of $\psi$.}
Let $x_0>0$. Then there exists $\delta>0$ such that
\begin{equation*}
[x_0-\delta,x_0+\delta]\subset (0,\infty).
\end{equation*}
For any $x,y$ satisfying
\begin{equation*}
x_0-\delta < x < x_0 < y < x_0+\delta,
\end{equation*}
convexity implies the monotonicity of secant slopes:
\begin{equation*}
\frac{f(x_0)-f(x)}{x_0-x}
\;\le\;
\frac{f(y)-f(x_0)}{y-x_0}.
\end{equation*}
Define the endpoint slopes
\begin{equation*}
m_- := \frac{f(x_0)-f(x_0-\delta)}{\delta},
\qquad
m_+ := \frac{f(x_0+\delta)-f(x_0)}{\delta}.
\end{equation*}
Then for every $y\in(x_0-\delta,x_0+\delta)$ with $y\neq x_0$ we have
\begin{equation*}
m_- \le \frac{f(y)-f(x_0)}{y-x_0} \le m_+.
\end{equation*}
Let $M:=|m_-|\vee |m_+|$. Taking absolute value on the above inequalities yields
\begin{equation*}
|f(y)-f(x_0)| \le M|y-x_0| \qquad \text{for all } y\in(x_0-\delta,x_0+\delta).
\end{equation*}
Hence $f$ is locally Lipschitz at $x_0$, and in particular continuous at $x_0$.
Since $x_0$ was arbitrary, $f$ is continuous on $(0,\infty)$.

Let $0\le s<t$. Write
\begin{equation*}
s = \frac{s}{t} t + \Big(1-\frac{s}{t}\Big) 0.
\end{equation*}
By convexity and $\psi(0)=0$,
\begin{equation*}
\psi(s)
\le \frac{s}{t} \psi(t) + \Big(1-\frac{s}{t}\Big) \psi(0)= 
\frac{s}{t}\psi(t).
\end{equation*}
If $\psi(t)=0$, then the positiveness of $\psi$ yields $\psi(s)=0$.
If $\psi(t)>0$, then $s<t$ yields
\begin{equation*}
	\psi(s)\leq \frac{s}{t}\psi(t)\Longrightarrow \frac{\psi(s)}{s}<\frac{\psi(t)}{t}.
\end{equation*}
Combining the above two situations, we have 
\begin{equation*}
\frac{\psi(s)}{s}<\frac{\psi(t)}{t}, \quad\forall s<t.
\end{equation*}
Hence $\frac{\psi(t)}{t}$ is nondecreasing on $[0,\infty)$. Moreover, since
\begin{equation}
    \frac{\psi(s)}{s}<\frac{\psi(t)}{t}\Longrightarrow \psi(s)\leq \frac{s}{t}\psi(t)<\psi(t),\quad \forall s<t.
\end{equation}
Hence $\psi$ is also nondecreasing on $[0,\infty)$.

\medskip
\textbf{Step 2: Divergence at infinity.}
Since $\psi\not\equiv 0$, there exists $t_0>0$ such that $\psi(t_0)>0$.
For any $t>t_0$, convexity implies the monotonicity of secant slopes:
\begin{equation*}
\frac{\psi(t)-\psi(t_0)}{t-t_0}
\ge
\frac{\psi(t_0)-\psi(0)}{t_0-0}
=
\frac{\psi(t_0)}{t_0}.
\end{equation*}
Thus
\begin{equation*}
\psi(t)
\ge
\psi(t_0)+\frac{\psi(t_0)}{t_0}(t-t_0)
\ge
\frac{\psi(t_0)}{t_0}\,t,
\end{equation*}
which shows that $\psi(t)\to\infty$ as $t\to\infty$.

\medskip
\textbf{Step 3: Positivity of the Luxemburg functional.}
Let $X$ be a random variable with $P(|X|>0)>0$.
Then there exists $\varepsilon>0$ such that $P(|X|\ge\varepsilon)>0$.
For any $C>0$, using the monotonicity of $\psi$,
\begin{equation*}
\mathbb E\,\psi(|X|/C)
\ge
\mathbb E\big[\psi(|X|/C)\mathbf 1_{\{|X|\ge\varepsilon\}}\big]
\ge
\psi(\varepsilon/C)\,\mathbb P(|X|\ge\varepsilon).
\end{equation*}
By Step~2, $\psi(\varepsilon/C)\to\infty$ as $C\downarrow 0$, hence
$\mathbb E\,\psi(|X|/C)\to\infty$.
Therefore $\|X\|_\psi>0$ whenever $X\neq 0$ almost surely.

\medskip
\textbf{Step 4: Triangle inequality.}
Let $a,b>0$ satisfy
\begin{equation*}
\mathbb E\,\psi(|X|/a)\le 1,
\qquad
\mathbb E\,\psi(|Y|/b)\le 1.
\end{equation*}
Since 
\begin{equation*}
	\frac{|X+Y|}{a+b}\leq \frac{|X|+|Y|}{a+b}=\frac{a}{a+b}\frac{|X|}{a}
+
\frac{b}{a+b}\frac{|Y|}{b},
\end{equation*}
by monotonicity and convexity of $\psi$,
\begin{equation*}
\psi\!\left(
\frac{|X+Y|}{a+b}
\right)
\le
\psi\!\left(
\frac{a}{a+b}\frac{|X|}{a}
+
\frac{b}{a+b}\frac{|Y|}{b}
\right)
\le
\frac{a}{a+b}\psi(|X|/a)
+
\frac{b}{a+b}\psi(|Y|/b).
\end{equation*}
Taking expectations yields
\begin{equation*}
\mathbb E\,\psi\!\left(\frac{|X+Y|}{a+b}\right)\le 1,
\end{equation*}
so $a+b$ is admissible for $X+Y$ and
\begin{equation*}
\|X+Y\|_\psi \le a+b.
\end{equation*}
Taking the infimum over all such $a,b$ gives
\begin{equation*}
\|X+Y\|_\psi \le \|X\|_\psi + \|Y\|_\psi.
\end{equation*}

\medskip
Combining homogeneity, positivity, and the triangle inequality, the Luxemburg functional defines a norm on $L_\psi$.
\end{proof}

\begin{question}
Let $\psi:[0,\infty)\to[0,\infty)$ be a Young function, i.e.\ $\psi$ is convex,
$\psi(0)=0$, and $\psi\not\equiv 0$. For $C>0$, define $(C\psi)(t):=C\,\psi(t)$.
Show that for any random variable $X$,
\begin{equation*}
\|X\|_{C\psi}\le C\|X\|_{\psi}\quad\text{if } C\ge 1,
\end{equation*}
and that the reverse inequality
\begin{equation*}
\|X\|_{C\psi}\ge C\|X\|_{\psi}
\end{equation*}
holds if $0<C\le 1$.
\end{question}

\begin{proof}
Recall that the Luxemburg norm is defined by
\begin{equation*}
\|X\|_{\psi}
:=
\inf\Big\{a>0:\ \mathbb E\,\psi(|X|/a)\le 1\Big\}.
\end{equation*}

\textbf{Step 1 (scaling inequalities from convexity).}
Fix $t\ge 0$.
If $0\le \lambda\le 1$, then by convexity and $\psi(0)=0$,
\begin{equation*}
\psi(\lambda t)
=
\psi\big(\lambda t+(1-\lambda)0\big)
\le
\lambda\psi(t)+(1-\lambda)\psi(0)
=
\lambda\psi(t).
\end{equation*}
If $\alpha\ge 1$, apply the previous inequality with $\lambda=1/\alpha\in(0,1]$ to the point $\alpha t$ to obtain
\begin{equation*}
\psi(t)
=
\psi\Big(\frac{1}{\alpha}\,(\alpha t)\Big)
\le
\frac{1}{\alpha}\,\psi(\alpha t),
\end{equation*}
which is equivalent to
\begin{equation*}
\psi(\alpha t)\ge \alpha\,\psi(t).
\end{equation*}

\medskip
\textbf{Step 2 (case $C\ge 1$).}
Let $a=\|X\|_{\psi}$. By definition,
\begin{equation*}
\mathbb E\,\psi(|X|/a)\le 1.
\end{equation*}
Using Step~1 with $\lambda=1/C\le 1$, we have
\begin{equation*}
\psi\!\left(\frac{|X|}{Ca}\right)
=
\psi\!\left(\frac1C\,\frac{|X|}{a}\right)
\le
\frac1C\,\psi\!\left(\frac{|X|}{a}\right).
\end{equation*}
Taking expectations yields
\begin{equation*}
\mathbb E\,\psi\!\left(\frac{|X|}{Ca}\right)\le \frac1C,
\end{equation*}
and therefore
\begin{equation*}
\mathbb E\,(C\psi)\!\left(\frac{|X|}{Ca}\right)
=
C\,\mathbb E\,\psi\!\left(\frac{|X|}{Ca}\right)
\le 1.
\end{equation*}
Thus $Ca$ is admissible for $\|X\|_{C\psi}$, so
\begin{equation*}
\|X\|_{C\psi}\le Ca = C\|X\|_{\psi}.
\end{equation*}

\medskip
\textbf{Step 3 (case $0<C\le 1$).}
Let $a=\|X\|_{\psi}$ and fix any $b<a$. By the definition of the infimum,
\begin{equation*}
\mathbb E\,\psi(|X|/b)>1.
\end{equation*}
Since $1/C\ge 1$, Step~1 with $\alpha=1/C$ gives
\begin{equation*}
\psi\!\left(\frac{|X|}{Cb}\right)
=
\psi\!\left(\frac1C\,\frac{|X|}{b}\right)
\ge
\frac1C\,\psi\!\left(\frac{|X|}{b}\right).
\end{equation*}
Taking expectations and multiplying by $C$ yields
\begin{equation*}
\mathbb E\,(C\psi)\!\left(\frac{|X|}{Cb}\right)
=
C\,\mathbb E\,\psi\!\left(\frac{|X|}{Cb}\right)
\ge
\mathbb E\,\psi(|X|/b)
>1.
\end{equation*}
Hence no constant smaller than $Ca$ can satisfy the Luxemburg constraint for $C\psi$, and therefore
\begin{equation*}
\|X\|_{C\psi}\ge Ca = C\|X\|_{\psi}.
\end{equation*}

\medskip
Combining Steps~2 and~3 completes the proof.
\end{proof}

\begin{question}[Comparison of $L_p$- and $\psi_\alpha$-norms]
Let $\alpha\ge 1$ and $\psi_\alpha(x)=e^{x^\alpha}-1$ for $x\ge 0$. Define the Luxemburg--Orlicz norm
\begin{equation*}
\|X\|_{\psi_\alpha}:=\inf\Big\{t>0:\ \mathbb{E}\,\psi_\alpha(|X|/t)\le 1\Big\}.
\end{equation*}
Show that there exist constants $0<c_\alpha<C_\alpha<\infty$ depending only on $\alpha$ such that for every random variable $X$,
\begin{equation*}
c_\alpha \|X\|_{\psi_\alpha}
\;\le\;
\sup_{p\ge 1}\frac{\|X\|_p}{p^{1/\alpha}}
\;\le\;
C_\alpha \|X\|_{\psi_\alpha}.
\end{equation*}
In particular, $\|X\|_1\le \|X\|_2\le \|X\|_{\psi_2}$.
\end{question}

\begin{proof}
Throughout, $\|X\|_p=(\mathbb{E}|X|^p)^{1/p}$. Note that
\begin{equation*}
\mathbb{E}\,\psi_\alpha(|X|/t)\le 1
\quad\Longleftrightarrow\quad
\mathbb{E}\,e^{(|X|/t)^\alpha}\le 2,
\end{equation*}
since $\psi_\alpha(u)=e^{u^\alpha}-1$.

\medskip
We first prove that $\|X\|_{\psi_\alpha}$ is controlled by $\sup_{p\ge 1}\|X\|_p/p^{1/\alpha}$. Set
\begin{equation*}
D:=\sup_{p\ge 1}\frac{\|X\|_p}{p^{1/\alpha}},
\end{equation*}
and assume $D<\infty$. Then for every $p\ge 1$,
\begin{equation*}
\|X\|_p\le D\,p^{1/\alpha},
\qquad
\mathbb{E}|X|^p\le D^p p^{p/\alpha}.
\end{equation*}
Fix an integer $j\ge 1$ and choose $p=\alpha j\ge 1$. This gives
\begin{equation*}
\mathbb{E}|X|^{\alpha j}\le D^{\alpha j}(\alpha j)^j.
\end{equation*}
For any $C>0$, the series $\sum_{j\ge 0}\frac{(|X|/C)^{\alpha j}}{j!}$ is nonnegative and increasing in the truncation index, hence by monotone convergence,
\begin{equation*}
\mathbb{E}e^{(|X|/C)^\alpha}
=\mathbb{E}\sum_{j=0}^\infty \frac{|X|^{\alpha j}}{C^{\alpha j}j!}
=\sum_{j=0}^\infty \frac{\mathbb{E}|X|^{\alpha j}}{C^{\alpha j}j!}
=1+\sum_{j=1}^\infty \frac{\mathbb{E}|X|^{\alpha j}}{C^{\alpha j}j!}.
\end{equation*}
Using the moment bound yields
\begin{equation*}
\mathbb{E}e^{(|X|/C)^\alpha}
\le 1+\sum_{j=1}^\infty \left(\frac{D^\alpha}{C^\alpha}\right)^j\frac{(\alpha j)^j}{j!}
=1+\sum_{j=1}^\infty \left(\frac{\alpha D^\alpha}{C^\alpha}\right)^j\frac{j^j}{j!}.
\end{equation*}
We now use a concrete Stirling inequality. For all integers $j\ge 1$,
\begin{equation*}
j!\ge \sqrt{2\pi j}\Big(\frac{j}{e}\Big)^j,
\end{equation*}
and therefore
\begin{equation*}
\frac{j^j}{j!}\le \frac{e^j}{\sqrt{2\pi j}}.
\end{equation*}
Substituting this estimate gives
\begin{equation*}
\mathbb{E}e^{(|X|/C)^\alpha}
\le 1+\sum_{j=1}^\infty \left(\frac{\alpha D^\alpha}{C^\alpha}\right)^j\frac{e^j}{\sqrt{2\pi j}}
\le 1+\frac{1}{\sqrt{2\pi}}\sum_{j=1}^\infty r^j,
\qquad
r:=\frac{e\alpha D^\alpha}{C^\alpha}.
\end{equation*}
Choose $C=(2e\alpha)^{1/\alpha}D$, so that $r=1/2$ and hence $\sum_{j=1}^\infty r^j=1$. Then
\begin{equation*}
\mathbb{E}e^{(|X|/C)^\alpha}\le 1+\frac{1}{\sqrt{2\pi}}<2,
\end{equation*}
so $\mathbb{E}\psi_\alpha(|X|/C)\le 1$, which implies
\begin{equation*}
\|X\|_{\psi_\alpha}\le C=(2e\alpha)^{1/\alpha}D.
\end{equation*}
Equivalently,
\begin{equation*}
(2e\alpha)^{-1/\alpha}\,\|X\|_{\psi_\alpha}
\le
\sup_{p\ge 1}\frac{\|X\|_p}{p^{1/\alpha}}.
\end{equation*}

\medskip
We now prove the reverse bound. Let $D=\|X\|_{\psi_\alpha}$. By definition,
\begin{equation*}
\mathbb{E}e^{(|X|/D)^\alpha}\le 2.
\end{equation*}
For any $x\ge 0$, Markov's inequality applied to $e^{(|X|/D)^\alpha}$ yields
\begin{equation*}
\mathbb{P}(|X|>x)
=\mathbb{P}\!\left(e^{(|X|/D)^\alpha}>e^{(x/D)^\alpha}\right)
\le e^{-(x/D)^\alpha}\,\mathbb{E}e^{(|X|/D)^\alpha}
\le 2e^{-(x/D)^\alpha}.
\end{equation*}
We next justify the tail integral identity for moments. For $p\ge 1$,
\begin{equation*}
|X|^p=\int_0^{|X|} p t^{p-1}\,dt,
\end{equation*}
hence by Tonelli's theorem (nonnegative integrand),
\begin{equation*}
\mathbb{E}|X|^p
=\mathbb{E}\int_0^{|X|} p t^{p-1}\,dt
=\int_0^\infty p t^{p-1}\,\mathbb{P}(|X|>t)\,dt.
\end{equation*}
Using the tail bound,
\begin{equation*}
\mathbb{E}|X|^p
\le \int_0^\infty p t^{p-1}\cdot 2e^{-(t/D)^\alpha}\,dt
=2p\int_0^\infty t^{p-1}e^{-(t/D)^\alpha}\,dt.
\end{equation*}
We compute this integral by a change of variables. Let $u=(t/D)^\alpha$, so that $t= D u^{1/\alpha}$ and
\begin{equation*}
dt = D\frac{1}{\alpha}u^{1/\alpha-1}\,du.
\end{equation*}
Then
\begin{equation*}
t^{p-1}dt
=\big(D u^{1/\alpha}\big)^{p-1}\cdot D\frac{1}{\alpha}u^{1/\alpha-1}\,du
= D^p\frac{1}{\alpha}u^{p/\alpha-1}\,du,
\end{equation*}
and therefore
\begin{equation*}
\mathbb{E}|X|^p
\le 2p\cdot D^p\frac{1}{\alpha}\int_0^\infty u^{p/\alpha-1}e^{-u}\,du
=\frac{2p}{\alpha}D^p\,\Gamma(p/\alpha).
\end{equation*}
Using $\Gamma(z+1)=z\Gamma(z)$ with $z=p/\alpha$ gives
\begin{equation*}
\frac{p}{\alpha}\Gamma(p/\alpha)=\Gamma(p/\alpha+1),
\qquad
\mathbb{E}|X|^p\le 2D^p\,\Gamma(p/\alpha+1).
\end{equation*}
Taking $p$-th roots,
\begin{equation*}
\|X\|_p\le D\cdot 2^{1/p}\,\Gamma(p/\alpha+1)^{1/p}.
\end{equation*}

It remains to bound $\Gamma(p/\alpha+1)^{1/p}$ by a constant multiple of $p^{1/\alpha}$. We use a concrete Stirling-type upper bound for the Gamma function: there exists an absolute constant $K>0$ such that for all $y\ge 1$,
\begin{equation*}
\Gamma(y+1)\le K\,y^{y+1/2}e^{-y}.
\end{equation*}
Apply this with $y=p/\alpha$. Since $p\ge 1$ and $\alpha\ge 1$, we have $y=p/\alpha\ge 1/\alpha$, and in particular for all $p\ge \alpha$ one has $y\ge 1$. (For $1\le p<\alpha$, we will absorb finitely many values of $p$ into the constant at the end.) Thus for $p\ge \alpha$,
\begin{equation*}
\Gamma(p/\alpha+1)
\le K\left(\frac{p}{\alpha}\right)^{p/\alpha+1/2}e^{-p/\alpha}.
\end{equation*}
Taking $p$-th roots yields
\begin{equation*}
\Gamma(p/\alpha+1)^{1/p}
\le K^{1/p}\left(\frac{p}{\alpha}\right)^{1/\alpha+1/(2p)}e^{-1/\alpha}.
\end{equation*}
Since $p\ge 1$, we have $K^{1/p}\le \max\{K,1\}$, and also
\begin{equation*}
\left(\frac{p}{\alpha}\right)^{1/(2p)}
= \exp\!\left(\frac{1}{2p}\log\frac{p}{\alpha}\right)
\le \exp\!\left(\frac{1}{2}\sup_{u\ge 1}\frac{\log u}{u}\right)
= \exp\!\left(\frac{1}{2e}\right),
\end{equation*}
because $\sup_{u\ge 1}(\log u)/u=1/e$. Consequently, for $p\ge \alpha$,
\begin{equation*}
\Gamma(p/\alpha+1)^{1/p}\le C'_\alpha\,p^{1/\alpha},
\end{equation*}
where one may take
\begin{equation*}
C'_\alpha:= e^{-1/\alpha}\,\alpha^{-1/\alpha}\,\max\{K,1\}\,\exp\!\left(\frac{1}{2e}\right).
\end{equation*}
For the remaining range $1\le p<\alpha$, note that there are only finitely many such $p$, and $\Gamma(p/\alpha+1)^{1/p}$ is finite for each; hence we may increase the constant if necessary so that the bound holds for all $p\ge 1$. Therefore there exists $C_\alpha<\infty$ depending only on $\alpha$ such that for every $p\ge 1$,
\begin{equation*}
\Gamma(p/\alpha+1)^{1/p}\le C_\alpha\,p^{1/\alpha}.
\end{equation*}
Combining with $\|X\|_p\le D\cdot 2^{1/p}\Gamma(p/\alpha+1)^{1/p}$ and using $2^{1/p}\le 2$ gives
\begin{equation*}
\|X\|_p\le 2C_\alpha\,D\,p^{1/\alpha}.
\end{equation*}
Taking the supremum over $p\ge 1$ yields
\begin{equation*}
\sup_{p\ge 1}\frac{\|X\|_p}{p^{1/\alpha}}
\le 2C_\alpha\,\|X\|_{\psi_\alpha}.
\end{equation*}
Renaming $2C_\alpha$ as $C_\alpha$ completes the two-sided comparison.

\medskip
Finally, the particular inequalities follow. By Cauchy--Schwarz,
\begin{equation*}
\|X\|_1=\mathbb{E}(|X|\cdot 1)\le (\mathbb{E}|X|^2)^{1/2}(\mathbb{E}1^2)^{1/2}=\|X\|_2.
\end{equation*}
Also, for every $u\ge 0$, $e^{u^2}\ge 1+u^2$ implies $u^2\le e^{u^2}-1=\psi_2(u)$. If $t>0$ satisfies $\mathbb{E}\psi_2(|X|/t)\le 1$, then
\begin{equation*}
\mathbb{E}\left(\frac{|X|}{t}\right)^2 \le \mathbb{E}\psi_2(|X|/t)\le 1,
\end{equation*}
so $\|X\|_2\le t$, and taking the infimum over such $t$ gives $\|X\|_2\le \|X\|_{\psi_2}$.
\end{proof}

\begin{question}[Comparing $\psi_p$-norms]
For $p\ge 1$, let $\psi_p(x)=e^{x^p}-1$ for $x\ge 0$, and define the Luxemburg--Orlicz norm
\begin{equation*}
\|X\|_{\psi_p}
:=\inf\Big\{t>0:\ \mathbb{E}\,\psi_p(|X|/t)\le 1\Big\}.
\end{equation*}
Show that the map $p\mapsto (\log 2)^{1/p}\|X\|_{\psi_p}$ is nondecreasing on $[1,\infty)$, i.e., for all $1\le p\le q<\infty$,
\begin{equation*}
(\log 2)^{1/p}\|X\|_{\psi_p}\le (\log 2)^{1/q}\|X\|_{\psi_q}.
\end{equation*}
\end{question}

\begin{proof}
Fix $1\le p\le q$ and set $\theta:=p/q\in(0,1]$. Let $B:=\|X\|_{\psi_q}$. By definition of the Luxemburg norm,
\begin{equation*}
\mathbb{E}\,\psi_q(|X|/B)\le 1,
\end{equation*}
that is,
\begin{equation*}
\mathbb{E}\big(e^{(|X|/B)^q}-1\big)\le 1
\qquad\Longleftrightarrow\qquad
\mathbb{E}e^{(|X|/B)^q}\le 2.
\end{equation*}
Define
\begin{equation*}
A:=(\log 2)^{1/q-1/p}\,B.
\end{equation*}
We will prove $\mathbb{E}\,\psi_p(|X|/A)\le 1$, which implies $\|X\|_{\psi_p}\le A$ and hence
\begin{equation*}
(\log 2)^{1/p}\|X\|_{\psi_p}\le (\log 2)^{1/p}A=(\log 2)^{1/q}\|X\|_{\psi_q}.
\end{equation*}

Introduce the auxiliary function $\phi:[0,\infty)\to[0,\infty)$ by
\begin{equation*}
\phi(u):=\exp\Big((\log 2)^{1-\theta}\big(\log(1+u)\big)^{\theta}\Big)-1.
\end{equation*}
First we record the identity linking $\psi_p$ and $\psi_q$ through $\phi$. For $x\ge 0$ set $u:=\psi_q(x)=e^{x^q}-1$. Then $1+u=e^{x^q}$, hence
\begin{equation*}
x^q=\log(1+u),
\qquad
x^p=(x^q)^{p/q}=\big(\log(1+u)\big)^{\theta}.
\end{equation*}
Moreover,
\begin{equation*}
\big((\log 2)^{1/p-1/q}\big)^p=(\log 2)^{1-p/q}=(\log 2)^{1-\theta}.
\end{equation*}
Therefore,
\begin{equation*}
\begin{split}
	\psi_p\big((\log 2)^{1/p-1/q}x\big)
&=\exp\Big((\log 2)^{1-\theta}x^p\Big)-1\\
&=\exp\Big((\log 2)^{1-\theta}\big(\log(1+u)\big)^\theta\Big)-1
=\phi(u).
\end{split}
\end{equation*}
Since $u=\psi_q(x)$, this gives
\begin{equation*}
\psi_p\big((\log 2)^{1/p-1/q}x\big)=\phi(\psi_q(x)),
\qquad x\ge 0.
\end{equation*}
In particular,
\begin{equation*}
\phi(1)
=\exp\Big((\log 2)^{1-\theta}\big(\log 2\big)^\theta\Big)-1
=\exp(\log 2)-1
=1.
\end{equation*}

We next prove that $\phi$ is concave on $[0,\infty)$. Let $u\ge 0$ and denote
\begin{equation*}
L:=\log(1+u),
\qquad
a:=(\log 2)^{1-\theta}.
\end{equation*}
Then
\begin{equation*}
\phi(u)=e^{aL^\theta}-1.
\end{equation*}
Differentiate using the chain rule. Since $dL/du=1/(1+u)$,
\begin{equation*}
\phi'(u)
=e^{aL^\theta}\cdot a\theta L^{\theta-1}\cdot \frac{1}{1+u}.
\end{equation*}
Write $E:=e^{aL^\theta}$ and $A_0:=a\theta L^{\theta-1}$, so that
\begin{equation*}
\phi'(u)=\frac{EA_0}{1+u}>0.
\end{equation*}
So, $\phi$ is increasing. Moreover, differentiate again:
\begin{equation*}
\phi''(u)=\frac{(EA_0)'}{1+u}-\frac{EA_0}{(1+u)^2}.
\end{equation*}
Compute $(EA_0)'=E'A_0+EA_0'$. Using $E'=E\cdot a\theta L^{\theta-1}\cdot \frac{1}{1+u}=E\cdot \frac{A_0}{1+u}$ and
\begin{equation*}
A_0' = a\theta(\theta-1)L^{\theta-2}\cdot \frac{1}{1+u},
\end{equation*}
we obtain
\begin{equation*}
(EA_0)'=E\cdot \frac{A_0}{1+u}\cdot A_0 + E\cdot a\theta(\theta-1)L^{\theta-2}\cdot \frac{1}{1+u}.
\end{equation*}
Substituting into the expression for $\phi''(u)$ gives
\begin{equation*}
\phi''(u)
=\frac{E}{(1+u)^2}\Big(A_0^2 + a\theta(\theta-1)L^{\theta-2} - A_0\Big).
\end{equation*}
Factor the bracket by substituting $A_0=a\theta L^{\theta-1}$:
\begin{equation*}
A_0^2 - A_0 + a\theta(\theta-1)L^{\theta-2}
=a\theta L^{\theta-2}\Big(a\theta L^\theta - L + (\theta-1)\Big).
\end{equation*}
Hence
\begin{equation*}
\phi''(u)
=
\frac{E}{(1+u)^2}\,a\theta L^{\theta-2}\,h(L),
\qquad
h(L):=a\theta L^\theta - L + (\theta-1).
\end{equation*}
All factors outside $h(L)$ are nonnegative for $u\ge 0$ (since then $E>0$, $1+u>0$, $a\ge 0$, $\theta>0$, and $L=\log(1+u)\ge 0$), so $\phi''(u)\le 0$ follows once we prove
\begin{equation*}
h(L)\le 0
\qquad\text{for all }L\ge 0.
\end{equation*}

To bound $h$, differentiate:
\begin{equation*}
h'(L)=a\theta^2 L^{\theta-1}-1.
\end{equation*}
If $\theta=1$ then $h(L)=a\cdot 1\cdot L - L + 0=0$ for all $L$ and hence $\phi''(u)=0$. Assume $0<\theta<1$.
Since $\theta-1<0$, the function $L\mapsto L^{\theta-1}$ is strictly decreasing on $(0,\infty)$, so $h'(L)$ is strictly decreasing on $(0,\infty)$. Therefore $h$ has at most one critical point, which (if it exists) must be its global maximizer. Solve $h'(L)=0$:
\begin{equation*}
a\theta^2 L^{\theta-1}=1
\qquad\Longleftrightarrow\qquad
L_0=(a\theta^2)^{1/(1-\theta)}.
\end{equation*}
Using $a\theta^2 L_0^{\theta-1}=1$, we have $L_0^\theta = L_0/(a\theta^2)$, hence
\begin{equation*}
h(L_0)
=a\theta\frac{L_0}{a\theta^2}-L_0+(\theta-1)
=\frac{L_0}{\theta}-L_0+(\theta-1)
=(1-\theta)\Big(\frac{L_0}{\theta}-1\Big).
\end{equation*}
Thus $h(L_0)\le 0$ is equivalent to $L_0\le \theta$, and this holds because
\begin{equation*}
L_0\le \theta
\Longleftrightarrow
(a\theta^2)^{1/(1-\theta)}\le \theta
\Longleftrightarrow
a\theta^2\le \theta^{1-\theta}
\Longleftrightarrow
a\le \theta^{-1-\theta}.
\end{equation*}
Since $\log 2<1$ and $1-\theta\ge 0$, we have $a=(\log 2)^{1-\theta}\le 1$, while $\theta\in(0,1]$ implies $\theta^{-1-\theta}\ge 1$. Therefore $a\le \theta^{-1-\theta}$, so $h(L_0)\le 0$. Because $L_0$ is the (unique) maximizer of $h$, it follows that $h(L)\le h(L_0)\le 0$ for all $L\ge 0$. Consequently $\phi''(u)\le 0$ for all $u\ge 0$, and $\phi$ is concave on $[0,\infty)$.

Now define the nonnegative random variable
\begin{equation*}
Y:=\psi_q(|X|/B)=e^{(|X|/B)^q}-1.
\end{equation*}
Then $\mathbb{E}Y\le 1$. Applying the identity $\psi_p((\log 2)^{1/p-1/q}x)=\phi(\psi_q(x))$ with $x=|X|/B$ gives
\begin{equation*}
\psi_p\big((\log 2)^{1/p-1/q}|X|/B\big)=\phi(Y).
\end{equation*}
Since $\phi$ is concave, Jensen's inequality (concave form) yields
\begin{equation*}
\mathbb{E}\,\phi(Y)\le \phi(\mathbb{E}Y)\le \phi(1)=1,
\end{equation*}
where we used $\mathbb{E}Y\le 1$ and the fact that $\phi$ is increasing (which follows from $\phi'(u)\ge 0$).
Therefore
\begin{equation*}
\mathbb{E}\,\psi_p\big((\log 2)^{1/p-1/q}|X|/B\big)\le 1.
\end{equation*}
Because $A=(\log 2)^{1/q-1/p}B$, the left-hand side equals $\mathbb{E}\,\psi_p(|X|/A)$. Hence
\begin{equation*}
\mathbb{E}\,\psi_p(|X|/A)\le 1,
\end{equation*}
and by the definition of the Luxemburg norm,
\begin{equation*}
\|X\|_{\psi_p}\le A=(\log 2)^{1/q-1/p}\|X\|_{\psi_q}.
\end{equation*}
Multiplying by $(\log 2)^{1/p}$ yields
\begin{equation*}
(\log 2)^{1/p}\|X\|_{\psi_p}\le (\log 2)^{1/q}\|X\|_{\psi_q},
\end{equation*}
which proves that $p\mapsto (\log 2)^{1/p}\|X\|_{\psi_p}$ is nondecreasing.
\end{proof}

\begin{proposition}[Orlicz norms and monotone convergence]
Let $\psi$ be a convex, nondecreasing, nonzero function on $[0,\infty)$ with
$\psi(0)=0$. Suppose that $0 \le X_n \uparrow X$ almost surely. Then
\begin{equation*}
	\|X_n\|_\psi \uparrow \|X\|_\psi .
\end{equation*}
\end{proposition}

\begin{proof}
Recall the Luxemburg--Orlicz norm
\begin{equation*}
\|Y\|_\psi
=
\inf\Bigl\{c>0:\ \mathbb E\,\psi(|Y|/c)\le 1\Bigr\}.
\end{equation*}
Since $0\le X_n \le X$ almost surely and $\psi$ is nondecreasing, for any $c>0$,
\begin{equation*}
\mathbb E\,\psi(X_n/c)\le \mathbb E\,\psi(X/c).
\end{equation*}
Hence, whenever $\mathbb E\,\psi(X/c)\le 1$, we also have
$\mathbb E\,\psi(X_n/c)\le 1$, which implies
\begin{equation*}
\|X_n\|_\psi \le \|X\|_\psi .
\end{equation*}
Moreover, since $X_n \uparrow X$ almost surely, the sequence
$\|X_n\|_\psi$ is nondecreasing. Therefore there exists a limit
\begin{equation*}
\|X_n\|_\psi \uparrow L \le \|X\|_\psi .
\end{equation*}


Fix any $r\in(0,1)$. By the definition of the Luxemburg
norm, we have
\begin{equation*}
\mathbb E\,\psi(X/\|X\|_\psi)\le 1.
\end{equation*}
On the other hand, since $ra<a$, the minimality of $a$ implies
\begin{equation*}
\mathbb E\,\psi\!\left(\frac{X}{r\|X\|_\psi}\right)>1.
\end{equation*}

Because $X_n \uparrow X$ almost surely and $\psi$ is nondecreasing, we have
\begin{equation*}
\psi\!\left(\frac{X_n}{r\|X\|_\psi}\right)
\uparrow
\psi\!\left(\frac{X}{r\|X\|_\psi}\right)
\quad\text{almost surely}.
\end{equation*}
By the monotone convergence theorem,
\begin{equation*}
\lim_{n\to\infty}
\mathbb E\,\psi\!\left(\frac{X_n}{r\|X\|_\psi}\right)
=
\mathbb E\,\psi\!\left(\frac{X}{r\|X\|_\psi}\right)
>1.
\end{equation*}
Hence, for all sufficiently large $n$,
\begin{equation*}
\mathbb E\,\psi\!\left(\frac{X_n}{r\|X\|_\psi}\right)>1,
\end{equation*}
which by definition of the Luxemburg norm yields
\begin{equation*}
\|X_n\|_\psi > r\|X\|_\psi.
\end{equation*}
Taking $n\to\infty$ gives
\begin{equation*}
L \ge r\|X\|_\psi.
\end{equation*}
Since $r\in(0,1)$ is arbitrary, letting $r\uparrow 1$ yields
\begin{equation*}
L \ge  \|X\|_\psi .
\end{equation*}
\end{proof}

\begin{question}
Let $(\Omega ,\f ,\p )$ and $( E ,\mathscr{E},\mathbb{Q})$ be probability spaces.
Let $\phi:\Omega \to E $ be a bijection such that both $\phi$ and $\phi^{-1}$ are measurable.
Assume moreover that $\phi$ preserves null sets, i.e.
\begin{equation*}
\mathbb{Q}(A)=0 \ \Longleftrightarrow\ \p (\phi^{-1}(A))=0
\qquad \forall A\in\mathscr{E} .
\end{equation*}
For an arbitrary mapping $T: E \to[-\infty,\infty]$ (not assumed measurable), write its measurable envelope as $T^*$ and similarly define $(T\circ\phi)^*$ on $(\Omega ,\f )$.
Prove that $\phi$ is \emph{perfect}, i.e.
\begin{equation*}
(T\circ\phi)^* = T^*\circ\phi \quad \p\text{-a.s. }.
\end{equation*}
\end{question}

\begin{proof}
Fix an arbitrary mapping $T: E \to[-\infty,\infty]$.

\medskip
\noindent
\textbf{(1) The inequality $(T\circ\phi)^*\le T^*\circ\phi$ holds $\p$-a.s..}
Since $T^*$ is $\mathscr{E}$-measurable and satisfies $T^*\ge T$, the composition
$T^*\circ\phi$ is $\f $-measurable (because $\phi$ is measurable) and satisfies
$T^*\circ\phi\ge T\circ\phi$ pointwise in $\Omega$. Hence, $T^*\circ\phi$ is an admissible measurable majorant
in the definition of $(T\circ\phi)^*$, which yields
\begin{equation*}
(T\circ\phi)^*\le T^*\circ\phi,\qquad \p\text{-a.s.}.
\end{equation*}

\medskip
\noindent
\textbf{(2) The reverse inequality $T^*\circ\phi\le (T\circ\phi)^*$ holds $\p $-a.s.}
Let $U:\Omega \to[-\infty,\infty]$ be any $\f $-measurable function with
$U\ge T\circ\phi$ pointwise.
Then $U\circ\phi^{-1}$ is $\mathscr{E}$-measurable (since $\phi^{-1}$ is measurable), and 
\begin{equation*}
U\circ\phi^{-1}\ge T\circ\phi\circ\phi^{-1}=T,\qquad\text{ pointwise in }E .
\end{equation*}
Thus $U\circ\phi^{-1}$ is a measurable majorant of $T$ on $( E ,\mathscr{E})$.
By minimality of $T^*$ in the $\mathbb Q$-a.s.\ sense, there exists $\mathscr{N}\in\mathscr{E}$ with
$\mathbb{Q}(\mathscr{N})=0$ such that
\begin{equation*}
T^*( e)\le (U\circ\phi^{-1})( e ),\qquad \forall  e \in E \setminus \mathscr{N}.
\end{equation*}
Composing with $\phi$ gives, 
\begin{equation*}
(T^*\circ\phi)(\omega )\le (U\circ\phi^{-1}\circ\phi)(\omega )=U(\omega ),\qquad\forall \omega \in\Omega \setminus\phi^{-1}(\mathscr{N}).
\end{equation*}
Since $\phi$ preserves null sets, $\p (\phi^{-1}(\mathscr{N}))=0$, hence
\begin{equation*}
T^*\circ\phi \le U \qquad \p \text{-a.s.}
\end{equation*}
Because this holds for every measurable $U\ge T\circ\phi$, taking $U$ to be $(T\circ \phi)^*$
yields
\begin{equation*}
T^*\circ\phi \le (T\circ\phi)^* \qquad \p \text{-a.s.}
\end{equation*}
\end{proof}
  
  \begin{question}
Let $(\Omega,\mathscr F,\mathbb P)$ be a probability space and let
$(\Omega^n,\mathscr F^{n},\mathbb P^{n})$ be the $n$-fold product space.
Fix a permutation $\sigma\in S_n$ and define the coordinate-permutation map
\begin{equation*}
\phi_\sigma:\Omega^n\to\Omega^n,\qquad
\phi_\sigma(\Omega ,\dots,\omega_n)=(\omega_{\sigma(1)},\dots,\omega_{\sigma(n)}).
\end{equation*}
Show that $\phi_\sigma$ is a measure-preserving isomorphism (hence an isomorphism mod $0$),
that it is \emph{perfect} in the sense that for every mapping $T:\Omega^n\to[-\infty,\infty]$,
\begin{equation*}
(T\circ\phi_\sigma)^* = T^*\circ\phi_\sigma\quad \text{a.s.},
\end{equation*}
and that it preserves outer expectation, i.e. if $\e^* T<\infty$, then
\begin{equation*}
\mathbb E^*(T\circ\phi_\sigma)=\mathbb E^*T.
\end{equation*}
Here $T^*$ denotes the (outer) measurable envelope (smallest measurable majorant) of $T$.
\end{question}

\begin{proof}
\noindent \textbf{(1) $\phi_\sigma$ is a bimeasurable bijection.}
For any measurable rectangle $R=\prod_{i=1}^n A_i$ with $A_i\in\mathscr F$, we have
\begin{equation*}
\phi_\sigma^{-1}(R)
=\{ \omega\in\Omega^n:\ \omega_{\sigma(i)}\in A_i\ \forall i\}
=\prod_{j=1}^n A_{\sigma^{-1}(j)}\in \mathscr F^{n}.
\end{equation*}
Since rectangles generate $\mathscr F^{n}$, this implies that $\phi_\sigma$ is
$\mathscr F^{n}/\mathscr F^{n}$-measurable. The map is bijective with
inverse $\phi_{\sigma^{-1}}$, and the same argument shows $\phi_{\sigma^{-1}}$ is measurable.
Hence $\phi_\sigma$ is a bimeasurable bijection.

\medskip
\noindent \textbf{(2) $\phi_\sigma$ is measure-preserving.}
First verify the identity on measurable rectangles:
for $R=\prod_{i=1}^n A_i$,
\begin{equation*}
\mathbb P^{n}(\phi_\sigma^{-1}(R))
=\mathbb P^{n}\Bigl(\prod_{j=1}^n A_{\sigma^{-1}(j)}\Bigr)
=\prod_{j=1}^n \mathbb P(A_{\sigma^{-1}(j)})
=\prod_{i=1}^n \mathbb P(A_i)
=\mathbb P^{n}(R).
\end{equation*}
Define
\begin{equation*}
\mathcal C:=\bigl\{B\in\mathscr F^{n}:\ \mathbb P^{n}(\phi_\sigma^{-1}(B))
=\mathbb P^{n}(B)\bigr\}.
\end{equation*}
Then $\mathcal C$ is a $\lambda$-system: $\Omega^n\in\mathcal C$, if $B\in\mathcal C$ then
$\Omega^n\setminus B\in\mathcal C$, and if $B_k\uparrow B$ with each $B_k\in\mathcal C$ then
by continuity from below and measurability of $\phi_\sigma$,
\begin{equation*}
\mathbb P^{n}(\phi_\sigma^{-1}(B))
=\lim_{k\to\infty}\mathbb P^{n}(\phi_\sigma^{-1}(B_k))
=\lim_{k\to\infty}\mathbb P^{n}(B_k)
=\mathbb P^{n}(B).
\end{equation*}
Let $\mathcal P$ be the $\pi$-system of measurable rectangles; we already showed
$\mathcal P\subset\mathcal C$. By the $\pi$-$\lambda$ theorem,
$\sigma(\mathcal P)=\mathscr F^{n}\subset\mathcal C$, hence $\mathcal C=\mathscr F^{n}$.
Therefore,
\begin{equation*}
\mathbb P^{n}(\phi_\sigma^{-1}(B))=\mathbb P^{n}(B)
\quad\forall B\in\mathscr F^{n}.
\end{equation*}
So $\phi_\sigma$ is measure-preserving. Combined with Step 1, $\phi_\sigma$ is a
measure-preserving isomorphism (hence also an isomorphism mod $0$).

\medskip
\noindent \textbf{(3) $\phi_\sigma$ is perfect, i.e. $(T\circ\phi_\sigma)^* = T^*\circ\phi_\sigma$ a.s.} Based on (1), $\phi:\Omega^n \to \Omega^n$ is a bijection such that both $\phi$ and $\phi^{-1}$ are measurable. Moreover, $\phi$ preserves null sets in {\bfseries both directions}, which follows from (2). So, the previous QUESTION shows that $\phi_\sigma$ is perfect.

\medskip
\noindent \textbf{(4) outer expectation is preserved.}
Outer expectation is defined by
\begin{equation*}
\mathbb E^*T:=\inf\{\mathbb EZ: Z\text{ is } \mathscr F^{n}\text{-measurable and } Z\ge T\}.
\end{equation*}
Now we show $\mathbb E^*(T\circ\phi_\sigma)=\mathbb E^*T$. The perfectness and the measure-preserving property of $\phi_\sigma$ yields
\begin{equation*}
	\begin{split}
		\mathbb E(T\circ\phi_\sigma)^*&=\mathbb E (T^*\circ \phi_\sigma)=\int T^*\circ \phi_\sigma d\p^n\\
		&=\int T^* d(\p^n \circ \phi_\sigma^{-1})=\int T^*d\p^n=\mathbb E T^*=\mathbb E ^*T<\infty.
	\end{split}
\end{equation*}
So, $\mathbb E^*(T\circ\phi_\sigma)= \mathbb E(T\circ\phi_\sigma)^*$ holds and therefore $\mathbb E^*(T\circ\phi_\sigma)=\mathbb E^*T$.
\end{proof}

\begin{question}
Let $T:\Omega\to\mathbb R$ be an arbitrary (not necessarily measurable) map defined on a probability space. Let $T^*$ denote the outer measurable envelope of $T$.  Let $g:\mathbb R\to\mathbb R$ be a nondecreasing function that is left-continuous on an interval $(a,b)$ containing the ranges of both $T$ and $T^*$. Prove that
\begin{equation*}
g(T)^* = g(T^*).
\end{equation*}
If instead $g$ is nondecreasing and right-continuous, prove that
\begin{equation*}
g(T)_* = g(T_*).
\end{equation*}
Show furthermore that the continuity assumptions on $g$ can be replaced by the condition that $g$ is one-to-one, measurable, and has a measurable inverse.
\end{question}

  \begin{proof}
We first prove the statement for a nondecreasing and left-continuous function $g$. By definition of the outer measurable envelope, $T^*$ is measurable and satisfies $T^*\ge T$ almost surely. Since $g$ is nondecreasing, this implies
\begin{equation*}
g(T^*) \ge g(T)
\end{equation*}
almost surely. Because $g(T^*)$ is measurable and dominates $g(T)$, it follows immediately from the definition of outer envelope that
\begin{equation*}
 g(T^*)\geq g(T)^*.
\end{equation*}

To prove the reverse inequality, let $U$ be an arbitrary measurable random variable such that
\begin{equation*}
g(T^*) \ge U \ge g(T)
\end{equation*}
almost surely. Define the generalized inverse of $g$ by
\begin{equation*}
g^{-1}(u) := \sup\{x\in\mathbb R : g(x)\le u\}.
\end{equation*}
Since $g$ is nondecreasing and left-continuous, we claim that this inverse satisfies
\begin{equation*}
g(x)\le u \quad\text{if and only if}\quad x\le g^{-1}(u).
\end{equation*}
In fact, if $g(x)\le u$ the definition of $g^{-1}$ yields $x\leq g^{-1}(u)$. On the contrary, by the definition of $g^{-1}(u)$, there exists a sequence $(x_n)_{n\geq 1}$ satisfying 
\begin{equation*}
	g(x_n)\leq u\tag{$*$}
\end{equation*}
 and $x_n\uparrow g^{-1}(u),n\to\infty$. Obviously, $g^{-1}$ is nondecreasing and has left limits. Letting $n\to\infty$ in $(*)$, we will obtain $g(g^{-1}(u))\leq u$. So, if $x\le g^{-1}(u)$, take $g(\cdot)$ on both sides to get $g(x)\leq g(g^{-1}(u))\leq u$.
 
 
Applying this pointwise with $u=U(\omega)$ yields
\begin{equation*}
g^{-1}(U) \ge T
\end{equation*}
almost surely. Because $g^{-1}(U)$ is measurable and dominates $T$, the definition of $T^*$ implies
\begin{equation*}
g^{-1}(U) \ge T^*
\end{equation*}
almost surely. Applying $g$ to both sides and using monotonicity again gives
\begin{equation*}
U \ge g(T^*)
\end{equation*}
almost surely. Since $U$ was an arbitrary measurable random variable satisfying $U\ge g(T)$, this shows that $g(T^*)$ is the smallest measurable majorant of $g(T)$. Hence
\begin{equation*}
g(T)^* = g(T^*).
\end{equation*}

For the case where $g$ is nonincreasing and right-continuous, apply the preceding result to the function $x\mapsto -g(-x)$, which is nondecreasing and left-continuous. Substitute $-T$ for $x$ and apply what we have proved to obtain
\begin{equation*}
(-g(-(-T)))^* = -g(-(-T)^*). \tag{$*$}
\end{equation*}
Since $X_*=-(-X)^*$ for any arbitraty mapping, we have 
\begin{equation*}
	(-g(-(-T)))^*=(-g(T))^*=-g(T)_*
\end{equation*}
and
\begin{equation*}
	-g(-(-T)^*)=-g(T_*).
\end{equation*}
Substituting the above two identities into ($*$), we will have $g(T)_* = g(T_*)$.

Finally, if $g$ is one-to-one, measurable, and has a measurable inverse, the same argument applies with $g^{-1}$ in place of the generalized inverse defined above, and the continuity assumptions are no longer needed. This completes the proof.
\end{proof}

\begin{question}
Let $(\mathcal X,\mathcal A,\p)$ be a probability space and let $\p^n$ be the product measure on $(\mathcal X^n,\mathcal A^n)$. 
Let $h:\mathcal X^n\to\mathbb R$ be \emph{permutation-symmetric}, i.e., for every permutation $\sigma$ of $\{1,\dots,n\}$,
$h(x_1,\dots,x_n)=h(x_{\sigma(1)},\dots,x_{\sigma(n)})$.
Show that there exists a version $h^*$ of the measurable cover of $h$ (with respect to $\p^n$) that is also permutation-symmetric.
\end{question}

\begin{proof}
Let $\tilde h$ be any minimal measurable cover of $h$ with respect to $\p^n$. For each permutation $\sigma$, define $\tilde h_\sigma:=\tilde h\circ\sigma$, i.e.,
\begin{equation*}
\tilde h_\sigma(x_1,\dots,x_n)=\tilde h(x_{\sigma(1)},\dots,x_{\sigma(n)}).
\end{equation*}
Then $\tilde h_\sigma$ is $\mathcal A^n$-measurable, and since $h$ is permutation-symmetric and $\tilde h\ge h$,
\begin{equation*}
\tilde h_\sigma=\tilde h\circ\sigma\ge h\circ\sigma= h,
\end{equation*}
so each $\tilde h_\sigma$ is a measurable majorant of $h$.

Define
\begin{equation*}
h^*(x):=\min_{\sigma}\tilde h_\sigma(x),
\end{equation*}
where the minimum is over the (finite) set of all $n!$ permutations. As a minimum of finitely many measurable functions, $h^*$ is measurable; and since each $\tilde h_\sigma\ge h$, we have $h^*\ge h$.

Fix a permutation $\tau$. Using that the permutation group is closed under composition and that $\sigma\mapsto\sigma\circ\tau$ is a bijection of the set of permutations, we obtain
\begin{equation*}
h^*\circ\tau
=\min_{\sigma} (\tilde h\circ \tau)\circ \sigma 
=\min_{\sigma} \tilde h\circ (\tau\circ \sigma)
=\min_{\rho}\tilde h\circ\rho
=h^*,
\end{equation*}
so $h^*$ is permutation-symmetric.

Since $h^*\geq h$ and $h^*$ is measurable, we have $h^*\geq \tilde h,\p^n\text{-a.s.}$. However, by definition $h^*=\min_\sigma \tilde h_\sigma\leq \tilde h_{\mathrm{id}}=\tilde h$. So, $h^*=\tilde h,\p^n\text{-a.s.}$.
\end{proof}

\begin{question}
Let $(\mathcal X^\infty,\mathcal A^\infty,\p^\infty)$ be the canonical product probability space, and let $X_i(\omega)=\omega_i$ be the coordinate projections.  For $n \ge 1$, let $\mathscr S_n$ be the $\sigma$-field generated by all measurable functions $h:\mathcal X^\infty \to \mathbb R$ that are permutation-symmetric in their first $n$ arguments.  Suppose $Z_1,\dots,Z_m$ are integrable random variables satisfying
\begin{equation*}
Z_j = Z_i \circ \sigma \quad \text{a.s.}
\end{equation*}
whenever $\sigma$ is a permutation sending $i$ to $j$, where $\sigma$ denotes the corresponding coordinate permutation map on $\mathcal X^\infty$ satisfying 
\begin{equation*}
	\sigma (\underbrace{ x_1,x_2,\cdots,x_i}_{i\mathrm{\ elements}},\cdots)=(\underbrace{ \cdots,x_j}_{i\mathrm{\ elements}},\cdots).
\end{equation*}
Prove that for every $i,j$,
\begin{equation*}
\mathbb{E}(Z_i \mid \mathscr S_n)
=
\mathbb{E}(Z_j \mid \mathscr S_n)
\quad \text{a.s.}
\end{equation*}
\end{question}

\begin{proof}

Let $\sigma$ be a permutation of the first $n$ coordinates.  
Define
\begin{equation*}
\mathcal C
=
\{A \in \mathcal A^\infty : \sigma^{-1}(A) = A \text{ for all } \sigma \in \mathfrak S_n \}.
\end{equation*}
It is straightforward to verify that $\mathcal C$ is a $\sigma$-algebra. If $h$ is permutation-symmetric in its first $n$ coordinates, then $h \circ \sigma = h$, hence for any Borel set $B \subset \mathbb R$,
\begin{equation*}
\sigma^{-1}(\{h \in B\})
=
\{h \circ \sigma \in B\}
=
\{h \in B\}.
\end{equation*}
Thus all generating sets of $\mathscr S_n$ lie in $\mathcal C$, and therefore
\begin{equation*}
\mathscr S_n \subset \mathcal C.
\end{equation*}
Hence every $A \in \mathscr S_n$ satisfies
\begin{equation*}
\sigma^{-1}(A) = A.
\end{equation*}

Fix $i$ and $j$, and let $\sigma$ be a permutation sending $i$ to $j$.  
By assumption,
\begin{equation*}
Z_j = Z_i \circ \sigma \quad \text{a.s.}
\end{equation*}
Let $A \in \mathscr S_n$. Then
\begin{equation*}
\begin{split}
    \int_A Z_j \, d\p^\infty
&=
\int_A Z_i \circ \sigma \, d\p^\infty=\int_{\sigma^{-1}( A)} Z_i \circ \sigma \, d\p^\infty\\
&=\int_{A} Z_i \, d(\p^\infty\circ \sigma^{-1}) =\int_{A} Z_i \, d\p^\infty,
\end{split}
\end{equation*} 
where the last equality follows from the fact that $\p^\infty\circ \sigma^{-1}=\p^\infty$.

By definition of conditional expectation, for every $A \in \mathscr S_n$,
\begin{equation*}
\int_A \mathbb{E}(Z_i \mid \mathscr S_n) \, d\p^\infty
=
\int_A Z_i \, d\p^\infty,
=\int_A Z_j \, d\p^\infty=
\int_A \mathbb{E}(Z_j \mid \mathscr S_n) \, d\p^\infty.
\end{equation*}
Since $\mathbb{E}(Z_i \mid \mathscr S_n)\in\mathscr{S}_n$ and $\mathbb{E}(Z_j \mid \mathscr S_n)\in\mathscr{S}_n$,
\begin{equation*}
\mathbb{E}(Z_i \mid \mathscr S_n)
=
\mathbb{E}(Z_j \mid \mathscr S_n)
\quad \text{a.s.}
\end{equation*}
This completes the proof.
\end{proof}

\begin{question}[Why symmetrization cannot be reversed in general; centered classes]
Let $X_1,\dots,X_n$ be i.i.d.\ real-valued random variables and let $\varepsilon_1,\dots,\varepsilon_n$
be i.i.d.\ Rademacher random variables (i.e.\ $\mathbb P(\varepsilon_i=\pm1)=1/2$), independent of the
$X_i$'s.

\begin{enumerate}
\item Show that there is \emph{no} universal constant $K<\infty$ such that for \emph{all} $n$ and \emph{all}
i.i.d.\ $X_1,\dots,X_n$,
\begin{equation*}
\mathbb E\Big|\sum_{i=1}^n \varepsilon_i X_i\Big|
\ \le\
K\,\mathbb E\Big|\sum_{i=1}^n (X_i-\mathbb E X_i)\Big|.
\end{equation*}

\item Let $\mathcal F$ be a class of measurable functions on $(\mathcal X,\mathcal A)$ and let
$\|\mathbb P_n-P\|_{\mathcal F}:=\sup_{f\in\mathcal F}|(\mathbb P_n-P)f|$ and
$\|\mathbb P_n^\circ\|_{\mathcal F}:=\sup_{f\in\mathcal F}|\mathbb P_n^\circ f|$ where
\begin{equation*}
\mathbb P_n^\circ f:=\frac1n\sum_{i=1}^n \varepsilon_i f(X_i).
\end{equation*}
Explain why the standard symmetrization inequality (for nondecreasing convex $\Phi$)
\begin{equation*}
\mathbb E^*\Phi\big(\|\mathbb P_n-P\|_{\mathcal F}\big)
\ \le\
\mathbb E^*\Phi\big(2\|\mathbb P_n^\circ\|_{\mathcal F}\big)
\end{equation*}
cannot, in general, be reversed up to universal constants. Conclude that any reverse inequality (up to
constants) must impose a \emph{centering} condition on $\mathcal F$, i.e., $Pf=0$ for any $f\in\mathcal F$.
\end{enumerate}
\end{question}

\begin{proof}
\textbf{1.}
It suffices to consider $n=1$. The desired inequality would then read
\begin{equation*}
\mathbb E|\varepsilon_1 X_1|
\ \le\
K\,\mathbb E|X_1-\mathbb E X_1|.
\end{equation*}
Since $|\varepsilon_1 X_1|=|X_1|$, this becomes
\begin{equation*}
\mathbb E|X_1|
\ \le\
K\,\mathbb E|X_1-\mathbb E X_1|
\qquad\text{for all laws of }X_1.
\end{equation*}
Take $X_1\sim N(\mu,1)$. Then $X_1-\mathbb E X_1=X_1-\mu\sim N(0,1)$, hence
\begin{equation*}
\begin{split}
	\mathbb E|X_1-\mathbb E X_1|&=\mathbb E|Z|=
\int_{-\infty}^{\infty} |z| \frac{1}{\sqrt{2\pi}} e^{-z^2/2}\,dz =
2 \int_{0}^{\infty} z \frac{1}{\sqrt{2\pi}} e^{-z^2/2}\,dz,\\
&=-\sqrt{\frac{2}{\pi}}\int_0^\infty e^{-z^2/2}d(z^2/2)=\sqrt{\frac{2}{\pi}},\quad Z\sim N(0,1).
\end{split}
\end{equation*}
On the other hand, writing $X_1=\mu+Z$ and using $|\mu+Z|\ge|\mu|-|Z|$,
\begin{equation*}
\mathbb E|X_1|
=\mathbb E|\mu+Z|
\ \ge\
|\mu|-\mathbb E|Z|
=
|\mu|-\sqrt{\frac{2}{\pi}}.
\end{equation*}
Therefore, the inequality would imply, for all $\mu\in\mathbb R$,
\begin{equation*}
|\mu|-\sqrt{\frac{2}{\pi}}
\ \le\
K\sqrt{\frac{2}{\pi}}.
\end{equation*}
Letting $|\mu|\to\infty$ yields a contradiction. Hence, no such universal $K$ can exist.

\medskip
\noindent
\textbf{2.}
A function class $\mathcal F$ is said to be \emph{$P$-centered} if
\begin{equation*}
Pf=0\qquad\text{for all }f\in\mathcal F.
\end{equation*}
Equivalently, $\mathcal F\subset L_0^2(P):=\{f\in L^2(P):Pf=0\}$. In this case $(\mathbb P_n-P)f=\mathbb P_n f$
for all $f\in\mathcal F$, so the ``mean/shift'' obstruction disappears, and reversing symmetrization up to
constants becomes plausible

Suppose, toward a contradiction, that there exist universal constants $C<\infty$ and (say) a nondecreasing
convex $\Phi(t)=t$ such that for \emph{all} function classes $\mathcal F$,
\begin{equation*}
\mathbb E^*\big\| \mathbb P_n^\circ \big\|_{\mathcal F}
\ \le\
C\,\mathbb E^*\big\| \mathbb P_n-P \big\|_{\mathcal F}.
\tag{$*$}
\end{equation*}
Take $n=1$ and take a singleton class $\mathcal F=\{f\}$ with $f(X_1)$ integrable. Then
\begin{equation*}
\|\mathbb P_1^\circ\|_{\mathcal F}
=
|\varepsilon_1 f(X_1)|
=
|f(X_1)|,
\qquad
\|\mathbb P_1-P\|_{\mathcal F}
=
|f(X_1)-Pf|.
\end{equation*}
Plugging this into $(*)$ gives
\begin{equation*}
\mathbb E|f(X_1)|
\ \le\
C\,\mathbb E|f(X_1)-Pf|.
\end{equation*}
Since $f(X_1)$ can be \emph{any} integrable random variable (by choosing $f$ appropriately on the support of
$P$), this contradicts Part 1. Hence a reverse inequality of symmetrization type cannot hold uniformly over
all classes $\mathcal F$.

The obstruction is the possible presence of a \emph{mean/shift component}. Indeed, write
$f(X)=Pf+(f(X)-Pf)$. The symmetrized quantity $\varepsilon f(X)$ retains the magnitude of the mean term
$Pf$ (it becomes $\varepsilon Pf$), while $(\mathbb P_n-P)f$ removes the mean and only measures
fluctuations. If $|Pf|$ can be arbitrarily large relative to the fluctuation scale, no universal reverse bound
is possible.
\end{proof} 

\begin{question}[Comparison between $\|\cdot\|_{2,1}$ and $L^p$ norms]
Let $\xi$ be a real-valued random variable and define, for $p \ge 1$,
\begin{equation*}
\|\xi\|_{p,1}
:=
\int_{0}^{\infty}
\mathbb P(|\xi| > x)^{1/p}\,dx.
\end{equation*}
Show that for every $r>2$,
\begin{equation*}
\|\xi\|_2 \le \|\xi\|_{2,1}
\le
\frac{r}{r-2}\|\xi\|_r.
\end{equation*}
\end{question}

\begin{proof}
Let $X := |\xi| \ge 0$. Then
\begin{equation*}
\|\xi\|_{2,1}
=
\int_{0}^{\infty}
\mathbb P(X>x)^{1/2}\,dx.
\end{equation*}

Set $a(x) := \mathbb P(X>x)^{1/2}.$
Then $a$ is nonincreasing and $\mathbb P(X>x) = a(x)^2.$
Using the layer-cake representation,
\begin{equation*}
\mathbb E[X^2]
=
\int_0^\infty 2x\,\mathbb P(X>x)\,dx
=
2\int_0^\infty x\,a(x)^2\,dx.
\end{equation*}

On the other hand,
\begin{equation*}
\Big(\int_0^\infty a(x)\,dx\Big)^2
=
\int_0^\infty\int_0^\infty a(x)a(y)\,dx\,dy.
\end{equation*}
By symmetry of the integrand,
\begin{equation*}
\int_0^\infty\int_0^\infty a(x)a(y)\,dx\,dy
=
2\int_0^\infty\int_0^x a(x)a(y)\,dy\,dx.
\end{equation*}
Since $a$ is nonincreasing, for $0\le y\le x$ we have $a(y) \ge a(x),$
and hence $a(x)a(y) \ge a(x)^2.$
Therefore,
\begin{equation*}
\Big(\int_0^\infty a(x)\,dx\Big)^2
\ge
2\int_0^\infty\int_0^x a(x)^2\,dy\,dx
=
2\int_0^\infty x\,a(x)^2\,dx
=
\mathbb E[X^2].
\end{equation*}
Taking square roots gives
\begin{equation*}
\|\xi\|_2 = (\mathbb E[X^2])^{1/2}
\le
\int_0^\infty a(x)\,dx
=
\|\xi\|_{2,1}.
\end{equation*}



Let $r>2$ and denote $A:=\|X\|_r=\|\xi\|_r$.
By Markov's inequality,
\begin{equation*}
\mathbb P(X>x)
\le
\frac{\mathbb E[X^r]}{x^r}
=
\Big(\frac{A}{x}\Big)^r.
\end{equation*}
Hence
\begin{equation*}
\mathbb P(X>x)^{1/2}
\le
\min\Big\{1,\Big(\frac{A}{x}\Big)^{r/2}\Big\}.
\end{equation*}
Thus
\begin{equation*}
\|\xi\|_{2,1}
\le
\int_0^\infty
\min\Big\{1,\Big(\frac{A}{x}\Big)^{r/2}\Big\}\,dx.
\end{equation*}
Split the integral at $x=A$:
\begin{equation*}
\int_0^\infty
\min\Big\{1,\Big(\frac{A}{x}\Big)^{r/2}\Big\}\,dx
=
\int_0^A 1\,dx
+
\int_A^\infty \Big(\frac{A}{x}\Big)^{r/2}\,dx.
\end{equation*}
The first term equals $A$. Since $r/2>1$,
\begin{equation*}
\int_A^\infty x^{-r/2}\,dx
=
\frac{A^{1-r/2}}{(r/2)-1}.
\end{equation*}
Hence
\begin{equation*}
\int_A^\infty \Big(\frac{A}{x}\Big)^{r/2}\,dx
=
\frac{A}{(r/2)-1}.
\end{equation*}
Therefore,
\begin{equation*}
\|\xi\|_{2,1}
\le
A + \frac{A}{(r/2)-1}
=
A\Big(1+\frac{1}{(r/2)-1}\Big)
=
\frac{r}{r-2}A.
\end{equation*}
This proves
\begin{equation*}
\|\xi\|_{2,1}
\le
\frac{r}{r-2}\|\xi\|_r.
\end{equation*}
\end{proof}

\begin{question}[Relative open sets as countable unions of compacts; tightness implies inner regularity]
Let $(S,d)$ be a metric space.

\begin{enumerate}
\item (\textbf{Relative open sets inside a compact set})  
Let $K\subset S$ be compact and let $G\subset S$ be open. Show that the relatively open set
\begin{equation*}
U := K\cap G
\end{equation*}
can be written as an increasing countable union of compact sets:
\begin{equation*}
U=\bigcup_{m=1}^\infty K_m,
\qquad
K_m\subset K_{m+1}\subset U,
\qquad
K_m \ \text{compact}.
\end{equation*}

\item (\textbf{Tightness $\Rightarrow$ inner regularity})  
Let $\mu$ be a Borel probability measure on $S$ that is tight, i.e.\ for every $\varepsilon>0$ there
exists compact $K\subset S$ with
\begin{equation*}
\mu(K)\ge 1-\varepsilon.
\end{equation*}
Show that $\mu$ is inner regular on open sets: for every open $G\subset S$,
\begin{equation*}
\mu(G)=\sup\{\mu(C): C\subset G,\ C \text{ compact}\}.
\end{equation*}
Conclude that for every open $G$ and every $\varepsilon>0$ there exists a compact $C\subset G$ such that
\begin{equation*}
\mu(G\setminus C)<\varepsilon.
\end{equation*}
\end{enumerate}
\end{question}

\begin{proof}
\textbf{1.}
Fix a compact $K\subset S$ and an open set $G\subset S$. Let $G^c$ denote the complement of $G$ and define
\begin{equation*}
\phi(x):=d(x,G^c)=\inf_{y\in G^c} d(x,y),\qquad x\in S.
\end{equation*}
Since $G^c$ is closed, the distance function $\phi$ is continuous. Moreover,
\begin{equation*}
x\in G \quad \Longleftrightarrow \quad \phi(x)>0.
\end{equation*}
For each integer $m\ge 1$, define
\begin{equation*}
K_m:=\{x\in K:\phi(x)\ge 1/m\}
=
K\cap \{x\in S: d(x,G^c)\ge 1/m\}.
\end{equation*}
Because $\phi$ is continuous, the set $\{x:\phi(x)\ge 1/m\}$ is closed; hence $K_m$, being a closed subset
of the compact set $K$, is compact. Also $K_m\subset K_{m+1}$ since $1/m\ge 1/(m+1)$.

Next, if $x\in K_m$, then $\phi(x)\ge 1/m>0$, so $x\in G$; thus $K_m\subset K\cap G=U$ and therefore
\begin{equation*}
\bigcup_{m=1}^\infty K_m \subset U.
\end{equation*}
Conversely, if $x\in U=K\cap G$, then $\phi(x)>0$, so choose $m$ with $1/m\le \phi(x)$; then $x\in K_m$.
Hence
\begin{equation*}
U \subset \bigcup_{m=1}^\infty K_m.
\end{equation*}
Combining the two inclusions yields
\begin{equation*}
U=\bigcup_{m=1}^\infty K_m,
\end{equation*}
with $K_m$ compact and increasing.

\medskip
\noindent
\textbf{2.}
Let $G\subset S$ be open and fix $\varepsilon>0$. By tightness, choose a compact $K\subset S$ such that
\begin{equation*}
\mu(K)\ge 1-\varepsilon,
\qquad \text{equivalently } \mu(K^c)\le \varepsilon.
\end{equation*}
Then
\begin{equation*}
\mu(G)=\mu(G\cap K)+\mu(G\cap K^c)\le \mu(G\cap K)+\mu(K^c)\le \mu(G\cap K)+\varepsilon,
\end{equation*}
so
\begin{equation*}
\mu(G\cap K)\ge \mu(G)-\varepsilon.
\end{equation*}
By Part 1, the relatively open set $U:=G\cap K$ can be written as an increasing union of compact sets
$U=\bigcup_{m=1}^\infty K_m$ with $K_m\subset U\subset G$. By continuity from below,
\begin{equation*}
\mu(U)=\mu\Big(\bigcup_{m=1}^\infty K_m\Big)=\lim_{m\to\infty}\mu(K_m).
\end{equation*}
Hence there exists $m$ such that
\begin{equation*}
\mu(K_m)\ge \mu(U)-\varepsilon.
\end{equation*}
Let $C:=K_m\subset G$ (compact). Then
\begin{equation*}
\mu(C)\ge \mu(U)-\varepsilon \ge (\mu(G)-\varepsilon)-\varepsilon=\mu(G)-2\varepsilon,
\end{equation*}
which implies
\begin{equation*}
\mu(G\setminus C)=\mu(G)-\mu(C)\le 2\varepsilon.
\end{equation*}
Since $\varepsilon>0$ is arbitrary, this proves the inner regularity identity
\begin{equation*}
\mu(G)=\sup\{\mu(C): C\subset G,\ C \text{ compact}\}.
\end{equation*}
In particular, given any $\eta>0$, applying the above argument with $\varepsilon=\eta/2$ yields a compact
$C\subset G$ such that $\mu(G\setminus C)<\eta$.
\end{proof}

\begin{question}[Union of compact sets inside shrinking neighborhoods of a compact set]
Let $(S,d)$ be a metric space and let $K_0\subset S$ be compact.
Let $\{K_n\}_{n\ge1}$ be a sequence of \emph{compact} subsets of $S$.
Assume that there exists a sequence $\delta_n\downarrow 0$ such that
\begin{equation*}
K_n \subset K_0^{\delta_n}
:=\{x\in S : d(x,K_0)<\delta_n\},
\qquad n\ge1.
\end{equation*}
Show that
\begin{equation*}
K:=\bigcup_{n=1}^\infty K_n
\end{equation*}
is relatively compact. Equivalently,
\begin{equation*}
\overline{K}=\overline{\bigcup_{n=1}^\infty K_n}
\end{equation*}
is compact.
\end{question}

\begin{proof}
Let $\{x_j\}$ be an arbitrary sequence in $K=\bigcup_{n\ge1}K_n$.
For each $j$ choose an index $n(j)$ such that $x_j\in K_{n(j)}$.

We distinguish two cases.

\medskip
\textbf{Case 1: Some index occurs infinitely often.}
Suppose there exists $m\ge1$ such that $n(j)=m$ for infinitely many $j$.
Then there is a subsequence $\{x_{j_k}\}$ with $x_{j_k}\in K_m$ for all $k$.
Since $K_m$ is compact, $\{x_{j_k}\}$ admits a further convergent subsequence.

\medskip
\textbf{Case 2: No index occurs infinitely often.}
Then $n(j)\to\infty$ along a subsequence. Re-label so that $n(j)\to\infty$.
Since $x_j\in K_{n(j)}\subset K_0^{\delta_{n(j)}}$, there exists $y_j\in K_0$ such that
\begin{equation*}
d(x_j,y_j)\le \delta_{n(j)}.
\end{equation*}
Because $K_0$ is compact, the sequence $\{y_j\}\subset K_0$ has a convergent subsequence
$y_{j_k}\to y\in K_0$. Along this subsequence we still have $n(j_k)\to\infty$, hence
$\delta_{n(j_k)}\to0$. By the triangle inequality,
\begin{equation*}
d(x_{j_k},y)\le d(x_{j_k},y_{j_k})+d(y_{j_k},y)
\le \delta_{n(j_k)}+d(y_{j_k},y)\to0,
\end{equation*}
so $x_{j_k}\to y$.

\medskip
In either case, every sequence in $K$ has a convergent subsequence in $S$.
Therefore $\overline{K}$ is sequentially compact, hence compact.
\end{proof}
  
  \begin{question}[Uniform tightness vs. asymptotic tightness]
Let $(S,d)$ be a metric space and let $X_n$ be Borel measurable maps into $S$.

\begin{enumerate}
\item Recall that $\{X_n\}$ is called \emph{uniformly tight} if for every $\varepsilon>0$
there exists a compact set $K\subset S$ such that
\begin{equation*}
\mathbb P(X_n\in K)\ge 1-\varepsilon \qquad \text{for every } n .
\end{equation*}

It is called \emph{asymptotically tight} if for every $\varepsilon>0$ there exists
a compact set $K_0$ such that
\begin{equation*}
\varliminf_{n\to\infty}\mathbb P(X_n\in K_0^\delta)\ge 1-\varepsilon
\qquad \text{for every } \delta>0 .
\end{equation*}

Show that a sequence $\{X_n\}$ is uniformly tight if and only if it is asymptotically
tight and each $X_n$ is tight.

\item Conclude that for Borel measurable random elements in a Polish space,
uniform tightness and asymptotic tightness coincide.
\end{enumerate}
\end{question}

\begin{proof}

\textbf{1.}
Suppose $\{X_n\}$ is uniformly tight. Then for every $\varepsilon>0$
there exists a compact $K$ such that
\begin{equation*}
\mathbb P(X_n\in K)\ge 1-\varepsilon \qquad \text{for all } n .
\end{equation*}
Since $K\subset K^\delta$ for every $\delta>0$,
\begin{equation*}
\mathbb P(X_n\in K^\delta)\ge \mathbb P(X_n\in K)\ge 1-\varepsilon .
\end{equation*}
Hence
\begin{equation*}
\varliminf_{n\to\infty}\mathbb P(X_n\in K^\delta)\ge 1-\varepsilon ,
\end{equation*}
so the sequence is asymptotically tight.

\medskip
\noindent
\textbf{2.}
Fix $\varepsilon>0$. By asymptotic tightness, there exists a compact set $K_0$ such that
\begin{equation*}
\varliminf_{n\to\infty}\mathbb P(X_n\in K_0^\delta)\ge 1-\varepsilon
\qquad \text{for every } \delta>0 .
\end{equation*}
Based on this, for each integer $m\ge1$ choose $n_m$ increasing so that
\begin{equation*}
\mathbb P(X_n\in K_0^{1/m})\ge 1-2\varepsilon
\qquad \text{for all } n\ge n_m .
\end{equation*}

Now consider the finitely many indices $n_m\le n\le n_{m+1}$.  
Since $X_n$ is tight, by the previous and previous question, there exists a compact set $K_n\subset K_0^{1/m}$ such that
\begin{equation*}
\mathbb P\bigl(X_n\in K_0^{1/m}\setminus K_n\bigr)<\varepsilon .
\end{equation*}

Define
\begin{equation*}
K:=\bigcup_{n=0}^{\infty}K_n .
\end{equation*}
By the previous question, $K$ is pre-compact. So, for every $n$,
\begin{equation*}
\begin{split}
	\mathbb P(X_n\in \overline{ K})&\geq\mathbb P(X_n\in K)
\ge \p(X_n\in K_n)\\
&=\mathbb P(X_n\in K_0^{1/m})-\mathbb P(X_n\in K_0^{1/m}\setminus K_n)
>
(1-2\varepsilon)-\varepsilon
=
1-3\varepsilon .
\end{split}
\end{equation*}
Since $\varepsilon$ is arbitrary, this proves uniform tightness.

\medskip
\noindent
\textbf{3.}
If $S$ is Polish, every Borel probability measure on $S$ is tight.
Hence the condition that each $X_n$ is tight holds automatically.
Therefore, in Polish spaces
\begin{equation*}
\text{uniform tightness}
\quad\Longleftrightarrow\quad
\text{asymptotic tightness}.
\end{equation*}
\end{proof}

\begin{question}[Uniform tightness and weak compactness of a set of Borel measures]

Let $\mathscr L$ be a collection of Borel probability measures on a complete metric space $(\mathbb D,d)$. 
Call $\mathscr L$ \emph{uniformly tight} if for every $\varepsilon>0$ there exists a compact set $K\subset\mathbb D$ such that
\begin{equation*}
L(K)\ge 1-\varepsilon \qquad \text{for every } L\in\mathscr L.
\end{equation*}

Show that the following statements are equivalent:

\begin{enumerate}
\item[(i)] $\mathscr L$ is totally bounded for the bounded Lipschitz metric $d_{BL}$;
\item[(ii)] $\mathscr L$ is uniformly tight;
\item[(iii)] $\mathscr L$ is relatively compact.
\end{enumerate}
\end{question}

\begin{proof}

\textbf{(ii) $\Rightarrow$ (iii).}  
If $\mathscr L$ is uniformly tight, then by Prokhorov's theorem every sequence in $\mathscr L$ has a weakly convergent subsequence. Hence $\mathscr L$ is relatively compact.

\medskip
\noindent
\textbf{(iii) $\Rightarrow$ (i).}  
In any semimetric space, relative compactness implies total boundedness. Therefore if $\mathscr L$ is relatively compact under $d_{BL}$, it must be totally bounded.

\medskip
\noindent
\textbf{(i) $\Rightarrow$ (ii).}
Let $\varepsilon>0$ be fixed.  
For any Borel set $B\subset\mathbb D$, define
\begin{equation*}
f(x) = \bigl(1-\varepsilon^{-1}d(x,B)\bigr)^+ .
\end{equation*}

This function is bounded by $1$ and has Lipschitz constant $\varepsilon^{-1}$. Moreover,
\begin{equation*}
1_B(x)\le f(x)\le 1_{B^\varepsilon}(x),
\end{equation*}
where
\begin{equation*}
B^\varepsilon=\{x:d(x,B)<\varepsilon\}.
\end{equation*}
Hence for any probability measures $L$ and $L'$,
\begin{equation*}
L(B)
\le \int f\,dL
\le \int f\,dL' + d_{BL}(L,L')
\le L'(B^\varepsilon) + d_{BL}(L,L').
\end{equation*}
Rearranging gives
\begin{equation*}
L'(B^\varepsilon) \ge L(B) - \varepsilon^{-1} d_{BL}(L,L').
\end{equation*}

\medskip

Because $\mathscr L$ is totally bounded under $d_{BL}$, for any $\delta>0$ there exists a finite set
\begin{equation*}
\mathscr L_0 = \{L_1,\dots,L_k\}\subset\mathscr L
\end{equation*}
such that
\begin{equation*}
\mathscr L \subset \bigcup_{i=1}^k B_{d_{BL}}(L_i,\delta^2).
\end{equation*}

For each $L_i$, since it is a probability measure on Polish space, there exists a compact set $K_i$ satisfying
\begin{equation*}
L_i(K_i)\ge 1-\delta.
\end{equation*}

Let
\begin{equation*}
F=\bigcup_{i=1}^k K_i .
\end{equation*}
The set $F$ is compact as a finite union of compact sets.  
Because compact sets in metric spaces are totally bounded, there exists a finite set $G$ such that $F \subset G^\delta.$
Then for any $L\in\mathscr L$ there exists $L_0\in\mathscr L_0$ with $d_{BL}(L,L_0)<\delta^2 .$

Applying the previous inequality with $B=F$ yields
\begin{equation*}
L(G^{2\delta}) \ge L_0(G^\delta) - \delta^{-1}d_{BL}(L,L_0).
\end{equation*}
Since $L_0(F)\ge1-\delta$ and $F\subset G^\delta$, we obtain
\begin{equation*}
L(G^{2\delta})\ge (1-\delta)-\delta^{-1} \delta^{2} =1-2\delta .
\end{equation*}

\medskip

Define
\begin{equation*}
A=\bigcap_{m=1}^{\infty} G^{2\varepsilon 2^{-m}},
\qquad
K=\overline{A}.
\end{equation*}
Because $\varepsilon2^{-m}\to0$, for any $\eta>0$ there exists $m$ such that
\begin{equation*}
2\varepsilon2^{-m}<\eta .
\end{equation*}
Since
\begin{equation*}
A \subset G^{2\varepsilon2^{-m}}
= \bigcup_{x\in G} B(x,2\varepsilon2^{-m}),
\end{equation*}
and $G$ is finite, $A$ can be covered by finitely many $\eta$-balls. Hence $A$ is totally bounded.  
The closure of a totally bounded set remains totally bounded, so $K=\overline A$ is also totally bounded.

Because $\mathbb D$ is complete and $K$ is closed. Therefore
\begin{equation*}
\text{close + totally bounded} \;\Rightarrow\; \text{compact},
\end{equation*}
and $K$ is compact.

Finally, by construction
\begin{equation*}
L(K)\ge 1-\varepsilon
\qquad \text{for every } L\in\mathscr L.
\end{equation*}

Thus $\mathscr L$ is uniformly tight.

\end{proof}


\begin{question}[Poisson splitting / multinomial thinning]
Let $(\Omega,\mathcal A,\mathbb P)$ be a probability space. Fix $k\ge2$, 
$\lambda>0$, and probabilities $p_1,\dots,p_k\in[0,1]$ satisfying
$\sum_{i=1}^k p_i=1$.
Let
\begin{equation*}
N \sim \mathrm{Poisson}(\lambda).
\end{equation*}
Let $\{Y_j\}_{j\ge1}$ be an i.i.d. sequence of random variables taking
values in $\{1,\dots,k\}$ such that
\begin{equation*}
\mathbb P(Y_j=i)=p_i,\qquad i=1,\dots,k,
\end{equation*}
and assume that $\{Y_j\}_{j\ge1}$ is independent of $N$.

Define, for each $i=1,\dots,k$,
\begin{equation*}
N_i:=\sum_{j=1}^{N} \mathbf 1\{Y_j=i\}.
\end{equation*}
Then the random variables $N_1,\dots,N_k$ are independent and
\begin{equation*}
N_i \sim \mathrm{Poisson}(\lambda p_i),\qquad i=1,\dots,k.
\end{equation*}
\end{question}

\begin{proof}
We use two ways to prove this.

\noindent
{\bfseries Proof via probability generating functions.} 
First condition on the value of $N$.
If $N=m$, then
\begin{equation*}
N_i=\sum_{j=1}^{m}\mathbf 1\{Y_j=i\}.
\end{equation*}
Hence
\begin{equation*}
\prod_{i=1}^k s_i^{N_i}
=
\prod_{i=1}^k s_i^{\sum_{j=1}^m 1\{Y_j=i\}}
=
\prod_{j=1}^m \prod_{i=1}^k s_i^{1\{Y_j=i\}}.
\end{equation*}
For each $j$, exactly one indicator equals $1$, hence
\begin{equation*}
\prod_{i=1}^k s_i^{1\{Y_j=i\}} = s_{Y_j}.
\end{equation*}
Therefore
\begin{equation*}
\prod_{i=1}^k s_i^{N_i}=\prod_{j=1}^m s_{Y_j}.
\end{equation*}
Taking the conditional expectation and using the i.i.d property of $Y_1,\dots,Y_m$ gives
\begin{equation*}
\begin{split}
	\mathbb E\!\left[\prod_{i=1}^k s_i^{N_i}\mid N=m\right]
&=
\mathbb E\!\left[\prod_{j=1}^m s_{Y_j}\right]=
\prod_{j=1}^m \mathbb E[s_{Y_j}]\\
&=\mathbb \e^m[s_{Y_j}]=\(\prod_{j=1}^m
 p_i s_i\)^m.
\end{split}
\end{equation*}

Now remove the conditioning:
\begin{equation*}
\mathbb E\!\left[\prod_{i=1}^k s_i^{N_i}\right]
=
\mathbb E\!\left[\(\prod_{j=1}^m
 p_i s_i\)^N\right].
\end{equation*}
Using the Poisson generating function,
\begin{equation*}
\mathbb E[a^N]=\exp\{\lambda(a-1)\},
\end{equation*}
we obtain
\begin{equation*}
\begin{split}
	\mathbb E\!\left[\prod_{i=1}^k s_i^{N_i}\right]
&=
\exp\!\left(\lambda\left(\sum_{i=1}^k p_i s_i-1\right)\right)\\
&=
\exp\!\left(\sum_{i=1}^k \lambda p_i (s_i-1)\right)
=
\prod_{i=1}^k \exp\!\left(\lambda p_i (s_i-1)\right)
\end{split}
.
\end{equation*}
This is the joint probability generating function of independent
Poisson random variables with parameters $\lambda p_1,\dots,\lambda p_k$.
Therefore
\begin{equation*}
N_i \sim \mathrm{Poisson}(\lambda p_i)
\end{equation*}
and the variables are independent.

\medskip
\noindent
{\bfseries Proof via direct computation of the joint pmf.}
Let $n_1,\dots,n_k\in\mathbb N_0$ and set
\begin{equation*}
m=\sum_{i=1}^k n_i.
\end{equation*}
Using the law of total probability,
\begin{equation*}
\mathbb P(N_1=n_1,\dots,N_k=n_k)
=
\sum_{r=0}^\infty
\mathbb P(N=r)
\mathbb P(N_1=n_1,\dots,N_k=n_k \mid N=r).
\end{equation*}
The conditional probability is zero unless $r=m$, because
\begin{equation*}
\sum_{i=1}^k N_i = N.
\end{equation*}
Thus
\begin{equation*}
\mathbb P(N_1=n_1,\dots,N_k=n_k)
=
\mathbb P(N=m)
\mathbb P(N_1=n_1,\dots,N_k=n_k \mid N=m).
\end{equation*}

Conditional on $N=m$, the vector $(N_1,\dots,N_k)$ follows a
multinomial distribution, hence
\begin{equation*}
\mathbb P(N_1=n_1,\dots,N_k=n_k \mid N=m)
=
\frac{m!}{n_1!\cdots n_k!}
\prod_{i=1}^k p_i^{n_i}.
\end{equation*}
Since $N\sim\mathrm{Poisson}(\lambda)$,
\begin{equation*}
\mathbb P(N=m)
=
e^{-\lambda}\frac{\lambda^m}{m!}.
\end{equation*}
Multiplying yields
\begin{equation*}
\mathbb P(N_1=n_1,\dots,N_k=n_k)
=
e^{-\lambda}
\frac{\lambda^m}{m!}
\frac{m!}{n_1!\cdots n_k!}
\prod_{i=1}^k p_i^{n_i}.
\end{equation*}
Since $\sum_{i=1}^k p_i=1$, we have
\begin{equation*}
e^{-\lambda}
=
\prod_{i=1}^k e^{-\lambda p_i}.
\end{equation*}
Hence
\begin{equation*}
\mathbb P(N_1=n_1,\dots,N_k=n_k)
=
\prod_{i=1}^k
\left(
e^{-\lambda p_i}
\frac{(\lambda p_i)^{n_i}}{n_i!}
\right).
\end{equation*}

This factorization shows that $N_1,\dots,N_k$ are independent
Poisson random variables with parameters $\lambda p_i$.
\end{proof}

\begin{question}[Poisson CLT, and an a.s.\ law via a Poisson process]
\leavevmode
\begin{enumerate}
\item Let $N_n\sim \mathrm{Poisson}(n)$. Prove the (standard) Poisson central limit theorem
\begin{equation*}
\frac{N_n-n}{\sqrt{n}}\ \Rightarrow\ \mathcal N(0,1).
\end{equation*}
Deduce that
\begin{equation*}
N_n-n=O_p(\sqrt n),
\qquad
\frac{N_n-n}{n}=O_p(n^{-1/2})=o_p(1),
\qquad
\frac{N_n}{n}=1+O_p(n^{-1/2}).
\end{equation*}

\item Let $\{N(t):t\ge0\}$ be a rate-$1$ Poisson process, and define $N_n:=N(n)$. Prove the almost sure law
\begin{equation*}
\frac{N_n}{n}\ \to\ 1
\qquad\text{a.s.}
\end{equation*}
\end{enumerate}
\end{question}

\begin{proof}
\textbf{Part 1: Poisson CLT via mgf.}
Let $N_n\sim \mathrm{Poisson}(n)$ and set
\begin{equation*}
Z_n:=\frac{N_n-n}{\sqrt n}.
\end{equation*}
For any $t\in\mathbb R$, the mgf of $N_n$ is finite and equals
\begin{equation*}
\mathbb E\big[e^{tN_n}\big]=\exp\big(n(e^{t}-1)\big).
\end{equation*}
Hence
\begin{equation*}
\begin{split}
	\mathbb E\big[e^{tZ_n}\big]
&=
\mathbb E\Big[\exp\Big(\frac{t}{\sqrt n}(N_n-n)\Big)\Big]=
\exp\Big(-t\sqrt n\Big)\,
\mathbb E\Big[\exp\Big(\frac{t}{\sqrt n}N_n\Big)\Big]\\
&
=
\exp\Big(n(e^{t/\sqrt n}-1)-t\sqrt n\Big)=
\exp\Big\{n\Big(e^{t/\sqrt n}-1-\frac{t}{\sqrt n}\Big)\Big\}.
\end{split}
\end{equation*}
Using the Taylor expansion $e^x=1+x+\frac{x^2}{2}+o(x^2)$ as $x\to0$ with $x=t/\sqrt n$, we have
\begin{equation*}
e^{t/\sqrt n}-1-\frac{t}{\sqrt n}
=
\frac{t^2}{2n}+o\Big(\frac{1}{n}\Big),
\end{equation*}
and thus
\begin{equation*}
n\Big(e^{t/\sqrt n}-1-\frac{t}{\sqrt n}\Big)
=
\frac{t^2}{2}+o(1).
\end{equation*}
Therefore
\begin{equation*}
\lim_{n\to\infty}\mathbb E\big[e^{tZ_n}\big]
=
\exp\Big(\frac{t^2}{2}\Big),
\qquad t\in\mathbb R.
\end{equation*}
Since $\exp(t^2/2)$ is the mgf of $\mathcal N(0,1)$, the mgf convergence theorem yields
\begin{equation*}
Z_n\Rightarrow \mathcal N(0,1),
\end{equation*}
i.e.\ the Poisson CLT.

\medskip
\textbf{Stochastic order consequences.}
From $Z_n\Rightarrow \mathcal N(0,1)$ it follows that $\{Z_n\}$ is tight, hence
\begin{equation*}
\frac{N_n-n}{\sqrt n}=O_p(1),
\qquad\text{equivalently}\qquad
N_n-n=O_p(\sqrt n).
\end{equation*}
Dividing by $n$ gives
\begin{equation*}
\frac{N_n-n}{n}
=
\frac{1}{\sqrt n}\cdot \frac{N_n-n}{\sqrt n}
=
O_p(n^{-1/2})=o_p(1),
\end{equation*}
and in particular $\frac{N_n-n}{n}=o_p(1)$. Finally,
\begin{equation*}
\frac{N_n}{n}
=
1+\frac{N_n-n}{n}
=
1+O_p(n^{-1/2}).
\end{equation*}

\medskip
\textbf{Part 2: a.s.\ convergence using a Poisson process.}
Let $\{N(t):t\ge0\}$ be a rate-$1$ Poisson process, and let $\{T_j\}_{j\ge1}$ be its i.i.d.\ interarrival times, so $T_j\sim \mathrm{Exp}(1)$ and the arrival epochs are
\begin{equation*}
S_k:=\sum_{j=1}^k T_j,
\qquad k\ge1.
\end{equation*}
Then
\begin{equation*}
N(t)=\max\{k\ge0:S_k\le t\},
\end{equation*}
and for every $t\ge0$ we have the fundamental inequalities
\begin{equation*}
S_{N(t)}\le t<S_{N(t)+1}.
\end{equation*}
By the strong law of large numbers applied to $\{T_j\}$,
\begin{equation*}
\frac{S_k}{k}\to \mathbb E[T_1]=1
\qquad\text{a.s.}
\end{equation*}
Moreover, $N(t)\to\infty$ as $t\to\infty$ a.s. (since $S_k\to\infty$ a.s.), hence substituting $k=N(t)$ and $k=N(t)+1$ into the SLLN yields
\begin{equation*}
\frac{S_{N(t)}}{N(t)}\to 1
\qquad\text{and}\qquad
\frac{S_{N(t)+1}}{N(t)+1}\to 1
\qquad\text{a.s.}
\end{equation*}
Divide the inequalities $S_{N(t)}\le t<S_{N(t)+1}$ by $N(t)$ to get
\begin{equation*}
\frac{S_{N(t)}}{N(t)}
\le
\frac{t}{N(t)}
<
\frac{S_{N(t)+1}}{N(t)}
=
\frac{S_{N(t)+1}}{N(t)+1}\cdot\frac{N(t)+1}{N(t)}.
\end{equation*}
The rightmost product converges to $1\cdot 1=1$ a.s., and the leftmost term converges to $1$ a.s., hence by squeezing
\begin{equation*}
\frac{N(t)}{t}\to 1
\qquad\text{a.s.}.
\end{equation*}
Evaluating at integer times $t=n$ gives
\begin{equation*}
\frac{N_n}{n}=\frac{N(n)}{n}\to 1
\qquad\text{a.s.}
\end{equation*}
This completes the proof.
\end{proof}

\begin{question}[Rademacher randomization of a symmetric distribution]
Let $\xi$ be a real-valued random variable whose distribution is symmetric about $0$, i.e.
\begin{equation*}
\xi \overset{d}{=} -\xi .
\end{equation*}
Let $\epsilon$ be a Rademacher random variable such that
\begin{equation*}
\mathbb P(\epsilon=1)=\mathbb P(\epsilon=-1)=\tfrac12,
\end{equation*}
and assume that $\epsilon$ is independent of $\xi$. Show that
\begin{equation*}
\epsilon |\xi| \overset{d}{=} \xi .
\end{equation*}
\end{question}

\begin{proof}
Let $f:\mathbb R\to\mathbb R$ be bounded and measurable. By the tower property
and the independence of $\epsilon$ and $\xi$,
\begin{equation*}
\mathbb E[f(\epsilon|\xi|)]
=
\mathbb E\!\left[\mathbb E\big(f(\epsilon|\xi|)\mid \xi\big)\right].
\end{equation*}
Conditional on $\xi$, the quantity $|\xi|$ is fixed while $\epsilon$
remains a Rademacher random variable. Hence
\begin{equation*}
\mathbb E\big(f(\epsilon|\xi|)\mid \xi\big)
=
\frac12 f(|\xi|)+\frac12 f(-|\xi|).
\end{equation*}
Taking expectation yields
\begin{equation*}
\mathbb E[f(\epsilon|\xi|)]
=
\frac12\,\mathbb E[f(|\xi|)]
+
\f rac12\,\mathbb E[f(-|\xi|)].
\end{equation*}

Define the function
\begin{equation*}
g(x):=\frac12\big(f(x)+f(-x)\big), \qquad x\in\mathbb R.
\end{equation*}
Then $g$ is bounded, measurable, and even, i.e.\ $g(x)=g(-x)$.
Observe that for every $x\in\mathbb R$,
\begin{equation*}
g(|x|)=g(x).
\end{equation*}
Consequently,
\begin{equation*}
\mathbb E[f(\epsilon|\xi|)]
=
\mathbb E[g(|\xi|)]
=
\mathbb E[g(\xi)].
\end{equation*}

Finally, using the symmetry of $\xi$, namely $\xi\overset d=-\xi$, we obtain
\begin{equation*}
\mathbb E[g(\xi)]
=
\frac12\mathbb E[f(\xi)]
+
\frac12\mathbb E[f(-\xi)]
=
\mathbb E[f(\xi)].
\end{equation*}
Therefore
\begin{equation*}
\mathbb E[f(\epsilon|\xi|)]=\mathbb E[f(\xi)].
\end{equation*}
Since this holds for every bounded measurable $f$, it follows that
\begin{equation*}
\epsilon|\xi|\overset{d}{=}\xi .
\end{equation*}
\end{proof}

  \begin{thebibliography}{99}
   \bibitem{ref1} Dachun Yang, Wen Yuan, {\itshape Selected lectures on functional analysis} [M]. Beijing Normal University Press, 2016.
  \end{thebibliography}



\end{document} 